% 第6章：动态场景与4D
\chapter{动态场景与4D}
\label{ch:06}

\paragraph*{本章阅读锚点}
本章的锚点是时间。原始 3DGS 表达静态场景；动态方法要么让高斯随时间移动，要么引入变形场、4D 表示或轨迹约束，关键是区分几何变化、外观变化和相机运动。


前五章建立了静态场景的 3DGS 技术栈：高斯表示（第~\ref{ch:02}~章）、渲染（第~\ref{ch:03}~章）、几何（第~\ref{ch:04}~章）、材质（第~\ref{ch:05}~章）。但现实世界不是静止的——人行走、车行驶、旗帜飘动、水在流动。如何用 3DGS 建模和渲染这些动态场景，是本章的核心问题。

3DGS 的静态表示假设场景不变：每个高斯的位置、颜色、协方差保持固定。动态场景打破了这一假设——高斯需要随时间变化。最直觉的方案是为每个时间步建立独立的 3DGS 模型，但这样参数量爆炸（$N$ 帧 $\times$ 每帧百万级高斯），而且帧间不连续。核心挑战在于：如何在保持渲染质量的同时，用紧凑的表示建模高斯的时变属性。

本章围绕动态场景建模展开。第~\ref{sec:06_1}~节介绍变形场（deformation field）与 4D 高斯表示——两条建模时变属性的主要技术路线。第~\ref{sec:06_2}~节讨论 4D 扩展与特殊动态重建。第~\ref{sec:06_3}~节讨论单目动态场景重建。第~\ref{sec:06_4}~节介绍运动建模与跟踪。第~\ref{sec:06_5}~节面向自由视角视频应用。第~\ref{sec:06_6}~节处理动态物体与场景分解。第~\ref{sec:06_7}~节覆盖其他动态场景方法。第~\ref{sec:06_8}~节介绍物理仿真与动力学。

\begin{figure}[htbp]
\centering
\resizebox{\textwidth}{!}{% 图6.1: 动态场景 3DGS 技术路线图
\begin{tikzpicture}[
  node distance=0.8cm and 1.8cm,
  block/.style={rectangle, rounded corners=3pt, draw=techgray, fill=techbg,
    text width=2.8cm, minimum height=1.0cm, align=center, font=\small},
  route/.style={rectangle, rounded corners=3pt, draw=#1, fill=techbg,
    text width=2.5cm, minimum height=0.8cm, align=center, font=\small\bfseries},
  arr/.style={-{Stealth[length=5pt]}, thick, techgray},
  label/.style={font=\scriptsize\color{techlightgray}, align=center},
]

% Central node: Static 3DGS
\node[block, text width=3.5cm, minimum height=1.2cm, draw=techblue, fill=techbluebg, font=\normalsize\bfseries] (gs)
  {静态 3DGS\\（位置+颜色固定）};

% Two main routes
\node[route=techblue, above left=1.8cm and 2.0cm of gs] (r1)
  {路线 A\\变形场};
\node[route=techgreen, above right=1.8cm and 2.0cm of gs] (r2)
  {路线 B\\原生 4D};

% Arrows to routes
\draw[arr, techblue] (gs.north west) -- (r1.south east)
  node[label, midway, above left] {静态 GS + 变形 MLP};
\draw[arr, techgreen] (gs.north east) -- (r2.south west)
  node[label, midway, above right] {时空联合高斯};

% Deformation details
\node[block, left=1.5cm of r1, text width=2.2cm, draw=techblue, font=\scriptsize] (d1)
  {逐高斯位移\\+ 旋转 + 缩放};
\draw[arr, techblue] (d1.east) -- (r1.west);

% 4D details
\node[block, right=1.5cm of r2, text width=2.2cm, draw=techgreen, font=\scriptsize] (d2)
  {HexPlane / K-plane\\张量分解};
\draw[arr, techgreen] (r2.east) -- (d2.west);

% Key encodings
\node[block, below=0.3cm of d1, text width=2.2cm, draw=techblue, font=\scriptsize] (e1)
  {嵌入 / 哈希\\B样条 / ODE};
\node[block, below=0.3cm of d2, text width=2.2cm, draw=techgreen, font=\scriptsize] (e2)
  {4D Scaffold\\锚点驱动};

% Section labels
\node[route=techorange, below left=2.5cm and 1.5cm of gs] (s2)
  {§6.2 单目重建};
\node[route=techpurple, left=2.5cm of gs] (s3)
  {§6.3 运动跟踪};
\node[route=techgray, below right=2.5cm and 1.5cm of gs] (s4)
  {§6.4 视频流};
\node[route=techorange, below=2.5cm of gs] (s5)
  {§6.5 场景分解};

% Arrows from center to sections
\draw[arr, techorange] (gs.south west) -- (s2.north east)
  node[label, midway, below left] {深度先验/正则化};
\draw[arr, techpurple] (gs.west) -- (s3.east)
  node[label, midway, above] {轨迹/速度场};
\draw[arr, techgray] (gs.south east) -- (s4.north west)
  node[label, midway, below right] {压缩/流式传输};
\draw[arr, techorange] (gs.south) -- (s5.north)
  node[label, midway, right] {静态-动态分离};

% Paper counts
\node[label, below=0.15cm of r1] {82篇};
\node[label, below=0.15cm of r2] {原生 4D};
\node[label, below=0.15cm of s2] {35篇};
\node[label, below=0.15cm of s3] {18篇};
\node[label, below=0.15cm of s4] {49篇};
\node[label, below=0.15cm of s5] {24篇};

\end{tikzpicture}
}
\caption{动态场景 3DGS 技术路线图。变形场（路线~A）和原生 4D 高斯（路线~B）分别建模高斯的时变属性。}
\label{fig:ch06_routes}
\end{figure}

%% ============================================================

%=====================================================================
\section{变形场与4D高斯表示}
\label{sec:06_1}

本节介绍动态场景建模的两条主要技术路线：变形场路线与原生 4D 高斯路线。内容按以下子主题组织：问题背景、变形场基础方法、原生 4D 高斯表示、高阶 4D 编码方案、变形场的正则化与优化策略、物理约束变形、事件相机辅助动态重建、前馈与大规模 4D 重建、动态表面重建、运动建模与操纵，以及分析与概率方法。

\subsection{问题背景}

3DGS 在静态场景中表现出色，但现实世界充满动态变化。将 3DGS 扩展到动态场景的核心问题在于：如何建模高斯基元随时间的属性变化？具体而言，动态场景中每个高斯的位置 $\boldsymbol{\mu}(t)$、旋转 $\mathbf{R}(t)$、缩放 $\mathbf{S}(t)$、协方差（covariance）$\mathbf{c}(t)$ 都可能随时间 $t$ 变化。如果为每个时间帧独立存储所有高斯参数，存储开销将线性增长（$N$ 帧 $\times$ 每帧百万级高斯 $\times$ 59 个参数），在实际应用中不可接受。

两条主要技术路线应运而生：

\textbf{变形场路线}（Deformation Field）保持一组规范空间（canonical space）中的静态高斯不变，由一个时间条件的变形网络预测每帧高斯的位移、旋转和缩放变化。这种方式的优势在于：静态部分的参数被所有帧共享，变形网络仅编码差异信息，整体参数量远小于逐帧独立建模。变形场的概念源自 NeRF 领域的 D-NeRF 和 Nerfies，但 3DGS 的显式表示特性使变形场的应用面临独特挑战——变形需要作用于离散的高斯基元，而非连续的体积密度场。

\textbf{原生 4D 路线}（Native 4D）不依赖规范空间，而是直接将时间作为第四维度引入高斯表示。每个 4D 高斯（4D Gaussian）定义在时空域 $(x, y, z, t)$ 中，其属性自然包含时间维度。4D 表示省去了变形场的中间环节，但需要设计高效的 4D 数据结构来避免参数爆炸。为此，研究者提出了 HexPlane（六面体平面）张量分解、K-平面（K-Planes）、4D 哈希编码等技术来压缩 4D 表示。

本节按以下子主题组织：首先介绍变形场基础方法，然后讨论原生 4D 高斯表示，接着深入高阶 4D 编码方案，再介绍变形场的正则化与优化策略，随后覆盖物理约束变形、事件相机辅助、前馈与大规模 4D 重建，最后讨论动态表面重建与运动操纵。

\subsection{变形场基础理论}

变形场路线的核心思想是：在规范空间维护一组静态 3D 高斯，由一个时间条件的变形网络预测每帧的属性偏移量。可以把它想象成一块``时空橡皮''——将所有时刻的 3D 高斯都``压''到一个共享的规范空间中优化，渲染时再``拉伸''回各时刻的位置。具体而言，给定时间 $t$，变形网络 $D_\theta$ 将规范空间中的高斯属性 $(\boldsymbol{\mu}, \mathbf{R}, \mathbf{S})$ 映射为时间 $t$ 处的变形属性：

\begin{equation}
(\Delta\boldsymbol{\mu}, \Delta\mathbf{R}, \Delta\mathbf{S}) = D_\theta(\boldsymbol{\mu}, t)
\end{equation}

变形后的高斯参数为 $\boldsymbol{\mu}(t) = \boldsymbol{\mu} + \Delta\boldsymbol{\mu}$，旋转和缩放类似。这种设计使静态场景参数被所有帧共享，变形网络仅编码帧间差异。

\paragraph{逐高斯嵌入与粗-细变形分解}

本方法修改了变形场的输入机制——用逐高斯可学习嵌入（per-Gaussian learnable embedding）替代 4D 坐标输入，使变形网络能区分"位置相近但身份不同"的高斯基元。其余变形场框架（规范空间 + 时间条件变形 + 加性偏移）保持不变。

\textbf{问题：坐标输入无法区分高斯身份。}
Baseline 做法（如 Deformable 3DGS）将变形建模为坐标函数 $D(\mathbf{x}, t)$，以高斯基元的 3D 坐标和时间组成 4D 输入，预测参数偏移量——也就是"问空间位置怎么动"。但 3DGS 的场景表示本质上是一个混合体（mixture）：场景由多个具有独立辐射场的各向异性高斯组成，不同高斯基元可能属于不同物体，运动模式截然不同。例如手和杯子紧挨着，一个快速运动，一个静止。坐标函数 $D(\mathbf{x}, t)$ 只看位置，无法区分"位置相近但身份不同"的高斯基元——即使将特征网格分辨率翻倍，改善也微乎其微。这一观察指向一个核心结论：变形的定义对象应该是高斯基元本身，而非它所处的空间坐标。

\textbf{方案：从坐标到嵌入。}
基于上述分析，该方法放弃坐标输入，改为为每个高斯基元 $i$ 分配一个可学习的嵌入向量 $z_g \in \mathbb{R}^{32}$——这个嵌入相当于高斯基元的"身份证"，携带了该基元属于哪个物体、应如何运动的身份信息。同时为每帧学习一个时间嵌入 $z_t \in \mathbb{R}^{256}$（从 1D 特征网格 $Z \in \mathbb{R}^{N \times 256}$ 中插值得到），编码"当前帧的运动状态"。嵌入 $z_g$ 回答"我是谁"，时间嵌入 $z_t$ 回答"现在是什么时刻"，两者组合即可唯一确定该高斯基元在该帧的变形。变形函数定义为：

\begin{equation}
(\Delta\boldsymbol{\mu}, \Delta\mathbf{q}, \Delta\mathbf{s}, \Delta o, \Delta\mathbf{Y}) = F_\theta(z_g, z_t)
\end{equation}

其中 $\mathbf{q}$ 为旋转四元数、$\mathbf{s}$ 为缩放向量、$o$ 为不透明度、$\mathbf{Y}$ 为球谐系数（spherical harmonics, SH，用于建模视角相关颜色）。$F_\theta$ 由一个浅层 MLP 后接多个参数专用 MLP 头实现，每个头独立输出一种属性偏移。也就是说，变形网络不再问"这个空间位置怎么动"，而是问"这个高斯基元在这个时刻怎么动"——从空间坐标驱动变为身份驱动。帧 $t$ 处的高斯参数由规范参数加上变形偏移确定。训练时联合优化嵌入 $z_g$、变形函数 $F_\theta$ 和规范高斯参数，使用 L1 损失加周期性 D-SSIM 损失作为渲染损失。

\textbf{粗-细变形分解：分离慢运动与快运动。}
为进一步提升变形精度，方法引入粗-细（coarse-fine）变形分解。场景中不同部分可能同时存在大尺度慢速运动（如人体缓慢从左向右走）和小尺度快速运动（如手快速搅拌锅），用单一变形函数同时捕捉两者很困难。因此方法将变形拆分为两个具有相同架构但输入不同时间嵌入的函数，最终变形为两者之和。具体而言，时间嵌入网格 $Z \in \mathbb{R}^{N \times 256}$ 为 $N$ 帧各存储一个 256 维特征向量。细变形直接使用该网格的原始特征，通过线性插值提取帧 $t$ 处的时间嵌入 $z_t^{(\text{fine})} = \text{Interp}(Z, t)$，保留高频细节，负责快速或精细运动。粗变形则先将 $Z$ 沿时间轴线性下采样 5 倍（例如 120 帧降为 24 帧），去除高频成分，再同样通过线性插值得到低频时间嵌入 $z_t^{(\text{coarse})} = \text{Interp}(\text{Downsample}(Z, 5), t)$，负责大尺度或慢速运动。最终变形为两者之和：$F_\theta(z_g, z_t^{(\text{coarse})}) + F_\theta(z_g, z_t^{(\text{fine})})$。这种设计的好处是：粗变形先建立整体运动趋势，细变形在此基础上补充细节，两者各司其职。

\textbf{局部平滑正则化：嵌入空间的一致性约束。}
物理世界中，相邻的高斯基元往往属于同一物体，应该具有相似的变形（例如手臂上的相邻高斯一起运动）。方法据此提出局部平滑正则化，约束空间上邻近的高斯具有相似的嵌入向量，从而间接诱导相似的变形：

\begin{equation}
\mathcal{L}^{\text{emb-reg}} = \frac{1}{k|S|} \sum_{i \in S} \sum_{j \in \text{KNN}_k(i)} w_{i,j} \|z_{g,i} - z_{g,j}\|_2, \quad w_{i,j} = \exp(-\sigma_w \|\boldsymbol{\mu}_j - \boldsymbol{\mu}_i\|_2^2)
\end{equation}

其中 $k=20$，$\sigma_w=2000$，仅在致密化时计算 KNN 以控制开销。权重 $w_{i,j}$ 越大（距离越近），正则化约束越强。该正则化不直接约束物理属性（如刚性或旋转），而是通过嵌入空间的平滑间接诱导相似的变形——这是一种更灵活的约束：它允许相邻高斯在运动上有微小差异（保持灵活性），但在大方向上保持一致。这种间接约束在保持灵活性的同时，改善了动态区域的纹理和细节捕捉。

\paragraph{稀疏体素 U-Net 提取几何感知变形特征}

前述 Per-Gaussian Embedding 方法通过局部平滑正则化间接约束相邻高斯的嵌入相似性，但正则化仅作用于嵌入空间，变形网络本身仍然逐点独立预测——网络看不到高斯之间的空间邻域关系，也就无法直接利用局部几何结构来约束变形的几何一致性。

Geometry-aware Deformable GS~\cite{2404_06270}（2024.04）从另一个角度解决同一问题：与其在嵌入空间间接施加平滑约束，不如让变形网络直接"看到"局部几何——即从高斯的空间分布中提取几何感知特征，使相邻高斯的变形自然趋于一致。

\textbf{动机。} 逐点独立预测变形无法利用局部几何结构来约束变形的一致性——相邻高斯的变形可能出现"撕裂"，破坏物体表面的连续性。

\textbf{方案：从可学习嵌入到几何感知特征。}
与 Per-Gaussian Embedding 为每个高斯分配独立嵌入 $z_g$（需额外正则化才能保证相邻相似）不同，本方法直接从高斯的空间分布推导特征——特征天然编码了几何邻域关系，无需额外正则化。具体而言，方法将规范空间的高斯位置视为 3D 点云 $P \in \mathbb{R}^{N \times 3}$，通过双分支结构提取特征。

\textbf{几何分支}用固定网格尺寸 $s$ 将高分辨率点集体素化为低分辨率稀疏体素 $V = \lfloor P / s \rfloor$（$M$ 个体素，$M \ll N$），再通过稀疏 3D U-Net 聚合局部几何特征 $F_v \in \mathbb{R}^{M \times C}$——同一体素内的多个高斯共享同一几何特征，自然保证了局部一致性。

\textbf{身份分支}用一个 MLP 将高斯 3D 坐标映射为逐点嵌入 $F_p \in \mathbb{R}^{N \times C}$，保持每个高斯的独立分辨率。两者通过体素特征回传与拼接融合为点级特征：

\begin{equation}
F_{\text{fuse}} = \phi([F_p', F_p])
\end{equation}

其中 $F_p'$ 为体素特征回传到对应点的结果（同一体素内的高斯获得相同的 $F_p'$），$[\cdot, \cdot]$ 为拼接操作，$\phi$ 为 2 层 MLP（宽度 64）。这种设计使融合特征同时具备局部几何一致性（来自几何分支）和逐点独立性（来自身份分支），消融实验表明几何分支是性能提升的最大贡献者。

\textbf{旋转表示：从四元数到连续 6D。}
另一个关键差异是旋转参数化。Per-Gaussian Embedding 沿用 3DGS 的四元数表示预测旋转偏移 $\Delta\mathbf{q}$，但四元数空间存在拓扑不连续性（$\mathbf{q}$ 和 $-\mathbf{q}$ 表示同一旋转），网络在学到平滑旋转变形时容易遇到梯度路径折断。本方法将旋转表示替换为连续 6D 表示（2019）~\cite{zhou2019continuity}，通过映射 $f_{V2M}(\mathbf{a}_1, \mathbf{a}_2)$ 将 6D 向量转化为 $SO(3)$ 旋转矩阵：

\begin{equation}
f_{V2M}(\mathbf{a}_1, \mathbf{a}_2) = [\mathbf{b}_1, \mathbf{b}_2, \mathbf{b}_3], \quad \mathbf{b}_i = \text{Normalize}\left(\mathbf{a}_i - \sum_{j=1}^{i-1} \langle \mathbf{a}_i, \mathbf{b}_j \rangle \mathbf{b}_j \right)
\end{equation}

其中 $\mathbf{a}_1, \mathbf{a}_2 \in \mathbb{R}^3$ 为可学习参数，初始化为 $[1,0,0]^\top$ 和 $[0,1,0]^\top$（对应单位旋转）。消融实验表明，将 6D 替换回四元数会导致明显的性能下降，印证了连续表示对学习时序平滑旋转变形的重要性。

\textbf{变形场与训练。}
变形场解码器 $G_\theta$（5 层 MLP，宽度 256，带 skip connection）以融合特征 $F_{\text{fuse}}$、位置编码 $\gamma(x)$（$L=10$）和时间编码 $\gamma(t)$（$L=6$）为输入，输出位置变形 $\Delta x_t$、旋转变形 $\Delta r_t \in \mathbb{R}^6$（6D 格式）和缩放变形 $\Delta s_t$。规范空间到时刻 $t$ 的变换为加性偏移：$x_t = \bar{x} + \Delta x_t$，$s_t = \bar{s} + \Delta s_t$，旋转通过 6D 矩阵乘法 $r_t = f_{V2M}(\Delta r_t) \cdot f_{V2M}(\bar{r})$ 实现。值得注意的是，颜色 $c$ 和不透明度 $o$ 保持时不变——这与 Per-Gaussian Embedding 预测 $\Delta o$ 和 $\Delta\mathbf{Y}$ 的做法不同。其设计动机是让每个高斯持续表示同一物体区域，避免颜色闪烁。训练总损失为光度损失（$L_1$ + D-SSIM）加运动最小化损失 $\mathcal{L}_{\text{motion}} = \|\Delta x_t\|_1$，隐式约束场景尽量保持静态。训练约 2 小时（D-NeRF 数据集平均），前 3000 次迭代仅优化规范空间，之后引入变形场。

前述方法从嵌入设计和几何感知角度改进变形场的输入特征，但变形场本身仍然是确定性的——输入坐标和时间直接映射为变形输出，没有任何去噪能力。当输入视频含噪或相机位姿不精确时，变形误差会被放大。

\paragraph{两阶段去噪变形与解耦时空聚合}

与 Per-Gaussian Embedding 为每个高斯基元引入可学习嵌入 $Y_i$ 以区分身份的做法一脉相承，DN-4DGS（2024.10）~\cite{2410_13607} 继承了这一设计，但聚焦于前者未触及的一个瓶颈：规范空间高斯的坐标本身携带显著噪声——这些坐标从多帧动态图像中联合优化得到，无法保证空间精度。

\textbf{动机。} 在带噪坐标上执行 $k$-最近邻聚合时，邻居关系本身就是错的，噪声被进一步放大而非抑制——需要先"清洗"坐标再做精细变形。DN-4DGS 的核心贡献正是针对这一"噪声级联放大"问题，提出两阶段去噪变形策略与解耦时空聚合机制。

\textbf{噪声抑制策略（NSS）：先修正坐标，再做精细变形。}
方法将变形拆为两阶段串行操作。第一阶段执行标准变形：输入原始规范坐标 $(x, y, z)$ 和时间 $t$ 的邻居帧集合 $T_n$，通过网络 $\Phi$ 输出坐标变形量 $(\Delta x, \Delta y, \Delta z) = \Phi(x, y, z, T_n)$，将其加回原坐标得到修正坐标 $(x', y', z') = (x + \Delta x, y + \Delta y, z + \Delta z)$。此时坐标分布已被"清洗"——第一阶段变形本质上是一次粗粒度的坐标修正。第二阶段以修正坐标为输入，通过 $\Phi'$ 计算最终变形。训练也采用两阶段策略：早期仅监督第一阶段变形，待其收敛、修正坐标足够精确后，切换到第二阶段监督。这种设计避免了在带噪坐标上直接做精细变形导致误差传递。

\textbf{解耦时空聚合模块（DTS）：时间聚合在前，空间聚合在后。}
在引入邻居聚合时，方法面临一个选择：何时执行空间聚合（KNN）和时序聚合？答案取决于坐标噪声。空间聚合依赖坐标精度——KNN 在带噪坐标上会找错邻居，因此空间聚合被安排在第二阶段变形中执行，此时输入坐标已由第一阶段去噪。时序聚合不依赖空间坐标（时间维度与规范空间噪声无关），因此被安全地安排在第一阶段变形中：对每个高斯 $G_t(i)$，提取当前帧特征 $F_t(i)$ 及相邻帧特征 $F_{t-1}(i)$、$F_{t+1}(i)$，沿时间维度 MaxPooling 得到 $F_{\max}(i)$，再与可学习嵌入 $Y_i$ 拼接生成变形特征 $F_t(i)' = [F_{\max}(i), F_t(i), Y_i]$。空间聚合则在第二阶段执行：基于修正坐标 $(x', y', z')$ 计算 $k$-最近邻，聚合邻居特征后与修正坐标拼接，输入第二阶段变形网络 $\Phi'$。这种"时间在前、空间在后"的解耦设计，从机制上杜绝了噪声通过空间聚合被放大的路径——这是对 Per-Gaussian Embedding 框架的实质性增强。

\paragraph{高斯运动预测与时间感知高斯缓存}

回顾上节，E-D3DGS 通过逐高斯嵌入 $z_g$ 和粗-细双时间嵌入 $(z_t^{(\text{coarse})}, z_t^{(\text{fine})})$ 驱动变形——用两次独立的变形前向传播分别捕捉大尺度慢速运动和精细快速运动，最终取两者之和。这一设计虽然有效，却引入了三个在实践中反复出现的瓶颈。Fast and Robust Deformable 3D Gaussian Splatting (FRoGS)（2026.03）~\cite{2603_20857} 系统分析了这三个瓶颈，并逐一提出针对性改进。

\textbf{动机。} 粗-细双变形导致渲染速度减半，稀疏初始化点云导致运动伪影，暗场景中前景阴影高斯引入色偏——这三个工程瓶颈严重限制了变形场方法的实用性。

\textbf{瓶颈一：粗-细双变形导致渲染速度减半。}
E-D3DGS 的粗-细分解要求每帧执行两次完整的前向传播（分别输入低频和高频时间嵌入），渲染速度因此减半。FRoGS 的核心改进是将后期融合（late fusion，粗-细时间嵌入分别输入两次独立变形）替换为早期融合（early fusion）：先在输入端将粗时间嵌入 $t_c$ 和细时间嵌入 $t_f$ 融合为单一嵌入，再只调用一次变形 MLP 计算全部属性偏移。关键在于融合方式的选择——FRoGS 系统比较了四种候选：加法将所有维度等权聚合，难以突出局部变化，削弱对动态区域的建模精度；注意力机制可学习加权分布，但依赖额外网络层，计算开销高；拼接直接堆叠嵌入，需额外参数学习融合关系，且对嵌入中的局部信号不敏感。相比之下，Hadamard 逐元素乘积 $t_{\text{fused}} = t_c \odot t_f$ 通过非线性逐元素交互高效捕获粗-细时间信号之间的局部乘性关系，既无额外参数也无额外计算，被选为最终方案。融合后的时间嵌入与逐高斯嵌入 $e$ 拼接后送入变形 MLP：

\begin{equation}
t_{\text{fused}} = t_c \odot t_f, \quad (\Delta\boldsymbol{\mu}, \Delta\mathbf{s}, \Delta\mathbf{q}, \Delta o, \Delta\mathbf{sh}) = F(\text{concat}(t_{\text{fused}}, e))
\end{equation}

这样，FRoGS 保留了 E-D3DGS 粗-细分解的多尺度时间信息优势，但将两次前向传播压缩为一次，渲染速度近乎翻倍。

\textbf{瓶颈二：稀疏初始化点云导致运动伪影。}
当初始 SfM 点云在某些区域稀疏甚至缺失时，标准致密化策略（分裂已有高斯）无法有效补充几何——因为根本没有足够的高斯可供分裂。此时变形场被迫学习远距离变换，将远处的现有高斯"拉"到空白区域，这种过程极难收敛，容易产生运动伪影和模糊细节。FRoGS 的思路是跳过变形场的补偿负担，直接在规范场（canonical field）中手动补充新高斯——具体分为定位问题区域和注入新锚点两步。定位阶段，FRoGS 用中值深度（而非均值深度）和高误差像素反向投影来确定缺失几何的 3D 坐标。选择中值深度的原因是：均值深度是沿光线上所有半透明高斯的加权平均，容易被浮游体（floater）严重干扰；而中值深度由累积透射率首次越过 0.5 阈值的那个表面型高斯决定，天然抵抗半透明噪声，能可靠地定位主表面。注入阶段，一个自然的想法是学习逆变形映射，将问题坐标逆映射回答范空间。但这一策略在稀疏区域根本不可行：最近邻高斯可能来自遥远的无关物体，拷贝其嵌入后求解逆变形（$p_c + \Delta p = p_d$）高度病态且对嵌入质量极度敏感。FRoGS 因此采用更直接鲁棒的策略：将问题坐标直接作为低偏差锚点注入规范场，后续由联合优化局部微调即可。新锚点的属性初始化为：缩放由最近邻距离的对数变换确定，旋转设为单位四元数，不透明度初始化为 0.1，嵌入从属性变形总量最小的高斯复制——选择变形最小的嵌入提供了稳定的优化起点。

\textbf{瓶颈三：暗场景中前景阴影高斯引入色偏。}
在暗场景实验中（如 N3DV 数据集的 cut roasted beef 场景），E-D3DGS 倾向于利用相机附近的高斯建模阴影。这些前景阴影高斯在测试视角引入显著色偏，导致 PSNR 下降。FRoGS 通过消融实验确认了两个事实：(1) 变形场确实对部分前景阴影高斯预测了不透明度下降；(2) 渲染时过滤掉不透明度低于 0.05 的高斯可显著改善视觉质量。这说明阴影高斯的不透明度本应被抑制，但标准变形的抑制力度不够。基于此观察，FRoGS 引入激进不透明度衰减函数，在标准变形更新的基础上叠加一个乘性惩罚：

\begin{equation}
o_{\text{final}} = (\sigma^{-1}(o + \Delta o) + o) \cdot \sigma(k)
\end{equation}

其中 $\sigma$ 为 Sigmoid 函数，$k$ 为超参数。直觉上，前景阴影高斯的特征是规范不透明度 $o$ 显著大于变形后不透明度——变形场对其预测大幅负变形 $\Delta o \ll 0$。此时 Sigmoid 项 $\sigma(k)$ 处于快速衰减区域，对最终不透明度施加乘性惩罚，使其渲染权重远低于标准变形更新的结果。这一机制在保持正常结构重建不受影响的前提下，有效抑制了阴影浮游体。

\paragraph{时间维度自注意力建模跨帧运动关系}

回顾 Per-Gaussian Embedding-Based Deformation 的设计思路：用逐高斯嵌入 $z_g$ 回答"我是谁"，用时间嵌入 $z_t$ 回答"现在是什么时刻"，两者组合驱动变形。这一策略成功解决了"坐标无法区分高斯身份"的问题，但其时间建模仍然是从单帧出发的——时间嵌入 $z_t$ 仅编码当前帧的运动状态，变形网络在每一帧独立预测偏移量，帧与帧之间没有显式信息交换。当运动模式简单时这不构成问题，但面对剧烈运动、物体突然出现或消失、动态反射表面等复杂场景时，仅凭当前帧的局部信息很难推断正确的运动。因为这类运动的正确解读往往依赖于"之前在哪里"和"之后去了哪里"。

\textbf{问题：局部时间约束难以应对复杂运动。}
更具体地说，现有变形方法（包括 Per-Gaussian Embedding）本质上是从独立时间戳学习运动模式：变形函数 $D(\cdot, t_i)$ 在每个时间戳 $t_i$ 独立预测偏移量，训练信号仅来自当前帧或相邻帧的渲染损失（及光流约束）。这意味着优化梯度 $\partial \mathbf{g}_i / \partial G = \partial D(G, t_i) / \partial G$ 只反映当前时间戳 $t_i$ 的信息——对于时间戳 $t_i$ 处的运动，网络无法利用 $t_j \neq t_i$ 处的观测来辅助理解。这种局部时间视角的局限在复杂运动中尤为突出：例如一个快速旋转的物体，仅看单帧无法判断旋转方向，但结合前后帧的上下文则一目了然。

\textbf{方案：沿时间维度施加自注意力，隐式学习跨时间运动关系。}
TimeFormer（2024.11）~\cite{2411_11941} 的核心设计是在变形场之前插入一个 Transformer 模块，沿时间维度对每个高斯施加自注意力，使其能"看到"多个时间戳处的运动信息。

\textbf{动机。} 现有变形方法本质上从独立时间戳学习运动模式——对于时间戳 $t_i$ 处的运动，网络无法利用 $t_j \neq t_i$ 处的观测来辅助理解，面对剧烈运动、物体突然出现或消失等复杂场景时表现不佳。灵感来源于 BatchFormer（2022）~\cite{hou2022batch}——注意力机制可以在批次维度上学习样本间关系，而时间戳也可以视为一种特殊的"时间批次"。具体而言，从全部时间帧中随机采样 $B$ 个时间戳组成时间批次 $T_s = \{t_i\}_{i=0}^{B-1}$，规范空间的高斯位置 $G \in \mathbb{R}^{N \times 3}$ 被复制 $B$ 份扩展为 $G_c \in \mathbb{R}^{B \times N \times 3}$（每个时间戳一份），采样时间戳同样扩展为 $T' \in \mathbb{R}^{B \times N \times 1}$。两者拼接后施加位置编码 $\gamma$（$L=6$）提取高频信息，构成 Transformer 的输入特征 $F_0 \in \mathbb{R}^{B \times N \times (3 \cdot 2 + 1 \cdot 2)}$。这可以理解为 $N$ 条长度为 $B$ 的序列，每条序列对应一个高斯基元在 $B$ 个时间戳处的特征——Transformer 的自注意力机制沿时间维度（即序列长度方向）聚合信息，使每个时间戳处的特征融合了其他时间戳的运动线索。经过 $M=4$ 层 Transformer 编码器后，最后一层特征由一个小型 MLP 映射为位置偏移量 $O \in \mathbb{R}^{B \times N \times 3}$：

\begin{equation}
O = \text{MLP}(F_{M-1}), \quad G_t = G_c + O
\end{equation}

该偏移被设计为残差修正项，而非直接替代原始位置，从而鼓励渐进式的运动模式学习。这种隐式学习的一个有趣特性是：具有相似运动变化模式的高斯会在优化过程中自动聚合，促进更高效的空间分布。

\textbf{梯度流视角：从局部时间到全局时间序列。}
从梯度流角度可以更清晰地理解 TimeFormer 的作用。设 $D$ 和 $P$ 分别为变形场和 TimeFormer，无 TimeFormer 时变形输出为 $\mathbf{g}_i = D(G, t_i)$，其对规范高斯 $G$ 的梯度仅依赖当前时间戳 $t_i$。引入 TimeFormer 后，数据流变为 $\mathbf{g}_i = D(G + P(G, T_s), t_i)$，对 $G$ 求偏导：

\begin{equation}
\frac{\partial \mathbf{g}_i}{\partial G} = \frac{\partial D}{\partial a}\bigg|_{a=G+P(G,T_s)} \cdot \left(1 + \frac{\partial P(G, T_s)}{\partial G}\right)
\end{equation}

其中新增的梯度项 $\partial P(G, T_s) / \partial G$ 使当前时间戳 $t_i$ 的优化同时受所有采样时间戳 $t_j \in T_s$（$j \neq i$）的影响——换言之，优化不再只看"这一帧怎么动"，而是从全局时间序列视角理解运动模式。这正是 TimeFormer 能处理复杂运动的根本原因。

\textbf{双流共享权重：知识迁移以保持推理速度。}
Transformer 的引入虽然增强了运动建模能力，但也带来了额外的计算开销，会显著降低渲染速度。TimeFormer 采用双流优化解决这一问题：原始分支直接将规范高斯送入变形场，TimeFormer 分支先经过 Transformer 修正位置再送入变形场，两个分支共享变形场权重。训练损失为 $\mathcal{L} = \lambda_c \mathcal{L}_c + \lambda_t \mathcal{L}_t$，其中 $\lambda_c=1 > \lambda_t=0.8$——TimeFormer 分支的权重较小，这是因为其增强的运动建模能力容易导致过拟合，从而降低推理质量。共享权重的效果是将 TimeFormer 学到的跨时间运动知识隐式迁移到原始分支的变形场中，使原始分支在不借助 Transformer 的情况下也能获得更强的运动理解。推理时移除 TimeFormer 分支，仅使用原始分支，从而保持与标准变形场相同的渲染速度。

\paragraph{光流引导运动解耦与异质 3D 控制点}

\textbf{从身份驱动到运动解耦：为什么嵌入还不够。}
Per-Gaussian Embedding 方法通过为每个高斯基元分配可学习嵌入，让变形网络"认识"每个基元——这解决了"谁在动"的身份问题。但嵌入驱动的变形本质上仍是一个纯学习过程：运动信息全部隐式编码在嵌入向量和 MLP 参数中，既没有利用现成的运动观测信号（如光流），也无法显式控制运动的物理分解。H3D-DGS（2024.08）~\cite{2408_13036} 从另一个角度切入：既然光流已经提供了场景运动的 2D 观测，为什么不直接利用这一信号，只让网络学习"光流看不到"的那部分运动？\textbf{动机。} 嵌入驱动的变形是完全隐式学习——既没有利用现成的运动观测信号（如光流），也无法显式控制运动的物理分解，将变形从"完全隐式学习"推向"观测引导 + 隐式补全"的混合范式可以显著降低优化难度。这一思路将变形从"完全隐式学习"推向"观测引导 + 隐式补全"的混合范式。

\textbf{动机：离散场景需要离散运动表示。}
3DGS 的离散高斯结构天然适合局部化运动表示——与 NeRF 的全局隐式场不同，3DGS 仅在有物体存在的区域分配表示容量，避免了为空白空间浪费容量。同理，场景中只有一小部分是动态的，大部分保持静态，因此运动表示也应离散化和局部化。H3D-DGS 据此提出异质 3D 控制点（H3D Control Points）作为运动表示——在场景中以均匀网格模式稀疏采样，每个控制点携带局部 3D 运动信息，周围高斯通过插值获取变形。

\textbf{运动解耦：可观测分量来自光流，不可观测分量来自梯度。}
H3D-DGS 的核心思想是将每个控制点的局部 3D 运动解耦为可观测分量和不可观测分量。直觉上，从相机视角看去，场景的 3D 运动有一部分投影到图像平面上（例如物体从左向右移动），这部分信息可以直接从光流提取；而深度方向的运动（例如物体远离相机）在光流中几乎不可见，必须通过优化学习。为此，方法为每个控制点建立一个射线坐标系（ray coordinate system）：z 轴沿相机到控制点的方向，x-y 平面与图像平面近似平行。在该坐标系下，2D 平面上的平移和绕 z 轴旋转属于可观测分量，可直接从光流推导；深度方向速度和绕 x、y 轴旋转属于不可观测分量，设为可学习参数。

具体地，光流 $(u_i, v_i)$ 首先通过相机内参逆变换 $K^{-1}$ 转换为物理单位（像素到米），再通过旋转矩阵 $R_{cr}$ 从相机坐标系变换到射线坐标系，得到射线坐标系下的 2D 速度 $v^r$。可观测的 2D 平移 $V^r$ 和绕 z 轴旋转 $\omega_z^r$ 由邻域 $N$ 内的光流聚合得到：

\begin{equation}
V^r = \frac{f}{Z^r} \cdot v^r, \quad \omega_z^r = \frac{1}{N} \sum_{x_i \in N} \frac{(x_i^r - x_0^r) \cdot v_i^r}{\|x_i^r - x_0^r\|^2}
\end{equation}

其中 $f$ 为焦距，$Z^r$ 为控制点深度。焦距的作用是将归一化坐标还原到物理空间。不可观测的深度速度 $v_z^r$ 和绕 x、y 轴旋转 $(\omega_x^r, \omega_y^r)$ 则为可学习参数，通过反向传播优化。最终，控制点的完整运动表示为平移 $t^r = (V^r_x, V^r_y, \hat{v}^r_z)$ 和旋转四元数 $q^r = \text{E2Quat}((\hat{\omega}^r_x, \hat{\omega}^r_y, \omega^r_z))$，其中 $\text{E2Quat}$ 将欧拉角转换为四元数（欧拉角便于解耦各旋转分量，四元数则适配底层 3DGS 框架的旋转表示）。这种异质设计将 6 自由度运动中"看得见的"3 自由度交由光流直接提供、"看不见的"3 自由度交由梯度优化，显著降低了优化难度——方法仅需约 100 次迭代即可收敛，单帧处理速度约 2 秒（单 NVIDIA 4070 GPU）。

\textbf{从控制点到高斯：距离加权插值。}
与 Per-Gaussian Embedding 方法中每个高斯基元直接拥有自己的嵌入向量不同，H3D-DGS 的运动信息由稀疏的控制点承载，高斯通过空间插值获取变形。具体地，每个运动高斯从其空间邻域 $N_G$ 中选取 $K=3$ 个最近控制点，以距离的倒数作为权重进行加权插值：

\begin{equation}
(\mathbf{t}, \mathbf{q}) = \sum_{C_i \in N_G} w_i \cdot (\mathbf{t}_i, \mathbf{q}_i), \quad w_i = \frac{1 / \|X_G - X_{C_i}\|}{\sum_{C_j \in N_G} 1 / \|X_G - X_{C_j}\|}
\end{equation}

高斯在时刻 $t$ 的位置和旋转通过累加计算：$X_t = \mathbf{t} + X_{t-1}$，$q_t = \mathbf{q} \cdot q_{t-1}$。这种插值策略的优势是运动场天然具有空间平滑性（相近的高斯由相近的控制点驱动），无需额外的正则化损失——对比 Per-Gaussian Embedding 方法需要显式的局部平滑正则化来约束嵌入空间的一致性。

\textbf{流式处理框架：视频编码启发的逐帧重建。}
方法还构建了一个流式处理框架（streaming framework），将动态重建分解为四个独立模块：（1）\textbf{3D 分割}——通过多视角掩码投票（将每个高斯投影到所有训练视角的图像平面上，统计各类别票数，取多数票确定类别）将高斯分为动态/静态，控制点仅操纵同类别的动态高斯，这一设计还天然支持空间位置相近但类别不同的物体之间的拓扑变换（如手抓取杯子）；（2）\textbf{H3D 控制点生成}——按上述运动解耦策略从光流构建控制点；（3）\textbf{逐物体运动操纵}——按类别独立地对各动态物体施加运动；（4）\textbf{残差补偿}——借鉴视频编码中的关键帧机制，以 Group of Scenes（GOS）为单位组织帧序列：GOS-N 表示一个关键帧后接 $N-1$ 个非关键帧。关键帧处全量优化高斯属性和控制点参数，非关键帧处高斯属性固定、仅更新控制点参数。这种设计在计算效率和时序稳定性之间取得平衡，确保长序列重建中误差不会持续累积。

\paragraph{预算约束下的可微分人口控制}

前述 Per-Gaussian Embedding-Based Deformation 通过逐高斯嵌入实现了身份驱动的变形，Gaussian-Flow 通过解析函数替代了神经网络——这些方法都在追求"更准确的变形"，但共享一个隐含假设：高斯基元可以自由增殖，致密化策略不受约束。Constrained Dynamic Gaussian Splatting (CDGS)（2026.02）~\cite{2602_03538} 转向一个被长期忽视的更高层次问题：如何在高斯基元数量受到严格预算约束的前提下，仍实现高质量的动态场景重建？\textbf{动机。} 标准 3DGS 的密度控制机制会无约束地增殖高斯——模型尺寸膨胀到无法部署在手机、AR 头显等边缘设备上，而训练后刚性剪枝会丢弃已学到的信息导致质量下降。

\textbf{问题：无约束增殖与边缘部署的矛盾。}
标准 3DGS 的自适应密度控制（Adaptive Density Control, ADC）机制会无约束地增殖高斯基元——训练越久、场景越复杂，高斯数量越多，模型尺寸随之膨胀。对于服务器端离线渲染这或许可以接受，但当目标是在手机、AR 头显等边缘设备上实时渲染时，内存和带宽限制使得数百万高斯基元的模型完全无法部署。现有应对方案各有局限：Ex4DGS 通过分解降低冗余，但不提供显式的容量控制——优化结束后高斯数量仍取决于分解策略，无法精确设定上限；Taming-3DGS 为静态场景引入了预算控制，但其策略是训练完成后的刚性剪枝——优化过程中高斯对容量限制始终无知，按"无约束最优"的轨迹收敛，剪枝时被迫丢弃已学到的信息，导致次优的重建质量。这些方案的共同缺陷是：容量约束与优化过程脱节，而非嵌入式协同。

\textbf{方案：约束优化 + 可微分人口控制。}
CDGS 从根本上重新审视问题结构，将动态场景重建形式化为约束优化问题：在固定高斯数量预算 $N_{\text{target}}$ 下最大化渲染质量，即 $\min_G \mathcal{L}_{\text{render}}(G) \;\text{s.t.}\; |G| \leq N_{\text{target}}$。这一形式化的核心洞察是：与其先无约束优化再事后剪枝，不如让优化从一开始就在预算约束下进行——高斯在"知道"容量有限的前提下学习如何最优地分配自身。

困难在于高斯数量是离散且不可微的，无法直接参与梯度优化。CDGS 引入可微分的人口控制器（Differentiable Population Controller）来弥合这一鸿沟：为每个高斯 $g_i$ 分配一个连续激活变量 $c_i \in [0,1]$，通过温度控制的硬 Sigmoid 门控 $c_i = \text{clip}((M_i - 0.5)\tau_c + 0.5,\, 0,\, 1)$ 实现，其中 $M_i$ 为可学习的重要性分数，温度 $\tau_c$ 从 1.0 指数退火至 0.01，逐步将软选择逼近为硬二值选择。在渲染管线中，$c_i$ 直接调制高斯的有效不透明度 $\alpha'_i = c_i \cdot \alpha_i$——当 $c_i \to 0$ 时，高斯变得透明并退出渲染，物理上等效于软删除。这种设计确保渲染损失的梯度能通过 $c_i$ 回传到重要性分数 $M_i$，使感知关键的高斯被保护、冗余高斯被抑制。有效活跃数量通过可微代理估计：$N_p = \sum_i c_i$。预算损失定义为二次惩罚 $\mathcal{L}_{\text{budget}} = (N_p - N_{\text{target}})^2$，驱动 $N_p$ 趋近目标容量。总损失为渲染损失、预算损失和正则化项的加权和：$\mathcal{L} = \mathcal{L}_{\text{render}} + \lambda_b \mathcal{L}_{\text{budget}} + \lambda_r \mathcal{L}_{\text{reg}}$。

\textbf{统一重要性分数：谁该留下？}
在严格预算下，必须精确评估每个高斯基元对最终重建的贡献，才能决定谁被保留、谁被淘汰。此前 Taming-3DGS 的剪枝指标仅关注静态几何属性（如不透明度和屏幕覆盖），无法处理动态场景中特有的时序冗余和运动复杂性。CDGS 提出统一重要性分数 $M_i$，融合几何稳定性、运动显著性和感知影响三个维度：

\begin{equation}
M_i = \mathcal{N}(w_1^T \mathcal{N}([F_{\text{geom/motion}}(g_i)]) + w_2^T \mathcal{N}([F_{\text{perceptual}}(g_i)]))
\end{equation}

其中 $\mathcal{N}(\cdot)$ 为 Min-Max 归一化，将各分量缩放到 $[0,1]$ 以平衡多模态信号的贡献。几何/运动分量 $F_{\text{geom/motion}}$ 捕捉高斯基元的结构与运动必要性，包含五个子指标：位置梯度 $\|\nabla_{\mu_i}\|_2$ 量化重建损失对高斯位置的敏感度，定位结构关键区域；峰值不透明度 $\max_t \alpha_i(t)$ 取时间序列上的最大值（而非均值），确保瞬态结构（如挥手瞬间）不被因低平均可见度而剪枝；逆深度 $d_i^{-1}$ 显式优先前景元素，因为前景占据更大的屏幕面积、需要更高的渲染精度；空间范围（协方差矩阵最大特征值）保护覆盖大面积区域的高斯，防止剪枝后出现几何空洞；运动幅度 $\|T_i\|_2$ 对静态和动态高斯采用分段策略——静态高斯保留用于校正全局平移的基元，动态高斯优先保留具有大幅位移和旋转的基元以保护高频运动细节。感知分量 $F_{\text{perceptual}}$ 从"移除后退化"的角度量化视觉影响，包含三个子指标：光度残差 $I_i$ 直接计算移除高斯 $g_i$ 后各训练视角的像素级 L1 误差，量化渲染退化程度；像素覆盖 $\rho_i$ 累积高斯在各视角的投影面积，足迹越大意味着对最终外观贡献越大；多视角一致性 $\sigma_i^{-1}$ 利用高斯跨视角贡献的方差倒数，优先保留视角一致性好的基元，抑制仅在某些角度可见的噪声或伪影。

\textbf{自适应动静分配：预算的最优分配。}
场景的时空动态性通常极不均匀——大部分区域保持静止或仅做规则刚体运动，显著的非刚性变形仅发生在特定物体的有限时段内。若均匀分配高斯预算，大量容量将浪费在静态区域。CDGS 引入自适应动静分配策略（Adaptive Static-Dynamic Allocation），将场景高斯集合 $G$ 分解为不相交的静态子集 $S$ 和动态子集 $D(t)$，并为两者设计不同的参数化策略以最大化预算利用率。

核心问题是如何确定静态与动态的分界。CDGS 不依赖预标注掩码或手动设定的比例参数，而是完全数据驱动地推断分解：为每个高斯分配平移属性 $T_i$ 表征运动强度，在预热阶段后分析运动幅度的统计分布。方法构建 $B$ 个 bin 的平滑直方图，该直方图通常呈双峰分布——对应静态和动态两个群体。通过识别两个最显著的峰值 $p_s$ 和 $p_d$，并定位两者之间的谷底，方法自动确定分离阈值 $\tau_{\text{motion}}$，无需任何手动调参。

基于此分类，静态高斯 $S$ 维护时间不变属性（位置、旋转、尺度、颜色），仅附加轻量的每帧全局变换 $T_i$ 以应对轻微的相机偏差或光照变化；动态高斯 $D(t)$ 则配备显式时间参数——位置向量 $f^\mu_i$ 和旋转向量 $f^R_i$ 通过关键帧插值确定当前时刻的瞬时状态，加上可学习的时间激活窗口 $[t^s_i, t^e_i]$，使高斯仅在有效贡献的时间段内消耗计算资源。预算根据动静分解比例分配：动态区域和高活动帧获得更密集的表示预算，静态区域保持紧凑但稳定的覆盖。

\textbf{三阶段训练策略。}
为确保约束优化稳定收敛，CDGS 设计了三阶段渐进训练流程：第一阶段为预热初始化，不加预算约束，仅用渲染损失训练静态高斯和每帧全局变换，建立基础几何先验；第二阶段激活人口控制器，引入预算损失和正则化，门控温度逐步退火，同时周期性地重新执行动静分配和致密化/剪枝操作，将预算动态重分配到最需要的时空区域；第三阶段为稳定微调，当有效计数 $N_p$ 收敛至 $N_{\text{target}}$ 附近时，将软掩码 $c_i$ 显式二值化，永久剔除 $c_i = 0$ 的冗余高斯，剩余紧凑集合仅用渲染损失和正则化微调以恢复掩码过程可能损失的细节。

此外，CDGS 还设计了双模式混合压缩策略以进一步降低存储和传输开销：静态高斯通过前景/背景分离（剔除距离场景中心超过 $\mu_d + 3\sigma_d$ 的异常点以缩小坐标范围）、KD-Tree 空间重排序和预测编码压缩残差；动态高斯则将量化属性重塑为 2D 图像格式，按属性分组后使用 H.264 编码器压缩。实验表明，CDGS 在严格预算约束下实现高保真动态渲染，最终高斯数量误差控制在 2\% 以内，模型尺寸相比先前最优方法减少约 3 倍。

以上方法都依赖神经网络预测变形。Gaussian-Flow 则走上了一条截然不同的道路——完全抛弃神经网络，用纯数学模型（时域多项式和频域傅里叶）直接建模高斯的运动轨迹。

\paragraph{解析双域变形替代神经网络}

Per-Gaussian Embedding-Based Deformation 通过逐高斯嵌入实现了身份驱动的变形场，有效解决了坐标输入无法区分高斯基元身份的问题，但其变形函数 $F_\theta$ 仍依赖 MLP 神经网络，推理开销拖慢了 3DGS 原本极快的渲染流水线。Gaussian-Flow（2023.12）~\cite{2312_03431} 沿用"每个高斯基元独立建模变形"的身份驱动思路，但彻底抛弃神经网络，改用纯解析函数直接拟合每个高斯基元属性的时间变化曲线。

\textbf{动机。} MLP 神经网络的推理开销拖慢了 3DGS 原本极快的渲染流水线——纯解析函数（多项式 + 傅里叶级数）只需计算多项式与三角函数，完全无需 MLP 前向推理，渲染速度与原始 3DGS 基本持平。这一设计选择使训练速度比逐帧 3DGS 快 5 倍，渲染速度与原始 3DGS 基本持平——因为渲染时只需计算多项式与三角函数，完全无需 MLP 前向推理。

\textbf{双域变形模型（Dual-Domain Deformation Model, DDDM）：用解析函数替代 MLP。}
与锚点方法用 MLP 学习从 $(z_g, z_t)$ 到属性偏移的映射不同，Gaussian-Flow 的核心问题是：用什么解析函数能以最小参数量精准逼近高斯基元的时间变化轨迹？关键洞察在于单一拟合函数难以兼顾两种运动模式——时域多项式 $P_N(t) = \sum_{n=0}^{N} a_n t^n$ 善于捕捉平滑运动（小阶数即可拟合），但面对剧烈运动时需要大阶数，导致轨迹出现不合理振荡；频域傅里叶级数 $F_L(t) = \sum_{l=1}^{L} (f^l_{\sin} \sin(lt) + f^l_{\cos} \cos(lt))$ 善于捕捉剧烈与周期性变化，但处理平滑运动时需手动降阶。DDDM 将两者统一：每个高斯基元在参考帧 $t_0$ 处维护基准属性 $S_0 = \{\mu_0, q_0, c_0\}$（位置、四元数旋转、颜色），时间变化通过残差建模：

\begin{equation}
S(t) = S_0 + D(t), \quad D(t) = P_N(t) + F_L(t)
\end{equation}

多项式分量 $P_N(t) = \sum_{n=0}^{N} a_n t^n$ 负责捕捉平滑运动趋势，傅里叶分量 $F_L(t) = \sum_{l=1}^{L} (f^l_{\sin} \sin(lt) + f^l_{\cos} \cos(lt))$ 则捕捉周期性和剧烈变化——两者各司其职，粗运动由多项式承担，细节与突变由傅里叶补充。值得注意的是，高斯的缩放 $s$ 和不透明度 $\alpha$ 在时间上保持不变，仅位置、旋转和颜色参与变形，且每个属性的每个维度独立拥有自己的 $D(t)$（例如三维位置 $\mu$ 使用三个独立的残差函数 $\{D_{\mu_i}(t)\}_{i=0}^{3}$）。消融实验表明，单独使用多项式拟合虽平滑但易模糊细节，单独使用傅里叶级数虽锐利但引入更多伪影，而 DDDM 的混合函数兼具锐利细节与低伪影。

\textbf{自适应时间戳缩放：防止短时间大幅运动导致优化崩溃。}
归一化帧索引 $t \in [0, 1]$ 直接作为 $D(t)$ 的输入时，若某高斯基元在极短时间窗内发生大幅运动，多项式或傅里叶系数必须取极大值才能拟合，导致优化不稳定甚至崩溃。为此方法为每个高斯基元引入独立学习的时间膨胀因子，将输入时间重映射为 $t_s = \lambda_s t + \lambda_b$，其中 $\lambda_s$（膨胀因子）和 $\lambda_b$（基准因子）初始化为 1 和 0，优化过程中自适应调整——效果类似于"拉伸或压缩该高斯基元的局部时间轴"，使系数保持在合理范围内。

\textbf{正则化：时空连续性约束。}
方法设计两种正则化确保运动连续性。(1) 时间平滑损失 $\mathcal{L}_t = \|D(t) - D(t+\epsilon)\|_2$，约束相邻时刻的变形一致，扰动幅度自适应设为 $\epsilon = 0.1/T$（$T$ 为总帧数）。(2) KNN 刚性损失 $\mathcal{L}_s = \sum_{j \in \mathcal{N}_i} \|D(t)_i - D(t)_j\|_2$，约束空间邻近高斯具有相似变形——与锚点方法的嵌入平滑正则化（约束嵌入向量相似）相比，这里直接约束变形残差本身一致，且为应对致密化过程中邻域不断变化的难题，训练交替进行"带密度控制的优化"和"无密度控制的 KNN 优化"两个阶段。

此外，SplatFields（2024.09）~\cite{2409_11211} 用神经高斯泼溅实现稀疏 3D 和 4D 重建。

\subsubsection{原生4D高斯表示}

本路线介绍直接在时空域中定义高斯基元的原生4D表示方法，使时间成为表示的内在维度。

以上变形场方法通过"规范空间+变形网络"间接建模动态。本节转向另一条更直接的路线——不借助中间的变形网络，而是直接在时空域中定义高斯基元，让时间成为表示的内在维度。原生 4D 路线不依赖规范空间和变形网络，4D 高斯的属性天然包含时间维度，无需额外的变形网络。这种表示更直接，但需要解决 4D 参数量爆炸的问题。

\paragraph{HexPlane 张量分解原生 4D 高斯}

4D Gaussian Splatting (4D-GS)（2023.10）~\cite{2310_08528} 是最早将 3DGS 扩展为变形场方法的工作之一。

\textbf{动机。} 逐帧独立 3DGS 的存储随帧数线性增长，纯隐式神经场方法的 MLP 前向推理拖慢渲染——需要一种既能紧凑存储又能实时渲染的动态表示。与逐帧独立 3DGS（存储随帧数线性增长）和纯隐式神经场方法（MLP 前向推理拖慢渲染速度）不同，4D-GS 提出一种高效的高斯变形场网络（Gaussian Deformation Field），仅维护一组规范 3D 高斯，通过时空结构编码器和多头变形解码器预测每帧的属性偏移。

\textbf{改动概要。}本方法修改了 3DGS 的动态表示环节——用"一组规范 3D 高斯 + 可学习的变形场网络"替代"逐时间戳独立 3D 高斯"——3DGS 的泼溅渲染管线、球谐函数颜色模型和基于 SfM 的点云初始化保持不变。直观地说，4D-GS 不为每一帧存储一整套高斯，而是只保留一帧的"模板高斯"，用一个轻量网络告诉每个高斯"在当前时刻应该移动到哪里、形状如何变化"。

\textbf{Baseline 的问题与设计动机。}在 3DGS 中，每个 3D 高斯由位置 $X \in \mathbb{R}^3$、球谐颜色系数 $C \in \mathbb{R}^k$（$k$ 为球谐系数数量）、不透明度 $\alpha \in \mathbb{R}$、旋转四元数 $r \in \mathbb{R}^4$ 和缩放因子 $s \in \mathbb{R}^3$ 定义，通过泼溅（splatting）渲染：$I = S(M, G)$，其中 $M = [R, T]$ 是相机矩阵，$S$ 是可微泼溅算子。对于动态场景，最直接的方案（如 Dynamic 3DGS）是为每个时间戳 $t_i$ 独立存储一套 3D 高斯参数，导致存储随帧数线性增长 $O(tN)$（$N$ 为高斯数量），长序列场景代价不可忽视。纯隐式神经场方法（如 TiNeuVox、K-Planes）虽可建模时空，但 MLP 前向推理引入不可忽视的渲染延迟。4D-GS 的核心动机是同时解决这两个问题：仅维护 $O(N+F)$ 的存储（$F$ 为变形场网络参数量），同时保留 3DGS 的实时泼溅渲染能力。

\textbf{时空结构编码器。}方法的核心组件是时空结构编码器 $H$。直觉上，它要回答的问题是：给定场景中一个高斯的位置和时间，这个位置附近正在发生怎样的运动？为此，$H$ 利用多分辨率 K-Planes 模块将 4D 时空体积分解为 6 个 2D 平面——3 个空间平面 $R_l(x,y), R_l(x,z), R_l(y,z)$ 和 3 个时间平面 $R_l(x,t), R_l(y,t), R_l(z,t)$，每个平面在多尺度 $l \in \{1,2\}$ 上定义。这类似于将一个 4D 立方体沿不同投影方向"压扁"为 2D 切片，每个切片只记录沿该投影的信息；拼接所有切片即可恢复近似完整的 4D 特征，但存储从 $O(N^4)$ 大幅降至 $O(6N^2)$。每个平面模块定义为 $R(i,j) \in \mathbb{R}^{h \times lN_i \times lN_j}$，其中 $h$ 为特征隐藏维度，$N$ 为基础分辨率（实验中设为 64），$l$ 为上采样倍率。给定高斯位置 $\mu = (x,y,z)$ 和时间 $t$，从各平面双线性插值取特征后拼接：

\begin{equation}
f_h = \text{concat}_l(\text{interp}(R_l(i,j))), \quad (i,j) \in \{(x,y),(x,z),(y,z),(x,t),(y,t),(z,t)\}
\end{equation}

即把高斯坐标投影到 6 个平面上，在每个平面的 4 个最近网格顶点间做双线性插值，得到 $f_h \in \mathbb{R}^{h \times l}$ 维特征向量。再经小型 MLP $\phi_d$ 融合为时空特征 $f_d = \phi_d(f_h)$。这种设计的关键优势是使邻近高斯能够共享时空上下文信息——落在同一网格单元内的高斯会查询到相似的特征——从而预测出更准确、更一致的运动和变形，避免相邻高斯的变形出现撕裂（avulsion）。

\textbf{多头变形解码器。}$D = \{\phi_x, \phi_r, \phi_s\}$ 由三个独立 MLP 组成，分别从时空特征 $f_d$ 解码位置偏移 $\Delta X = \phi_x(f_d)$、旋转偏移 $\Delta r = \phi_r(f_d)$ 和缩放偏移 $\Delta s = \phi_s(f_d)$。变形后的高斯属性通过残差连接得到：

\begin{equation}
(X', r', s') = (X + \Delta X, \; r + \Delta r, \; s + \Delta s)
\end{equation}

即将规范空间中的高斯"搬运"到当前时刻的位置和形状，残差设计保证了变形学习的稳定性——网络只需预测偏移量而非绝对值。变形后的 3D 高斯 $G^* = \{X', s', r', \alpha, C\}$ 直接通过标准 3DGS 的泼溅管线渲染：$I = S(M, G^*)$，颜色 $C$ 和不透明度 $\alpha$ 在变形过程中保持不变。论文还指出，对于流体或非刚性运动等外观会发生显著变化的场景，可额外引入颜色偏移解码器 $\phi_C$ 和不透明度偏移解码器 $\phi_\alpha$，使变形场也能建模外观变化。

\textbf{优化策略。}训练分为两个阶段。前 3000 次迭代为热身阶段，仅使用静态 3DGS 渲染（$I = S(M, G)$），不使用变形场。这一设计有三个好处：(a) 让部分 3D 高斯自然停留在动态区域，减轻变形场需要学习大幅运动的压力；(b) 学习到合理的静态高斯分布，引导变形场更关注动态部分；(c) 避免变形场网络训练初期的数值不稳定。热身之后，变形场与 3D 高斯联合训练，总迭代次数约为 20000 次（合成数据集）。损失函数使用 L1 颜色损失加全变分正则化：$\mathcal{L} = |I - \hat{I}| + \lambda_{tv} \mathcal{L}_{tv}$。K-Planes 模块的学习率设为 $1.6 \times 10^{-3}$（训练结束时衰减至 $1.6 \times 10^{-4}$），变形解码器的学习率设为 $1.6 \times 10^{-4}$（衰减至 $1.6 \times 10^{-5}$）。

\textbf{效率与性能。}K-Planes 分解将存储复杂度从 $O(N^4)$ 降低到 $O(6N^2)$，整个方法在单张 RTX 3090 上训练，合成数据集（D-NeRF）达到 82 FPS（$800 \times 800$ 分辨率）、真实数据集（Neu3D）达到 30 FPS（$1352 \times 1014$ 分辨率）的实时渲染速度。当高斯数量不超过 30000 时，渲染速度可达 90 FPS。方法总存储开销仅为 3D 高斯约 10 MB + 变形场网络约 8 MB，远低于逐帧独立方法。

\paragraph{时空高斯基元与 K-平面分解}

回顾前一方法：4D Gaussian Splatting（4D-GS）采用"规范 3D 高斯 + 变形场网络"的范式——一组静态高斯作为模板，由 K-Planes 编码器和多头解码器预测每帧的属性偏移。这种范式保留了 3DGS 的实时泼溅能力，但变形场网络引入了额外的中间表示：高斯本身不包含时间信息，运动完全依赖外接网络的预测。一个自然的问题是：能否让高斯基元本身天然包含时间维度，从而完全跳过变形场网络这一中间环节？

Real-time Photorealistic Dynamic Scene Representation and Rendering with 4D Gaussian Splatting（2023.10）~\cite{2310_10642} 给出了肯定的回答。

\textbf{动机。} 变形场方法依赖规范空间到各帧的显式运动追踪，在复杂真实场景中面临全局运动追踪的歧义——让高斯基元本身天然包含时间维度，可以完全跳过变形场这一中间环节。方法从 6D 全光函数（plenoptic function）的视角出发，指出已有动态场景方法可以归为两类：第一类使用 MLP 或网格直接参数化 6D 全光函数，不显式建模场景运动，其表达能力受限于数据结构的内在特性——参数共享导致时空位置之间的干扰，或者过于独立而无法利用物体运动带来的关联性。第二类——即前一方法所属的变形场路线——基于规范 3D 模型的运动或变形来显式学习场景运动，但在复杂真实场景（如突然出现/消失的物体）中面临全局运动追踪的歧义和奇异性问题。为同时突破这两类方法的局限，本方法提出将动态场景建模重新形式化为底层时空 4D 体积的直接拟合问题，并设计原生 4D 高斯基元，无需任何变形网络。

\textbf{改动概要。}与前一方法的核心区别在于表示范式：前一方法维护一组规范 3D 高斯，通过网络预测偏移实现动态；本方法直接在 4D 时空 $(x,y,z,t)$ 中定义高斯基元，将"规范空间 + 变形"的两步过程合并为"一步到位"的 4D 建模。具体而言，前一方法的 K-Planes 编码器和多头变形解码器被完全移除，取而代之的是每个高斯自带 4D 均值和 $4\times 4$ 协方差矩阵；渲染不再经过"变形再泼溅"，而是通过条件概率将 4D 高斯切片为 3D 高斯后直接泼溅；球谐函数被推广为 4D 球面柱谐函数（4DSH），使颜色同时依赖视角和时间。3DGS 的泼溅渲染管线和基于 SfM 的点云初始化保持不变。

\textbf{原生 4D 高斯表示。}每个 4D 高斯具有 4D 均值 $\boldsymbol{\mu}_4 = (\mu_x, \mu_y, \mu_z, \mu_t) \in \mathbb{R}^4$ 和 $4 \times 4$ 协方差矩阵 $\boldsymbol{\Sigma}_4 = R S S^T R^T$。与前一方法中 3D 高斯的 $\boldsymbol{\Sigma} \in \mathbb{R}^{3\times 3}$ 不同，这里 $\boldsymbol{\Sigma}_4 \in \mathbb{R}^{4\times 4}$ 直接编码了空间维度与时间维度之间的耦合关系。缩放矩阵 $S = \text{diag}(s_x, s_y, s_z, s_t)$ 将缩放从 3 个维度扩展到 4 个维度，新增的 $s_t$ 控制高斯在时间轴上的"宽度"——快速运动的物体需要在时间维度上更"细"（$s_t$ 较小）以捕捉瞬间状态，静止物体则在时间上更"胖"（$s_t$ 较大），在多帧中持续存在。

\textbf{4D 旋转与运动建模。}这是本方法与前述方法的关键区别之一。前一方法通过 MLP 预测显式的位置、旋转、缩放偏移来建模运动；本方法则通过 4D 旋转直接在时空结构中隐式编码运动。直觉上，在 3D 空间中，旋转改变物体朝向；但在 4D 时空 $(x,y,z,t)$ 中，旋转可以将空间位移转化为时间演化——这恰好对应了物体运动的本质。4D 旋转由一对四元数 $q_l = (a,b,c,d)$ 和 $q_r = (p,q,r,s)$ 构造：$R = L(q_l) R(q_r)$，其中 $L(q_l)$ 和 $R(q_r)$ 分别为左右等距旋转矩阵。这对四元数共 8 个参数（扣除 2 个归一化约束后为 6 个自由度），对应 4D 旋转的 6 个自由度。初始化时 $q_l$ 和 $q_r$ 均设为单位四元数（恒等旋转），训练过程中通过渲染损失的反向传播自动学习。论文通过消融实验验证了 4D 旋转的核心作用：去掉 4D 旋转（将 $\boldsymbol{\Sigma}_4$ 限制为空间-时间分块对角，即假设空间与时间独立）后，渲染质量显著下降，证明时空耦合的协方差结构对捕捉场景运动至关重要。

\textbf{条件概率切片渲染。}前一方法在变形后直接得到 3D 高斯，通过标准泼溅渲染；本方法需要先将 4D 高斯"投影"到给定时刻的 3D 空间。这一步骤通过多元高斯分布的条件概率公式完成——给定查询时间 $t$，从 4D 高斯中提取条件 3D 高斯：

\begin{align}
\boldsymbol{\mu}_{xyz|t} &= \boldsymbol{\mu}_{1:3} + \boldsymbol{\Sigma}_{1:3,4} \boldsymbol{\Sigma}_{4,4}^{-1}(t - \mu_t) \\
\boldsymbol{\Sigma}_{xyz|t} &= \boldsymbol{\Sigma}_{1:3,1:3} - \boldsymbol{\Sigma}_{1:3,4} \boldsymbol{\Sigma}_{4,4}^{-1} \boldsymbol{\Sigma}_{4,1:3}
\end{align}

这两条公式的物理含义值得仔细解读。条件均值 $\boldsymbol{\mu}_{xyz|t}$ 表明：3D 位置不再是固定的，而是随时间 $t$ 线性变化，变化速率由时空交叉协方差 $\boldsymbol{\Sigma}_{1:3,4}$ 决定。这正是 4D 高斯自然编码运动的机制——无需额外的变形网络，运动信息已经内嵌在协方差矩阵的时空交叉项中。条件协方差 $\boldsymbol{\Sigma}_{xyz|t}$ 表明：确定了时间 $t$ 后，空间上的不确定性会减小，减小量恰好等于时间维度提供的"额外信息"。同时，边际概率 $p(t) = \mathcal{N}(t; \mu_t, \boldsymbol{\Sigma}_{4,4})$ 作为 1D 高斯控制每个 4D 高斯在时间 $t$ 处的不透明度权重——当查询时间远离 $\mu_t$ 时，高斯的贡献自动衰减。切片得到的 3D 高斯 $\mathcal{N}(\boldsymbol{\mu}_{xyz|t}, \boldsymbol{\Sigma}_{xyz|t})$ 随后通过标准 3DGS 泼溅管线渲染，只需对原始 tile-based 光栅化器做最小修改（增加边际概率 $p(t)$ 的累积）。

\textbf{4D 球面柱谐函数（4DSH）。}前一方法中，规范 3D 高斯的颜色通过标准球谐函数建模，仅依赖视角方向，颜色偏移是可选的附加模块。本方法需要颜色同时随视角和时间变化——因为每个 4D 高斯覆盖一个时间窗口，其间物体外观可能因光照变化、遮挡关系变化等原因而改变。一个直观的方案是为同一空间位置在不同时刻分配不同的高斯，但这导致表示冗余且优化困难。为此，方法提出 4D 球面柱谐函数（4D Spherindrical Harmonics, 4DSH），将球谐函数与傅里叶级数融合为统一的正交基：

\begin{equation}
Z^m_{nl}(t, \theta, \phi) = \sin\!\Big(\frac{2\pi n}{T} t\Big) \, Y^m_l(\theta, \phi)
\end{equation}

其中 $Y^m_l$ 是标准 3D 球谐函数（$l \geq 0$ 为阶数，$-l \leq m \leq l$ 为次数），$n$ 是傅里叶级数的阶数，$T$ 为场景时长。直觉上，4DSH 将"视角怎么影响颜色"（球谐函数）和"时间怎么影响颜色"（傅里叶级数）解耦为两个独立的基函数族，再通过张量积组合——每个 4DSH 基函数同时捕捉一种空间方向频率和一种时间频率的外观变化模式。论文通过消融实验表明，移除 4DSH（退化为标准 SH）会导致渲染质量下降，证实了显式建模颜色时间演化的必要性。

\textbf{时空密度控制。}前一方法沿用标准 3DGS 的密度控制策略，仅在空间维度上操作。本方法将密度控制扩展到时间维度，增加两个关键机制。第一，将时间维度的位置梯度 $\nabla_{\mu_t}$ 的均值作为额外的密度控制指标：如果某个区域的 4D 高斯在时间维度上的重建不充分（$\nabla_{\mu_t}$ 过大），则触发克隆或分裂操作。第二，在过度重建区域进行高斯分裂时，采用联合空间-时间位置采样：新高斯的位置不再仅在 3D 空间中随机采样，而是在完整的 4D 高斯分布（作为概率密度函数）中采样，使新高斯在空间和时间两个维度上同时分散。这确保了动态区域的 4D 高斯在时空两个维度上都有足够的覆盖。

\textbf{端到端优化与效率。}与前一方法的两阶段训练（先静态热身再联合训练变形场）不同，本方法的优化过程是完全端到端的——所有 4D 高斯参数（均值、协方差、4DSH 系数、不透明度）从第一个迭代起就同时优化，无需任何热身或分阶段策略。训练采用标准 30000 次迭代，前半段进行密度控制，后半段停止密度控制仅优化参数。对于变化剧烈的场景，论文发现简单的时间维度批量采样（batch sampling in time）优于引入额外正则化，因为它保证了时间维度上的均匀监督覆盖。渲染时，对边际概率 $p(t) < 0.05$ 的高斯施加高斯滤波截断，避免时间上过远的高斯影响当前帧。在单张 RTX 3090 上，方法在 Plenoptic Video 数据集（真实多视角）上实现实时渲染速度，显著超过所有前方法；在 D-NeRF 数据集（合成单目）上同样取得最优质量。

\paragraph{4D 椭球基元与 4D 旋转-缩放参数化}

前一方法（Real-time Photorealistic 4DGS）首次提出了原生 4D 高斯基元的概念与完整实现；本方法——4D Gaussian Splatting: Modeling Dynamic Scenes with Native 4D Primitives（2024.12）~\cite{2412_20720}——是其期刊扩展版本。核心表示（4D 均值与协方差、条件概率切片、4DSH 颜色建模、时空密度控制）与前一方法完全一致，扩展版本不再赘述。期刊版的新贡献围绕一个实际问题展开：原生 4D 高斯基元虽然免去了变形场网络，但用大量分段线性基元逼近动态场景运动带来了显著的存储冗余，一个典型序列的模型体积可达 1183 MB。扩展版本从三个方向回应了这一挑战。

\textbf{紧凑变体：4DGS$_C$。}存储压缩从两条路径同时入手——降低单个高斯的属性编码代价，以及削减高斯总数。在属性编码方面，论文观察到不同属性对渲染质量的敏感度差异显著：位置参数最敏感，仅施加半精度量化；不透明度最不敏感，使用 8-bit Min-Max 量化即可；高阶 4DSH 系数虽然重要但高度冗余，对基础颜色和剩余分量分别使用残差向量量化（Residual Vector Quantization, R-VQ）——缩放属性用 6 个码本（每码本 64 条目），旋转属性共享另一组同样规模的 R-VQ 模型，颜色分量的基础部分和剩余部分分别用 6 个码本（64 条目）和 8 个码本（256 条目）压缩；码本在训练后的高斯参数上构建，量化后额外微调码本条目以恢复精度损失，最后对码本索引施加 Huffman 编码完成无损压缩。实验表明，仅属性压缩即可将形状相关参数从 88.19 MB 压缩至 24.20 MB（约 3.5 倍压缩率），量化导致的 0.49 dB PSNR 下降可通过码本微调完全恢复（+0.51 dB）。在削减高斯总数方面，论文引入可学习的显著性掩码（significance mask）：为每个高斯附加参数 $m \in \mathbb{R}$，显著性得分 $M = \sigma(m)$，在训练中通过掩码压缩损失 $\mathcal{L}_{\text{mask}} = \frac{1}{N}\sum_i^N M_i$ 推动低显著性高斯趋向零；同时修改高斯的缩放矩阵和不透明度，使低显著性高斯自动"淡出"。实验表明，掩码修剪可将高斯点数从 1.26M 降至 0.51M（2.47 倍削减），PSNR 基本不受影响。两条路径组合后，完整模型从 1183.21 MB 压缩至 56.74 MB（约 20 倍压缩），渲染质量几乎不变。

\textbf{应用扩展。}扩展版本将 4DGS 的适用范围从受控多视角场景拓展至三个新领域。第一，大规模动态城市场景——驾驶视频具有无界范围和稀疏视角两个独特挑战。方法利用 LiDAR 点云（带时间戳聚合）初始化 4D 高斯位置，引入可学习天空立方体贴图（cube map）建模远景，并设计了一套针对性正则化：$L_{\text{sparse}}$ 鼓励左右四元数 $q_l, q_r$ 接近以注入静态先验；$L_{\text{covt}}$ 鼓励高斯在时间维度上扩展寿命以利用多视角监督；$L_{\text{rigid}}$ 约束空间邻域高斯的运动一致性（刚体先验）；LiDAR 深度监督弥补纹理匮乏区域的几何重建。训练采用热身策略——前 3000 步保持 $q_l = q_r$ 模拟静态场景，并添加时间随机扰动以抑制背景区域的不必要运动。在 Waymo Open Dataset 上，4DGS 在重建和新视角合成两个任务上均达到与专用城市场景流水线可比甚至更优的性能，同时保持高效渲染。第二，生成式 4D 内容创建——从单目视频通过生成先验创建 4D 内容。

方法利用多视角扩散模型（SyncDreamer）为每帧生成伪多视角图像，结合 Stable-Zero123 的 Score Distillation Sampling（SDS）损失提供 2D 扩散先验，整体损失函数联合参考视角重建损失、伪多视角重建损失和 SDS 损失优化 4DGS。在 Consistent4D 数据集上，4DGS 的生成时间和质量均显著优于 C-DyNeRF 和 Deform-GS。第三，4D 时空场景理解——受 SAGA（3D 高斯分割）启发，为每个 4D 高斯附加尺度门控特征向量 $f \in \mathbb{R}^D$（$D=32$），通过尺度映射 $S: [0,1] \to [0,1]^D$ 实现多粒度控制，在预训练 4DGS 上利用 2D SAM 掩码进行尺度感知对比学习训练特征。

推理阶段支持两种模式：基于 HDBSCAN（Hierarchical Density-Based Spatial Clustering of Applications with Noise，层次密度聚类）的自动场景分解（无需任何提示），以及基于 2D 点提示的 4D 交互式分割——用户在某一帧点击，系统通过特征匹配在全部时空视角上定位同一物体。

\textbf{全面评估。}扩展版本在原始 Plenoptic Video 和 D-NeRF 数据集之外，新增 Technicolor 光场视频数据集（5 场景，$2560\times 1920$ 分辨率）、Waymo 动态城市场景数据集（32 序列）和 Consistent4D 生成基准（7 视频）的评测。在所有基准上，4DGS 在渲染质量和/或效率方面均达到最优或可比水平。紧凑变体 4DGS$_C$ 在 Technicolor 数据集上验证了压缩策略的通用性——大幅缩减存储开销的同时保持与未压缩版本几乎一致的渲染质量。

\paragraph{时空高斯特征泼溅与时间感知不透明度}

Spacetime Gaussian Feature Splatting（2023.12）~\cite{2312_16812} 提出时空高斯（Spacetime Gaussian, STG），同样将 3D 高斯扩展到 4D 时空域以建模动态场景，但选择了与 4D-GS 不同的技术路线。回顾 4D-GS 的设计：它维护一组规范 3D 高斯，通过 K-Planes 时空编码器和多头变形解码器（小型 MLP）预测每帧的位置、旋转、缩放偏移。这种"规范空间 + 变形网络"的策略将动态建模的负担集中在变形场上——变形场需要学会每个高斯在任意时刻应该移动到哪里、如何变形。然而，变形场本质上是连续运动建模：它假设高斯始终存在于场景中，只是位置和形状随时间变化。对于突然出现或消失的瞬态内容（例如走过的行人、飞溅的水花），变形场难以优雅地表达"这个高斯在某个时间点之前/之后不存在"。

\textbf{改动概要。}STG 的核心思路转变是：不再用额外的神经网络来回答"高斯如何随时间变化"，而是把时间直接嵌入高斯基元本身。具体而言，4D-GS 的变形场网络（K-Planes 编码器 + 多头 MLP 解码器）被完全移除，取而代之的是：时间径向基函数控制不透明度的"生命周期"、多项式轨迹参数化运动、9 维特征向量替代球谐函数。3DGS 的泼溅渲染管线保持不变，但渲染输入从球谐颜色变为紧凑特征向量 + 轻量 MLP。直观地说，如果说 4D-GS 让高斯"站在原地、由网络告诉它们怎么动"，STG 则让每个高斯自己"携带一份时间表"——它知道自己什么时候该出现（时间中心 $\mu_i^{\tau}$）、存活多久（时间缩放 $s_i^{\tau}$）、沿着什么路径移动（多项式系数）。

\textbf{时间径向基函数与瞬态建模。}4D-GS 中，每个高斯的不透明度 $\alpha$ 在变形过程中保持不变——高斯永远存在，只是位置和形状变了。STG 则为每个高斯引入了独立的时间生命周期。具体地，每个 STG 的不透明度由空间高斯函数和时间径向基函数联合定义：

\begin{equation}
\alpha_i(t) = \alpha_i^s \exp(-s_i^{\tau} |t - \mu_i^{\tau}|^2)
\end{equation}

其中 $\mu_i^{\tau}$ 是时间中心（STG 最可见的时间戳），$s_i^{\tau}$ 是时间缩放因子（控制有效持续时间，即时间不透明度保持较高的时段），$\alpha_i^s$ 是与时间无关的空间不透明度。这一设计的核心优势在于：瞬态内容（仅在特定时间段出现/消失的物体）可以被建模为一组在各自时间窗口内"激活"的 STG，无需引入额外的可见性机制。当 $t$ 远离 $\mu_i^{\tau}$ 时，$\alpha_i(t)$ 自然衰减至接近零，该 STG 在渲染中自动"消失"。这种参数化也带来存储优势——不需要为每个时间步存储独立的高斯集，只需为每个高斯增加两个额外参数（$\mu_i^{\tau}$ 和 $s_i^{\tau}$）。

\textbf{多项式运动轨迹。}4D-GS 使用 MLP 变形解码器从时空特征预测位置偏移 $\Delta X = \phi_x(f_d)$。STG 采用了更简洁的方案——多项式轨迹参数化：

\begin{equation}
\mu_i(t) = \sum_{k=0}^{n_p} b_{i,k}(t - \mu_i^{\tau})^k
\end{equation}

其中 $\{b_{i,k}\}_{k=0}^{n_p}$ 是每个 STG 各自的多项式系数，在训练中直接优化（不需要经过任何神经网络）。实践中使用 $n_p=3$ 次多项式，在表示能力和模型尺寸之间取得平衡。旋转也采用类似的多项式参数化（$n_q=1$ 次），缩放保持时间不变（论文实验发现对缩放添加时间依赖并未带来质量提升）。与 4D-GS 的 MLP 变形场相比，多项式参数化有两个特点：(a) 每个高斯的运动是独立参数化的，不共享网络权重，因此不受邻近高斯运动模式的影响，但也无法像 4D-GS 的 K-Planes 编码那样利用邻近高斯的时空上下文来正则化运动；(b) 计算开销极低——仅需多项式求值，无需 MLP 前向推理，有利于渲染速度。

\textbf{特征泼溅渲染：替代球谐函数。}4D-GS 保留了 3DGS 的球谐函数颜色模型（3 阶 SH 需要 48 个参数），颜色 $C$ 在变形过程中保持不变。STG 选择了完全不同的颜色表示方案——用紧凑的 9 维特征向量替代球谐函数：

\begin{equation}
f_i(t) = [f_i^{\text{base}}, \; f_i^{\text{dir}}, \; (t-\mu_i^{\tau}) f_i^{\text{time}}]^T
\end{equation}

其中 $f_i^{\text{base}} \in \mathbb{R}^3$ 编码基础 RGB 颜色，$f_i^{\text{dir}}, f_i^{\text{time}} \in \mathbb{R}^3$ 分别编码与视角和时间相关的信息。泼溅过程与 3DGS 相同，只是用特征 $f_i(t)$ 替代颜色 $c_i$。泼溅到屏幕空间后，在每个像素处将特征拆分为三部分，通过一个 2 层 MLP $\Phi$ 解码为最终颜色：

\begin{equation}
I = F^{\text{base}} + \Phi(F^{\text{dir}}, F^{\text{time}}, r)
\end{equation}

其中 $r$ 是该像素的视角方向。相比球谐函数的 48 个参数，9 维特征将每个高斯的存储开销降低约 5 倍。论文还提供一个轻量版本：完全去掉 MLP $\Phi$，仅保留 $F^{\text{base}}$（即 $I = F^{\text{base}}$），以牺牲视角和时间依赖性换取极致速度——8K 分辨率渲染可达到 60 FPS。

\textbf{引导采样策略。}4D-GS 沿用 3DGS 的标准密度控制（基于位置梯度的分裂和克隆），在已有高斯附近逐步扩展。STG 观察到这一策略的盲区：初始化时高斯稀疏覆盖的远处区域（尤其是远离训练相机的区域）渲染质量差，而标准密度控制只在已有高斯附近生长，难以"凭空"在远处创建新高斯。为此，STG 引入引导采样策略：在训练误差稳定后，以训练误差和粗深度图为引导，在误差大的区域沿射线均匀采样新高斯。具体流程为：(a) 将训练误差按 patch 聚合（避免噪声像素的干扰），选中误差大的 patch；(b) 从选中 patch 的中心像素发射射线，利用高斯中心的粗深度图（在特征泼溅过程中已自然生成，几乎无额外开销）确定更精确的深度范围；(c) 在该深度范围内沿射线均匀采样新高斯，并添加少量位置噪声。不必要的新高斯在后续训练中不透明度会自然降低并被剪枝。整个采样过程在实验中仅需执行不超过 3 次，即可显著改善远距离区域的渲染质量。这一策略与 3DGS 的标准密度控制是互补的：标准密度控制通过分裂在已有高斯附近逐步生长，引导采样则在完全缺乏高斯的区域"播种"。

\paragraph{神经 ODE 连续时间动力学系统}

EvoGS（2025.12）~\cite{2512_19648} 在上述 4D-GS 的基础上做了一次根本性的范式转换：将离散时间的残差变形场替换为连续时间的神经速度场。回顾前文，4D-GS 用 K-Planes 六平面编码时空特征，再通过多头 MLP 直接预测每帧的位移偏移 $\Delta X, \Delta r, \Delta s$——这隐含了一个离散时间假设：变形只在训练帧对应的时间点被定义，帧与帧之间的运动轨迹未被建模。当训练帧密集时这不构成问题，但若帧间隔拉大（如仅保留每第 8 帧），变形场就缺少足够的时间监督来推断中间运动，导致严重的轨迹漂移和伪影。EvoGS 的核心思想是：不再让网络预测``高斯在哪里''（离散位移），而是学习``高斯往哪走''（连续速度），从而通过 ODE 积分自然插值出任意时刻的状态。这使得重建和外推被统一在同一个连续动力空间中——积分步长可以任意加密或外延，无需额外的网络查询。

\textbf{改动概要。}与 4D-GS 相比，EvoGS 在三个环节做出关键改变：(1) 将变形解码器的输出从残差偏移 $\Delta X, \Delta r, \Delta s$ 改为瞬时导数 $d\mathbf{x}/dt$，并将离散查表替换为 ODE 积分；(2) 将旋转的参数化从简单的残差加法 $r' = r + \Delta r$ 改为指数映射更新（exponential-map update），保证旋转始终位于流形 $\text{SO}(3)$ 上；(3) 引入高斯航点机制和轨迹一致性正则，专门应对 ODE 积分的长程漂移问题。时空特征编码仍然采用与 4D-GS 相同的六平面因子化结构（HexPlane），但下游的使用方式从``查表取偏移''变为``查表驱动速度场''。3DGS 的泼溅渲染管线、SfM 点云初始化和密度控制策略保持不变。

\textbf{连续时间 ODE 动力学。}在这一框架下，每个高斯基元拥有完整状态向量 $\mathbf{x}_i(t) = [\mathbf{p}_i(t), R_i(t), S_i(t), c_i(t), \alpha_i(t)]$，涵盖位置、旋转、缩放、颜色和不透明度。该状态通过连续时间 ODE 演化：

\begin{equation}
\frac{d\mathbf{x}_i}{dt} = \mathbf{v}_\theta(\mathbf{x}_i(t), \mathbf{f}_i(t), t),
\end{equation}

其中 $\mathbf{v}_\theta$ 是一个轻量级 MLP，负责预测位置和缩放的瞬时导数以及旋转的指数映射增量（颜色和不透明度不参与动力学演化）。时空特征 $\mathbf{f}_i(t)$ 仍由六平面因子化结构通过双线性插值获得，但与 4D-GS 不同的是，这里特征不再直接解码为变形偏移，而是驱动一个连续速度场。积分采用四阶 Runge--Kutta（RK4）求解器（消融实验表明一阶 Euler 积分器在长程外推中迅速积累漂移），支持双向传播——既可以从过去积分到未来，也可以从未来反向积分回过去：

\begin{equation}
\mathbf{x}_i(t_1) = \text{RK4}(\mathbf{x}_i(t_0),\, t_0,\, \Delta t,\, \mathbf{v}_\theta), \quad \mathbf{x}_i(t_0) = \text{RK4}(\mathbf{x}_i(t_1),\, t_1,\, -\Delta t,\, \mathbf{v}_\theta).
\end{equation}

双向积分的可逆性使 EvoGS 能够穿越缺失帧或观测模糊区域传播运动信息。

\textbf{高斯航点与轨迹一致性正则。}长时间跨度的 ODE 积分会因数值误差而累积漂移。EvoGS 借鉴经典 Kalman 滤波中``预测-校正''的思想，引入高斯航点（Gaussian Waypoint）机制：训练期间在时间轴上存储 3 个锚点快照 $A = \{t^{(a)}_1, t^{(a)}_2, t^{(a)}_3\}$（设在序列的起始、中点和末尾），每个锚点保存该时刻的高斯状态。渲染目标帧时，系统不从起始时刻全程积分，而是从最近的锚点 $t^{(a)}$ 重新初始化 ODE 状态，仅向前积分到目标时刻 $t$，从而将有效积分区间缩短至相邻锚点之间，大幅抑制漂移累积。锚点一致性损失 $\mathcal{L}_{\text{anchor}} = \sum_{t^{(a)} \in A} \|\mathbf{x}(t^{(a)}) - \hat{\mathbf{x}}(t^{(a)})\|_2^2$ 鼓励积分状态与锚点快照保持一致，其中 $\hat{\mathbf{x}}(t^{(a)})$ 是存储的锚点状态，$\mathbf{x}(t^{(a)})$ 是从前一个锚点积分得到的状态——这与经典 Kalman 滤波的不同之处在于，EvoGS 不维护显式的速度估计或协方差，锚点仅作为稀疏的固定参考状态来约束长程积分。此外，EvoGS 引入轨迹一致性正则项 $\mathcal{L}_{\text{coh}}$，对规范空间中每个高斯的 $K$ 个近邻，将它们前向积分一个小步长 $h$ 后，施加距离加权的散度惩罚：$\mathcal{L}_{\text{coh}} = \sum_{i} \sum_{j \in N(i)} w_{ij} \|\tilde{\mathbf{x}}_i - \tilde{\mathbf{x}}_j\|_2^2$，其中 $w_{ij} = \exp(-\|\mathbf{x}_i - \mathbf{x}_j\|)$（$\sigma$ 设为 12 倍近邻距离）。关键之处在于，这项正则比较的是积分后的轨迹而非瞬时速度，直接约束了长程运动的一致性。完整的训练目标为 $\mathcal{L} = \mathcal{L}_{\text{photo}} + \lambda_{\text{coh}} \mathcal{L}_{\text{coh}} + \lambda_{\text{anchor}} \mathcal{L}_{\text{anchor}} + \lambda_{\text{tv}} \mathcal{L}_{\text{tv}}$，其中 $\mathcal{L}_{\text{tv}}$ 平滑时空特征平面。整个方法在单张 NVIDIA L40 GPU 上训练，采用 4D-GS 的训练设置并做少量调整以适配连续时间动力学。

\textbf{连续时间建模的独特能力。}这种 ODE 范式赋予了 EvoGS 三个离散变形场无法实现的能力：(1) \textbf{稀疏帧重建}——仅用三分之一训练帧即可学习一致运动，在 D-NeRF 基准上以 stride-8 训练时 PSNR 显著优于离散方法，因为 ODE 积分在帧间提供了物理上合理的插值，而非让变形网络在缺乏监督的时间点上``猜测''偏移量；(2) \textbf{未来外推}——通过继续前向积分预测训练窗口之外的帧，实验中在仅用前 75\% 帧训练时仍能合理预测后续运动（此时最后一个锚点被移除以评估纯外推能力），反向积分同样可行；(3) \textbf{可组合运动编辑}——由于动力学编码为向量场，可以通过空间掩码混合不同速度场实现局部运动注入：$\mathbf{v}_{\text{mix}}(\mathbf{x},t) = m(\mathbf{x})\,\mathbf{v}_\theta(\mathbf{x},t) + (1 - m(\mathbf{x}))\,\mathbf{v}_{\text{inj}}(\mathbf{x},t)$，其中 $m(\mathbf{x}) \in [0,1]$ 控制每个空间区域跟随哪个动力学场，无需重新训练即可为选中区域注入新的运动。论文进一步指出，通过 3D 高斯分割掩码隔离目标物体、利用 Zero123 补全未观测几何后重新插入场景，可以解决注入运动时因视角不全导致的物体不完整问题。不过，EvoGS 的外推能力本质上依赖数据驱动——在需要真正因果推理或物理推理的场景（如预判水杯中水的流动），学到的速度场可能无法泛化；在极端时间稀疏（如 stride-20）下，动力学约束不足，优势也随之消退。

\paragraph{7D 时空-角度联合高斯泼溅}

4D-GS 的变形场仅建模时空动态，视角相关效果仍依赖球谐函数颜色模型——对于强镜面反射、半透明折射等需要视角方向与场景动态联合建模的现象，球谐函数的表达能力不足。7DGS（2025.03）~\cite{2503_07946} 的核心思路是不再用"规范高斯 + 变形场网络"的间接建模方式，而是将空间 $(x, y, z)$、时间 $t$ 和视角方向 $\mathbf{d} \in \mathbb{R}^3$ 直接统一为一个 7D 高斯随机变量 $X = [X_p, X_t, X_d]^\top \sim \mathcal{N}(\boldsymbol{\mu}_{7D}, \boldsymbol{\Sigma}_{7D})$。与 4D-GS 的 K-Planes 分解式编码不同，7DGS 显式维护完整的 7D 协方差矩阵，其跨维度协方差块 $\boldsymbol{\Sigma}_{pt}$、$\boldsymbol{\Sigma}_{pd}$、$\boldsymbol{\Sigma}_{td}$ 直接捕获空间-时间、空间-方向、时间-方向之间的关联性——例如 $\boldsymbol{\Sigma}_{pd}$ 可以表达"从某个方向观察时该高斯的几何形状如何偏移"。7D 协方差通过 Cholesky 分解 $\boldsymbol{\Sigma}_{7D} = LL^\top$ 参数化（灵感来自 6DGS），保证优化过程中协方差矩阵始终正定。颜色表示方面，7DGS 沿用 3DGS 的球谐函数而不引入额外时间依赖，因为动态信息已编码在高斯参数本身中。

\textbf{条件切片：从 7D 到可渲染 3D。}4D-GS 用变形场网络将规范高斯"搬运"到当前时刻；7DGS 则采用纯解析的条件切片（conditional slicing）机制——无需网络前向推理。给定观测时间 $t$ 和视角方向 $\mathbf{d}$，利用多元高斯条件分布的性质将 7D 高斯直接投影为条件 3D 高斯。将 7D 协方差划分为空间块与跨域关联块，条件均值和协方差为：

\begin{equation}
\boldsymbol{\mu}_{\text{cond}} = \boldsymbol{\mu}_p + \boldsymbol{\Sigma}_{p,(t,d)}\,\boldsymbol{\Sigma}_{(t,d)}^{-1}\begin{bmatrix} t - \mu_t \\ \mathbf{d} - \boldsymbol{\mu}_d \end{bmatrix}, \quad
\boldsymbol{\Sigma}_{\text{cond}} = \boldsymbol{\Sigma}_p - \boldsymbol{\Sigma}_{p,(t,d)}\,\boldsymbol{\Sigma}_{(t,d)}^{-1}\,\boldsymbol{\Sigma}_{p,(t,d)}^\top,
\end{equation}

其中修正项 $\boldsymbol{\Sigma}_{p,(t,d)}\,\boldsymbol{\Sigma}_{(t,d)}^{-1}[t - \mu_t;\, \mathbf{d} - \boldsymbol{\mu}_d]$ 根据时间和方向的偏差动态调整空间均值，$\boldsymbol{\Sigma}_{\text{cond}}$ 则移除已被时间和方向解释的空间方差部分。条件不透明度通过时间和方向调制因子联合控制：

\begin{align}
\alpha_{\text{cond}} &= \alpha \cdot f_{\text{temp}} \cdot f_{\text{dir}}, \\
f_{\text{temp}} &= \exp\!\bigl(-\tfrac{1}{2}\gamma_t\,(t - \mu_t)^2\,\sigma_t^{-1}\bigr), \quad
f_{\text{dir}} = \exp\!\bigl(-\tfrac{1}{2}\gamma_d\,(\mathbf{d} - \boldsymbol{\mu}_d)^\top \boldsymbol{\Sigma}_d^{-1}(\mathbf{d} - \boldsymbol{\mu}_d)\bigr),
\end{align}

其中 $\gamma_t$、$\gamma_d$ 控制衰减灵敏度（如 D-NeRF 数据集设为 0.05，7DGS-PBR 数据集设为 0.5），在 15000 次迭代后由固定值变为可学习参数，让网络自主决定每个高斯对时间/方向偏差的容忍度。当观测时间或视角方向偏离高斯期望值时，该高斯对渲染的贡献自动衰减——这是一种高斯概率框架下的"软排斥"，无需显式的可见性判断。

\textbf{自适应高斯精化。}条件切片仅调整均值和不透明度而保持协方差不变，无法表达非刚性形变等复杂动态。与 4D-GS 用多头 MLP 从时空特征解码位置/旋转/缩放偏移的思路类似，7DGS 引入自适应高斯精化（Adaptive Gaussian Refinement），但作用于 7D 参数而非 3D 参数：构造综合特征向量 $\mathbf{f} = [\boldsymbol{\mu}_p; \mu_t; \boldsymbol{\mu}_d; \gamma(t)]$（其中 $\gamma(t)$ 为 $K{=}10$ 级多频率时间编码，捕获低频趋势到高频细节），通过多个两层 MLP（架构 $C_{\text{in}} \to 64 \to C_{\text{out}}$）预测残差修正 $\Delta\boldsymbol{\mu}_p$、$\Delta\mu_t$、$\Delta\boldsymbol{\mu}_d$、$\Delta\mathbf{l}$（$L$ 的下三角元素），使空间位置、时间坐标、方向均值和协方差都能随时间动态调整。精化在条件切片之前执行，确保后续切片产生的 3D 高斯能更准确地反映演化中的场景动态。

\textbf{优化与分裂策略。}7DGS 的渲染管线完全兼容 3DGS：条件切片将 7D 表示转化为标准 3D 高斯后，直接使用 3DGS 的透视投影和可微光栅化。在自适应密度控制中，除了标准的空间梯度分裂准则外，7DGS 利用时间-空间跨维度协方差 $\boldsymbol{\Sigma}_{pt}$ 的幅值作为额外分裂信号：当该相关超过屏幕范围的 0.05 且归一化时间尺度（从 $\sigma_t$ 导出）大于 0.25 时触发分裂，确保具有显著运动动态的区域被充分表示。此外，最小不透明度剪枝阈值提高到 $\alpha_{\min} = 0.01$（高于 3DGS 默认值），以补偿条件不透明度被时间和方向调制因子衰减后整体偏低的效应。

\textbf{效率与性能。}在自建的 7DGS-PBR 数据集上（含心脏 CT、体积云、半透明材质等含强视角依赖的动态场景），7DGS 在 PSNR 上以平均 +4.71 dB 的优势超过 4D-GS，同时仅使用 15.3\% 的高斯数量，渲染速度超过 400 FPS。这得益于 7D 表示的信息密度：每个高斯同时编码几何、动态和视角依赖，使得更少的高斯即可覆盖完整的时空-角度空间。

\subsubsection{高阶4D编码方案}

本路线介绍通过方向感知平面分解、张量混合和时空哈希等高阶编码方案来压缩4D表示的方法。

无论是变形场还是原生 4D 高斯，都需要在时空域中查询特征，而直接使用 4D 体素网格的存储开销为 $O(N^4)$，完全不可接受。本节讨论如何通过更高效的编码方案来压缩 4D 表示。

\paragraph{方向感知平面分解 4D 表示}

DaRePlane（2024.10）~\cite{2410_14169} 修改了 4D 平面分解的平面参数化方式——用方向感知小波表示替代无方向性的平面系数——4D 分解的整体架构（HexPlane 的三对平面乘法）保持不变。

\textbf{Baseline 做法。}当前 4D 表示的主流做法是将 4D 时空体 $D$ 分解为三对 2D 平面的乘积。以 HexPlane（2024.03）~\cite{2403_09742} 为例，baseline 将 4D 体表示为 $D = \sum_{r=1}^{R_1} M_r^{XY} M_r^{ZT} v_r^1 + \sum_{r=1}^{R_2} M_r^{XZ} M_r^{YT} v_r^2 + \sum_{r=1}^{R_3} M_r^{YZ} M_r^{XT} v_r^3$，其中每个 $M_r^{AB}$ 是一个可学习的 2D 平面（分辨率为 $m \times n$），$v_r^i \in \mathbb{R}^F$ 是沿特征轴的可学习向量，$R = R_1 + R_2 + R_3$ 为低秩分量数。这组平面在 4D-GS 框架中的作用是：给定时空点 $(x,y,z,t)$，从配对平面（如 $XY$ 和 $ZT$）中提取特征向量并相乘、拼接，再由多头 MLP 预测高斯的变形量（位置、旋转、缩放、不透明度）。该分解将内存复杂度从 $O(N^3TF)$ 降至 $O(RN^2F)$。然而，每个平面 $M_r^{AB}$ 本质上是一张无方向性的灰度图——所有方向的特征被混在一起编码，无法区分场景中的运动方向。

\textbf{问题与动机。}传统做法中，平面系数通常直接在空间域中优化，或通过标准离散小波变换（DWT）在频域中表示。但标准 2D DWT 存在两个根本缺陷：(1) \textbf{移位方差}（shift variance）——输入信号的微小空间偏移会显著改变小波系数的振荡模式，导致动态场景中运动物体周围出现鬼影伪影（ghosting artifacts）；(2) \textbf{方向模糊}（direction ambiguity）——2D DWT 仅产生三个子带系数（水平、垂直、对角），其中对角子带将 $+45^\circ$ 和 $-45^\circ$ 两个方向的信息混合为棋盘状混叠模式，无法有效区分不同方向的边缘和运动。直觉上，标准 DWT 就像只用三把刀（水平、垂直、对角）切图像，对角方向的两种信息被混在一起；理想情况下需要六把不同角度的刀，让每个方向的信息独立编码。

\textbf{实现原理。}DaRePlane 受双树复小波变换（Dual-Tree Complex Wavelet Transform, DTCWT）启发，提出方向感知表示（direction-aware representation）。DTCWT 的核心思想是使用两套独立的滤波器组并行处理信号：一对低/高通滤波器 $(h_0, h_1)$ 构造实部小波 $\psi_h$，另一对 $(g_0, g_1)$ 构造虚部小波 $\psi_g$，组合形成复小波 $\psi(x) = \psi_h(x) + j\psi_g(x)$。两套滤波器之间满足 Hilbert 变换关系（即 $\psi_g \approx \mathcal{H}(\psi_h)$），这保证了无论输入信号如何平移，复小波系数的幅值基本不变——这就是移位不变性的来源。在 2D 情况下，对图像的行和列分别施加复小波变换，产生 $(\psi_h(x)\psi_h(y))$、$(\psi_h(x)\psi_g(y))$、$(\psi_g(x)\psi_h(y))$、$(\psi_g(x)\psi_g(y))$ 四组系数，再利用 Hilbert 关系取复共轭，最终获得六组方向选择性系数——分别对应 $\pm 15^\circ$、$\pm 45^\circ$、$\pm 75^\circ$ 六个方向，而非 DWT 的三个（水平、垂直、单一对角）。

基于上述原理，DaRePlane 将每个 2D 平面 $M_r^{AB}$ 的参数化从直接的空间域系数改为 DTCWT 域系数。具体而言，引入 12 个可学习的小波系数——6 个实部 $R_{i,r}^{AB} \in \mathbb{R}^{m/2^l \times n/2^l}$ 和 6 个虚部 $I_{i,r}^{AB} \in \mathbb{R}^{m/2^l \times n/2^l}$（$i=1,\ldots,6$，$l$ 为 DTCWT 变换层级）——以及 1 个近似系数 $W_{a,r}^{AB} \in \mathbb{R}^{m/2^{l-1} \times n/2^{l-1}}$，通过逆 DTCWT 重建平面：

\begin{equation}
M_r^{AB} = \text{IDTCWT}\bigl([W_{a,r}^{AB}, \{R_{i,r}^{AB}\}_{i=1}^6, \{I_{i,r}^{AB}\}_{i=1}^6]\bigr).
\end{equation}

换言之，每个平面不再是一张无方向性的灰度图，而是由沿六个特定方向编码特征的各向异性小波系数组合而成，自然消除了棋盘效应并获得移位不变性。这些方向感知平面直接替换 HexPlane 公式中的原始平面 $M_r^{AB}$，整体 4D 分解架构不变。

\textbf{存储稀疏化与压缩。}方向感知表示引入了 12 个方向系数（vs. DWT 的 3 个），虽然换来了更好的方向区分能力，但也带来了约 $4\times$ 的参数量增加。为此，DaRePlane 引入可学习掩码进行稀疏化：为每组小波系数 $W^{AB} \in \{R_i^{AB}, I_i^{AB}, W_a^{AB}\}$ 配备一个掩码 $M^{AB}$，通过直通估计器（straight-through estimator）实现可微二值化：

\begin{equation}
\hat{W}^{AB} = \bigl(\text{sg}(\sigma(M^{AB})) - \sigma(M^{AB})\bigr) \cdot M^{AB} + \sigma(M^{AB}) \odot W^{AB},
\end{equation}

其中 $\sigma$ 为 Sigmoid 函数，$\text{sg}$ 为停止梯度算子，前向传播中使用 Heaviside 阶跃函数 $H(\cdot)$ 将掩码二值化（$0$ 或 $1$），反向传播中梯度通过 Sigmoid 近似传递。掩码稀疏性通过额外损失项 $\mathcal{L}_m$（所有掩码值之和）控制，权重 $\lambda_m$ 调节稀疏程度。稀疏化后，零值元素占大多数，方法采用行程编码（RLE）将二值掩码转为 8 位无符号整数流，再用 Huffman 编码对非零元素进行熵编码，最终实现紧凑存储。

\textbf{在 GS 框架中的应用。}DaRePlane 可同时应用于 NeRF 和 GS 框架（分别称为 DaReNeRF 和 DaReGS）。在 GS 框架中，给定高斯点的 4D 坐标 $(x,y,z,t)$，从方向感知平面中提取时空特征 $f$，再通过四个小型 MLP $F = \{F_\mu, F_R, F_S, F_o\}$ 分别预测位置、旋转、缩放和不透明度的变形量：

\begin{equation}
G_t = G_0 + \Delta G = (\mu + \Delta\mu,\; R + \Delta R,\; S + \Delta S,\; o + \Delta o),
\end{equation}

其中 $\Delta\mu = F_\mu(f)$，$\Delta R = F_R(f)$，$\Delta S = F_S(f)$，$\Delta o = F_o(f)$。总损失函数为 $\mathcal{L} = \lambda_{\text{photo}}\mathcal{L}_{\text{photo}} + \lambda_{\text{reg}}\mathcal{L}_{\text{reg}} + \lambda_m\mathcal{L}_m$，其中 $\mathcal{L}_{\text{photo}}$ 为颜色和深度的光度损失（使用 L1 范数和 Pearson 相关系数），$\mathcal{L}_{\text{reg}}$ 为全变分（TV）正则化损失，$\mathcal{L}_m$ 为掩码稀疏性损失。训练采用由粗到细策略（coarse-to-fine），分辨率逐步增长以加速收敛并提供隐式正则化。此外，方法还维护一个小型 3D 空体素（emptiness voxel），跨时间步聚合最大不透明度，跳过空白区域的采样以加速训练。

方向感知表示在复杂运动场景中显著优于标准平面分解，在手术场景（显微镜、内窥镜、腹腔镜）中也表现出色。值得一提的是，该方向感知表示同样适用于静态 3D 场景——采用 TensorRF 风格的分解（$D = \sum_r M_r^{AB} v_r^C$），将平面系数同样替换为 DTCWT 域表示即可。

\paragraph{4D 哈希编码与动静自监督分解}

\textbf{从平面分解到哈希编码的范式转移。}前述 DaRePlane 在平面分解（HexPlane）框架内改进平面参数化——用方向感知小波系数替代无方向性的空间域系数——但整体架构仍基于低秩假设：将 4D 时空体分解为三对 2D 平面的乘积，认为变形特征具有大量共性、可被因子化为紧凑表示。DASH（2025.07）~\cite{2507_19141} 指出，这一低秩假设在面对空间坐标大量重叠的高斯时会产生\textbf{特征重叠}（feature overlap）：多个高斯共享同一平面网格单元，各自的变形模式被混杂编码，直接限制变形预测的区分度和渲染质量。Grid4D 虽将 4D 编码拆分为四个独立的 3D 哈希编码以回避低秩假设，但各 3D 哈希表之间仍然存在特征重叠，且时间信息被分散编码而非统一建模。DASH 的根本思路是：放弃平面分解，改用\textbf{多分辨率 4D 哈希编码}直接在统一的时空框架中表征动态高斯的变形，同时通过\textbf{自监督动静分解}将 4D 哈希编码的范围限制在动态区域，从而兼顾表达力和效率。

\textbf{问题：为什么不能直接用全场景 4D 哈希编码？}多分辨率哈希编码（Instant-NGP）在静态 3D 场景中已证明高效，但直接扩展到 4D 会带来 $O(n^4)$ 的空间复杂度——时间维度的加入使哈希表的碰撞率急剧上升，全场景编码需要远更大的哈希表才能维持足够的特征分辨率。DASH 的核心洞察是：\textbf{大多数场景包含大面积静态区域，这些区域完全不需要存储在 4D 哈希表中}——静态高斯只需 3D 空间编码即可，仅真正运动的部分才需要 4D 时空编码。

\textbf{自监督动静分解。}DASH 提出一种无需额外标注（如手动掩码或预计算分割）的动静分解策略，利用静态和动态高斯在优化过程中的\textbf{固有行为差异}进行区分。具体分三步：

\begin{enumerate}
\item \textbf{运动感知特征提取。}方法使用 3D 哈希编码器 $G_{3D}$（Instant-NGP）和一个小型 MLP $\phi_p$ 提取每个高斯的空间特征 $f_{3D} = \phi_p(G_{3D}(\mathbf{p}))$。关键在于：在 4D 空间中，3D 哈希编码的各向同性采样假设不再成立（4D 采样通常是各向异性的（2024.03）~\cite{2403_19198}），因此当 3D 哈希编码被用于动态场景时，\textbf{空间特征会隐式捕获运动模式}——同一空间位置在不同时刻查到的哈希特征不同，这种时间驱动的特征分布变化自然编码了运动信息。随后，MLP $\phi_l$ 将空间特征解码为位置变形量 $\Delta\mathbf{p} = \phi_l(f_{3D})$。
\item \textbf{一阶线性运动建模。}方法将每个高斯的时变轨迹近似为一阶线性运动模型 $\mathbf{p}(t) = \mathbf{p} + t\Delta\mathbf{p}$。尽管真实运动通常是非线性的，但这一简化带来两个关键性质：(a) 对于\textbf{静态高斯} $G_s$，真实变形满足 $\Delta\mathbf{p}_s \approx 0$，线性模型能准确捕捉；(b) 对于\textbf{动态高斯} $G_d$，线性模型无法完全表达复杂轨迹，$\Delta\mathbf{p}_d$ 在优化过程中呈现\textbf{非收敛振荡}（non-convergent oscillation），这一行为特征天然区分了动静两类。变形幅值 $\|\Delta\mathbf{p}\|$ 直接对应高斯的平均速度，提供了物理上有意义的运动强度度量。
\item \textbf{自适应阈值分割。}计算所有高斯的运动幅值 $\|\Delta\mathbf{p}\|$，设置阈值 $\tau$ 为这些值的 top-$k\%$ 分位数，将高斯分为动态 $G_d$（$\|\Delta\mathbf{p}\| > \tau$）和静态 $G_s$（$\|\Delta\mathbf{p}\| \le \tau$）两组。静态高斯受约束损失推动变形归零：$\mathcal{L}_s = \frac{1}{N_s}\sum_{\mathbf{p} \in G_s}\|\Delta\mathbf{p}\|$。此外，方法对静态和动态高斯\textbf{分别}执行自适应密度控制（ADC），使用不同的配置参数，让两组高斯根据各自的运动特性独立优化空间分布。
\end{enumerate}

分解阶段的损失函数为 $\mathcal{L}_d = (1 - \lambda_c)\mathcal{L}_1 + \lambda_c\mathcal{L}_{\text{D-SSIM}} + \lambda_s\mathcal{L}_s$，其中前两项为标准的颜色和结构相似性损失，第三项为静态约束。

\textbf{4D 多分辨率哈希编码。}动静分解完成后，DASH 对动态高斯使用专门的 4D 多分辨率哈希编码 $G_{4D}$。与平面分解（将 4D 体拆为三对 2D 平面乘积）截然不同，4D 哈希编码\textbf{直接在统一的时空网格}中对动态高斯进行编码，无需低秩假设。具体而言，哈希编码在每一维度按几何级数设置多级分辨率 $N_l = \lfloor N_{\min} \cdot b^l\rfloor$（$l \in [0, L-1]$），其中 $b = \exp\bigl(\frac{\ln N_{\max} - \ln N_{\min}}{L-1}\bigr)$ 为各级之间的增长率。对于 4D 输入坐标 $\mathbf{x} = (\mathbf{p}, t)$，在每一级 $l$ 通过向下和向上取整计算超体素（hypervoxel）位置，再通过哈希函数 $h_l(\mathbf{x}_l) = \bigl(\bigoplus_{i=1}^{d} x_{l,i} \cdot \pi_i\bigr) \mod T_l$（其中 $\oplus$ 为逐位异或，$\pi_i$ 为大素数，$T_l$ 为哈希表大小）索引可学习特征向量，最后通过四线性插值获得编码特征 $f_h = G_{4D}(\mathbf{x}) \in \mathbb{R}^{L \times F}$，再经小型 MLP $\phi_d$ 得到最终 4D 特征 $f_{4D} = \phi_d(f_h)$。

\textbf{多头变形解码器。}4D 特征 $f_{4D}$ 被送入多头变形解码器 $D = \{\phi_x, \phi_r, \phi_s\}$，三个独立 MLP 分支分别预测位置、旋转和缩放的变形量：

\begin{equation}
(\mathbf{x}', r', s') = (\mathbf{x} + \Delta\mathbf{x},\; r + \Delta r,\; s + \Delta s),
\end{equation}

其中 $\Delta\mathbf{x} = \phi_x(f_{4D})$，$\Delta r = \phi_r(f_{4D})$，$\Delta s = \phi_s(f_{4D})$。变形后的动态高斯 $G_d'$ 与静态高斯 $G_s$ 合并，构成完整的场景表示。

\textbf{与 DaRePlane 的平滑策略对比。}DaRePlane 通过 DTCWT 域参数化获得隐式的方向选择性平滑——六方向小波系数天然抑制了棋盘伪影和移位方差。DASH 则采用了完全不同的路线：由于 4D 哈希编码是显式表示，缺乏 MLP 隐式表示的固有平滑性（哈希表的特征 lookup 在空间上是离散的），DASH 提出显式的\textbf{时空平滑正则化}（spatio-temporal smoothness regularization）——在特征空间中对输入坐标施加小随机扰动，强制邻域特征保持一致：

\begin{equation}
\mathcal{L}_r = \|G_{4D}(\mathbf{x}) - G_{4D}(\mathbf{x} + \delta_\mathbf{x})\|_2^2,
\end{equation}

其中 $\delta_\mathbf{x} = (\delta_\mathbf{p}, \delta_t)$ 为小随机扰动向量。该正则化在局部邻域建立时空一致性约束，确保相邻高斯在相邻时刻具有相似的变形模式，从而抑制混沌伪影。总损失为 $\mathcal{L} = (1 - \lambda_c)\mathcal{L}_1 + \lambda_c\mathcal{L}_{\text{D-SSIM}} + \lambda_r\mathcal{L}_r$。

\textbf{效率与性能。}DASH 的动静分解策略将 4D 哈希编码仅聚焦于动态高斯——动态高斯的数量通常远小于全场景高斯总数，这使 4D 哈希表的存储需求和碰撞率大幅降低，同时保持了快速的查询速度。方法在保持实时渲染速度（264 FPS）的同时，通过统一的时空编码避免了平面分解中的特征重叠问题。

\paragraph{4D 分解哈希编码与体素化变形}

回顾上文，DaRePlane 在 HexPlane 三对平面乘法的框架内优化了每个 2D 平面的参数化方式（用方向感知小波系数替代灰度系数），但保留了将 4D 时空编码分解为六对 2D 平面的基本结构。Grid4D（2024.10）~\cite{2410_20815} 则提出一个更根本的问题：平面分解本身的低秩假设是否合理？

\textbf{问题：平面分解的特征重叠。}Grid4D 指出，平面方法将 4D 坐标 $(x,y,z,t)$ 分解为六对 2D 平面（如 $(x,y)\times(z,t)$、$(y,z)\times(x,t)$ 等），暗含的假设是 4D 变形特征可以被低秩因子化。但实际场景中，空间坐标大量重叠的高斯极为常见。一个直观的例子：考虑两个仅 $x$ 坐标不同的高斯 A 和 B（共享 $y$、$z$），平面方法在 $(y,z)$、$(y,t)$、$(z,t)$ 三个平面上为它们产生完全相同的编码——整体 6 个平面中有 3 个重复，高达 50\% 的特征重叠率。这意味着即使 A 和 B 在时间轴上有完全不同的变形行为，它们在这些平面上的特征表示却无法区分。DaRePlane 的方向感知小波提升了每个平面内部的方向编码质量，但无法解决平面间特征重叠这一结构性问题——因为重叠的根源在于分解维度本身，而非平面系数的表示方式。

\textbf{动机与方案：从 2D 平面到 3D 网格。}直觉上，既然 2D 平面分解导致了过多的维度丢弃（每次只保留两个坐标），一个自然的方向是提高分解粒度——保留更多维度。直接使用完整 4D 超网格哈希编码可以彻底消除重叠，但其 $O(n^4)$ 空间复杂度导致极高的碰撞率和不可接受的存储开销。Grid4D 提出的 Tri-axial 4D 分解方案处于两者之间：将 4D 编码分解为四个 3D 哈希编码——一个空间编码 $G_{xyz}(x,y,z)$ 和三个时间编码 $G_{xyt}(x,y,t)$、$G_{yzt}(y,z,t)$、$G_{xzt}(x,z,t)$。与平面方法的区别是本质性的：每个子编码保留了三个维度的信息（而非两个），两个高斯需要同时共享三个坐标才会产生完全重叠，重叠率从 50\% 降至 25\%。空间复杂度从完整 4D 网格的 $O(n^4)$ 降至 $O(n^3)$，同时不依赖任何低秩假设——每个 3D 编码本身就是完整的低维特征，而非平面方法中两个因子的乘积。此外，考虑到 4D 空间中的采样通常是各向异性的（时间维度远比空间维度稀疏），时间编码在 $t$ 维度使用与空间维度不同的分辨率。具体而言，每个分辨率层级的网格大小按几何级数设置 $N_l = \lfloor N_{\min} \cdot b^l\rfloor$（$l \in [0, L{-}1]$），网格位置通过哈希函数 $h_l(\mathbf{x}_l) = (\bigoplus_{i=1}^{d} x_i \Pi_l^{x_i}) \mod T_l$ 映射到哈希表，再通过三线性插值获取特征。

\textbf{方向注意力聚合。}Grid4D 的编码器产生两类特征：空间特征 $G_{xyz}$（跨时间线的公共信息）和时间特征 $G_{xyt}, G_{yzt}, G_{xzt}$（时间戳相关的动态信息）。一个关键观察是：不同场景组件往往具有截然不同的变形模式。例如，阴影的运动方向通常与投射它的物体相反——如果仅用一个标量权重将空间与时间特征组合，就无法表达这种方向对立。为此，Grid4D 设计方向注意力（directional attention）模块：空间特征经小型 MLP $f_s$ 生成注意力分数 $a = 2\sigma(h_{xyz}) - 1$，其中 $h_{xyz} = f_s(G_{xyz}(x,y,z))$。与传统注意力范围 $(0, 1)$ 不同，Grid4D 将分数缩放到有向范围 $(-1, 1)$——正值表示空间特征增强时间变形的方向，负值表示抑制甚至反转该方向，从而让具有相反变形趋势的邻近高斯（如物体与阴影）获得有区分性的特征表示。最终变形特征通过注意力加权的时间特征得到：$h = a \odot f_t(G_{xyt}, G_{yzt}, G_{xzt})$。多头变形解码器 $D$ 将 $h$ 解码为位置变形（以四元数旋转矩阵 $R_x$ 和平移向量 $T_x$ 表示，四元数保证更精确的插值和稳定的优化）、缩放偏移和旋转偏移：$\boldsymbol{\mu}' = R_x\boldsymbol{\mu} + T_x$，$S' = S + \Delta s$，$R' = R + \Delta r$。

\textbf{平滑正则化。}与 DaRePlane 通过小波系数的方向感知隐式获得平滑性不同，Grid4D 的 3D 哈希编码作为显式表示同样缺乏固有平滑性（不同于 MLP 的隐式平滑）。方法提出特征空间平滑正则化：$\mathcal{L}_r = \|G_{xyzt}(x,y,z,t) - G_{xyzt}(x{+}\delta_x, y{+}\delta_y, z{+}\delta_z, t{+}\delta_t)\|_2^2$，其中 $(\delta_x, \delta_y, \delta_z, \delta_t)$ 为小随机扰动，$G_{xyzt}$ 为四个网格哈希编码的拼接。正则化直接作用于特征空间而非 MLP 输出，强制邻域特征相似以抑制混乱变形预测。空间和时间编码使用不同的正则化设置以适应各自特性。为提高效率，每次迭代仅随机选择部分高斯进行正则化，而非全部。总损失为 $\mathcal{L} = (1 - \lambda_c)\mathcal{L}_1 + \lambda_c\mathcal{L}_{\text{D-SSIM}} + \lambda_r\mathcal{L}_r$。此外，Grid4D 使用分离的高斯位置（detached positions）进行变形预测以获得更好的性能，初始化时使用 SfM 点云构建静态规范场，并复用 3DGS 的自适应密度控制和不透明度重置机制。Grid4D 在高保真动态场景渲染中表现出色，特别是在运动幅度大、物体形变复杂的场景中显著优于平面方法。

\paragraph{自适应分治 4D 高斯管线}

前述 DASH 和 Grid4D 都采用了动静分离策略来缓解 4D 哈希编码的碰撞与存储问题，但它们的动静分类方式各不相同。DASH 通过一阶线性运动模型的位置变形量做阈值分割（top-$k\%$ 分位数），Grid4D 则完全不区分动静而依赖分解结构本身缓解碰撞。Swift4D（2025.03）~\cite{2503_12307} 沿着动静分离的思路进一步推进，但其核心贡献不在于分离策略本身，而在于提出了一种基于渲染的动静分类机制——将分类问题嵌入 3DGS 的 $\alpha$-混合渲染管线中，利用多视角的自然遮挡关系自动纠正误判，使得分类过程既快速（约 1 分钟）又鲁棒。

\textbf{渲染式动静分类。}Swift4D 为每个高斯点引入一个可学习的动态参数 $d$（初始化为 0），并通过 3DGS 标准的 $\alpha$-混合将其投影到 2D：

\begin{equation}
D(\mathbf{x}) = \sigma\biggl(\sum_i d_i \alpha_i(\mathbf{x})\prod_{j=1}^{i-1}(1-\alpha_j(\mathbf{x}))\biggr),
\end{equation}

其中 $\sigma$ 为 Sigmoid 函数，将动态值映射到 $(0,1)$。这个设计的直觉在于：如果一个高斯覆盖了动态像素，反向传播会赋予其正梯度（推动 $d$ 增大）；如果覆盖了静态像素，则赋予负梯度（推动 $d$ 减小）。由于 $\alpha$-混合本身包含深度排序和透明度累积，梯度会随着遮挡关系和距离自然衰减——远处的误判高斯不会获得强梯度。监督信号来自训练视频的 2D 动静像素掩码，通过对每个像素计算时间标准差 $S(\mathbf{x}) = \sqrt{\frac{1}{T}\sum_{t=1}^{T}(C(\mathbf{x},t) - \bar{C}(\mathbf{x}))^2}$ 并以阈值 $\tau = 0.02$ 二值化得到。整个优化通过二值交叉熵损失 $\mathcal{L}_d = \mathbb{E}_\mathbf{x}[-\hat{D}(\mathbf{x})\log D(\mathbf{x}) - (1-\hat{D}(\mathbf{x}))\log(1-D(\mathbf{x}))]$ 驱动，将动静分类转化为可微的二元分类问题。这种方法的一个关键优势是多视角一致性：即使某个视角下静态高斯被动态像素错误覆盖（因而获得正向梯度），其他视角中同一高斯覆盖静态像素时会产生更强的负向梯度，最终纠正分类——这本质上是利用了多视角之间的互补遮挡信息。优化完成后（约 1 分钟），$d > \tau_d = 7.0$ 的高斯被分类为动态，其余为静态。

\textbf{4D 哈希编码与变形解码。}动静分离后，Swift4D 仅对动态高斯施加 4D 多分辨率哈希编码 $h_{4d}$。给定 4D 坐标 $(\boldsymbol{\mu}, t)$，在每个分辨率层级 $l$ 通过线性插值获取特征 $h_{4d}(\boldsymbol{\mu}, t; l) \in \mathbb{R}^F$，拼接后经小型 MLP $\phi_d$ 融合为 $f_d = \phi_d([h_{4d}(\boldsymbol{\mu},t;0), \ldots, h_{4d}(\boldsymbol{\mu},t;L-1)])$。与 DASH 的三点差异值得注意：(1) Swift4D 的哈希编码在完整的 4D 空间 $(x,y,z,t)$ 上操作，而非 DASH 那样将 3D 和 4D 哈希分离处理；(2) 动静分离后动态高斯数量远小于全场景，使 4D 哈希表保持紧凑（可小至 8 MB），无需 Grid4D 的三轴分解来降低复杂度；(3) 多头变形解码器 $D = \{\phi_\mu, \phi_s, \phi_q, \phi_o, \phi_{sh}\}$ 同时预测位置、缩放、旋转、不透明度和球谐系数的变形量（5 个属性），比 DASH 的 3 属性解码器（位置、旋转、缩放）更完整。变形后参数为 $\boldsymbol{\mu}' = \boldsymbol{\mu} + \Delta\boldsymbol{\mu}$ 等，静态高斯则直接与变形后的动态高斯合并渲染。

\textbf{时间重要性剪枝。}Swift4D 进一步引入时间重要性剪枝（Temporal Importance Pruning）来消除浮动伪影。原始 3DGS 的剪枝策略基于不透明度，直接沿用到动态场景中会导致规范空间（canonical space）与变形空间的过度耦合。为此，Swift4D 在每个时间戳 $t$ 下计算每个高斯对全部训练视角的渲染贡献：

\begin{equation}
w_i = \sum_{I,\mathbf{x}}\sum_{t \in T}\biggl(\alpha_i(\mathbf{x}|t)\prod_{j=1}^{i-1}(1-\alpha_j(\mathbf{x}|t))\biggr),
\end{equation}

即跨所有时间步和训练视角累积 $\alpha$ 权重。剪除 $w_i < 0.02$ 的高斯可有效消除那些悬浮在空中、从未被任何训练视角观测到的浮动伪影。子高斯在克隆和分裂时继承父高斯的动态属性 $d$。

\textbf{训练流程。}Swift4D 的训练分为三个阶段：(1) 用首帧图像通过 SfM 初始化点云，训练 5000 次迭代建立稳定的规范空间；(2) 冻结规范空间，用上述渲染式分类优化动态参数 $d$（约 1 分钟）；(3) 冻结 $d$，联合优化高斯属性与时空结构。损失函数保持简洁 $\mathcal{L}_{\text{rec}} = (1 - \lambda_1)\mathcal{L}_1 + \lambda_1\mathcal{L}_{\text{SSIM}}$，无需 DASH 的时空平滑正则化或 DaRePlane 的掩码稀疏化等额外正则项。Swift4D 在训练速度、渲染速度和存储效率上均达到 SOTA 水平。

\paragraph{4D 神经体素网格与变形查询}

前述 DaRePlane、DASH、Grid4D、Swift4D 等方法虽然改进了 4D 时空特征的编码方式（方向感知小波、动静分解哈希、三轴哈希、渲染式动静分类），但共享一个共同前提：时间维度的建模通过变形场施加在已有的高斯点上，随时间轴扩展时高斯数量线性增长，内存以 $O(NT)$ 增长（$N$ 为高斯数量，$T$ 为时间戳数）。对于边缘设备（如 8 GB 内存的 NVIDIA Jetson）或长序列场景，这一瓶颈尤为突出。4D Neural Voxel Splatting (4D-NVS)（2025.11）~\cite{2511_00560} 从根本上改变了这一前提——不再为每个时间戳维护独立的高斯云，而是维护一组持久的体素锚点（persistent voxel anchors），渲染时按需从中生成神经高斯，再通过统一的变形网络施加时间形变。这一设计使内存复杂度降至 $O(fV + F)$（$V$ 为体素数量、$f$ 为每个体素的特征维度、$F$ 为变形网络参数量），与时间戳数量 $T$ 完全解耦。

\textbf{统一 4D 体素架构。}4D-NVS 以 Scaffold-GS 的体素生成机制为基础，将其从 3D 扩展到 4D。具体而言，初始阶段从 SfM 点云创建体素中心 $\mathbf{V} = \{\lfloor \mathbf{P}/\epsilon \rceil\} \cdot \epsilon$，每个体素维护 32 维特征 $f_v$、缩放因子 $l_v$ 和 $k$ 个可学习偏移 $O_v \in \mathbb{R}^{k \times 3}$。渲染时，首先通过视锥体裁剪（visibility culling）将计算从 $O(N)$ 降至 $O(V_{\text{visible}})$，再对每个可见体素按需生成 $k$ 个神经高斯，位置为 $\mu_i = x_v + O_i \cdot l_v$，属性（不透明度、颜色、缩放、旋转）由专用 MLP 从体素特征、视距和视角方向解码。与传统方法的关键区别在于：这些体素锚点在整个时间线上持久存在，不随时间戳数量增加而膨胀。

\textbf{HexPlane 变形与选择性变形策略。}与前述 DaRePlane 修改 HexPlane 的平面参数化不同，4D-NVS 直接使用标准 HexPlane 分解——将 4D 时空体因子化为六个 2D 平面（XY、XZ、YZ、XT、YT、ZT），通过双线性插值提取多分辨率时间特征 $f_h$，再由浅层 MLP $\phi_d$ 聚合为变形特征 $f_d$。然而，4D-NVS 在变形范围上做出了一个关键的设计选择。前述方法（如 DaRePlane 的 $G_t = G_0 + \Delta G$）对高斯的所有属性施加变形——位置、旋转、缩放、颜色和不透明度。4D-NVS 通过大量实验发现：外观属性（颜色、不透明度）对变形误差极为敏感，反向传播中微小的变形误差会沿颜色通道级联放大，导致训练不稳定。因此方法引入选择性变形（selective deformation）——仅对几何属性施加时间变形：

\begin{equation}
\mu' = \mu + \varphi_x(f_d), \quad r' = r + \varphi_r(f_d), \quad s' = s + \varphi_s(f_d),
\end{equation}

其中 $\varphi_x$、$\varphi_r$、$\varphi_s$ 为单层 MLP，分别预测位置偏移、四元数旋转和缩放因子。颜色和不透明度保持初始值不变，切断了误差沿外观通道的级联传播，显著改善了收敛性和渲染质量。

\textbf{视角自适应精化。}训练过程中，作者观察到某些视角持续表现不佳——由大幅变形、复杂遮挡或快速时间变化导致。与 DaRePlane 的由粗到细策略（全局分辨率逐步增长）不同，4D-NVS 提出局部化的视角自适应精化（view-adaptive refinement）：方法通过双重检测机制识别粗差视角（crude views）——(1) 基于 PSNR 的检测：跟踪每个渲染视角的 PSNR，与指数加权移动平均（EMA）比较，识别统计异常值 $\text{PSNR}_i < (1{+}\gamma) \cdot \text{EMA}_{\text{PSNR}}$；(2) 基于梯度的检测：跟踪视图空间点梯度幅值，识别收敛困难的视角 $\|\nabla \mathcal{L}_{\text{view}}\| > (1{+}\gamma) \cdot \text{EMA}_{\text{grad}}$。检测阈值 $\gamma$ 从 0.05 逐步衰减至 0.02，随训练推进变得更具选择性。识别出的粗差视角进入专门的精化阶段（14K 次迭代），使用更激进的致密化参数（梯度阈值从 0.0002 降至 0.0001，不透明度阈值从 0.05 降至 0.03）仅在问题区域创建额外高斯，不干扰已良好重建的区域。这种靶向策略在粗差视角上实现 0.5--1.2 dB 的 PSNR 提升，同时保持全局渲染效率。

\textbf{训练流程与损失函数。}4D-NVS 的训练分为三个阶段：(1) 粗初始化——使用混合时间戳样本通过 Scaffold-GS 训练 3K 次迭代，建立基础几何表示，不建模时间变形；(2) 精细时间训练——启用 HexPlane 变形模块，学习完整的时空动态（14K 次迭代），同时通过上述双重检测机制识别粗差视角；(3) 视角精化——仅对粗差视角进行靶向优化（14K 次迭代）。总损失函数为 $\mathcal{L} = \mathcal{L}_{\text{color}} + \lambda_{tv}\mathcal{L}_{tv} + \lambda_{vol}\mathcal{L}_{vol}$，其中 $\mathcal{L}_{\text{color}} = \mathcal{L}_1 + \lambda_{\text{SSIM}}\mathcal{L}_{\text{SSIM}}$ 保证像素精度和结构一致性，$\mathcal{L}_{tv}$ 为 HexPlane 的全变分正则化损失（空间平滑性），$\mathcal{L}_{vol} = \sum_{i=1}^{N_{ng}} s_{x,i} \cdot s_{y,i} \cdot s_{z,i}$ 为体积正则化损失，防止高斯过度膨胀（消融实验表明去掉该项会出现漂浮在相机前方的巨型高斯导致画面模糊）。4D-NVS 在 HyperNeRF 数据集上仅需 3,050 MiB 内存和 13 分钟训练时间即达到 28.5 dB PSNR（相比 TiNeuVox-B 的 21,558 MiB 和 4D-GS 的 4,500 MiB），内存缩减达 86\%，同时保持 44 FPS 的实时渲染速度。

\subsubsection{变形场的正则化与优化策略}

本路线介绍针对变形场训练中过拟合、变形不连续和高斯增殖等优化难题的正则化与训练策略。

前面介绍了变形场和 4D 表示的基本方法。然而，将这些方法训练到实际可用的质量并不容易——变形场容易过拟合训练帧、变形不连续导致闪烁、高斯无控制增殖导致存储爆炸。本节聚焦于这些训练过程中的优化难题及其解决方案。

多项研究提出了针对性的正则化和优化策略。

\paragraph{变形非一致性诊断：过拟合的真正来源}

变形场方法存在一个普遍却鲜被深究的问题：训练视角 PSNR 超过 40 dB，但测试视角骤降 6--11 dB。Incoherent Deformation, Not Capacity（2026.04）~\cite{2604_16747} 通过 D-NeRF 基准上 260 余次训练的系统消融，揭示了过拟合的真正根源——不是参数容量过大，而是变形场在相邻高斯之间的不一致性（incoherent deformation）。这一诊断颠覆了"更多高斯 = 更严重过拟合"的朴素直觉，并催生了以平滑先验替代容量削减的新策略。

\textbf{Baseline 做法与过拟合假说。}变形场方法的典型训练流程包含自适应密度控制（ADC）：每 100 次迭代从第 500 次到第 15K 次，对视图空间梯度超过阈值 $\tau_0 = 2 \times 10^{-4}$ 的高斯执行分裂（split，大高斯一分为二）或克隆（clone，小高斯添加偏移副本），并剪枝低不透明度高斯。过拟合的直接表现是训练-测试 PSNR 差距（gap）：在 D-NeRF 基准上，基线方法的平均差距为 $6.18 \pm 1.15$ dB。直觉上，高斯数量越多，模型越容易记忆训练视角的逐像素残差而非学习泛化的场景结构——这构成了"容量假说"。

\textbf{消融实验：split 是过拟合的瓶颈但不是原因。}作者逐个关闭 ADC 的子操作，发现关闭 split 操作可将 PSNR 差距从 6.18 dB 缩小到 1.15 dB（缩小 81.4\%），但同时测试 PSNR 从 34.11 dB 暴跌到 24.18 dB（下降 9.93 dB）。换句话说，没有 split，模型无法重建场景，也无法过拟合——split 创造了后续过程将其转化为记忆的工作容量。其他操作的影响微乎其微：关闭 prune 仅改变 2\% 的高斯数量和 1\% 的差距；将 split 频率减半仅缩小 9\% 的差距；提前在 7,500 次迭代终止 ADC 仅缩小 2\%——因为 84--89\% 的高斯增长发生在 7,500 次迭代之前。真正有意义的调度参数是阈值大小：$\tau_0 \times 2$ 可减少 63\% 的高斯数量和 26\% 的差距（代价仅 $-0.76$ dB）；$\tau_0 \times 0.5$ 则使高斯数量增至 126K 并扩大差距 23\%。

\textbf{容量假说的证伪。}上述消融中一个引人注目的模式是：PSNR 差距与最终高斯数量 $K$ 呈近似对数线性关系——在 9 种消融条件下，Pearson 相关系数 $r(\log K, \text{gap}) = 0.995$（bootstrap 95\% CI $[0.993, 1.000]$），Spearman 秩相关 $\rho = 1.000$。表面上看，"更多高斯 = 更严重过拟合"似乎成立。但作者引入了一个简单的局部平滑先验 EER（Elastic Energy Regularization）后，这一假说被直接证伪。

\textbf{EER：局部变形平滑正则化。}直觉上，EER 要求相邻高斯在变形后仍然保持相邻——就像一根橡皮筋被拉伸时，相邻点的位移应该连续变化。具体地，在训练的每个步骤中，随机采样一个小批量 $G$ 的高斯，对每个高斯 $i$ 在规范空间（变形前）中找到其 $k$ 个最近邻 $N_k(i)$，施加以下惩罚：

\begin{equation}
L_{\text{EER}} = \frac{\lambda_{\text{EER}}}{|G| \cdot k} \sum_{i \in G} \sum_{j \in N_k(i)} \frac{\|u(\mathbf{x}_i, t) - u(\mathbf{x}_j, t)\|^2}{\|\mathbf{x}_i - \mathbf{x}_j\|^2 + \epsilon}
\end{equation}

其中 $u(\mathbf{x}, t)$ 是时间 $t$ 处的变形位移，分母中的 $\|\mathbf{x}_i - \mathbf{x}_j\|^2$ 是规范空间中的邻域距离，用于归一化——惩罚的是相对变形（应变），而非绝对运动幅度。换言之，该公式鼓励变形场在局部满足 Lipschitz 连续性：一个高斯的变形不应与它的邻居显著不同。实现上，$k = 8$，邻域图每 500 次迭代重建一次（因为规范位置在训练中缓慢漂移），每步采样 2,048 个高斯以保证计算效率。EER 仅在精细阶段激活，权重 $\lambda_{\text{EER}}$ 从 0 按半余弦曲线在 3K--10K 迭代间升至目标值，避免干扰早期结构学习。

\textbf{EER 的效果：更少过拟合，更多高斯。}在 $\lambda_{\text{EER}} = 0.05$ 时，EER 将 PSNR 差距从 6.18 dB 缩小到 3.66 dB（缩小 40.8\%），代价仅 $-0.49$ dB 测试 PSNR。但关键的反直觉发现是：最终高斯数量从基线的 44,516 增至 82,498——增加了 85\%。在 $\lambda_{\text{EER}} = 0.1$ 时，差距进一步降至 3.45 dB（缩小 44.2\%），高斯数量为 86K。即使最温和的 $\lambda_{\text{EER}} = 0.01$，也能缩小 30.1\% 的差距（仅 $-0.19$ dB 代价），优于所有非 EER 正则化方案。这些数据点远偏离 $\log K$--gap 的对数线性回归线——容量假说无法解释 EER。

\textbf{为什么高斯数量反而增加？}直观解释是：在基线训练中，一个高斯可以在不同帧中"游走"到不同位置来逐帧记忆训练残差——EER 禁止了这种"游走"行为。于是，那些区域的重建损失保持较高，视图空间梯度持续较大，ADC 的 split 判据更频繁地触发。EER 本质上用"更多高斯"替代了"每个高斯更大的变形自由度"——用种群数量换取每个个体的变形约束。这也解释了为什么需要在 EER 基础上引入 GAD 来控制高斯增殖。

\textbf{直接机制验证。}为确认 EER 确实通过平滑变形场（而非其他隐含机制）起作用，作者在每个训练好的检查点上直接测量 per-Gaussian $k$-NN 应变 $\varepsilon_i = \frac{1}{k}\sum_j \|u_i - u_j\|^2 / \|\mathbf{x}_i - \mathbf{x}_j\|^2$。在 8/8 场景上一致地，EER 将应变均值降低了 99.72\%（中位数 99.80\%，最小 99.58\%，最大 99.90\%）。更关键的是，在测试相机对应的、EER 训练过程中从未采样过的时间步上，应变降低幅度为 99.73\%（std 0.11\%）——与固定时间步上的测量几乎相同。这说明平滑性是结构性的，而非优化器的采样偏差。分布分析进一步显示：EER 下每个场景的中位数高斯的应变，低于基线下第 1 百分位（最好的）高斯的应变；EER 的 p99 应变低于基线的中位数应变，比率从 4.6 到 26 倍不等。变形场从基线时的混沌状态（重尾应变分布）变为 EER 下的局部平滑状态（集中在 $< 10^{-2}$ 的紧密分布）。

\textbf{结论：不一致变形是过拟合的根源。}综合以上证据，过拟合的真正机制是变形场的 incoherent deformation：相邻高斯在未见过的时间帧上变形方向不协调，各自独立记忆训练视角的逐像素残差，导致新视角渲染崩溃。高斯数量与过拟合的对数线性关系是表象，底层原因是 ADC 在高应变区域产生更多高斯，更多高斯又有更大的总体变形自由度——形成正反馈环路。

\textbf{GAD：损失率感知的自适应致密化阈值。}基于上述诊断，作者提出 GAD（Growth-Aware Densification），在 EER 控制应变的同时控制高斯增殖。标准 ADC 使用全局固定阈值 $\tau_0$，这在训练早期合理（一切均欠表示），但在训练后期——损失已平台化时——每个新高斯主要在记忆噪声残差。GAD 将阈值改为自适应的：

\begin{equation}
\tau_{\text{GAD}}(t) = \tau_0 \left(1 + \frac{K(t) \cdot \beta}{N \cdot \sigma_{\text{ema}}(t)}\right)
\end{equation}

其中 $K(t)$ 是当前高斯数量，$N$ 是训练像素数，$\sigma_{\text{ema}}(t)$ 是每次迭代损失改善量的指数移动平均（EMA，衰减率 $\alpha = 0.99$），$\beta$ 是唯一的超参数。直觉上：当高斯数量大（$K$ 大）且损失已平台化（$\sigma_{\text{ema}}$ 小）时，阈值大幅上升——此时触发 split 的梯度更有可能是噪声而非信号。$\beta$ 需要根据宿主代码库的损失量级调节（因为 $K/(N\sigma)$ 的量纲为 $[\sigma]^{-1}$），但作者在 D-NeRF 上对 $\beta \in \{10^{-6}, 10^{-5}, 10^{-4}\}$ 做了扫描。

\textbf{PTDrop：抖动加权的高斯丢弃。}PTDrop 是 DropGaussian 在动态场景上的适配版本，同样无需架构修改。它做了两处关键改动：（1）余弦调度丢弃率——每次前向传播中，随机排除一部分高斯不参与渲染，全局丢弃概率从 0 按余弦曲线在 5K--12K 迭代间升至 0.3，允许训练在正则化开始前先稳定下来；（2）抖动加权选择——每个高斯的实际丢弃概率按其变形轨迹在时间步上的方差缩放。变形轨迹越不一致的高斯（即 Finding~2 所识别的 incoherent deformation 行为）被丢弃得越频繁，而变形轨迹平滑的高斯则基本保留。

\textbf{推荐配置与总结。}GAD + EER 在 D-NeRF 上将 PSNR 差距缩小 48.2\%，PTDrop 在此基础上额外贡献约 10 个百分点的差距缩减，且与 GAD、EER 干净组合。三者的协同逻辑是：EER 约束变形平滑性（消除过拟合的根源），GAD 控制高斯增殖（吸收 EER 带来的数量膨胀），PTDrop 随机丢弃不一致高斯（提供额外的正则化压力）。

\paragraph{时间自由度高斯与时变属性建模}

前述 Incoherent Deformation, Not Capacity 从过拟合的诊断出发，揭示了变形场一致性问题。FreeTimeGS++（2026.05）~\cite{2605_03337} 则采取了正交的路线：不从单一问题入手，而是对整个 4DGS 流水线做系统性拆解，回答一个更根本的问题——不同 4DGS 方法在结构上差异显著（变形场 vs.\ 4D 高斯基元、不同运动表示、不同生命周期控制），为何报告了相似的重建指标？这些高分到底来自结构性的表示优势，还是来自利用数据集特性的启发式技巧？

\textbf{实验方法论。}作者选择模块化且性能领先的 FreeTimeGS 作为代表性基线（因其公开训练代码不可用，作者通过技术通信对其进行了忠实复现，所有结果在六次独立运行上取平均），将 4DGS 框架分解为表示（representation）、初始化（initialization）、训练（training）和评估（evaluation）四个基本组件，逐个隔离其影响。与 Incoherent Deformation, Not Capacity 的消融策略不同——后者聚焦变形场一致性——FreeTimeGS++ 的消融覆盖整个流水线，旨在区分"真正驱动性能的设计选择"与"表面上的架构差异"。

\textbf{五个隐藏秘密。}通过上述系统性消融，作者揭示了五个此前未被充分阐明的发现：

\textbf{秘密 1（时间划分的涌现性）：}即使 4DGS 的公式中未显式区分静态与动态高斯，优化过程会自发涌现出功能性的时间划分——部分高斯退化为"永久存在"的静态背景，另一部分则局限在短时间窗口内作为瞬态基元。这种涌现行为未被先前文献充分讨论，却对理解模型的实际工作机制至关重要。

\textbf{秘密 2（运动噪声）：}逐高斯显式速度参数在优化后期产生严重的运动噪声——相邻高斯的运动方向不一致，导致物理上不合理的渲染结果。这与前述 Incoherent Deformation, Not Capacity 对不一致变形的诊断遥相呼应，但 FreeTimeGS++ 的视角是"表示层面"的：问题根源在于每个高斯独立拥有速度参数，缺乏共享的运动先验。

\textbf{秘密 3（初始化的决定性影响）：}初始点云的密度和质量对最终性能有决定性影响。FreeTimeGS 原文提到使用 RoMa 进行 2D 匹配，但匹配策略（如关键帧步长 keyframe stride 的选取：DyNeRF 设为 10、SelfCap 设为 1）、几何有效性检查（chirality 和坐标范围约束）等细节均未充分说明——这些看似琐碎的实现选择实际上决定了性能的上限。

\textbf{秘密 4（MCMC 密度控制的隐含运动效应）：}FreeTimeGS 采用 MCMC 启发的克隆策略进行密度控制，但该策略在重定位高斯时未重新初始化运动参数——新克隆的高斯继承了源高斯的速度和时序参数。这种隐含的参数继承对动态建模有不可忽视的影响，却未被先前工作讨论。

\textbf{秘密 5（可重复性危机）：}评估结果对数据集和评估帧的选择高度敏感。不同方法可能在不同帧子集上报告结果，导致看似合理的横向比较实际并不可比。六次独立运行的方差分析表明，部分配置下单次运行的 PSNR 波动可达 $0.5$ dB 量级。

\textbf{FreeTimeGS++ 的改进方案。}基于上述五个秘密的诊断，作者提出了一系列针对性的改进——每个改进直接对应一个或多个秘密：

\textbf{门控边缘化（Gated Marginalization，针对秘密 1）：}为每个高斯 $i$ 引入可学习的持久性门控 $g_i = \sigma(\gamma \cdot \hat{g}_i)$（$\gamma = 20$ 为锐度因子），其控制时间权重函数 $\alpha_i(t) = g_i + (1 - g_i)\exp(-\frac{1}{2}(\frac{t - \mu_i}{s_i/6})^2)$ 在永久（$g_i \to 1$，静态背景）和瞬态（$g_i \to 0$，短生命周期动态物体）之间平滑切换。这将秘密 1 中观察到的涌现行为形式化为可微的显式机制。

\textbf{神经速度场（Neural Velocity Field，NVF，针对秘密 2）：}取代逐高斯显式速度参数，改为共享的连续速度场 $v(\mathbf{x}, t) = f_\theta(\phi_x(\mathbf{x}), \psi_t(t))$，其中 $\phi_x$ 为多级哈希网格空间编码，$\psi_t$ 为时间编码，$f_\theta$ 为 HyperMLP。关键差异在于：所有高斯共享同一速度场，相邻空间位置的速度自然连续——从表示层面消除了运动噪声的根源。训练早期使用 warm-start 蒸馏损失 $L_{\text{distill}}$ 将初始化速度蒸馏到 NVF 中，桥接离散初始化与连续建模。

\textbf{UFM 引导的时空初始化（针对秘密 3）：}引入预训练的 Unified Flow \& Matching 模型（UFM）提供运动感知的几何先验。将初始 3D 点投影到图像平面，采样多视角光流并反投影得到 UFM 驱动的初始速度 $v_i^{\text{ufm}}$，再与 KNN 速度融合：$v_i^0 = m_i v_i^{\text{ufm}} + (1 - m_i) v_i^{\text{knn}}$，其中 $m_i \in \{0, 1\}$ 是基于 UFM 置信度和多视角一致性的有效性掩码。这使得高斯从训练开始就被赋予物理意义的运动向量。

\textbf{MCMC 密度控制（针对秘密 4）：}形式化了 FreeTimeGS 中未明确的重定位策略——重定位的高斯继承目标活跃高斯的所有参数，但不透明度和尺度按 3DGS-MCMC 的更新规则重新计算，保证优化稳定性。

\textbf{仿射颜色校正（针对秘密 5）：}集成仿射颜色校正（CC）模块吸收全局光照漂移，最大化训练的可重复性。总损失为 $L = \lambda_1 L_1 + \lambda_s L_{\text{ssim}} + \lambda_l L_{\text{lpips}} + L_{\text{reg}} + L_{\text{gate}} + L_{\text{cc}} + L_{\text{distill}}$。

\textbf{与前述锚点方法的互补关系。}Incoherent Deformation, Not Capacity 通过 EER 在优化层面约束变形一致性（局部平滑正则化），而 FreeTimeGS++ 的 NVF 从表示层面消除了不一致运动的根源（共享连续速度场）。两者分别从"正则化已有表示"和"改用更合理的表示"两个角度解决同一类问题——前者轻量且可即插即用，后者需要架构修改但提供更强的归纳偏置。在实验效果上，FreeTimeGS++ 在 DyNeRF 上取得 0.89 dB 的 PSNR 提升，在 SelfCap 上取得 0.67 dB 的提升，同时显著提高了六次运行的可重复性和物理一致的运动场。

\paragraph{高斯生命周期调制与时间锐度控制}

上文 Incoherent Deformation, Not Capacity 从正则化角度（EER 平滑约束 + GAD 致密化控制）抑制变形不一致性，但并未改变 4D 高斯基元本身的运动和不透明度建模方式。SharpTimeGS（2026.02）~\cite{2602_02989} 采取了根本不同的策略：不依赖事后正则化来约束变形，而是在表示层面引入逐基元的生命期（lifespan）参数，使运动幅度和时间可见性自动适应基元的"存续时长"——静态基元自然冻结，动态短寿命基元自然拥有大运动自由度。

\textbf{核心问题与设计动机。}现有 4DGS 方法用高斯曲线建模时间不透明度，其钟形轮廓导致两个结构性缺陷：(1) 长寿命基元逐渐衰减——要表示一个平坦的、时间不变的不透明度，需要多个重叠高斯来近似真实曲线，优化过程会反复插入新基元（这正是 EER 实验中"高斯数量反而增加"现象的另一种成因）；(2) 静态基元需要学习一个极小的速度来保持稳定，但优化无法收敛到绝对零速度，微小的残余运动随时间累积，导致可观测的空间漂移和模糊。SharpTimeGS 的关键洞察是：基元的运动行为和时间可见性都应该与其生命期耦合——这不通过外部正则化实现，而是直接嵌入表示公式。

\textbf{生命期参数与运动调制。}方法引入两个可学习的逐基元生命期参数——生命期方差 $\sigma_t$（控制基元衰减的快慢）和时间半径 $r$（基元保持全活跃的时间窗口）——将其同时嵌入运动和不透明度公式。与 EER 的 $k$-NN 应变惩罚不同，生命期调制不需要构建邻域图或计算高斯间距离，而是完全解耦的逐基元操作。运动调制函数为：

\begin{equation}
\mathbf{X}_t = \mathbf{X} + \frac{\mathbf{v}}{f(\sigma_t, r)}(t - T), \quad f(\sigma_t, r) = 1.0 + \frac{1.0}{(\sigma_t + r)^2}
\end{equation}

其中 $T$ 为时间锚点，$\mathbf{v}$ 为速度。关键设计在于分母 $f(\sigma_t, r)$：当 $\sigma_t + r$ 大时（长寿命/静态基元），$f$ 趋于无穷，$\mathbf{v}/f \to 0$，基元位置自动冻结——无需正则化即可实现 EER 花费大量计算量约束的"静态稳定性"。当 $\sigma_t + r$ 小时（短寿命/动态基元），$f$ 接近 1，允许完整速度发挥作用。两个参数均可通过反向传播自动学习。

\textbf{平顶不透明度轮廓。}在时间不透明度方面，方法用平顶轮廓（flat-top profile）替换高斯衰减：

\begin{equation}
O_t = O \cdot l(t), \quad l(t) = \begin{cases} 1, & |t - T| \leq r \\ \exp\!\left(-\frac{(|t-T|-r)^2}{\sigma_t^2}\right), & |t - T| > r \end{cases}
\end{equation}

静态基元（大 $r$）在平坦区间 $|t - T| \leq r$ 内保持不透明度恒为 1，仅在边界处急剧下降，消除了钟形曲线的渐进衰减问题——这使得一个静态基元即可覆盖整个视频序列，无需多个重叠基元的冗余表示。动态基元（小 $r$、小 $\sigma_t$）则保持快速进出衰减，精确建模瞬态事件。

\textbf{速度-生命期感知的致密化。}与 GAD 用损失率自适应调整全局致密化阈值不同，SharpTimeGS 在致密化阶段引入逐基元的评分机制：

\begin{equation}
s = \lambda_e E + \lambda_o O + \lambda_l\left(1 - e^{-|\mathbf{v}| + 1/f(\sigma_t, r)}\right)
\end{equation}

其中 $E$ 为累积重建误差，$O$ 为不透明度。最后一项是关键差异：短寿命（小 $f$）、快运动（大 $|\mathbf{v}|$）的基元获得更高评分，被更频繁地克隆。训练分为两阶段：前 $1/3$ 迭代使用 AbsGS 的标准致密化扩展场景覆盖，记录高斯数量 $N$ 并固定；后 $2/3$ 采用上述评分驱动的替换策略——每步移除低不透明度高斯，从评分最高的基元克隆等量新高斯，在保持总数不变的前提下将表示容量动态分配给快速瞬态区域。

\textbf{速度感知初始化。}方法进一步引入基于运动先验的初始化策略：通过 RAFT 光流检测运动区域，用 SAM2 获取完整物体掩码，再经 COLMAP 重建逐帧 3D 点云并通过 KNN 匹配估计初始速度 $\mathbf{v}_0$。动态点被赋予短生命期（$\sigma_t$ 覆盖 3 帧，$r = 10^{-6}$）和初始速度；静态点被赋予长生命期（$\sigma_t$ 覆盖 3 倍总帧数，$r = 10^{-6}$）和零速度。这一初始化为生命期参数提供了物理合理的起点，加速收敛。

\paragraph{全局-局部运动分层分解策略}

前述 Incoherent Deformation, Not Capacity 从正则化角度解决变形不一致性——用 EER 施加局部平滑先验，用 GAD（损失率感知致密化）控制高斯增殖。但它并未改变变形场本身的表示结构：一个六面体平面仍然需要同时编码全局刚体运动和局部非刚性变形，容量瓶颈依然存在。MoDec-GS（2025.01）~\cite{2501_03714} 采取了正交路线：不做正则化，而是将变形场拆成两级，让全局运动和局部运动由不同容量的表示分别承担。

\textbf{核心设计：两级变形替代单一六面体平面。}现实视频中的运动通常是全局运动（如物体整体位移、摄像机运动）和局部运动（如手指弯曲、衣角飘动）的复杂组合。MoDec-GS 同样以锚点表示（Scaffold-GS）作为规范空间基础，但提出全局到局部运动分解（Global-to-Local Motion Decomposition, GLMD），将变形分为两个串行阶段。第一阶段是全局锚点变形（Global Anchor Deformation, GAD）：直接变形锚点本身的位置和属性，而非变形重建后的每个 3D 高斯。锚点 $v$ 的位置 $(x_v, y_v, z_v)$ 连同规范时间 $t_c$ 输入一个小型全局六面体平面 $H_G$，查询的特征经 MLP 解码后得到锚点属性偏移量（如局部上下文特征偏移 $\Delta f_v = \lambda_f[F_G(H_G(x_v, y_v, z_v, t_c))]$），直接加到锚点属性上。每个锚点还附加两个可学习的动态参数 $d_G$（全局）和 $d_L$（局部），$d_G$ 通过可微二值掩码 $M(d_G) = (\mathbb{I}[\sigma(d_G) > \tau] - \sigma(d_G)) \cdot \text{sg}(\sigma(d_G)) + \sigma(d_G)$ 控制锚点是否参与全局运动，实现静态/动态区域的自动分离，无需预计算的外部掩码。消融实验表明，加入这对可学习动态参数对存储的增加几乎可以忽略，但能带来一致的质量提升。

\textbf{第二阶段精化局部运动。}在 GAD 捕获全局运动后，局部高斯变形（Local Gaussian Deformation, LGD）处理残余的非刚性细节：从变形后的锚点通过神经高斯推导生成 $k$ 个 3D 高斯，然后用共享的局部六面体平面 $H_L$ 逐高斯变形到目标时间戳（如位置偏移 $\Delta x_k, \Delta y_k, \Delta z_k = \lambda_p[F_L(H_L(x_k, y_k, z_k, t_L))]$）。这一设计的关键洞察是：经过 GAD 提取全局运动后，残余的局部运动幅度大幅减小、方向更分散，此时逐高斯的显式变形比继续锚点级变形更合适——因为局部变形需要的是每个高斯独立的小幅调整，而非锚点级的特征重新生成。消融实验中，将单阶段锚点变形扩展为两阶段 GLMD 后，PSNR 显著提升且存储反而降低（因为相邻高斯不再需要各自占用锚点）。局部六面体平面 $H_L$ 在每个局部规范空间（Local CS）内共享，进一步保证了紧凑性。

\textbf{TIA：自适应时间分辨率。}方法还引入时间间隔调整（Temporal Interval Adjustment, TIA）：初始均匀划分时间轴为 $l$ 段，训练过程中积累各段的位置梯度，梯度显著高于阈值 $\tau_{TIA}$ 的段被认为表达力不足，其时间间隔向两侧各缩小步长 $s_{TIA}$，实现运动剧烈段自动获得更细的时间分辨率。可视化实验验证了 TIA 的行为：初始均匀的时间间隔在训练过程中自动向光流幅度较大的区域收缩和偏移，使每个 Local CS 的表示能力被更高效地利用。

\textbf{与前述锚点方法的关系。}Incoherent Deformation, Not Capacity 的 EER 在优化层面约束变形一致性（局部平滑正则化），而 MoDec-GS 的 GLMD 从表示层面消除了单一变形场的容量瓶颈（两级分层表示）。两者分别从"约束已有表示"和"改用更合理的表示结构"两个角度解决同一类问题——前者轻量且可即插即用，后者需要架构修改但提供更强的归纳偏置。在 iPhone 数据集上，MoDec-GS 相比第二名 SC-GS，PSNR 提高 +0.7 dB，同时存储仅为 SC-GS 的约 6\%（降低 94\%）。消融实验还表明，将全局和局部六面体平面的尺寸缩小 55\% 对视觉质量的影响极小，证实了 GLMD 的两阶段分解使有限容量的表示被更高效地利用。

\paragraph{锚点驱动变形场与 4D 压缩}

前述 MoDec-GS 同样采用锚点结构（Scaffold-GS），但其核心目标是运动分解——将全局刚体运动与局部非刚性运动分离建模。ADC-GS（2025.05）~\cite{2505_08196} 则从存储压缩的角度出发，直击一个被忽视的事实：现有 4DGS 方法对每个高斯基元独立施加变形，但局部场景往往可近似为刚体运动，相邻高斯的变形高度相似——这种冗余不仅导致存储膨胀，还拖慢渲染速度。例如，RealTime4DGS 重建一段 10 秒视频需超过 2 GB 存储和 9 小时训练。MEGA 等方法虽通过分解球谐系数来压缩存储，但仍在逐高斯独立压缩，忽略了局部属性之间（除不透明度外）的强相似性。

\textbf{规范空间的锚点表示。}与 MoDec-GS 从 Scaffold-GS 继承锚点不同，ADC-GS 直接对 COLMAP 下采样点云初始化锚点：每个锚点 $v$ 关联 $K$ 个神经高斯，通过共享残差网络 $F_\theta$ 从锚点属性推导生成。锚点包含参考特征 $f_v \in \mathbb{R}^{N_v}$（捕捉 $K$ 个子高斯的共性）和残差特征 $f_g \in \mathbb{R}^{K \times N_g}$（捕捉个体差异），子高斯属性以残差形式生成：位置 $X_k = X_v + \Delta X_k$、协方差 $\Sigma_k = \Sigma_v \cdot \Delta\Sigma_k$、颜色 $C_k = C_v + \Delta C_k$；不透明度 $O_k$ 因缺乏空间相关性而直接预测。这种锚点-子高斯的层级关系，正是后续变形压缩效率的来源——变形锚点即可带动其所有子高斯。

\textbf{粗到细变形：利用锚点层级的变形冗余消除。}基于上述锚点结构，ADC-GS 提出粗到细（coarse-to-fine）变形策略，核心洞察是：锚点级别的变形已经捕捉了局部刚体运动（全局位置、协方差、颜色的整体偏移），只有不透明度和精细颜色残差需要逐高斯精化。

粗阶段仅变形锚点的显式属性，将参考特征 $f_v$ 与时间嵌入 $f_t = F_s(Z, t)$（$Z \in \mathbb{R}^{256}$ 为可学习插值权重，$F_s$ 为网格采样函数）拼接后输入一个小型 MLP：

\begin{equation}
(\Delta X_{v,t},\; \Delta\Sigma_{v,t},\; \Delta C_{v,t}) = F_\text{coarse}(f_v,\; f_t)
\end{equation}

$K$ 个子高斯自动跟随锚点变形，变形调用量从逐高斯降低到逐锚点。细阶段则利用变形后锚点的位置嵌入 $f_p$ 和 $f_t$，通过精化网络 $F_\text{fine}$ 逐高斯调整不透明度和颜色：$(\Delta O_{k,t},\; \Delta C_{k,t}) = F_\text{fine}(f_p, f_t)$。最终帧 $T$ 处的子高斯属性为：

\begin{equation}
X_{k,t} = X_k + \Delta X_{v,t}, \quad \Sigma_{k,t} = \Sigma_k \cdot \Delta\Sigma_{v,t}, \quad O_{k,t} = O_k + \Delta O_{k,t}, \quad C_{k,t} = C_k + \Delta C_{v,t} + \Delta C_{k,t}
\end{equation}

注意颜色 $C_{k,t}$ 同时继承了锚点级粗变形 $\Delta C_{v,t}$ 和逐高斯精细变形 $\Delta C_{k,t}$，体现了粗到细的层级叠加。

\textbf{多维熵模型：锚点特征的率失真压缩。}锚点结构天然适合压缩——$K$ 个子高斯的共性被浓缩到单个参考特征 $f_v$ 中，而非 $K$ 份独立存储。ADC-GS 进一步引入多维熵模型（Multi-dimension Entropy Model, MEM），对四种锚点特征（$f_v$、$f_g$、$\Sigma_v$、$C_v$）分别建模概率分布，实现熵编码压缩。锚点位置 $X_v$ 使用 G-PCC 编码。

核心设计是自适应量化与分块上下文编码。首先，对特征 $f_i$ 使用自适应量化 $\hat{f}_i = f_i + \mathcal{U}(-\frac{1}{2}, \frac{1}{2}) + Q_i(1 + \tanh(F_q(f_i)))$，其中 $Q_i$ 为预设量化步长（四种特征各自不同）、$F_q$ 为 MLP 调节器，使不同尺度的特征获得灵活的量化精度；均匀噪声在训练时注入以模拟量化损失。概率建模采用层级上下文结构：参考特征的概率分布由超先验预测为高斯 $p(\hat{f}_v) = \mathcal{N}(\mu_{f_v}, \sigma_{f_v})$；协方差和颜色的分布类似建模，均以已解码的 $\hat{f}_v$ 作为上下文。残差特征 $f_g$ 占码流最大份额，被分块为 $M$ 段，已解码的参考特征和前序块作为上下文引导后续块的编码（$p(\hat{f}_g^{ch}) = \mathcal{N}(\mu_{f_g}, \sigma_{f_g})$，参数由 $E_{f_g}(\hat{f}_v \| \bigcup_{m=1}^{ch-1} \hat{f}_g^m)$ 预测），有效消除通道间冗余。总码率估计为 $R = -\sum \log_2 p(\hat{f}_v) - \sum \log_2 p(\hat{f}_g^{ch}) - \sum \log_2 p(\hat{\Sigma}_v) - \sum \log_2 p(\hat{C}_v)$，通过率失真损失 $L_{\text{loss}} = (1-\lambda_{\text{ssim}})L_1 + \lambda_{\text{ssim}}L_{\text{ssim}} + \lambda_e R$ 联合优化。

\textbf{时间显著性驱动的锚点精炼。}前述 Incoherent Deformation, Not Capacity 中的 GAD 通过损失率感知阈值控制静态场景中的高斯增殖，但动态场景中高斯的时间重要性跨帧变化，静态场景的自适应密度控制不适用。ADC-GS 提出时间显著性（temporal significance）驱动的锚点增长和剪枝：按渲染权重加权累积 $N$ 帧的位置梯度 $s_k = \frac{\sum_{t}^{N} \omega(k,t) \|g_{k,t}\|}{\sum_{t}^{N} \omega(k,t)}$，其中 $\omega(k,t)$ 是 $\alpha$-混合 的渲染权重。超过阈值的子高斯提升为新锚点（继承所属锚点的隐特征，保留自身高斯属性），不透明度低于阈值的锚点被剪枝。相比现有 4DGS 方法，ADC-GS 实现了高达 200 倍的存储压缩。

\paragraph{混合专家融合多种变形策略}

前述 Incoherent Deformation, Not Capacity、FreeTimeGS++、SharpTimeGS 等方法共享一个隐含假设：存在某种"正确的"单一变形先验，只要找到并施加它，就能解决动态场景重建中的过拟合、不一致变形和运动噪声等问题。MoE-GS（2025.10）~\cite{2510_19210} 从一个截然不同的角度切入这一困局——也许问题不在于某种先验"不够好"，而在于真实动态场景同时包含多种运动模式，任何单一变形先验注定无法全局最优。

\textbf{动机：单一先验的结构性失败。}作者在 N3V 和 Technicolor 基准上系统比较了四类代表性方法，发现性能退化呈三级空间粒度：（1）场景级——不同方法在不同场景上表现差异巨大，没有"万能赢家"；（2）空间级——同一场景内，不同空间区域由不同方法主导；（3）时间级——同一视频中，最优方法逐帧切换。这一现象的根源在于各方法不同的归纳偏置：HexPlane 规范变形（如 4DGaussians）生成短而光滑的轨迹，因其全局位置--时间网格使空间相邻高斯接收几乎相同的变形信号，在静态或低运动区域表现优异，但在快速运动区域因规范特征无法足够快地变化而过度平滑；逐高斯体积变形（如 E-D3DGS）的双分支网络产生方向一致且高速的轨迹，擅长捕捉快速连贯流；关键帧插值方法（如 Ex4DGS）让每个高斯独立更新位置，轨迹多样且自由，能捕捉不规则快速位移，但缺乏空间连贯性；多项式轨迹模型（如 STG）偏向全局平滑低曲率运动，适合刚体运动但难以表示高频变化。因为真实场景同时包含上述多种运动模式——一个厨房场景中，水龙头区域是刚体运动、水流是不规则高速流动、背景是静态——所以任何单一变形先验必然在某些区域失败。这正是 Incoherent Deformation 所揭示的不一致变形问题的另一面：当用不适合该区域的先验去建模时，变形场被迫在不匹配的约束下"扭曲"自己，结果就是相邻高斯之间出现不协调。

\textbf{MoE 框架：异构专家的自适应融合。}基于上述观察，MoE-GS 提出混合专家（Mixture of Experts, MoE）框架，集成 $N$ 个异构动态 GS 方法作为专家网络 $E_1, \dots, E_N$，由路由器逐像素决定各专家的权重。与 LLM 中以稀疏激活降低计算量为目标的 MoE 不同，MoE-GS 的首要目标是增加表示容量——通过组合不同归纳偏置来覆盖异构运动模式；效率优化是独立的后续问题。

\textbf{Volume-aware Pixel Router：在 3D 编码与 2D 混合之间架桥。}路由器设计面临一个核心张力：纯像素级路由（由 MLP 直接根据时间和视角输出权重）缺乏体积感知，无法区分被遮挡的前后物体——实验中产生过度平滑的结果，边界模糊；纯高斯级路由（在 3D 空间中直接调制高斯不透明度）虽保留几何结构，但因高斯参数到像素贡献的映射间接且复杂，优化过程中频繁出现不稳定和过平滑伪影——Volume-aware Pixel Router 在 N3V 上的 PSNR 分别比这两种方案高出 2.11 dB 和 1.18 dB。解决方案是"高斯级编码 + 泼溅投影 + 像素级混合"：在每个高斯 $G_i$ 上附加可学习的时-视权重 $\mathbf{w}_i^{\text{per}} = [w_i, w_i^{\text{dir}}, (t \cdot w_i^{\text{time}})]^\top$，通过可微高斯泼溅将这些 3D 权重投影到像素空间，获得聚合了所有重叠高斯贡献的像素级嵌入 $\mathbf{w}_{2D}(u)$、$\mathbf{w}_{2D}^{\text{dir}}(u)$、$\mathbf{w}_{2D}^{\text{time}}(u)$，再经轻量 MLP $\phi$ 计算路由权重：

\begin{equation}
R'(u) = w_{2D}(u) + \phi(\mathbf{w}_{2D}^{\text{dir}}(u), \mathbf{w}_{2D}^{\text{time}}(u), \mathbf{r}(u))
\end{equation}

其中 $\mathbf{r}(u)$ 是像素观察方向。路由权重经 softmax 归一化后，逐像素混合各专家的渲染输出：

\begin{equation}
I_{\text{MoE}}(u) = \sum_{k=1}^{N} G'_k(u) \cdot I_{E_k}(u), \quad G'_k(u) = \text{softmax}(R'_k(u))
\end{equation}

直觉上，$w_{2D}(u)$ 捕获体积结构的全局信息（深度、可见性、变形），$\phi$ 项则注入视角和时间的细粒度变化。这种设计的优势是双向的：权重在泼溅前定义于高斯空间，因此可通过责任加权（responsibility-weighted）方式自然地反向提升到 3D 域，支持事后的高斯级分析和剪枝；而混合在 2D 像素空间进行，训练稳定且不干扰各专家的内部优化。训练采用两阶段策略：先独立优化各专家至收敛，再冻结专家参数、仅训练路由器的权重 $\mathbf{w}_i^{\text{per}}$ 和 MLP $\phi$——路由器本身的开销不到专家训练的 5\%。

\textbf{训练开销远低于线性增长。}使用 $N$ 个专家直觉上意味着 $N$ 倍训练时间，但实验揭示了一个有趣的特性：在 $N=3$ 配置下，即使每个专家仅用 20\% 的训练预算（如将 E-D3DGS 的 4.2 小时缩减至约 50 分钟），MoE-GS 的 PSNR（32.60 dB）仍高于所有完全训练的单一专家（最高 32.33 dB）。甚至 10\% 预算的 MoE-GS（32.14 dB）仍可与完全训练的单一方法竞争。这说明 MoE-GS 不依赖完全收敛的专家——多个"半成品"专家的组合足以超越任何单个"成品"。

\textbf{效率路径一：单次渲染与门控剪枝。}为解决多专家推理的计算开销，MoE-GS 提出 Single-Pass Multi-Expert Rendering：将所有专家的高斯合并为一个批次，每个高斯携带独热专家身份标记 $\mathbf{e}_j \in \mathbb{R}^K$，一次投影计算所有高斯的可见性和排序，仅在 $\alpha$-混合阶段通过 $(\mathbf{e}_j)_k$ 选择属于专家 $k$ 的高斯。这消除了 $K$ 次独立光栅化的冗余核启动和内存遍历。在此基础上，Gate-Aware Gaussian Pruning 通过累计门控权重对每个高斯权重的梯度来衡量其全局重要性：

\begin{equation}
E_i = \frac{1}{|D|}\sum_{v \in D}\|\nabla_{\mathbf{w}_i^{\text{per}}} G'_k(v)\|
\end{equation}

梯度接近零的高斯对路由决策几乎无影响，剪除 $E_i < \tau$ 的高斯。在 $N=2$ 配置下，55\% 剪枝率的 MoE-GS 几乎不损失 PSNR（32.80 vs.\ 32.82 dB），但将存储从 878.7 MB 降至 351.2 MB、FPS 从 44 提升至 83；75\% 剪枝率下 PSNR 仅降 0.37 dB，但 FPS 达到 101、存储降至 281.3 MB。

\textbf{效率路径二：知识蒸馏。}当专家数量较多（$N \geq 4$）时，直接推理全部专家代价过高。MoE-GS 提出知识蒸馏策略：用完全训练的 MoE 模型的渲染输出 $I_{\text{MoE}}$ 作为伪标签，结合路由权重 $G'_k$ 作为置信度分数，从头训练单个专家。蒸馏损失为：

\begin{equation}
L_k^{KD} = \beta \, L(G'_k I_{E_k}, G'_k I_{GT}) + (1-\beta) \, L((1-G'_k) I_{E_k}, (1-G'_k) I_{\text{MoE}})
\end{equation}

其中 $L$ 结合 L1 和 SSIM 损失，$\beta$ 平衡真实标签与 MoE 伪标签的权重。核心思想是：$G'_k$ 高的区域说明专家 $k$ 本已擅长，直接用真实标签监督即可；$G'_k$ 低的区域说明专家 $k$ 原本不擅长，此时用 MoE 输出作为软标签，让该专家学习其他专家在这些区域的知识。实验表明，蒸馏后的单个专家在所有指标上一致优于其未蒸馏版本——例如在 Technicolor 上，蒸馏后的 E-D3DGS PSNR 从 32.88 提升至 33.67（+0.79 dB），Ex4DGS 从 33.57 提升至 33.91（+0.34 dB）——且无需修改专家架构即可逼近完整 MoE 的性能。

\textbf{与前述锚点方法的关系。}Incoherent Deformation, Not Capacity 的 EER 试图用局部平滑正则化使单一变形场变得一致，FreeTimeGS++ 试图用共享速度场从表示层面消除不一致的根源，而 MoE-GS 采取了第三条路线：承认不同变形先验在不同运动模式下的不可替代性，用路由器自适应选择"局部最优先验"。从这个角度看，MoE-GS 不是在修复某种变形场的不一致性，而是在设计一个更高层次的选择机制，让每种先验只在它最擅长的区域发挥作用——这同时也避免了 Incoherent Deformation 所揭示的"不适合的先验导致不协调变形"问题。

\paragraph{时空解耦渲染管线}

前述 Incoherent Deformation, Not Capacity 通过 EER 在优化层面约束变形一致性，揭示了"相邻高斯变形不协调"是过拟合的根源。但 EER 的工作前提是变形场能够正确地学会"把哪个高斯变到哪里"——如果初始化阶段就已经把来自不同时间步的高斯叠加在同一个规范空间位置上，变形场连"这个高斯属于哪个时刻"都无法判定，再强的平滑约束也无济于事。STDR（2025.05）~\cite{2505_22400}（Spatio-Temporal Decoupling for Real-time rendering）正是从这一观察出发，将问题定位到更早的阶段：不是变形场学不好，而是初始化造成了不可逆的时空混淆（spatio-temporal incoherence）。

\textbf{时空混淆的诊断。}现有动态 3DGS 方法普遍采用两阶段流水线：先用所有输入帧在共享时间戳 $t_\gamma$ 下初始化规范高斯，再学习变形场建模动态。问题在于，warm-up 阶段把来自不同时刻的观测当作同一时刻来处理——例如一个移动的挖掘机臂，其轨迹上多个位置的高斯被叠加到同一规范空间中，形成"鬼影"重叠。变形场随后必须区分这些空间上靠近但时间上远离的高斯，却没有任何显式的时间引导。这种初始化引起的歧义具体表现为两类困难：(1) 身份歧义——变形场无法判定哪些高斯属于哪个时间步，多个高斯占据相似位置却代表不同时间状态；(2) 目标冲突——同一空间区域需要根据时间戳映射到多个不同位置，产生相互矛盾的变形目标。最终呈现为鬼影效应、运动模糊和时间不一致的变形。与 Incoherent Deformation 诊断的"相邻高斯变形方向不协调"不同，STDR 诊断的是更上游的"初始化阶段就丢失了时间身份"——前者是变形场优化中的不一致性，后者是表示层面的结构性缺陷。

\textbf{组件 1：时空掩码——为每个高斯赋予时间身份。}STDR 的核心操作是为每个高斯 $i$ 添加可学习的时间掩码 $\mathbf{m}_i \in \mathbb{R}^K$（$K$ 为总帧数），在渲染时调制原始不透明度：

\begin{equation}
\alpha_i^{\text{new}} = m_i^j \cdot \alpha_i
\end{equation}

其中 $j$ 为当前时间戳。这一设计的精妙之处在于：掩码不需要手工标注或外部监督——训练过程中，静态区域的高斯在各时间步获得一致的梯度，掩码趋于均匀分布；动态区域的高斯因时间错位产生较大的渲染误差，反向传播驱动掩码向正确的时间表示收敛。掩码经 softmax 归一化后得到时空概率分布：

\begin{equation}
\tilde{m}_i^j = \frac{\exp(m_i^j)}{\sum_{k=1}^{K} \exp(m_i^k)}
\end{equation}

该分布编码了"高斯 $i$ 在各个时间步的激活强度"，相当于为每个高斯赋予了一个可学习的时间身份。与 EER 的 $k$-NN 应变惩罚不同，时空掩码不需要构建邻域图，而是完全解耦的逐高斯操作——但它隐含地解决了 EER 无法触及的问题：即使两个高斯在规范空间中位置重合，只要它们的时间概率分布不同，变形场就能将它们区分开。

\textbf{组件 2：分离变形场——基于时间身份的解耦变形。}有了时空概率分布后，STDR 不再直接基于位置变形，而是先提取解耦的空间和时间特征：

\begin{equation}
z_s, z_t = f_{\text{sep}}(\mathbf{x}_\gamma, \tilde{\mathbf{m}}_i), \quad \mathbf{G}_\tau = f_{\text{def}}(\mathbf{x}_\gamma, z_s, z_t, t_\tau)
\end{equation}

其中 $f_{\text{sep}}$ 为 6 层 MLP，$\mathbf{x}_\gamma$ 为规范空间位置，$\tilde{\mathbf{m}}_i$ 为归一化的时空概率分布。空间特征 $z_s$ 捕获结构上下文以区分静态背景和动态物体，时间特征 $z_t$ 编码时变运动模式以赋予高斯时间身份——相同时间身份的高斯表现出一致的运动行为。与前述方法的对比：EER 通过约束"变形后的位移差"来保证一致性，FreeTimeGS++ 的 NVF 通过共享速度场消除运动噪声的根源，而 STDR 的分离变形场通过"时间身份 + 空间上下文"的显式解耦来解决初始化歧义——三者分别从"优化正则化"、"表示共享"和"时空解耦"三个层面切入同一族问题。

\textbf{组件 3：时空一致性正则化。}为进一步稳定学习过程，STDR 引入两项互补的正则化。时间平滑正则化鼓励掩码在相邻帧间平滑变化，防止突变：

\begin{equation}
L_{\text{temp}} = \lambda_1 \sum_{i} \sum_{t} \|m_i^{t} - m_i^{t+1}\|_2^2
\end{equation}

空间感知正则化基于 KL 散度，鼓励空间相邻高斯共享相似的时间行为（实现上采样 $M$ 个高斯，为每个高斯找 $K$ 个最近邻）：

\begin{equation}
L_{\text{spatial}} = \frac{\lambda_2}{MK} \sum_{i=1}^{M} \sum_{j=1}^{K} \sum_{d=1}^{D} m_i^{d} \log\left(\frac{m_i^{d}}{m_j^{d}}\right)
\end{equation}

其中 $D$ 为总时间步数。总训练损失为 $L = \lambda L_1 + (1 - \lambda) L_{\text{D-SSIM}} + \lambda_1 L_{\text{temp}} + \lambda_2 L_{\text{spatial}}$。在 SC-GS 上的消融表明，移除 $L_{\text{temp}}$ 或 $L_{\text{spatial}}$ 均导致 PSNR 下降（分别为 $-0.33$ dB 和 $-0.13$ dB），两项缺一不可。

\textbf{即插即用与实验效果。}STDR 不修改宿主方法的架构，仅需添加掩码参数和分离变形场，可无缝集成到 4DGS、DeformGS、SPGS、SC-GS 等基线中。在 D-NeRF 数据集上，STDR + SC-GS 将 PSNR 从 41.66 dB 提升至 42.24 dB（+0.58 dB），SSIM 从 0.993 提升至 0.997，LPIPS 从 0.0090 降至 0.0059。在 NeRF-DS 和 HyperNeRF 等真实世界数据集上同样带来一致的提升。定性上，STDR 有效消除了 T-Rex 场景中的脚部鬼影和 HyperNeRF 扫帚场景中的严重伪影。

\paragraph{可扩展运动层次化 4D 高斯}

前述 Incoherent Deformation, Not Capacity 揭示了变形不一致性是过拟合的根源，并通过 EER 平滑正则化从优化层面约束相邻高斯的变形行为。然而，EER 的策略本质上是"事后修补"——变形场本身仍然给予每个高斯独立的变形自由度，一致性仅靠损失函数的外部压力维持。MotionScale（2026.03）~\cite{2603_29296} 从一个根本不同的角度切入：如果不一致性来自变形场对每个高斯过于自由的参数化，那么能否设计一种运动表示，使一致性成为架构的天然属性而非优化目标？此外，该方法还瞄准了一个 EER 未涉及的挑战——可扩展性。当前 4DGS 方法存在两个系统性瓶颈：(1) 欠约束几何——监督主要依赖视角相关的外观信号，缺乏强制 3D 结构一致性的能力；(2) 累积时间漂移——运动模型依赖缺乏 3D 感知的 2D 跟踪先验，长序列中误差不可避免地累积。这两个瓶颈在大规模、长序列场景中尤为突出。

\textbf{基于聚类的可扩展运动场。}针对 MotionScale 聚类运动场在大规模场景中的局限——划分 $K$ 个不重叠聚类本身依赖预先设定的聚类数量，若场景实际运动模式数超过 $K$ 则局部细节丢失，若远少于 $K$ 则部分聚类退化为空簇引入参数开销；且聚类划分基于高斯的初始空间分布而非运动相似性，相邻高斯若被误分入不同聚类，其共享运动基的架构假设将与实际运动轨迹产生系统性偏差——MotionScale 的核心设计是将全局变形场替换为聚类层次化运动场（cluster-centric motion field）。直觉上，现实场景中不同区域的运动模式差异很大——一个行走的人，躯干的刚体运动和衣角的非刚性飘动需要不同粒度的建模能力。使用固定容量的全局 MLP 或预定义时间基，要么因容量不足丢失局部细节，要么因参数过多过拟合训练帧。MotionScale 将动态高斯 $\mathcal{G}_d = \{g_i\}_{i=1}^{N_d}$ 划分为 $K$ 个不重叠聚类 $\{\mathcal{C}_k\}_{k=1}^K$，每个聚类定义一个层次化运动模型：全局刚体变换 $\mathbf{G}_k^t = [\mathbf{R}_{k,g}^t \mid \mathbf{t}_{k,g}^t] \in SE(3)$ 捕获聚类的整体刚性运动，以及 $B$ 个细粒度基变换 $\mathcal{B}_k^t = \{(\mathbf{r}_{k,b}^t, \mathbf{t}_{k,b}^t)\}_{b=1}^B$ 建模局部非刚性变形。这种设计与 EER 形成鲜明对比：EER 通过 $k$-NN 邻域图和应变惩罚在优化中强制变形一致，而聚类运动场在架构上保证了同一聚类内的高斯共享同一组运动基——一致性不再需要额外的正则化项来维持。每个高斯 $g_i \in \mathcal{C}_k$ 被分配可学习系数向量 $\mathbf{w}_i = [w_{i,1}, \dots, w_{i,B}]^\top$（满足 $\sum_{b=1}^B w_{i,b} = 1$），通过混合聚类基计算局部变换 $\mathbf{L}_i^t$：

\begin{equation}
\mathbf{R}_{i,\ell}^t = \mathcal{R}\!\left(\sum_{b=1}^B w_{i,b}\, \mathbf{r}_{k,b}^t\right), \quad \mathbf{t}_{i,\ell}^t = \sum_{b=1}^B w_{i,b}\, \mathbf{t}_{k,b}^t
\end{equation}

其中 $\mathcal{R}(\cdot)$ 映射到 $SO(3)$（使用 6D 连续旋转表示）。最终高斯状态由全局和局部变换的复合得到：

\begin{equation}
\boldsymbol{\mu}_i^t = \mathbf{R}_{k,g}^t(\mathbf{R}_{i,\ell}^t \boldsymbol{\mu}_i^0 + \mathbf{t}_{i,\ell}^t) + \mathbf{t}_{k,g}^t, \quad \mathbf{R}_i^t = \mathbf{R}_{k,g}^t \mathbf{R}_{i,\ell}^t \mathbf{R}_i^0
\end{equation}

每个高斯仅受一个聚类影响，因此计算成本近乎恒定，且运动场的规模随聚类数量自然扩展——这与 EER 训练中高斯数量反而增加 85\% 的副作用形成对比。

\textbf{自适应控制：运动不一致检测与聚类分裂。}在架构层面保证同一聚类内的一致性之后，一个自然的问题是：如果一个聚类内部确实包含了运动模式不同的高斯怎么办？MotionScale 引入自适应控制机制来检测和修复这种不一致——这与 Incoherent Deformation, Not Capacity 对不一致变形的诊断形成了有趣的互补关系：后者通过应变 $\varepsilon_i$ 测量训练完成后的变形不一致性，而 MotionScale 在优化过程中主动检测并分裂不一致的聚类。具体地，在长期细化阶段，方法提取聚类内高斯在传播窗口上的 3D 轨迹作为特征描述符，先经 HDBSCAN（Hierarchical Density-Based Spatial Clustering of Applications with Noise，层次密度聚类）密度聚类识别子结构，再通过层次聚类将运动不一致的聚类分裂为两个子聚类。分裂时复制原始运动参数给两个新聚类，保证分裂瞬间运动状态完全一致，允许两簇在后续优化中独立发散。这种设计与 EER 的 GAD 致密化阈值策略不同——GAD 通过自适应阈值防止高斯增殖，而 MotionScale 的自适应控制是在运动基的层面扩展容量，不是增加高斯数量，而是增加运动模式的表达力。同时，过小的聚类被剪枝，保持表示紧凑。

\textbf{渐进式三阶段优化。}为解决长序列中的累积时间漂移，MotionScale 采用渐进式传播策略，将新帧分批融入已优化的窗口。背景扩展阶段（background extension）处理新帧中暴露的未见区域：将已有背景高斯投影到新图像平面，标记无覆盖像素，用单目深度图初始化补充高斯，并通过光度损失端到端优化联合精化相机外参。前景传播阶段（foreground propagation）通过三阶段精化确保运动的时间连续性：(1) 初始对齐——在短期传播窗口内，仅用单向跟踪损失 $L_{\text{track}}$（从已优化帧指向新帧）优化新帧的运动基，冻结所有其他参数，避免未对齐的新数据污染已收敛的旧帧；(2) 短期一致性——切换为窗口内任意帧对之间的双向跟踪损失，允许梯度回流到近期帧，细化局部时间一致性；(3) 长期细化——跨全序列采样帧对联合优化所有参数并允许高斯致密化，结合跟踪损失 $L_{\text{track}} = \frac{1}{|I_t|}\sum_{p \in I_t}\|\hat{\mathbf{U}}_{t \to t'}(p) - \mathbf{U}_{t \to t'}(p)\|$、深度一致性损失 $L_{\text{depth}}$、光度损失和 ARAP（as-rigid-as-possible）正则化项，消除累积漂移和遮挡。阶段 (1) 和 (2) 随新帧不断纳入而重复迭代，阶段 (3) 在积累足够帧数后执行一次全局精化。

\textbf{阴影高斯。}方法还引入一组专用的"阴影高斯"（shadow Gaussians），用于建模移动物体投射的瞬态光照效果。这些高斯附加在背景表示中，由前景分割掩码引导初始化，仅通过光度损失和分割约束优化，不参与运动场的传播。这种设计与 STDR 的时空掩码有相似的思路——都通过显式分离静态和瞬态成分来提高渲染质量——但 MotionScale 的阴影高斯是几何层面的专用基元，而非调制不透明度的掩码参数。

\textbf{与前述锚点方法的关系总结。}Incoherent Deformation, Not Capacity 从诊断出发，揭示了变形不一致性是过拟合的根源，并通过 EER（平滑正则化）和 GAD（自适应致密化阈值）在优化层面抑制不一致性。SharpTimeGS 从表示层面引入生命期参数，通过运动调制函数使静态基元自动冻结、动态基元自由运动。MotionScale 则提供了第三条路径：用聚类层次化运动场替代全局变形场，使一致性成为架构的天然属性。三者的策略差异可以用一个类比来理解——EER 像是给每辆车安装速度限制器（外部约束），SharpTimeGS 像是根据车型预设不同的速度上限（参数耦合），而 MotionScale 像是把车编队行驶（共享运动基），一致性由编队结构本身保证。在实验效果上，MotionScale 在具有挑战性的真实世界基准上显著超越现有 4DGS 方法，同时展示了对大规模场景和长序列的可扩展性。

\subsubsection{高效与紧凑表示}

本路线介绍通过属性解耦压缩、结构化锚点表示和层级混合表示来设计更紧凑4D表示结构的方法。

正则化和优化策略改善了训练质量，但 4D 表示的存储开销仍然巨大。本节不改善训练过程本身，而是设计更紧凑的 4D 表示结构，从根本上压缩存储和计算开销。动态场景的 4D 表示天然比静态 3D 表示更加庞大——每个时间步都需要独立或半独立的参数来编码变形信息。如果静态场景需要 $N$ 个 3D 高斯，那么包含 $T$ 个时间步的动态场景理论上需要 $O(N \times T)$ 的参数量。如何在保持渲染质量的前提下压缩这一表示、提升训练和推理效率，是一个关键的研究方向。本节的方法围绕三个核心策略展开：属性解耦压缩（MEGA、EDGS）、结构化锚点表示（SD-GS、4D Scaffold GS、Superpoint GS）和层级/混合表示（CTRL-GS、Hybrid 3D-4D、SaRO-GS）。同时期还有 Scale-GS（2025.08）~\cite{2508_21444} 通过冗余过滤训练实现高效可扩展高斯泼溅。

\paragraph{运动感知压缩的内存高效 4D 高斯}

前述 Incoherent Deformation, Not Capacity 通过 EER 平滑正则化和 GAD 致密化控制，从优化层面抑制了变形不一致性带来的过拟合。但它留下的一个根本问题未被触及：即使变形场完全一致，4D 高斯基元本身的参数结构仍存在巨大的存储浪费——每个 4D 高斯 161 个参数中球谐系数占 144 个，且时间衰减机制使超过 90\% 的高斯在任何单帧中不可见。换言之，过拟合被控制住了，但存储效率和参数利用率依然低下。MEGA（2024.10）~\cite{2410_13613}（Memory-Efficient 4D Gaussian Splatting，ICCV 2025）从表示结构设计层面切入这一瓶颈，而非继续在优化正则化上做文章——用 DC-AC 解耦颜色替代 4D 球谐系数来压缩参数冗余，用时-视感知变形预测器替代固定线性运动假设来提高每个高斯的有效利用率，同时引入不透明度熵损失主动控制高斯增殖。4D 高斯的时间切片和可微光栅化流程保持不变。

\textbf{Baseline 的内存瓶颈——两重浪费叠加。}上述低效来自两个相互叠加的问题。第一个问题出在颜色参数上：标准 4DGS 用 4D 球谐系数（spherical harmonics）建模随视角和时间变化的颜色，自由度 $k_v$ 和 $k_t$ 的组合使每个高斯需要 144 个球谐参数，占其总参数量 161 个的 $89.4\%$。这些系数存在大量冗余——一个物体的"基础颜色"通常稳定，但每个高斯仍被分配完整的 4D 球谐展开。渲染 Birthday 场景时约 1300 万个高斯点的存储开销高达约 7.79 GB，其中绝大部分来自球谐系数。

第二个问题出在运动建模上：4DGS 假设切片后的 3D 高斯仅做线性运动（$\boldsymbol{\mu}_{3D}(t) = \boldsymbol{\mu}_{3D} + (t - \mu_t) \cdot \boldsymbol{\delta}$）且旋转和缩放在时间切片后保持恒定，复杂非线性运动必须依赖大量额外高斯组合表示。再加上时间衰减不透明度 $\sigma_t(t) = e^{-(t-\mu_t)^2/2W}$ 使每个高斯仅在其时间中心附近可见，每帧超过 90\% 的高斯被排除在渲染之外（参与率低于 50\%），迫使模型持续分裂和克隆高斯以弥补覆盖不足。这两重问题——颜色参数冗余与高斯利用率低下——共同导致 4DGS 的存储膨胀。MEGA 的三条策略分别精确对应这三个瓶颈：DAC 压缩颜色冗余，变形预测器提高高斯利用率，不透明度熵损失控制增殖。

\textbf{策略一：DC-AC 颜色解耦（DAC）。}直觉上，一个物体的"基础颜色"（如苹果的红色）不应随视角和时间剧烈变化，而高光、阴影等"变化部分"可以用一个轻量网络按需生成。MEGA 受电气工程中直流/交流分量概念的启发，将颜色属性解耦为逐高斯 DC 颜色分量 $\mathbf{c}_{dc} \in \mathbb{R}^3$（仅需 3 个参数，替代原 144 个球谐系数）和一个时-视感知 AC 颜色预测器 $F_\theta$。最终颜色通过混合预测得到：

\begin{equation}
\mathbf{c}_{t,v} = \text{sigmoid}(\mathbf{c}_{dc} + F_\theta(\text{sg}(\boldsymbol{\mu}_{3D}), \text{sg}(\mathbf{d}_v), t, \mathbf{c}_{dc}))
\end{equation}

即最终颜色等于 DC 基础颜色加上 AC 预测的视角-时间修正项，再经 sigmoid 映射到 $[0,1]$。其中 $\text{sg}(\cdot)$ 为停止梯度操作（切断 $\boldsymbol{\mu}_{3D}$ 和 $\mathbf{d}_v$ 对颜色预测的反向传播，防止颜色损失干扰几何优化），$F_\theta$ 是一个仅含三层线性层的轻量 MLP。DC 分量保留每个高斯的个体颜色信息，AC 预测器补充缺失的视角和时间信息，在等效高斯点数下实现约 8 倍压缩比。值得注意的是，直接使用 grid-based 神经场替换球谐系数虽然也能减少参数，但实验表明会导致严重的渲染质量下降——因为网格无法保留足够的个体细节信息。DAC 的解耦设计恰好解决了这一问题：DC 分量显式保留每高斯的固有颜色，AC 预测器专注于建模视角和时间的动态变化。

\textbf{策略二：熵约束高斯变形场。}Baseline 中每个高斯仅在线性运动假设和时间衰减机制下有很小的时间作用范围，导致每帧超过 90\% 的高斯处于"不可见"状态。MEGA 引入时-视感知变形预测器来扩大每个高斯的有效作用范围。具体而言，将 4D 高斯中心 $\boldsymbol{\mu}_{4D}$、视方向 $\mathbf{d}_v$ 和时间 $t$ 经频率位置编码映射到高维空间后输入 MLP $F_d$，预测位置、缩放和旋转的变形量：

\begin{equation}
(\Delta\boldsymbol{\mu}_{4D}^{t,v},\; \Delta\mathbf{s}_{4D}^{t,v},\; \Delta\mathbf{q}_l^{t,v},\; \Delta\mathbf{q}_r^{t,v}) = F_d(\gamma(\text{sg}(\boldsymbol{\mu}_{4D})),\; \gamma(\text{sg}(\mathbf{d}_v)),\; \gamma(t))
\end{equation}

其中频率位置编码 $\gamma(\mathbf{p}) = (\sin(2^l \mathbf{p}), \cos(2^l \mathbf{p}))_{l=0}^{L-1}$ 将低维输入映射到高维频率空间，使 MLP 能捕捉多尺度的时空变化。预测得到的变形量以乘法方式作用于原始 4D 高斯属性：

\begin{equation}
\boldsymbol{\mu}_{4D}^{t,v} = \boldsymbol{\mu}_{4D} \cdot \Delta\boldsymbol{\mu}_{4D}^{t,v}, \quad \mathbf{s}_{4D}^{t,v} = \mathbf{s}_{4D} \cdot \Delta\mathbf{s}_{4D}^{t,v}, \quad \mathbf{q}_l^{t,v} = \mathbf{q}_l \cdot \Delta\mathbf{q}_l^{t,v}, \quad \mathbf{q}_r^{t,v} = \mathbf{q}_r \cdot \Delta\mathbf{q}_r^{t,v}
\end{equation}

即变形后的 4D 高斯在原始属性基础上按预测的缩放因子进行调整——位置被缩放平移，尺度和旋转被缩放修正，从而显著扩大每个高斯的时空作用范围。

然而，单纯引入变形预测器而不约束高斯数量会导致高斯在致密化过程中持续分裂和克隆，产生数量膨胀。为此，MEGA 引入空间不透明度熵损失来迫使每个高斯"做出选择"——要么高不透明度（积极参与渲染），要么低不透明度（被剪枝淘汰）：

\begin{equation}
L_{\text{opa}} = -\frac{1}{N}\sum_{j=1}^{N} o_j \log(o_j)
\end{equation}

该损失函数当 $o_j$ 接近 0 或 1 时取最小值（熵最小），推动每个高斯的空间不透明度趋向二值状态。训练中配合每 $K$ 次迭代主动剪枝近零不透明度的高斯，将高斯参与率从不到 50\% 提升至约 75\%，同时有效控制高斯总量。完整训练损失为 $L = (1-\lambda)\mathcal{L}_1 + \lambda \mathcal{L}_{\text{SSIM}} + \beta L_{\text{opa}}$，其中 $\lambda=0.2$, $\beta=0.0005$。

\textbf{策略三：存储后处理。}在优化阶段采用 FP16 半精度训练，优化完成后以 FP16 格式存储所有可学习参数，再应用 zip delta 压缩算法（一种无损压缩，通常额外降低约 10\% 存储开销）。

\textbf{压缩效果与定位。}三种策略联合使用，在 Technicolor 和 Neural 3D Video 数据集上分别实现约 190 倍和 125 倍的存储压缩。在 Technicolor 上，MEGA 相比原始 4DGS 的 PSNR 提升 1.2 dB、DSSIM 改善 0.01、LPIPS 改善 0.018，同时渲染速度提升 50\%。回顾本节开头的 Incoherent Deformation, Not Capacity——它用 EER 和 GAD 从优化层面约束变形一致性，但保留了 4D 高斯基元的完整参数结构。MEGA 采取了正交的策略：不约束变形行为，而是重新设计基元本身的参数化（DC-AC 解耦）和时空作用范围（变形预测器 + 不透明度熵），从表示结构层面消除冗余。两方面的改进可以叠加：优化层面的平滑正则化确保变形的泛化性，表示层面的紧凑设计压缩存储开销。

\paragraph{结构化变形图与刚性约束}

前述 MEGA 从 4D 高斯基元的参数化层面消除冗余（DC-AC 解耦、变形预测器），但它仍然保留了"逐高斯独立存储"的基本范式——每个高斯都有自己的一套参数，高斯之间不共享任何信息。这意味着 MEGA 的压缩上限受制于基元数量本身：即使每个高斯的参数再紧凑，数百万个高斯的总量仍然带来存储膨胀。一个自然的追问是：能否让多个高斯共享一部分参数，从根本上减少需要独立存储的单元数量？SD-GS（2025.07）~\cite{2507_07465}（Structured Deformable 3D Gaussians）正是沿着这一方向，将静态场景中已被验证高效的 Scaffold-GS 锚点层级表示引入动态场景重建——用少量结构化锚点替代海量独立高斯，每个锚点管理其局部时空邻域内的多个神经高斯，从表示拓扑层面压缩存储。

\textbf{为什么锚点结构适合动态场景。}在讨论具体设计前，值得先理解为什么锚点结构在动态场景中特别有效。显式 4D 高斯方法（如前述 MEGA 所优化的原始 4DGS）将每个高斯扩展到 4D 时空体，渲染速度快，但无法利用时空相关性——相距很远的时间步上的同一物体表面需要完全独立的高斯集合，导致存储随时间步线性增长。隐式变形方法（如 Deformable 3DGS、4DGS）通过变形场共享时空信息，表示更紧凑，但逐高斯计算变形带来显著的计算开销，且 MLP 难以建模快速非线性运动。SD-GS 的核心洞察是：在时空局部邻域内，高斯之间往往具有高度相似的属性（同一物体表面附近的点颜色、运动模式相近），因此可以用一个锚点统一管理多个局部高斯，锚点之间通过变形场建模全局运动——这既利用了局部冗余实现压缩，又保留了显式表示的渲染速度优势。

\textbf{组件一：规范高斯模型（Canonical Gaussian Model）$\mathcal{M}$。}SD-GS 将场景表示从独立 3D 高斯替换为 Scaffold-GS 式的锚点层级结构。每个锚点由局部上下文特征 $\mathbf{f}_v \in \mathbb{R}^{32}$、3D 位置 $\mathbf{p} \in \mathbb{R}^3$、缩放因子 $\mathbf{s} \in \mathbb{R}^6$、旋转四元数 $\mathbf{q} \in \mathbb{R}^4$ 和 $k$ 个可学习偏移 $\mathbf{o} \in \mathbb{R}^{k \times 3}$ 组成。缩放因子的前三个维度与偏移共同决定 $k$ 个神经高斯的位置：$\{\mathbf{p}_0, \cdots, \mathbf{p}_{k-1}\} = \mathbf{p} + \mathbf{s}_{1:3} \odot \{\mathbf{o}_0, \cdots, \mathbf{o}_{k-1}\}$，后三个维度支持各向异性缩放。旋转四元数 $\mathbf{q}$ 主要用于视锥体计算（推理时仅处理视锥体内的锚点以加速渲染）。与原始 Scaffold-GS 不同的是，SD-GS 用 COLMAP 稀疏点云初始化锚点后，先在粗训练阶段用所有帧的图像训练一个纯静态的规范模型（不引入变形），持续 3000 次迭代，目的是学习高质量的锚点空间分布——这为后续动态建模提供了稳定的几何骨架。

\textbf{组件二：锚点变形场（Anchor Deformation Field）$\mathcal{F}$。}粗训练阶段得到的规范模型 $\mathcal{M}$ 编码了场景的静态几何结构。在精训练阶段，SD-GS 引入锚点变形场将规范模型变换到任意时间戳 $t$ 下的变形模型。具体而言，时空结构编码器 $H$ 从规范模型和时间戳中联合提取时空特征 $\mathbf{f}_d = H(\mathcal{M}, t)$，随后多头锚点变形解码器 $\mathcal{D}$ 通过三个独立 MLP 分别预测锚点的位置、缩放和旋转变形：

\begin{equation}
\Delta = (\Delta \mathbf{x}, \Delta \mathbf{s}, \Delta \mathbf{q}) = (\psi_x(\mathbf{f}_d),\; \psi_s(\mathbf{f}_d),\; \psi_q(\mathbf{f}_d))
\end{equation}

将变形叠加到规范模型上得到变形模型 $\mathcal{M}' = \Delta + \mathcal{M}$。三个独立 MLP 的设计各有侧重：消融实验表明，位置变形对拟合动态场景最为关键（物体运动的核心），缩放变形模拟非刚性运动中表面的拉伸和扭转，旋转变形则建模锚点可见性随运动的变化。变形在锚点层级而非逐高斯层级进行，这正是存储效率的来源——$N_a$ 个锚点的变形参数远少于 $N_a \times k$ 个独立高斯的变形参数。

\textbf{组件三：时间感知神经高斯生成网络（Neural Gaussians Generation Network）$\mathcal{N}$。}得到变形后的锚点后，需要从中解码出 $k$ 个神经 3D 高斯的具体属性（不透明度、颜色、旋转、缩放）。SD-GS 的关键改进在于将时间信息同时注入锚点变形场和神经高斯生成网络——条件向量由锚点特征 $\mathbf{f}_v$、视角方向 $\mathbf{d}_{vc}$ 和时间嵌入 $\gamma(t)$ 组合而成 $[\mathbf{f}_v, \mathbf{d}_{vc}, \gamma(t)]$，送入四个独立全连接网络 $F_\alpha, F_c, F_r, F_s$ 分别解码每个神经高斯的不透明度、颜色、旋转和缩放。这意味着同一锚点在不同时间戳下可以生成属性不同的神经高斯，赋予模型逐高斯级别的时变表达能力，而非仅依赖锚点层级的全局变形。消融实验验证了这种"双重时间注入"设计的有效性：仅在锚点变形场中注入时间信息时，动态区域的渲染质量明显低于同时在生成网络中注入时间信息的版本。

\textbf{组件四：变形感知致密化策略（Deformation-aware Densification）。}锚点层级表示的一个关键问题是锚点分配：场景的哪些区域需要更多锚点？标准 Scaffold-GS 按梯度均值 $\bar{g} = \sum_{i=1}^N \|\nabla_{2D}^i\| / N$ 决定锚点增长，但这一策略在动态场景中存在系统性偏差——动态区域的瞬态运动因出现时间短，累积的梯度被分母 $N$（总迭代次数）稀释，无法达到增长阈值，导致锚点不足；而静态背景的持续存在使其稳定累积梯度，反而获得过多锚点。SD-GS 的核心改进是以锚点的变形幅度对梯度进行加权，使梯度信号与运动幅度正相关：

\begin{equation}
\bar{g} = \frac{\sum_{i=1}^N w_i \|\nabla_{2D}^i\|}{\sum_{i=1}^N w_i}, \quad w_i = \|\Delta \mathbf{x}_i\|_{s_x} + \|\Delta \mathbf{s}_i\|_{s_s} + \|\Delta \mathbf{q}_i\|_{s_q}
\end{equation}

其中 $s_x, s_s, s_q$ 分别为位置、缩放和旋转变形的归一化参考值（取每次迭代中各变形类型的第 90 百分位值，以突出显著变形的锚点），$\|\Delta \mathbf{x}_i\|$ 和 $\|\Delta \mathbf{s}_i\|$ 为变形前后位置和缩放参数的欧氏距离，旋转变形幅度通过四元数点积计算：

\begin{equation}
\|\Delta \mathbf{q}_i\| = 2\arccos\left(\text{clip}\left(\left|\sum_{k=1}^4 q_{\text{orig}}^{(k)} q_{\text{def}}^{(k)}\right|,\; 0,\; 1\right)\right)
\end{equation}

这里取绝对值确保旋转方向不变性，裁剪到 $[0,1]$ 避免浮点误差导致的数值不稳定，$2\arccos(\cdot)$ 将四元数点积映射到 $[0, \pi]$ 的旋转角。直观理解：如果一个锚点在位置、缩放和旋转上均发生了大幅变形（即它在建模动态运动），那么它周围的神经高斯更有可能欠重建，因此应该获得更大的梯度权重来触发锚点增长。实验可视化证实了这一效果——传统致密化策略在静态背景和动态区域保持相似的锚点密度，而 SD-GS 的锚点密度在动态区域显著高于静态区域。

\textbf{训练与优化。}训练损失函数为 $L = \mathcal{L}_1 + \lambda_{\text{SSIM}} \mathcal{L}_{\text{SSIM}} + \lambda_{\text{tv}} \mathcal{L}_{\text{tv}} + \lambda_{\text{vol}} \mathcal{L}_{\text{vol}}$，其中 $\lambda_{\text{SSIM}} = 0.2$，$\lambda_{\text{tv}} = 0.01$（基于网格的时空总变差正则化，平滑变形场），$\lambda_{\text{vol}} = 0.01$（体积正则化，防止高斯过度膨胀）。粗训练 3000 次迭代后进入精训练阶段，总计 140k 次迭代，在单张 RTX 3090 上完成。

\textbf{压缩效果与定位。}在 N3V（Neural 3D Video）数据集上，SD-GS 相比先前 SOTA 方法平均减少约 65\% 存储空间，FPS 提升约 100\%（即 2 倍渲染加速），同时 PSNR、SSIM、LPIPS 指标保持或超越基线。与同为显式方法的 Real-time 4DGS 相比，SD-GS 实现了 362 倍存储压缩和 6.4 倍训练加速，同时 FPS 仍提升约 82\%。在 HyperNeRF 数据集上，SD-GS 同样在渲染速度、训练时间和存储需求上取得最优，对快速运动区域（如手部运动、刀具操作）的重建质量明显优于已有方法。回顾本节的方法脉络：MEGA 从基元参数化层面（DC-AC 解耦）压缩每个 4D 高斯的存储开销，但未改变"逐高斯独立"的表示拓扑；SD-GS 则从表示拓扑层面引入锚点层级结构，让多个局部高斯共享锚点参数，从根上减少需要独立存储的单元数量。两者的压缩机制是正交的——理论上可以在锚点框架内进一步对每个锚点的参数做 DC-AC 式解耦，实现叠加压缩。

\paragraph{动态感知锚点增长与运动引导致密化}

前述 SD-GS 将 Scaffold-GS 的锚点结构引入动态场景，用变形幅度加权梯度来引导锚点向高动态区域生长——它解决的是"哪些区域有运动"的问题。但 SD-GS 的锚点仍停留在 3D 空间（位置 $\mathbf{p} \in \mathbb{R}^3$），锚点增长策略也仅按变形幅度这一个维度加权，并未考虑动态元素在时间轴上的覆盖特性——一个仅在视频中闪现 0.5 秒的物体与一个持续运动的物体，即使变形幅度相似，其锚点需求截然不同。4D Scaffold Gaussian Splatting（2024.11）~\cite{2411_17044} 沿着这条路线走得更远：它将锚点直接提升到 4D 时空，并围绕时间覆盖特性重新设计锚点增长策略，从根本上解决"哪些高斯仅短暂出现"这一更细粒度的问题。

\textbf{从 3D 锚点到 4D 锚点。}与 SD-GS 在 3D 空间构建锚点不同，4D Scaffold GS 直接在 4D 时空体素网格上构建锚点。从 SfM 获取静态点云后，将其体素化为网格 $V \in \mathbb{R}^{N \times 3}$，每个体素中心 $(x_v, y_v, z_v)$ 配合初始时间 $t_0$ 构成 4D 锚点位置 $\mathbf{p} = (x_v, y_v, z_v, t_0)$，并附带压缩特征向量 $\mathbf{f} \in \mathbb{R}^C$ 和 $K$ 个 4D 偏移 $\{\mathbf{o}_k \in \mathbb{R}^4\}$。共享 MLP 从锚点特征解码 $K$ 个神经 4D 高斯的属性：基础不透明度 $\alpha$、四元数 $\mathbf{q}$、缩放 $\mathbf{s}$、视角相关颜色 $\mathbf{c}$、时间尺度 $\sigma$ 以及神经速度 $\mathbf{u} \in \mathbb{R}^3$。这一 4D 锚点设计的关键优势在于：锚点特征 $\mathbf{f}$ 是被多个神经 4D 高斯共享的压缩表示，而非逐高斯独立存储——这正是该方法"不减少高斯数量，而压缩高斯特征的存储开销"的核心思路。当渲染时间 $t_r$ 的帧时，第 $k$ 个 4D 高斯被切片为 3D 高斯：位置 $\boldsymbol{\mu}_k = \mathbf{x}_k^{xyz} + (t_r - x_k^t)\mathbf{u}$，不透明度 $\alpha_k = \alpha_k \cdot g(t_r, x_k^t, \sigma_k)$，然后送入标准 3DGS 光栅化管线。

\textbf{紧凑时间参数化：线性运动 + 广义高斯时间不透明度。}每个神经 4D 高斯采用极简的时间参数化。时变位置通过线性运动建模 $h(t, x_k^t, \mathbf{u}) = (t - x_k^t)\mathbf{u}$，仅需 3 个速度参数。时间不透明度则采用广义高斯函数而非标准高斯函数：

\begin{equation}
g(t, x_k^t, \sigma_k) = \exp\left(-(|t - x_k^t|/\sigma_k)^\gamma\right)
\end{equation}

其中 $\gamma > 2$ 控制激活曲线的陡峭程度（实践中设 $\gamma = 2^\nu$，$\nu$ 为超参数）。当 $\gamma = 2$ 时退化为标准高斯衰减；$\gamma$ 越大，曲线在激活区间内越平坦、在边界处越陡峭，更接近理想阶跃函数，从而更好地拟合真实场景中物体突然出现/消失的"分段持续性"模式——一个人走进画面后在某个区域停留一段时间再离开，其时间不透明度应近似方波而非钟形曲线。实践中将 $1/\sigma$ 设为可学习参数（比直接优化 $\sigma$ 更稳定）。

\textbf{核心创新：时间覆盖感知锚点增长。}至此，4D 锚点框架和紧凑时间参数化只是搭建了表示结构，而锚点增长策略才是 4D Scaffold GS 区别于 SD-GS 的关键。SD-GS 用变形幅度 $w_i = \|\Delta \mathbf{x}_i\| + \|\Delta \mathbf{s}_i\| + \|\Delta \mathbf{q}_i\|$ 对梯度加权——这回答了"运动幅度大的区域需要更多锚点"。但它未考虑时间维度：一个仅出现 5 帧的高速物体和一个出现 100 帧的慢速物体，如果变形幅度相似，SD-GS 会给它们分配相似的梯度权重，导致短时动态元素被严重忽视。

4D Scaffold GS 的核心观察是：原始 Scaffold-GS 的梯度累积 $\bar{g} = \sum^N \|\nabla_{2D}\| / N$ 对所有帧一视同仁，但 4D 高斯仅在时间不透明度大于零时才参与渲染，短时出现的高斯在大部分迭代中处于"休眠"状态，其梯度被分母 $N$ 严重稀释。4D Scaffold GS 以时间覆盖加权的梯度累积替代：

\begin{equation}
\bar{g} = \frac{\sum^N w(\sigma', \sigma)\|\nabla_{2D}\|}{\sum^N w(\sigma', \sigma)}, \quad w(\sigma', \sigma) = \sigma' \cdot (1/\sigma)^\beta
\end{equation}

其中 $\sigma' = g(t_r, x^t, \sigma)$ 是高斯在当前渲染时间 $t_r$ 的时间不透明度分量，$\sigma$ 是该高斯的时间尺度参数，$\beta$ 为超参数（实践中设为 1）。对比 SD-GS 的变形幅度加权，这里的时间覆盖加权从完全不同的角度切入——不是"运动多剧烈"，而是"高斯被激活了吗，它的作用时间有多短"。

具体而言，这一加权带来两个关键机制：(1) \textbf{激活门控}——仅当高斯的时间不透明度 $\sigma' > 0$ 时（即高斯在当前时间步被激活），梯度才被累积。对于在大部分时间步中处于休眠状态的短时高斯，这一门控精确捕获其在少数活跃帧上的重建误差，避免被大量"非活跃"迭代的零梯度稀释。(2) \textbf{短时补偿}——权重中的 $(1/\sigma)^\beta$ 项使时间尺度 $\sigma$ 越小（作用范围越短）的高斯获得越大的梯度权重。直觉上，一个仅在视频中存在 5 帧的物体与一个存在 100 帧的物体，即使单位时间内的重建误差相同，前者的梯度总和天然小 20 倍——$(1/\sigma)^\beta$ 正是为了补偿这一劣势。新锚点被放置在梯度超过阈值的体素中心处，配合训练后剪枝所有高斯不透明度为负的无效锚点，以及缓存 MLP 的时间不变和视角不变输出以加速推理。在 Neural 3D Video 数据集上，4D Scaffold GS 在动态区域的重建质量上以显著优势超越所有基线方法，同时维持实用的存储开销。

\paragraph{级联时间残差渐进变形}

前述 MEGA 从表示结构层面压缩 4D 高斯的参数冗余——用 DC-AC 解耦替代 4D 球谐系数，用变形预测器扩大每个高斯的时空作用范围。但 MEGA 的变形预测器 $F_d$ 对所有时间步使用同一个 MLP，隐含假设是变形模式在整个视频中相对均匀。然而真实动态场景往往呈现层级时间结构：全局几何和光照在整个视频中基本稳定，中等尺度的运动模式（如肢体运动、曲面变形）以中等速率演化，而精细变化（如阴影、遮挡边界）则在帧级别波动。当场景包含大幅运动或突然的场景转换时，单一变形场难以同时捕捉慢变和快变分量，导致重建伪影。CTRL-GS（2025.05）~\cite{2505_18306}（Cascaded Temporal Residue Learning，4D Vision Workshop @ CVPR 2025）从这一瓶颈出发，在标准 4D-GS 的变形管线基础上引入级联残差分解——不改变 4D 高斯基元的参数化方式，而是重新设计变形场的组织结构，将其拆分为多个时间尺度的残差之和。

\textbf{锚点回顾：标准 4D-GS 的变形管线。}CTRL-GS 直接构建在标准 4D-GS 的框架之上。回顾其变形流程：时空结构编码器 $H$（多分辨率 HexPlane + 微型 MLP）将 3D 高斯属性与时间戳编码为时空特征，多头变形解码器 $D = \{\varphi_x, \varphi_r, \varphi_s\}$ 分别预测位置、旋转和缩放的变形量，最终变形后的特征为 $(X', r', s') = (X + \delta X, r + \delta r, s + \delta s)$，其中 $\delta X = \varphi_x(X)$ 等为帧级变形。关键点在于：$\varphi_x, \varphi_r, \varphi_s$ 是整个视频共享的单一 MLP，它们必须在一个统一的参数空间中同时建模所有时间尺度的变化——从稳定的背景到剧烈的运动转换。CTRL-GS 的核心改动正是针对这一"单一 MLP 建模所有时间尺度"的局限。

\textbf{核心设计：三级级联残差分解。}CTRL-GS 的核心思想来自残差学习——不直接预测完整的时间信号，而是将其分解为多个尺度的残差之和。具体而言，CTRL-GS 将变形拆分为三个层级：

第一层是视频级常量特征 $(X, r, s)$，即原始 3D 高斯的规范空间属性，在整个视频中保持不变，编码静态几何和外观。第二层是片段级均值，通过为每个时间窗口创建独立的中间 MLP $\varphi_x^T, \varphi_r^T, \varphi_s^T$ 来预测，分别学习位置、旋转和缩放在该窗口内的平均偏移。第三层是帧级变形 $\delta X, \delta r, \delta s$，即原始 4D-GS 的变形解码器输出的逐帧细粒度修正。三层叠加后，最终变形特征为：

\begin{equation}
  (X', r', s') = (X + \Delta_t X + \delta X, \; r + \Delta_t r + \delta r, \; s + \Delta_t s + \delta s),
\end{equation}

其中 $\Delta_t X = \varphi_x^T(X_T)$ 为片段级残差，$\delta X = \varphi_x(X)$ 为帧级变形。最终变形后的 4D 高斯 $G'$ 保留原始的不透明度和球谐系数：$G' = \{X', s', r', \alpha, C\}$。直觉上，这一分解将"一个 MLP 学会所有变化"的难题转化为"三个层次各司其职"——视频层承载静态内容，片段层捕捉中等尺度的运动模式（如一个人从坐到站的整体位移），帧层补充瞬时细节（如衣物的褶皱变化）。片段级 MLP 在不同窗口间共享架构但独立参数，既利用了跨窗口的结构共享，又保留了对每个时间段运动特征的独立建模能力。

\textbf{时间窗口构造与量化插值。}级联残差的前提是将视频划分为时间窗口。论文提供了三种窗口构造方法：(1) 等长时间划分——将时间轴 $[0,1]$ 均匀分为 $N$ 个窗口，窗口长度 $I = 1/N$；(2) 基于光流峰值的自适应划分（N-highest-flow）——用预训练 RAFT 模型计算相邻帧间的光流 $f_{ij}$，选取光流最大的 $N-1$ 个帧对作为窗口边界，使高运动区间成为自然的分段点；(3) 基于光流累积阈值的贪心划分——设阈值 $\tau = \sum f_{ij} / N$，逐帧累加光流，当累积值超过阈值时开启新窗口并重置累加器，实现高运动区间短窗口、低运动区间长窗口的自适应粒度。三种方法中，自适应划分在大幅运动场景上优势最为明显。

窗口划分完成后，时间特征需要同时携带片段级和帧级信息。CTRL-GS 设计了量化时间特征来实现两级之间的平滑插值：

\begin{equation}
T = \left\lfloor\frac{t + 10^{-9}}{I}\right\rfloor \cdot I + q \cdot I
\end{equation}

其中 $t \in [0,1]$ 为归一化时间戳，$I = 1/N$ 为窗口长度，$10^{-9}$ 确保 floor 函数正确执行时间戳量化，$q \in \mathbb{Q}^+$ 为量化系数——一个控制片段级和帧级信息混合比例的超参数。$q$ 越大，时间特征越偏向当前帧在窗口内的精确位置，帧级残差的权重隐式增大；$q$ 越小，时间特征越接近窗口起点，片段级均值主导变形。构造好的时间特征 $T$ 与原始网格特征拼接后输入片段级 MLP $\varphi_x^T, \varphi_r^T, \varphi_s^T$，使片段级预测能够区分同一窗口内的不同时间位置。

\textbf{与 MEGA 的定位对比。}回顾本节开头的 MEGA——它从表示结构层面压缩 4D 高斯的参数冗余（DC-AC 解耦颜色、变形预测器扩大高斯利用率）。CTRL-GS 则从变形场的组织结构层面入手，不改变 4D 高斯基元本身的参数化方式，而是将单一变形 MLP 分解为视频级 + 片段级 + 帧级的三级级联。两者正交互补：MEGA 压缩"每个高斯存多少参数"，CTRL-GS 改善"变形信息如何跨时间组织"。CTRL-GS 在多个已建立的数据集上实现了 SOTA 视觉质量和实时渲染，尤其在当前方法退化最严重的大幅运动、遮挡和精细细节场景上改进最为显著——这正是多尺度残差分解的优势所在：片段级 MLP 专门学习大幅运动的全局偏移，帧级 MLP 再在此基础上微调瞬时细节，避免了单一 MLP 在大幅运动和精细细节之间的建模冲突。

\paragraph{3D-4D 混合高斯表示}

Hybrid 3D-4D Gaussian Splatting（2025.05）~\cite{2505_13215} 从前述方法的 4D 高斯表示中识别出一个被忽视的结构性浪费：标准 4DGS 理论上可以通过在时间轴上设置大尺度 $s_t$ 来让单个 4D 高斯覆盖整个时间跨度（从而表示静态内容），但实际训练中绝大多数高斯的时间尺度都很小。论文分析了完整训练后的 4DGS 模型，发现大部分高斯的 $s_t$ 集中在 0.5 以下——这意味着对静态背景区域，4DGS 被迫分配大量时间尺度很小的 4D 高斯来跨时间步"拼凑"覆盖，每个高斯携带完整的时间参数（时间中心 $\mu_t$、时间尺度 $s_t$、4D 旋转矩阵 $R_{4D}$），但这些参数对静态区域而言完全冗余。MEGA 和 SD-GS 从表示压缩和锚点结构化角度缓解了这一低效，但它们仍假设所有高斯统一使用 4D 参数化。Hybrid 3D-4D GS 提出了一个更根本的问题：如果静态区域不需要时间维度，为什么不直接把它们变成 3D 高斯？

\textbf{动静态自适应分类。}方法的起点是将所有高斯基元初始化为 4D 高斯。在训练过程中，每当致密化阶段执行时（每 100 次迭代），方法逐一检查每个 4D 高斯的有效时间尺度 $\sigma(s_{t,i})$，其中 $\sigma(\cdot)$ 为指数激活函数，$s_{t,i} \in \mathbb{R}$ 为第 $i$ 个高斯的时间尺度参数。若 $\sigma(s_{t,i})$ 超过预设阈值 $\Theta$，则该高斯被分类为静态。阈值 $\Theta$ 的确定基于对完整训练的 4DGS 模型中时间尺度分布的实证分析——论文观察到分布呈双峰结构：大量高斯集中在较小的时间尺度（动态部分），少量高斯具有较大时间尺度（静态部分），$\Theta$ 取在两个峰之间的"谷底"位置。直觉上，较大的时间尺度意味着该高斯在整个时间跨度上几乎均匀可见，没有高频时间变化，因此不需要时间维度的建模能力。

这种分类不是一次性预处理，而是在每次致密化阶段迭代执行。一个高斯可以在早期迭代中保持 4D 状态，随着训练进行其时间尺度逐渐扩大，一旦超过 $\Theta$ 便被转换为 3D。这一渐进过程持续精炼哪些区域真正需要 4D 建模——最终只有包含真实运动（物体移动、形变、出现/消失）的动态区域保留 4D 高斯，而占据场景大部分面积的静态背景（地面、墙壁、天花板）全部转换为 3D。

\textbf{4D 到 3D 的精确转换。}将一个 4D 高斯降维为 3D 高斯需要丢弃其全部时间参数并提取等效的空间参数。4D 高斯的均值 $\boldsymbol{\mu}_{4D} = (\boldsymbol{\mu}_x, \mu_t)$ 中，空间中心 $\boldsymbol{\mu}_x \in \mathbb{R}^3$ 被保留作为 3D 高斯的位置，时间坐标 $\mu_t$ 被丢弃。4D 高斯拥有一个 $4 \times 4$ 旋转矩阵 $R_{4D}$，原则上它可以混合空间轴和时间轴（使高斯在时空中"倾斜"）。但对于跨越整个时间序列的静态高斯，$R_{4D}$ 退化为近似块对角形式：

\begin{equation}
R_{4D} \approx \begin{pmatrix} R_{3D} & \mathbf{0} \\ \mathbf{0}^T & 1 \end{pmatrix}
\end{equation}

其中 $R_{3D} \in SO(3)$ 为纯空间旋转，$\mathbf{0}$ 为三维零向量。转换时仅保留左上角的 $3 \times 3$ 子块 $R_{3D}$，并从中提取等效的单位四元数 $q_{3D} = (w, x, y, z)$：

\begin{align}
w &= \frac{1}{2}\sqrt{1 + \text{tr}(R_{3D})}, \quad x = \frac{R_{3D}(3,2) - R_{3D}(2,3)}{4w}, \\
y &= \frac{R_{3D}(1,3) - R_{3D}(3,1)}{4w}, \quad z = \frac{R_{3D}(2,1) - R_{3D}(1,2)}{4w}
\end{align}

其中 $\text{tr}(\cdot)$ 为矩阵迹运算。空间缩放 $(s_x, s_y, s_z)$ 直接保留。不透明度 $\alpha$ 和球谐颜色系数也保持不变——静态内容不需要时间依赖的颜色更新。转换后的 3D 高斯完全由 $(\boldsymbol{\mu}_x, q_{3D}, s_x, s_y, s_z, \alpha, \mathbf{c})$ 指定，不再依赖时间变量 $t$。在运行时，每个 3D 静态高斯在所有帧中保持恒定，而每个 4D 动态高斯根据当前时间戳条件化计算。

\textbf{分离优化与加速收敛。}分类完成后，3D 和 4D 高斯进入各自独立的优化管线。这一分离带来关键的速度优势：原始 4DGS 在每次训练迭代中，仅更新对当前时间步有显著贡献的那一小部分 4D 高斯——大部分高斯因时间衰减不透明度极低而被剔除。而 Hybrid 3D-4D GS 的 3D 高斯在\textbf{每次}训练迭代中都被更新，无论当前渲染哪一帧——因为它们不依赖时间，始终参与渲染和梯度回传。这使得 3D 高斯的优化效率远高于 4DGS 中对应的静态部分。论文还取消了训练中的不透明度重置（opacity reset）操作——原始 3DGS 用周期性的不透明度重置清除场景中的"浮游"伪影高斯，但在动态场景的短训练周期中，这种重置会破坏联合时空优化的稳定性。由于混合表示天然减少了总高斯数量，不透明度饱和问题得到缓解，无需重置即可实现稳定收敛。整个模型在 10 秒动态场景上约 6K 次迭代即可收敛，而标准 4DGS 通常需要 20K--30K 次迭代。

\textbf{统一渲染管线。}在渲染阶段，4D 高斯在当前时间 $t$ 处被切片为瞬时 3D 高斯（均值 $\boldsymbol{\mu}_{xyz|t}$ 和协方差 $\Sigma_{xyz|t}$ 按标准 4DGS 条件化公式计算），然后与所有 3D 静态高斯合并为单一列表，统一执行投影、瓦片排序和从后到前的 $\alpha$ 合成。3D 和 4D 高斯在渲染过程中不做区分——光栅化器看到的都是标准 3D 高斯。

在 N3V 和 Technicolor 数据集上，Hybrid 3D-4D GS 实现了与纯 4DGS 相当或更优的 PSNR/SSIM，同时训练速度提升约 3--5 倍，内存占用显著降低。回顾本节的方法脉络：MEGA 在 4D 高斯内部做参数压缩（DC-AC 解耦），SD-GS 和 4D Scaffold GS 用锚点结构化减少高斯总数，CTRL-GS 用级联残差分解时间尺度的复杂性——Hybrid 3D-4D GS 则提出了一个更直接的思路：如果静态区域不需要时间维度，就直接移除它，将 4D 高斯降维为 3D 高斯。

\paragraph{尺度感知残差场与自适应时间间隔}

前述 MEGA 从表示结构层面重新设计 4D 高斯基元的参数化（DC-AC 解耦）和时空作用范围（变形预测器），有效压缩了存储冗余。但它的变形预测器仅建模已有高斯的运动，无法处理一种更棘手的场景：物体的出现与消失——新物体突然进入画面或旧物体永久移除时，不存在"可以被预测变形"的已有高斯。此外，MEGA 和大多数基于 HexPlane 的方法在从特征平面提取残差时，仅按高斯中心点坐标做点查询，完全忽略了高斯在空间中占据的是一个椭球区域而非一个点——大高斯和小高斯提取到同样粒度的特征，这在高斯自我分裂后会导致子高斯的残差特征与父高斯严重偏离。SaRO-GS（2024.12）~\cite{2412_06299}（Scale-aware Residual Field）正是从这两个缺口切入：引入寿命机制让高斯可以"生灭"以建模物体的出现和消失，同时设计尺度感知的残差场让不同大小的高斯获取匹配其空间尺度的特征。

\textbf{表示：寿命赋予高斯"生灭"能力。}SaRO-GS 使用一组 4D 高斯基元 $G^{4D}$，每个基元拥有 4D 位置 $\boldsymbol{\mu}^{4D} = (x,y,z,\tau_i)$（含时间位置 $\tau_i$）、初始协方差 $\Sigma^{4D}$、颜色 $c^{4D}$ 和不透明度 $\alpha^{4D}$。与前述方法的关键区别在于：每个高斯额外拥有一个从残差特征预测的寿命 $\lambda_i = F_w(M(G_i^{4D}))$，其中 $M$ 为尺度感知残差场，$F_w$ 为一个小型 MLP。寿命的存在使高斯可以在时间域中自然消亡，而非像 MEGA 那样仅通过不透明度熵来控制高斯的"活跃/淘汰"。高斯的时间活跃度用类高斯状态函数建模：

\begin{equation}
  \omega_i(t) = e^{-k(t - \tau_i)^2 / \lambda_i^2},
\end{equation}

其中 $k=4$。当采样时间 $t$ 偏离 $\tau_i$ 超过 $\lambda_i$ 时，$\omega_i(t)$ 降至约 0.01，表明该高斯几乎不可见。给定采样时间 $t_0$，投影后的 3D 不透明度为 $\alpha_i^{3D} = \alpha_i^{4D} \cdot \omega_i(t_0)$。除不透明度外，其他属性（位置、协方差、颜色）也随采样时间变化——SaRO-GS 用一组 MLP 解码器 $F_\Delta$ 从残差特征和时间偏移量预测投影残差：

\begin{equation}
  \Delta\boldsymbol{\mu}_i(t),\; \Delta\Sigma_i(t),\; \Delta c_i(t) = F_\Delta(f_i,\; t - \tau_i)
\end{equation}

注意这里解码器输入的是 $t - \tau_i$ 而非 $t$ 本身——目的是获取相对于高斯初始时间位置的残差量，而非绝对时间。最终投影后的 3D 属性为初始值加上残差：$\boldsymbol{\mu}_i^{3D} = \boldsymbol{\mu}_i^{4D}[:3] + \Delta\boldsymbol{\mu}_i(t_0)$，$\Sigma_i^{3D} = \Sigma_i^{4D} + \Delta\Sigma_i(t_0)$，$c_i^{3D} = c_i^{4D} + \Delta c_i(t_0)$。

\textbf{尺度感知残差场：椭球投影而非点查询。}上述残差特征 $f_i = M(G_i^{4D})$ 的提取方式是 SaRO-GS 的核心创新。残差场 $M$ 采用 HexPlane 结构，由三个空间平面 $P_{so}$ 和三个时空平面 $P_{st}$ 组成。问题出在空间平面上：高斯在空间上占据一个椭球区域，投影到空间平面（如 $P_{x,y}$）上得到的是一个椭圆区域而非一个点，但标准方法仅按中心点坐标做点查询。这带来两个具体后果：第一，大高斯的残差特征仅反映其中心点附近的局部信息，而非它实际覆盖的整个区域；第二，当大高斯分裂为多个小高斯时，这些子高斯在同一个点上查询特征，但它们的残差特征应与父高斯一致，分裂后却突然发生了偏离。

SaRO-GS 的解决方案是在每个空间平面 $P_{x,y}$ 上构建 MipMap 金字塔。第 0 级 $P_{x,y}^0$ 是分辨率为 $N \times N$ 的特征图，对应最小空间尺度 $\delta_s^0 = (B_{\max} - B_{\min}) / N$；后续层级 $P_{x,y}^l$ 由上一级的特征做 $2\times$ 下采样得到，空间尺度按 $2^l$ 递增。实际训练中只存储和优化第 0 级特征，其余层级在前向推理时动态计算。给定一个缩放为 $(s_x, s_y)$ 的高斯，它在 $P_{x,y}$ 上的投影轴为 $(s_x, s_y)$，对应的空间层级为：

\begin{equation}
  l^x = \log_2(s_x / \delta_s^0), \quad l^y = \log_2(s_y / \delta_s^0), \quad l = \min(\lfloor l^x \rfloor, \lfloor l^y \rfloor)
\end{equation}

取最小值确保不超出高斯的实际覆盖范围。然后在相邻两个层级 $P_{x,y}^{\lfloor l \rfloor}$ 和 $P_{x,y}^{\lceil l \rceil}$ 之间做三线性插值获取特征 $f_{x,y} = \text{tri}(\pi_{x,y}(\boldsymbol{\mu}^{4D}), l;\; P_{x,y}^{\lfloor l \rfloor}, P_{x,y}^{\lceil l \rceil})$。最终残差特征为空间平面和时空平面特征的求和：

\begin{equation}
  f(G^{4D}) = \sum_{(i,j) \in C_{so}} \text{tri}(\pi_{i,j}(\boldsymbol{\mu}^{4D}), l;\; P_{i,j}^{\lfloor l \rfloor}, P_{i,j}^{\lceil l \rceil}) + \sum_{(i,j) \in C_{st}} \text{bi}(\pi_{i,j}(\boldsymbol{\mu}^{4D});\; P_{i,j})
\end{equation}

其中 $C_{so} = \{(x,y),(x,z),(y,z)\}$ 为空间平面组合，$C_{st} = \{(x,t),(y,t),(z,t)\}$ 为时空平面组合。空间平面使用尺度感知的三线性插值（tri），时空平面使用常规的双线性插值（bi）。这使得大高斯自动获取粗粒度特征（高层级），小高斯获取精细特征（低层级），与 3DGS 的自我分裂行为保持一致——分裂后的子高斯因缩放更小而查询更低层级的特征，但该层级是由父高斯的高层级特征下采样得到，保证了父子特征的一致性。

\textbf{自适应优化：按时间覆盖率调整学习率和致密化阈值。}寿命机制虽然赋予了高斯建模物体出现/消失的能力，但也引入了新的优化不平衡问题：动态高斯因寿命短，在观测时间范围内被采样的概率低，反向传播时累积的梯度偏小。如果直接沿用 MEGA 和标准 3DGS 的统一致密化阈值，这些动态高斯将因梯度不足而难以被致密化，导致动态区域重建不充分。SaRO-GS 利用高斯的状态函数 $\omega_i(t)$ 计算其在观测时间范围 $[t_{\text{start}}, t_{\text{end}}]$ 上的积分 $I_i$ 来衡量其采样概率。由于 $\omega_i(t)$ 是类高斯函数，精确积分困难，SaRO-GS 借助近似累积分布函数 $Q(t)$ 来高效计算：

\begin{equation}
  F(t) = \frac{\lambda_i}{\sqrt{k}} Q\!\left(\frac{\sqrt{2k}(t - \tau_i)}{\lambda_i}\right), \quad Q(t) = \frac{1}{\sqrt{2\pi}} \int_{-\infty}^t e^{-x^2/2}\, dx \approx 1 - \frac{1}{e^{1 + \alpha_1 t^3 + \alpha_2 t}}
\end{equation}

其中 $\alpha_1 = 0.0706$, $\alpha_2 = 1.5976$ 为近似常数，$I_i = F(t_{\text{end}}) - F(t_{\text{start}})$。积分越大说明高斯在观测时间范围内被采样的概率越高。基于此，SaRO-GS 动态调整每个高斯的致密化梯度阈值和学习率：

\begin{equation}
  \tau_i = \tau_{\text{base}} \cdot \frac{I_i}{I_{\max}}, \quad \text{lr}_i = \text{lr}_{\text{base}} \cdot \frac{I_{\max}}{I_i}
\end{equation}

阈值 $\tau_i$ 正比于采样概率（动态高斯阈值更低，更容易触发致密化），学习率 $\text{lr}_i$ 反比于采样概率（动态高斯学习率更高，加速收敛）。此外，SaRO-GS 引入缩放残差正则化 $L_{SR} = \frac{1}{n}\sum_i \|\Delta s_i(\tau_i)[:3]\|_2$，约束高斯在自身时间位置处的缩放残差最小化，避免投影过程过度依赖残差修正而忽视初始属性的优化。完整训练损失为 $L = (1-\lambda_1)L_1 + \lambda_1 L_{D\text{-SSIM}} + \lambda_2 L_{SR}$。

\paragraph{稀疏时变属性建模降低 4D 开销}

前述 MEGA 从 4D 高斯基元的参数结构层面解耦颜色属性（DC-AC），压缩了逐高斯存储冗余。但它和 Deformable 3DGS、4DGS 等方法共享一个更深层的效率瓶颈：所有高斯——无论静态还是动态——的时间变属性（位置、缩放、旋转）都要逐高斯、逐时间步通过 MLP 查询，而渲染速度与高斯数量线性相关。换言之，场景中大量静态区域的高斯本不需要任何变形计算，却仍占用 MLP 查询时间和内存。EDGS（2025.02）~\cite{2502_20378}（Efficient Dynamic Gaussian Splatting）正是从这一角度切入：如果 MEGA 解耦的是"哪些属性随视角-时间变化"，那么 EDGS 解耦的就是"哪些高斯随时间变化"——只对真正需要变形的高斯查询时间变属性，静态区域的高斯完全跳过变形计算。

\textbf{锚点网格初始化与属性分离。}EDGS 建立在 Scaffold-GS 的锚点结构之上：将 COLMAP 稀疏点云通过体素化为结构化锚点 $A = \lfloor P / d \rfloor$（$d=0.6$ 为体素大小，floor 操作移除冗余过密的点），每个锚点关联 $K=10$ 个子高斯。关键设计是将高斯属性分离为时间不变属性（不透明度 $\alpha$、颜色 $c$）和时间变属性（位置、缩放、四元数旋转）。时间不变属性直接由两层轻量 MLP（中间层维度 64）从锚点特征 $f_a \in \mathbb{R}^8$ 解码：$\alpha = F_\alpha(f_a)$，$c = F_c(f_a)$，无需任何时间输入，训练和推理过程中保持恒定。

\textbf{时间掩码：无监督静态-动态分离。}EDGS 的核心创新是时间掩码（Time Mask）机制——一个二分类 MLP $F_{\text{mask}}$ 以锚点特征 $f_a$ 为输入，输出 0（静态）或 1（动态）。整个分类过程完全无监督：不需要任何 2D 或 3D 掩码作为监督信号，分类能力纯粹从渲染损失的反向传播中涌现。被分类为静态的锚点直接跳过变形 MLP，其子高斯的位置、缩放和旋转在所有时间步保持不变；只有动态锚点才进入变形网络计算时间变属性。这意味着场景中椅子、墙壁、地面等静态区域的锚点在推理时完全不需要 MLP 查询，渲染速度直接取决于动态锚点的数量。为防止所有锚点都被分类为动态（失去稀疏性），EDGS 引入时间掩码正则化项：

\begin{equation}
  L_{t\text{-mask}} = \sum_{n=0}^{N-1} F_{\text{mask}}(f_a^n) / N
\end{equation}

该正则化鼓励尽可能多的锚点被分类为静态——每个锚点被分为动态时贡献 $1/N$ 的正惩罚，推动 $F_{\text{mask}}$ 在不影响渲染质量的前提下尽可能多地"关掉"锚点的变形计算。消融实验表明，引入时间掩码后不仅加速渲染，还提升了渲染质量：在 Sieve 场景上 PSNR 从 26.18 提升至 27.12 dB，可能原因是通过减少需要优化的锚点数量，变形 MLP 的优化更加充分有效。

\textbf{RBF 核驱动的语义运动传播。}对于被分类为动态的锚点，其子高斯的运动建模面临一个关键选择：刚体假设（锚点与其子高斯运动完全相同）过于严格，KNN 方法按空间距离分配运动权重在碰撞后分离的物体上性能退化。EDGS 的方案是使用径向基函数（RBF）核在特征空间中进行语义运动传播。具体而言，锚点 $a$ 在时间 $t$ 的运动由四层变形 MLP（中间层维度 128）预测：$\Delta x_a^t = F(\text{PE}(x_a, t))$，其中 $\text{PE}$ 为位置编码。其第 $k$ 个子高斯的运动通过 RBF 核加权计算：

\begin{equation}
  \Delta x_o^{k,t} = K(f_a, f_o^k) \cdot \Delta x_a^t, \quad K(f_a, f_o^k) = \exp\!\left(-\frac{\|f_a - f_o^k\|^2}{2\sigma^2}\right),
\end{equation}

其中 $f_a \in \mathbb{R}^8$ 和 $f_o^k \in \mathbb{R}^8$ 分别为锚点和子高斯的可学习特征，$\sigma = 1$ 为 RBF 核带宽。RBF 核根据特征相似性（而非空间距离）分配运动幅度：语义相近的子高斯获得相近的变形，即使空间上因碰撞而分离。子高斯在时间 $t$ 的最终位置为 $x_o^{k,t} = x_a + x_o^k + \Delta x_o^{k,t}$，即锚点位置加原始偏移加 RBF 加权运动量。缩放和旋转的时间变属性通过类似方式由 MLP $F_s$、$F_r$ 预测偏移量 $\Delta s_o^{k,t}$、$\Delta r_o^{k,t}$。

\textbf{锚点致密化与剪枝。}训练中 EDGS 沿用 Scaffold-GS 的梯度致密化策略：每 100 次迭代计算锚点的梯度，超过阈值 $\tau$ 的锚点所在体素被细分为 $m$ 个子体素，每个子体素继承父锚点的特征并初始化 $K$ 个新偏移。剪枝则按每个锚点关联高斯的平均不透明度执行——低于预设阈值的锚点被移除。值得注意的是，每个锚点 $a$ 的空间位置在训练全过程中保持冻结，不参与梯度更新，只有其特征 $f_a$ 和关联偏移 $o_a$ 被优化。

\textbf{整体损失与效果定位。}整体损失函数为 $L = (1-\lambda)L_1 + \lambda L_{\text{SSIM}} + \lambda_t L_{t\text{-mask}}$，其中 $\lambda = 0.2, \lambda_t = 0.2$。在 NeRF-DS 数据集上，EDGS 取得最高平均 PSNR（24.39 dB）和 SSIM（0.8873），在 HyperNeRF 数据集上达到 25.7 dB PSNR、0.860 MS-SSIM，渲染速度 117 FPS，仅需约 7K 个需要查询时间变属性的高斯——在所有对比方法中最少。与 MEGA 的定位对比：MEGA 解耦的是颜色属性内部的静态/动态成分（DC-AC），压缩逐高斯存储；EDGS 解耦的是高斯群体内部的静态/动态成员（Time Mask），减少需要查询变形 MLP 的高斯数量。两者从不同维度压缩 4D 表示——存储层面 vs 计算层面——理论上可以互补叠加。

\paragraph{超点聚类分组与共享变形}

前述 SD-GS 将 Scaffold-GS 的锚点层级表示引入动态场景，从表示拓扑层面实现了存储压缩——锚点之间通过变形场建模全局运动，每个锚点管理其局部邻域内的多个神经高斯。但 SD-GS 的锚点结构沿用了 Scaffold-GS 的设计范式：锚点属性（位置、缩放、旋转）在训练中固定为场景几何的静态骨架，仅通过变形场对其进行全局偏移，锚点内部的神经高斯属性由小型 MLP 从锚点特征解码而来。换言之，SD-GS 的锚点本质上是一种"预定义的层级结构"——锚点数量和位置由致密化策略决定，高斯与锚点的归属关系通过空间邻域隐式确定。这引出一个更灵活的问题：能否让高斯自己"选择"归属哪个聚类，并且让聚类的运动模式而非空间距离成为归属依据？Superpoint Gaussian Splatting（2024.06）~\cite{2406_03697}（SP-GS，ICML 2024）正是从这一角度切入——借用计算机视觉中超像素/超点过分割（over-segmentation）的经典思想，将属性相似（尤其是运动模式相似）的高斯显式聚类为超点，让变形预测在超点级别进行而非逐高斯进行，从而将变形 MLP 的查询次数从高斯总数 $P$ 降低到超点数量 $M$（通常减少一个数量级以上），同时保持与标准 3DGS 相当的渲染速度。

\textbf{为什么超点聚类有效：ARAP 先验。}SP-GS 的核心动机来自 As-Rigid-As-Possible（ARAP）正则化在 3D 重建中的经典洞察：同一刚体上的相邻点在所有时间步应具有相似的平移和旋转变换。标准 Deformable 3DGS 忽略了这一先验，为每个高斯独立查询变形 MLP——当高斯数量达数十万时，绝大多数查询本质上是冗余的（同一物体表面的高斯运动模式几乎相同）。SP-GS 的做法是显式利用这一冗余：将运动模式相似的高斯聚类为超点，只在超点级别预测刚体变换（相对旋转 $\Delta\mathbf{R}_j^t$ 和相对平移 $\Delta\bm{t}_j^t$），组内高斯通过刚体变换继承超点的运动。值得注意的是，SP-GS 假设每个高斯仅经历刚体运动，因此只有位置 $\bm{\mu}$ 和旋转 $\mathbf{R}$ 随时间变化，其他属性（不透明度、缩放、球谐系数）在规范空间中保持不变——这与 SD-GS 锚点变形场同时预测位置、缩放和旋转变形的设计不同。

\textbf{组件一：可学习超点关联。}与 SD-GS 通过空间邻域隐式确定锚点-高斯归属不同，SP-GS 引入可学习关联矩阵 $\mathbf{A} \in \mathbb{R}^{P \times M}$（$P$ 为高斯总数，$M$ 为超点数量）显式建模每个高斯与 $K$ 个最近超点之间的软关联。训练分两阶段：首先进行 3000 次迭代的预热训练（warm-up），在不使用变形网络的情况下让 3D 高斯在规范空间中稳定，获得可靠的初始位置和形状。然后用最远点采样（Farthest Point Sampling）从稳定后的高斯中选取 $M$ 个作为超点中心，其规范位置 $\bm{p}_j^c$ 等于被采样高斯的中心。关联矩阵的初始化带有先验偏好：若第 $j$ 个超点由第 $i$ 个高斯初始化，则 $A_{ij} = 0.9$（强关联），否则 $A_{ij} = 0.1$（弱关联）。关联概率通过 softmax 归一化：

\begin{equation}
  a_{ij} = \frac{\exp(A_{ij})}{\sum_{j' \in \mathcal{N}_i} \exp(A_{ij'})}, \quad j \in \mathcal{N}_i,
\end{equation}

其中 $\mathcal{N}_i$ 为第 $i$ 个高斯的 $K$ 个最近超点集合。每个高斯的最终归属超点为 $j^* = \arg\max_{j \in \mathcal{N}_i} a_{ij}$。训练过程中关联矩阵随梯度更新，使高斯能自主"迁移"到运动模式更匹配的超点——这比 SD-GS 的固定空间归属更灵活。

\textbf{组件二：超点级刚体变形。}变形网络 $\mathcal{F}$ 预测每个超点的相对旋转 $\Delta\mathbf{R}_j^t$ 和相对平移 $\Delta\bm{t}_j^t$，组内高斯通过刚体变换计算变形后的位置和旋转：

\begin{equation}
  \bm{\mu}_i^t = \Delta\mathbf{R}_j^t \cdot \bm{\mu}_i^c + \Delta\bm{t}_j^t, \quad \mathbf{R}_i^t = \Delta\mathbf{R}_j^t \cdot \mathbf{R}_i^c,
\end{equation}

其中上标 $c$ 表示规范空间中的初始值，$j$ 为第 $i$ 个高斯的归属超点。变形网络 $\mathcal{F}$ 为 8 层 MLP、256 个隐藏神经元，输入为超点规范位置 $\bm{p}_j^c$ 和时间步 $t$ 的位置编码（位置 $L=10$，时间 $L=6$），输出超点级别的相对变换：

\begin{equation}
  (\Delta\mathbf{R}_j^t, \Delta\bm{t}_j^t) = \mathcal{F}(\gamma(\bm{p}_j^c), \gamma(t)).
\end{equation}

推理时还可以预计算所有训练时间步的超点变形，对中间帧通过线性插值获取变形参数，避免重复查询 MLP，进一步加速渲染。这使得变形预测的计算量从 $O(P)$ 降低到 $O(M)$——在 D-NeRF 数据集上，超点数量通常仅为高斯数量的 $1/10$ 到 $1/20$。

\textbf{组件三：属性重建损失。}仅靠刚体变换假设不足以保证超点内高斯运动的一致性——如果关联矩阵将运动模式差异较大的高斯分到同一超点，刚体变换假设就会被违反。为此，SP-GS 引入双向属性重建损失来显式约束超点内一致性。具体而言，定义两种重建路径：从高斯到超点的重建 $R_{g \to sp}$，将归属同一超点的所有高斯的属性加权平均回传给超点；以及从超点到高斯的重建 $R_{sp \to g}$，将超点属性按关联概率分配给高斯。属性重建损失衡量原始属性与"高斯 $\to$ 超点 $\to$ 高斯"往返重建后属性之间的差异：

\begin{equation}
  \mathcal{L}_{\bm{v}} = \frac{1}{P}\sum_{i}\|\bm{v}_i - \bm{v}_i'\|, \quad \bm{v}_i' = R_{sp \to g}(R_{g \to sp}(\bm{v}_i)),
\end{equation}

其中 $\bm{v} \in \{\bm{\mu}^t, \Delta\mathbf{R}^t, \Delta\bm{t}^t\}$ 分别对位置、旋转和平移施加约束。直觉上：同一超点内高斯的属性越相似，往返重建的误差越小，该损失就越小——这迫使关联矩阵将运动模式相似的高斯聚合到同一超点，充分利用 ARAP 先验。旋转使用李代数 $\mathfrak{se}(3)$ 表示以支持线性插值。整体损失为 $\mathcal{L} = \mathcal{L}_{img} + \sum_{\bm{v}} \lambda_{\bm{v}} \mathcal{L}_{\bm{v}}$，其中 $\lambda_{\bm{\mu}^t} = 10^{-3}$，$\lambda_{\Delta\mathbf{R}^t} = \lambda_{\Delta\bm{t}^t} = 1$。

\textbf{可扩展性：可选非刚体变形网络。}刚体变换假设在物体发生弹性变形或局部非刚性运动时会成为瓶颈。SP-GS 提供了一个可选的非刚体变形网络 $\mathcal{G}$（3 层 MLP、64 个隐藏神经元），在超点级刚体变换之上为每个高斯额外预测逐高斯的残余变形 $(\hat{\Delta\mathbf{R}}_i^t, \hat{\Delta\bm{t}}_i^t) = \mathcal{G}(\gamma(\bm{\mu}_i^t), \gamma(t))$。最终位置和旋转为两级变形的复合：

\begin{equation}
  \bm{\mu}_i^t = \hat{\Delta\mathbf{R}}_i^t(\Delta\mathbf{R}_j^t \bm{\mu}_i^c + \Delta\bm{t}_j^t) + \hat{\Delta\bm{t}}_i^t, \quad \mathbf{R}_i^t = \hat{\Delta\mathbf{R}}_i^t \Delta\mathbf{R}_j^t \mathbf{R}_i^c.
\end{equation}

该版本（SP-GS+NG）用纯刚体版本的预训练模型初始化，额外训练 20000 次迭代，仅使用图像重建损失 $\mathcal{L}_{img}$。此外，超点结构还天然支持场景编辑——同一超点内的高斯属于同一刚体部件，可以整体移动或变换，提供比逐高斯操作更有语义意义的操控粒度。

\begin{takeaway}
变形场与 4D 高斯表示是动态 3DGS 的两大核心路线：变形场在规范空间 3DGS 上叠加时间条件的变形网络，成熟且兼容静态管线；原生 4D 表示直接在时空域定义基元，自然支持连续时间建模。HexPlane 张量分解和 4D 哈希编码将存储开销降低数倍到数十倍，前馈方法（DGS-LRM、L4GM）通过大规模预训练实现了秒级 4D 重建。
\end{takeaway}

\textbf{效果与定位。}在 D-NeRF 合成数据集上，SP-GS 以 $800 \times 800$ 分辨率渲染达到 227 FPS，SP-GS+NG 达到 119 FPS，同时 PSNR、SSIM 和 LPIPS 与先前 SOTA 方法相当或更优。在 HyperNeRF 真实数据集上以 $536 \times 960$ 分辨率达到 117 FPS。与 D-3D-GS（Deformable 3DGS，为每个高斯独立查询变形 MLP）相比，SP-GS 的渲染速度是其 5 倍，因为变形查询次数从数十万降低到数千。回顾本节的方法脉络：MEGA 从基元参数化层面解耦颜色属性（DC-AC）来压缩逐高斯存储；SD-GS 从表示拓扑层面引入锚点层级结构来共享参数；SP-GS 则从计算效率层面将变形预测从逐高斯提升到超点级别——三者的压缩机制分别作用于存储、拓扑和计算三个正交维度。SP-GS 与 SD-GS 虽然都采用"高层级代理管理低层级高斯"的思想，但实现路径不同：SD-GS 的锚点继承自 Scaffold-GS 的空间层级结构，锚点-高斯归属由空间邻域隐式确定；SP-GS 的超点来自运动模式的显式聚类，归属关系通过可学习关联矩阵端到端优化。

%=====================================================================

\subsection{近期方法脉络：动态表示设计、静动态分解、动态3DGS范式比较等}

这一组近期工作延续 dynamic 3DGS、Deformable 3DGS、dynamic Gaussian splatting，主要把改动集中在增量改进、正则化、应用、扩展等环节。下面只保留相对既有路线的变化，便于和前文方法直接对照。

\paragraph{Multi4D}
Multi4D（2026.06）~\cite{2606_22197} 延续 dynamic 3DGS 在“动态表示设计”上的基本设定，主要改变是：Multi4D：多层竞争分配，将建模容量分配到不同层级以平衡形变和4D primitive路线。高斯基元仍承担可渲染动态场景表示。
其设计动机是：提出非单一表示的容量分配机制，把不同动态因素交给不同层级建模，避免形变方法过平滑和4D primitive过参数化。
具体而言，框架在多层表示间竞争性分配建模容量，保持对象身份和时间一致的同时捕获高频动态细节。
适用边界是：多层分配策略复杂；极端长序列和多对象交互仍可能需要额外约束。

\paragraph{DSD-GS}
DSD-GS（2026.05）~\cite{2605_30863} 延续 dynamic 3DGS 在“静动态分解”上的基本设定，主要改变是：DSD-GS：用前馈Gaussian Splatting编码器和光流模型进行动态-静态分解，只对动态区域建模时变分量。高斯基元仍承担可渲染动态场景表示。
其设计动机是：提出基于前馈GS编码器和光流的动态/静态分解策略，把静态区域从昂贵的动态建模中剥离，提升训练和渲染效率。
具体而言，方法先估计场景中静态与动态区域，再用静态高斯承载不变部分，用动态分支和光流约束处理运动区域，从而避免所有高斯都携带时间相关参数。
适用边界是：分解质量依赖光流与前馈编码器；复杂遮挡和非刚性运动可能导致静动态边界错误。

\paragraph{Beyond Static Gaussians}
Beyond Static Gaussians（2026.05）~\cite{2606_00452} 可作为“动态3DGS范式比较”中的新锚点候选。相较于 dynamic 3DGS，主要改变是：动态3DGS经验研究，将方法分为结构引导和高斯中心两类并在D-NeRF上评估。高斯基元仍承担可渲染动态场景表示。
其设计动机是：提供动态3DGS方法的系统分类和实证比较，帮助读者理解形变场、规范空间、网格/4D高斯等不同路线的取舍。
具体而言，论文选取代表性动态3DGS方法，按是否使用辅助结构表示组织，对重建保真度、模型大小、渲染速度和训练特性做横向比较。
这一锚点判断来自当前分类中的机制关系：候选锚点：路线变化，且相关性 0.80；后续方法可围绕 比较动态3DGS架构范式 继续比较和扩展。
适用边界是：评估范围受所选方法和D-NeRF数据集限制；不是新建模方法。

\paragraph{SemDynReg}
SemDynReg（2026.06）~\cite{2606_28656} 延续 Deformable 3DGS 在“Semantic regularization for dynamic Gaussians”上的基本设定，主要改变是：用 SAM 分割和语义特征约束对象级一致形变。高斯基元仍承担可渲染动态场景表示。
其设计动机是：把动态形变约束从纯局部场提升到语义对象层面，改善动态 Gaussian 场景的时序和结构一致性。
具体而言，先用 SAM 提取分割 mask 并获得语义特征，再将其用于形变场正则，使同一对象内 Gaussian 运动更一致，不同对象间耦合更少。
适用边界是：依赖初始分割和语义特征质量；遮挡、细粒度非刚性和错误 mask 会影响正则效果。

\paragraph{TIDES}
TIDES（2026.06）~\cite{2606_02058} 延续 dynamic Gaussian splatting 在“事件与连续时间建模”上的基本设定，主要改变是：TIDES：基于dynamic Gaussian splatting的连续时间事件模拟器，从显式3D场景和运动直接计算像素强度导数。高斯基元仍承担可渲染动态场景表示。
其设计动机是：提出基于动态Gaussian splatting的事件仿真框架，避免传统帧序列差分造成的timestamp batching，并能处理遮挡下的多次阈值穿越。
具体而言，方法在显式3D高斯场景中学习几何和运动，直接推导每个像素的强度时间导数，预测事件相机阈值交叉时间，而无需时间上采样。
适用边界是：依赖动态重建质量；事件噪声、传感器非理想性和复杂材质可能需要额外建模。

\paragraph{Learning Efficient 4D Gaussian Representations from Monocular Videos with Flow Splatting}
Learning Efficient 4D Gaussian Representations from Monocular Videos with Flow Splatting（2026.06）~\cite{2606_29976} 延续 4D Gaussian Splatting 在“Flow-guided 4D Gaussian representations”上的基本设定，主要改变是：Flow Splatting：构建 velocity field 并用 splatting 渲染 optical flow。高斯基元仍承担可渲染动态场景表示。
其设计动机是：把光流预测纳入 splatting 表示，使 4D Gaussian 不只依赖 deformation field、轨迹或逐帧体积。
具体而言，构建 velocity field，用传统 splatting 技术渲染 optical flow，利用视频中的 dense motion information 约束动态 3D 场景重建。
适用边界是：光流估计错误、遮挡和非刚性大变形会影响速度场学习；复杂长视频仍可能累积误差。

\paragraph{Liquid Neural Networks as a Drop-in Continuous-Time Deformation Field for Dynamic 3D Gaussian Splatting}
Liquid Neural Networks as a Drop-in Continuous-Time Deformation Field for Dynamic 3D Gaussian Splatting（2026.06）~\cite{2606_07670} 延续 D-3DGS 在“连续时间形变场”上的基本设定，主要改变是：用Closed-form Continuous-time Liquid Neural Network替代D-3DGS中的时间位置编码MLP形变场。高斯基元仍承担可渲染动态场景表示。
其设计动机是：将形变场架构从独立预测每个时间的MLP改为带时间门控的闭式连续时间单元，在损失景观中显式注入平滑时间响应。
具体而言，保留D-3DGS其他管线，仅替换deformation field为CfC cells；每个cell通过sigmoidal time gate在候选隐状态之间插值，近似液体时间常数ODE的闭式解。
适用边界是：LNN表达能力和训练稳定性需调参；对突变运动可能过度平滑。

\paragraph{Temporally Aware Densification for Dynamic 3D Gaussian Splatting}
Temporally Aware Densification for Dynamic 3D Gaussian Splatting（2026.06）~\cite{2606_23212} 延续 dynamic 3DGS 在“动态致密化”上的基本设定，主要改变是：Visibility-Aware Densification(VAD)与Temporally-Adaptive Thresholding(TAT)，按高斯实际时间可见性调整致密化。高斯基元仍承担可渲染动态场景表示。
其设计动机是：指出动态场景不能直接继承静态ADC策略，提出按时间存在和可见性决定高斯分裂/致密化阈值。
具体而言，VAD统计高斯在时间维度中的实际可见性，TAT为每个高斯自适应调整densification threshold，以避免短生命周期区域监督不足。
适用边界是：需要可靠时间可见性估计；快速遮挡和噪声运动可能影响阈值。

\paragraph{MVFusion-GS}
MVFusion-GS（2026.07）~\cite{2607_01578} 延续 Dynamic 3DGS 在“Motion-aware temporal attention for dynamic Gaussians”上的基本设定，主要改变是：Motion-Variance Guided Refinement 与 temporal coherence attention。高斯基元仍承担可渲染动态场景表示。
其设计动机是：补足现有 deformation networks 缺少显式运动感知的问题，联合建模长期运动强度和短期 temporal coherence。
具体而言，对每个 Gaussian 聚合 motion variance 以引导 refinement，并利用 temporal attention 捕捉局部时间一致性，减少前景形变误差和 pseudo-static background residuals。
适用边界是：运动估计和注意力窗口选择可能影响结果；复杂遮挡和长周期运动仍需验证。

\paragraph{Progressive Pose-Guided 4D Animal Reconstruction from Monocular Video}
Progressive Pose-Guided 4D Animal Reconstruction from Monocular Video（2026.06）~\cite{2607_00157} 延续 3DGS 在“Pose-guided 4D Gaussian reconstruction beyond humans”上的基本设定，主要改变是：基于 3DGS 的渐进式 test-time optimization，解耦 articulated pose 和 non-rigid deformation。高斯基元仍承担可渲染动态场景表示。
其设计动机是：在缺少可靠动物模板的情况下，用渐进策略连接粗形状先验、关节姿态和非刚性变形。
具体而言，从单目视频出发，先获得粗形状和姿态引导，再逐步优化 Gaussian 表示，使 articulated pose 与 non-rigid deformation 分离并提升外观保真。
适用边界是：依赖粗形状先验和姿态估计；极端动物形态、遮挡和毛发细节仍具有挑战。

\paragraph{Lift4D}
Lift4D（2026.06）~\cite{2606_23688} 延续 deformable 3DGS 在“单目4D先验”上的基本设定，主要改变是：Lift4D：用 causal latent conditioning 生成时序一致单帧 3D 先验，并优化 canonical 3D Gaussians 与形变参数。高斯基元仍承担可渲染动态场景表示。
其设计动机是：针对单目 4D 重建中仅初始化使用先验或直接预测 4D 表示的局限，将时序一致的单视图 3D 预测持续接入可形变 3DGS 优化。
具体而言，方法把每帧单视图 3D 估计作为初始化和约束，随后在 canonical 3D Gaussians、稀疏形变控制节点和遮挡感知重建损失下进行测试时优化，并用 view-conditioned diffusion prior 补全不可见区域。
适用边界是：依赖单视图 3D/SAM3D 初始化质量；遮挡严重、快速非刚性变形和长视频仍可能带来漂移或补全误差。
\section{4D扩展与特殊动态重建}
\label{sec:06_2}

本节讨论 4D 高斯表示在物理约束、事件相机、前馈重建、动态表面重建和运动建模等方向的扩展应用。

\subsection{物理约束与事件驱动动态重建}

前面讨论的方法主要追求视觉保真度，但变形本身可能是物理上不合理的——物体会穿越、拉伸或违反质量守恒。本节将物理先验引入变形场，让重建结果不仅"看起来对"，而且"物理上对"。多项研究将物理先验引入变形场，使重建结果不仅视觉逼真，而且物理合理。

\paragraph{物理先验约束的变形场}

PIDG（2025.11）~\cite{2511_06299}（Physics-Informed Deformable Gaussian Splatting）修改了动态 3DGS 的变形场模块——用物理定律约束的连续介质力学变形替代纯数据驱动的 MLP 变形——渲染管线和哈希编码框架保持不变。针对纯 MLP 变形场忽略材料本构律、仅靠光度损失无法约束三维粒子运动的问题，PIDG 将每个高斯建模为携带速度与应力张量的拉格朗日粒子，通过 Cauchy 动量残差显式约束动量平衡。然而，这一物理约束存在内在局限：依赖事先指定的本构律（流体/弹性/刚体）而非从数据中发现，对多材料耦合界面的应力连续性建模能力有限；Cauchy 约束本质上是"惩罚偏离"而非保证守恒结构，无法严格保证能量守恒；在光流监督失效的场景（遮挡、大运动、快速形变）下，物理约束缺乏独立的观测锚定。

\textbf{Baseline 问题。}已有的变形场方法（如 Deformable 3DGS）用 MLP 从时空特征直接预测高斯基元的平移、旋转和缩放增量，仅依赖单目视频的光度损失（$L_1 + L_\text{D-SSIM}$）进行端到端监督。这种纯数据驱动的做法存在两个根本缺陷：(1) 将运动简化为刚体变换，忽略了不同材料（流体、弹性体、刚体）各自的运动本构律（constitutive law，即描述材料应力-应变关系的物理定律）；(2) 仅靠二维视觉投影无法统一约束三维粒子的运动状态，导致高斯粒子偏离物理上合理的运动轨迹——出现穿越、非物理拉伸、质量不守恒等伪影。

\textbf{核心思路。}PIDG 的关键洞察是：每个高斯基元天然可以被视为一个拉格朗日物质粒子——就像在流体力学中跟踪流体微团的运动轨迹一样。因此，PIDG 将每个高斯的状态向量扩展为 $G_i = (\mu_i, q_i, s_i, c_i, o_i, e_i, v_i, \sigma_i)$，其中前三项是常规的位置、旋转（四元数）、缩放，$c_i, o_i$ 是颜色和不透明度，而 $e_i \in \mathbb{R}^H$（嵌入向量）、$v_i \in \mathbb{R}^3$（速度向量）和 $\sigma_i \in \mathbb{R}^6$（应力张量的六个独立分量）是新引入的独立物理属性。这意味着每个粒子不仅知道自己"在哪里"，还知道自己"运动多快"和"承受多大的内部力"。

\textbf{4D 分解哈希编码。}在物理建模之前，PIDG 首先需要一个高效的时空表示来支撑后续的物理场预测。它将 4D 坐标 $(x,y,z,t)$ 映射到四个独立的 3D 哈希网格 $G_{xyz}, G_{xyt}, G_{yzt}, G_{xzt}$，将内存从 $O(n^4)$ 降至 $O(n^3)$。空间 MLP $f_s$ 从 $G_{xyz}$ 提取空间上下文特征并转换为方向性注意力权重 $\alpha = 2\sigma(f_s(G_{xyz})) - 1 \in (-1,1)$，用于调制时间 MLP $f_t$ 产生的原始时间编码：$h = \alpha \odot f_t(G_{xyt}, G_{yzt}, G_{xzt})$。直觉上，空间注意力权重决定了"这个位置的粒子应该对时间变化有多敏感"，使运动预测同时兼顾空间结构和时间变化。解码 MLP $D(h)$ 将特征 $h$ 映射为平移 $T_x$、旋转增量 $\Delta r$（四元数表示，转换为旋转矩阵以实现平滑插值和稳定优化）和缩放更新 $\Delta s$，更新规范空间的高斯参数：$\mu' = R_x\mu + T_x$，$S' = S + \Delta s$，$R' = R + \Delta r$。

\textbf{时间演化物质场与 Cauchy 动量约束。}在获得了高效的时空编码后，PIDG 构建时间演化物质场来预测每个粒子的物理状态。特征向量 $F$ 由三部分组成：六个可学习的轴对齐张量（空间平面 $XZ, XY, YZ$ 和时间平面 $XT, YT, ZT$，捕获轴对齐的空间和时间相关性）、傅里叶时间编码 $T(t)$（捕获高频时间变化）和可学习粒子索引嵌入 $e_i \in \mathbb{R}^H$（编码每个粒子的内在物理属性），总维度为 $|F| = 6 + 2n + H$。多头 MLP $f_\phi$ 从 $F$ 联合预测速度向量 $v(x,t) \in \mathbb{R}^3$ 和六个独立的应力张量分量 $\sigma(x,t) \in \mathbb{R}^6$，紧凑地写作 $(v, \sigma) = f_\phi(F)$。

物理约束的核心思想来自连续介质力学：在一个控制体积中，牛顿第二定律 $F = ma$ 推广到连续介质后，动量的物质导数（即随流体运动的观察者感受到的加速度）应等于体力与面力之和。当忽略体力时，这一平衡等价于 Cauchy 动量残差为零。PIDG 将此残差定义为：

\begin{equation}
  r(x,t) = \rho\!\left(\frac{\partial v}{\partial t} + (v \cdot \nabla)v\right) - \nabla \cdot \sigma,
\end{equation}

其中 $\rho(\partial v/\partial t + (v \cdot \nabla)v)$ 是惯性项（随物质运动的加速度乘以密度），$\nabla \cdot \sigma$ 是应力散度（内部力密度），残差 $r$ 衡量偏离动量平衡的程度。物理约束损失惩罚残差的 $L_2$ 范数：

\begin{equation}
  L_{\text{CMR}} = \frac{1}{M}\sum_{i=1}^{M}\|r(x_i, t_i)\|_2^2,
\end{equation}

通过选择不同的本构律（流体、弹性固体、刚体），同一框架可适应多种材料类型。这意味着 PIDG 将每个粒子的速度和应力视为独立的内在属性，使其不仅被编码在变形场中，还能随时间连续演化，以物理一致的方式更新运动和内部状态。

\textbf{拉格朗日粒子流匹配。}仅靠 Cauchy 动量残差约束往往无法让速度场和应力场收敛到物理上有意义的解——这是因为 PDE 约束本身是一个弱约束，缺少直接的观测锚定。为此，PIDG 引入拉格朗日粒子流匹配（Lagrangian Particle Flow Matching, LPFM）作为数据拟合项，直接用观测光流监督粒子运动。具体而言，PIDG 采用后向光流分解策略：从第 $t{+}1$ 帧的像素 $p_4$ 反向投影到三维再重投影到第 $t$ 帧，得到后向运动流 $flow_\text{gt} = p_2 - p_1$。与此前的前向策略相比，后向运动流已在参考帧 $I_t$ 中表达，因此可以直接在该帧上应用运动掩码、去除条带伪影，得到干净的光流监督信号。在此基础上，PIDG 计算两种流：高斯流（根据预测的位移传播高斯并加权求和得到像素位移）和速度流（根据预测的速度平流高斯并投影到图像平面），分别与观测运动流计算 $L_1$ 损失：

\begin{equation}
  L_{\text{LPFM}} = \lambda_g \|flow_g - flow_\text{gt}\|_1 + \lambda_v \|flow_v - flow_\text{gt}\|_1,
\end{equation}

其中 $\lambda_g = \lambda_v = 0.5$。这一数据拟合项将每个粒子的拉格朗日轨迹锚定到其物理状态，与基于边界条件的约束不同，它提供了更通用、更具物理基础的正则化，促使速度场和应力场快速收敛到物理一致的解。

\textbf{训练策略。}整体训练采用两阶段策略。第一阶段：渲染损失 $L_\text{render} = (1-\lambda_c)L_1 + \lambda_c L_\text{D-SSIM}$ 驱动所有高斯进行自适应致密化，联合优化静态几何和动态变形。第二阶段：暂停致密化，应用动态掩码隔离运动区域——被掩码覆盖的动态高斯继续更新变形场，未被掩码的静态高斯冻结变形、仅在规范哈希空间中精炼渲染参数。这一过程实现了静态和动态组件的高效分离，在最大程度保留静态表示的同时，为后续物理建模提供了干净的动态区域。完整的训练损失为 $L = L_\text{render} + \lambda_\text{CMR} L_\text{CMR} + \lambda_\text{LPFM} L_\text{LPFM}$，将视觉保真度和物理一致性统一在同一个可微管线中。

\paragraph{哈密顿力学保证能量守恒的变形场}

上一节 PIDG 用 Cauchy 动量残差显式约束速度场和应力场，实现了连续介质力学层面的物理一致性。但 PIDG 的物理约束本质上是"惩罚偏离"——它通过残差损失强制运动满足 PDE，而非让运动从守恒结构中自然涌现。NeHaD（2025.12）~\cite{2512_10424} 走了另一条路：不显式预测速度和应力，也不构造 PDE 残差，而是用一个学习守恒律的神经网络替代普通 MLP 变形解码器，让变形自动满足能量守恒和辛结构。两者共享同一基础设施——hex-plane 编码器提取时空特征——但变形解码器的设计哲学截然不同。

\textbf{从 MLP 到哈密顿神经网络。}PIDG 的变形解码器是标准 MLP，直接输出平移、旋转和缩放增量；NeHaD 将其替换为哈密顿神经网络（HNN）。直觉上，PIDG 的 MLP 问的是"这个高斯应该移动多少"，HNN 问的是"这个高斯所处的能量景观长什么样"——回答后者自然就能推导出移动方向。具体而言，hex-plane 编码器 $\mathcal{E}$ 提取时空特征 $\boldsymbol{f}_i$（与 4DGS（2024.04）~\cite{2404_03210} 相同），经可微 MLP $\mathcal{M}$（深度 $D$、宽度 $W$）映射为隐表示 $\boldsymbol{h}_i = \mathcal{M}(\boldsymbol{f}_i) \in \mathbb{R}^W$。关键区别在于，NeHaD 不直接从 $\boldsymbol{h}_i$ 预测变形向量，而是让 HNN 学习两个标量函数 $F_1(\boldsymbol{h}_i)$ 和 $F_2(\boldsymbol{h}_i)$，通过自动微分从标量势生成向量场：

\begin{equation}
  \boldsymbol{v}_c = \nabla_{\boldsymbol{h}_i} F_1(\boldsymbol{h}_i) \cdot \boldsymbol{I}, \quad \boldsymbol{v}_s = \nabla_{\boldsymbol{h}_i} F_2(\boldsymbol{h}_i) \cdot \boldsymbol{M}^\top,
\end{equation}

其中 $\boldsymbol{v}_c$ 是保守场（$\nabla \times \boldsymbol{v}_c = \boldsymbol{0}$，保能量，驱动平移运动），$\boldsymbol{v}_s$ 是旋度场（$\nabla \cdot \boldsymbol{v}_s = 0$，保体积，驱动旋转运动），$\boldsymbol{M}$ 是保持辛结构的排列张量。这一设计的深层动机是：直接学习高维向量场容易导致模式坍缩，且难以保证能量守恒——因为能量是标量，而向量场是高维的。通过先学习标量势函数再用自动微分生成向量场，物理性质（保守性、无散度性）由构造保证，而非靠损失函数惩罚。最终，轻量属性适配器 $\mathcal{A}_\mu, \mathcal{A}_s, \mathcal{A}_r$（线性层）将组合向量场 $\boldsymbol{v} = \boldsymbol{v}_c + \boldsymbol{v}_s$ 投射到位置、缩放和旋转各自的变形空间：$\Delta\boldsymbol{\mu}_i = \mathcal{A}_\mu(\boldsymbol{v})$，$\Delta\boldsymbol{s}_i = \mathcal{A}_s(\boldsymbol{v})$，$\Delta\boldsymbol{r}_i = \mathcal{A}_r(\boldsymbol{v})$。注意，这里不需要像 PIDG 那样为每个高斯显式存储速度 $v_i$ 和应力 $\sigma_i$ 作为独立属性——物理信息隐含在 HNN 的参数和前向传播中。

\textbf{Boltzmann 平衡分解：软掩码式静态/动态分离。}PIDG 用两阶段硬分割隔离运动区域（动态掩码覆盖的高斯继续更新，其余冻结），NeHaD 提出了一种更灵活的连续替代方案。其核心思想是：只有偏离平衡态的高斯基元才需要动态建模。对第 $i$ 个高斯，计算其空间偏离 $\Delta d_i = \|\boldsymbol{\mu}_i - \boldsymbol{\mu}_{eq}^{(i)}\|_2 / \sigma_s$ 和时间偏离 $\Delta\tau_i = (t - t_{eq}^{(i)}) / \sigma_t$（其中 $\boldsymbol{\mu}_{eq}^{(i)}$ 和 $t_{eq}^{(i)}$ 均初始化为高斯属性并随训练优化），组合为谐振子模型的能量：

\begin{equation}
  E_{st}^{(i)} = \frac{1}{2}(\Delta d_i^2 + \Delta\tau_i^2) + \lambda \, \Delta d_i \, \Delta\tau_i,
\end{equation}

其中耦合项 $\lambda \, \Delta d_i \, \Delta\tau_i$ 捕获空间运动与时间演化的相互依赖。位置平衡掩码遵循 Boltzmann 分布：

\begin{equation}
  M_{pos}^{(i)} = (1 - \gamma) \cdot \exp(-\beta \, E_{st}^{(i)}) + \gamma,
\end{equation}

其中 $\beta = 1/T$ 控制能量分布的锐度，$\gamma$ 确保最小动态响应以防止完全失活。距平衡态越近的基元掩码值越大（趋向静态），越远的掩码值越小（趋向动态）。与 PIDG 的硬二值掩码不同，这是一个连续软掩码，变形后的位置为 $\boldsymbol{\mu}_i' = \boldsymbol{\mu}_i + \Delta\boldsymbol{\mu}_i \odot (1 - M_{pos}^{(i)})$。对缩放属性，NeHaD 采用仅时间维度的分解（因为缩放变化影响的是全场景表面细节而非全局结构），掩码为 $M_{scale}^{(i)} = (1 - \gamma) \cdot \exp(-\beta \, E_t^{(i)}) + \gamma$。

\textbf{物理信息约束：辛积分与局部刚性。}HNN 保证了变形场的辛结构，但标准 HNN 在真实耗散场景（如摩擦）中仍有局限——它学到的守恒量接近但并非精确等于总能量。为此，NeHaD 引入两个物理信息约束来增强鲁棒性。第一，位置更新采用二阶辛积分（Position Verlet），而非简单的线性叠加。HNN 的变形输出 $\Delta\boldsymbol{\mu}_i$ 被解释为瞬时速度，力 $\boldsymbol{F}_i$ 从保守分量 $\boldsymbol{v}_c$ 通过自动微分获得（负势能梯度），积分公式为：

\begin{equation}
  \tilde{\boldsymbol{\mu}}_i = \boldsymbol{\mu}_i + \Delta t \cdot \Delta\boldsymbol{\mu}_i + \frac{(\Delta t)^2}{2} \boldsymbol{F}_i.
\end{equation}

相比 PIDG 直接用速度平流更新位置（一阶），二阶辛积分在长时间序列中保持更好的能量守恒和数值稳定性，且本身保持辛结构。辛积分在 BED 掩码之前应用：$\boldsymbol{\mu}_i' = \tilde{\boldsymbol{\mu}}_i \odot (1 - M_{pos}^{(i)}) + \boldsymbol{\mu}_i \odot M_{pos}^{(i)}$。第二，旋转更新受局部刚性约束（受 ARAP 启发）：将 HNN 预测的四元数增量 $\Delta\boldsymbol{r}_i = [\Delta w_i, \Delta\boldsymbol{g}_i]^\top$ 转换为轴角表示 $\phi_i = 2 \cdot \mathrm{atan2}(\|\Delta\boldsymbol{g}_i\|, \Delta w_i)$，经平滑限幅 $\phi_i^* = \phi_{\max} \cdot \tanh(\phi_i / \phi_{\max})$ 抑制大角度突变，再重构为约束后的四元数增量 $\Delta\boldsymbol{r}_i'$ 并通过四元数乘法更新朝向 $\boldsymbol{r}_i' = \mathcal{N}(\boldsymbol{r}_i \otimes \Delta\boldsymbol{r}_i')$。训练损失为标准渲染损失加全变分正则：$\mathcal{L}_{total} = (1-\lambda)\mathcal{L}_1 + \lambda\mathcal{L}_{DSSIM} + \mathcal{L}_{TV}$（$\lambda = 0.2$），无需像 PIDG 那样附加 PDE 残差损失和光流损失——物理一致性已内置于网络架构和积分方案中。

\paragraph{物理信息动态凸辐射场}

\textbf{Baseline 问题。}前述 PIDG 和 NeHaD 都在椭球高斯基元上叠加物理约束——PIDG 用 Cauchy 动量残差约束速度场和应力场，NeHaD 用哈密顿动力学保证能量守恒。然而，无论哪种方法，底层的椭球基元本身并未改变。事实上，已有的物理感知方法所使用的动态基元——体素（NeRF）或椭球核（3DGS）——本质上都是为渲染效率而设计的几何表示，在物理模拟中暴露出四个根本缺陷：(1) 在混合欧拉-拉格朗日框架中，中心驱动的动力学仅用中心粒子或池化邻域更新基元状态，缺乏对非均匀变形的空间敏感性——一个椭球内部的不同位置可能发生截然不同的变形，但椭球只能整体移动；(2) 预定义的粒子-基元绑定忽略材料和几何变异，导致物理耦合不准确——柔软材料和刚性材料的粒子不应被绑定到相同的变形模式；(3) 基于网格的无网格 MPM 方法难以表示尖锐或演化的边界，阻碍物理上有意义的曲面变形；(4) 固定的椭球核受限于 splatting 效率，无法建模各向异性或非线性材料，且在平面或棱角区域空间覆盖率差。本质上，渲染和物理模拟对动态基元提出了矛盾的要求：渲染需要紧凑的外观编码，物理模拟需要丰富的几何和物理表达能力——现有方法未能调和这一矛盾。

\textbf{核心思路。}PhysConvex（2026.02）~\cite{2602_18886} 的关键洞察是：凸多面体（convex polytope）天然比椭球更适合同时充当渲染基元和物理模拟单元。凸体通过顶点和半平面（half-space）定义，天然支持两种互补的变形表示——顶点动力学提供空间灵活性和非均匀变形的敏感性，曲面动力学捕获演化的边界和棱角——这正是椭球做不到的。PhysConvex 将动态场景表示为 $N$ 个时变凸体的集合：

\begin{equation}
C(t) = \{K_n(t),\, c_n,\, \alpha_n,\, \sigma_n,\, \gamma_n\}_{n=1}^N,
\end{equation}

其中 $K_n(t)$ 是第 $n$ 个动态凸体在时刻 $t$ 的凸包，$c_n, \alpha_n$ 是时间不变的颜色和不透明度，$\sigma_n, \gamma_n$ 分别控制凸体的平滑度和锐度。凸体之间以分片（piecewise）方式逼近复杂动态几何，用紧凑的基元集合实现密集的体积覆盖。

\textbf{边界驱动凸变形。}凸变形同时从两个层面驱动——这构成了 PhysConvex 的表示核心。显式层面，凸包的每个顶点 $x_k(t)$ 在物理动力学下独立演化，提供空间自适应的顶点变形：不同顶点可以以不同方式响应力，实现非均匀变形。隐式层面，凸体等价地由 $H$ 个时变半平面 $H(t) = \{(n_h(t), d_h(t))\}_{h=1}^H$ 表示，其中 $n_h(t)$ 和 $d_h(t)$ 分别是法向量和偏移量。点 $x$ 到第 $h$ 个时变平面的符号距离为 $f_h(x,t) = n_h(t) \cdot x + d_h(t)$，由此定义时变占据函数：

\begin{equation}
O(x,t) = \sigma\!\left(-\gamma(x,t)\right), \quad \gamma(x,t) = \sigma\!\left(\sum_{h=1}^{H} \sigma(f_h(x,t))\right),
\end{equation}

其中外层 $\sigma$ 控制凸体平滑度，内层 $\sigma$ 近似凸包的符号距离函数。直觉上，显式顶点动力学适合物理模拟（牛顿运动定律直接作用于顶点），隐式曲面动力学适合可微渲染（半平面表示支持解析梯度）——两种表示通过凸包的对偶性自然统一。

\textbf{降阶凸模拟。}在获得了凸体表示后，PhysConvex 开发了一种无网格（mesh-free）的降阶凸模拟方法，无需预定义绑定即可捕获复杂几何和异质材料的丰富变形行为。每个凸体的顶点位移通过空间变化的神经蒙皮特征模态（neural skinning eigenmodes）$W^\phi$ 参数化：

\begin{equation}
x_k(t) = X_k + \sum_{m=1}^{M} W_m^\phi(X_k) \, Z_m(t) \, [X_k, 1]^T,
\end{equation}

其中 $X_k$ 是参考构型中的顶点位置，$M$ 是降阶自由度数，$Z_m(t) \in \mathbb{R}^{3 \times 4}$ 是第 $m$ 个时变手柄变换矩阵。蒙皮特征模态 $W^\phi(X) \in \mathbb{R}^{3 \times M}$ 由小型 MLP 表示（参数 $\phi$），作为物理先验的变形基函数。与 PIDG 的连续体力学场不同，$W^\phi$ 不是直接预测速度或应力，而是学习编码形状和材料感知的低维物理变形模式——其训练目标是最小化弹性变形能量与正交约束的联合目标：

\begin{equation}
\phi^* = \arg\min_\phi \int_{\mathbb{R}^3} \rho(X)\, \Psi(\epsilon(X, z, W^\phi)) + \lambda_\text{ortho} \sum_{i=1}^{M}\sum_{j=1}^{M}\left(\int_{\mathbb{R}^3}\rho(X)(W_i^\phi W_j^\phi - \delta_{ij})^2 \,dX\right) \,dX,
\end{equation}

其中 $\rho(X)$ 是密度函数，$\Psi(\cdot)$ 是应变能密度函数，$\delta_{ij}$ 是 Kronecker delta。弹性能量项鼓励模态编码物理上显著的低维变形行为，正交项确保不同模态之间不冗余——这一设计使 $W^\phi$ 兼具材料感知和形状感知。直觉上，这些模态就像一组预定义的变形"按钮"，每个按钮控制一种变形模式（如弯曲、扭转、膨胀），优化过程只需学习每个时刻该按下哪些按钮以及按多深。时间演化的降阶自由度 $\mathbf{z}(t) = \text{flat}(Z(t)) \in \mathbb{R}^{12M}$（将变换矩阵展平）初始化为零（保证 $x_k(0) = X_k$），在每个时间步通过隐式时间积分的增量势公式（incremental potential formulation）更新，由标准的牛顿求解器求解。

\textbf{可微渲染与端到端物理优化。}渲染时，3D 时变凸体投影到 2D 并构建动态凸包，通过 $\alpha$-混合 光栅化：

\begin{equation}
C(q,t) = \sum_{i=1}^{N} c_i \, \alpha_i \, O_i(q,t) \prod_{j=1}^{i-1}(1 - \alpha_j \, O_j(q,t)),
\end{equation}

其中 $O_i(q,t)$ 是基于半平面符号距离的时变占据函数。整个流程分两阶段训练。第一阶段：从输入视频的首帧多视角出发，优化未变形的凸辐射场 $C(0)$，最小化 $L = (1-\lambda)L_1 + \lambda L_\text{D-SSIM} + L_m$（其中 $L_m$ 正则化凸体数量），得到材料空间中的初始几何和外观。第二阶段：通过联合可微模拟与渲染，仅用单视角前端视频监督，端到端优化弹性参数（杨氏模量 $E$、泊松比 $\nu$）和蒙皮场 $W^\phi$：

\begin{equation}
E^*, \nu^*, \phi^* = \arg\min_{E,\nu,\phi} L_\text{sim}, \quad L_\text{sim} = \frac{1}{SP}\sum_{p=1}^{P}\sum_{s=1}^{S} \|C(q_p, t_s) - C_\text{gt}(q_p, t_s)\|_2^2,
\end{equation}

其中 $S$ 是时间步数，$P$ 是像素数。为增强物理先验的泛化能力，PhysConvex 采用前馈物理系统辨识框架：预训练的视频 Vision Transformer 从大规模视频数据中提取运动线索和先验物理知识，小型 MLP 回归头估计 $E, \nu, \phi$ 的粗预测值，再通过可微模拟精修以更好地匹配参考视频。

\textbf{与 PIDG / NeHaD 的对比。}三种方法代表了物理感知 3DGS 的三条技术路线。PIDG 在椭球基元上叠加连续体力学约束（Cauchy 动量残差），通过显式的速度和应力属性让每个粒子"知道自己运动多快、承受多大的力"。NeHaD 用哈密顿神经网络替代普通 MLP，将物理守恒嵌入网络架构——不再需要附加 PDE 残差损失，物理一致性由辛结构保证。PhysConvex 则走了第三条路：从基元本身的几何表达能力出发，用凸多面体替代椭球——凸体的顶点和曲面天然具有空间自适应变形能力，无需像 PIDG 那样为每个粒子显式存储速度和应力，也无需像 NeHaD 那样在相空间中学习保守场。三种方法共同指向一个结论：将物理定律引入神经辐射场/高斯 splatting 的方式，正从"事后约束"（损失函数惩罚 PDE 残差）走向"内置设计"（基元选择、网络架构本身满足物理结构）。

\paragraph{多模态物理场景联合模拟}

前述 PIDG 用 Cauchy 动量方程约束\emph{已有的}多视角视频中的高斯变形场——它的目标是重建，即从观测中恢复物理一致的运动。然而在实际应用中，用户往往只有一张图片或一段短视频，却希望看到物体在施加外力后的真实物理响应（如杯子被撞倒、布料被拉扯）——这是一个正向模拟问题，而非重建问题。PhysFlow（2024.11）~\cite{2411_14423} 正是从这一目标出发：给定一张或多张输入图像，自动推断材料属性并驱动物理模拟，生成新视角下的 4D 动态场景。

\textbf{从 PDE 约束到粒子模拟。}PIDG 将物理定律编码为损失函数中的 PDE 残差项，通过反向传播间接引导变形场收敛——这是一种"软约束"策略，物理一致性依赖于网络的表达能力。PhysFlow 则走了一条更直接的路径：采用可微物质点法（MLS-MPM）作为物理引擎，将每个高斯基元视为拉格朗日粒子 $P = \{(x_p, v_p, F_p)\}$，在欧拉网格上求解动量守恒方程 $\rho \frac{Dv}{Dt} = \nabla \cdot \sigma + f^\text{ext}$，通过 B 样条核函数在粒子和网格之间传递物理量。变形梯度 $F$ 被分解为弹性部分 $F^E$ 和塑性部分 $F^P$，使框架能自然处理永久变形（如橡皮泥被压扁后不回弹）。与 PIDG 的 PDE 残差约束不同，MPM 是真正的粒子-网格耦合模拟器——物理定律在离散化层面被严格满足，而非作为损失项被"鼓励"。

\textbf{多模态材料推理与场景构建。}PIDG 假设材料类型已知或隐式学习，不显式优化材料参数。PhysFlow 面对的挑战更根本：从单张图片出发，既不知道"这是什么材料"，也不知道"它有多硬"。为此，PhysFlow 利用 GPT-4 进行视觉-语言联合推理——给定输入图像和文本提示，推断材料类型 $M$（弹性、橡皮泥、金属、泡沫、沙子、牛顿流体、非牛顿流体共七类）、密度 $\rho$ 和初始材料参数 $M_p^{(0)}$。不同材料类型对应不同的参数化：弹性体由杨氏模量 $E$ 和泊松比 $\nu$ 描述；橡皮泥和金属额外需要屈服应力 $\sigma_Y$（区分弹性变形和塑性流动）；泡沫引入塑性粘度 $\eta$；沙子由摩擦角 $\phi_\text{fric}$ 控制；牛顿流体由粘度 $\mu$ 和体积模量 $\kappa$ 刻画；非牛顿流体则组合剪切模量 $G$、体积模量 $\kappa$、屈服应力 $\sigma_Y$ 和塑性粘度 $\eta$。这种显式的材料分类确保每个物体使用正确的本构模型，而非用同一个通用方程拟合所有材料。

在获得材料属性后，PhysFlow 根据输入类型灵活选择 3DGS 重建方法构建初始高斯场景：多视角图像用 PGSR 构建静态表示，动态视频用 Grid4D 生成时变表示，单张图像用 Splatt3R 从学习先验推断场景结构。场景中的前景物体通过 SAM2Point 分割，作为外力施加区域。

\textbf{光流引导的材料参数优化。}有了初始场景和材料参数后，PhysFlow 对场景施加外力，通过可微 MPM 模拟物体运动，渲染模拟结果得到光流 $\hat{U}$。与此同时，利用视频扩散模型从输入图像生成参考视频并计算光流 $U$——扩散模型的选择同样适配输入条件：仅有图像时用 SVD，有图像和文本时用 CogVideoX，涉及轨迹或相机运动时用 MotionCtrl，更复杂的笔触运动则用 MOFA-Video 或 Kling 1.5。优化目标为光流损失：

\begin{equation}
\mathcal{L}_{flow} = \sum_t \|U(I_t, I_{t+1}) - \hat{U}(\hat{I}_t, \hat{I}_{t+1})\|^2,
\end{equation}

其中 $\hat{I}_t$ 为模拟渲染帧，$U$ 由 RAFT 光流网络从参考视频帧对中提取。这一设计的关键洞察在于损失函数的选择：PhysDreamer 使用的渲染损失 $L_\text{render} = \sum_t \lambda L_1(I_t, \hat{I}_t) + (1-\lambda)L_\text{D-SSIM}(I_t, \hat{I}_t)$ 直接比较像素值，但视频扩散模型生成的帧包含大量噪声，像素级匹配仅适用于微幅运动；SDS 损失 $L_\text{SDS}$ 虽然能绕过直接渲染，但对视频扩散模型需要巨大显存。光流则在运动层面比较——它对大幅复杂运动更鲁棒、对生成噪声更不敏感、且内存效率远高于 SDS。消融实验证实：在合成数据集上，光流引导在粘度 $\mu$、体积模量 $\kappa$ 和屈服应力 $\sigma_Y$ 等参数上的相对误差（RE）和端点误差（EPE）均低于渲染损失和 SDS 损失。

\textbf{优化后 4D 重建。}优化后的材料参数 $M_p$ 被整合回 MLS-MPM 框架，驱动 3D 高斯粒子在新视角 $v_1, v_2, \cdots, v_k$ 下渲染，实现完整的 4D 重建。32 人人类评估实验表明，PhysFlow 在物理真实感和照片真实感上的平均得分均超过 3 分（5 分制），显著高于 PhysGaussian（手动参数）、PhysDreamer（仅弹性）和 Physics3D（弹性+粘弹性）三个基线方法。

\paragraph{单目视频物理参数估计与物理驱动变形}

前述 PhysFlow 同样使用可微 MPM 和光流优化材料参数，但它假设几何已由独立重建方法提供——先重建再估计物理参数的两步策略在多视角输入下行得通，因为三角测量已经约束了 3D 结构和绝对尺度。MonoPhysics（2026.05）~\cite{2605_30320} 面对的是更根本的单目设定：几何必须在缺少多视角约束的条件下推断，且尺度天然歧义——而尺度是物理模拟的关键输入（同样外观的物体，真实尺寸差十倍意味着完全不同的物理行为）。因此，几何重建和物理参数估计无法像多视角方法那样解耦进行，必须联合优化。

\textbf{管线概述。}MonoPhysics 的整体管线也基于 3DGS + 可微 MPM，与 PhysFlow 共享"高斯既做渲染基元又做 MPM 粒子"的双重身份设计。但 PhysFlow 的光流损失仅用于优化材料参数（几何已经固定），而 MonoPhysics 需要同时优化几何、外观和物理参数，且初始化来自 3D 基础模型（如 SAM 3D）预测的单图几何——这类模型能恢复视觉上合理的形状，但输出缺乏绝对尺度且在遮挡区域不可靠。MonoPhysics 识别出单目设定下视觉侧（渲染、光流）与物理侧（MPM 模拟）之间存在三条"断裂"，并分别为每条断裂设计了一座桥梁。

\textbf{桥梁一：全局尺度校正。}单目重建只能恢复到未知尺度因子的几何——沿同一投影射线上的点对相机不可区分。物理模拟却要求绝对尺度：世界空间位置、速度、接触时间、体积和应力幅值全部取决于尺度。PhysFlow 等多视角方法将粒子直接存储在世界空间，通过多视角三角测量自然获得绝对尺度。MonoPhysics 则将粒子存储在相机空间，引入单个可学习标量 $s$ 控制全局缩放，通过已知相机外参转换到世界空间：

\begin{equation}
  \mathbf{x}_i^w = \mathbf{R}(s \cdot \mathbf{x}_i) + \mathbf{t},
\end{equation}

其中 $\mathbf{R}, \mathbf{t}$ 是相机到世界的旋转和平移，$\mathbf{x}_i$ 是相机空间位置。关键洞察在于：在相机空间中缩放会使粒子沿其投影射线移动，2D 投影保持不变但所有物理量（体积、速度、接触力）随之改变。因为每个粒子的物理行为都依赖 $s$，图像空间损失通过可微模拟反传时会从所有粒子和所有帧中聚合出一致的尺度信号——这是逐粒子位置更新做不到的（后者只能产生局部修正，无法产生物体整体远离/靠近相机的一致运动）。

\textbf{桥梁二：物理感知几何精化。}3D 基础模型提供的初始几何在遮挡区域不可靠，单视角下视觉观测无法修正这些区域。MonoPhysics 的关键洞察是：可微物理模拟提供了互补的信号——物理行为约束了视觉观测无法解析的区域。与 PhysFlow 在初始化后固定粒子体积不同，MonoPhysics 使粒子体积随优化自适应调整。具体而言，标准 MPM 初始化要求每个网格节点的粒子体积贡献之和等于网格单元体积 $\Delta x^3$：

\begin{equation}
  \sum_i V_i w_{i\ell} = \Delta x^3,
\end{equation}

其中 $\Delta x$ 是网格边长，$w_{i\ell}$ 是粒子 $i$ 对网格节点 $\ell$ 的贡献权重。已有方法（如 GIC、MASIV）在初始化时计算一次后固定体积，MonoPhysics 在每步优化中从当前粒子分布迭代求解体积，使几何表示能根据模拟反馈持续精化。在此基础上，MonoPhysics 提出物理感知粒子管理方案：定义视觉重要性 $p_i^\text{vis} \propto \det(\Sigma_i)^{1/2} \cdot o_i$（协方差行列式乘以不透明度）和物理重要性 $p_i^\text{phys} \propto V_i$（粒子体积），综合重要性为两者均值 $p_i = \frac{1}{2}p_i^\text{vis} + \frac{1}{2}p_i^\text{phys}$。在粒子重定位时，视觉重要性最低或物理重要性最低的粒子被移除，综合重要性最高的粒子被分裂——这使粒子分布同时跟踪视觉欠解析区域和物理欠解析区域，而非仅关注渲染质量。

\textbf{桥梁三：可微位置图。}标准 3DGS 的 $\alpha$-compositing 渲染仅通过颜色和不透明度传递梯度，PhysFlow 的光流损失也只在相邻帧的形状已大致对齐时才有效。当预测形状与目标形状尚未重叠时（例如初始几何偏差很大），没有任何梯度路径能将"你偏移了多远"的信号传递给高斯——这是标准渲染和光流损失的盲区。MonoPhysics 构建可微位置图来填补这一空白。对于被高斯 $i$ 覆盖的像素 $\mathbf{p}$，首先定义重新参数化的位置：

\begin{equation}
  \tilde{\mathbf{p}}_i = \mathbf{x}_i^{2D} + (\boldsymbol{\Sigma}_i^{2D})^{1/2} z, \quad z = \mathrm{sg}\!\left((\boldsymbol{\Sigma}_i^{2D})^{-1/2}(\mathbf{p} - \mathbf{x}_i^{2D})\right),
\end{equation}

其中 $\mathbf{x}_i^{2D}, \boldsymbol{\Sigma}_i^{2D}$ 是高斯的 2D 图像空间均值和协方差，$\mathrm{sg}(\cdot)$ 是停止梯度操作（将参数视为常数）。然后用与标准渲染相同的 $\alpha$-compositing 公式将 $\tilde{\mathbf{p}}_i$ 混合为位置图：

\begin{equation}
  \tilde{\mathbf{p}} = \mathbf{p}\prod_i(1-\alpha_i) + \sum_i \alpha_i \, \tilde{\mathbf{p}}_i \prod_{j<i}(1-\alpha_j).
\end{equation}

数值上 $\tilde{\mathbf{p}}$ 始终等于实际像素坐标 $\mathbf{p}$，但梯度可以沿 $\mathbf{x}_i^{2D}$ 和 $\boldsymbol{\Sigma}_i^{2D}$ 反传到 3D 高斯参数。基于此，MonoPhysics 使用非平衡去偏 Sinkhorn 散度（unbalanced debiased Sinkhorn divergence）对齐渲染位置图和目标轮廓的前景像素点云：

\begin{equation}
  \mathcal{L}_\text{sil}^t = \mathrm{SD}\!\left(\{\tilde{\mathbf{p}}_k, \tilde{a}_k\}_{k \in \hat{\Omega}_t}, \{\mathbf{p}_j, a_j\}_{j \in \Omega_t}\right),
\end{equation}

其中 $\hat{\Omega}_t$ 和 $\Omega_t$ 分别是渲染和目标的前景像素集合。最优传输提供的梯度即使在两个轮廓完全不重叠时也能指示每个渲染像素应该移动的方向——这正是标准渲染损失的盲区。

\textbf{光流损失与粒子分布正则化。}MonoPhysics 同样使用光流损失提供运动层面的监督（与 PhysFlow 的设计理念一致），但采用基于概率的光流损失，同时监督两种互补信号：从首帧到当前帧的全局位移和从上一帧到当前帧的瞬时速度：

\begin{equation}
  \mathcal{L}_\text{flow} = -\frac{1}{|\Omega_t|}\sum_{j \in \Omega_t}\left[\log p\!\left(\mathbf{f}_{0 \to t}(\mathbf{p}_j) \mid \hat{\mathbf{f}}_{0 \to t}(\mathbf{p}_j)\right) + \log p\!\left(\mathbf{f}_{t\text{-}1 \to t}(\mathbf{p}_j) \mid \hat{\mathbf{f}}_{t\text{-}1 \to t}(\mathbf{p}_j)\right)\right],
\end{equation}

其中 $\hat{\mathbf{f}}$ 是通过高斯光栅化将粒子世界空间速度投影到图像平面得到的渲染光流。概率形式（来自 SEA-RAFT）能自适应地降低噪声和不确定区域的光流权重。此外，为防止粒子在优化中脱离物体漂移或过度聚集，MonoPhysics 引入 $K$-近邻粒子分布正则化：

\begin{equation}
  \mathcal{L}_\text{distr} = \sum_{i=1}^{N}\sum_{j \in \mathcal{N}_K(i)}\max\!\left(0, \|\mathbf{x}_i - \mathbf{x}_j\| - \tau_\text{max}\Delta x\right) + \max\!\left(0, \tau_\text{min}\Delta x - \|\mathbf{x}_i - \mathbf{x}_j\|\right),
\end{equation}

其中 $\tau_\text{min}, \tau_\text{max}$ 是以模拟网格尺寸 $\Delta x$ 为单位的距离下界和上界。第一项在粒子漂移过远时将其拉回，第二项在粒子过度聚集时将其推开，两项在分布良好的区域自动消失。

\textbf{训练策略。}整体优化分两阶段。第一阶段联合优化物理参数（材料参数和初始速度）、高斯参数（位置、不透明度、协方差）和全局尺度 $s$，通过可微 MPM 模拟，使用全部四种损失：

\begin{equation}
  \mathcal{L} = \sum_{t=0}^{T}\left[\mathcal{L}_\alpha^t + \mathcal{L}_\text{sil}^t\right] + \sum_{t=1}^{T}\mathcal{L}_\text{flow}^t + \mathcal{L}_\text{distr}.
\end{equation}

注意此阶段不优化颜色——遵循 GIC 的实践，聚焦物理动态并在后续精炼外观。每次模拟反传后额外执行若干步 alpha-map 损失（无需重新运行模拟），保持视觉表示与演化几何的一致性。第二阶段冻结所有物理参数，仅用图像损失 $\mathcal{L}_\text{color}$（$L_1$ + SSIM）和 $\mathcal{L}_\alpha$ 精炼外观。在 Vid2Sim 基准上，MonoPhysics 仅用单目相机即达到杨氏模量 MAE $\log E = 0.52$、泊松比 MAE $\nu = 0.06$——与 Vid2Sim 使用 12 个相机的多视角结果（$\log E = 0.51, \nu = 0.06$）持平，远超 PAC-NeRF、SpringGaus 和 GIC 的单视角表现。

\paragraph{网格吸附高斯与骨骼驱动变形}

前面讨论的 PIDG、NeHaD、PhysConvex、PhysFlow 等方法都在点云层面上引入物理先验——高斯基元在欧氏空间中自由运动，缺乏显式的曲面结构。另一条路线则尝试引入三角网格作为几何先验：DG-Mesh 和 SplattingAvatar 将高斯固定锚定在网格面片上，虽然获得了曲面信息，却牺牲了高斯本身的灵活性——高斯无法离开所在面片，也无法在面片间滑动；而 SC-GS 等纯点云方法在大幅编辑时则缺乏曲面连续性信息。MaGS（2024.06）~\cite{2406_01593}（Mesh-adsorbed Gaussian Splatting）的核心设计思想介于两者之间：让高斯在网格面上”漫游”（roam）而非”锚定”（anchor），既保留网格的曲面连续性和模拟通用性，又不丧失高斯的灵活表达能力。

\textbf{Mesh-adsorbed Gaussian 表示。}MaGS 的混合表示由三角网格和 3D 高斯两部分组成。三角网格具有 $N$ 个顶点和 $M$ 个面片，定义了顶点坐标 $v_i \in \mathbb{R}^3$、面片拓扑（每面片三个顶点索引）和法向量 $n_j \in \mathbb{R}^3$（可从顶点和面片计算得出，面片内任意位置的法向量通过三个顶点法向的线性插值获得（2024.06）~\cite{2406_01593}）。每个高斯在标准 3DGS 属性之外，额外存储两个关键属性：逻辑坐标 $w \in [0,1]^3$ 和所在面片编号 $\phi$。初始化时，$w$ 随机生成，$\phi$ 均匀分配到各面片以确保每个面片上高斯数量相等。高斯的空间位置通过面片顶点插值加法向偏移计算：

\begin{equation}
\mu_i = f(w_i, v_{j_1}, v_{j_2}, v_{j_3}) + w_{i,3} \, n_j \, l_j,
\end{equation}

其中 $f$ 是重心插值函数 $f(w, x, y, z) = w_1 x + w_2 y + (1 - w_1 - w_2)z$，$l_j = |v_{j_1}v_{j_2} \times v_{j_1}v_{j_3}| \, / \, |v_{j_1}v_{j_2} \cdot v_{j_1}v_{j_3}|$ 是面片 $j$ 的缩放因子（反映面片面积的度量）。当 $w_{i,3} = 0$ 时高斯精确落在网格面上；当 $w_{i,3} \neq 0$ 时高斯悬浮于网格面上方或下方，增加了高斯沿法向的自由度。这一设计称为 Gaussian Hover，使得高斯不必严格贴合网格表面就能捕捉毛发、绒毛等超出网格表面的细节。

\textbf{三个协同变形网络。}与 PIDG 用 MLP 变形场从时空特征直接预测高斯增量不同，MaGS 的变形驱动完全绕开时间戳，改为从粗网格（guide mesh）的几何状态提取信息——这是因为模拟阶段没有视频时间信息可用，依赖时间戳的方法无法泛化到训练外的新变形。具体而言，三个网络各司其职：

\textit{MPE-Net（Mesh Pose Embedding）。}从粗网格的 handle 顶点提取姿态信息。Handle 顶点通过 Poisson 盘采样（2024.06）~\cite{2406_01593}在网格上均匀采样后用 KNN 映射到最近网格顶点，确保空间覆盖均匀。MPE-Net 以 handle 顶点集 $H$ 及其法向量集 $N$ 为输入，输出姿态嵌入 $E_p$（编码整体姿态）和顶点变形嵌入 $E_d$（编码每个顶点的局部变形特征）。这两个嵌入替代时间戳，成为后续两个网络的驱动信号。

\textit{RMD-Net（Relative Mesh Deformation）。}以 $E_p$ 和顶点 $i$ 对应的 $E_{d,i}$ 为输入，预测每个顶点的相对位移 $\Delta v_i$ 和旋转 $\Delta q_i$，精化粗网格的变形。此外还预测辅助量缩放 $\Delta s_i$、不透明度 $\Delta\alpha_i$ 和颜色 $\Delta c_i$——这些量虽不直接对应网格变形，但用于后续高斯属性的渲染插值。RMD-Net 的本质是从视频中学习运动先验（motion prior），使精化后的网格变形比粗网格更合理。

\textit{RGD-Net（Relative Gaussian Deformation）。}这是 MaGS 区别于所有锚定方法的关键网络。它预测每个高斯逻辑坐标的变化 $\Delta w_i$，输入为 $E_p$、该高斯所在面片三个顶点的变形嵌入 $E_{d,j_1}, E_{d,j_2}, E_{d,j_3}$ 以及当前逻辑坐标 $w_i$。$\Delta w_i$ 使高斯能随网格变形在面片间滑动——这正是"漫游"的含义。

\textbf{变形高斯的属性更新。}获得 $\Delta v$（顶点位移）、$\Delta w_i$（逻辑坐标增量）后，更新过程如下。首先计算变形后面片的缩放因子变化 $\Delta l_j = l_j' / l_j$ 和法向量变化 $\Delta n_j$。逻辑坐标和顶点位置同步更新：$w_i' = w_i + \Delta w_i$，$v_j' = v_j + \Delta v_j$。变形高斯的属性全部通过面片插值获得——位置、不透明度、颜色、缩放和旋转分别更新为：

\begin{align}
\mu_i' &= f(w_i', v_{j_1}', v_{j_2}', v_{j_3}') + w_{i,3}' \, n_j' \, l_j', \\
\alpha_i' &= f(w_i', \alpha_{j_1}, \alpha_{j_2}, \alpha_{j_3}) \cdot \alpha_i, \\
c_i' &= f(w_i', c_{j_1}, c_{j_2}, c_{j_3}) + \Delta c_i, \\
s_i' &= \Delta l_j \cdot f(w_i', s_{j_1}, s_{j_2}, s_{j_3}) \cdot s_i, \\
q_i' &= \Delta n_j \cdot f(w_i', q_{j_1}, q_{j_2}, q_{j_3}) \cdot q_i,
\end{align}

其中 $s_i, q_i, \alpha_i, c_i$ 是变形前的属性，$\Delta l_j$ 和 $\Delta n_j$ 由 RMD-Net 的顶点变形间接决定。直觉上，位置通过重心插值加 Gaussian Hover 偏移获得，不透明度和颜色通过面片顶点值的加权混合并叠加局部残差，缩放和旋转则通过面片顶点值加权后乘以面片变形引起的变化量。整个过程可形式化为 $(w + \Delta w) \otimes (v + \Delta v) \to \text{Final Deformed Gaussians}$，其中 $\otimes$ 表示 MaGS 的插值操作。最终变形高斯具有标准 3DGS 的全部参数，可直接通过可微渲染器进行 splatting 渲染。

\textbf{优化与模拟。}整个流程端到端可微，通过 $L = (1 - \lambda_\text{ssim})L_1 + \lambda_\text{ssim}L_\text{SSIM}$ 渲染损失联合优化 MPE-Net、RMD-Net、RGD-Net 和所有 Mesh-adsorbed Gaussian 参数。模拟时，MaGS 的关键优势在于网格的跨帧连续性——顶点-面片关系在所有帧中保持一致，因此可直接利用网格的通用性接入 ARAP（As-Rigid-As-Possible）编辑、SMPL 参数模型或软体物理引擎（Blender/Taichi）。RMD-Net 已在训练视频中学习到的运动先验可进一步精化模拟中的网格变形，重建结果无需重新训练即可用于布料模拟、碰撞响应等下游任务。

\subsubsection{事件相机辅助动态重建}

本路线介绍利用事件相机微秒级时间分辨率和高动态范围特性来辅助动态场景重建的方法。

以上所有方法都依赖标准 RGB 相机的输入。然而，标准相机的帧率有限（通常 30--60 fps），难以捕获快速运动。事件相机（event camera，又称神经形态相机）以微秒级时间分辨率异步捕获亮度变化，与标准相机以固定帧率捕获完整图像形成鲜明对比。在动态场景重建中，事件相机可以提供两个关键能力：(1) 高时间分辨率的运动信息——填补 RGB 帧之间的运动空白；(2) 高动态范围（HDR，High Dynamic Range）感知——不受标准相机的过曝/欠曝限制。

\paragraph{事件相机增强高时间分辨率变形}

标准变形场方法仅使用 RGB 帧作为监督信号，帧率通常为 30--60 fps，帧间快速运动的细节完全丢失。引入事件相机可以弥补这一信息缺口——事件相机在像素亮度对数变化超过触发阈值 $C$ 时异步输出事件 $e_{x,y}(t) = p \cdot C \cdot \delta(t)$，时间分辨率达微秒级。然而，先前的事件-3DGS 方法（如 EventNeRF、Deformable Event-3DGS）均假设阈值为全局常数，而真实事件相机的阈值在极性（正/负）、空间位置和时间上均存在复杂变化。常数阈值假设导致事件积分与实际亮度变化严重不一致，直接削弱了事件监督的有效性。

Event-boosted Deformable 3D Gaussians（2024.11）~\cite{2411_16180} 的核心改进是将事件阈值从"固定常数"变为"可学习参数"，并与 3DGS 重建形成互相增强的循环：更准确的阈值提升事件监督质量，更好的 3DGS 反过来又帮助精化阈值估计。方法提出 GS-阈值联合建模策略（GS-Threshold Joint Modeling, GTJM），包含三个阶段。

第一阶段是 RGB 辅助阈值估计。方法首先将事件按时间箱和极性累积为事件计数图（Event Count Map）$ECM \in \mathbb{R}^{B \times P \times H \times W}$，其中 $B$ 个时间箱均匀划分帧间区间，$P{=}2$ 对应正/负极性。与此同时，引入可学习的阈值参数 $C \in \mathbb{R}^{B \times P \times H \times W}$，与计数图同维度，使阈值能逐像素、逐极性、逐时间箱地建模变化。直觉上，事件计数图记录了"每个像素在每个时间窗口内触发了多少次事件"，而阈值参数将计数值转换为实际的亮度变化量。阈值建模损失为
\begin{equation}
\mathcal{L}_{thres} = \|E_{thres}(t,f) - \hat{E}_{thres}(t,f)\|_2^2,
\end{equation}

其中 $E_{thres}(t,f) = \log I(f) - \log I(t)$ 是两帧 RGB 图像之间的对数亮度差（作为真值），$\hat{E}_{thres}(t,f) = \sum_{b=1}^{B} \sum_{p=1}^{P} (ECM_{t,f} \odot C_{t,f})_{b,p,:,:}$ 是事件按阈值加权后的累积亮度差（$\odot$ 表示逐元素乘积）。该损失要求阈值加权后的事件累积与 RGB 帧间实际亮度变化保持一致。

第二阶段是 3DGS 辅助阈值精化。仅靠稀疏 RGB 帧监督阈值存在精度瓶颈——帧率越低，积分区间越长，事件累积误差越大。方法利用一个关键观察：一旦 3DGS 初步训练完成，它可以渲染任意时刻的中间帧，为阈值估计提供密集的伪监督。具体做法是冻结已训练的 3DGS，利用渲染的中间帧缩短事件积分区间，以 $\mathcal{L}_{thres} + \mathcal{L}_{event}$ 优化阈值 $C$。由于积分区间缩短，事件累积更可靠，阈值估计 MSE 降低约 50\%。

第三阶段是阈值与 3DGS 联合优化。前两阶段验证了阈值建模与 3DGS 重建之间的互利关系，因此最终将两者联合优化：

\begin{equation}
\{C^*, GS^*\} = \arg\min_{C, GS}(\mathcal{L}_{thres} + \mathcal{L}_{event} + \mathcal{L}_{rgb}),
\end{equation}

其中事件渲染损失
\begin{equation}
\mathcal{L}_{event} = \left\|\int_t^{t+\Delta t} C \, e(\tau) \, d\tau - [\log\hat{I}(t{+}\Delta t) - \log\hat{I}(t)]\right\|_2^2
\end{equation}

约束事件加权积分与 3DGS 渲染帧间的对数亮度变化一致，$\mathcal{L}_{rgb}$ 为标准的 $L_1$ + D-SSIM 组合损失。联合优化形成正向循环：准确的阈值提升事件监督质量，改善 3DGS 重建；更精确的 3DGS 又通过几何一致性纠正阈值误差。RGB 辅助阈值估计为 3DGS 重建带来 2.17 dB PSNR 提升。

此外，方法提出动态-静态分解策略（Dynamic-Static Decomposition, DSD）以进一步提升效率和质量。核心观察是动态场景中大量区域（如桌面、墙壁）完全静止，无需变形处理。具体流程为：前 3000 次迭代仅用静态高斯重建整个场景，动态区域因无法被静态模型拟合而自然呈现低重建质量。随后利用预训练 VGG19 提取渲染图 $\hat{I}(t)$ 与真值 $I(t)$ 的多尺度特征，计算余弦相似度图并上采样至统一分辨率取平均，该相似度图呈现明显的双峰分布——静态区域高相似度，动态区域低相似度。通过 Otsu 阈值法自动分割生成 2D 动态区域掩码。为将 2D 掩码扩展到 3D，方法利用 3DGS 可微渲染生成的深度图，将掩码像素反投影为 3D 点，再基于空间邻近性将 3D 点映射到高斯。为缓解分割边界的不确定性，采用双半径缓冲区策略：半径 $r_1$ 以内的高斯标记为动态，$r_2$ 以外的标记为静态，$r_1$ 与 $r_2$ 之间的高斯被剪枝形成缓冲带，让 3DGS 通过自适应密度控制（ADC）自动优化分解边界。最终，仅动态高斯经过变形场处理（变形场以高斯中心位置和时间戳为输入，输出位置、旋转、缩放的增量），静态高斯直接光栅化，同时加速渲染并提升重建质量。整个 DSD 流程在训练中仅执行一次，耗时约一分钟。

\paragraph{事件驱动变形与 RGB-事件双分支}

前述 Event-boosted Deformable 3DGS 关注的核心问题是事件阈值建模——将固定阈值变为可学习参数以提升事件监督精度。DEGS（2025.10）~\cite{2510_07752}（Deformable Event-based 3D Gaussian Splatting）则从另一个角度切入同一问题：即使事件阈值准确，低帧率 RGB 视频重建动态 3DGS 仍然面临一个根本困难——大帧间运动增加了解空间的歧义。一帧中的一个像素在下一帧中有更多候选匹配位置，变形场 $F_\theta$ 难以在相邻帧之间优化出唯一正确的变形路径。Event-boosted 通过更精确的事件监督来约束亮度一致性，而 DEGS 则直接利用事件流提供帧间运动轨迹先验来引导变形场优化。

DEGS 的解决思路是利用事件相机提供帧间运动先验。事件相机在像素亮度对数变化超过阈值时异步触发事件 $e_k=(x_k, t_k, p_k)$，其时间分辨率远高于 RGB 帧率。与 Event-boosted 将事件积分用于亮度监督不同，DEGS 将事件视为运动轨迹的载体——同一运动边缘在不同时刻触发的事件构成一条轨迹，这条轨迹就编码了帧间的运动信息。围绕这一思路，DEGS 面临两个关键技术问题：(1) 如何从事件流中提取可靠的运动轨迹；(2) 如何建立 2D 事件与 3D 高斯之间的运动对应关系。

\textbf{运动轨迹提取：LoCM 框架。}
对于运动轨迹提取，DEGS 提出 LoCM（LoRA + Contrast Maximization）框架。直觉上，同一运动边缘触发的事件如果沿正确的轨迹对齐，会在参考时刻形成锐利的图像；反之则模糊——这一观察是对比度最大化（Contrast Maximization, CM）的基础。具体地，给定事件集 $E=\{e_k\}_{k=1}^N$，通过事件光流将事件扭曲到参考时刻得到扭曲事件图像（Image of Warped Events, IWE）：

\begin{equation}
x_k' = W(x_k, t_k;\theta) = x_k - (t_k - t_\text{ref})\,v_\theta(x_k, t_k),
\end{equation}

其中 $W$ 是基于事件光流场 $v_\theta$ 的扭曲函数。若光流正确，扭曲后的 IWE 应具有最大的对比度（梯度方差）。然而，DEGS 没有直接用参数化的运动场作为 $v_\theta$，而是使用预训练的事件光流估计器 $f_\phi$——后者已在大量数据集上学习了多种运动模式的先验知识，包括相机自运动和多物体运动，比简单参数化运动场具有更强的建模能力。

直接在目标场景上全参数微调估计器会因两个原因失效：(1) 缺乏事件光流真值，无法进行有监督微调；(2) 在小数据集上全参数微调会过拟合，破坏原有的运动先验。受 NeuMA 启发，LoCM 将估计器参数分解为冻结的初始权重 $W_0$ 与低秩适配增量 $\Delta W = BA$（$B$ 零初始化、$A$ 随机高斯初始化），更新后的参数为 $W_\phi := W_0 + \Delta W$。零初始化保证训练开始时 $\Delta W = 0$，不偏离原始预测；低秩结构限制可学习自由度，防止微调颠覆已有先验。微调以无监督方式进行，目标函数为：

\begin{equation}
\mathcal{L}_\text{flow} = \frac{1}{\text{Var}(f_\phi(E(t_1, t_\text{ref})))} + \lambda \cdot TV(\phi),
\end{equation}

其中第一项最大化 IWE 的梯度方差（等价于最小化其倒数），$TV(\phi)$ 为全变分正则化以缓解过拟合。此外，为增强鲁棒性，DEGS 将图像平面划分为多个不重叠的 patch，并在多尺度分支上分别优化。通过 LoCM 框架，预训练估计器被高效地适配到目标场景，微调后的事件光流在后续用于引导变形场优化。

\textbf{事件-高斯运动对应。}
对于事件-高斯运动对应，DEGS 利用 3DGS 可微渲染器生成深度图，将每个 2D 事件反投影到 3D 空间，找到最近的 3D 高斯作为关联。为实现鲁棒绑定，每个高斯对应其最近的至多 $k{=}3$ 个反投影事件，并通过逆距离加权函数生成对应权重。一个关键实践细节是：初始深度图的质量直接影响绑定准确性，因此 DEGS 先以稀疏 RGB 用 $\mathcal{L}_\text{rgb}$ 预热 3500 次迭代获得粗糙但可用的深度图来完成初始绑定；在后续联合训练中，深度图不断优化，绑定关系也周期性更新以保持事件运动与高斯运动的对齐。

由于事件同时包含高斯变形和相机自身运动两部分贡献，DEGS 将运动信号分解为两个分量。相机自运动由 PoseNet 建模——PoseNet 接收插值的相机位姿作为输入，输出修正项生成连续的位姿函数 $(R, t)$，由此可导出瞬时平移/旋转速度 $(v_c, \omega_c)$。像素级自运动速度可由刚体运动和对应 3D 位置推导：

\begin{equation}
F_\text{ego} = \frac{1}{Z}\begin{pmatrix} -1 & 0 & x \\ 0 & -1 & y \end{pmatrix}v_c + \begin{pmatrix} xy & -(1{+}x^2) & y \\ (1{+}y^2) & -xy & -x \end{pmatrix}\omega_c,
\end{equation}

其中 $1/Z$ 为 3DGS 渲染的逆深度。高斯场景流由变形场预测，通过将变形前后高斯中心投影到图像平面得到：

\begin{equation}
F_\text{gs}^i = \text{Proj}(x_i + F_\theta(\gamma(x_i), \gamma(t_1))[0]) - \text{Proj}(x_i + F_\theta(\gamma(x_i), \gamma(t_0))[0]),
\end{equation}

其中 $[0]$ 表示取变形场的第一个输出（位置增量），$\text{Proj}$ 为 3D 到 2D 的相机投影，$\gamma$ 为位置编码。两部分的像素级速度之和应与 LoCM 微调后的事件光流一致，运动对齐损失为：

\begin{equation}
\mathcal{L}_\text{motion} = \sum_{i=1}^{N} \sum_{k \in B(i)} \| f_\phi(E(t_0, t_2))_k - (F_\text{ego} + F_\text{gs}^i) \|,
\end{equation}

其中 $B(i)$ 为第 $i$ 个高斯的事件绑定集合（仅包含时间戳在 $t_0 \pm \Delta t$ 邻域内的事件，以排除不同物体在不同时刻触发的无关事件）。

\textbf{两阶段训练流程。}
训练分两阶段。第一阶段为预热：仅用稀疏 RGB 以 $\mathcal{L}_\text{rgb} = \|I_c - \hat{R}(c)\|_1$ 训练 3500 次迭代，建立粗糙但可用的规范 3DGS 和深度图。第二阶段为联合训练：在 RGB 重建损失之外，引入事件伪图像损失和运动对齐损失。事件伪图像损失利用事件在帧间生成伪标签：在每个优化步中，随机选择帧间时间戳 $t_e$，用事件累积生成 $t_e$ 时刻的伪图像，作为 3DGS 渲染的监督信号。为缓解遮挡影响，使用最近的 RGB 图像估计 $t_e$ 处的颜色。总损失为：

\begin{equation}
\mathcal{L}_\text{total} = \mathcal{L}_\text{rgb} + \lambda_1 \mathcal{L}_\text{event} + \lambda_2(\text{iter})\, \mathcal{L}_\text{motion},
\end{equation}

其中运动损失权重采用退火策略 $\lambda_2 = 1 - e^{-\text{iter}/14000}$——训练初期变形场尚未收敛，其预测的场景流不准确，过强的运动约束会误导优化；随着训练推进，变形场逐步精确，运动约束逐步增强。这种 RGB 重建 + 事件伪图像 + 运动对齐的三重监督使变形场在低帧率条件下也能准确预测中间帧的变形。

\paragraph{单目事件流非刚性场景渲染}

前面讨论的 Event-boosted Deformable 3DGS 和 DEGS 都依赖 RGB 图像作为核心监督信号——前者用 RGB 帧校准可学习事件阈值，后者用 RGB 损失预热变形场。然而在许多实际场景中（无人机、AR 眼镜、高速运动），持续获取高质量 RGB 帧并不可行。Ev4DGS（2025.10）~\cite{2510_11717}（Event-based 4D Gaussian Splats）向前迈出更激进的一步：仅从单目事件流（无任何 RGB 输入）对非刚性物体进行新视角渲染。此前，基于事件的刚性场景新视角合成已有若干成果（EventNeRF 等），但非刚性场景的最新方法 DE-NeRF 仍需要额外的 RGB 图像输入，这在功耗受限和高速运动场景中是一个根本性障碍。Ev4DGS 是首个证明纯事件流足以驱动非刚性 4D 重建的工作。

核心困难在于问题的严重病态性。单目设置下，每个时刻只能从单一视角观测非刚性物体，且事件数据本质上是空间稀疏的——只有亮度变化超过阈值的像素才会触发事件，大量像素区域没有信号。直接端到端学习高维的非刚性 4D 高斯表示在这种弱监督下极易陷入局部最优。Ev4DGS 的关键洞察是将运动（motion）和外观（appearance）解耦为粗、精两阶段独立训练：先学习低秩变形模型捕获大尺度运动，再用事件流监督精化外观表示。

\textbf{掩码生成。}由于完全没有 RGB 输入，Ev4DGS 需要从事件流中提取最基本的物体轮廓信息。方法将事件按位置映射为事件位置图（event location map），然后用经典的 Snakes 活动轮廓算法自动提取二值分割掩码 $\{\hat{\mathbf{M}}_i\}_{i=1}^{N_v}$。这些掩码是整个框架唯一的几何监督来源，替代了 RGB 方法中的图像分割步骤。

\textbf{粗阶段：低秩变形模型。}受 NPGs 的参数化点模型启发，Ev4DGS 用一组时间依赖的稀疏粗点 $\mathbf{P}(t) \in \mathbb{R}^{N_c \times 3}$（$N_c \approx 1000$）表示物体的整体形状演化。为控制变形自由度、避免过拟合稀疏监督信号，模型采用低秩分解表示：

\begin{equation}
\mathbf{P}(t) = \sum_{k=1}^{K} \alpha_k(t) \mathbf{B}_k,
\end{equation}

其中 $K$ 个基形状 $\{\mathbf{B}_k \in \mathbb{R}^{N_c \times 3}\}_{k=1}^{K}$ 为可学习的低秩变形基向量，基权重 $\alpha_k(t)$ 由 MLP $f_\theta$ 从正弦位置编码的时间输入 $\phi(t)$ 预测。$K$ 是控制刚性的关键超参数：$K$ 越小，模型自由度越低，越趋向刚性变形；$K$ 越大则能表达更复杂的非刚性运动。直觉上，基形状 $\mathbf{B}_k$ 编码了物体变形的主要模式（如人体的摆臂、屈膝），而时间权重 $\alpha_k(t)$ 控制各模式在不同时刻的激活强度。粗阶段仅用事件流生成的二值分割掩码监督，损失为 Chamfer 距离：

\begin{equation}
\mathcal{L}_\text{mask} = \sum_{i=1}^{N_v} CD(\mathbf{P}_{\hat{\mathbf{M}}_i}, \pi(\mathbf{P}(t_i))),
\end{equation}

其中 $\mathbf{P}_{\hat{\mathbf{M}}_i}$ 是从掩码 $\hat{\mathbf{M}}_i$ 均匀采样的 2D 点集，$\pi(\cdot)$ 为投影函数。Chamfer 距离约束粗点模型的投影轮廓与掩码对齐，使模型从极稀疏的监督中学习物体的瞬时形状。粗阶段仅优化 MLP 权重 $\theta$ 和基形状 $\{\mathbf{B}_k\}$，参数量远小于完整的 4D 高斯表示。

\textbf{精阶段：4D 高斯外观重建。}在粗阶段捕获了大尺度运动后，精阶段构建完整的 4D 高斯表示。与 Event-boosted D3DGS 为每个高斯学习独立变形场不同，Ev4DGS 利用粗阶段的变形模型驱动所有高斯的运动——每个高斯的中心由粗阶段中 $k$ 个最近邻粗点的重心插值确定：

\begin{equation}
\bm{\mu}(t_i) = \sum_{j=1}^{k} w_j \mathbf{P}_j(t_i),
\end{equation}

其中 $\mathbf{w} = \{w_j\}_{j=1}^{k}$ 为可学习的重心权重。其余高斯参数（旋转 $\mathbf{R}$、缩放 $\mathbf{S}$、不透明度 $\mathbf{o}$、颜色 $c$）在时间上共享，完全由粗阶段的变形模型驱动。这种设计确保高斯运动与粗形状模型一致，大幅减少精阶段需要优化的参数量。

精阶段的训练无需任何 RGB 图像，核心的事件损失将渲染结果与实际事件流的亮度变化对齐：

\begin{equation}
\mathcal{L}_\text{event} = \frac{1}{N_v} \sum_{i=1}^{N_v} \| (L(t_j) - L(t_i)) - E(t_j, t_i) \|_1, \quad t_j \sim U[t_i, t_i + L_\text{max}),
\end{equation}

其中 $L(t_i) = L(\mathcal{R}(\mathcal{G}(t_i), \mathbf{V}_i))$ 将渲染像素强度经 gamma 校正转换为对数亮度，$E(t_j, t_i) = \sigma \sum_{t_i < t_k \leq t_j} p_k$ 为事件累积亮度差（按极性 $p_k$ 求和乘以阈值 $\sigma$），$t_j$ 在以 $t_i$ 为起点、$L_\text{max}=0.02$ 秒的窗口内均匀随机采样。该损失要求从任意两个时空邻近视点渲染的亮度差与事件流记录的实际亮度变化保持一致。值得注意的是，事件损失在所有像素上计算——不仅包括触发事件的像素，也包括未触发事件的像素（隐含约束"无事件即无变化"），这为大面积静止区域提供了有效监督。

由于事件在空间上的稀疏性（仅运动边缘触发），Ev4DGS 额外引入剪影损失抑制背景冗余高斯：

\begin{equation}
\mathcal{L}_\text{sil} = \sum_{i=1}^{N_v} \| \mathbf{M}_i - G(\hat{\mathbf{M}}_i) \|_1,
\end{equation}

其中 $\mathbf{M}_i$ 为高斯光栅化器渲染的二值掩码，$G(\cdot)$ 为高斯滤波函数（对 Snakes 提取的掩码 $\hat{\mathbf{M}}_i$ 进行平滑以缓解边界噪声）。精阶段的联合优化目标为 $\mathcal{G}^*, \hat{\theta}, \{\hat{\mathbf{B}}_k\} = \arg\min_{\mathcal{G}, \theta, \{\mathbf{B}_k\}} (\lambda_e \mathcal{L}_\text{event} + \lambda_s \mathcal{L}_\text{sil})$——值得注意的是，精阶段在主要学习高斯参数的同时，还对粗阶段的 MLP 权重和基形状进行微调，使变形模型和外观表示相互适配。整个框架的自监督训练仅依赖事件流和相机位姿，不需要任何 RGB 图像或预定义形状模板，这是首个证明纯事件流可以实现非刚性物体新视角渲染的工作。

\paragraph{时空解耦帧-事件协同重建}

前述 Event-boosted Deformable 3DGS 和 DEGS 虽然通过可学习阈值和事件光流分别提升了事件监督的精度和变形场的运动引导，但其核心范式仍然是"统一表示"——所有高斯都采用相同形式，动态/静态的区分（如 Event-boosted 的 DSD 策略）发生在训练后期，是对已有统一表示的事后分割。STD-GS（2025.06）~\cite{2506_23157}（SpatioTemporal-Disentangled Gaussian Splatting）从根本思路上转变：不在统一表示上做事后分离，而是先解耦场景的时空特征，再为背景和物体分别构建不同维度的高斯表示。

这一思路转变的动机来自对高动态场景中时空特征错配的深入分析。帧相机提供密集空间外观但时间运动不连续（长曝光导致纹理模糊、低帧率导致大位移），而事件相机提供稀疏空间外观但连续时间运动（微秒级异步触发）。两种模态在各自的优势维度上恰好互补——帧擅长空间、事件擅长时间。更关键的是，STD-GS 发现背景和物体在这两种模态中呈现系统性的时空差异（spatiotemporal discrepancy）：在帧图像的空间特征中，背景外观紧凑聚类、物体外观相对分散；在事件流的时间特征中，背景运动全局离散、物体运动局部聚类。这意味着统一建模会强制将两种本质不同的特征分布映射到同一表示空间，加剧错配。

基于上述分析，STD-GS 提出累积式场景-物体解耦（cumulative scene-object disentanglement）框架，包含动态场景解耦（DSD）和高斯表示融合（GRF）两个模块，从场景级到物体级逐步解耦时空特征。

\textbf{动态场景解耦（DSD）。}
与 Event-boosted 在训练后期通过 VGG19 特征相似度和 Otsu 阈值分割动态/静态区域不同，STD-GS 的 DSD 在训练前端就利用帧-事件双模态互补知识进行可学习的时空聚类。具体流程为：先用超像素将帧图像划分为局部块作为空间外观特征，用欧氏距离匹配不同帧间的外观特征；同时对事件流计算内积生成相关性体（correlation volume）作为时间运动特征，利用局部一致性（短时间内的匀速运动假设）匹配运动特征。然后以 K-Means 聚类获得背景和物体的时空特征分布作为伪标签，再引入 ClusterFormer 网络实现端到端的可学习外观-运动聚类。聚类一致性损失为
\begin{equation}
\mathcal{L}_{clu} = \sum_{i,j} \|x_i - c_{y_i}\| + \|f(x_i) - f(x_j)\| \cdot S(y_i, y_j),
\end{equation}

其中 $x_i$ 是第 $i$ 个样本，$c_{y_i}$ 是其聚类中心，$y_i$ 是 K-Means 生成的聚类标签，$f(x)$ 是 ClusterFormer 提取的特征，$S(\cdot)$ 定义两个样本标签之间的关系。该损失将先验知识（K-Means 伪标签）与深度模型（ClusterFormer）结合，使背景和物体的时空特征分离过程可端到端优化且具有可解释性。此外，方法先参照 GEM 利用事件数据恢复帧图像中高动态物体的模糊纹理，为后续空间特征提取提供更清晰的输入。

\textbf{高斯表示融合（GRF）。}
DSD 分离出的动态物体同时包含帧空间特征和事件时间特征，两种不同模态的特征无法直接渲染密集连续的物体，因此 GRF 进一步将物体解耦到统一的表示空间中。为此，GRF 首先用事件相关性体作为匹配模板，通过扩展卡尔曼滤波（EKF）沿时间维度跟踪动态物体，更新一系列连续跟踪点。然后分别构建两种高斯：静态背景映射为 3D 高斯 $\mathcal{G}_s$（仅空间维度，颜色 $\mathcal{C}_s$），动态物体映射为 4D 高斯 $\mathcal{G}_d$（包含时间维度，颜色 $\mathcal{C}_d$），以 EKF 更新的跟踪点为中心。这种维度分离的设计让背景只承担静态空间建模、物体承担时空联合建模，避免了统一表示中的特征错配。

GRF 的关键设计是高斯-事件时空一致性（spatiotemporal consistency）。STD-GS 发现一个重要现象：高斯表示和事件数据虽然形式不同，但对同一动态物体的时空特征是一致的——高斯的空间颜色直方图与事件的空间积分亮度直方图形状吻合，高斯的时间变形轨迹与事件的时间光流轨迹对齐。基于这一观察，方法将事件数据的时空特征作为伪标签，构造时空一致性损失约束物体高斯的学习：

\begin{equation}
\mathcal{L}_{consis} = \sum \|\mathcal{C}_d - L_{ev}\|_1 + \|\mathcal{F}(\mathcal{G}_d) - F_{ev}\|_1,
\end{equation}

其中 $\mathcal{F}$ 是高斯变形算子，$L_{ev}$ 和 $F_{ev}$ 分别是事件数据的空间亮度和时间光流。第一项约束颜色-亮度一致性（空间维度），第二项约束变形-光流一致性（时间维度）。最后，方法用基于阴影的权重融合背景高斯和物体高斯进行渲染。整个框架从场景级（DSD 区分背景/物体）到物体级（GRF 解耦物体高斯的时空特征并施加一致性约束）逐步精细化时空建模，实现时间连续的高动态场景重建。

\subsection{前馈重建与动态几何}

本节介绍两条技术路线：（1）前馈与大规模4D重建；（2）动态表面重建。

\subsubsection{前馈与大规模4D重建}

本路线介绍通过前馈模型实现秒级4D推理以及面向稀疏视角和大规模场景的高效重建方法。

以上方法均为逐场景优化（per-scene optimization）——针对每个动态场景单独训练一个模型，耗时数十分钟到数小时。前馈方法（feed-forward）则训练一个通用模型，能够一次前向传播从输入视频直接预测 4D 高斯参数，将推理时间压缩到秒级。本节同时讨论从稀疏视角和大场景中高效重建 4D 的方法。

\paragraph{大型前馈 Transformer 实时变形高斯}

前面介绍的方法（如 Deformable 3DGS 及其改进）都遵循\textbf{逐场景优化}范式——对每个动态场景单独迭代数千步以拟合高斯参数和变形场，耗时数十分钟到数小时。DGS-LRM（2025.06）~\cite{2506_09997}（Deformable Gaussian Splats Large Reconstruction Model）提出了一条完全不同的路线：训练一个通用的前馈（feed-forward）模型，在单次前向传播中从单目视频直接预测变形 3DGS 的全部参数，将推理时间从小时级压缩到约一秒。此前的前馈重建方法（如 GS-LRM）仅处理静态场景，无法预测动态变形；DGS-LRM 是首个将前馈范式扩展到动态场景的工作，其核心挑战是：如何在单次前向传播中\textbf{同时}预测规范空间中的 3D 高斯及其跨所有时间步的变形向量。

\textbf{变形表示。}DGS-LRM 的变形表示是逐像素的：输入视频的每个像素 $p$ 对应一个 3D 高斯 $g_p$，其参数包含深度 $d_p$、RGB 颜色、四元数旋转、缩放和不透明度。除此之外，每个像素还预测一组 3D 变形向量 $f_p = \{f_0, f_1, \ldots, f_N\}$，将 $g_p$ 从其原始时间步平移到任意目标时间步 $n$。直观地说，每个像素的 3D 高斯在规范空间中有一个"锚定位置"，变形向量则编码了它在每个时刻应该移动到哪里。渲染时间步 $n$ 的新视角时，将所有高斯按对应的变形向量平移后组成新集合：

\begin{equation}
W_n = \mathcal{W}^G_p(g_p, f_p)
\end{equation}

再经光栅化输出图像。实验发现仅变形平移即可获得高质量结果，变形旋转或不透明度无额外增益——这说明在逐像素表示下，平移已足以捕捉大部分运动模式。

\textbf{架构与时间 token 化。}将多帧视频输入标准 Transformer 需要海量 GPU 显存，而减少输入帧数（即缩小时间窗口 $N$）又会缩短相机基线，加剧几何歧义。DGS-LRM 用\textbf{时间 token 化}（temporal tokenization）同时解决这两个矛盾：不将每帧独立切分为 $s \times s$ 的图像块（patch），而是将视频视为三维体积，切分为 $s \times s \times l$ 的时空立方体，再通过线性层映射为单个 token。实践取 $l{=}4$，即每 4 帧共享一组 token，token 数量减少为四分之一，大幅降低显存和计算开销。此外，为在不增加 token 数的前提下扩大相机基线、减少几何歧义，模型还引入时间上远距离的\textbf{参考帧} $R$（仅提供几何线索，不对其预测高斯）。编码后的 token 送入 24 层多头自注意力 Transformer，所有参数施加权重归一化（weight normalization）以稳定训练，最后经两层 MLP 投影为变形高斯参数。

\textbf{训练数据与监督。}由于真实世界缺乏带有多视角真值和 3D 场景流标注的动态视频，DGS-LRM 基于 Kubric 引擎（MOVi-E 配置）生成大规模合成动态场景。每场景使用 4 台同步相机从不同视角渲染（相机间距 4--16 米，共享注视点以保证共视性）。为缩小合成-真实域差距，对其中一台相机添加运动模糊（仅用作输入、不作监督），焦距在 25mm--55mm 之间随机采样。训练时采用\textbf{双视角采样}（dual-view sampling）策略：从另一台同步相机同时采样两个输出视角（共享时间戳），经验性地发现这比单视角监督显著提升训练收敛速度。

损失函数包含四项：

\begin{equation}
\mathcal{L}_\text{total} = \mathcal{L}_\text{mse} + \lambda_d \mathcal{L}_d + \lambda_f \mathcal{L}_f + \lambda_p \mathcal{L}_\text{lpips}
\end{equation}

其中 $\mathcal{L}_\text{mse}$ 和 $\mathcal{L}_\text{lpips}$（$\lambda_p{=}0.5$）分别是像素级 MSE 损失和感知损失，用于监督渲染质量。深度损失 $\mathcal{L}_d = \sum_{g_p}\|d^*_p - d_p\|_1$ 直接用像素对齐的真值深度 $d^*_p$ 监督预测深度（比渲染深度监督更有效）。场景流损失 $\mathcal{L}_f = \sum_{g_p}\|f^*_p - f_p\|_1$ 直接用真值 3D 场景流 $f^*_p$ 监督变形向量（$\lambda_d{=}10$, $\lambda_f{=}10$），解决了仅用渲染损失时变形轨迹不连续甚至高斯飞出相机视锥的问题。场景流以稀疏张量存储，节省 80\% 内存和 I/O 开销。

此外，为消除训练和推理时场景尺度不一致导致的泛化问题，DGS-LRM 引入场景归一化：用单目度量深度估计器估计场景尺度，将深度第 20 百分位的视差归一化为 2，统一应用到相机位姿、真值深度和场景流上。

推理仅需约一秒，比逐场景优化方法快数千倍，在 DyCheck 数据集上比 L4GM 高 3 dB PSNR。此外，通过将相邻片段的场景流拼接（flow chaining），DGS-LRM 可将 1 秒片段扩展到长视频的稠密 3D 跟踪。

\paragraph{大型 4D 高斯重建 Transformer}

DGS-LRM 采用逐像素变形表示（规范空间高斯 + 变形向量），通过时间 token 化的 Transformer 同时预测静态结构和动态变形，核心挑战在于变形场的精确学习。一个自然的问题是：能否绕过变形场的分解，直接为每个时间步独立预测一组完整的 3D 高斯？这样做的好处是避免了变形场的误差累积，代价则是需要更强大的 4D 先验来保证时间一致性。L4GM（2024.06）~\cite{2406_10324}（Large 4D Gaussian Reconstruction Model）给出了肯定的回答——它放弃了 DGS-LRM 的"规范空间 + 变形"分解策略，转而为每个时间步独立输出一组 3D 高斯，4D 表示即为 $T$ 组高斯的集合 $P = \{P_t\}_{t=1}^T$。时间一致性不是通过变形场的显式约束实现的，而是通过在架构中嵌入时间自注意力来隐式学习。这一设计选择使得 L4GM 可以直接复用已有的 3D 大型重建模型作为预训练起点。

\textbf{预训练策略：从 3D 到 4D 的迁移。}L4GM 的核心灵感来自视频扩散模型的训练范式：先在丰富得多的静态图像/3D 数据上预训练，再扩展到视频域微调以建模时间一致性。类似地，L4GM 以预训练的 3D 大型多视角高斯模型 LGM 为起点——LGM 本身已在大规模静态 3D 物体数据上学会了强大的多视角 3D 重建能力，能从多视角图像直接输出 3D 高斯表示（每个高斯含 14 个参数：中心 $z \in \mathbb{R}^3$、缩放 $s \in \mathbb{R}^3$、四元数旋转 $q \in \mathbb{R}^4$、不透明度 $o$、颜色特征 $c \in \mathbb{R}^3$）。消融实验表明，若不使用 LGM 预训练权重初始化（即完全从头训练），模型甚至无法收敛——这说明在有限的 4D 数据上同时学习 3D 理解和时间动态过于困难，利用 3D 预训练是不可或缺的基础。

\textbf{架构：时间自注意力驱动的 U-Net。}L4GM 保留 LGM 的非对称 U-Net 骨干，关键修改是在每个跨视角自注意力层之后插入时间自注意力层。输入构造如下：用 ImageDream 从单目视频第一帧 $I_1$ 生成 $V{=}4$ 个正交视角的多视角图像 $J_1$，然后将 $J_1$ 中除 $I_1$ 外的生成视角复制到所有时间步（仅首帧使用真实视频帧），形成 $T \times V$ 网格输入。这里隐含一个简化假设：参考视频中的相机是静止的、仅物体运动，因此除首帧外的相机位姿 $O$ 可跨时间步复制。位姿经 Pl\"ucker 射线嵌入编码后与 RGB 通道拼接。时间自注意力的关键操作是将视角轴 $V$ 视为独立视频的批量维度：

\begin{equation}
x = \text{rearrange}(x, (B\, T\, V)\, H\, W\, C \to (B\, V)(T\, H\, W)\, C), \quad x = x + \text{Attn}(x)
\end{equation}

这使得不同时间步的特征能够在共享空间特征的基础上交换时间信息，隐式地学习时间一致的 3D 表示。消融实验表明，若移除时间自注意力（退化为逐帧独立重建），模型仍可从多视角输入获得不错的单帧质量，但新视角视频会出现明显的闪烁——时间一致性无法仅靠输入中的时序信息隐式获得。U-Net 输出 14 通道特征图，每像素对应一个高斯参数集，沿视角维度拼接后形成每时间步一组 3D 高斯，共 $T$ 组构成 4D 表示。训练时取 $T{=}8$，推理时使用 $T{=}16$。

\textbf{训练数据。}L4GM 的训练数据集规模空前：从 Objaverse 1.0 策划约 44K 物体、110K 动画，每个动画从 48 个视角渲染，总计约 1200 万段多视角视频。如此大规模的合成数据是前馈 4D 重建泛化能力的关键基础——与 DGS-LRM 使用 Kubric 合成数据不同，L4GM 的数据来自真实的 3D 动画资产库，物体种类和运动模式更为丰富。

\textbf{损失与端到端微调。}训练监督使用额外一组相机位姿 $O'$ 的多视角视频渲染损失（MSE + LPIPS）。值得注意的是，L4GM 选择端到端微调整个 U-Net（而非仅训练新增的时间注意力层）——消融表明冻结 LGM 原始层仅训练时间注意力虽然前期收敛更快，但最终收敛到更低的 PSNR，说明 4D 任务需要整体调整 3D 先验以适应动态场景。此外，L4GM 探索了变形场（HexPlane）替代方案，发现在大规模训练下变形场分支完全无法学到合理的形变，输出始终静态——这与逐场景优化方法中变形场广泛成功的现象形成鲜明对比。推测原因是 SDS 方法依赖隐式表示的平滑性进行正则化，而 L4GM 的显式高斯表示已从 4D 训练数据中直接学到了平滑输出的能力。

\textbf{长视频与 4D 插值。}L4GM 还支持自回归长视频重建和 4D 插值。自回归模式下，将上一段最后一帧的高斯从四个正交角度渲染为下一段的多视角输入（替代 ImageDream 重新生成），依次处理直到所有帧完成。由于不做跨帧高斯跟踪，直接插值高斯轨迹不可行，因此 L4GM 训练了独立的 4D 插值模型：输入相邻两组高斯，输出中间帧的高斯，新帧的颜色取相邻多视角图像的 RGB 像素加权平均。推理时每前向传播约 0.3 秒，插值模型每两帧仅需 0.065 秒，可在约 10 次自回归后仍保持质量。

尽管仅在合成数据上训练，L4GM 对 Sora 生成的视频和真实视频均展现了出色的泛化能力，在 Consistent4D 基准上达到最优质量，在 LPIPS、CLIP 相似度和 FVD 指标上全面领先，同时速度比基于优化的方法快 100--1000 倍。用户研究进一步证实了 L4GM 在整体质量、3D 外观、视频对齐和运动真实性上的优势。一个显著差异是分辨率：基于优化的方法通常仅使用数千个高斯以保持优化可行，而 L4GM 的前馈方式可轻松为每帧输出超过 60,000 个高斯，在 3D 分辨率上大幅领先。

\paragraph{稀疏 4 相机 4D 高斯泼溅}

前面的 DGS-LRM 和 L4GM 通过前馈范式将 4D 重建压缩到秒级，但它们的泛化能力依赖大规模合成预训练数据，对真实世界复杂场景的适应性仍有限。另一个同样关键但方向不同的问题是\textbf{硬件门槛}：无论逐场景优化还是前馈方法，已有 4DGS 框架几乎都假设密集多视角捕捉（18--48 台同步摄像机），将动态场景建模限制在实验室或大型公司。4C4D（2026.04）~\cite{2604_04063}（4 Camera 4D Gaussian Splatting）回到了逐场景优化范式，但将视角需求从数十台压缩到仅四台消费级摄像机——整个捕捉系统的成本低于 1500 美元，普通用户用手机或 GoPro 即可部署。核心挑战是：在极端稀疏视角下，4DGS 优化会遭遇严重的\textbf{有偏优化}（biased optimization）。

\textbf{稀疏视角下的有偏优化问题。}4C4D 的出发点是一个细致的诊断实验：用标准 4DGS 在仅有四台相机的设定下训练场景，然后对比训练视角和测试视角的渲染质量。结果显示，训练视角的外观（RGB 颜色）能被准确还原，但深度几何却充满噪声和不一致——在测试视角上产生明显的伪影。这说明在空间监督不足时，4DGS 的优化天然偏向更容易的外观拟合，而忽略更困难的几何恢复。根本原因在于：4DGS 原有的高斯不透明度参数在几何学习中扮演关键角色（不透明度决定高斯对深度图贡献的权重），但其优化被外观损失主导，几何梯度被淹没。

\textbf{神经衰减函数。}为纠正这种偏差，4C4D 提出\textbf{神经衰减函数}（Neural Decaying Function）——一个可学习的小型神经网络 $f_\theta$，用于自适应控制每个 4D 高斯的不透明度衰减。其输入为高斯的关键属性（中心 $(x,y,z)$、不透明度 $o$、旋转 $r$），输出衰减因子 $\sigma$。经神经衰减调制后，4D 高斯在时间步 $t$ 的最终不透明度为：

\begin{equation}
o(t) = \sigma \cdot \exp\!\left(-\frac{1}{2}\left(t - t_t\right)^2 \sigma_{4,4}^{-1}\right) \cdot o
\end{equation}

其中 $t_t$ 是 4D 高斯的时间中心，$\sigma_{4,4}$ 是时间协方差。直觉上，原始 4DGS 中不透明度的衰减由固定的指数时间核控制，而 4C4D 用可学习的 $\sigma$ 替代了固定衰减——通过神经控制不透明度这一几何学习的关键参数，将优化梯度重新分配到几何方向，从机制层面纠正外观-几何的优化不平衡。消融实验证实，可学习的神经衰减函数显著优于常数衰减（$\sigma{=}0.9$）和手动设计的指数/幂函数衰减。

\textbf{分离衰减与可见性检测。}然而，简单地将神经衰减函数均匀应用于所有 4D 高斯会导致性能下降。原因在于：4DGS 的梯度仅存在于当前视角和时间步下\textbf{参与渲染的可见高斯}区域，对不可见区域施加相同的可学习衰减会扭曲梯度方向。因此 4C4D 引入\textbf{时空可见性检测}策略，联合视角交叉 $Z_V$ 和时间交叉 $Z_T$ 筛选出当前视角和时间步下的可见高斯集合 $G_m$：

\begin{equation}
G_m = Z_V(\pi, t, Z_T(t, s_t, G))
\end{equation}

其中 $Z_T$ 根据每个高斯的时间跨度 $s_t$ 过滤掉时间上不活跃的高斯，$Z_V$ 根据相机视角 $\pi$ 过滤掉空间上不可见的高斯。对可见高斯使用神经衰减函数 $f_\theta$，对不可见高斯使用小常数衰减，构成\textbf{分离衰减策略}：

\begin{equation}
\sigma(g) = \begin{cases} f_\theta(x,y,z,o,r) & g \in G_m \\ 0.999 & g \in G_m^* \end{cases}
\end{equation}

其中 $G_m^* = G - G_m$ 为不可见高斯集合。常数 $\sigma{=}0.999$ 为不可见高斯提供轻微的稳定性衰减（与 Binocular3DGS 的观察一致），避免其参数在无梯度时漂移。这种分离设计确保可见区域的梯度正确更新神经衰减函数和 4D 高斯参数，不可见区域保持稳定。

\textbf{实现与实验。}4C4D 使用 MAST3R 进行稀疏视角初始化（比传统 COLMAP 更稳定），神经衰减函数在训练 500 次迭代后引入以提供可靠的热身监督。方法在四个数据集上验证：Neural3DV（6 个动态场景，21 台相机中选 4 台）、ENeRF-Outdoor（真实室外场景，含光照变化和背景杂乱）、Mobile-Stage（大规模多人舞蹈场景，24 台相机中选 4 台），以及自采集的 Dyn4Cam（仅用四台 GoPro 相机拍摄，60 FPS，每场景 300 帧）。4C4D 在所有数据集和所有指标（PSNR、DSSIM、LPIPS）上均显著优于 4DGaussians、Ex4DGS、ST-GS 和 4DGS 等基线方法，尤其在快速运动和多人交互等困难场景中展现出更干净的几何和更稳定的时间一致性，大幅降低了动态场景捕捉的硬件门槛。

\paragraph{整流流无监督 4D 重建}

前面的 DGS-LRM 和 L4GM 证明了前馈范式可以将 4D 重建从小时级压缩到秒级，但它们面向的是桌面级物体或合成场景，且依赖大规模合成监督数据（DGS-LRM 需要 Kubric 渲染的场景流真值，L4GM 需要 Objaverse 的 1200 万段多视角视频）。自动驾驶等大规模室外场景面临更严苛的约束：既无法获取密集多视角真值标注，也不能依赖预训练视觉模型（其预测噪声会注入重建管线并难以恢复）。Flux4D（2025.12）~\cite{2512_03210}（Flow-based Unsupervised 4D Reconstruction）从 DGS-LRM 的前馈路线出发，但将目标场景从桌面级合成数据转向真实室外驾驶场景，同时彻底消除对任何标注和预训练模型的依赖。其核心动机是一个反直觉的发现：\textbf{跨场景多样数据本身就是最有效的隐式正则化}——无需精心设计的几何约束、循环一致性或语义先验，仅在足够多样化的场景上联合训练，网络就能自动学会区分静态背景和动态前景。

\textbf{场景表示与初始化。}Flux4D 的输入是自动驾驶平台采集的多相机图像 $I$ 和 LiDAR 点云 $P$ 序列。与 DGS-LRM 从视频像素逐像素预测高斯不同，Flux4D 直接在 3D 空间中初始化高斯：从每帧 LiDAR 点云反投影得到初始位置，投影到对应相机图像采样颜色，缩放初始化为最近三个邻域点的平均距离（经对数变换和 softplus 稳定训练），不透明度初始化为 0.5。关键地，每个高斯被赋予两个额外参数——可学习的瞬时速度 $v_i \in \mathbb{R}^3$ 和固定的捕获时间戳 $t_i$（归一化到 $[0,1]$）。此外，为建模天空和远景，在场景包围盒上方的半球面上随机采样 100 万个点作为补充高斯。

\textbf{速度场与运动传播。}初始高斯集合 $G_\text{init}$ 送入前馈网络 $f_\theta$，网络同时输出高斯参数的残差修正 $\Delta G$ 和速度场 $V$。其中网络采用稀疏 3D U-Net（基于 TorchSparse 稀疏卷积实现），输入为 $N \times 14$ 的张量（3D 位置、四元数旋转、缩放、不透明度、RGB 颜色），输出残差叠加到初始高斯上得到精修后的高斯 $G = G_\text{init} + \Delta G$。有了预测速度后，每个高斯可通过线性运动模型传播到任意查询时刻：

\begin{equation}
p_i^{t'} = p_i^{t_i} + v_i (t' - t_i)
\end{equation}

其中 $p_i^{t'}$ 是高斯在时刻 $t'$ 的位置。这种恒速假设虽然简单，但对于自动驾驶场景中约 1 秒的短时间窗口已足够有效——这与 DGS-LRM 的逐像素变形向量方案形成对比：DGS-LRM 为每个时间步预测独立的变形向量（参数量与时间步数线性增长），而 Flux4D 仅预测一个瞬时速度（参数量固定为 3 维），通过时间插值传播到所有时刻。当场景需要建模更复杂的运动（如加速、转弯）时，Flux4D 进一步支持多项式运动模型，预测速度参数 $\{v_0, v_1, \ldots, v_d\}$，位移变为 $p_i^{t'} = p_i^{t_i} + \sum_{j=0}^{d} v_j \frac{({t'}^{j+1} - {t_i}^{j+1})}{j+1}$，实践中取 $d{=}1$（恒加速假设）即可获得显著提升，更高阶项无额外增益。

\textbf{完全自监督训练。}Flux4D 的训练无需任何标注，总损失仅包含两项：

\begin{equation}
\mathcal{L} = \mathcal{L}_\text{recon} + \lambda_\text{vel} \mathcal{L}_\text{vel}
\end{equation}

其中重建损失 $\mathcal{L}_\text{recon} = \lambda_\text{rgb} \mathcal{L}_1^\text{img} + \lambda_\text{SSIM} \mathcal{L}_\text{SSIM} + \lambda_\text{depth} \mathcal{L}_1^\text{depth}$ 包含图像 $L_1$ 损失、SSIM 损失和深度 $L_1$ 损失（将 LiDAR 点云投影到图像平面作为深度监督）。速度正则化 $\mathcal{L}_\text{vel} = \frac{1}{M}\sum_i \|v_i\|_2$ 鼓励高斯尽量保持静止——即``尽量不动的先验''（as-static-as-possible）。权重设置为 $\lambda_\text{rgb}{=}0.8$、$\lambda_\text{SSIM}{=}0.2$、$\lambda_\text{depth}{=}0.01$、$\lambda_\text{vel}{=}5 \times 10^{-3}$。与 DGS-LRM 需要真值场景流监督不同，Flux4D 仅用渲染损失和速度正则化就实现了动静态分解——关键在于跨场景联合训练提供了充分的隐式正则化。实验的缩放分析证实：随着训练场景数量增加，重建质量和运动估计精度持续单调提升。

\textbf{迭代精修。}基础模型 Flux4D-base 的单步前馈能力有限（因网络容量和遮挡导致的不完美初始化），Flux4D 借鉴 G3R 的思路引入迭代精修机制：在首轮前向传播后，计算高斯参数关于损失的 3D 梯度 $\nabla G$，将其与当前高斯一同送入精修网络 $r_\phi$（与 $f_\theta$ 共享相同架构但独立参数），预测残量修正。整个过程迭代 $N{=}3$ 轮（而 G3R 需要 24 轮），逐步修正颜色不一致和细节缺失。最终高斯为 $G = G_\text{init} + \sum_{j=1}^{N} \Delta G^{(j)}$。

\textbf{速度重加权。}为改善动态区域的重建质量，Flux4D 对光度损失施加像素级重加权：$\mathcal{L}_\text{recon}^\text{vw} = (1 + \|\text{sg}(v_r)\|) \cdot \mathcal{L}_\text{recon}$，其中 $\text{sg}(\cdot)$ 为停止梯度操作，$v_r$ 是渲染到图像平面的速度。这使得运动更快的区域（通常占据更少像素、对总损失贡献更小）获得更高的训练权重，有效缓解了动态前景被静态背景淹没的问题。

\textbf{实验表现。}Flux4D 处理全分辨率图像（$1920 \times 1080$）并支持 20+ 帧输入，在 PandaSet 和 Waymo Open Dataset 上均大幅超越无监督基线：在 WOD 上比 DrivingRecon 高 5.99 dB PSNR，在 PandaSet 上甚至在多项指标上与需要标注的监督方法竞争力相当。推理在单块 RTX A5000 上约 3.9 秒完成（含 3 轮迭代精修），比需要预训练视觉模型的 DrivingRecon（480 A100 GPU 时训练）和 STORM（768 A100 GPU 时）高效得多——Flux4D 仅需 140 L40S GPU 时。此外，Flux4D 的速度场可直接用于场景流估计，在 PandaSet 和 WOD 的标准场景流指标上超越了专门的无监督场景流方法（如 NSFP）。通过 DBSCAN 对速度场聚类，还可实现无标注的实例分割和可控仿真编辑（物体移动、删除、相机轨迹修改）。Flux4D 也展示了未来帧预测能力——在观测帧之外外推最多 5 帧，且质量超过依赖标注的监督方法，验证了``跨场景数据即隐式正则化''的核心洞察。

\paragraph{稀疏帧动态场景联合重建}

前面的 DGS-LRM、L4GM、Flux4D 分别从前馈预测、大规模预训练、无监督学习等维度推进了 4D 重建，但它们（及本章此前介绍的所有变形高斯方法）都隐含一个共同假设：输入视频提供足够密集的时间采样。当设备限制导致帧率降至 5 FPS 甚至更低时，变形网络面临的困难不是简单的信息不足——而是一个更微妙的\textbf{选择性退化}：纹理丰富区域（物体边缘、高频纹理）的几何最先塌陷。这些区域的高频信号本应是变形建模的关键线索，却在稀疏帧条件下最容易被变形网络忽略，因为跨帧外观变化太大，变形场难以建立可靠的时间对应。Sparse4DGS（2025.11）~\cite{2511_07122} 回到逐场景优化范式，沿袭 Deformable 3DGS 的"规范空间 + 变形场"框架，但针对稀疏帧条件的特殊退化模式提出\textbf{纹理感知}（texture-aware）优化策略——其核心洞察是：在信息稀缺的稀疏帧场景下，高频纹理信号不是噪声，而是变形建模最可靠的细节来源和动态线索，应主动引导高斯聚焦于纹理丰富区域。方法提出三个协同的纹理感知组件。

\textbf{纹理强度高斯场。}为在 3D 空间中显式表示纹理丰富度，Sparse4DGS 引入纹理强度（Texture Intensity, TI）高斯场。首先对每帧输入图像 $I \in \mathbb{R}^{H \times W \times 3}$，分别用水平和垂直 Sobel 算子计算梯度图，再取逐像素梯度幅值：

\begin{equation}
\text{TI}_\text{gt}(i,j) = \sqrt{\text{TI}_x(i,j)^2 + \text{TI}_y(i,j)^2}
\end{equation}

作为该像素处局部纹理丰富度的显式度量。然后将 TI 作为新属性嵌入每个 3D 高斯，通过可微光栅化将其渲染为 $\text{TI}_\text{render}$。由于纹理提取是对每帧独立进行的，不同视角的 TI 图之间存在空间不一致性，传统 $L_1$ 损失会放大这种不一致。Sparse4DGS 改用 Pearson 相关系数（PCC）度量 GT 与渲染 TI 图的相对变化率一致性：

\begin{equation}
\mathcal{L}_\text{tex} = 1 - \rho(\text{TI}_\text{gt},\, \text{TI}_\text{render})
\end{equation}

PCC 关注变量间的协方差与标准差之比，忽略绝对数值差异，天然容忍多帧间的空间不一致性。消融实验表明，将 PCC 替换为 $L_1$ 损失导致 PSNR 下降 0.6 dB，验证了这一设计选择的有效性。

\textbf{纹理感知变形正则化（TADR）。}变形场的质量直接决定稀疏帧下的时间插值效果。已有深度正则化方法直接计算渲染深度图 $D_\text{render}$ 与单目深度估计器（DPT）预测深度图 $D_\text{dpt}$ 之间的 PCC，但这类图像级约束在捕捉局部深度变化方面能力有限。TADR 的思路是：纹理变化通常与深度变化高度相关（物体边缘既是纹理突变处，也是深度不连续处），因此从深度图中提取纹理强度图，比直接比较深度图本身能提供更精确的局部几何约束。具体地，TADR 对 $D_\text{render}$ 和 $D_\text{dpt}$ 分别用 Sobel 算子提取纹理强度图 $\text{TI}_\text{render}^\text{depth}$ 和 $\text{TI}_\text{gt}^\text{depth}$，再用 PCC 约束两者一致：

\begin{equation}
\mathcal{L}_\text{tadr} = 1 - \rho(\text{TI}_\text{gt}^\text{depth},\, \text{TI}_\text{render}^\text{depth})
\end{equation}

这种纹理感知的深度约束比直接 PCC 深度损失更精确地捕捉局部几何变化。总体训练损失为：

\begin{equation}
\mathcal{L} = \mathcal{L}_\text{rgb} + \lambda_1 \mathcal{L}_\text{tex} + \lambda_2 \mathcal{L}_\text{tadr}
\end{equation}

其中 $\mathcal{L}_\text{rgb} = (1{-}\lambda)\mathcal{L}_1(\hat{I}, I) + \lambda \mathcal{L}_\text{SSIM}(\hat{I}, I)$ 为标准 RGB 损失，超参数 $\lambda_1{=}\lambda_2{=}0.01$。

\textbf{纹理感知规范空间优化（TACO）。}前两个组件从损失函数层面引导纹理一致性，TACO 则从优化器层面直接控制高斯的分布——将高斯推向纹理丰富区域并使其稳定在那里。其理论基础是随机梯度 Langevin 动力学（SGLD）：传统 SGD 的更新规则为 $\theta_g \leftarrow \theta_g - \eta_g \nabla_{\theta_g} E_{II}$，SGLD 在此基础上引入随机噪声项以实现探索与收敛的平衡。此前 3DGS 中的 SGLD 应用仅使用不透明度相关的噪声 $\sigma_o = \sigma(-k(o - t))$（$k{=}100$，$t{=}0.995$）来消除低不透明度的模糊高斯。TACO 在保留 $\sigma_o$ 的基础上，增加纹理感知噪声项：

\begin{equation}
\theta_g \leftarrow \theta_g - \eta_g \nabla_{\theta_g} E_{II} + \epsilon_\text{noise}(\sigma_\text{tex} + \sigma_o)
\end{equation}

其中 $\sigma_\text{tex} = \sigma(-k(\text{TI} - t))$。关键机制是：TI 值越大（纹理越丰富），$\sigma_\text{tex}$ 越接近零，高斯受到的扰动噪声越小——这意味着当高斯迁移到纹理丰富区域后，其位置会被噪声"锁定"。反之，TI 值低的高斯持续受到较大噪声扰动，被不断推离纹理贫乏区域。由于 TI 的值域与不透明度相同，TACO 直接复用与 $\sigma_o$ 相同的超参数 $k$ 和 $t$，无需额外调参。整个过程持续进行，直到高斯同时获得高不透明度和高 TI 值——即同时满足几何稳定性和纹理丰富性。消融实验证实：移除 TACO 导致 PSNR 下降 1.01 dB，是三个组件中贡献最大的；移除 $\sigma_\text{tex}$ 项则导致高斯无法在纹理丰富区域聚集。

\textbf{实验表现。}Sparse4DGS 在四个数据集上验证：NeRF-Synthetic（8 个合成场景）、NeRF-DS（7 个真实场景）、HyperNeRF（5 个真实场景），以及自采集的 iPhone-4D（4 个真实场景，30 FPS 采集）。训练时从原始帧中均匀采样稀疏子集（20/30/40 帧），评估集保持不变。在 NeRF-Synthetic 和 NeRF-DS 上，Sparse4DGS 比次优方法提升超过 1 dB PSNR，在纹理丰富的物体边缘区域尤其显著——其他基线在动态物体边缘产生明显模糊和伪影，而 Sparse4DGS 保持清晰的结构细节。iPhone-4D 数据集进一步验证了真实帧率应用场景：无论 5 FPS（极低帧率）还是 30 FPS（标准帧率）输入，Sparse4DGS 均一致优于 Deformable 3DGS、4DGaussians 和 CoRGS 等基线，在 MacBook 场景中成功重建了清晰的 Apple 标志，而其他方法结果模糊。这些结果表明，纹理感知策略有效弥补了稀疏帧条件下的信息缺失，使稀疏帧动态场景重建成为可能。

\paragraph{骨架驱动稀疏视角 4D 重建}

前面讨论的 DGS-LRM、L4GM、4C4D 和 Flux4D 分别从不同角度解决动态场景重建——前馈范式将推理压缩到秒级，稀疏视角优化降低硬件门槛，无监督训练消除标注依赖。但它们共享一个隐含假设：输入视频在\textbf{时间维度上是密集的}（连续帧或短间隔帧），从而提供充足的运动线索和时序对应。现实中另一种更具挑战性的场景是\textbf{稀疏时间观测}：例如监控摄像头往往只在少数时间步从完全不同的视角各捕捉一张图片，帧间外观变化剧烈，时序对应几乎无法建立。在这种极端设定下，DGS-LRM 的逐像素变形向量缺乏足够的帧间信息来学习有意义的轨迹，Flux4D 的恒速假设无法捕捉大时间间隔间的复杂运动，4C4D 的多视角几何约束在每时间步仅一图时也完全失效。SV-GS（2026.01）~\cite{2601_00285}（Sparse View 4D Reconstruction with Skeleton-Driven Gaussian Splatting）的核心洞察是：当时间信息极度匮乏时，必须引入\textbf{骨架结构的低维先验}弥补——骨架将高维的逐点变形压缩为少量关节旋转参数（$J$ 个关节 $\times$ 4 维四元数，而非 $N$ 个高斯 $\times$ 3 维位移），大幅降低稀疏观测下的优化自由度，使变形推断从病态变为可解。

\textbf{问题设定与输入。}SV-GS 面向铰接式动态物体的 4D 重建：输入稀疏时间观测 $\mathcal{I}=\{I_t\}_{t \in [0,1]}$，每个时间步仅有\textbf{一张}来自任意视角的位姿已知图像。方法需要两个额外输入：（1）粗略骨架图 $\mathcal{F}$——指定 $J$ 个节点的 3D 位置和父子连接关系，可通过人工标注或现成工具（如 RigNet、RigAnything）获取，但\textbf{不含}关节参数、蒙皮权重或点-部件关联；（2）第一帧的静态 3DGS 重建 $\mathcal{G}$，可通过多视角图像的标准 3DGS 流程获得，或后续用预训练扩散模型从单图生成（放松多视角需求）。注意：输入骨架可以是有噪声的，且不构成完整的绑定模型（rigged model），关节姿态、蒙皮权重和逐点-骨骼关联均需在优化中学习。

\textbf{粗到精的骨架驱动变形场。}SV-GS 的变形场采用三级递进结构，将骨架约束从粗粒度传播到逐高斯细粒度：

\textit{第一级：关节姿态预测器}（$MLP_\Theta$，唯一依赖时间的模块）。以时间 $t$ 的位置编码 $\gamma(t)$ 为输入，预测每个关节 $j$ 的局部旋转 $q_j^t$（四元数）和根关节平移 $p^t \in \mathbb{R}^3$：

\begin{equation}
\{q^t, p^t\} = MLP_\Theta(\gamma(t))
\end{equation}

局部旋转定义在各关节的局部坐标系中，通过正向运动学 $fk(\cdot)$ 沿骨架的父子层级传播，得到每个关节在全局坐标系中的变换：

\begin{equation}
\{\hat{\mathbf{R}_j^t}, \hat{T_j^t}\} = fk(\mathcal{F}, q^t, p^t)
\end{equation}

其中 $\hat{\mathbf{R}_j^t} \in \mathbb{R}^{3\times3}$ 和 $\hat{T_j^t}$ 分别是关节 $j$ 在时刻 $t$ 的全局旋转矩阵和平移。关键设计：$MLP_\Theta$ 是整个框架中\textbf{唯一}以时间为条件的模块——这意味着运动被参数化为一个关于时间的连续函数，推理时可直接在未观测时间步插值关节姿态（只需前向传播新的 $t$），实现平滑运动合成，而无需学习任何离散的帧间对应。

\textit{第二级：可学习线性混合蒙皮（LBS）。}骨架中 $J$ 个关节构成 $B$ 条骨骼（每条骨骼 $b_j$ 对应关节 $j$ 与其父节点的连线）。由于输入骨架缺乏蒙皮信息，SV-GS 为每条骨骼建模可学习的蒙皮权重——以径向基函数（RBF）为基础、MLP 修正场补偿骨架噪声。具体地，未归一化权重为：

\begin{equation}
\hat{w}_{i,j} = \Delta w_{i,j} \cdot \exp\!\left(-\frac{d_{i,j}^2}{2r_j^2}\right)
\end{equation}

其中 $d_{i,j}$ 是规范空间中高斯中心 $\mu_i$ 到骨骼 $b_j$ 的距离，$r_j$ 是每条骨骼的可学习影响半径，$\Delta w_{i,j} = MLP_\Phi(\gamma(\mu_i))$ 是位置依赖的修正场（$\gamma(\cdot)$ 为位置编码）。归一化后 $w_{i,j} = \hat{w}_{i,j} / \sum_{j=1}^{B} \hat{w}_{i,j}$ 满足 $\sum_j w_{i,j} = 1$。高斯中心通过线性混合蒙皮变形：

\begin{equation}
\mu_i^t = \sum_{j=1}^{B} w_{i,j}(\hat{\mathbf{R}_j^t}\, \mu_i + \hat{T_j^t})
\end{equation}

高斯旋转同理通过权重加权的旋转混合近似：$\sum_{j=1}^B w_{i,j} \hat{\mathbf{R}_j^t}\, \mathbf{R}_i$。与 DGS-LRM 的逐像素变形向量（每个时间步独立预测 3D 位移，参数量与时间步数线性增长）不同，SV-GS 的 LBS 参数（蒙皮权重和骨骼影响半径）完全\textbf{不依赖时间}——它们在规范空间中学习一次，所有时刻的变形均由关节姿态驱动，参数效率远高于逐时间步变形方案。

\textit{第三级：姿态相关细节变形场}（$MLP_\Psi$）。骨架驱动变形仅捕捉粗粒度的铰接运动，无法表达布料飘动、肌肉形变等细粒度非刚性效应。细节变形场为每个高斯预测一个偏移量：

\begin{equation}
\mu_i^t = \bar{\mu}_i^t + MLP_\Psi(\gamma(\mu_i), \hat{q}^t)
\end{equation}

其中 $\bar{\mu}_i^t$ 是经 LBS 变形后的位置，$\hat{q}^t = \{q_j^t\}$ 是关节局部旋转的拼接。细节变形场以关节姿态为条件（而非直接以时间为条件），确保学习到的非刚性变形与姿态一致——当两个不同时间步有相似的关节姿态时，细节变形也保持一致，避免时间维度上的过拟合。

\textbf{训练与监督。}总损失包含三项：

\begin{equation}
\mathcal{L} = \lambda_1 \mathcal{L}_\text{perceptual} + \lambda_2 \mathcal{L}_\text{motion} + \lambda_3 \mathcal{L}_\text{detail}
\end{equation}

渲染感知损失 $\mathcal{L}_\text{perceptual}$ 为 $L_1$ + D-SSIM 的组合，监督稀疏观测时刻的渲染质量。关节运动时序平滑正则 $\mathcal{L}_\text{motion} = \frac{1}{T}\sum_t \sum_j |q_j^{t-1} - 2q_j^t + q_j^{t+1}|$ 施加时间拉普拉斯约束，确保关节轨迹在时间上平滑（二阶差分趋于零）——这在稀疏观测（可能仅有 5--10 个时间步）下尤为关键，防止 $MLP_\Theta$ 在无监督的中间时刻产生不合理的姿态跳变。细节变形 $\mathcal{L}_2$ 正则约束偏移量幅度，防止细节场退化为弥补骨架误差的``补丁''。

\textbf{放松多视角初始化需求。}SV-GS 进一步展示，第一帧的多视角静态重建可被预训练扩散模型（如 Zero-1-to-3、DreamGaussian）从单图生成的 3DGS 替代——仅需一张图片即可完成初始化，使方法在真实场景中更为实用。生成的初始高斯虽质量低于多视角重建，但骨架驱动变形场仍能通过稀疏时序观测逐步精修几何细节。

\textbf{实验表现。}SV-GS 在合成数据集上比现有密集观测方法 PSNR 提升高达 34\%（在仅 5--10 个稀疏时间步的设定下），大幅超过 SK-GS 和 RigGS 等同样使用骨架驱动的视频方法（它们在稀疏设定下因缺乏连续帧的时序对应而严重退化）。在真实数据上，SV-GS 用显著更少的帧即可达到与密集单目视频方法可比的性能。消融实验证实三个组件（RBF 蒙皮权重 + MLP 修正场、关节运动正则、细节变形场）各自不可或缺：去掉 MLP 修正场时 PSNR 下降约 1.5 dB（验证了骨架噪声补偿的必要性），去掉运动正则时出现明显的姿态抖动，去掉细节变形场时非刚性区域质量显著下降。

\paragraph{4D 运动脚手架与高斯绑定}

前面的方法（Deformable 3DGS 及其改进、DGS-LRM、L4GM 等）或假设已知相机位姿，或依赖大规模合成预训练数据。然而，真实世界中最普遍的视频形式——随意拍摄的单目视频（casual monocular video）——既没有已知相机参数，也不具备密集多视角线索。MoSca（2024.05）~\cite{2405_17421}（Dynamic Gaussian Fusion from Casual Videos via 4D Motion Scaffolds）正是面向这一最弱假设的全自动 4D 重建系统，无需任何相机先验即可从随意拍摄的单目视频恢复动态场景的完整几何与外观。MoSca 的核心洞察有两层：第一，虽然单目视频本身提供的约束极为有限，但预训练视觉基础模型（DPT 深度估计、CoTracker 长程 2D 跟踪、RAFT 光流）可以在大规模互联网数据上获得可靠的几何和对应关系先验，尽管这些先验是噪声的、局部的；第二，真实世界的运动通常是低秩且光滑的——少数``本征运动''（eigen motion）即可驱动整个场景的变形，因此可以用稀疏的图结构紧凑编码，避免在高维解空间中直接优化。

\textbf{运动脚手架表示。}MoSca 提出的核心表示是 4D 运动脚手架（Motion Scaffold，简称 $\mathcal{M}$）——一个稀疏轨迹图 $(V, E)$，它将密集的高维变形场压缩为少量节点的 6-DoF 轨迹。具体地，每个图节点 $v^{(m)}$ 定义为一条时空轨迹：$v^{(m)} = (\{Q_t^{(m)}\}_{t=1}^T, r^{(m)})$，其中 $Q_t^{(m)} = [R_t^{(m)}, t_t^{(m)}] \in SE(3)$ 是逐时间步的刚体变换（旋转 + 平移），$r^{(m)} \in \mathbb{R}^+$ 是 RBF 控制半径，决定该节点对周围空间的影响范围。节点数 $M$ 远小于场景所需的点数，实现了变形场的低秩压缩。图边 $E$ 基于曲线距离（curve distance）$D_\text{curve}(m,n) = \sum_t \|t_t^{(m)} - t_t^{(n)}\|$ 的 K 近邻构建——该度量在所有时间步上累积轨迹间的平移差，捕捉全局邻近性，同时能处理拓扑变化（例如开门时门与墙在空间中相邻但轨迹分离，曲线距离正确反映了这种``不连接''关系）。从图节点到密集变形场的插值采用 Dual Quaternion Blending（DQB）——与线性混合蒙皮（LBS）不同，DQB 在 $SE(3)$ 流形上插值，保证结果始终是合法的刚体变换。给定查询点 $\mathbf{x}$ 在源时刻 $t_\text{src}$ 的最近节点 $v^{(m^*)}$（$m^* = \arg\min_m \|t_{t_\text{src}}^{(m)} - \mathbf{x}\|$），变形到目标时刻 $t_\text{dst}$ 定义为：

\begin{equation}
\mathcal{W}(\mathbf{x}, \mathbf{w}; t_\text{src}, t_\text{dst}) = \text{DQB}(\{w_i, Q^{(i)}\}_{i \in E(m^*)})
\end{equation}

其中 $Q^{(i)} = Q^{(i)}_{t_\text{dst}} \cdot (Q^{(i)}_{t_\text{src}})^{-1}$ 是从源时刻到目标时刻的相对变换，蒙皮权重 $w_i = \exp(-\|\mathbf{x} - t_{t_\text{src}}^{(i)}\|_2^2 / 2r^{(i)})$ 由 RBF 距离衰减定义——离节点平移越近的点受该节点影响越大。

\textbf{从基础模型提取先验。}4D 重建从单目视频出发是高度不适定的问题，MoSca 的第一步是利用预训练视觉基础模型获取噪声但宝贵的几何与对应关系先验。具体提取三类信息：（1）逐帧深度估计 $\mathcal{D} = [D_1, D_2, \ldots, D_T]$（来自 DPT/Metric3D/DepthCrafter），提供相对一致的帧间深度；（2）长程 2D 像素轨迹 $\mathcal{T} = \{\tau^{(i)}\}$（来自 CoTracker/SpatialTracker），每条轨迹 $\tau^{(i)} = [(p^{(i)}_1, v^{(i)}_1), \ldots, (p^{(i)}_T, v^{(i)}_T)]$ 记录像素 $p^{(i)}_t$ 在各帧的 2D 位置和可见性 $v^{(i)}_t$；（3）基于 RAFT 光流计算的逐帧对极误差图 $\mathcal{E} = [E_1, \ldots, E_T]$，指示动态前景的似然——对极误差大的轨迹更可能属于运动前景。这些原始先验虽然不完整、有噪声且仅覆盖局部，但为后续的全局融合优化提供了关键起点。

\textbf{相机初始化。}为实现``野外''（in-the-wild）4D 重建，MoSca 必须在相机参数未知的情况下工作。第二步是基于轨迹的 bundle adjustment，鲁棒地初始化相机焦距和位姿。首先，对每条轨迹计算最大对极误差 $e(\tau) = \sum_{t=1}^T E_t[p_t] v_t$，以预设阈值筛出置信度高的静态背景轨迹。以预设初始焦距为起点，通过最小化静态轨迹上的重投影误差联合优化相机位姿和内参：

\begin{equation}
L_\text{proj} = \sum_i \sum_{a,b} \mathbb{1}(v^{(i)}_a \wedge v^{(i)}_b) \|\pi_K(W_b^{-1} W_a \pi_K^{-1}(p^{(i)}_a, D_a[p^{(i)}_a])) - p^{(i)}_b\|
\end{equation}

其中 $\pi_K$ 表示用内参 $K$ 投影，$W_t$ 是时刻 $t$ 的相机位姿。为补偿深度估计中的尺度误差（特别是帧间尺度不一致），同时优化深度修正因子（含逐帧全局缩放和逐像素小修正），深度对齐损失为：

\begin{equation}
L_z = \sum_i \sum_{a,b} \mathbb{1}(v^{(i)}_a \wedge v^{(i)}_b) D_\text{-inv}([\cdot]_z, D_b[p^{(i)}_b])
\end{equation}

其中 $[\cdot]_z$ 取 $z$ 坐标，$D_\text{-inv}(x,y) = |x/y - 1| + |y/x - 1|$ 是对称逆距离。总体 BA 损失 $L_\text{BA} = \lambda_p L_\text{proj} + \lambda_z L_z$。虽然相机求解并非 MoSca 的主要贡献，但该步骤使系统在动态视频上达到了最先进的相机位姿精度。

\textbf{运动脚手架的几何优化。}第三步是 $\mathcal{M}$ 的几何构建——将 2D 基础模型先验无缝集成到 4D 图结构中。首先，通过对极误差阈值识别前景轨迹，利用深度估计在可见帧反投影、在遮挡帧线性插值，将 2D 轨迹提升为 3D：

\begin{equation}
h_t = \begin{cases} W_t^{-1} \pi_K^{-1}(p_t, D_t[p_t]), & v_t = 1 \\ \text{interp}(h_\text{left}, h_\text{right}), & v_t = 0 \end{cases}
\end{equation}

每条提升后的 3D 轨迹初始化一个节点 $v^{(i)}$：平移部分取 $h_t^{(i)}$，旋转部分初始化为单位矩阵（即 $Q_t^{(i)} = [I, h_t^{(i)}$]，控制半径取预设值 $r_\text{init}$。实践中通过曲线距离均匀下采样，保留一个子集作为节点。初始化后，可见帧的 3D 位置保持固定，需要优化的``缺口''是：不可见帧的旋转和位置。MoSca 通过图拓扑 $E$ 上传播可见信息来填补这些缺口，优化物理启发的正则化目标：

\begin{equation}
L_\text{geo} = \lambda_\text{arap} L_\text{arap} + \lambda_\text{acc} L_\text{acc} + \lambda_\text{vel} L_\text{vel}
\end{equation}

其中 ARAP 损失 $L_\text{arap}$ 在邻域上同时约束局部距离保持（$\|t_t^{(m)} - t_t^{(n)}\| \approx \|t_{t+\Delta}^{(m)} - t_{t+\Delta}^{(n)}\|$）和局部坐标系保持（$Q^{-1} t^{(m)}$ 在时间上一致）；速度正则 $L_\text{vel}$ 和加速度正则 $L_\text{acc}$ 约束节点运动的时间平滑性，同时惩罚平移速度和旋转速度（旋转部分用 $\|(\cdot)\|_F$ 取旋转对数的 Frobenius 范数）。几何优化仅更新旋转和不可见帧的平移，保持可见帧位置不变以防止退化。

\textbf{动态场景表示与光度优化。}运动脚手架 $\mathcal{M}$ 的全局变形场可以将任意时间步的点变换到任意其他时间步，这一性质使得 MoSca 可以将所有帧的观测融合为单一的连贯表示——这是区别于逐帧独立重建的关键优势。第四步从所有帧的前景深度反投影点初始化动态高斯集合：

\begin{equation}
\mathcal{G} = \{(\mu_j, R_j, s_j, o_j, c_j; t_j^\text{ref}, \mathbf{w}_j)\}_{j=1}^N
\end{equation}

前五个属性是标准 3DGS 参数（中心、旋转、缩放、不透明度、球谐颜色），后两个是 MoSca 特有的：$t_j^\text{ref}$ 是该高斯被初始化时的参考时间步，$\mathbf{w}_j \in \mathbb{R}^K$ 是逐高斯可学习的蒙皮权重修正。查询时刻 $t$ 时，来自所有参考时间步的高斯统一变形到 $t$ 后融合为完整场景：

\begin{equation}
\mathcal{G}(t) = \{(T_j(t)\mu_j, T_j(t)R_j, s_j, o_j, c_j) \mid T_j(t) = \mathcal{W}(\mu_j, w(\mu_j, t_j^\text{ref}) + \mathbf{w}_j; t_j^\text{ref}, t)\}_{j=1}^N
\end{equation}

其中 $w(\cdot)$ 是基础 RBF 蒙皮权重，$\mathbf{w}_j$ 为可学习修正，$\mathcal{W}$ 是运动脚手架定义的变形场。静态背景则用标准 3DGS $\mathcal{H}$ 单独建模（已知相机参数下直接反投影深度图初始化）。最终可渲染的动态场景为 $\mathcal{G}(t) \cup \mathcal{H}$。

光度优化的总目标包含多项损失：

\begin{equation}
L = \lambda_\text{rgb} L_\text{rgb} + \lambda_d L_\text{depth} + \lambda_\text{track} L_\text{track} + \lambda_\text{arap} L_\text{arap} + \lambda_\text{acc} L_\text{acc} + \lambda_\text{vel} L_\text{vel}
\end{equation}

其中 $L_\text{rgb}$ 和 $L_\text{depth}$ 是标准的光度和深度渲染损失。跟踪损失 $L_\text{track}$ 是 MoSca 独特的监督信号：将每个高斯的 3D 坐标（替代 RGB 颜色）光栅化为 XYZ 图，与基础模型推断的 2D 轨迹监督对齐，使变形场在长程对应关系上保持一致。几何正则项 $L_\text{arap}, L_\text{acc}, L_\text{vel}$ 继承自几何优化阶段，在光度优化中继续约束变形场的合理性。

\textbf{节点控制策略。}类似标准 3DGS 的梯度驱动致密化和剪枝，MoSca 为运动脚手架节点设计了专门的控制策略。

\textbf{致密化}：周期性地选取跟踪损失 $L_\text{track}$ 梯度最高的高斯，下采样后转化为新的图节点，增加运动表示的密度。

\textbf{剪枝}：周期性地将动态前景高斯随机拷贝到静态背景集合中，然后将前景高斯的不透明度重置为低值——这使冗余的前景高斯被简化。此外，还剪枝蒙皮权重对所有高斯均低于阈值的节点（表明其对变形建模贡献有限）。这种节点级别的控制使运动脚手架能自适应地调整稀疏性，在运动复杂区域增加节点密度，在运动简单区域保持紧凑。

\paragraph{UFO-4D：2026.02 同期工作}

UFO-4D（2026.02）~\cite{2602_24290} 实现从两张未定位图像的前馈 4D 重建。

\paragraph{No Pose No Problem：2026.05 同期工作}

No Pose No Problem（2026.05）~\cite{2605_22190} 进一步去除了 4D 前馈方法中的位姿先验需求。

\paragraph{ViewSplat：2026.03 同期工作}

ViewSplat（2026.03）~\cite{2603_25265} 提出视角自适应的前馈动态高斯合成。

\subsubsection{动态表面重建}

本路线介绍从动态4D高斯中提取时间一致的显式表面网格的方法。

前面讨论的方法都聚焦于新视角合成的渲染质量。然而，许多下游应用（如物理仿真、3D打印）需要显式的表面网格。标准 3DGS 的点云表示没有显式表面，将其扩展到动态场景后更是难以提取时间一致的表面。本节介绍从动态 4D 高斯中重建高质量时变表面的方法。

\paragraph{2D 高斯面元平面约束动态表面重建}

DynaSurfGS（2024.08）~\cite{2408_13972}（Dynamic Surface Reconstruction with Planar-based Gaussian Splatting）是本小节的代表性工作，它将两个已有方案——动态场景的 4D-GS 变形场和静态场景的 PGSR 无偏深度渲染——桥接到统一框架中，并新增 ARAP 时序正则化以应对动态特有挑战。

\textbf{两个 baseline 的互补局限。}回顾前文，4D-GS 的 HexPlane 变形场能将规范高斯映射到任意时间步，实现高质量动态渲染，但优化后的高斯因无序不规则分布而无法贴合场景表面——直接可视化会出现伪影和浮游体。另一方面，PGSR（见第四章表面重建小节）的无偏深度渲染在静态场景中表现出色，能精确提取表面网格，但完全缺乏时间建模能力。DynaSurfGS 的核心动机正是将两者结合：用 4D-GS 的变形场处理时间维度，用 PGSR 的平面化深度渲染处理几何维度。

\textbf{改动概要。}与 4D-GS 相比，DynaSurfGS 做了两个关键改变：(1) 将标准 3DGS 的 $\alpha$-混合 深度渲染替换为 PGSR 的无偏深度渲染方案，使变形后的高斯被约束贴合场景表面；(2) 新增 ARAP 正则化，约束相邻时间步之间的局部刚性。HexPlane 变形场的架构保持不变——仍将 4D 时空分解为六个 2D 平面（三个空间-空间平面 + 三个空间-时间平面），编码时空特征 $F = E(\mathcal{G})$ 后由紧凑 MLP 解码器预测逐高斯变形 $\delta\mu, \delta r, \delta s = D(F)$，变形后高斯属性为 $(\mu', r', s') = (\mu + \delta\mu, r + \delta r, s + \delta s)$。

\textbf{从 PGSR 迁移的无偏深度渲染。}DynaSurfGS 直接采用了 PGSR 的平面化深度渲染方案（详见第四章）。核心流程是：分别渲染法线图 $N = \sum_i T_i R_c^T n_i \alpha_i$ 和距离图 $L = \sum_i T_i l_i \alpha_i$（其中 $l_i = (R_c^T(\mu_i - T_c)) R_c^T n_i^T$ 是相机中心到高斯平面的距离），再通过 $L(m) = \frac{L}{N(m) \cdot K^{-1}\tilde{m}}$ 将两者转换为无偏深度图。这一除法消除了 $\alpha$-混合 权重总和不等于 1 带来的深度偏差，使深度值真正落在高斯平面上。

\textbf{法线正则化：表面光滑性约束。}在 PGSR 的单视角几何一致性基础上，DynaSurfGS 引入法线正则化 $L_\text{normal}$ 约束动态物体的表面光滑性。具体而言，从无偏深度图中采样四邻域点，投影到相机坐标系后拟合局部平面法线：

\begin{equation}
N_l(m) = \frac{(M_1-M_0) \times (M_3-M_2)}{|(M_1-M_0) \times (M_3-M_2)|},
\end{equation}

然后用图像梯度五次方加权的 $L_1$ 损失 $L_\text{normal} = \frac{1}{W}\sum_{m \in W} |\overline{\nabla I}|^5 \|N_l(m) - N(m)\|_1$ 约束其与渲染法线一致。梯度加权的作用是：纹理丰富区域（梯度大）对应几何细节区域，法线对齐约束更强；平坦区域（梯度小）则放松约束以避免过拟合噪声。对于边缘区域，局部平面假设可能失效，方法用图像边缘近似几何边缘以避免错误惩罚。

\textbf{ARAP 正则化：动态场景的新增约束。}法线正则化解决了单帧内的表面光滑性，但动态场景还面临跨帧几何一致性问题——相邻高斯在不同时间步之间的相对位置应保持近似刚性（除非发生非刚性形变）。这是 DynaSurfGS 相对于两个 baseline 的真正增量贡献。ARAP（As-Rigid-As-Possible）正则化 $L_\text{arap}$ 的实现分为三步。第一步，对每个高斯 $i$ 在规范空间找 $k=10$ 个最近邻，计算距离衰减权重：

\begin{equation}
w_{ij} = \frac{\tilde{w}_{ij}}{\sum_{j \in K}\tilde{w}_{ij}}, \quad \tilde{w}_{ij} = \exp\left(-\frac{d_{ij}^2}{2\sigma_j^2}\right),
\end{equation}

其中 $d_{ij}$ 是高斯 $i$ 与邻居 $j$ 的中心距离，$\sigma_j$ 是高斯 $j$ 的可学习缩放参数。第二步，随机采样两个时间步 $t_1, t_2$，用变形场得到各高斯在这两个时刻的空间位置，然后估计局部刚体旋转：

\begin{equation}
R_i = \arg\min_{R \in \text{SO}(3)} \sum_{j \in K} w_{ij} \|(\mu_i^{t_1} - \mu_j^{t_1}) - R(\mu_i^{t_2} - \mu_j^{t_2})\|^2.
\end{equation}

这一步通过 Procrustes 分析找到最优旋转，使 $t_2$ 时刻的局部邻域与 $t_1$ 时刻尽可能对齐。第三步，用估计的旋转计算 ARAP 损失：

\begin{equation}
L_\text{arap}(\mu_i, t_1, t_2) = \sum_{j \in K} w_{ij} \|(\mu_i^{t_1} - \mu_j^{t_1}) - R_i(\mu_i^{t_2} - \mu_j^{t_2})\|^2.
\end{equation}

直觉上，如果物体在 $t_1$ 到 $t_2$ 之间做刚体运动（旋转+平移），那么邻域内的相对位置关系经旋转 $R_i$ 修正后应完全不变；对于非刚性区域（如弯曲的手臂），损失值会反映变形的程度，软性约束变形场不要产生过大的非刚性形变。

\textbf{训练策略。}总训练损失为 $L = L_\text{photo} + L_\text{tv} + \lambda_1 L_\text{normal} + \lambda_2 L_\text{arap}$，其中 $L_\text{photo}$ 是标准 $L_1$ 光度损失，$L_\text{tv}$ 是 HexPlane 的全变分正则化。两个正则化系数采用渐进预热策略：$\lambda_1, \lambda_2$ 在第 7000--9000 迭代间从 0 线性渐增至目标值 0.05 和 0.02，之后保持固定。这种设计的原因是：训练初期变形场尚未收敛，高斯位置不稳定，过早施加强法线和刚性约束会干扰变形场学习；等变形场基本收敛后再逐步引入几何约束，使其精修表面质量而不破坏运动建模。网格提取采用 TSDF 融合算法，从各时间步的无偏深度图中融合时变三角网格。DynaSurfGS 在 D-NeRF、DG-Mesh 和 Ub4D 数据集上同时实现了高质量渲染（D-NeRF 上 PSNR 34.31 dB，超过 4D-GS 的 34.06 dB）和光滑动态表面重建（DG-Mesh 上 Chamfer Distance 0.910，EMD 0.123，显著优于 D-NeRF 的 1.252/0.184 和 DG-Mesh 的 0.770/0.128），验证了"4D 变形 + 无偏深度 + ARAP"三组件协同的有效性。

\paragraph{显式高斯与隐式函数的桥接}

DySurface（2026.05）~\cite{2605_10360}（Consistent 4D Surface Reconstruction via Bridging Explicit Gaussians and Implicit Functions）指出从预训练的可变形高斯泼溅模型中提取网格并不简单——缺乏显式表面约束导致几何不一致，产生噪声和模糊的网格。前述 DynaSurfGS 通过无偏深度渲染与法线/ARAP 正则化实现了基于平面的表面重建，但其表面质量仍受限于高斯基元本身的离散分布，无法直接获得连续水密表面。DG-Mesh 等方法通过 Poisson 重建等 ad-hoc 算法从离散高斯中提取网格，天然缺乏连续表面约束，在严重非刚性变形下产生破碎不准确的几何。DySurface 的设计动机正是针对这些局限：将离散高斯基元、连续隐式 SDF 和显式网格三种互补表示紧密耦合到统一管线中，同时解决一个被忽视的结构性矛盾——3DGS 的前向映射（canonical $\rightarrow$ dynamic）与体渲染 SDF 所需的后向映射（dynamic $\rightarrow$ canonical）方向相反，直接嫁接 SDF 约束到动态 3DGS 上并不平凡。

方法分三阶段递进优化。第一阶段是高斯泼溅（GS）分支，直接沿用 GaGS 的前向变形框架：在规范空间定义 3D 高斯基元 $\mathcal{X}_c$，通过 3D 稀疏卷积编码器 $\mathcal{H}_\phi$ 将其离散化为体素网格提取全局特征体积，对每个坐标 $x$ 查询全局特征 $\mathbf{f}_\text{global}(x)$ 并拼接逐点局部特征 $\mathbf{f}_\text{local}(x)$，由 MLP 解码器 $\Phi_\psi$ 预测时间相关残差，形成前向变换场 $\mathcal{T}_\theta$（$\theta = \{\phi, \psi\}$）：$\mathcal{T}_\theta(x,t) = \Phi_\psi(\text{Concat}(\mathbf{f}_\text{global}(x), \mathbf{f}_\text{local}(x)), t)$，输出 $(\Delta\mu, \Delta s, \Delta r)$ 叠加到规范属性上得到变形后高斯 $\mathcal{X}_d$。GS 分支的训练损失为 $L_\text{GS} = L_\text{rgb} + \lambda_m L_\text{mask} + \lambda_\text{reg}^\text{gs} L_\text{reg}^\text{gs}$，其中 $L_\text{reg}^\text{gs} = \|\Delta\mu\|_1$ 鼓励最小空间变形。

第二阶段是 DySurface 的核心创新——VoxGS-DSDF（体素化高斯驱动动态 SDF）分支。与 DynaSurfGS 仅在高斯层面添加表面正则不同，DySurface 构建了一个完整的隐式 SDF 场，并通过显式高斯体素对其进行锚定。具体而言，将变形后高斯 $\mathcal{X}_d(t)$ 构建动态稀疏体素网格 $\nu(\mathcal{X}_d)$，对沿射线采样的动态空间查询点 $q_d$，先用 KNN 搜索 $\mathcal{N}_K(q_d; \mathcal{X}_d)$ 找到附近 $K$ 个高斯体素，经距离阈值 $\tau$ 过滤后得到有效邻居集 $\mathcal{N}'(q_d; \nu(\mathcal{X}_d)) = \{k \in \mathcal{N}_K \mid \|v_k - x_d\|_2^2 < \tau\}$，从中提取几何特征 $\mathbf{f}_\text{geo,k}$ 和外观特征 $\mathbf{f}_\text{app,k}$，取均值后与位置编码 $\gamma(q_d)$ 拼接形成空间条件向量：$\eta_d = \text{Concat}(\gamma(q_d), \frac{1}{K'}\sum_k \mathbf{f}_\text{geo,k}, \frac{1}{K'}\sum_k \mathbf{f}_\text{app,k})$。这一设计的核心意义在于：将 GS 分支已学习的显式几何先验注入隐式 SDF 优化过程，使 SDF 场不再是在真空中优化，而是被离散高斯基元显式引导。条件化于时间嵌入 $t$，MLP $\Phi_\omega$ 将 $\eta_d$ 映射为后向变形 $\Delta q_\text{bwd} = \Phi_\omega(\eta_d, t)$，将动态查询点映射回答范空间 $q_c = q_d + \Delta q_\text{bwd}$，几何网络 $G$ 预测 SDF 值 $s_i = G(q_{c,i})$，颜色网络 $C$ 预测视角相关辐射 $c_i$，SDF 值通过可学习逻辑分布 $\mathcal{G}(x) = (1 + e^{\beta x})^{-1}$ 转换为离散 alpha 值 $\alpha_i = \max(\frac{\mathcal{G}(s_i) - \mathcal{G}(s_{i+1})}{\mathcal{G}(s_i)}, 0)$，参数 $\beta$ 控制表面锐度，最终通过标准体积分积 $\hat{C}(\mathbf{r}) = \sum_{i=1}^N T_i \alpha_i c_i$ 合成像素颜色。

VoxGS-DSDF 分支的关键损失设计体现了显式-隐式桥接的思想。SDF-高斯锚定损失 $L_\text{SDF-GS} = \text{BCE}(|s(q_c)| > \epsilon,\, d_\text{NN}(q_c; \mathcal{X}_c) > \epsilon)$ 用二元交叉熵引导连续零水平集与离散基元对齐——在规范空间中，SDF 值接近零（靠近表面）的点应同时靠近高斯基元中心，反之亦然。消融实验表明移除此损失后 vIoU 从 0.3928 骤降至 0.3388，表面出现严重碎裂和浮游伪影，证实显式锚定对水密边界的必要性。循环一致性损失 $L_\text{cycle} = \|\Delta\overleftarrow{q}_\text{fwd} + \overrightarrow{q}_\text{bwd}\|_1$ 确保前向变形（GS 分支，canonical $\rightarrow$ dynamic）与后向变形（SDF 分支，dynamic $\rightarrow$ canonical）构成双射，消除两方向映射之间的矛盾。消融显示加入此损失使 vIoU 从 0.3519 提升至 0.3928，PSNR 从 27.50 提升至 31.09。此外还包含：从 GS 分支蒸馏深度和法线图的蒸馏损失 $L_\text{distill}$，将显式高斯的几何知识迁移到隐式场；Eikonal 正则 $L_\text{eik} = (\|\nabla s(q_c)\|_2 - 1)^2$ 约束 SDF 梯度范数为 1；以及时序平滑正则 $L_\text{smooth}$。SDF 总损失为 $L_\text{SDF} = L_\text{rgb}^\text{sdf} + \lambda_c L_\text{cycle} + \lambda_\text{SG} L_\text{SDF-GS} + \lambda_\text{dist} L_\text{distill} + \lambda_e L_\text{eik} + \lambda_s L_\text{smooth}$。消融还表明，移除整个 VoxGS-DSDF 锚定机制（仅用位置编码替代高斯体素特征）虽几何指标下降最少（vIoU 0.3726），但 PSNR 从 31.09 降至 27.22，说明显式锚定对渲染保真度的贡献甚至大于对几何精度的贡献。

第三阶段是动态网格精修（MeshGS 分支）：从收敛的 SDF 零水平集通过 Marching Cubes 提取规范网格 $\mathcal{M}_c$，直接复用 GS 分支已学习的前向变换场 $\mathcal{T}_\theta$ 初始化精修变换场 $\mathcal{T}'_\theta$，驱动网格顶点变形 $v_d(t) = v_c + [\mathcal{T}'_\theta(v_c, t)]_\mu$。这一设计的精妙之处在于：3D 高斯作为显式离散基元，擅长在初始优化阶段追踪快速复杂的非刚性动态，而精修阶段将这一已学习的显式运动先验无缝迁移到网格表面，桥接了体渲染和显式表面渲染之间的域差异。精修损失 $L_\text{Mesh} = L_\text{rgb}^\text{mesh} + \lambda_m^\text{mesh} L_\text{mask}^\text{mesh} + \lambda_\text{lap} L_\text{lap} + \lambda_\text{reg}^\text{mesh} L_\text{reg}^\text{mesh}$，其中 Laplacian 平滑仅应用于最低 25\% 曲率误差的顶点以保持稳定几何，$L_\text{reg}^\text{mesh} = \|v_d - v_c\|_1$ 正则化最小变形。

在 D-NeRF 数据集上，DySurface 的平均 vIoU（0.3928）和 Chamfer Distance（0.0102）均显著优于 DG-Mesh（0.3814/0.0137）和 D-2DGS（0.3426/0.0111），同时 PSNR（31.09）接近纯渲染方法 GaGS（32.24），证明了连续表面约束与渲染质量可以同时实现。在 Nerfies 真实世界数据集上，PSNR（27.91）和 LPIPS（0.0815）均优于 DG-Mesh（27.23/0.0927），在复杂背景和高度纹理化物体上捕获了更稳定的高频网格细节。

\paragraph{时空 2D 高斯与运动正则化}

在动态场景表面重建中，现有 4D 高斯方法（如 Deformable-GS、4D-Gaussians）聚焦于新视角合成而非几何重建，而静态场景中的 2D Gaussian Splatting（2DGS）虽能将高斯圆盘贴合物体表面实现高质量表面提取，却未设计动态支持。Space-time 2D Gaussian Splatting（2024.09）~\cite{2409_18852} 的核心动机正是将 2DGS 的表面重建优势迁移到动态域。方法学习一组规范 2D 高斯（canonical 2DGS），对每个时间戳 $t$，一个变形网络将规范高斯调整到该时刻的状态，使变形后的 2D 高斯圆盘紧密贴合物体表面。具体而言，规范 2D 高斯经变形网络 $F$ 映射后得到时变高斯，其位置偏移 $\Delta\mu_t$ 和协方差变形 $\Delta\Sigma_t$ 由网络预测：$\mu_t^{Ob} = \mu^c + \Delta\mu_t$，$\Sigma_t^{Ob} = \Sigma^c + \Delta\Sigma_t$——即变形网络输出平移偏移和协方差变形，叠加到规范参数上。为确保变形后表面质量，方法施加深度一致性和法线一致性正则化，在时空域联合约束表面。此外，方法引入时变不透明度模型（time-varying opacity），对每个时间步学习独立的不透明度参数而非共享静态值，使被动态遮挡区域的高斯能自动降低不透明度以减少几何噪声。为去除纹理缺失的背景区域，方法还设计了一个将渲染 $\alpha$ 遮罩与前景遮罩对齐的正则化项，加速学习收敛。最终表面通过 TSDF 融合从变形后的 2D 高斯点云中提取三角网格，在 D-NeRF 和 CMU Panoptic 等含多人活动和人物交互的复杂动态数据集上，几何精度和训练效率均超越已有方法。

\paragraph{动态高斯到网格重建}

前述方法聚焦于动态场景的新视角合成渲染质量，而 DG-Mesh（2024.04）~\cite{2404_12379}（Dynamic Gaussians Mesh）向前迈出关键一步：从动态观测中提取高质量、时间一致且具有跨帧顶点对应的显式网格——这是首个从动态观测中恢复具有跨帧对应关系的高保真网格的框架。现代 3D 引擎和图形管线需要网格作为内存高效的表示，但动态场景的网格提取面临细节结构和时间一致性的双重挑战。

\textbf{变形基础：规范空间范式。}DG-Mesh 的变形基础与 Deformable-GS 完全相同。它维护一组规范 3D 高斯，并通过变形网络 $\mathcal{F}$ 将它们变换到不同时间步进行渲染。变形网络以经过位置编码的高斯参数与时间标签作为输入，输出的变形量涵盖四个分量：位置平移 $\delta x$、朝向旋转 $\delta r$、尺寸缩放 $\delta s$ 以及透明度调整 $\delta\alpha$。

其中位置编码采用多频率的 Fourier 特征映射——对每个输入坐标 $p$ 计算 $\sin$ 和 $\cos$ 在 $2^0, 2^1, \ldots, 2^k$ 等 $k+1$ 个频率上的响应——用以增强对高频细节的重建能力。经过变形之后，高斯的新表示为 $G'(x+\delta x;\, \mu,\, r+\delta r,\, s+\delta s,\, \alpha+\delta\alpha)$。

然而，直接训练得到的 3D 高斯会在空间中出现不均匀分布——有些区域密集堆叠、有些区域则相对稀疏——这带来了两个关键问题：（1）违反了 Poisson 重建假设的均匀采样前提，导致重建的网格质量明显下降；（2）长时间跨度的运动追踪变得困难，因为不均匀分布的高斯缺乏稳定的空间对应结构。

\textbf{表面提取：从高斯到网格的可微管线。}DG-Mesh 将变形后的高斯视为有向点云（oriented point cloud），通过两步可微操作提取网格。首先，用可微 Poisson 求解器（DPSR）将高斯法向量场转化为空间中的指示函数 $\chi(x)$，该函数描述几何的内外关系。然后，用可微 Marching Cubes 从指示函数中提取水密流形网格。为保持网格表面细分质量，引入 Laplacian 正则化 $L_{lap} = \frac{1}{n}\sum_{i=0}^{n}\|\delta_i\|^2$，其中 $\delta_i = v_i - \frac{1}{|N_i|}\sum_{k \in N_i} v_k$ 是顶点 $v_i$ 与其一环邻域 $N_i$ 平均位置的偏差，惩罚局部曲率突变。

\textbf{Gaussian-Mesh Anchoring：网格引导的高斯均匀化。}DG-Mesh 的核心创新是 Gaussian-Mesh Anchoring（高斯-网格锚定）机制——用已提取网格的均匀面分布反过来指导高斯的密度控制，将不均匀的高斯重新调整为与网格面一一对应的均匀分布。具体操作如下：给定变形后的高斯集合 $\{G(x_i)\}_{i=0}^{N}$ 和重建网格的面质心集合 $\{f_j\}_{j=0}^{M}$（$N \neq M$），对每个高斯 $x_i$ 找其最近邻网格面 $n_{x_i} = \arg\min_f \|x_i - f\|$。对于被多个（$K$ 个）高斯最近邻的网格面 $f_j$，剪除原始的 $K$ 个高斯，取其属性平均值创建一个新高斯：$G' = \frac{1}{K}\sum_{k=0}^{K} G(x_k)$；对于没有被任何高斯覆盖的网格面 $f_j$，在该面质心处直接创建新高斯：$G' = G(f_j)$。锚定后的高斯与网格面形成一一对应，分布均匀。此外，锚定损失 $L_{anchor} = \frac{1}{n}\sum_{i=0}^{n}\|x_i - n_{x_i}\|^2$ 约束每个高斯到其最近网格面质心的距离，维持对应关系的紧密性。

\textbf{跨帧对应：两级分解与循环一致变形。}为建立跨帧网格的顶点对应关系，DG-Mesh 将网格到网格的对应（mesh-to-mesh correspondence）分解为两级：（1）帧内的网格到高斯点对应（mesh-to-point correspondence），由锚定的一一映射保证；（2）跨帧的高斯点到规范高斯点的对应（point-to-canonical-point correspondence），由变形网络保证。两级分解的关键洞察是：只需分别解决两个简单的对应问题，即可间接获得复杂的网格间对应。

锚定操作在每帧的变形空间中修改了高斯（合并、新增），因此需要将修改后的锚定高斯投影回答范空间以保持一致性。DG-Mesh 引入循环一致变形（Cycle-Consistent Deformation）：定义前向变形 $\mathcal{F}_f$ 和后向变形 $\mathcal{F}_b$，给定规范高斯 $G$ 及其在变形空间中的锚定形式 $G'$，约束 $\mathcal{F}_f(G, t) = -\mathcal{F}_b(G', t)$，用 $L_1$ 损失作为循环一致损失 $L_{cycle}$。这条循环路径是：规范高斯 $\xrightarrow{\mathcal{F}_f}$ 变形高斯 $\xrightarrow{\text{锚定}}$ 锚定高斯 $\xrightarrow{\mathcal{F}_b}$ 规范高斯。循环约束确保锚定操作引起的修改能正确传播回答范空间，规范空间中的高斯可以自动适应新增或移除的高斯。

\textbf{双路渲染监督与训练策略。}DG-Mesh 采用双路渲染监督：3DGS 高斯泼溅产生高斯图像 $I_{gs}$，同时通过可微光栅化器（Nvdiffrast）渲染提取的网格产生网格图像 $I_{mesh}$，两者均与真实图像 $I_{gt}$ 计算渲染损失。其中网格的顶点颜色并非在变形空间直接优化，而是将变形后的顶点位置投影回答范空间，查询时间相关的外观模块获取颜色——实验表明在规范空间优化顶点颜色比在各变形空间独立优化效果更好。此外，还加入光栅化掩码损失 $L_{mask}$ 约束前景区域。

总训练损失为 $L = L_{gs} + L_{mesh} + L_{mask} + L_{lap} + L_{anchor} + L_{cycle}$，端到端优化高斯参数、网格几何和跨帧对应。

训练采用精心设计的分阶段策略：先冻结变形网络，仅训练规范高斯 3k 迭代以稳定其位置和形状；5k 迭代后引入 DPSR 和可微 Marching Cubes 提取网格几何；Gaussian-Mesh Anchoring 每 100 次迭代执行一次。全部训练在单张 RTX 3090Ti 上进行 50k 迭代。

\textbf{实验表现。}在 D-NeRF 合成数据集、DG-Mesh 自建数据集（含 ground truth 网格的六段动态场景）以及 Nerfies、Dycheck、Unbiased4D、NeuralActor 等真实数据集上，DG-Mesh 的 Chamfer Distance（CD）和 Earth Mover Distance（EMD）均显著优于所有基线方法（D-NeRF、K-Plane、HexPlane、TiNeuVox、4D-GS、Deformable-GS、4DGS、SC-GS），尤其在鸟翼等薄结构的重建上，其他方法产生浮游体或过大的网格，DG-Mesh 则能准确恢复精细几何。DG-Mesh 的网格可处理拓扑变化（如球体到环面的变换，亏格从 0 变为 1），并可直接导入 3D 引擎进行纹理编辑和物理仿真。

\paragraph{自适应时序一致高斯面元}

Adaptive and Temporally Consistent Gaussian Surfels（AT-GS）（2024.11）~\cite{2411_06602} 解决动态场景表面重建中拓扑变化（物体出现/消失）的难题。已有方法（如 Vidu4D、DG-Mesh）通过全局优化整个序列来重建动态表面，将所有帧统一在单一规范空间中，但面对显著拓扑变化（新物体出现或旧物体消失）或快速运动的长序列时表现不佳。AT-GS 提出 per-frame 增量优化策略：不维护全局规范空间，而是逐帧处理动态场景。首帧使用标准静态 Gaussian Surfels 从 SfM 稀疏点云进行完整训练；对后续每帧 $t$，采用 coarse-to-fine 两阶段策略。粗阶段仅更新高斯中心位置和旋转：训练一个逐帧 Neural Transformation Cache（NTC），由多分辨率哈希编码和浅层 MLP 组成，将空间位置映射到 SE(3) 变换。NTC 利用体素网格的空间平滑性和线性插值实现快速收敛。

细阶段优化全部可学习参数（中心、旋转、缩放、视角相关颜色、不透明度），同时允许剪枝和自适应梯度引导的密度控制。密度控制的关键改进是将传统的 clone 和 split 统一为单一步骤：原始高斯沿位置梯度方向移动，同时在梯度引导的位置采样生成新高斯，采样标准差由梯度幅度自适应调整：$\sigma = \frac{\min(\|\nabla L\|, 2\overline{\|\nabla L\|})}{2\overline{\|\nabla L\|}} \cdot \overline{\|\nabla L\|}$，使大梯度高斯移动更远、小梯度高斯做精细校正，保持原始大小不缩分裂变以加速收敛。为解决帧间时序抖动问题，AT-GS 引入基于曲率的时序一致性约束：用 RAFT 估计相邻帧光流，将前一帧渲染法线图 warp 到当前帧，计算曲率图 $C = \|(\partial_x N + \partial_y N)\|_2$，然后最小化相邻帧曲率图的均方误差 $L_t = \text{MSE}(\hat{C}_t, C_t)$，间接确保深度图和高斯朝向在时间上的一致性。每帧粗阶段训练 200 次迭代、细阶段 800 次迭代，单帧训练约 30 秒，最终通过 screened Poisson 重构提取表面网格。

此外，4DSurf（2026.03）~\cite{2603_28064} 实现高保真动态场景表面重建，DyGASR（2024.11）~\cite{2411_09156} 通过广义指数泼溅和表面对齐提升动态几何精度。

\subsection{运动建模与扩展应用}

本节介绍三条技术路线：（1）运动建模与操纵；（2）分析与概率方法；（3）其他变形与4D方法。

\subsubsection{运动建模与操纵}

本路线介绍对重建后的运动进行显式建模、编辑和交互式操纵的方法，使4D场景变为可交互资产。

前面所有方法都聚焦于"重建"——从观测数据恢复动态场景。本节转向一个更主动的目标：不仅重建运动，还要显式地建模、编辑和操纵运动。这些方法让重建后的 4D 场景变成可交互的资产，而非只能回放的录像。

\paragraph{Wasserstein 距离约束 4D 高斯传输}

前述 AT-GS 通过逐帧增量优化和曲率时序约束处理动态表面重建，但其帧间一致性仍然依赖欧氏光流匹配，本质上只约束了高斯中心位置和法线方向，无法感知协方差矩阵（形状、朝向、尺度）的帧间突变。现有 4DGS 方法（如 Deformable-GS、4D-Gaussians）依赖神经变形场预测高斯时变参数，但帧间参数突变导致闪烁伪影，且欧氏距离正则化仅约束均值位置或协方差矩阵的 Frobenius 范数，无法捕捉高斯分布作为整体的位置-形变联合变化。Gaussians on their Way（2024.12）~\cite{2412_00333} 从一个完全不同的角度切入：不从重建流程改进帧间传递，而是从最优传输（Optimal Transport）理论出发，用 Wasserstein 距离替代欧氏度量，将每个高斯分布视为一个完整的概率实体而非零散参数的拼凑，从而在度量层面根本解决帧间参数突变的闪烁问题。方法框架受控制论中的 Kalman 滤波启发，包含三个协同组件：观察器（Observer）、预测器（Predictor）和融合器（Merger）。

\textbf{观察器：神经高斯变形场。}观察器是一个 MLP 变形网络 $f_\theta$，输入规范高斯参数 $N^c = N(\mu^c, \Sigma^c)$ 和时间 $t$（经位置编码），输出观测状态 $N_t^{Ob}(\mu_t^{Ob}, \Sigma_t^{Ob})$：

\begin{equation}
\mu_t^{Ob} = \mu^c + \Delta\mu_t, \quad \Sigma_t^{Ob} = \Sigma^c + \Delta\Sigma_t,
\end{equation}

其中 $\Delta\mu_t$ 为平移偏移，$\Delta\Sigma_t$ 为协方差变形增量。这一设计与 Deformable-GS 的变形场结构类似，但关键区别在于观察器的输出不直接用于渲染，而是作为"观测数据"送入后续的状态一致性滤波器——观测本身不可避免地带有离散时间采样引入的噪声和跳变。

\textbf{Wasserstein 距离：高斯分布的整体度量。}现有方法（如 SC-GS、Deformable-GS）对帧间变化的正则化使用欧氏度量：均值用 $\ell_2$ 距离，协方差用 Frobenius 范数。这存在根本缺陷——协方差矩阵是 SPD(3) 流形上的元素，而非 $\mathbb{R}^{3\times3}$ 中的普通向量，Frobenius 范数忽略了 SPD 流形的几何结构，也无法同时感知位置变化和形变变化之间的耦合关系。Gaussians on their Way 引入 2-Wasserstein 距离作为高斯分布间的整体度量：

\begin{equation}
W_2^2(N_1, N_2) = \|\mu_1 - \mu_2\|^2 + \text{tr}\!\left(\Sigma_1 + \Sigma_2 - 2\left(\Sigma_1^{1/2}\,\Sigma_2\,\Sigma_1^{1/2}\right)^{\!1/2}\right),
\end{equation}

其中第一项量化均值位移，第二项中的迹运算在 SPD(3) 流形上度量协方差差异。利用 3DGS 的旋转-缩放分解 $\Sigma = RSS^TR^T$，迹项可进一步化简为 $D_\Sigma = \text{tr}(S_1 + S_2 - 2(S_1^{1/2}E_{12}S_1^{1/2})^{1/2})$，其中 $E_{12} = R_1^TR_2 S_2 R_2^T R_1$ 为旋转对齐后的交叉协方差，避免了直接计算 $3\times3$ 矩阵平方根的数值不稳定性。

\textbf{预测器：Wasserstein 几何上的高斯动力学。}在欧氏度量下，高斯速度被简单地定义为 $v_t^{\text{Eucl}} = (\mu_t - \mu_{t-1},\, \Sigma_t - \Sigma_{t-1})$，将均值和协方差视为独立分量分别做差。Wasserstein 几何则通过流形上的对数/指数映射定义内在速度，正确处理高斯分布作为完整实体的非线性演化。速度通过对数映射计算：

\begin{equation}
v_t = -\log_{N_t} N_{t-1}, \quad
v_\mu = \mu_t - \mu_{t-1}, \quad
\Omega_{\Sigma_t \Sigma_{t-1}} = 2\Sigma_{t-1} - P - P^T,
\end{equation}

其中 $P = \Sigma_{t-1}^{1/2}(\Sigma_{t-1}^{1/2}\,\Sigma_t\,\Sigma_{t-1}^{1/2})^{1/2}\,\Sigma_{t-1}^{-1/2}$。均值速度退化为欧氏差，但协方差速度 $\Omega_{\Sigma_t\Sigma_{t-1}}$ 捕捉了 SPD 流形上协方差矩阵之间的最优传输方向。下一帧预测状态通过指数映射 $N_{t+1}^P = \exp_{N_t}(v_t)$ 获得：均值预测 $\mu_{t+1}^P = \mu_t + v_\mu$ 为标准欧氏外推；协方差预测 $\Sigma_{t+1}^P = \Sigma_t + \Omega_t + \Phi_t\,\Sigma_t\,\Phi_t^T$ 需解 Sylvester 方程 $\Phi_t\Sigma_t + \Sigma_t\Phi_t = \Omega_t$，确保预测结果仍在 SPD 流形上。整个对数-指数映射在假设零加速度（匀速运动）的简化动力学模型下运作——虽然远非真实物理动力学，但提供了比欧氏线性外推更合理的几何先验。

\textbf{融合器：类 Kalman 增益的状态更新。}融合器借鉴 Kalman 滤波的核心思想，将预测状态 $N_t^P$（来自 Wasserstein 动力学）与观测状态 $N_t^{Ob}$（来自变形网络）加权合并：

\begin{equation}
K = \Sigma_t^{Ob}\left(\Sigma_t^{Ob} + \Sigma_t^P\right)^{-1}, \quad \hat{N}_t = N_t^{Ob} + K\left(N_t^P - N_t^{Ob}\right),
\end{equation}

其中 $K$ 为类 Kalman 增益矩阵。与传统 Kalman 滤波不同，这里不引入独立的过程噪声和测量噪声分布，而是将 Kalman 增益纯粹用作预测与观测之间的平衡权重——协方差较大的路径获得较低权重。最终渲染使用融合后的高斯 $\hat{N}_t$，而非直接使用变形网络输出。

\textbf{双重 Wasserstein 损失与训练策略。}方法引入两个互补的 Wasserstein 损失。状态-观测对齐损失 $L_{SOA} = W_2^2(N_t^P, N_t^{Ob})$ 鼓励 Wasserstein 预测与神经变形场观测在物理上保持一致——虽然观测本身有离散采样误差，但预测融合了运动学先验，二者的 Wasserstein 距离衡量了"物理合理性"与"数据拟合度"之间的偏差。Wasserstein 正则化 $L_{WR} = \sum_t \sum_{i=1}^N W_2^2(\hat{N}_t^{(i)}, \hat{N}_{t-1}^{(i)})$ 直接惩罚帧间参数突变，针对闪烁伪影的核心成因。总损失为
\begin{equation}
L = L_{\text{render}} + \lambda_{SOA}\,L_{SOA} + \lambda_{WR}\,L_{WR},
\end{equation}

其中 $L_{\text{render}}$ 结合 $L_1$ 和 D-SSIM。训练采用分阶段策略：前 3k 迭代仅优化静态 3D 高斯以稳定初始化；6k 迭代后引入滤波器模块；20k 迭代后激活 SOA 和 Wasserstein 正则化损失（$\lambda_{SOA} = 0.1$，$\lambda_{WR} = 0.01$）；同时延迟状态一致性更新直到高斯数量稳定（3k 迭代后），且仅对预测-观测中心偏差 $\|\mu_{pred} - \mu_{obs}\| < 0.1$ 的高斯启用滤波。高斯数量变化（clone/split）期间，偏差过大的高斯（$>3$）回退到标准 3DGS 优化以避免拓扑冲突。

\textbf{实验表现。}在 D-NeRF 合成数据集（$800\times800$，白背景）上，方法达到 34.45 dB PSNR、0.970 SSIM，保持 45.5 FPS，显著优于同期方法。在快速运动场景上表现尤为突出：Bouncing Balls 场景达到 42.79 dB（超过 Def-3D-Gauss 的 41.01 dB），Jumping Jacks 场景达到 37.91 dB；Lego 场景达到 34.74 dB，超过 4D-Rotor-Gauss 约 9.5 dB。在 Plenoptic Video 真实数据集（$1352\times1014$）上达到 32.79 dB PSNR、0.945 SSIM，37 FPS，在非刚性变形场景上展现出较强的鲁棒性。全量训练在单张 NVIDIA A800 上进行 150k 迭代。Wasserstein 几何确保高斯的运动轨迹和形变在最优传输意义下最自然平滑，避免了欧氏度量下的闪烁和跳跃。

\paragraph{运动-几何解耦与关节变形}

前面的 Gaussians on their Way 用 Wasserstein 最优传输度量约束高斯帧间演化，确保运动轨迹和形变在整体分布意义下平滑。但它沿用的 4D 协方差参数化本身仍将运动与几何耦合在单一矩阵中——协方差的更新同时影响轨迹和形状，优化时二者相互干扰。VeGaS（2026.03）~\cite{2603_00952}（Decoupling Motion and Geometry in 4D Gaussian Splatting）从参数化层面根本解决这一耦合问题：不改变度量方式，而是重新设计 4D 协方差的构造，使运动参数和几何参数在数学上完全解耦。

\textbf{耦合诊断。}在标准 4DGS 中，4D 高斯在时间 $t$ 的条件 3D 分布为
\begin{equation}
\mu_{xyz|t} = \mu_{1:3} + (t - \mu_t)\,\Sigma_{1:3,4}\,\Sigma_{4,4}^{-1}, \quad
\Sigma_{xyz|t} = \Sigma_{1:3,1:3} - \Sigma_{1:3,4}\,\Sigma_{4,4}^{-1}\,\Sigma_{4,1:3},
\end{equation}

问题一目了然：均值以恒定速度 $\Sigma_{1:3,4}\Sigma_{4,4}^{-1}$ 线性演化，且条件协方差与时间无关——几何形状不随时间变化，运动仅为匀速直线。$\Sigma$ 的更新同时改变速度项和形状项，运动与几何在参数空间中纠缠，优化时互相干扰，常导致视觉伪影（如窗户外背景扭曲、牛排截面纹理退化）。

\textbf{Galilean 剪切解耦。}VeGaS 的核心改进是引入基于 Galilean 变换的剪切矩阵（shearing matrix）实现运动-几何解耦。剪切矩阵
\begin{equation}
V = \begin{pmatrix} I_3 & v(t) \\ 0 & 1 \end{pmatrix}
\end{equation}

将时变速度 $v(t)$ 引入 4D 协方差，通过合同变换得到 $\Sigma' = V \Sigma V^T$。关键理论保证是 Schur 补不变性：

\begin{equation}
\mathcal{S}_{4,4}(V\Sigma V^T) = \mathcal{S}_{4,4}(\Sigma),
\end{equation}

即剪切变换不影响任何时刻 $t$ 的条件 3D 协方差——高斯的 3D 形状和朝向完全由几何参数决定，不受运动干扰。由于 $\det(V)=1$，该变换还保持了 4D 高斯体积和概率质量守恒。变换后的条件均值为 $\mu_{xyz|t}' = \mu_{1:3} + (\Sigma_{1:3,4}\Sigma_{4,4}^{-1} + v(t))(t - \mu_t)$，速度项变为原始常速度与可学习时变速度之和。

\textbf{非线性轨迹积分。}时变速度 $v(t)$ 由 $N_v$ 个等间距速度锚点参数化，通过线性插值获得任意时刻的瞬时速度。累积位移 $\Delta\mu_{3D}$ 通过对总瞬时速度的时间积分计算：

\begin{equation}
\Delta\mu_{3D} = \int_{\mu_t}^{t}(v(\tau) + v_0)\,d\tau = \int_{\mu_t}^{t}v(\tau)\,d\tau + v_0(t - \mu_t),
\end{equation}

其中 $v_0 = \Sigma_{1:3,4}\Sigma_{4,4}^{-1}$ 为时间无关分量。对于跨越多个锚点区间的积分，方法利用前缀和将内部完整区间的累积位移预计算，实现任意查询时刻的 $O(1)$ 位移查询。

\textbf{几何形变网络。}运动解耦之后，复杂非刚性形变（如衣物褶皱、肌肉运动）仍需额外建模。几何形变由独立的轻量变形网络 $F_\theta$ 处理：

\begin{equation}
\Delta s,\;\Delta q,\;q_r = F_\theta\bigl(\gamma(\mu_{3D}),\;\gamma(\mu_t),\;\gamma(t_q),\;\gamma(V)\bigr),
\end{equation}

其中 $\mu_{3D} \in \mathbb{R}^3$ 为规范中心，$\mu_t$ 为时间均值，$t_q$ 为查询时间，$V \in \mathbb{R}^{N_v \times 3}$ 为所有锚点速度的拼接特征。各输入经位置编码 $\gamma(x) = \bigl[\sin(2^i x),\,\cos(2^i x)\bigr]_{i=0}^{L_x-1}$ 映射到高维空间后拼接送入 MLP（含 Layer Normalization 和 ReLU），输出缩放残差 $\Delta s \in \mathbb{R}^4$ 和四元数残差 $q_r$。更新后的缩放和旋转通过双四元数策略构建 4D 旋转矩阵 $R' = q' \otimes q_r'$。值得注意的是，变形网络同时接收速度特征 $V$ 作为输入，使几何形变能够感知运动模式——同一高斯在不同运动状态下可以产生不同的形状调整。

\textbf{实验验证。}VeGaS 在 Neu3DV 多视角真实数据集上全面超越基线：相比 4DGS，PSNR 从 32.01 提升至 32.68 dB（+0.67 dB），LPIPS 从 0.10 降至 0.09（相对提升超过 10\%）。消融实验揭示了运动建模与几何建模的互补性：仅保留时变速度分量时，刚性运动物体的轨迹更清晰（如肉夹运动轨迹更锐利）；仅保留几何变形网络时，高变形物体（如火焰）的形状还原更忠实。二者联合后，VeGaS 不仅视觉质量提升，还在效率上表现出显著优势——所需高斯数量从 4DGS 的 3572K 降至 1690K，节省约 1000 MB 存储，同时保持相当的渲染速度。速度锚点数量的消融表明 6 个锚点为最优平衡：更多锚点（如 8 个）不再带来性能提升，甚至因引入冗余运动基而降低重建质量。在 D-NeRF 单目合成数据集上，VeGaS 同样取得所有指标的最优，验证了该方法在缺乏多视角约束的单目设定下仍能可靠利用时间相干性进行重建。最终损失函数为 $L = L_{L_1} + (1 - \lambda)L_{D\text{-}SSIM}$。

\paragraph{运动混合器插值多层次变形}

现有动态高斯方法将运动编码在神经网络（如 Deformable-GS 的变形场）或逐高斯参数（如 Shape-of-Motion 的运动系数）中，这种隐式表示使得重建后的运动无法被直接操纵——用户只能重放而不能编辑运动轨迹，也无法将运动迁移到新物体。Motion Blender Gaussian Splatting（MBGS）（2025.03）~\cite{2503_09040} 的核心动机是引入显式且稀疏的运动表示，使运动在 3D 空间中可可视化、可编辑、可迁移。方法的设计灵感源自经典的网格动画技术（变形图、调和坐标、包围盒变形），但将它们从手工操作的网格表面推广到可微渲染的 3D 高斯表示。

\textbf{核心框架。}方法基于两个关键设计：第一，用稀疏运动图（motion graph）$\mathcal{G}_t = (\mathcal{J}_t, \mathcal{L})$ 替代逐高斯运动参数，其中 $\mathcal{J}_t$ 为关节集合、$\mathcal{L}$ 为边集合，满足 $\max(|\mathcal{J}_t|, |\mathcal{L}|) \ll |\mathcal{G}|$（通常不到 100 个节点，而高斯可达数十万）；第二，通过可学习的权重绘制函数（weight painting）$\mathcal{W}$ 将链接运动传播到每个高斯，而非将运动隐式编码在网络中。形式化地，高斯在时间 $t$ 的位姿由运动混合公式给出：

\begin{equation}
\mathbf{p}_t = \mathcal{B}\big(\mathcal{R}(P_{\mathcal{L},0}, P_{\mathcal{L},t}),\; \mathcal{W}(\mathbf{x}_0, X_{\mathcal{L},0})\big) \cdot \mathbf{p}_0
\end{equation}

其中 $P_{\mathcal{L},t} \in SE(3)^{|\mathcal{L}|}$ 是所有链接在时间 $t$ 的位姿矩阵，$\mathcal{R}: SE(3)^{|\mathcal{L}|} \times SE(3)^{|\mathcal{L}|} \mapsto SE(3)^{|\mathcal{L}|}$ 计算链接从初始帧到当前帧的相对 SE(3) 变换，$\mathcal{W}: \mathbb{R}^3 \times \mathbb{R}^{|\mathcal{L}| \times 6} \mapsto \Delta^{|\mathcal{L}|-1}$ 基于高斯初始位置 $\mathbf{x}_0$ 和链接初始位置 $X_{\mathcal{L},0}$ 估计每个链接对每个高斯的影响力（输出为概率单纯形），$\mathcal{B}: SE(3)^{|\mathcal{L}|} \times \Delta^{|\mathcal{L}|-1} \mapsto SE(3)$ 通过双四元数蒙皮（dual quaternion skinning, DQS）将链接运动加权混合后传播到高斯。DQS 用对偶四元数表示变换并在该空间中计算加权平均，保证混合结果始终是合法的 SE(3) 变换。

\textbf{权重绘制函数。}权重绘制函数定义为 $\mathcal{W}(\mathbf{x}_0, X_{\mathcal{L},0}) = \text{softmax}(\{K(\mathbf{x}_0, X_{\mathcal{L},0,i}) \mid \forall i \in [1,|\mathcal{L}|]\})$，其中核函数
\begin{equation}
K(\mathbf{x}_0, X_{\mathcal{L},0,i}) = \exp(-\gamma \cdot \text{dist}(\mathbf{x}_0, X_{\mathcal{L},0,i}))
\end{equation}

度量高斯与链接的亲和力，$\text{dist}(a, b)$ 是点到线段的距离，$\gamma$ 是可学习的核半径参数。每个链接学习独立的 $\gamma$，使运动图的影响区域自适应覆盖物体表面。在 $t=0$ 时刻，3D 高斯重建物体的初始几何，通过权重绘制与动态运动图绑定；绘制的权重随后用于 $t > 0$ 时刻的运动传播和动态重建。

\textbf{运动学树（kinematic tree）。}运动学树是一种层次化的无环图，适用于关节物体（人体、机器人手臂等）。参数化为 $\theta = (\mathbf{r}, \mathbf{X})$、$\phi = (\ell_i)_{i=1}^{L}$，其中 $\ell_i \in \mathbb{R}^+$ 为第 $i$ 个链接的长度，$\mathbf{r} \in SO(3)^{|\mathcal{J}_t|-1}$ 为关节旋转，$\mathbf{X} \in SE(3)$ 为根节点位姿。链接位姿 $P_{\mathcal{L},t} = \mathcal{P}_\mathcal{L}(\theta_t, \phi)$ 通过可微正向运动学算法计算：变换从根节点沿树结构逐层传播，逐关节累积旋转和平移，推导出每个链接的世界坐标位姿。正向运动学过程完全可微，支持通过反向传播学习 $\theta$ 和 $\phi$。

\textbf{可变形图（deformable graph）。}可变形图不施加拓扑约束，关节可在 3D 空间中自由移动，适用于非刚性变形物体（软体）。参数化为 $\theta = (\mathbf{n}_j \in \mathbb{R}^3)_{j=1}^{|\mathcal{J}_t|}$，其中 $\mathbf{n}_j$ 为第 $j$ 个关节的 3D 位置。由于缺少拓扑约束，链接可以拉伸或压缩，无法直接用刚体变换描述链接位姿。为此，MBGS 将每个高斯 $\mathbf{x}$ 在初始帧投影到每个链接上，并在后续帧中追踪投影点位置。具体地，对于链接 $l$ 的投影点在时间 $t$ 的位置为
\begin{equation}
\mathbf{n}_{\mathbf{x}_t,l} = \mathbf{n}_{s_l,t} + \frac{|\mathbf{n}_{\mathbf{x}_0,l} - \mathbf{n}_{s_l,0}|}{|\mathbf{n}_{s_l,0} - \mathbf{n}_{e_l,0}|} (\mathbf{n}_{s_l,t} - \mathbf{n}_{e_l,t})
\end{equation}

即投影点随链接拉伸按比例移动。链接位姿通过 look-at 变换获得：$P_\mathcal{L}(\theta_t, \mathbf{x}_t) = \{\mathscr{A}(\mathbf{n}_{\mathbf{x}_t,l}, \text{ray}(\mathbf{n}_{s_l,t}, \mathbf{n}_{e_l,t})) \mid \forall l \in \mathcal{L}\}$，其中 $\text{ray}(\mathbf{a}, \mathbf{b}) = (\mathbf{b} - \mathbf{a})/|\mathbf{b} - \mathbf{a}|$ 为单位方向向量，$\mathscr{A}: \mathbb{R}^3 \times S^2 \mapsto SE(3)$ 将观测位置和方向映射为 SE(3) 位姿。需要指出，MBGS 的可变形图与 Dynamic Fusion 的变形图有本质区别：Dynamic Fusion 仅在关节上附着 SE(3) 位姿，而 MBGS 将关节表示为 $\mathbb{R}^3$ 位置并在链接上导出 SE(3) 位姿，简化了运动图的操纵。

\textbf{运动图初始化与实例分割。}运动图的连接性 $\mathcal{L}$ 作为固定参数，不接收梯度也不随时间变化。运动学树的连接由领域先验确定（如从 2D 人体骨架提升到 3D、或从机器人 URDF 模型提取）；可变形图通过最远点采样和随机连接初始化，确保均匀覆盖物体表面。值得注意的是，MBGS 利用 Grounding SAM2 进行实例分割、SAPIENS 估计 2D 人体骨架，在规范帧（$t=0$）初始化运动图，使框架能够实现实例级重建——高斯按实例显式分组以保持精确的实例几何，这一能力在此前的动态高斯文献中尚未被探索。

\textbf{运动操纵与下游应用。}由于运动图是显式的参数化模型，用户可以在 3D 空间中直接可视化和操纵运动。在动画新姿态方面，保持高斯固定、调整运动图的关节旋转或关节位置，即可渲染出训练视频中未出现的新姿态；创建短动画时，通过增量插值关节角度或位置变化并逐帧渲染。在机器人演示合成方面，分别从人类和机器人视频重建场景，以机器人的预定义运动学链作为运动图，移除人类高斯并添加机器人高斯到工作空间，通过逆运动学将关键帧映射到机器人关节角度，生成机器人模仿人类操作的演示视频。在视觉规划方面，从人类视频重建物体高斯后，保持高斯和权重固定，通过最小化 L1 渲染损失快速学习运动图（通常一分钟内收敛），然后模拟多种图轨迹并渲染结果图像，选择与目标图像 PSNR 最高的轨迹作为最优动作——奖励函数仅需一张目标图像即可指定。

\textbf{实验评估。}在 iPhone 数据集（5 个多相机场景 + 7 个单相机场景，$960 \times 720$）上，MBGS 在主要指标 LPIPS 上达到 0.37，超越 Shape-of-Motion 的 0.39；PSNR 为 16.79 vs 16.67，SSIM 均为 0.65。在 HyperNeRF vrig 数据集（4 个场景，$960 \times 540$）上，MBGS 的 LPIPS 为 0.35，与 4DGaussians 的 0.36 和 Shape-of-Motion 的 0.34 处于竞争水平；虽然 PSNR（20.60）低于 4DGaussians（25.19），但在 chicken、3D printer、broom 三个场景上 MBGS 渲染质量更清晰——例如在 3D printer 场景中准确渲染了打印机电机上的文字，而 4DGaussians 的输出模糊。渲染速度方面，在 teddy bear 场景（200 万高斯，可变形图 + 运动学树）上达到 18 FPS（缓存运动图后 25 FPS），在 HyperNeRF 的 chicken 场景（30 万高斯）上达到 32 FPS（缓存后 46 FPS）。训练在 40GB A100 GPU 上收敛需 10--30 小时。

\paragraph{时间刚体运动约束与射线分组}

回顾上文 FreeTimeGS 的运动模型：$\boldsymbol{\mu}_i(t) = \boldsymbol{\mu}_i + \mathbf{v}_i(t - \mu_{t,i})$，每个高斯以恒定速度沿直线运动，时间可见性由高斯形状的时间窗口控制。这一设计在短序列上效果良好，但面对弯曲轨迹时暴露出根本缺陷——线性假设逐渐偏离真实几何，迫使系统反复消除、分裂和重建高斯基元。TRiGS（2026.04）~\cite{2604_00538}（Temporal Rigid-Body Motion for Scalable 4D Gaussian Splatting）将这一现象命名为\emph{时间碎片化}（temporal fragmentation）：高斯基元丧失了长期时间身份（temporal identity），无法沿弧形轨迹稳定追踪，系统只能通过不断增殖新基元来覆盖运动误差，导致高斯数量激增和内存膨胀。TRiGS 的核心思路是用连续刚体变换替代分段线性近似，使单个高斯能够沿完整弧形轨迹稳定运动而无需反复重建。该方法通过三个协同组件实现这一目标。

\textbf{耦合刚体运动参数化。} TRiGS 在李代数 $\mathfrak{se}(3)$ 中参数化运动，而非独立优化旋转和平移参数。具体而言，对每个高斯基元 $i$ 定义时间条件的李代数系数 $\boldsymbol{\zeta}_i(t) = (\boldsymbol{\omega}_i(t), \boldsymbol{\nu}_i(t)) \in \mathfrak{se}(3)$，构造相对对数参数 $\mathbf{u}_i(t) := \boldsymbol{\zeta}_i(t)\,\Delta t$（其中 $\Delta t := t - \mu_{t,i}$），然后通过指数映射将其映射到 $SE(3)$ 刚体变换：

\begin{equation}
\mathbf{T}_i(\mu_{t,i} \!\to\! t) = \exp(\hat{\mathbf{u}}_i(t)) = \begin{bmatrix} R_i(t) & \mathbf{p}_i(t) \\ \mathbf{0}^\top & 1 \end{bmatrix}, \quad \mathbf{u}_i(\mu_{t,i}) = \mathbf{0}.
\end{equation}

这一构造满足 $\mathbf{T}_i(\mu_{t,i}) = \mathbf{I}$（中心时刻恒等变换）。$R_i(t)$ 通过 Rodrigues 公式计算，$\mathbf{p}_i(t) = J(\boldsymbol{\phi}_i(t))\,\boldsymbol{\upsilon}_i(t)$ 由 $SO(3)$ 左 Jacobian $J(\cdot)$ 耦合确定，而非独立优化平移。由于 $\boldsymbol{\omega}$ 和 $\boldsymbol{\nu}$ 共享同一指数映射，旋转和平移被天然耦合，避免了两者互相补偿导致的优化歧义。论文的消融实验有力验证了这一设计：当采用独立的 $SO(3) \times \mathbb{R}^3$ 参数化（Naive SE(3)）时，在 1200 帧序列上 PSNR 仅为 23.95，甚至低于纯线性基线的 24.43——因为独立优化允许旋转和平移互相补偿，产生连续的近等价解，导致运动分解纠缠和时间累积漂移。

\textbf{层次分解与 B\'ezier 残差。} 为了在耦合框架内进一步提升非线性运动的表达能力，TRiGS 对李代数系数进行层次分解：$\boldsymbol{\zeta}_i(t) = \boldsymbol{\zeta}_i^{\text{base}} + \boldsymbol{\zeta}_i^{\text{res}}(\tau)$，其中基础项 $\boldsymbol{\zeta}_i^{\text{base}} = (\boldsymbol{\omega}_i^{\text{base}}, \boldsymbol{\nu}_i^{\text{base}})$ 为可学习的常值，残差项 $\boldsymbol{\zeta}_i^{\text{res}}(\tau)$ 由二次 B\'ezier 曲线
\begin{equation}
\boldsymbol{\zeta}_i^{\text{res}}(\tau) = (1-\tau)^2\,\boldsymbol{\zeta}_i^{(0)} + 2(1-\tau)\tau\,\boldsymbol{\zeta}_i^{(1)} + \tau^2\,\boldsymbol{\zeta}_i^{(2)}
\end{equation}

参数化，$\tau \in [0,1]$ 为有效时间窗口 $\mathcal{W}_i = [\mu_{t,i} - 2s_{t,i},\, \mu_{t,i} + 2s_{t,i}]$ 内的归一化坐标。将 B\'ezier 残差限定在有效窗口内的设计动机是：由于时间可见性 $\gamma_i(t)$ 在远离 $\mu_{t,i}$ 时快速衰减，基元在远处时刻几乎不接收梯度，无约束的参数化会导致弱监督区域的优化漂移。此外，B\'ezier 曲线的二阶差分 $\mathbf{b}_i = \boldsymbol{\zeta}_i^{(0)} - 2\boldsymbol{\zeta}_i^{(1)} + \boldsymbol{\zeta}_i^{(2)}$ 对应加速度项，通过运动平滑正则 $\mathcal{L}_{\text{motion}} = \|\mathbf{b}_i\|_2^2$ 抑制轨迹的急剧弯曲，促进时间平滑性。论文消融表明，仅对无约束的平移施加非线性 B\'ezier 曲线（Linear + B\'ezier 配置）仅将 PSNR 从 24.43 提升至 24.56，远不如结合局部锚点的耦合 SE(3) 变换（25.52），证实了几何一致性约束的关键作用。

\textbf{局部锚点中心变形。} 上述 $SE(3)$ 变换若以全局原点为中心，则隐式假设所有基元围绕同一中心旋转，无法捕捉场景中多个部件的独立局部运动。TRiGS 为每个高斯引入可学习的局部锚点 $\mathbf{a}_i \in \mathbb{R}^3$，将变换围绕锚点执行。然而锚点与平移之间存在不可辨识性：同一变形可由多组 $(\mathbf{a}_i, \mathbf{p}_i(t))$ 等价解释。为此，TRiGS 将锚点投影到旋转轴正交平面：$\mathbf{a}_{i,\perp}(t) = \mathbf{a}_i - \langle \mathbf{a}_i, \bar{\boldsymbol{\omega}}_i(t)\rangle\,\bar{\boldsymbol{\omega}}_i(t)$，其中 $\bar{\boldsymbol{\omega}}_i(t)$ 是单位旋转轴方向。这一 gauge fixing 利用了 $(I - R_i(t))\bar{\boldsymbol{\omega}}_i(t) = \mathbf{0}$ 的性质——锚点沿旋转轴方向的分量位于 $(I - R_i(t))$ 的零空间中，对变形无贡献，因此可以安全消除。最终的变形公式为：

\begin{equation}
\boldsymbol{\mu}_i(t) = R_i(t)(\boldsymbol{\mu}_i - \mathbf{a}_{i,\perp}(t)) + \mathbf{a}_{i,\perp}(t) + \mathbf{p}_{i,\parallel}(t), \quad \Sigma_i(t) = R_i(t)\,\Sigma_i\,R_i(t)^\top,
\end{equation}

其中 $\mathbf{p}_{i,\parallel}(t) = J(\boldsymbol{\phi}_i(t))\,\boldsymbol{\upsilon}_{i,\parallel}(t)$ 仅保留沿旋转轴方向的平移分量。这一机制使不同部件可以围绕各自的旋转中心独立运动。消融实验证实 gauge fixing 的关键性：移除它导致 PSNR 从 26.05 下降至 25.15，是所有优化组件中影响最大的因素。

\textbf{训练目标与运动引导重定位。} TRiGS 的完整训练目标为：

\begin{equation}
\mathcal{L} = \lambda_{\text{img}}\mathcal{L}_{\text{img}} + \lambda_{\text{ssim}}\mathcal{L}_{\text{ssim}} + \lambda_{\text{reg}}\mathcal{L}_{\text{reg}} + \lambda_{\text{motion}}\mathcal{L}_{\text{motion}} + \lambda_{\text{rigid}}\mathcal{L}_{\text{rigid}},
\end{equation}

其中刚性正则 $\mathcal{L}_{\text{rigid}} = \sum_i \sum_{j \in \mathcal{N}(i)} K_{ij}(\|\boldsymbol{\nu}_i^{\text{base}} - \boldsymbol{\nu}_j^{\text{base}}\|_2^2 + \|\boldsymbol{\omega}_i^{\text{base}} - \boldsymbol{\omega}_j^{\text{base}}\|_2^2)$ 促进空间邻域的基础运动一致性，$K_{ij} = \exp(-\lambda_c \|\mathbf{c}_i - \mathbf{c}_j\|_2^2)$ 基于球谐 DC 颜色的外观亲和度权重抑制跨物体边界的错误耦合。为了在固定基元预算下重用容量，TRiGS 每隔 100 次迭代执行运动引导重定位（motion-guided relocation）：将低不透明度基元替换为高信息量活跃基元的克隆，采样概率基于不透明度梯度、2D 梯度、时间梯度和速度变化等多维难度线索的加权和，保持总基元数不变的同时将容量动态分配到复杂动态区域。

\textbf{实验验证。} TRiGS 的核心优势在长序列上尤为突出。在 SelfCap 数据集的扩展序列上（600/900/1200 帧），FTGS 的 PSNR 从 26.23 下降至 25.41，高斯数量翻倍，内存达 977 MB；而 TRiGS 始终保持 0.5M 高斯和 160 MB 内存在 1200 帧时仍维持 26.05 PSNR，帧率超过 110 FPS。在标准长度的 N3V 数据集（约 300 帧）上，TRiGS 同样达到最优 PSNR 33.36 和最低 LPIPS 0.031。即使在消融实验中允许基线方法自由增殖（Baseline 配置增长至 2.85M 高斯），其渲染质量仍无法匹敌仅用 0.5M 高斯的 TRiGS，有力证明了结构化运动参数化优于暴力容量扩展。

\paragraph{松弛刚性与射线分组去噪}

前述 TRiGS 通过连续刚体变换增强单个高斯的运动表达能力，但仍从基元内部视角出发建模运动。当视角转向基元间的运动协调性时，核心问题变为：如何在无需外部先验的条件下，让相邻高斯的运动保持物理上的一致性？Relaxed Rigidity（2026.03）~\cite{2603_24994} 正是从这一角度出发，提出一种模型无关的正则化框架，可即插即用地集成到任意动态 3DGS 架构中。

已有方法面临两类局限：基于外部先验（光流、深度）的方法将高斯运动投影到 2D 图像平面进行监督，但由于监督定义在屏幕空间而非底层 3D 几何中，先验信号中的误差和歧义会传播到优化过程；基于严格刚性的方法用 KNN 距离分组高斯并强制组内共享相同变换，但忽略了真实运动的非刚性和交互本质，也无法适应高斯原始体的不同空间尺度。Relaxed Rigidity 的设计动机是从渲染管线本身获取分组线索，而非依赖外部先验或距离度量。

\textbf{基于射线的分组。}方法的核心是两个互补机制，首先是基于射线的分组（ray-based grouping）。在 3DGS 的光栅化管线中，每个像素射线已经对高斯进行了排序和 $\alpha$-混合，方法直接复用这一过程，对每个像素 $p_j$ 定义射线组 $N_j = \{G_i \mid w_i > \tau\}$，其中 $w_i = T_i(1 - e^{-o_i})$ 是高斯的混合权重（contribution weight），$\tau$ 为阈值。混合权重 $w_i$ 天然充当了遮挡感知的过滤器——由于累积透射率在不透明遮挡物后方急剧衰减，按 $\tau$ 阈值筛选自然隔离了同一未遮挡表面上空间连续的高斯，而不会将前景和背景的高斯混淆。这一分组无需额外计算开销，且组大小根据局部复杂度自适应变化（论文在 D-NeRF 和 HyperNeRF 上的分组大小分布显示，组大小从几个到数十个高斯不等，宽分布特性使方法能动态适应从细薄结构到密集体素的不同局部复杂度）。相比 KNN 基于欧氏距离的分组忽略高斯的尺度和不透明度属性，射线分组通过可见性建立了 2D 观测到 3D 几何的直接桥梁，避免了物理距离近但结构独立的部件被错误聚类的问题。

\textbf{运动一致性正则（MCR）。}其次是松弛刚性约束，包含两项正则化。MCR 鼓励组内高斯运动方向一致但允许幅度差异：对于时间偏移 $\Delta t$（从区间 $[t-m, t+m]$ 均匀采样），计算高斯位移 $\mathbf{d}_{i,t} = \boldsymbol{\mu}_{i,t+\Delta t} - \boldsymbol{\mu}_{i,t}$ 和组内平均位移 $\bar{\mathbf{d}}_{N_j,t} = \frac{1}{|N_j|}\sum_{G_i \in N_j} \mathbf{d}_{i,t}$，然后通过余弦相似度损失
\begin{equation}
L_{\text{MCR}} = 1 - \frac{\mathbf{d}_{i,t} \cdot \bar{\mathbf{d}}_{N_j,t}}{\|\mathbf{d}_{i,t}\|\|\bar{\mathbf{d}}_{N_j,t}\| + \epsilon}
\end{equation}

惩罚方向不一致（仅对位移 $\|\mathbf{d}_{i,t}\| > 0.0001$ 的运动高斯施加）。关键设计是不惩罚位移幅度差异——如果强制所有高斯具有相同位移幅度，就等价于严格刚性平移，会错误地阻碍自然的非刚性变形。

\textbf{谱正则化（SR）。}SR 在 MCR 对齐运动方向的基础上进一步保持局部形状结构：计算射线组在 $t$ 和 $t+\Delta t$ 时刻高斯位置的协方差矩阵 $K_t$ 和 $K_{t+\Delta t}$，惩罚其特征值谱的差异
\begin{equation}
L_{\text{SR}} = \sum_{r \in \{1,2,3\}} \text{Huber}(\sigma_{t,r},\, \sigma_{t+\Delta t, r})
\end{equation}

其中 $\sigma_{t,r}$ 和 $\sigma_{t+\Delta t,r}$ 分别为 $K_t$ 和 $K_{t+\Delta t}$ 按升序排列的第 $r$ 个特征值，Huber 为标准 Huber 损失。SR 仅应用于包含多于一个高斯的组。匹配特征值谱保持局部形状统计量随时间不变，同时天然对刚体旋转不变（因为旋转不改变特征值），且允许灵活的非刚性运动（因为只要形状统计量不变，高斯可以自由重新排列）。此外，由于损失函数不最小化特征值谱的整体尺度，正则化不会迫使组收缩——即使组内包含距离较远的高斯也是如此。正则化仅选择性地约束那些位移大到足以改变物体形状的高斯，引导其维持与组内邻居的空间邻近性。

\textbf{Welford 算法与高效实现。}为避免传统两遍算法需要存储所有中间值并多次运行光栅化管线的问题，方法采用 Welford 算法实现协方差的在线单遍递推计算：$K_i = \frac{i-1}{i}K_{i-1} + \frac{1}{i}P_i$，其中 $P_i = (\boldsymbol{\mu}_i - \bar{\boldsymbol{\mu}}_{i-1})(\boldsymbol{\mu}_i - \bar{\boldsymbol{\mu}}_i)^\top$，$\bar{\boldsymbol{\mu}}_i$ 为运行均值。该递推关系可归纳为 $K_N = \frac{1}{N}\sum_{i=1}^N P_i$，使得沿每条射线的协方差计算只需顺序引用每个高斯一次，直接集成到 CUDA 光栅化管线中，支持前向和反向传播且内存开销有限。

\textbf{最终训练目标。}最终的训练目标为
\begin{equation}
L = (1-\lambda_{\text{dssim}})L_1 + \lambda_{\text{dssim}}L_{\text{dssim}} + \lambda_{\text{MCR}}L_{\text{MCR}} + \lambda_{\text{SR}}L_{\text{SR}}
\end{equation}

其中可见性阈值 $\tau \in [10^{-4}, 10^{-3}]$，$\lambda_{\text{MCR}} \in [10^{-4}, 5 \times 10^{-2}]$，$\lambda_{\text{SR}} \in [10^{-4}, 1]$，时间采样范围 $m \in [2, 20]$ 均按数据集调优。该方法不修改底层模型架构，因此在渲染阶段不引入额外开销，训练时间增加约 2--3 倍。

\textbf{实验验证与模型泛化性。}方法的模型无关性通过集成到四类代表性基线得到验证：变形场类的 RTD（HexPlane 参数化）、MoDec-GS（脚手架表示）、Grid4D（哈希网格编码）和基函数轨迹类的 Ex4DGS（样条参数化）。在 D-NeRF 数据集上，PSNR 平均提升 1.19 dB，其中 MoDec-GS + Ours 提升 +2.35 dB（29.33 $\to$ 31.68），Ex4DGS + Ours 提升 +1.11 dB（29.93 $\to$ 31.04），Grid4D + Ours 达到最高 PSNR 42.20。在 HyperNeRF 数据集上，Grid4D + Ours 获得最高 MS-SSIM 0.856。在更具挑战性的 NeRF-DS 数据集上，MoDec-GS + Ours 提升 +0.83 dB，Grid4D + Ours 提升 +0.32 dB。

消融实验揭示了三个关键发现：(1) 射线分组在 Full 配置（MCR + SR）下比 KNN 分组在 D-NeRF 上高 3.17 dB（36.46 vs. 33.29），证明基于渲染可见性的分组远优于基于欧氏距离的分组；(2) 将 SR 替换为传统 ARAP 刚性约束（RTD+RG\&ARAP，D-NeRF 35.27 dB）相比 Full 方案（36.46 dB）性能下降 1.19 dB，因为 ARAP 的逐点刚性限制了细粒度形状变化，阻碍了对详细几何变化的捕捉；(3) MCR 和 SR 单独使用均不足——MCR 在 D-NeRF 上效果显著（36.23 dB）但在 HyperNeRF 的复杂动力学中效果减弱（24.23 MS-SSIM），SR 单独在两个数据集上表现均不充分，两者互补才能实现一致的性能提升。此外，射线分组在 Full 配置下比 KNN 分组训练速度快 6\%--25\%。

\paragraph{中继高斯级联重建大幅运动}

RelayGS（2024.12）~\cite{2412_02493} 专门解决包含大幅运动和复杂运动的动态场景重建。已有 4DGS 方法（如 4D-GS）将规范高斯与神经运动场耦合优化，神经网络不仅难以预测大幅运动，还倾向于过拟合场景中占主导的小幅运动，限制了其对复杂大幅运动的建模能力。RelayGS 的核心动机是通过渐进式三阶段策略简化复杂运动轨迹的学习，将大幅运动分解为多个可控的小段。这种设计可视为一种\textbf{时间致密化}（temporal densification）策略——类比 3DGS 中通过分裂和克隆实现空间致密化，RelayGS 通过将前景高斯复制到多个时间片段实现沿时间轴的致密化，使每个副本只需覆盖短段轨迹。

\textbf{阶段一：初始表示与前背景解耦。}用所有帧训练一个静态 3DGS 以捕获共享背景，同时为每个高斯引入可学习掩码将高动态前景与低动态背景分离。掩码通过直通估计器（straight-through estimator）实现二值化的同时保持梯度传播：

\begin{equation}
M_n = \text{sg}\bigl(\mathbb{1}[\sigma(m_n) > \epsilon] - \sigma(m_n)\bigr) + \sigma(m_n) = \begin{cases} 1, & \text{if } \sigma(m_n) > \epsilon \\ 0, & \text{otherwise} \end{cases}
\end{equation}

其中 $m_n \in \mathbb{R}$ 为可学习掩码参数，$\sigma(\cdot)$ 为 sigmoid 函数，$\epsilon$ 为阈值，$\text{sg}(\cdot)$ 为停止梯度算子，$\mathbb{1}[\cdot]$ 为指示函数。尽管 $M_n$ 输出为二值，梯度仍可通过 $\sigma(m_n)$ 分支回传至 $m_n$。渲染时，所有高斯用于渲染第一帧（忽略掩码），其他帧将每个高斯的不透明度替换为 $\hat{o}_n = M_n \cdot o_n$。前景高斯在第一帧位置正确但在其他帧因大幅运动产生较大误差，导致高梯度持续降低 $M_n$，最终 $M_n \to 0$ 实现前背景解耦。掩码参数初始化为 $m_n = 2$（经 sigmoid 后接近 1），经 3000 步优化与周期性致密化完成初始表示学习。

\textbf{阶段二：中继高斯分解大幅运动轨迹。}将解耦的前景高斯复制多份，每份对应一个时间片段（连续 $k=16$ 帧为一个片段），这些副本作为大幅运动轨迹上的显式中继节点，将完整的复杂轨迹分解为多个短段。为了优化和加密中继高斯，方法构建伪视图：从每个片段中均匀选取 $p=3$ 帧，按混合权重
\begin{equation}
I_{\text{pseudo}} = \beta_1 I_{t_1} + \beta_2 I_{t_2} + \beta_3 I_{t_3}, \quad \beta_1 + \beta_2 + \beta_3 = 1
\end{equation}

生成伪视图作为监督信号。默认采用均匀混合 $\beta_1 = \beta_2 = \beta_3 = \frac{1}{3}$，这些伪视图捕获前景在不同时刻的快照，提供更丰富的表示用于优化中继高斯，确保其准确捕获片段内的运动轨迹。此阶段训练 14000 步，消融实验表明去掉此阶段（直接将前景高斯与隐式运动场联合学习）会导致性能显著下降，尤其在复杂场景中。

\textbf{阶段三：4D 时空联合建模与分治运动场。}采用 HexPlane 加轻量 MLP 的运动场，但采用分治策略——背景和前景使用独立的 MLP 分支，避免共享 MLP 导致的小幅运动过拟合。为进一步增强大幅运动的捕获能力，方法为每个中继高斯引入可学习位置变形缩放因子 $\boldsymbol{\gamma} \in \mathbb{R}^3$，调整 MLP 预测的位置偏移量以适应更大的运动范围：

\begin{equation}
\boldsymbol{\mu} \leftarrow \boldsymbol{\mu} + (1 + e^{\boldsymbol{\gamma}}) \cdot \Delta\boldsymbol{\mu}
\end{equation}

其中 $\Delta\boldsymbol{\mu}$ 为 MLP 预测的位置变化，$\boldsymbol{\gamma}$ 初始化为 $\mathbf{0}$。背景 MLP 分支和前景 MLP 分支各自独立预测位置、缩放、旋转、不透明度四个属性的时变变化（不包含球谐函数 MLP，因其增大模型尺寸且无显著性能增益）。此阶段训练 20000 步。消融实验表明，去掉前景-背景 MLP 隔离（共享 MLP）或去掉缩放因子 $\boldsymbol{\gamma}$ 均导致性能下降。

\textbf{实验验证。}在 PanopticSports 数据集（6 个动态运动场景，$640 \times 360$，150 帧，31 相机中 27 训练/4 测试）上，RelayGS 在六个场景上分别比次优方法提升 0.27 dB、0.53 dB、1.60 dB、1.19 dB、1.24 dB、1.28 dB PSNR，平均提升 1.02 dB；相比基线 4D-GS 平均提升 2.47 dB。在 VRU 篮球赛数据集（真实篮球赛场景，$1920 \times 1080$，250 帧，34 相机中 30 训练/4 测试）上，GZ 场景达到 28.06 dB（ST-GS: 27.32 dB，4D-GS: 25.83 dB），DG4 场景达到 26.94 dB（ST-GS: 26.79 dB，4D-GS: 25.17 dB）。虽然与 ST-GS 的 PSNR 差值看似不大，但 VRU 视角中静态地板约占 70\% 像素，动态前景区域的实际质量提升更为显著。方法在 RTX 4090 上实现约 70 fps 的实时渲染速度。消融实验还表明，片段长度 $k$ 的选取影响性能：$k$ 增大（片段减少）时片段内运动轨迹增大，性能逐渐下降；$k$ 过小则增加训练成本且性能提升有限，$k=16$ 为经验最优值。方法的主要局限在于：(1) 小而快速运动的物体（如篮球）因像素覆盖率低和周围规范高斯稀疏，运动建模仍不充分；(2) 时间分段和伪视图构建策略较为简单，未来可探索基于运动复杂度的自适应时间分段。

\paragraph{无模板关节高斯与骨骼发现}

Template-free Articulated Gaussian Splatting（2024.12）~\cite{2412_05570} 解决从视频中自动发现动态物体的骨架模型并实现重姿态（re-posing）的挑战。与前面讨论的 RelayGS 关注多相机场景的复杂动态重建不同，本方法聚焦于单物体关节动态的骨架发现——核心差异在于不仅重建外观和运动，更要从运动中自动提取可编辑的骨架结构。已有方法需要预定义的物体模板（如人体 SMPL 模型），限制了泛化到野外场景和人造关节物体的能力；无模板方法 AP-NeRF 虽从预训练动态 NeRF 中采样特征点云再初始化骨架，但无法实时渲染。该方法的核心动机是利用 3DGS 的显式点云表示实现类别无关的自动骨架发现，同时保持实时渲染速度。方法分三阶段。

\textbf{阶段一：超点驱动的动态重建。}引入超点（superpoint）概念——将运动相似的高斯聚合为超点 $p_j \in \mathbb{R}^3$，每个超点代表物体的一个刚性部件。变形场 $f$ 将每个超点映射到其时变 6-DOF 变换：$f(p_j, t) \mapsto (R_j^t, \tau_j^t)$，其中 $R_j^t$ 为局部旋转矩阵，$\tau_j^t \in \mathbb{R}^3$ 为平移向量。高斯的运动通过线性混合蒙皮（LBS）由相邻超点的 6-DOF 变换插值得到：

\begin{equation}
\hat{\mu}_i^t = \sum_{j \in N_i} w_{ij} (R_j^t \mu_i + \tau_j^t), \quad \hat{q}_i^t = \left(\sum_{j \in N_i} w_{ij} \hat{q}_j^t\right) \otimes q_i
\end{equation}

其中 $w_{ij} = \frac{\exp(W_{ij})}{\sum_{k \in N_i}\exp(W_{ik})}$ 是由可学习参数 $W \in \mathbb{R}^{M \times N}$ 计算的 LBS 权重，$N_i$ 是高斯 $G_i$ 的 $K$ 个最近超点，$\hat{q}_j^t \in \mathbb{R}^4$ 是 $R_j^t$ 的四元数表示，$\otimes$ 为四元数乘法。变形场采用 8 层 MLP（256 维隐藏层，ReLU 激活）配合位置编码实现，同时预测位置和旋转的变形。高斯的其余属性（不透明度、缩放、球谐系数）保持规范空间值不变。

\textbf{阶段二：基于运动一致性的骨架发现。}基于超点运动自动提取关节和骨骼连接。核心观察是：若两个超点 $p_a$、$p_b$ 之间存在关节，则该关节位置在相对姿态变化中保持不变。方法计算每对超点的关节距离度量：

\begin{equation}
d_{ab} = \sum_t \| \tau_r^t - (j_{ab} - R_{ab}^t j_{ab}) \|_2^2 + \lambda_d \| j_{ab} - j_{ba} \|_2^2
\end{equation}

其中 $R_{ab}^t = (R_b^t)^{-1}R_a^t$ 是相对旋转，$\tau_r^t$ 是相对平移，$j_{ab}$ 和 $j_{ba}$ 分别是从 $p_a$ 和 $p_b$ 视角估计的关节位置。距离 $d_{ab}$ 通过动量平滑 $\hat{d}_{ab}^{(i+1)} = (1-\beta)\hat{d}_{ab}^{(i)} + d_{ab}^{(i)}$（$\beta=0.1$）防止剧烈变化，然后按最小生成树算法选择关节距离最小的超点对构建无环骨架结构 $\mathcal{S}$。根节点选择为骨架中最长路径长度的最小值对应的超点。

\textbf{阶段三：运动学优化与正向运动学链。}在发现的骨架上，变形场 $g$ 为每个关节 $k$ 预测时变旋转 $\theta_k^t$：$g(j_k, t) \mapsto \theta_k^t$。每个关节的局部变换矩阵由绕父关节的旋转 $R_k^t(\theta_k^t)$ 定义：

\begin{equation}
T_k^t = \begin{pmatrix} R_k^t(\theta_k^t) & j_k - R_k^t(\theta_k^t) j_k \\ 0 & 1 \end{pmatrix}
\end{equation}

通过正向运动学链 $T_j^t = T_{\text{root}}^t \prod_{k \in C_j} T_k^t$ 将超点从规范空间变换到观测空间，其中 $C_j$ 是超点 $j$ 在骨架中的祖先关节列表。此阶段实现端到端的骨架和外观联合优化。

\textbf{超点自适应控制策略。}为避免训练过程中产生过多冗余超点并简化最终骨架结构（便于姿态编辑），方法借鉴 3D-GS 的自适应密度控制思想，设计了三重控制机制。定义超点 $p_j$ 的影响度 $W_j = \sum_{i \in \mathcal{N}_j} w_{ij}$，其中 $\mathcal{N}_j = \{i \mid j \in N_i\}$ 是以 $p_j$ 为 $K$ 近邻之一的高斯集合。

\textbf{剪枝}：当 $W_j < \sigma_{\text{prune}}$ 时移除该超点，因其对周围高斯的运动贡献极小。

\textbf{克隆}：当 $W_j > \sigma_{\text{clone}}$（表示该超点绑定了过多高斯）或加权梯度范数 $g_j > \sigma_{\text{grad}}$ 时触发，其中加权梯度范数定义为 $g_j = \frac{\sum_{i \in \mathcal{N}_j} w_{ij} \|\nabla_{\mu_i} \mathcal{L}\|_2^2}{\sum_{k \in \mathcal{N}_j} w_{kj}}$，衡量该超点所绑定高斯的优化压力。

\textbf{合并}：对于具有相似运动模式的超点对 $(a, b)$，计算平均相对变换距离 $D_{a,b} = \frac{1}{N_t}\sum_t \|(T_b^t)^{-1}T_a^t\|$（$\|\cdot\|$ 将刚体变换矩阵转换为其李代数表示的范数），当 $D_{a,b} < \sigma_{\text{merge}}$ 时合并两个超点。消融实验表明，自适应控制将超点数量从 608 个降至 97 个，超点在运动区域的分布更加均匀，最终简化了骨架结构并提高了姿态编辑的可用性。

\textbf{实验验证。}方法在三个数据集上进行评估。在 D-NeRF 数据集（5 个人形物体 + 2 个关节物体 + 1 个多组件场景，每场景 50--200 帧）上，选取 6 个场景进行测试，方法在 PSNR 上显著优于无模板关节方法 WIM 和 AP-NeRF（不仅在视觉质量上，渲染速度也有数量级提升），同时与 3DGS 动态方法（Deformable-3D-GS、4D-GS、SP-GS）的渲染质量相当或更优，并额外具备骨架发现和重姿态能力。在 Robots 数据集（7 个拓扑结构各异的机器人，18 训练视角 / 2 测试视角）上，方法同样展现出高保真的新视角合成质量。在 ZJU-MoCap 真实数据集（5 个序列，每序列 6 个训练视角）上，虽然由于相机标定不完美导致 PSNR 略低于 WIM 和 AP-NeRF，但方法实现了高达 198.23 FPS 的实时渲染速度，而 WIM 和 AP-NeRF 的渲染速度极慢（秒级）。所有实验在单块 NVIDIA Tesla V100（32 GB）上完成。方法的主要局限在于：(1) 相机标定精度对结果影响较大，标定误差会导致真实场景的视觉质量下降；(2) 自动发现的骨架结构依赖于超点初始化质量和训练收敛情况，对于拓扑非常复杂的关节物体可能需要更多先验引导。

\paragraph{可微基元结构化动态建模}

D-Prism（2026.04）~\cite{2604_17082}（Differentiable Primitives for Structured Dynamic Modeling）关注同时捕捉结构化动态物体（如多部件组装体、关节机构）的几何和刚体运动。已有动态方法存在根本局限：面向 NVS 的 4DGS 方法无法提供显式几何表示；将整个动态物体表示为单一连续变形网格的方法缺少固有的部件级分解，无法建模接触、分离和铰接运动——例如旋转中的魔方需要被建模为多个独立部件，而非单一变形体。D-Prism 的设计动机是将基元化表示的部件级分解能力引入动态场景，使基元既能建模独立的部件几何，又能捕捉部件间的刚体运动关系。

\textbf{可微基元参数化。}具体而言，D-Prism 用一组 $K$ 个可微基元表示物体，每个基元 $P_i$ 定义为超二次曲面网格（superquadric mesh），由映射函数 $F(\theta, \phi) = [s_1 \cos^{\epsilon_1}\theta \cos^{\epsilon_2}\phi,\, s_2 \cos^{\epsilon_1}\theta \sin^{\epsilon_2}\phi,\, s_3 \sin^{\epsilon_1}\theta]$ 参数化，其中 $\epsilon_1, \epsilon_2$ 控制形状（从球体到立方体），$s_1, s_2, s_3$ 控制缩放——仅需 5 个连续参数即可表达从球体到方块、圆柱的广泛形状族。每个基元还具有运动参数：6D 旋转表示 $R_i \in \mathbb{R}^6$（可转换为旋转矩阵）和平移 $T_i \in \mathbb{R}^3$，通过变换 $x_{\text{world}} = (R_i)x + T_i$ 将基元从局部空间映射到世界空间。此外，每个基元还具有不透明度参数 $o_i$，所有参数均可通过梯度优化。

\textbf{表面绑定 3D 高斯。}为了优化外观质量，D-Prism 在每个基元表面绑定 3D 高斯。高斯中心通过重心坐标分布在基元表面上，在主训练阶段这些重心坐标保持固定。旋转矩阵 $R_v = [r_1, r_2, r_3]$ 由局部面法线和切向量构造：$r_1$ 对齐面法线，$r_2$ 从质心 $m = (v_1+v_2+v_3)/3$ 指向最近顶点 $v_1$，$r_3$ 通过对 $v_2 - m$ 相对于 $r_1, r_2$ 正交化得到。缩放 $S_v = (\epsilon_s, d_{\max}, d_{\max})$ 中法线方向缩放 $\epsilon_s = 10^{-8}$ 确保薄表面表示，面内缩放 $d_{\max} = \max_{k \in \{1,2,3\}} \|v_k - g_c\|$ 为高斯中心 $g_c$ 到三个最近高斯中心的最大距离。高斯的不透明度继承自其宿主基元，仅 SH 外观系数独立优化。

\textbf{变形网络与运动建模。}基元的变形由变形网络 $D(\gamma(T), \gamma(t)) = (\Delta T, \Delta R)$ 驱动，其中 $\gamma$ 表示位置编码，将平移 $T$ 和时间戳 $t$ 投影到高维傅里叶空间。假设每个基元对应固定的物体部件、其形状和尺度在运动中不变，变形仅更新运动参数，变形后的基元表示为 $P(T+\Delta T, R+\Delta R, \epsilon, s)$。D-Prism 还引入一个结构相同的逆变形网络，确保位置变化可从当前空间映射回规范空间。

\textbf{训练损失函数。}D-Prism 的完整损失函数整合了渲染质量、几何约束和运动正则化：

\begin{equation}
L = \lambda_1 L_{\text{gs}} + \lambda_2 L_{\text{mask}} + \lambda_{\text{over}} L_{\text{over}} + \lambda_{\text{parsi}} L_{\text{parsi}} + \lambda_{\text{vol}} L_{\text{vol}} + L_{\text{deform}}
\end{equation}

其中 $L_{\text{gs}}$ 为 3DGS 标准渲染损失，$L_{\text{mask}}$ 为基元网格光栅化得到的 mask 与真实 mask 之间的误差。$L_{\text{over}}$ 防止基元间过度重叠，$L_{\text{parsi}}$ 降低冗余基元的不透明度以利于后续移除。$L_{\text{vol}}$ 对体积过小的基元施加惩罚，由超二次曲面体积公式 $V = 4 s_1 s_2 s_3 \frac{\Gamma(\epsilon_1/2)\,\Gamma(\epsilon_2/2)}{\Gamma(\frac{\epsilon_1+\epsilon_2}{2}) \cdot B(\frac{\epsilon_1}{2}, \frac{\epsilon_1+1}{2})}$ 计算并取倒数，在优化早期鼓励基元增长、减少冗余。运动正则化 $L_{\text{deform}}$ 包含三项：$L_{\text{smooth}}$ 惩罚相邻帧的平移加速度 $a_t = \Delta T_t - \Delta T_{t-1}$ 和 SO(3) 旋转加速度 $\delta R_t = (\Delta R_t)(\Delta R_{t-1})^{-1}$ 的突变；$L_{\text{trans}} = \frac{1}{3}|\Delta T|$ 在训练早期限制基元过度运动，防止其移出边界；$L_{\text{back}} = \frac{1}{3}\|\Delta T_{\text{forw}} - \Delta T_{\text{backw}}\|$ 通过逆变形网络确保前向-后向一致性。各损失权重设置为 $\lambda_{\text{over}}=1$、$\lambda_{\text{parsi}}=0.1$、$\lambda_{\text{vol}}=0.05$、$\lambda_{\text{smooth}}=1$、$\lambda_{\text{trans}}=0.5$、$\lambda_{\text{back}}=0.01$。

\textbf{面向动态场景的自适应控制。}D-Prism 设计了针对动态基元的自适应控制策略，每 2000 次迭代执行以下三种操作：（1）克隆——监控每个基元上 3D 高斯的梯度，当梯度超过 $\sigma_g = 4 \times 10^{-5}$ 的高斯比例超过 $\sigma_p = 0.15$ 时触发克隆，使基元更好地占据物体区域；若基元尺度过大则先缩小再克隆，否则在原位克隆相同基元。（2）合并——通过计算基元间多时间步的平均重叠率构建图，每个基元仅向其最高重叠邻居（且重叠率超过阈值 $\sigma_o$）添加边，通过连通分量定义组；组内先剪枝移除体积小于最大体积 $1/3$ 或与组内其他基元重叠超过 80\% 的基元，剩余基元合并为新基元（位置取体加权平均，其余属性取最大基元的值），再通过逆变形网络映射回规范空间。合并强度通过阈值 $\sigma_o$ 调节：低强度（$\sigma_o=1$）仅合并完全重叠的基元，高强度（$\sigma_o=0.3$）在 30\% 重叠时即合并，预设合适范围为 $\sigma_o \in [0.7, 0.8]$。（3）剪枝——移除不透明度低于 $\sigma_o = 0.3$ 的基元，同时基于体积排序移除异常小的基元（相邻体积值出现大于 10 倍跳跃时，移除跳跃下界以下的所有基元）。

\textbf{训练后精调阶段。}主训练完成后，D-Prism 进入精调阶段以进一步提升运动重建精度和外观质量。此阶段降低所有学习率（基元参数学习率从 0.001 缩小至 0.01 倍），冻结所有基元的规范空间 $T$ 和 $R$ 参数，同时解冻绑定 3D 高斯的所有参数——包括此前固定的重心坐标（高斯中心仍约束在基元表面上），基元数量也固定不变。主训练和精调各运行 60k 次迭代，在单块 NVIDIA 3090 GPU 上完成。

\textbf{结构化运动追踪精度指标。}为准确评估动态几何的运动建模质量，D-Prism 提出了结构化运动追踪精度（structured motion tracking accuracy）这一新指标。传统动态重建指标（如逐帧 Chamfer Distance 和 Earth Mover's Distance）仅比较单帧网格，无法反映预测的几何部件是否与真实部件以相同方式运动。新指标首先在规范网格上采样 $N$ 个点并计算其在所有时间戳的位置作为真实追踪值，然后使用待评估方法预测这 $N$ 个点的运动，最后通过 3D End-Point Error（EPE）和阈值内点比例 $\sigma_{3D}^{0.05}$、$\sigma_{3D}^{0.10}$ 量化运动相似度，直接反映方法对动态几何运动的建模能力。

\textbf{实验结果。}D-Prism 在自建 Dynamic Primitive Dataset（含 6 个结构化物体：宝箱、门、钳子、折叠椅、太阳镜、魔方）和 D-NeRF 数据集上验证。在运动追踪精度方面，D-Prism 在大规模、长时间运动（如魔方部件 360 度旋转、宝箱盖 180 度旋转）场景下显著优于 DG-Mesh、Ub4d 等基于几何的方法以及 MovingParts、Shape of Motion、SP-GS 等 NVS 追踪方法。在传统动态重建指标 CD$_d$ 和 EMD$_d$ 上也取得了强竞争力。在 D-NeRF 数据集的人形场景上，D-Prism 在提供结构化几何表示的同时达到了与专用 NVS 方法相当的 PSNR、SSIM、LPIPS 渲染质量。消融实验验证了三个核心设计：变形网络相比直接优化基元姿态大幅提升运动精度；合并操作以适当强度有效增强基元集的紧凑性和可解释性；克隆操作在复杂形状场景中显著提升基元表达力和渲染质量。

\textbf{下游应用。}D-Prism 的结构化表示天然支持多项下游应用：获取结构化表示后，可轻松进行部件级编辑（添加、删除、缩放特定组件）；可在 Blender 等软件中对基元进行关节动画，赋予物体全新运动。这使得 D-Prism 特别适用于关节物体（如机器人、家具、魔方）的动态建模与交互式操控。

\subsubsection{分析与概率方法}

本路线介绍从概率推断视角理解动态高斯泼溅的方法，包括变分推断、不确定性量化和理论分析。

前面的方法都是确定性的——给定输入，输出唯一的重建结果。本节转向概率视角：将动态重建建模为概率推断问题，使方法能够量化不确定性、从数据中学习运动先验，并进行时间外推。以下工作从概率或理论分析的角度理解动态高斯泼溅，为方法的诊断和改进提供了新视角。

\paragraph{变分高斯过程概率 4D 高斯}

GP-4DGS（2026.04）~\cite{2604_02915}（Probabilistic 4D Gaussian Splatting via Variational Gaussian Processes）首次将高斯过程（Gaussian Processes, GP）集成到 4DGS 中，实现动态场景的原则性概率建模。已有 4DGS 方法将运动视为确定性优化问题，使用手工设计的运动先验（多项式变形、刚性约束等）统一施加于所有基元，不从数据中学习。当基元观测不足或被遮挡时，这些固定先验变得不恰当或过于限制性，且缺乏原理性的不确定性估计和时间外推机制。GP-4DGS 的核心洞察是将基元变形表示为 GP 后验，使框架能够根据观测置信度自动调整正则化强度，并自然处理稀疏或未观测区域。

\textbf{概率变形建模。}对于 4D 输入 $\mathbf{x} = (\mathbf{p}, t) \in \mathbb{R}^4$（$\mathbf{p}$ 为规范空间 3D 位置，$t$ 为目标时间），GP 输出 $d=9$ 维变形向量（3 维平移变形 + 6 维连续旋转表示（2019）~\cite{zhou2019continuity}），每个输出维度 $f_i(\mathbf{x}) \sim \mathcal{GP}(m_i(\mathbf{x}), k_i(\mathbf{x}, \mathbf{x}'))$ 独立建模为一个 GP。这使得任意输入点的预测服从高斯分布 $f_i(\mathbf{x}) \sim \mathcal{N}(\mu_i(\mathbf{x}), \sigma_i^2(\mathbf{x}))$，均值和方差均由 GP 后验导出，为后续的不确定性量化奠定了概率基础。

\textbf{空间-时间复合核设计。}核函数定义 GP 输入之间的相关性结构，决定模型如何从已观测区域泛化至未观测区域。由于输入空间的 3 个空间维度 $p_x, p_y, p_z$ 与时间维度 $t$ 具有根本不同的相关性特征，GP-4DGS 采用解耦的复合核：

\begin{equation}
k_i(\mathbf{x}, \mathbf{x}') = k_i^{\text{spatial}}(\mathbf{p}, \mathbf{p}') + k_i^{\text{temporal}}(\mathbf{x}, \mathbf{x}')
\end{equation}

对于空间分量，GP-4DGS 选用 Mat\'ern 核 $k_i^{\text{spatial}}(\mathbf{p}, \mathbf{p}') = \sigma_{s,i}^2 \frac{2^{1-\nu_i}}{\Gamma(\nu_i)} r_{s,i}^{\nu_i} K_{\nu_i}(r_{s,i})$。其中 $r_{s,i} = \sqrt{2\nu_i \sum_{j \in \{x,y,z\}} (p_j - p_j')^2 / l_{s,i,j}^2}$ 是经过各向异性缩放后的距离度量，编码了"规范空间中相邻基元变形相似"的先验知识。之所以选择 Mat\'ern 核而非 RBF 核，是为了更好地处理动态场景中存在空间不连续物体时的建模需求。时间分量则采用各轴 Mat\'ern 核与周期核的加权组合 $k_i^{\text{temporal}}(\mathbf{x}, \mathbf{x}') = \sum_{j \in \{x,y,z\}} k_{i,j}(\mathbf{p}_j, \mathbf{p}_j') \cdot k_{i,j}^{\text{periodic}}(t, t')$，其中周期核 $k_{i,j}^{\text{periodic}}(t, t') = \sigma_{i,j}^2 \exp(-\frac{2\sin^2(\pi|t-t'|/p_{i,j})}{l_{i,j}^2})$ 为时间外推提供了强归纳偏置。此外，GP-4DGS 还设计了周期均值函数 $m_i(\mathbf{x}) = c_i + A_i \sin(2\pi t / T_i + \phi_i)$，其中 $c_i$ 为可学习的均值偏移，$A_i$ 为振幅，$T_i$ 为周期，$\phi_i$ 为相移。消融实验表明周期均值对长程外推至关重要（有效捕捉周期性运动），而周期核则在短期外推和非周期运动上提供更通用的性能——二者互补：周期均值提供整体运动趋势，周期核捕获局部时间模式。无任何周期结构时，GP 在观测时间范围边界处预测完全停止。

\textbf{变分推断与诱导点。}由于精确 GP 推断需要 $O(N^3)$ 的核矩阵求逆（$N$ 为高斯数量，通常数万），GP-4DGS 采用变分 GP（2009）~\cite{titsias2009variational}：用 $M \ll N$ 个诱导点 $\mathbf{Z} = \{\mathbf{z}_m\}_{m=1}^M$（$\mathbf{z}_m \in \mathbb{R}^4$）近似完整 GP，只需计算诱导点间的核矩阵 $K_{ZZ}^{(i)} \in \mathbb{R}^{M \times M}$ 及基元与诱导点之间的核矩阵 $K_{XZ}^{(i)} \in \mathbb{R}^{N \times M}$，将计算复杂度降至 $O(NM^2 + M^3)$。诱导点的初始化采用两阶段策略：空间位置通过 Chronos（2024）~\cite{ansari2024chronos} 时序基础模型提取轨迹特征——对每个基元 $k$ 的轨迹 $\mu_k \in \mathbb{R}^{T \times 3}$，按 $x, y, z$ 轴独立提取 256 维嵌入并拼接为 $\mathbf{F}_k = [\mathbf{f}_{k,x}; \mathbf{f}_{k,y}; \mathbf{f}_{k,z}] \in \mathbb{R}^{768}$（捕获周期模式和非线性变化等复杂运动特征），再通过 $k$-means 聚类选择 $M_{\text{spatial}}$ 个代表性规范位置；时间位置在观测范围内均匀采样 $M_{\text{time}}$ 个点。最终诱导点数量为笛卡尔积 $M = M_{\text{spatial}} \times M_{\text{time}}$。相比速度差异等启发式特征，Chronos 特征能编码长期时序依赖和高阶动力学，消融实验显示其一致获得更高的 ELBO。训练通过 ELBO 目标联合优化核超参数（长度尺度、信号方差、周期）、诱导点位置和变分参数 $(\mathbf{m}_i, \mathbf{S}_i)$，其中 $\mathbf{m}_i \in \mathbb{R}^M$ 和 $\mathbf{S}_i \in \mathbb{R}^{M \times M}$ 参数化诱导点上函数值的变分后验 $q(\mathbf{u}_i) = \mathcal{N}(\mathbf{m}_i, \mathbf{S}_i)$。ELBO 形式为 $\mathcal{L}_{\text{ELBO}} = \sum_{i=1}^d \left[\mathbb{E}_{q(f_i)}[\log p(\mathcal{D}_i | f_i)] - \text{KL}(q(\mathbf{u}_i) \| p(\mathbf{u}_i))\right]$，其中期望项鼓励拟合观测变形，KL 项将后验正则化向 GP 先验 $p(\mathbf{u}_i) = \mathcal{N}(\mathbf{0}, K_{ZZ}^{(i)})$ 收敛。推理时，任意查询点 $\mathbf{x}^*$ 的变形均值和方差分别为：

\begin{equation}
\mu_i^* = \mathbf{k}_*^{(i)\top} (K_{ZZ}^{(i)})^{-1} \mathbf{m}_i, \quad {\sigma_i^*}^2 = k_i^* - \mathbf{k}_*^{(i)\top} \mathbf{\Lambda}_i \mathbf{k}_*^{(i)}
\end{equation}

其中 $\mathbf{k}_*^{(i)} = k_i(\mathbf{Z}, \mathbf{x}^*) \in \mathbb{R}^M$ 是诱导点与查询点的互协方差向量，$k_i^* = k_i(\mathbf{x}^*, \mathbf{x}^*)$ 是查询点的核方差，$\mathbf{\Lambda}_i = (K_{ZZ}^{(i)})^{-1} - (K_{ZZ}^{(i)})^{-1}\mathbf{S}_i(K_{ZZ}^{(i)})^{-1}$。每个查询的计算开销仅 $O(M)$，无需访问全部训练基元。

\textbf{GP-GS 交替优化。}GP-4DGS 设计了 GP 训练与 4DGS 优化的交替策略，形成逐步精化两个表示的反馈循环。\emph{阶段一（GP 训练）：}通过累积渲染贡献衡量每个高斯基元的信息量 $C_k = \sum_{I^T} \sum_{r \in I} \alpha_{k,t}(r)$（$\alpha_{k,t}(r)$ 为基元 $k$ 在像素 $r$ 处的 $\alpha$-混合 权重），选择置信度高于阈值 $\tau_C$ 的基元子集 $\mathcal{D}_C$ 进行 GP 训练。由于规范空间坐标 $\mathbf{p}$ 本身也是优化参数而非精确位置，训练时对每个空间维度注入高斯噪声 $\tilde{\mathbf{x}} = (p_x + \epsilon_x, p_y + \epsilon_y, p_z + \epsilon_z, t)$（$\epsilon \sim \mathcal{N}(0, 0.02)$），作为正则化使 GP 能够处理规范位置的不确定性，产生更鲁棒的预测。\emph{阶段二（GS 优化）：}在 $\mathcal{D}_C$ 上训练完成后，对所有基元运行 GP 推断，得到每个基元 $k$ 在时间 $t$ 的后验均值 $\hat{\mathbf{y}}_{(k,t)}^*$ 作为伪引导信号。为平衡计算效率和引导信号的新鲜度，GP 预测每 $N_{\text{GP}} = 2000$ 次迭代缓存一次。GP 引导损失定义为：

\begin{equation}
L_{\text{GP}} = \frac{1}{NT} \sum_{k=1}^{N} \sum_{t=1}^{T} \gamma_{(k,t)} \| \mathbf{y}_{(k,t)} - \hat{\mathbf{y}}_{(k,t)}^* \|^2
\end{equation}

其中 $\gamma_{(k,t)} = \mathbb{1}(\|\mathbf{y}_{(k,t)} - \hat{\mathbf{y}}_{(k,t)}^*\| > \tau)$ 仅选择偏离 GP 运动预测的基元，阈值 $\tau$ 在优化过程中从 $\tau_{\text{start}}$ 退火至 $\tau_{\text{end}}$，随两个表示逐步收敛而收紧约束。总训练损失为 $L_{\text{total}} = L_{\text{recon}} + \lambda_{\text{GP}} L_{\text{GP}}$，其中 $L_{\text{recon}}$ 包含光度损失、D-SSIM、光流和平滑正则项，$\lambda_{\text{GP}} = 0.1$。

\textbf{不确定性量化。}GP 天然提供方差估计，平移变形的不确定性直接来自预测方差的前三个维度。旋转的 6D 表示到旋转矩阵的 Gram-Schmidt 正交化是非线性操作，无法直接传播方差，因此采用蒙特卡洛采样近似最终位置的不确定性：$U_{k,t} = \text{Var}(\{p_{k,t}^{(s)}\}_{s=1}^S)$，其中 $p_{k,t}^{(s)}$ 是第 $s$ 次采样变形后的位置。运动不确定性图通过将 $U_{k,t}$ 投影至目标视角渲染：$\mathcal{U}(r) = \sum_{k=1}^{N} U_{k,t} \cdot \alpha_{k,t}(r)$，与颜色渲染类似。在不确定性评估中，GP-4DGS 采用稀疏化误差下面积（AUSE）衡量估计不确定性与实际重建误差的对齐程度，在 DyCheck 上一致优于随机基线和 UA-4DGS（2024）~\cite{kim20244d}，且在高质量帧（如 top-20 和 top-40）上的优势更加显著，表明模型能够识别重建良好区域中的细微残差误差。

\textbf{时间外推与实验。}GP 基运动推理的关键优势是能够预测训练序列之外的未来运动。给定规范位置和未来时间戳 $t_f$，直接以 $\mathbf{x}_f^* = (\mathbf{p}, t_f)$ 查询训练好的 GP 即可获得该时刻的变形预测，无需额外训练或架构修改。在 DyCheck 上，训练时分别保留最后 5 帧和 15 帧用于评估外推质量，GP-4DGS 显著优于朴素线性外推基线，尤其在周期性场景中表现突出——周期核和周期均值协同捕获了循环动力学。在 DyCheck 基准的动态新视角合成评估中，GP-4DGS 在 mPSNR、mSSIM 和 mLPIPS 上一致超越所有基线，且在 Challenging 子集（使用扩散方法的子采样帧，具有更少的视角重叠）上的性能差距进一步扩大，证明 GP 优化通过将高置信区域的时空相关性传播至未观测区域，有效缓解了稀疏观测带来的伪影。在 DAVIS 数据集上的极端视角外推测试中，GP-4DGS 比基线更忠实地保持了几何和清晰边缘，验证了 GP-GS 优化在极端视角偏移下的结构完整性。

\textbf{贡献总结。}这种概率框架同时赋予了三个现有 4DGS 方法不具备的能力：不确定性量化（识别重建不可靠的区域）、时间外推（超越观测帧预测未来运动）和从数据学习的运动先验（通过 GP 核函数自动适应观测模式）。GP-GS 交替策略将 GP 的结构化运动先验与 4DGS 的渲染能力结合，置信度加权的训练数据选择和噪声注入正则化确保了概率建模的鲁棒性。

\paragraph{潜在 ODE 动态外推与长序列生成}

ODE-GS（2025.06）~\cite{2506_05480}（Latent ODEs for Dynamic Scene Extrapolation with 3D Gaussian Splatting，ICLR 2026）由 Yale、TU Delft、LSU 和 RTX 联合提出，引入潜在神经常微分方程实现动态 3D 场景的未来外推。已有动态场景重建方法（如 Deformable 3D Gaussians、4DGS）依赖时间条件的变形网络，本质上是在观测时间窗口内做插值，当时间戳超出训练分布时会产生 OOD 失效。ODE-GS 的核心思路是将动态场景外推重新表述为序列到序列预测问题，从而消除对显式时间戳的依赖。

\textbf{解耦的两阶段架构。}方法分两阶段：首先训练一个标准的规范帧加变形网络插值模型，在观测窗口内生成高斯参数的密集时序轨迹，训练完成后冻结该模型作为数据生成器；然后在生成的轨迹上训练 Transformer-Latent ODE 外推模块。形式化地，连续时间渲染算子分解为 $\mathcal{F}(t,V) = \mathcal{R}(\mathcal{G}(t), V)$，其中高斯状态在观测窗口内由规范高斯加变形网络生成（$\mathcal{G}(t) = \overline{\mathcal{G}} + \mathcal{D}_\omega(t, \overline{\mathcal{G}})$），在未来时刻由外推模块生成（$\mathcal{G}(t) = \mathcal{E}_\phi(\gamma, t)$）。每个规范高斯 $G_k = (\mu_k, q_k, s_k, c_k, \alpha_k)$ 中仅位置 $\mu_k \in \mathbb{R}^3$、旋转 $q_k \in \mathbb{R}^4$（四元数）和缩放 $s_k \in \mathbb{R}^3$ 是时间相关的，不透明度和球谐系数跨时间保持一致。

\textbf{Transformer-Latent ODE 外推模块。}外推模块 $\mathcal{E}_\phi$ 将时序预测建模为序列到序列映射。每个高斯基元 $k$ 在上下文窗口 $N_c$ 内均匀采样的参数序列 $\{G_k(t_j)\}_{j=1}^{N_c}$（包含位置、旋转、缩放共 10 维），加入正弦位置编码后经 Transformer 编码器 $\mathcal{F}_\phi: \mathbb{R}^{N_c \times 10} \to \mathbb{R}^d$ 映射为潜在状态 $\mathbf{z}(t_0) \in \mathbb{R}^d$，再由神经 ODE（参数化 MLP）建模其连续演化：

\begin{equation}
\dot{\mathbf{z}} = \frac{\mathrm{d}\mathbf{z}}{\mathrm{d}t} = f_\theta(\mathbf{z}(t))
\end{equation}

通过数值积分（ODE solver）产生任意未来时刻 $t > t_{\max}$ 的连续潜在轨迹 $\mathbf{z}(t)$，最终由解码器 $\delta_\psi: \mathbb{R}^d \to \mathbb{R}^{10}$ 映射回高斯参数 $\hat{G}_k(t) = \delta_\psi(\mathbf{z}(t))$。由于外推过程不涉及时间戳嵌入，未见时间戳不再是 OOD 问题——外推能力完全来自潜在空间中的连续动力学演化。

\textbf{动态轨迹采样。}ODE-GS 设计了一种动态轨迹采样策略以提升外推泛化性：不同于传统方法中采样轨迹始终占据相同的时间跨度，该方法从预训练插值模型生成的连续轨迹中提取可变长度的观测前缀（固定间隔采样以保证输入维度一致）和未来后缀（起始时间不同则时间跨度不同）。训练集构建为所有高斯基元在所有可能前缀-后缀分割上的并集，迫使模型在统一训练过程中同时学习短期和长期预测，从而增强超越观测窗口的泛化能力。

\textbf{自适应加权正则化。}训练目标包含外推损失 $\mathcal{L}_e = \frac{1}{N_e}\sum_{j=1}^{N_e} \|\hat{G}_k(t_j) - G_k(t_j)\|_1$ 和两项互补正则化。潜在正则化惩罚 ODE 向量场的高频振荡——即潜在加速度：

\begin{equation}
\mathcal{R}_{\text{latent}} = \frac{1}{N_e-1}\sum_{j=1}^{N_e-1} \left\| \frac{f_\theta(\mathbf{z}(t_{j+1})) - f_\theta(\mathbf{z}(t_j))}{\Delta t_j} \right\|_2^2
\end{equation}

轨迹正则化直接在 3D 空间中惩罚高斯位置加速度：

\begin{equation}
\mathcal{R}_{\text{traj}} = \frac{1}{MN_e}\sum_{k=1}^{M}\sum_{j=1}^{N_e-2} \left\| \frac{v_k(t_{j+1}) - v_k(t_j)}{\Delta t_j} \right\|_2^2
\end{equation}

其中速度 $v_k(t_j) = \frac{\mu_k(t_{j+1}) - \mu_k(t_j)}{\Delta t_j}$。自适应加权机制以 EMA 监测轨迹预测损失的收敛状态，在训练早期降低正则化权重以允许学习有意义的动力学，随着外推损失下降逐步增大正则化权重，在众多合理解中逐步偏向更平滑的轨迹。最终训练损失为 $\mathcal{L} = \mathcal{L}_e + s_t(\lambda_{\text{latent}}\mathcal{R}_{\text{latent}} + \lambda_{\text{traj}}\mathcal{R}_{\text{traj}})$，其中 $s_t$ 为自适应权重。

\textbf{实验表现。}ODE-GS 在 D-NeRF、NVFi 和 HyperNeRF 三个基准上达到最优外推性能。在 D-NeRF 上平均 PSNR 27.30、SSIM 0.9497、LPIPS 0.0467，相对此前最优外推方法 GaussianPrediction 提升 18.6\%，在运动平滑的场景（Mutant +10 dB、Standup +7 dB）优势尤为显著。在 NVFi 上平均 PSNR 33.43、SSIM 0.9471、LPIPS 0.0603，相对 NVFi 方法提升 20\%，相对 GaussianPrediction 提升 39.6\%，在遮挡密集场景（factory、darkroom）中 LPIPS 降低超 40\%。在真实 HyperNeRF 场景上同样持续优于基线，LPIPS 在 split-cookie、slice-banana、cut-lemon 上均为最优。消融实验显示：去除 ODE（纯自回归 Transformer）使 LPIPS 几乎翻倍（0.113 vs.\ 0.060），验证了连续 ODE 动力学相对离散自回归的平滑性优势；去除正则化使 PSNR 下降 0.53 dB；去除动态采样使 PSNR 下降 2.08 dB。性能退化分析表明 ODE-GS 的 SSIM/LPIPS 退化速率远慢于基线，随外推时间增长仍能维持合理的运动和外观。ODE-GS 相对最佳外推基线在三组数据集上平均提升 PSNR 21.4\%、SSIM 7.4\%、LPIPS 30.5\%。

\paragraph{运动感知高斯未来帧预测}

MoGaF（2026.02）~\cite{2602_21668}（Motion Group-aware Gaussian Forecasting，CVPR 2026）解决动态场景的长期未来预测问题。已有逐高斯预测方法（如 GSPred）中高斯独立运动导致空间不连贯的误差随时间累积，而短时预测器（如 ODE-GS）在长程推演中产生冻结或坍缩轨迹。MoGaF 的设计动机是：只有在物体级别建模运动，才能实现空间连贯且时间稳定的长期外推。方法构建于 Shape-of-Motion（SoM）的 4DGS 表示之上，其中每个高斯的运动由 $B$ 个共享运动基 $\{T_{c \to t}^{(b)}\}_{b=1}^{B}$ 的线性混合表示：$T_{c \to t} = \sum_{b=1}^{B} w^{(b)} T_{c \to t}^{(b)}$，$\|w\|=1$，运动基与运动系数 $w$ 在优化中联合更新。

\textbf{运动感知高斯分组。}方法分三阶段，第一阶段将高斯聚类为具有连贯运动的物体组。利用 Grounding DINO 等 2D 分割模型同时获取 $K$ 个物体掩码及其刚性/非刚性标签 $\tau^{(k)} \in \{0, 1\}$（1 为刚性、0 为非刚性）。直接将静态 3DGS 分组方法（如 Gaga）扩展到动态场景——仅用变形后的 2D 投影重叠度做帧间合并——在遮挡和运动导致的空间漂移下容易误分组。MoGaF 改用迭代区域增长策略，交替执行关键帧种子选取和特征空间扩展：在每个关键帧上选取渲染贡献最大的前向高斯作为可靠种子，然后以每个高斯的时空特征 $\mathbf{f}_g = [\mu_{c,g},\, w'_g]$（融合规范空间坐标 $\mu_{c,g}$ 和 PCA 降维的运动系数 $w'_g$）为度量，向满足 $|\mathbf{f}_g - \mathbf{f}_{g'}| < \epsilon_k$ 的近邻扩展组边界，其中自适应阈值 $\epsilon_k = \alpha \cdot \text{mean}_{g \in G^{(k)}}[\text{KNN}_K(\mathbf{f}_g)]$ 根据组内高斯的特征分布动态调整。迭代交替使分组逐步融合关键帧的空间约束和运动系数的时序一致性，最终产生 $K$ 个运动组 $\mathcal{M}^{(k)} = (G^{(k)}, \tau^{(k)})$。未分配的高斯通过 KNN 投票归入最近邻主导的组。

\textbf{分组约束优化。}第二阶段利用分组信息对 4DGS 表示做运动感知的精调。对刚性组（$\tau^{(k)}=1$），组内所有高斯共享单一 SE(3) 变换 $\Phi_t^{(k)} = [R_{c \to t}^{(k)} \mid t_{c \to t}^{(k)}]$，将规范空间位置映射到变形空间：$\Phi_t^{(k)}(g) = R_{c \to t}^{(k)} \mu_{c,g} + t_{c \to t}^{(k)}$。刚性约束损失衡量逐高斯运动与组级刚性变换的偏差：

\begin{equation*}
\mathcal{L}_{\text{rigid}}^{(k)} = \sum_{t} \sum_{g \in G^{(k)}} \|\mu_{t,g} - \Phi_t^{(k)}(g)\|_2^2
\end{equation*}

对非刚性组（$\tau^{(k)}=0$），在规范空间中对每个高斯的可学习运动系数 $w_g \in \mathbb{R}^B$ 施加空间邻域平滑正则：

\begin{equation*}
\mathcal{L}_{\text{nr}}^{(k)} = \sum_{g \in G^{(k)}} \sum_{g' \in \text{NN}(g)} \|w_g - w_{g'}\|_2^2
\end{equation*}

其中 NN($g$) 为规范空间中固定半径邻域内的空间邻居。整体运动目标按刚性标志自适应组合：

\begin{equation*}
\mathcal{L}_{\text{motion}} = \sum_{k=1}^{K} \big[\tau^{(k)} \mathcal{L}_{\text{rigid}}^{(k)} + (1 - \tau^{(k)}) \mathcal{L}_{\text{nr}}^{(k)}\big]
\end{equation*}

\textbf{分组运动预测。}第三阶段为每个运动组 $G^{(k)}$ 训练独立的浅层 Transformer 编码器作为运动预测器（单层、8 头注意力、32 维特征嵌入、64 维前馈扩展）。预测器以优化后的 SE(3) 运动序列 $\{T_{t,g}\}_{t=0}^{T}$（其中 $T_{t,g} = [R_{t,g} \mid \mu_{t,g}]$）为输入，自回归推演未来运动 $\{\hat{T}_{t,g}\}_{t>T}$，每次用最近 $T-1$ 帧（含预测帧）生成下一帧。训练采用段级掩码策略（segment-level masking）：随机遮蔽连续时间段的运动输入，迫使模型从可见上下文推断缺失动力学，掩码比例在训练中逐步退火以匹配推理条件。每个预测器的训练目标包含运动重建损失和加速度正则化：

\begin{equation*}
\mathcal{L}_{\text{pred}}^{(k)} = \frac{1}{|G^{(k)}|}\sum_{g \in G^{(k)}} \|T_{T,g} - \hat{T}_{T,g}\|_2^2, \quad
\mathcal{L}_{\text{acc}}^{(k)} = \frac{1}{|G^{(k)}|}\sum_{g \in G^{(k)}} \|\hat{\mu}_{T,g} - 2\mu_{T-1,g} + \mu_{T-2,g}\|_2^2
\end{equation*}
\begin{equation*}
\mathcal{L}_{\text{group}}^{(k)} = \mathcal{L}_{\text{pred}}^{(k)} + \lambda_{\text{acc}} \mathcal{L}_{\text{acc}}^{(k)}
\end{equation*}

各组的预测器独立训练，解耦异构运动动力学，确保同一物体的所有高斯遵循一致的运动模式，避免了逐高斯独立预测的空间不连贯问题。

\textbf{实验验证。}在 iPhone 真实数据集上（80\% 观测 / 20\% 预测），MoGaF 平均 mPSNR 15.58、mSSIM 0.6395、mLPIPS 0.4227，相对 GSPred-SoM 基线分别提升 +0.59 dB、-0.001、-0.0255；在更具挑战的 60\% 观测 / 40\% 预测设定下仍保持 mPSNR 15.51、mLPIPS 0.4245，长程预测退化极小。在 D-NeRF 合成数据集上（60\% 观测 / 40\% 预测），MoGaF 平均 PSNR 23.37、SSIM 0.9147、LPIPS 0.0746，相比 GSPred 的 21.78 / 0.9011 / 0.0919 全面领先，其中 Lego 场景提升最为显著（PSNR 21.61 vs 12.65，+8.96 dB）。消融实验证实分组策略的关键作用：移除分组后 3D 追踪 EPE 从 0.236 恶化至 0.296，$\delta_{3D}^{0.10}$ 从 44.8\% 降至 35.6\%；移除段级掩码训练导致 D-NeRF 上 PSNR 下降 1.19 dB（25.87 vs 24.68）。定性分析显示 MoGaF 在快速运动的 apple 场景中准确预测手的复杂运动轨迹，而基线方法产生手部冻结或几何崩塌。

\paragraph{拉普拉斯分析约束动力学建模}

Laplacian Analysis Meets Dynamics Modelling（2025.08）~\cite{2508_04966} 将拉普拉斯谱分析引入动态 3DGS，旨在解决动态场景重建中的固有谱冲突：低秩分解（如 HexPlane 的平面外积假设）导致运动细节过度平滑，高维网格采样则因重叠区域的特征碰撞产生渲染伪影，而标准 3DGS 的固定阈值分裂策略忽略了动态区域对基元密度和各向异性的差异化需求。方法的核心洞察在于：动态场景具有各向异性的时空采样特性——刚体部件呈现低频一致性运动，非刚性区域则需要高频变形建模，不同频率的运动需要区别对待。为此方法提出混合显式-隐式编码框架，包含三项关键创新。

\textbf{谱感知拉普拉斯编码模块。}方法首先通过多尺度 Hash 编码对 4D 时空 $(x,y,z,t)$ 进行各向异性多分辨率分解，将传统 InstantNGP 的 Hash 编码扩展到四个三维子空间：

\begin{equation*}
\mathcal{H}^l = \{\mathcal{H}_{xyz}^l,\; \mathcal{H}_{xyt}^l,\; \mathcal{H}_{yzt}^l,\; \mathcal{H}_{xzt}^l\},
\quad l \in \{1, \ldots, L\}
\end{equation*}

每个层级 $l$ 均通过几何级数计算维度特定的分辨率参数。在此基础上，拉普拉斯运动模块将时间序列分解为可学习的频率分量，从拉普拉斯级数分解出发：

\begin{equation*}
L(t) = \sum_{k=-K}^{K} \mathbf{c}_k \cdot e^{2\pi i kt/T}
\end{equation*}

经 Euler 公式展开并简化为实值可学习形式：

\begin{equation*}
L(t) = \sum_{k=0}^{K-1} \bigl[\alpha_k \cos(2\pi f_k t)
      + \beta_k \sin(2\pi f_k t)\bigr]
\end{equation*}

其中系数 $(\alpha_k, \beta_k)$ 和频率 $f_k$ 均为可学习参数。在物理含义上，低频分量负责捕捉整体刚体运动，而高频分量则保留局部变形细节。频率 $f_k$ 的梯度通过链式法则自动计算：$\frac{\partial \mathcal{L}}{\partial f_k} = \frac{1}{\sigma_k^2}\sum_t \left(\frac{\partial \mathcal{L}}{\partial L(t)} \cdot t \cdot [-\alpha_k \sin(2\pi f_k t) + \beta_k \cos(2\pi f_k t)]\right)$，其中 $\sigma_k$ 为时间方差，该机制可自动平衡频率保持与运动稳定性。拉普拉斯时序特征与 Hash 空间特征通过注意力机制融合：$\mathcal{A}_L(t) = L(t) \cdot \text{MLP}(\mathcal{H}_s)$，将自适应频率分析与空间上下文结合。

\textbf{增强高斯动态属性与自适应正则化。}方法为每个高斯 $\mathcal{G}_i$ 引入可学习的动态属性向量 $\mathbf{d}_i \in \mathbb{R}^{D_d}$，用于显式编码高斯的潜在时变特征。动态属性与 Hash 时序特征拼接形成增强表示 $\tilde{\mathcal{H}_t} = \text{Concatenate}(\mathbf{d}_i, \mathcal{H}_t)$，再通过注意力机制聚合时空信息 $\mathcal{A}_h(t) = \tilde{\mathcal{H}_t} \cdot \mathcal{H}_s$。自适应动态正则化采用数据驱动的选择策略：计算每个高斯的动态偏差 $d_i = \|\mathbf{d}_i - \bar{\mathbf{d}}\|_2$，以均值加一倍标准差为阈值标记"离群"高斯 $\text{mask}_i = \mathbb{1}[d_i > \mu_{\text{dist}} + \sigma_{\text{dist}}]$，仅对离群高斯施加额外梯度惩罚：

\begin{equation*}
\mathcal{L}_{\text{dy}} = \frac{1}{N}\sum_{i=1}^{N} \text{mask}_i \cdot \|\mathbf{d}_i - \bar{\mathbf{d}}\|_2^2
\end{equation*}

该设计将优化资源集中于高动态区域，同时避免对静态或低动态高斯的冗余约束。

\textbf{自适应高斯分裂策略。}针对标准 3DGS 固定阈值分裂在动态区域的不足，方法提出 KDTree 引导的谱分析分裂策略。首先通过 KDTree 查询每个高斯 $\mathcal{G}_i$ 的 $K$ 近邻，计算邻域内的协方差差异 $\Delta\Sigma_{ij} = \|\Sigma_i - \Sigma_j\|_2$ 以检测运动特征显著不同的区域。然而仅基于形状差异的硬阈值分裂会遗漏邻域内仍对动态建模有贡献的高斯，因此进一步引入 KL 散度度量邻域分布的均匀性 $D_{\text{KL}}(P \| Q) = \sum_{k=1}^{K} P(k) \log \frac{P(k)}{Q(k)}$，动态调整分裂阈值：

\begin{equation*}
\tau = \Delta\Sigma + D_{\text{KL}} \cdot \tau_{\text{base}}
\end{equation*}

该机制根据观察到的运动模式自适应调整模型复杂度，在捕获运动细节的同时维持计算效率。

\textbf{多尺度拉普拉斯金字塔监督。}方法在渲染图像和真实图像的拉普拉斯金字塔分解之间施加多尺度一致性约束：

\begin{equation*}
\mathcal{L}_{\text{lap}} = \sum_{l=1}^{L} \lambda_l \|\mathcal{L}_l(I_{\text{render}}) - \mathcal{L}_l(I_{\text{gt}})\|_1
\end{equation*}

权重 $\lambda_l$ 按指数衰减以强调细节层，确保不同频段的重建一致性。完整训练损失包含四项：$\mathcal{L} = \mathcal{L}_{\text{orig}} + \lambda_{\text{NCC}}\mathcal{L}_{\text{NCC}} + \lambda_{\text{lap}}\mathcal{L}_{\text{lap}} + \lambda_{\text{dy}}\mathcal{L}_{\text{dy}}$，其中 $\mathcal{L}_{\text{orig}}$ 为 3DGS 原始损失（$\mathcal{L}_1$ + SSIM），$\mathcal{L}_{\text{NCC}}$ 为归一化互相关损失以增强局部亮度不变性下的对齐精度。

\textbf{实验表现。}该方法在三个数据集上进行了全面评估：D-NeRF 合成数据集上达到 PSNR 42.17\,dB、SSIM 0.994、LPIPS 0.007，超越 Grid4D（41.99/0.994/0.008）和 SC-GS（41.66/0.993/0.009）。HyperNeRF 真实数据集上平均 PSNR 25.83\,dB、SSIM 0.720，优于 Grid4D（25.46/0.715）和 ED3DGS（25.43/0.697）。Neu3D 多视角数据集上达到 PSNR 32.12\,dB、SSIM 0.944、LPIPS 0.134，超越 ED3DGS（31.92/0.943/0.139）和 Spacetime（31.46/0.944/0.142）。消融实验表明：移除拉普拉斯模块后 PSNR 下降 0.48 dB，移除动态属性后 PSNR 下降 0.40 dB，移除自适应分裂策略后 PSNR 下降 0.16 dB 且高斯数量从 433k 增至 589k（减少 26\% 的高斯用量），移除拉普拉斯金字塔损失后 PSNR 下降 0.42 dB，验证了每个组件的贡献。

\subsubsection{其他变形与4D方法}

本路线介绍全显式动态表示、镜面动态场景、抗锯齿4D高斯等不属于前述类别但具有独特贡献的方法。

概率方法为 4D 重建提供了理论框架，但实际应用中还有许多特殊场景需要专门的处理策略。本节收集了不属于前述类别但具有独特贡献的方法，涵盖全显式动态表示、镜面动态场景、抗锯齿 4D 高斯等方向。

以下方法从不同角度补充了变形场和 4D 高斯表示的研究。

\paragraph{精细化 3D 高斯动态表示}

A Refined 3D Gaussian Representation for High-Quality Dynamic Scene Reconstruction（2024.05）~\cite{2405_17891} 提出 3DGS 在动态场景中的精细化策略。该方法沿袭 D-NeRF 的变形场范式：在规范空间中建立 3D 高斯表示，通过变形网络预测时变偏移量，将不变特征与时变特征解耦。方法观察到标准 3DGS 在动态场景中面临两大问题：点云初始化不足（COLMAP 在运动模糊区域难以提取匹配点）和密度控制不当（标准克隆/分裂在动态区域产生的伪影点无法被有效清除）。为此方法引入三项关键技术，形成``变形场 + Hash 编码颜色 + 去噪掩码 + 运动/静态约束''的完整框架。

\textbf{变形 Hash 高斯场。}方法的核心表示为变形 3D Hash 高斯场。变形场 MLP（$8 \times 256$ 层）以频率位置编码 $\gamma(\mathbf{p}) = (\sin(2^k \mathbf{p}), \cos(2^k \mathbf{p}))_{k=0}^{L-1}$ 和时间编码 $t$ 为输入，输出位置、旋转、尺度、颜色的偏移量：

\begin{equation}
(\Delta\mu, \Delta q, \Delta s, \Delta c) = F_\theta(\gamma(\mathbf{p}), \tau(t))
\end{equation}

其中 $F_\theta$ 为变形场 MLP，$\gamma(\mathbf{p})$ 为频率位置编码，$\tau(t)$ 为时间编码。在颜色表示方面，原版 3D-GS 使用三阶球谐函数表示每个高斯的颜色，每个 RGB 通道需要 16 个独立 32 位浮点数，颜色信息占总存储量的 81\%。方法借鉴 Instant-NGP（2022）~\cite{mueller2022instant} 的思路，用 Hash 编码（hash 表大小 $2^{20}$、分辨率范围 16--2048、每层特征维度 2）加微型 MLP（$64 \times 2$ 结构，输出维度 3）替代球谐函数。最终颜色表示为 Hash 编码颜色与变形场颜色偏移量的直接相加：

\begin{equation}
c_i = \phi_c(h(\mathbf{x}_i)) + \Delta c_i
\end{equation}

其中 $h(\cdot)$ 为 Hash 网格特征编码，$\phi_c$ 为微型 MLP 解码器。这种设计不仅加速了变形场与 Hash 编码之间的联合收敛，还将存储需求压缩至原始的 20\%。综合变形场和去噪掩码，最终的变形 3D Hash 高斯表示为：

\begin{equation}
G_i: \{\mu_i + \Delta\mu_i,\; q_i + \Delta q_i,\; M_i s_i + \text{sg}(M_i) s_i,\; c_i + \Delta c_i,\; M_i o_i\}
\end{equation}

\textbf{可学习去噪掩码。}标准 3DGS 的密度控制（基于梯度变化的克隆/分裂和低不透明度剔除）虽在静态场景中有效，但在动态场景中仍会产生大量噪声点，导致渲染伪影和存储浪费。方法提出可学习掩码 $M_i$ 并结合变形场消除噪声点。掩码的基本形式为连续 sigmoid 映射 $M_i = \sigma(m_i)$，将可学习参数 $m_i$ 映射到 $(0, 1)$。为提升去噪效果，方法参考紧凑 3D 高斯（2023.11）~\cite{2311_13681} 的设计，将连续掩码转化为二值掩码：

\begin{equation}
M_i = \sigma(m_i) + \mathbb{1}[\sigma(m_i) > \epsilon] - \sigma(m_i)
\end{equation}

其中 $\text{sg}(\cdot)$ 为停止梯度操作，$\epsilon$ 为预设阈值。通过该式将连续掩码二值化，虽可能导致渲染质量轻微下降，但能更彻底地移除无效点。二值掩码用于屏蔽噪声点的尺度 $s_i = M_i s_i + \text{sg}(M_i) s_i$ 和不透明度 $o_i = M_i o_i$。为训练掩码，方法引入基于滑动窗口期望值的去噪约束损失：设定尺度期望 $E(s)$ 和不透明度期望 $E(o)$（通过训练过程中的历史滑动窗口平均值计算），鼓励掩码学习识别并移除偏离期望的噪声点：

\begin{equation}
L_{dn} = \frac{1}{N}\sum_{i=1}^{N}\|E(s_i) - (M_i s_i + \text{sg}(M_i) s_i)\| + \frac{1}{N}\sum_{i=1}^{N}\|E(o_i) - M_i o_i\|
\end{equation}

此外附加掩码参数正则化 $L_m = \frac{1}{N}\sum_{i=1}^{N} m_i$ 以鼓励稀疏性。去噪约束 $L_{dn}$ 从第 5{,}000 次迭代开始加入训练。

\textbf{静态/运动一致性约束。}变形场学习所有高斯点的时变偏移量，但场景中存在大量静态点（时间不变），无差别优化会干扰变形场对真正动态区域的建模。方法通过 L1 范数阈值 $\|\Delta\mu_i\| < 0.1$ 将点分为静态和动态两类。对静态点，基于 $\|\Delta\mu_i\|$ 倒数计算归一化权重 $w_i = \frac{1/\|\Delta\mu_i\|}{\sum_{\|\Delta\mu_j\| < 0.1} 1/\|\Delta\mu_j\|}$（确保动态点不受约束影响），施加加权位移惩罚：

\begin{equation}
L_s = \sum_{i=1}^{N} w_i \|\Delta\mu_i\|, \quad \|\Delta\mu_i\| < \text{threshold}
\end{equation}

对动态点，方法基于先验假设——同一时刻大部分动态点应呈现相似运动趋势——将坐标偏移 $\Delta\mu$ 分解为三轴运动 $\Delta\mu_x, \Delta\mu_y, \Delta\mu_z$，每轴进一步按正/负方向分组（如 $\Delta\mu_x > 0$ 和 $\Delta\mu_x < 0$），约束同向运动一致性：

\begin{equation}
L_{\text{con}} = \sum_{\text{axis}} \sum_{\text{sign}} \frac{1}{N}\sum_{i} \|\Delta\mu_i^{\text{axis,sign}} - \overline{\Delta\mu}^{\text{axis,sign}}\|
\end{equation}

其中 $\overline{\Delta\mu}^{\text{axis,sign}}$ 为该分组内偏移量的均值。这一看似简单的约束能有效区分不同运动趋势并促进运动一致性。

\textbf{训练策略与实验表现。}训练共 40{,}000 次迭代：前 1{,}500 次仅优化 3D 高斯以获得稳定初始化，之后联合训练高斯、变形 Hash 场和去噪掩码。损失函数 $L_s$ 和 $L_{\text{con}}$ 从第 3{,}000 次迭代开始贡献，$L_{dn}$ 从第 5{,}000 次开始贡献；各损失权重分别为 $\lambda_{dn} = 1 \times 10^{-2}$、$\lambda_s = 1 \times 10^{-3}$、$\lambda_{\text{con}} = 1 \times 10^{-3}$、$\lambda_m = 5 \times 10^{-4}$。变形场学习率从 $8 \times 10^{-4}$ 指数衰减至 $1.6 \times 10^{-6}$，Hash 编码和微型 MLP 学习率从 $8 \times 10^{-4}$ 衰减至 $3.2 \times 10^{-4}$。在 D-NeRF 合成数据集上，方法平均 PSNR 达到 39.26 dB，显著超越 4D-GS 的 35.71 dB 和 Deformable-GS 的 37.56 dB。在 NeRF-DS 真实数据集上，方法平均 PSNR 为 24.26 dB、SSIM 为 0.8544，优于 Deformable-GS 的 24.11 dB / 0.8524，且使用更少的高斯点实现更高渲染速度。消融实验证实三个模块均有正向贡献：去噪掩码在移除无效点的同时提升渲染精度，静态约束限制静态点对变形场训练的干扰，运动一致性约束改善同一时刻点云运动的规整性。

\paragraph{动态敏感多尺度流形增强}

GS-DMSR（2026.01）~\cite{2601_05584}（Dynamic Sensitive Multi-scale Manifold Enhancement）解决 3D 动态场景重建中模型收敛速度和渲染质量之间的平衡问题。方法构建在 4D-GS 框架之上，基于一个关键观察：4D-GS 虽通过层次化时空耦合表示在动态建模上取得突破，但全局优化百万级高斯参数时，静态或低频动态参数的冗余更新严重拖慢收敛。整体管线如下：对一组 3D 高斯 $G$，首先通过多分辨率体素平面（multi-resolution voxel planes）查询每个高斯的中心坐标和时间戳 $t$ 以提取时空特征，随后进入两项核心模块的协同处理。

GS-DMSR 提出运动显著性驱动的动态高斯优化（MS-DGO）：为每个高斯计算运动显著性系数，量化其属性（位置、旋转、缩放）的动态变化幅度，以预设显著性阈值将高斯动态划分为高/低显著性两类——高显著性高斯（偏离目标场景状态较大）持续通过变形场迭代精化，低显著性高斯（已收敛到逼真近似、满足收敛准则）则免除后续变形场更新。这一机制可理解为对高斯运动状态变化的参数化筛选，通过动态分配计算资源和梯度更新，持续将优化聚焦于最有潜力的高斯单元，在加速收敛的同时降低存储冗余。消融实验表明，移除 MS-DGO（仅用变形场持续优化大量低显著性高斯）会导致存储开销显著增加、收敛速度明显变慢。

在表示增强方面，方法提出多尺度解耦流形变形模块：多头变形解码器 $\mathcal{D} = (\phi_x, \phi_r, \phi_s)$ 分别用独立 MLP 从特征向量 $\mathbf{f}_d$ 预测位置偏移 $\Delta\mathbf{X} = \phi_x(\mathbf{f}_d)$、旋转偏移 $\Delta\mathbf{r} = \phi_r(\mathbf{f}_d)$ 和缩放偏移 $\Delta\mathbf{s} = \phi_s(\mathbf{f}_d)$，变形后的高斯参数通过残差组合得到：$\mathcal{X'} = \mathcal{X} + \Delta\mathcal{X}$、$r' = r + \Delta r$、$s' = s + \Delta s$，最终输出变形 3D 高斯 $G' = \{x', s', r', \sigma, C\}$。在此基础上，通过隐式非线性解码器与显式变形场的协同优化实现分层特征增强——动态特征维度对齐（将隐式解码器各层输出与变形场输出对齐）、跨尺度特征金字塔（融合低频几何结构 $\mathcal{L}_c$ 与高频细节 $\mathcal{H}$）以及非线性流形映射，在增强复杂变形建模能力的同时保持时空连续性。消融实验表明，移除该模块会导致渲染质量显著下降；而 MS-DGO 引入的渲染质量损失恰好由该多尺度流形增强模块部分补偿，二者协同实现收敛速度与渲染质量的整体平衡。

在实验方面，方法在 D-NeRF 合成数据集（渲染分辨率 $800 \times 800$）上达到 PSNR 34.56 dB、SSIM 0.98、LPIPS 0.02，训练时间 8 分钟，渲染速度 96 FPS，相比 4D-GS（PSNR 34.05 dB，82 FPS）在渲染质量和速度上均有提升。在 HyperNeRF 真实数据集（渲染分辨率 $960 \times 540$）上达到 PSNR 25.5 dB、MS-SSIM 0.853，训练时间 26 分钟，FPS 28，同样优于 4D-GS（PSNR 25.2 dB，MS-SSIM 0.845）。方法在 Dynamic Object 数据集上的消融实验进一步验证了两个模块的互补性：移除 MS-DGO（$\phi_r$）时 PSNR 24.429 dB、训练 12.5 分钟、72 FPS，移除多尺度流形增强（$\phi_s$）时 PSNR 24.518 dB、训练 14 分钟、63 FPS，完整模型 PSNR 24.534 dB、训练 13.5 分钟、66 FPS。

\paragraph{全显式动态高斯泼溅}

Fully Explicit Dynamic Gaussian Splatting（2024.10）~\cite{2410_15629} 指出 3DGS 在静态场景中受益于完全显式的点基表示和光栅化渲染管线，但现有动态扩展方法（如 4D-GS）引入隐式编码（如 MLP 解码 4D 基函数），丧失了 3DGS 的核心优势——完全显式性。这些方法还存在对密集点云的依赖问题，需要多帧 COLMAP 重建才能正常初始化，限制了在真实场景中的适用性。方法提出全显式动态高斯泼溅（Ex4DGS），核心思路是将时域表示从隐式编码改为关键帧插值，只在稀疏选定的关键帧（间隔 $I$ 的整数倍时间戳 $nI$）上存储高斯的额外时变信息（位置、旋转、不透明度），中间帧通过显式插值计算。整个流程中所有时域信息直接存储，不经过任何编码过程。

具体而言，位置使用三次 Hermite 样条插值器 CHip（1995）~\cite{bartels1995introduction}：$\mu(t) = \phi(\mu_n, m_n, \mu_{n+1}, m_{n+1}; t')$，其中 $t' = (t - nI) / I \in [0,1]$ 为归一化时间，切线 $m_n = (\mu_{n+1} - \mu_{n-1}) / (2I)$ 和 $m_{n+1} = (\mu_{n+2} - \mu_n) / (2I)$ 由相邻关键帧位置直接计算，无需额外存储。CHip 的展开形式为三阶多项式 $\phi = (2t'^3 - 3t'^2 + 1)\mu_n + (t'^3 - 2t'^2 + t')m_n + (-2t'^3 + 3t'^2)\mu_{n+1} + (t'^3 - t'^2)m_{n+1}$，选择 CHip 而非线性插值或分段三次 Hermite 的原因是：它能以低阶多项式逼近复杂运动而无需额外计算开销或存储。旋转使用球面线性插值 Slerp（1985）~\cite{Animating85ken}：$q(t) = \text{Slerp}(q_n, q_{n+1}; t') = \frac{\sin(\theta(1-t'))}{\sin\theta}q_n + \frac{\sin(\theta t')}{\sin\theta}q_{n+1}$，其中 $\theta = q_n \cdot q_{n+1}$，避免角度值的线性插值偏差。Slerp 直接应用于四元数旋转，无需针对维度做任何修改。不透明度采用简化高斯混合模型处理物体的出现/消失——用两个高斯分量分别控制出现和消失的时序行为，三段式定义：出现段 $e^{-(t-a_s^o)^2/b_s^{o2}}$（$t < a_s^o$）、可见段 $1$（$a_s^o \leq t \leq a_f^o$）、消失段 $e^{-(t-a_f^o)^2/b_f^{o2}}$（$t > a_f^o$），其中 $g_s^o(a_s^o, b_s^o)$ 控制出现，$g_f^o(a_f^o, b_f^o)$ 控制消失。单一高斯只能建模对称的出现/消失行为，无法同时处理突然出现、缓慢消失、长时间可见等多种时序模式，而双高斯模型中两个分量的间隔即表示物体完全可见的持续时长，因此一个高斯点即可表示任意持续时间的可见行为，避免了单一高斯对长时间可见物体需要多个高斯点来拼接的存储浪费。

方法还将高斯分为静态和动态两类：静态高斯 $G_s$ 仅位置随时间线性变化 $\mu(t) = x + t'd$（$t' = t/l \in [0,1]$，$x$ 为枢轴位置，$d$ 为平移向量，$l$ 为场景时长），其余属性（旋转、缩放、颜色）保持不变；动态高斯 $G_d$ 才使用上述关键帧插值。分离过程完全无监督，不需要物体掩码等额外输入。分类通过运动触发器实现：为避免远距离点的偏差，在图像空间测量位移并按距离归一化（即 $d^2/Z^2$，其中 $Z$ 为点到最后一帧相机中心的距离），将所有点按归一化位移排序后取前 $\beta=2\%$ 转换为动态点——转换后的动态点位置由原静态参数 $x$ 和 $d$ 在各关键帧估计，旋转初始化为所有关键帧相同值，不透明度初始化为所有帧可见。消融实验证实，移除动态点提取后 PSNR 从 32.11 降至 28.58（下降 3.53 dB），是所有消融中影响最大的组件。

渐进训练方案从短时视频片段开始逐步增加时长（每 400 次迭代增加一个关键帧间隔 $I=10$），使模型仅需第一帧的 COLMAP 点云即可学习长视频。当关键帧数量增加时，新关键帧的位置通过最后 $I$ 帧的线性回归估计，使前一帧的运动信息得以传递到新关键帧。这一策略防止了因稀疏初始化与未来帧物体的错位而陷入局部最优——消融表明去除渐进训练后 PSNR 下降约 1 dB（32.11 → 31.02）。训练配合点回溯剪枝技术：在单次反向传播中追踪每个高斯在图像空间上的累积误差 $E_v = \sum_k w_k (\prod_{j=1}^{i-1}(1-\alpha_j) q_k)$（$w_k$ 为像素权重，$q_k$ 为像素误差，$i$、$j$ 为深度排序的高斯索引），然后在所有训练视角中聚合 $E_{\text{total}} = \sum_{v \in V} E_v / |V|$，在预定义步骤剪枝高误差点。此外，优化过程中引入运动正则化——对静态点最小化 $\|d\|$，对动态点最小化 $\|\mu_{n+1} - \mu_n\|$——减少累积运动误差，防止高斯过度偏移。

在实验方面，Ex4DGS 在 Neural 3D Video 数据集（18--21 相机、$2704 \times 2028$ 原始分辨率、300 帧，半分辨率评估）上平均 PSNR/SSIM$_1$/LPIPS 达到 33.04/0.956/0.041（Cook Spinach 场景），在 Technicolor 数据集上达到 33.62/0.916/0.088，优于同期 STG、4DGS、4DGaussians 等方法。在仅使用第一帧稀疏点云的条件下，Ex4DGS 对比其他方法优势尤为显著——其他方法因依赖密集点云初始化，无法正确学习视频中的运动。在极端长序列（1000 帧、20000 张图）的 Flame Salmon 场景上，Ex4DGS 仍保持 28.77 dB 的 PSNR 和 392 MB 的合理存储，而 4DGS 虽达 28.37 dB 却需要 6331 MB 存储。消融实验中各插值器的对比显示，对位置和旋转同时使用同等阶数的多项式基至关重要——若二者使用不同阶数，其中一个将因表示能力不足而出现过拟合或过平滑。整个流程无需任何 MLP 推理，在 300 帧、$1352 \times 1014$ 分辨率的场景上于单张 2080Ti GPU 达到 62 fps。

\paragraph{多帧节点引导 4D 高斯}

Mango-GS（2026.03）~\cite{2603_11543}（Multi-Frame Node-Guided 4D Gaussian Splatting）指出已有 4D-GS 方法（如 Deformable 3DGS、SC-GS）采用逐帧变形策略：对每一帧独立预测高斯偏移量，模型倾向于记忆各帧的瞬时状态而非学习底层运动规律，导致快速运动场景中出现模糊和闪烁。虽然时序 Transformer 天然适合捕获多帧间的运动趋势，但典型场景包含百万级高斯，直接对所有高斯施加时序注意力的计算和显存代价过高。Mango-GS 的核心思路是将百万级高斯的运动建模分解为稀疏控制节点的运动建模——用少量节点（$M \ll N$，默认 $M=2048$）捕获场景运动骨架，再通过加权和传播到全部高斯。方法引入两项关键技术。

首先是解耦控制节点表示（decoupled control node）。此前 SC-GS 仅用规范空间中的欧氏距离做 k-NN 来确定节点-高斯的对应关系，但当运动幅度较大时，初始帧中邻近的点可能被带到不同运动区域，静态空间邻域不再有效（邻域漂移）。Mango-GS 将每个控制节点 $n_i$ 解耦为规范位置 $p_i \in \mathbb{R}^3$ 和可学习特征码 $f_i \in \mathbb{R}^D$ 两部分。对于每个高斯 $g_j$，用轻量级 MLP 将其规范参数 $\phi_j = (x_j, q_j, s_j)$（位置、旋转、缩放）映射到嵌入空间，再与节点特征码比较计算亲和度分数，取 top-$k$ 并经 softmax 归一化得到影响权重：

\begin{equation*}
w_{ij} = \frac{\exp(-\mathcal{D}(g_j, n_i))}{\sum_{i' \in \mathcal{K}(j)} \exp(-\mathcal{D}(g_j, n_{i'}))}, \quad \forall i \in \mathcal{K}(j)
\end{equation*}

其中 $\mathcal{D}(\cdot)$ 衡量高斯规范参数与节点特征码在联合位置-特征空间中的距离，$\mathcal{K}(j)$ 为 $k$ 近邻节点集合（默认 $k=3$）。这种学习式亲和度产生的邻域是语义感知的——即使两点在空间上已分离，只要属于同一运动模式仍可保持关联。此外，该设计还具有参数效率优势：直接复用高斯的规范属性（位置、旋转、缩放）计算嵌入，无需为每个高斯额外引入潜在特征，避免了额外的显存开销。

其次是多帧时序注意力网络。网络接收所有控制节点的规范位置和 $T$ 帧时间戳（默认 $T=6$）作为输入，一次性预测 $T$ 帧内所有节点的变形序列：

\begin{equation*}
\{\Delta p_i(t), \Delta q_i(t), \Delta s_i(t)\}_{t=0}^{T-1} = \Phi(\{p_i\}_{i=1}^M, \{t\}_{t=0}^{T-1}; \Theta)
\end{equation*}

其中 $\Phi$ 为时序注意力网络，$\Theta$ 为其可学习参数。网络由深度 MLP 层与沿时间轴的多头自注意力块交错组成。节点位置先经位置编码映射为特征 $x_{\text{emb}}$，时间戳经时间嵌入映射为 $t_{\text{emb}} \in \mathbb{R}^{T \times C_t}$，两者拼接形成初始特征张量 $H^{(0)} \in \mathbb{R}^{N \times T \times C_{\text{in}}}$。在特定层 $l$，特征张量 $H^{(l)} \in \mathbb{R}^{N \times T \times W}$ 首先重排为 $H_{\text{in}}^{(l)} \in \mathbb{R}^{N \times W \times T}$（交换时间轴和特征轴），再沿时间轴施加多头自注意力：

\begin{equation*}
H_{\text{attn}} = \text{MHA}(H_{\text{in}}^{(l)}, H_{\text{in}}^{(l)}, H_{\text{in}}^{(l)})
\end{equation*}

注意力输出通过轻量级门控模块融合回主干流，门控模块从 $H_{\text{attn}}$ 生成门控信号 $(w_{\text{gate}}, w_{\text{bias}}) \in \mathbb{R}^{N \times W \times T}$：

\begin{equation*}
H^{(l+1)} = H^{(l)} \odot \sigma(w_{\text{gate}}) + w_{\text{bias}}
\end{equation*}

其中 $\sigma$ 为 sigmoid 函数，$\odot$ 为逐元素乘法，使网络可自适应控制时序信息的注入量。该门控更新比简单残差连接更具表达能力。最后一层的输出 $H^{(L)}$ 通过独立的线性解码头解码为每个节点在每个时间步的平移、旋转和缩放变形。最终，节点变形通过预计算的 k-NN 权重传播到高斯：

\begin{equation*}
\Delta(x_j)(t) = \sum_{i \in \mathcal{K}(j)} w_{ij} \cdot \Delta p_i(t)
\end{equation*}

旋转和缩放偏移同理加权混合。

训练采用两项策略。第一项是时序输入掩码：随机遮蔽部分时间嵌入，迫使网络仅从可见时序上下文推理全部窗口的变形，防止记忆帧号与瞬时外观。第二项是复合损失函数：

\begin{equation*}
\mathcal{L} = 0.8 \times \mathcal{L}_{\text{frame}} + 0.2 \times \mathcal{L}_{\text{motion}}
\end{equation*}

其中 $\mathcal{L}_{\text{frame}}$ 是 top-$k$ 难帧挖掘光度损失——对批次内所有帧计算 L1 与 DSSIM 加权光度损失，仅取误差最高的 60\% 帧反向传播，且每次迭代动态重新计算 top-$k$，使梯度始终聚焦于当前最难帧。$\mathcal{L}_{\text{motion}}$ 是运动感知损失，在相邻帧差上操作，定义为三项加权和：

\begin{equation*}
\mathcal{L}_{\text{motion}} = \lambda_{\text{diff}} \mathcal{L}_{\text{diff}} + \lambda_{\text{amp}} \mathcal{L}_{\text{amp}} + \lambda_{\text{dir}} \mathcal{L}_{\text{dir}}, \quad (\lambda_{\text{diff}}, \lambda_{\text{amp}}, \lambda_{\text{dir}}) = (0.7, 0.2, 0.1)
\end{equation*}

令 $\delta\hat{I}_t = \hat{I}_t - \hat{I}_{t-1}$ 和 $\delta I_t = I_t - I_{t-1}$ 分别为渲染帧差和真实帧差，三项分别为：帧差幅度一致性 $\mathcal{L}_{\text{diff}} = \sum_{t,x} \|\delta\hat{I}_t(x) - \delta I_t(x)\|_1$，确保场景运动的空间范围正确；运动幅度下界惩罚 $\mathcal{L}_{\text{amp}} = \sum_{t,x} \max(0, \|\delta I_t(x)\|_1 - \|\delta\hat{I}_t(x)\|_1)$，防止网络低估运动幅度导致过度平滑；时序梯度方向余弦相似性 $\mathcal{L}_{\text{dir}} = \sum_{t,x} (1 - \cos(\delta\hat{I}_t(x), \delta I_t(x)))$，惩罚渲染运动方向偏离真实方向的情况。

训练采用两阶段优化：前 5000 次迭代为节点初始化预热阶段，后 35000 次迭代为主训练阶段，时序注意力网络与规范高斯参数联合优化。消融实验表明各组件的渐进贡献：从单帧基线（PSNR 25.15, tLPIPS 0.0250）出发，引入无学习式亲和度的节点反而降低性能（PSNR 24.52），但解耦节点与学习式亲和度恢复并提升了对应关系（PSNR 25.31, tLPIPS 0.0223）；多帧设置加简单时序 MLP 后进一步提升（PSNR 27.30），加入时序自注意力继续改善（PSNR 27.78, tLPIPS 0.0196）；top-$k$ 难帧损失锐化困难区域（PSNR 28.05），运动感知损失最终提供最连贯的时序过渡（PSNR 28.32, tLPIPS 0.0192）。在时间窗口 $T$ 的消融中，$T=6$ 达到最佳平衡——更小的窗口不足以捕获长程运动依赖，更大的窗口因注意力序列长度增加而略微降低质量和渲染速度；近邻数 $K=3$ 效果最优，过小的 $K$ 信息不足导致变形不稳定，过大的 $K$ 导致运动过度平滑丢失细节。

在实验方面，方法在 HyperNeRF-vrig 数据集上达到 PSNR 26.2 dB、MS-SSIM 0.78、tLPIPS 0.0196，渲染速度 149.5 FPS、存储 60 MB，相比 MotionGS（PSNR 24.6, 39.9 FPS）和 TimeFormer（PSNR 24.3, 40.9 FPS）在渲染质量和速度上均有显著提升，速度提升超过 3 倍。在 Neural 3D Video 数据集上达到 PSNR 31.89 dB、SSIM 0.942、LPIPS 0.049，均优于其他方法。

\paragraph{三平面变形与潜在扩散 4D 生成}

SHaDe（2025.05）~\cite{2505_16535} 指出已有动态 NeRF 扩展（如 HexPlane、4D-GS）依赖深度变形 MLP 和密集 3D 结构，导致显存开销大、推理速度慢，难以扩展到长序列或实时应用。SHaDe 的设计围绕三个核心模块展开，每个模块针对一个瓶颈。

第一项是全显式三平面变形场。与此前动态 NeRF 用 MLP 建模运动不同，SHaDe 的变形场完全显式：用三个正交 2D 特征平面（$F_{xy}, F_{yz}, F_{xz}$，分辨率 $256 \times 256$、32 通道）编码时空场景，时间信息通过学习式调制注入。给定查询点 $x = (x,y,z)$ 和时间 $t$，先对三平面做三线性插值获取特征，再通过固定（非学习式）线性投影直接输出 3D 偏移量：

\begin{equation*}
\Delta x = W(f_{xy} + f_{yz} + f_{xz}) + b, \quad x_c = x + \Delta x
\end{equation*}

其中 $W \in \mathbb{R}^{3 \times 32}$ 和 $b \in \mathbb{R}^3$ 为固定投影矩阵和偏置，全程无 MLP 或非线性操作。这种设计将变形计算简化为查表加线性变换，推理极快。变形后的规范点 $x_c$ 随后进入两个并行分支：(1) 基于球谐注意力的辐射场解码器（用于颜色和密度预测）；(2) 潜在扩散精炼路径（通过 T3 编码器压缩到低维潜在空间，再由扩散模型精炼三平面特征和变形偏移）。

第二项是球谐注意力辐射场。传统 MLP 解码器在处理视角相关外观时计算量大且不易解释。SHaDe 在规范空间三平面上存储球谐系数 $\{\alpha_{lm}\}$（最高阶 $L=4$），颜色由视角 $d$ 和时间 $t$ 通过注意力加权的球谐基函数计算：

\begin{equation*}
c(d, t) = \sum_{l=0}^{L} \sum_{m=-l}^{l} \alpha_{lm}(d, t) \cdot c_{lm} \cdot Y_{lm}(d)
\end{equation*}

其中 $Y_{lm}(\cdot)$ 为球谐基函数，$\alpha_{lm}(d,t)$ 是由轻量 MLP（仅以视线方向 $d$ 和正弦时间嵌入 $\gamma(t) \in \mathbb{R}^{64}$ 为输入）预测的各频段注意力权重。值得注意的是，该 MLP 不使用全局相机位姿信息——这是一个有意识的设计选择：球谐基函数本身具有方向性，归一化射线方向 $d$ 即足以捕获视角相关辐射变化，而使用相机位姿可能限制新视角的泛化能力。密度 $\sigma$ 则由独立的三平面场 $F_\sigma$ 经三线性插值和线性投影独立预测：$\sigma = w^\top \text{trilinear}(F_\sigma, x_c) + b$。注意力权重使不同频段的球谐基可随视角和时间动态变化——例如高光仅在特定视角和时间步被强调——比静态 SH 解码器更灵活地建模镜面反射和时变外观。与 Gaussian Splatting 的静态 SH 解码相比，SHaDe 的注意力权重同时适应视角和时间，且针对稀疏输入和规范空间变形进行了优化。

第三项是时序感知潜在扩散先验。为应对快速运动、遮挡和观测模糊导致的重建歧义，SHaDe 引入基于 Transformer 的三平面标记编码器（Tri-Plane Token Transformer, T3）：将三个 $256 \times 256$ 平面切割为 $16 \times 16$ 块、投射为 128 维标记，共 768 个标记经 4 层 Transformer（4 头、隐藏维度 128，辅以正弦时间标记 $\tau(t)$ 增强时序感知）编码为 512 维潜在向量 $z = \mathcal{T}(\mathcal{F}, \Delta, t)$。解码器将 $z$ 映射回精炼的三平面特征和变形偏移 $(\hat{\mathcal{F}}, \hat{\Delta}) = \mathcal{D}(z)$。在此基础上训练 DDPM 去噪扩散模型——去噪网络为 3D U-Net 架构，通过 FiLM 层注入时间条件——使用 DDPM 式训练目标 $L_{\text{diff}} = \mathbb{E}_{t,\epsilon}[\|\epsilon - \epsilon_\theta(z_t, t)\|^2]$（仅 10 步去噪，推理时采用 DDIM 采样，运行时开销约 20\%），为三平面特征提供数据驱动的正则化先验，纠正稀疏视角或模糊运动导致的欠约束重建。时序一致性通过时间正则化 $L_{\text{temporal}} = \|\Delta z_{t \to t+1}\|^2$ 和场景自适应噪声调度（2 层 MLP 预测噪声水平 $\beta_t$，以三平面特征的全局统计量为输入）进一步保障。

训练分两阶段进行：扩散模块先独立预训练 500k 步，再与变形场和辐射场联合精调。联合训练总损失为三项加权和 $L = \lambda_{rec} L_{rec} + \lambda_{diff} L_{diff} + \lambda_{temporal} L_{temporal}$，使用 Adam 优化器（学习率 $5 \times 10^{-4}$，批次大小 4，总步数 200k）。扩散模块不依赖外部数据或预训练模型，仅使用与主流水线相同的合成多视角数据（如 D-NeRF 数据集）从零训练。推理时，向初始潜在向量 $z_0$ 注入受控高斯噪声，再经 DDIM 去噪，使扩散模块充当学习到的生成先验，纠正欠约束或噪声潜在表示。

在效率方面，SHaDe 的参数量为 27.1M（对比 HexPlane 38.7M、4D-GS 54.0M），峰值 GPU 显存 6.4GB（对比 10.2GB、14.3GB），推理速度 1.2s/帧（对比 2.2s、0.08s）。虽然 4D-GS 帧率更高，但它依赖密集视角训练且场景一致性较差；SHaDe 在质量、显存和稀疏视角鲁棒性之间取得了更好的平衡。在 D-NeRF 基准上，SHaDe 在 Lego、T-Rex、Stand Up、Jumping Jacks 四个场景中 PSNR 全面领先（例如 T-Rex: 35.8 vs HexPlane 34.3 vs 4D-GS 33.0），SSIM 和 LPIPS 同样最优。稀疏视角实验（3/5/10/20 视角）表明 SHaDe 在极端稀疏（仅 3 个视角）下仍保持高保真度，而此前方法如 D-NeRF 通常需要 100+ 个密集视角。

消融实验（以 T-Rex 场景为例）进一步量化各模块贡献：完整模型 PSNR 为 35.8；去除变形模块降至 31.3（$-4.5$ dB，模型无法学习非刚性运动，输出模糊或静态）；去除扩散模块降至 33.1（$-2.7$ dB，缺乏生成式精炼，时间平滑度和细节保真度下降）；两者都去除降至 28.7（$-7.1$ dB，退化为静态三平面 NeRF）。消融结果说明变形场和扩散模块互补且不可替代：变形场提供结构感知的运动先验，锚定场景几何并跟踪非刚性运动；扩散模块作为生成式正则化器，在潜在空间中去噪时序编码特征，尤其弥补稀疏或模糊输入下的细节损失。

\paragraph{镜面反射动态场景重建}

SpectroMotion（2024.10）~\cite{2410_17249}（CVPR 2025，Fan et al.）指出现有 3DGS 动态扩展在镜面表面表示上存在根本性缺陷：动态场景方法（Deformable 3DGS、4DGS）依赖低阶球谐表示颜色，其固有的低频特性无法捕捉镜面高光等复杂视觉效果；而镜面渲染方法（如 GaussianShader、GS-IR）虽然引入了环境贴图与 BRDF 模拟镜面反射，但标准环境贴图无法建模时变光照，因此仅适用于静态场景。镜面反射对法线极度敏感：法线微小偏差 $\delta n$ 可导致反射方向 $\omega_r^t = 2(\omega_o^t \cdot \mathbf{n}^t)\mathbf{n}^t - \omega_o^t$ 的剧烈偏转，因此精确的法线估计对动态镜面场景至关重要。SpectroMotion 提出三项核心创新：残差法线校正（residual correction）用于变形过程中的精确法线计算、可变形环境贴图（deformable environment map）适应时变光照条件、三阶段粗到精训练策略显著增强场景几何和镜面颜色预测。它是首个也是目前唯一能够合成真实世界动态镜面场景的 3DGS 方法。

在变形场建模方面，SpectroMotion 沿用 Deformable 3DGS 的范式：将规范高斯集合 $\mathbf{G}$ 中的每个高斯的位置 $\bm{x}$ 和时间 $t$ 输入变形高斯 MLP $F_{\theta_G}$（8 层全连接、256 维隐藏层、第 4 层跳跃连接），预测位置、旋转和缩放残差 $(\Delta\bm{x}^t, \Delta\bm{r}^t, \Delta\bm{s}^t) = F_{\theta_G}(\gamma(\bm{x}), \gamma(t))$，变形后属性通过直接相加获得 $(\bm{x}^t, \bm{r}^t, \bm{s}^t) = (\Delta\bm{x}^t, \Delta\bm{r}^t, \Delta\bm{s}^t) + (\bm{x}, \bm{r}, \bm{s})$，其中 $\gamma$ 为位置编码。该设计将运动学习与几何结构学习分离，在保持稳定几何参考的同时准确模拟动态行为。

在法线估计方面，SpectroMotion 先采用 GaussianShader 的思路，用规范高斯的最短轴 $v_s$ 加残差法线 $\Delta\mathbf{n}$ 估计表面法线。3D 高斯在训练中会趋于扁平化，越扁平的高斯其最短轴越贴近真实表面法线，残差法线用于补偿最短轴近似误差。但直接将其推广到动态场景存在两个关键问题：第一，变形改变高斯的轴长关系，原本的最短轴变形后可能不再是短轴，基于四元数差异的简单旋转在此情形下会导致错误的旋转矩阵；第二，残差法线是静态的，不随变形更新。SpectroMotion 提出基于旋转矩阵的残差校正，同时利用规范高斯和变形高斯的最短轴与最长轴来计算旋转，兼顾变形过程中的旋转和形状变化。首先进行符号一致性校正：对最短轴 $\mathbf{v}_s^t$ 和最长轴 $\mathbf{v}_l^t$，若与规范空间对应轴的点积为负则翻转方向（即 $\mathbf{v}_s^t \leftarrow \mathbf{v}_s^t$ 当 $\mathbf{v}_s \cdot \mathbf{v}_s^t > 0$，否则取负），确保轴方向一致。然后构造正交矩阵 $U = [\mathbf{v}_s, \mathbf{v}_l, \mathbf{v}_s \times \mathbf{v}_l]$ 和 $V^t = [\mathbf{v}_s^t, \mathbf{v}_l^t, \mathbf{v}_s^t \times \mathbf{v}_l^t]$，计算旋转矩阵 $R^t = V^t U^\top$，再将规范空间的残差法线旋转到变形空间。

同时引入扁率比例因子 $\beta^t = (e_l^t - e_s^t) / (e_l - e_s)$（$e_l, e_s$ 为最长/最短轴对应的缩放值），根据变形后高斯的扁平程度调整残差幅度：变形后的残差法线为 $\Delta\mathbf{n}^t = (\beta / \beta^t) R^t \Delta\mathbf{n}$，其中规范空间扁率 $\beta = (e_l - e_s) / (e_l - e_s)$。越扁平的高斯其最短轴越贴近真实表面，所需残差校正越小；反之，不够扁平的高斯需要更大的残差校正来弥补最短轴偏离表面法线的误差。最终法线 $\mathbf{n}^t = \Delta\mathbf{n}^t + \mathbf{v}_s^t$ 并确保朝向观察者方向。法线正则化损失 $L_{\text{reg}} = \beta^k \|\Delta\mathbf{n}\|_2^2$（$k=5$）根据扁率自适应调整惩罚：扁平高斯（$\beta^k \to 1$）惩罚大，迫使残差保持小幅度（因为最短轴已足够准确）；非扁平高斯（$\beta^k \to 0$）惩罚小，允许更大残差以弥补最短轴的偏差。消融实验表明 $k=5$ 为最优指数：$k$ 过小导致 $\mathcal{L}_{\text{reg}}$ 过大、残差被过度抑制、最短轴偏离表面法线；$k$ 过大则 $\mathcal{L}_{\text{reg}}$ 过小、残差主导法线优化、最短轴无法充分对齐真实法线。法线损失 $\mathcal{L}_{\text{normal}} = 1 - \mathbf{N}^t \cdot \hat{\mathbf{N}}^t$ 通过对比渲染法线图 $\mathbf{N}^t$ 和从渲染深度图 $\mathbf{D}^t$ 计算的梯度法线图 $\hat{\mathbf{N}}^t$，改善 3D 高斯之间的局部关联，优化整个场景的深度和法线信息。

在 PBR 渲染方面，SpectroMotion 采用基于图像的光照（IBL）模型和可学习立方体贴图。基于渲染方程，采用 split-sum 近似和 Cook-Torrance 微面元模型将镜面出射辐射分解为两项积分：第一项为纯白环境光下的镜面 BSDF（由 GGX 法线分布函数 $D$、几何衰减因子 $G$ 和菲涅尔项 $F$ 组成），预计算并存储在查找表中；第二项为预积分滤波立方体贴图，每个 mip 层级对应特定粗糙度值。镜面颜色 $L_s$ 由每个高斯可学习的粗糙度 $\rho \in [0,1]$ 控制查询的 mip 层级。在静态和动态阶段完成后几何已稳定，可准确计算反射方向 $\omega_r^t = 2(\omega_o^t \cdot \mathbf{n}^t)\mathbf{n}^t - \omega_o^t$。为处理时变光照，方法引入可变形环境贴图：将标准立方体贴图视为规范环境贴图（使用 $6 \times 128 \times 128$ 个可学习参数），变形反射 MLP $F_{\theta_R}$（与变形高斯 MLP 相同架构：8 层全连接、256 维隐藏层）以位置编码的反射方向 $\gamma(\omega_r^t)$ 和时间 $\gamma(t)$ 为输入，预测反射方向残差 $\Delta\bar{\omega}_r^t = F_{\theta_R}(\gamma(\omega_r^t), \gamma(t))$，变形后反射方向 $\bar{\omega}_r^t = \Delta\bar{\omega}_r^t + \omega_r^t$ 用于查询规范环境贴图获得时变镜面颜色 $\mathbf{c}_s^t$。该方法有效捕捉场景中光照的动态特性，同时保持稳定的规范参考。

颜色分解为漫反射与镜面分量 $\mathbf{c}_{\text{final}}^t = \mathbf{c}_d + \mathbf{c}_s^t \odot \mathbf{s}_{\text{tint}}$，其中 $\mathbf{c}_d$ 为漫反射颜色（使用零阶球谐作为视角无关颜色），$\mathbf{s}_{\text{tint}} \in [0,1]^3$ 为每个 3D 高斯存储的可学习镜面色调，$\odot$ 为逐元素乘法。三阶段粗到精训练策略如下：\textbf{静态阶段}（3000 迭代）使用原始 3DGS 优化位置、缩放、旋转、不透明度和球谐系数，通过光度损失 $L_{\text{color}} = (1-\lambda_{\text{D-SSIM}})L_1 + \lambda_{\text{D-SSIM}} L_{\text{D-SSIM}}$ 稳定静态几何。

\textbf{动态阶段}（6000 迭代）引入变形高斯 MLP 建模动态物体：前 3000 迭代优化规范高斯和变形高斯 MLP 直到场景达到相对稳定状态，然后引入法线损失 $\mathcal{L}_{\text{normal}}$ 同步优化法线和深度，再迭代 3000 步进一步精化几何。

\textbf{镜面阶段}（31000 迭代）将颜色表示从完整球谐切换到 $\mathbf{c}_{\text{final}}^t$：为缓解初始不完整 $\mathbf{c}_{\text{final}}^t$ 可能造成的几何扰动，前 6000 迭代固定变形高斯 MLP 和除零阶球谐、镜面色调及粗糙度外的所有高斯属性，同时暂停致密化；待 $\mathbf{c}_{\text{final}}^t$ 逐渐完整后恢复全部参数优化和致密化，再迭代 3000 步后停止致密化。环境贴图优化也分两步：前 2000 迭代仅优化规范环境贴图以学习时不变光照，之后开始优化变形反射 MLP 捕获时变光照效果直至训练完成。整个训练共 40000 迭代，使用 Adam 优化器。总损失 $L = L_{\text{color}} + \lambda_{\text{normal}} L_{\text{normal}} + L_{\text{reg}}$，其中 $\lambda_{\text{normal}} = 0.01$ 为最优权重——消融实验表明 $\lambda_{\text{normal}} \geq 0.1$ 时法线损失过大导致训练不稳定甚至场景崩溃。

在实验方面，SpectroMotion 在 NeRF-DS 数据集（7 个含动态镜面物体的单目序列）上全面超越 Deformable 3DGS、4DGS、GaussianShader、GS-IR、NeRF-DS 和 HyperNeRF。在动态镜面物体区域的评估（使用 Track Anything 生成动态镜面 mask）中，平均 PSNR/SSIM/LPIPS 达到 20.08/0.5909/0.3094，显著优于次优方法 Deformable 3DGS 的 19.66/0.5826/0.3181。在 HyperNeRF 数据集（无镜面物体的动态场景）上也保持竞争力，平均 LPIPS 达到 0.3143 为最优。方法无需 mask 监督即可自动区分静态和动态物体（变形幅度可视化显示动态区域高亮），并能实现镜面/漫反射分量的逼真分解。在 NVIDIA RTX 4090 上的训练效率为每场景约 1.0--1.9 小时，高斯数量 170k--335k，对于高斯数低于 178k 的场景可实现 $\geq 30$ FPS 的实时渲染。方法的局限性在于：对于手臂或身体进出画面等剧烈场景变化，仅靠变形 MLP 和粗到精策略不足以处理，会产生浮游物伪影。

\paragraph{无混叠 4D 高斯低通滤波}

Alias-free 4D Gaussian Splatting（2025.11）~\cite{2511_18367} 指出 4DGS 在调整相机焦距或高斯与相机距离（即改变采样率）时会出现严重伪影。问题根源在于 2D 膨胀滤波器（dilation filter）：该滤波器在屏幕空间扩展高斯以模拟像素覆盖范围，但当渲染分辨率变化时像素尺寸改变而高斯尺度固定，破坏了膨胀尺度与高斯尺度的比值，引入伪影。此外，现有 4DGS 缺乏对高斯频率的约束，尤其在多视角视频中，运动建模会产生冗余小高斯，这些小高斯在采样率增大时产生高频伪影。然而，直接将 3DGS 的 3D 平滑滤波器（Mip-Splatting）用于单目动态重建会导致重建质量下降——因为 4DGS 中高斯的尺度和位置随时间变化，固定膨胀尺度会使本不可见的小高斯变为可见，并扭曲各向异性高斯的尺度比。针对上述问题，该方法提出三条技术路线：基于 Nyquist-Shannon 采样定理推导 4D 高斯的最大采样频率并引入 \textbf{4D 尺度自适应滤波器}灵活约束高斯频率，设计 \textbf{尺度正则化损失}将频率控制责任在滤波器与高斯自身尺度之间分配，同时引入 Mip-Splatting 的 \textbf{2D Mip 滤波器}处理 4DGS 中的尺度失配，构成完整的抗锯齿方案。

\textbf{最大采样频率推导。}方法首先基于 Nyquist-Shannon 采样定理推导 4D 高斯的最大采样频率。对高斯基元 $k$，采样间隔 $\tau = d/(pf)$（$d$ 为到相机距离、$p$ 为像素尺寸、$f$ 为焦距），采样频率为 $1/\tau$。在所有相机和时间步上取最小采样间隔即可得到最大采样频率：

\begin{equation*}
\tau_k^{-1} = \min_{n \in \{1,\ldots,N\}} \min_{t} \frac{1}{d_n(p_k(t))\, p_n(t)\, f_n(t)}
\end{equation*}

其中 $N$ 为相机总数，$d_n(p_k(t))$ 为基元 $k$ 在时刻 $t$ 对相机 $n$ 的深度，通常用基元中心位置近似。由于 $p_k(t) = p_k + \Delta p_k(t)$ 随时间变化，在训练中反复计算变形后位置代价高昂。因此方法采用\textbf{动量更新}策略：热身阶段用初始位置 $p_k$ 估计粗略最小采样间隔 $\tau_k = \min_n \frac{1}{d_n(p_k)\, f_n}$，变形场稳定后以动量 $\tau_k = (1-\mu_v)\tau_k + \mu_v \cdot \text{current}$ 在线更新。实际实现中 $d_n(t)$ 可直接从 CUDA 渲染过程获取，额外开销可忽略。

\textbf{4D 尺度自适应滤波器。}这是方法的核心贡献，设计基于两个观察：第一，4D 高斯在某些时间步可能不可见，不应在此时膨胀；第二，当某维度的时间变化后尺度 $s_t = s + \Delta s(t)$ 小于或接近膨胀系数 $\sqrt{s}/\tau_k$ 时，固定膨胀会扭曲维度间的尺度比。自适应膨胀系数 $\gamma_{\text{adapt}}$ 首先根据尺度变化比动态调整：

\begin{equation*}
\gamma_{\text{adapt}} = \text{clamp}\!\left(\frac{s_t^2}{s^2},\; \gamma_{\min},\; \gamma_{\max}\right)
\end{equation*}

$\gamma_{\min}$ 约束最大频率，$\gamma_{\max}$ 在一定范围内扩大膨胀尺度以保持膨胀前后的维度比例，减轻滤波对高斯各向异性的影响。对可见高斯，实际膨胀尺度取 $\gamma_{\text{adapt}} \cdot \sqrt{s}$；当 $s_t^2 \leq \gamma_{\text{thre}} \cdot \sqrt{s}/\tau_k^2$ 时，高斯尺度低于阈值，被赋予固定最小膨胀尺度 $\sqrt{s}$ 以防止不可见高斯变为可见。滤波后的 4D 高斯为：

\begin{equation*}
G_k(x, t) = e^{-\frac{1}{2}(x - p_k(t))^\top \hat{\Sigma}_k(t)^{-1}(x - p_k(t))}
\end{equation*}

其中 $\hat{\Sigma}_k(t) = \Sigma_k(t) + \gamma_{\text{adapt}} / \tau_k^2 \cdot I$。

\textbf{尺度正则化。}当高斯尺度小于滤波器的膨胀尺度时，滤波误差会显著放大。方法引入尺度正则化损失 $L_{\text{scales}}$，将频率控制责任在滤波器和高斯自身尺度之间分配：

\begin{equation*}
L_{\text{scales}} = \sum_k \left(\gamma_{\min}\,\frac{\sqrt{s}}{\tau_k^2} - s_k(t)^2\right)^+ \cdot M_1
\end{equation*}

其中 $s_k(t) = s_k + \Delta s_k(t)$ 为基元 $k$ 在时刻 $t$ 的尺度，$M_1$ 为指示函数，仅在 $\gamma_{\text{thre}} \cdot \sqrt{s}/\tau_k^2 < s_k(t)^2 < \gamma_{\min} \cdot \sqrt{s}/\tau_k^2$ 时激活。总训练损失为 $L = L_{\text{color}} + \lambda_1 L_{\text{scales}}$，其中 $L_{\text{color}}$ 为集成了滤波机制的渲染颜色损失。训练策略上，前 3000 次迭代仅优化静态高斯作为热身，之后联合优化高斯与变形场。

\textbf{实验验证。}方法集成到两个框架中进行验证：单目视频重建框架 D3DGS 和多视角视频重建框架 4DGaussians，均同时使用 2D Mip 滤波器和 4D 尺度自适应滤波器。在 D-NeRF 数据集上，多尺度训练与多尺度测试（缩放因子 2/4/8）的实验表明该方法在所有指标上一致超越先前方法，包括使用 Mipmapped Tri-Plane 表示的 DMiT。在单尺度训练、多尺度测试（渲染分辨率 1/1/2/1/4/1/8）的设置下，D3DGS 在全分辨率下能有效捕获变形细节，但低采样率下出现严重膨胀伪影；集成 Mip-Splatting 后全分辨率重建质量下降；而该方法在全分辨率保持原始重建质量的同时，在低分辨率下保留更多真实场景细节。在 N3DV 多视角数据集上，以 4 倍下采样训练并在逐步更高分辨率（1$\times$、2$\times$、4$\times$）渲染的设置下，该方法达到与 Mip-Splatting 相当的抗锯齿性能，同时有效约束 4DGS 频率、减少冗余高斯并提升重建质量。消融实验验证了 4D 尺度自适应滤波器相对 3D 平滑滤波器的两项优势：避免不可见高斯的不必要膨胀，以及通过自适应调整膨胀尺度保持高斯各向异性；同时也确认了尺度正则化损失对于在更小最小膨胀比下仍能保持有效频率约束的必要性。这是将 3DGS 抗锯齿研究扩展到动态场景的首次尝试。

\paragraph{大拓扑变化的 4D 演化建模}

Neural 4D Evolution（2024.11）~\cite{2411_15018} 研究从 2D 图像序列重建拓扑大幅变化的 4D 形状（如球体破裂、物体分裂/合并）。已有工作 NIE（Neural Implicit Evolution）展示了通过瞬态流场从显式 3D 表面演化到目标表面，但直接扩展到 4D 面临两个根本性挑战：一是 NIE 使用 SIREN 架构，其超参数对每个场景高度敏感，且在最佳配置下（将时间位置编码拼接为 $3+64$ 维输入）仍只能学到所有帧的平均形状而非真正的时序变形；二是拓扑变化意味着连通性改变，连续流场无法处理这种不连续性。Neural 4D Evolution 从隐式-显式互补的视角提出三组件框架。

第一组件是基于 HashGrid 编码的动态 SDF 模块。方法定义时间条件 SDF $f_\theta(x, t) = \phi(E(x), \gamma(t))$，其中 $E$ 为 HashGrid 坐标编码器（$F$ 个特征、$L$ 个层级，输出编码维度 $F \times L$），$\gamma(t)$ 为时间位置编码（时间缩放到 $[0,1]$），$\phi$ 为 4 层 128 神经元 MLP（使用双曲正切激活以保证 SDF 的平滑性）。关键设计是只用一个 3D HashGrid（而非 4D 网格或每帧独立网格），通过时间编码共享参数——不同时间步的 SDF 由同一组 HashGrid 权重表示，参数量可控且帧间信息自动流动。

第二组件是隐式表面演化。每轮迭代先用 Marching Cubes 从当前 SDF 提取显式网格 $S_t$，然后在显式网格上最小化能量 $\mathcal{E}$（多尺度光度损失 + Laplacian 正则化，权重 $\lambda_l = 0.0002$），该能量诱导出流场 $V(x)$。通过 PDE $\partial f / \partial t = -\nabla_x f \cdot V$ 更新 SDF 参数，其中非参数下一最优水平集估计为 $s^{i+1}(x) = f_\theta^i(x,t) - \Delta t \cdot \nabla_x f_\theta^i(x,t) \cdot V^i(x)$。演化损失 $L_{\text{evo}} = \frac{1}{|S_t|} \sum_{x_j \in S_t} \|s^{i+1}(x_j) - f_\theta^i(x_j, t)\|^2$ 驱动 SDF 参数更新。由于 HashGrid 的三线性插值保证零水平集在体素间连续移动，即使拓扑发生变化（如球体从一处分裂到两处），SDF 的演化也是连续的。

第三组件是连续式高斯泼溅渲染模块。不同于标准 3DGS 的显式高斯管理（初始化、克隆、分裂、删除），方法将高斯属性隐式化为 MLP 的输出：在每个 epoch，将高斯放置在 SDF 模块提取的网格顶点上，用第二个 HashGrid 编码器对顶点坐标做连续编码，再由 MLP 回归不透明度 $\alpha$、球谐系数 $SH$ 和旋转 $R$。高斯尺度固定为体素距离 $d_{\text{vox}}$ 以防止单个高斯跨越多个表面点。这种隐式参数化使高斯属性随网格拓扑变化自动适应——无需手动管理高斯的增删。时间一致性通过正则化 $\lambda_t |\partial f_\theta / \partial t|$（初始 $\lambda_t = 0.05$，指数衰减）保障。

\textbf{拓扑感知4D重建的两条互补路线。}Neural 4D Evolution 通过隐式 SDF 演化处理拓扑发生变化的场景（物体分裂、合并），代表了"允许拓扑变化"的技术路线。与之互补的另一条路线则旨在"保持拓扑一致"——即使场景发生变形，也要求帧间网格具有相同的顶点数量和连接关系。

\paragraph{Topo4D：2024.06 同期工作}

Topo4D（2024.06）~\cite{2406_00440}（Topology-Preserving Gaussian Splatting）是后一路线的代表性工作，专注于从标定多视角视频中生成拓扑一致的 4D 头部资产（含毛孔级 8K 纹理）。其核心设计是将 3D 高斯中心绑定到预定义拓扑的网格顶点上，形成"Gaussian Mesh"混合表示——每个高斯的位置完全由其关联顶点的位置决定，拓扑信息通过固定的面片连接关系显式维护。训练策略采用交替式几何-纹理优化：几何阶段在物理先验（法向一致性、Laplacian 平滑）约束下优化顶点位置以保持网格规则性，纹理阶段通过 UV 致密化模块（将高斯颜色逆映射到 UV 空间）学习超高分辨率纹理。最终通过法向对齐从高斯表面提取高质量网格。与 Neural 4D Evolution 的隐式 SDF 不同，Topo4D 的显式拓扑绑定天然保证了帧间网格的一致性，使其输出可直接用于表情重定向和重光照等下游应用，但代价是拓扑结构必须预定义（如固定的头部模板），无法处理拓扑发生变化的场景。

在生成式场景中，拓扑变化的处理同样面临挑战。

\paragraph{Helix4D：2026.05 同期工作}

Helix4D（2026.05）~\cite{2605_26109} 通过将强图像到 3D 生成模型 Trellis2 扩展到视频条件下的 4D 生成，首次实现了同时处理拓扑变化（物体破碎、熔化）、透明材质、薄结构和内表面的动态网格生成。其关键设计解决了两个核心问题。第一，跨帧信息共享：采用滑动窗口跨帧注意力机制，并以第一帧为锚点——第一帧由基础 Trellis2 模型生成后注入 Helix4D，使得模型通过跨帧注意力继承 Trellis2 对稀有案例（透明物体、内表面）的生成能力，即使 4D 训练数据中缺乏此类样本。第二，时序编码注入：提出基于 ReRoPE 的时空旋转位置编码，将冗余的低频空间 RoPE 频段重新分配给时间维度，在零额外参数的条件下将 3D 编码扩展到 4D，同时保持与预训练权重的兼容性——相比 SS4D 在每层注意力上添加额外时序 RoPE 的做法，避免了预训练 key/query 投影从未见过的新相位干扰。实验表明 Helix4D 在 ActionBench 上 CD-3D 指标相比 ActionMesh 提升 3.8\%，在包含 52 个视频的复杂动态基准上，ULIP-2 和 Uni3D 分别比最强基线提升 5.7\% 和 7.8\%，用户研究中 67.9\% 的对比偏好 Helix4D。\paragraph{朝向锚定超高斯 5D 表示}

OriGS（2025.09）~\cite{2509_23492}（Orientation-anchored Gaussian Splatting）指出已有动态 3DGS 方法使用运动锚点（图节点、样条控制点、运动基等）建模变形，但普遍依赖低秩假设——相邻区域具有相似运动。这一假设在真实随意拍摄视频中往往不成立：同一物体的不同部位可能呈现截然不同的运动模式（如人体的关节运动），导致空间漂移、结构碎片化或时序不一致。OriGS 的核心洞察是：虽然绝对位置对噪声和尺度歧义敏感，但朝向（orientation）在局部角动量和惯性的平滑控制下演化，是更稳定的运动代理信号。方法提出两个协同组件：全局朝向场和朝向感知超高斯。

\textbf{3D 轨迹恢复。}全局朝向场的输入是场景点的长程 3D 轨迹。由于单目视频天然存在深度歧义，OriGS 首先结合度量深度估计（Metric3D、DepthCrafter、UniDepth）与长程像素追踪（SpatialTracker、CoTracker3），分别预测逐帧深度图 $d^t: \mathbb{R}^2 \rightarrow \mathbb{R}_+$ 和 2D 轨迹 $\{u^t\}_{t=1}^T$，然后借助深度先验将 2D 轨迹提升到 3D 世界空间：$\tau^t = W^t \pi_K^{-1}(u^t, d^t(u^t))$，其中 $W^t$ 为相机位姿，$\pi_K$ 为相机内参投影函数。当相机参数不可用时，通过 bundle adjustment 估计。

\textbf{全局朝向场的构建}分两步。首先，从场景点的 3D 轨迹 $\{\tau^t\}_{t=1}^T$ 中提取主方向：对每个锚点取前 $W$ 帧轨迹，中心化后计算协方差矩阵 $C = \frac{1}{W} \sum_{t=1}^W \hat{\tau}^t \hat{\tau}^{t\top}$，PCA 分解的主特征向量 $v_1$ 反映锚点的主运动轴，即为初始朝向 $O^1 \in SO(3)$。然后用 Procrustes 对齐将朝向传播到所有时间步：在每个 $t$ 时刻，找到锚点的 $K$ 个局部邻域 $\{\tau_k^t\}_{k=1}^K$，通过优化 $\min_{T^t \in SO(3)} \sum_{k=1}^K \|(\tau_k^t - \bar{\tau}^t) - T^t(\tau_k^1 - \bar{\tau}^1)\|^2$ 求解最优旋转，其中 SVD 分解给出 $T^{t*} = VU^\top$（并根据行列式符号校正最后一列），当前朝向 $O^t = T^{t*} O^1$。这产生覆盖全场景、全时序的朝向场，锚点间的相对变换可用 $SE(3)$ 表示为 $\Delta Q_i^{t \rightarrow t'} = (O_i^{t'}, \tau_i^{t'})(O_i^{t}, \tau_i^{t})^{-1}$。

\textbf{朝向感知超高斯}的核心思想是将每个高斯基元与一个定义在时空-几何-朝向联合流形 $\mathcal{M}$ 上的多元动态状态关联。动态状态 $\xi = (\Delta p, \Delta g, t, O) \in \mathcal{M}$ 编码位置偏移 $\Delta p \in \mathbb{R}^3$、几何变化 $\Delta g \in \mathbb{R}^3_+ \times SO(3)$（缩放和旋转）、时间坐标 $\mu_t \in \mathbb{R}$ 和局部朝向 $\mu_O \in SO(3)$。该状态被建模为多元高斯分布 $\mathcal{N}(\mu_\xi, \Sigma_\xi)$，其中均值 $\mu_\xi = (\mu_{\Delta p}, \mu_{\Delta g}, \mu_t, \mu_O)$，协方差矩阵具有分块结构：

\begin{equation*}
\Sigma_\xi = \begin{bmatrix} \Sigma_{(\Delta p, \Delta g)} & \Sigma_{(\Delta p, \Delta g),(t,O)} \\ \Sigma_{(\Delta p, \Delta g),(t,O)}^\top & \Sigma_{(t,O)} \end{bmatrix}
\end{equation*}

捕获组内（如时间或几何自身的）和跨组（如空间-朝向间）相关性。为提高计算效率并避免构建完整协方差矩阵，方法将 $\Sigma_\xi$ 进一步分解为两个可学习子组件：边缘协方差 $\Sigma_{(t,O)}$（使用 Cholesky 分解 $\Sigma_\xi = L_\xi L_\xi^\top$ 参数化以保证正半定性和数值稳定性）和跨协方差项 $\Sigma_{(\Delta p, \Delta g),(t,O)}$，既保留条件切片所需的统计耦合，又规避全矩阵参数化的开销。$\mu_\xi$ 和 $L_\xi$ 作为可学习参数与原始 3DGS 属性联合优化。

\textbf{变形推理}通过条件切片（conditioned slicing）实现：给定查询 $(t', O^{t'})$，利用高斯条件分布公式推导空间-几何子空间的后验分布 $P(\Delta p, \Delta g \mid t', O^{t'}) = \mathcal{N}(\mu_{\text{cond}}, \Sigma_{\text{cond}})$，其中条件均值和条件协方差分别为：

\begin{equation*}
\mu_{\text{cond}} = (\mu_{\Delta p}, \mu_{\Delta g}) + \Sigma_{(\Delta p, \Delta g), (t,O)} \Sigma_{(t,O)}^{-1} (t' - \mu_t, O^{t'} \ominus \mu_O)
\end{equation*}
\begin{equation*}
\Sigma_{\text{cond}} = \Sigma_{(\Delta p, \Delta g)} - \Sigma_{(\Delta p, \Delta g), (t,O)} \Sigma_{(t,O)}^{-1} \Sigma_{(t,O), (\Delta p, \Delta g)}
\end{equation*}

其中 $\ominus$ 表示 $SO(3)$ 上的相对旋转算子。为避免采样引入的优化不稳定性，方法采用最大似然估计原则，取条件均值 $\hat{\mu}_{\text{cond}} = \mu_{\text{cond}}$ 作为确定性预测。这本质上是利用朝向和时间作为结构信号来推断该区域应如何变形——朝向不同的区域自然产生不同的变形模式，无需全局低秩假设。

\textbf{属性调制。}推断得到的动态状态直接用于调制 3D 高斯基元属性以适配渲染管线：位置 $\hat{\mu}_p^{t'} = \mu_p^{t'} + \hat{\mu}_{\Delta p}$，缩放 $\hat{S}^{t'} = S + \hat{\mu}_{\Delta S}$，旋转 $\hat{R}^{t'} = R^{t'} \oplus \hat{\mu}_{\Delta R}$（$\oplus$ 为四元数旋转复合）。此外，方法还根据高斯基元在时间和朝向上偏离其规范状态的距离对其不透明度进行置信度加权调制：

\begin{equation*}
\hat{\sigma}^{t'} = \sigma^{t'} \cdot \exp\left(-\frac{1}{2}\left[\Sigma_t^{-1}(t' - \mu_t)^2 + (O^{t'} \ominus \mu_O)^\top \Sigma_O^{-1}(O^{t'} \ominus \mu_O)\right]\right)
\end{equation*}

这一机制鼓励每个超高斯自适应地表达与当前时间和朝向一致的几何与外观，同时抑制远离其支撑域区域的噪声贡献。

\textbf{锚点驱动变形与朝向查询。}为获取条件切片所需的局部朝向 $O^{t'}$，方法借鉴嵌入式变形图（Embedded Deformation Graphs）的思想，为每个高斯基元 $\mathcal{G}_j$ 分配 $K$ 个邻近朝向锚点的蒙皮权重 $\{w_{ij}\}$，将每个锚点的 $SE(3)$ 相对变换转换为双四元数 $q_i^{1 \rightarrow t'} \in \mathbb{D}$，通过加权双四元数混合 $\hat{q}_j^{1 \rightarrow t'} = \text{Blend}_\mathbb{D}(\{w_{ij}, q_i^{1 \rightarrow t'}\}_{i=1}^K)$ 得到变形后的位置和旋转。由于变形后的高斯不一定恰好在某个锚点上，局部朝向通过周围锚点的 $SO(3)$ 加权插值获得：$O^{t'} = \text{Blend}_{SO(3)}(\{w_{ij}, O_i^{t'}\}_{i=1}^K)$，产生在全局结构与局部动力学之间平滑过渡的朝向场。

\textbf{优化与训练细节。}方法使用 3DGS 的可微光栅化渲染管线，适配高维切片和动态调制。损失函数组合三部分：（1）光度重建损失；（2）2D 对应损失，对齐长程 2D 轨迹和深度先验；（3）变形正则化，对锚点驱动的变换施加 as-rigid-as-possible 约束。方法还设计了剪枝-致密化方案：低不透明度高斯被剪枝，对 $\mu_p$ 和 $\mu_{\Delta p}$ 响应梯度高的空间区域通过局部复制进行致密化。完整优化在单块 NVIDIA RTX A6000 上进行，典型场景约需 0.5--2 小时，取决于视频长度和场景复杂度。

\textbf{实验验证。}在 DyCheck 基准的七个场景上，OriGS 进行了全面的定量评估（PSNR/SSIM/LPIPS），对比涵盖 NeRF 类方法（T-NeRF、NSFF、Nerfies、HyperNeRF 等）和 3DGS 类方法（Dynamic 3DGS、4DGS、D-NPC、Marbles、MoSca 等）。在不使用真实相机位姿的设定下，OriGS 的平均 PSNR/SSIM/LPIPS 为 19.30/0.685/0.270，均超过此前最优方法 MoSca 的 18.84/0.676/0.289；在使用真实位姿时进一步提升至 19.69/0.716/0.256。在涉及旋转变形的场景（如 ``Paper Windmill'' 和 ``Apple''）上优势尤为明显：Paper Windmill 场景 OriGS 达到 22.46 dB（+GT 位姿），比 MoSca 的 21.37 dB 提升约 1 dB，因为朝向场能更好地组织旋转运动的空间结构。

在 DAVIS、YouTube-VOS 等随意拍摄视频上的定性比较显示 OriGS 在快速运动场景（如旋转的纸风车、飞行的猪、鹦鹉）中恢复出更清晰的几何和更连贯的运动。消融实验从 3DGS-MLP 基线（Apple 场景 13.47 dB）逐步叠加组件：引入全局朝向场后提升至 16.28 dB，加入时间条件的超高斯后提升至 18.71 dB，完整 OriGS（加入朝向条件）达到 19.46 dB，验证了每个组件的独立贡献。在 DyCheck 点追踪评估中，OriGS 的 PCK-T@0.05\% 达到 0.851，超过 MoSca 的 0.824 和 CoTracker3 的 0.803。

\textbf{局限性。}OriGS 存在两方面局限：（1）在简单运动主导的场景（如沿固定轴的近刚性平移）中，朝向随时间近似恒定，朝向条件化提供的额外增益有限，但超高斯的高维建模仍引入不必要的计算开销——未来可探索根据场景运动复杂度自适应调整朝向信息的使用或维度的机制；（2）方法依赖 2D 视觉基础模型提供的点轨迹和深度先验来初始化和优化朝向场，这些先验的可靠性直接影响 4D 重建质量。

\subsection{小结}

本节系统介绍了变形场与 4D 高斯表示这一动态 3DGS 的核心研究方向的 78 篇论文。两条主要技术路线——变形场（在规范空间 3DGS 上叠加时间条件的变形网络）和原生 4D（直接在时空域中定义高斯基元）——各有优劣：变形场更加成熟、与静态 3DGS 兼容性好，但受限于变形网络的表达能力和优化难度；原生 4D 表示更加直接、自然支持连续时间建模，但面临参数爆炸的挑战。

效率优化方面，HexPlane 张量分解、4D 哈希编码、锚点驱动变形等技术将 4D 表示的存储开销降低了数倍到数十倍。物理约束方面，哈密顿力学、刚体运动、本构律等物理先验使变形更加物理合理。事件相机提供了高时间分辨率的运动信息，填补了 RGB 帧间的运动空白。前馈方法（DGS-LRM、L4GM）通过大规模预训练实现了秒级的 4D 重建，为实际应用打开了大门。

动态场景重建依然面临开放挑战：大幅运动和拓扑变化的处理、单目设置下的深度歧义、长时间一致性维护，以及存储效率的进一步突破。第~\ref{sec:06_2}~节将深入单目动态场景重建这一更具挑战性但更实用的设定。

\begin{takeaway}
4D 扩展与特殊动态重建涵盖从大规模动态场景到人体/手部/动物等特定目标的重建。语义驱动的弹珠对齐（Marble-SGS）通过刚性约束和 Chamfer 对齐实现多视角时序一致，TraceFlow 用 2DGS 精确法线加残差材质增强实现动态镜面场景渲染。核心趋势是从通用变形场走向结构感知的动态建模。
\end{takeaway}

%% ============================================================

%=====================================================================


\section{单目动态场景重建}
\label{sec:06_3}

单目动态场景重建是 3DGS 动态扩展中最具挑战性也最具实用价值的研究方向。与多视角设置不同，单目输入天然缺乏多视角几何约束，面临三大核心困难：（1）\textbf{深度歧义}（depth ambiguity）——单张图像无法唯一确定场景点的深度值，同一像素可能对应近处小物体或远处大物体；（2）\textbf{尺度不确定性}（scale ambiguity）——单目系统无法恢复绝对尺度，不同帧之间的运动幅度估计可能存在系统性偏差；（3）\textbf{时序一致性}（temporal consistency）——缺乏跨视角监督使得变形场容易过拟合训练帧，导致新视角渲染出现闪烁和伪影。此外，实际单目视频还面临运动模糊、散焦模糊、相机标定缺失等额外挑战。本节按照方法范型将这些工作分为若干子主题进行讨论。

\subsection{变形场与几何引导}
\label{ss:06_3a}

本节介绍两条技术路线：（1）变形场与正则化；（2）深度先验与几何引导。

\subsubsection{变形场与正则化}

本路线介绍在规范空间中定义3D高斯并通过变形场映射到各帧的单目动态重建方法及其正则化策略。

单目动态场景重建的主流范式是在规范空间（canonical space）中定义 3D 高斯，并通过变形场（deformation field）将规范表示映射到各帧。核心挑战在于：如何在不引入过多先验约束的前提下，学习出物理合理的变形，避免过拟合并实现良好的泛化。

\paragraph{高保真单目动态场景重建基线}

Deformable 3D Gaussians（2023.09）~\cite{2309_13101} 是最早将 3DGS 扩展至单目动态场景的方法之一，提出了将规范空间（canonical space）3D 高斯与可学习变形场相结合的基本范式，后续许多单目动态方法都沿用了这一架构。

\textbf{从 Baseline 到新范式。}在 Baseline 层面，一种朴素的做法是为每个时间步独立训练一组 3DGS 模型，然后对相邻时间步的结果进行插值——但这种逐帧策略仅适用于多视角离散捕获（Multi-View Stereo），无法处理连续的单目时间序列。另一种直觉方案是为每个 3D 高斯直接添加时间条件参数（time-conditioned learnable parameter），但这不仅与原始可微分高斯光栅化管线的设计初衷相矛盾，还会导致运动的时空连续性丧失——每个高斯的运动变成彼此独立的逐帧拟合，而非物理上连贯的轨迹。因此，该方法的核心设计是将运动与几何解耦：在规范空间中优化一组与时间无关的 3D 高斯 $G(x, r, s, \alpha)$，同时训练一个隐式变形场，以高斯中心位置 $x$ 和时间 $t$ 为输入，输出位置偏移 $\Delta x$、旋转偏移 $\Delta r$ 和缩放偏移 $\Delta s$，从而将规范高斯变换到各帧对应的变形空间。这种解耦设计的额外好处是引入了场景的几何先验——3D 高斯的位置变化同时受时间和空间坐标的约束，而非各自独立运动。

\textbf{变形场架构选择。}在介绍变形场的具体公式前，值得先解释一个关键的设计选择：为什么变形场采用纯 MLP 而非在静态 NeRF 中广泛使用的 grid/plane 结构（如 TensoRF、K-Planes）？原因在于 grid/plane 结构依赖于场景的低秩假设（low-rank assumption），而动态场景本质上具有更高的秩——同一空间区域在不同时间呈现截然不同的外观和几何。加之 3DGS 的显式点基渲染进一步提升了场景的秩（每个高斯是独立的表示单元），这使得低秩分解结构难以有效捕捉动态变化的复杂性。因此，纯 MLP 成为此处更合理的变形场骨架。

\textbf{变形场公式。}变形场的核心公式为：

\begin{equation}
(\Delta x, \Delta r, \Delta s) = F_\theta\bigl(\gamma(\mathrm{sg}[x]), \gamma(t)\bigr),
\end{equation}

其中 $\gamma(\cdot)$ 为位置编码（positional encoding），形式为 $\gamma(p) = (\sin(2^k p), \cos(2^k p))_{k=0}^{L-1}$，对位置 $x$ 取 $L=10$，对时间 $t$ 取 $L=6$（合成场景）或 $L=10$（真实场景）。实验表明，位置编码可以显著增强渲染结果的细节表现力。变形网络采用深度 $D=8$、隐藏维度 $W=256$ 的 MLP 结构。注意输入位置 $x$ 在送入网络前经过 stop-gradient 操作（即 $\mathrm{sg}[x]$），防止梯度从变形场回传到高斯位置参数，从而将几何学习（高斯位置、旋转、缩放、球谐系数的优化）和运动学习（变形网络的优化）的梯度信号解耦——这一设计确保了规范空间中的高斯表示在训练过程中保持稳定。

在规范空间中，原始 3DGS 的高斯密度自适应控制（adaptive density control）机制仍然生效：位置梯度阈值 $t_{\text{pos}} = 0.0002$ 用于判断哪些高斯需要调整，过小的高斯被克隆（clone）并沿梯度方向移动，过大的高斯被分裂（split）并将尺度除以超参数 $\phi = 1.6$。实验观察到经过约 30k 次训练迭代后，高斯形状趋于稳定，规范空间收敛，间接证明了该解耦设计的有效性。

\textbf{退火平滑训练（AST）。}然而，真实数据集中 SfM 估计的相机位姿不可避免地存在误差（如 HyperNeRF 也曾指出 COLMAP 在动态场景中的位姿偏差问题），而显式点基渲染（point-based rendering）会放大这种位姿偏差——与隐式表示中 MLP 固有的平滑性不同，显式高斯在微小位姿偏移下会产生明显的空间抖动（spatial jitter）。这种抖动在固定时间的新视角渲染中不明显，但在时间插值任务中（即在不同时间步之间进行插值渲染）会导致不规则的闪烁。为解决这一问题，该方法提出了退火平滑训练（Annealing Smooth Training, AST）机制。其核心思想是在训练早期向时间编码注入高斯噪声以增强变形场的时间泛化能力，随着训练推进逐步衰减噪声以保留动态细节：

\begin{equation}
\hat{t} = t + \mathcal{X}(i), \quad \mathcal{X}(i) = \mathcal{N}(0, 1) \cdot \sigma \cdot \tau \cdot \left(1 - \frac{i}{N}\right),
\end{equation}

其中 $\mathcal{X}(i)$ 表示在第 $i$ 次训练迭代中加入的线性衰减高斯噪声，$\mathcal{N}(0,1)$ 为标准正态分布，$\sigma = 0.1$ 为经验缩放因子，$\tau$ 为平均时间间隔，$N = 20\text{k}$ 为退火的截止迭代次数。这意味着在前 20k 次迭代中，变形场的输入时间编码被随机扰动，迫使网络学习对时间扰动鲁棒的变形映射；超过 20k 次后噪声完全衰减，网络在精确时间输入下精细化渲染细节。与 Tensor4D 和 D-NeRF 引入的平滑损失（smoothness loss）相比，AST 不引入额外的损失项和计算开销，而是直接在输入端施加扰动。从正则化的视角来看，AST 本质上是一种数据增强策略：通过对时间维度的随机扰动，扩大了变形场的有效训练分布，使其不仅仅拟合离散的训练时间点，而是学习一个在时间轴上连续平滑的变形映射。

\textbf{实验结果。}该方法在渲染质量和速度上均显著优于当时的隐式神经表示方法。在 D-NeRF 合成数据集上（$800 \times 800$ 全分辨率），Deformable 3DGS 在所有 8 个场景中均取得最佳结果：Hell Warrior（41.54 dB vs.\ D-NeRF 的 24.06 dB）、Mutant（42.63 dB）、Hook（37.42 dB）、Bouncing Balls（41.01 dB）、Lego（33.07 dB）、T-Rex（38.10 dB）、Stand Up（44.62 dB）、Jumping Jacks（37.72 dB），相比 D-NeRF 平均提升约 10--17 dB，同时保持实时渲染速度。

\paragraph{不确定性感知野外视频正则化}

4DGS in the Wild（2024.11）~\cite{2411_08879} 针对手持单目相机随意拍摄的真实视频，首次关注了自然条件下的单目动态重建问题。与此前方法多假设受控采集环境（如 HyperNeRF 使用不切实际的训练-测试划分，将同一视频轨迹中的帧同时用于训练和测试，使得任务更接近视频插值而非真正的新视角合成）不同，该方法面向的是相机围绕快速运动物体拍摄的真实手持视频场景。在这种场景中，4DGS 模型容易过拟合训练帧，而朴素地施加正则化虽能改善测试视角的渲染质量，却会牺牲训练图像的重建精度——二者之间存在根本性的平衡难题（trade-off）。

\textbf{从 Deformable 3DGS 到 4DGS in the Wild。}在方法架构上，4DGS in the Wild 与前述 Deformable 3DGS 的核心区别在于变形场的骨架选择。Deformable 3DGS 采用纯 MLP 变形场以避免 grid/plane 结构的低秩假设与动态场景高秩本质的冲突；而 4DGS in the Wild 选择 HexPlane（2023）~\cite{wu20234d} 作为基线变形场——HexPlane 通过六个特征平面（三个空间平面 + 三个时空平面）的外积分解来参数化时空变形，在计算效率上优于纯 MLP。然而，两者的核心差异并不在变形场本身，而在于正则化策略：Deformable 3DGS 通过 AST 在时间编码输入端施加扰动来间接正则化变形场，本质上是一种时间维度的数据增强；4DGS in the Wild 则直面正则化的根本困境——不加区分的正则化会同时作用于已充分观测和未充分观测的区域，在改善测试视角的同时损害训练精度——从而提出了完全不同的\textbf{不确定性感知正则化}（uncertainty-aware regularization）范式。

\textbf{不确定性量化。}该方法的核心思想是：不是对所有区域施加均匀正则化，而是首先量化每个高斯基元的不确定性，然后仅对不确定区域施加先验约束，低不确定性区域（即训练数据已提供充足重建细节的区域）则不加正则化。具体而言，每个高斯基元 $g_k$ 的信息量通过其对所有训练像素的 $\alpha$-混合贡献累积计算：

\begin{equation}
C_k = \sum_{I \in \mathcal{T}} \sum_{\mathbf{r} \in I} \alpha^\text{render}_k(\mathbf{r}),
\end{equation}

其中 $\mathcal{T}$ 为训练图像集合，$\alpha^\text{render}_k(\mathbf{r})$ 是高斯基元 $g_k$ 在像素 $\mathbf{r}$ 处的 $\alpha$-混合贡献。信息量高的高斯被更多像素以更高的混合权重观测到，其参数估计准确、不确定性低；信息量低的高斯则缺乏充分观测、不确定性高。不确定性通过 Sigmoid 函数归一化：

\begin{equation}
U_k = 1 - \mathrm{Sigmoid}(C_k;\, c_0, c_1),
\end{equation}

其中 $c_0, c_1$ 分别控制 Sigmoid 函数的拐点偏移和斜率。对于任意未见视角，通过 $\alpha$-混合将逐高斯不确定性投影为 2D 不确定性图：

\begin{equation}
U(\mathbf{r}) = \sum_{k=1}^K \alpha^\text{render}_k(\mathbf{r}) \cdot U_k.
\end{equation}

该不确定性图为后续两种正则化提供了空间自适应的权重信号。

\textbf{不确定性感知扩散正则化。}方法设计了第一种自适应正则化——不确定性感知扩散损失 $L_\text{UA-diff}$。其流程为：首先利用视觉-语言模型 BLIP 从训练帧自动生成文本提示，然后通过 DreamBooth 在训练图像上微调 Stable Diffusion 模型，使其理解特定场景内容。对于未见视角的渲染图像 $\hat{I}$，先用潜在扩散编码器 $\mathcal{E}$ 将其编码到潜在空间，再按噪声调度注入噪声得到加噪潜在表示：

\begin{equation}
\mathbf{x}_t = \sqrt{\bar{\alpha}_t}\,\mathcal{E}(\hat{I}) + \sqrt{1 - \bar{\alpha}_t}\,\boldsymbol{\epsilon}, \quad \boldsymbol{\epsilon} \sim \mathcal{N}(\mathbf{0}, \mathbf{I}), \quad t \in [0, T],
\end{equation}

其中 $\bar{\alpha}_t$ 控制噪声水平。随后以 BLIP 生成的文本嵌入为条件，通过 50 步 DDIM 采样从加噪潜在表示中生成精炼图像 $I_\text{DDIM} = \mathcal{D}(\mathbf{x}_t, \mathbf{v})$，其中 $\mathcal{D}$ 为扩散解码器，$\mathbf{v}$ 为文本嵌入。最终的不确定性感知扩散损失为：

\begin{equation}
L_\text{UA-diff} = \frac{\lVert U \odot (\hat{I} - I_\text{DDIM}) \rVert_2}{\lVert U \rVert_2} + \frac{\lVert U \odot (\hat{I} - I_\text{DDIM}) \rVert_1}{\lVert U \rVert_1},
\end{equation}

其中 $\odot$ 表示逐元素乘积，$U$ 为该视角的不确定性图。不确定区域（$U$ 大）受到更强的扩散先验约束，低不确定区域则几乎不受影响。为避免每轮迭代都运行 DDIM 采样带来的计算开销，方法每 2000 次迭代随机采样 200 张未见视角图像，预计算并缓存其 DDIM 结果，在后续 2000 次迭代中复用缓存结果。

\textbf{不确定性感知深度平滑正则化。}方法设计了第二种自适应正则化——不确定性感知全变分损失 $L_\text{UA-TV}$。对未见视角的深度图 $D(\mathbf{r}) = \sum_{k=1}^K \alpha^\text{render}_k(\mathbf{r}) \cdot d_k$（其中 $d_k$ 为高斯基元中心到相机的深度值），施加以不确定性图加权的全变分正则化：

\begin{equation}
L_\text{UA-TV} = \frac{1}{u_r}\sum_{i,j}\frac{U_{i,j} + U_{i+1,j}}{2}\,|D_{i,j} - D_{i+1,j}| + \frac{1}{u_c}\sum_{i,j}\frac{U_{i,j} + U_{i,j+1}}{2}\,|D_{i,j} - D_{i,j+1}|,
\end{equation}

其中 $u_r = \sum_{i,j}\frac{U_{i,j}+U_{i+1,j}}{2}$ 和 $u_c = \sum_{i,j}\frac{U_{i,j}+U_{i,j+1}}{2}$ 分别为垂直和水平方向相邻像素的平均不确定性之和，起到归一化作用。该损失对高不确定区域施加更强的深度平滑约束以抑制深度噪声，同时保留已准确重建区域的几何细节。消融实验证实：不使用不确定性加权的朴素 $L_\text{TV}$ 虽能改善测试视角质量，但会在训练图像上产生欠拟合；而 $L_\text{UA-TV}$ 则在训练重建精度与测试泛化质量之间实现了更优的平衡。

\textbf{动态区域初始化。}此外，该方法还解决了一个被此前工作忽视的\textbf{动态区域初始化缺失}问题。在快速运动场景中，SfM 假设场景静态，将动态区域视为噪声并拒绝重建特征点，导致这些区域完全没有高斯基元初始化。这种不完整的初始化会扰乱训练过程：静态区域的高斯被反复克隆和分裂以试图填充动态区域，不仅效率低下，还可能导致内存溢出。方法提出利用场景流（scene flow）（2021）~\cite{zhang2021consistent} 识别动态像素，并随机选取子集进行逆投影初始化。对选中的动态像素 $\mathbf{r}$，对应高斯基元 $g'_k$ 的初始位置和特征向量通过逆投影计算：

\begin{equation}
\boldsymbol{\mu}_k = \pi^{-1}(\mathbf{r},\, D'(\mathbf{r})), \quad \mathbf{f}_k = I(\mathbf{r}),
\end{equation}

其中 $D'(\mathbf{r})$ 为场景流算法估计的深度值，$\pi^{-1}$ 为逆投影函数将像素重投影到 3D 空间，特征向量 $\mathbf{f}_k$ 的前三个分量设为该像素的 RGB 颜色值，其余维度（用于球谐表示）填充零。这为动态区域提供了合理的初始高斯分布。

\textbf{数据驱动损失。}由于随意拍摄的单目动态重建是高度病态的问题，方法还引入了基于深度图和光流的辅助监督。深度驱动损失为 $L_\text{depth} = \lVert D - D' \rVert_1$，其中 $D$ 为 $\alpha$-混合渲染深度图，$D'$ 为场景流算法估计的深度图。光流驱动损失类似：对帧 $t$ 和 $t'$，通过 $\alpha$-混合渲染光流 $\hat{F}^{t \to t'}(\mathbf{r}) = \sum_{k=1}^K \alpha^\text{render}_k(\mathbf{r}) F_k^{t \to t'}$，其中 $F_k^{t \to t'} = \pi(\boldsymbol{\mu}_k^{t'}) - \pi(\boldsymbol{\mu}_k^t)$ 表示第 $k$ 个高斯基元在两帧间的投影位移，然后与 RAFT 光流估计对齐：$L_\text{flow} = \lVert \hat{F}^{t \to t'} - F'^{t \to t'} \rVert_1$。总数据驱动损失为 $L_\text{data} = L_\text{depth} + L_\text{flow}$。

\textbf{完整训练目标。}方法的总损失函数为：

\begin{equation}
L = L_\text{recon} + \lambda_\text{grid}\, L_\text{grid} + \lambda_\text{data}\, L_\text{data} + \lambda_\text{UA-diff}\, L_\text{UA-diff} + \lambda_\text{UA-TV}\, L_\text{UA-TV},
\end{equation}

其中 $L_\text{recon}$ 为标准训练图像重建损失，$L_\text{grid}$ 为 HexPlane 变形场的正则化损失，$\{\lambda_\text{grid}, \lambda_\text{data}, \lambda_\text{UA-diff}, \lambda_\text{UA-TV}\}$ 为各项的平衡超参数。

\textbf{实验结果。}方法在 DyCheck 数据集（手持单目相机拍摄的真实动态视频）上进行评估，该数据集包含 apple、block、paper-windmill、teddy、space-out、spin、wheel 等 7 个可评估场景，具有 held-out 测试视角。实验表明 UA-4DGS（即该方法的简称）在所有指标上均显著优于 D-3DGS（即 Deformable 3DGS）、ST-4DGS 和 4DGS 等 4D 高斯 Splatting 方法。值得注意的是，在这些随意拍摄的单目场景中，基于 MLP 的隐式方法（如 HyperNeRF）在整体表现上反而优于所有基于 3DGS 的方法——这暗示基于高斯的显式方法更倾向于对单个基元进行局部优化，在面对训练视角覆盖不足的 in-the-wild 场景时更容易过拟合。消融实验证实了每个组件的贡献：动态区域初始化通过将高斯更好地对齐到场景几何来提升数据驱动损失的效力；$L_\text{UA-diff}$ 带来了最显著的增益；在此基础上叠加 $L_\text{UA-TV}$ 可获得额外提升。此外，该方法还可以推广到静态场景的少视角重建：将不确定性感知正则化应用于 FSGS（Few-Shot Gaussian Splatting）后，在 LLFF 数据集上使用 3 张训练图像的设置下也取得了改进，验证了该方法框架的通用性。

\paragraph{鲁棒动态高斯物体分离}

前述 4DGS in the Wild 假设已知相机位姿（通过 SfM 或外部设备获取），并通过不确定性感知正则化缓解手持单目视频中的过拟合问题。然而，在更一般的场景中——随意拍摄的单目视频既无已知相机位姿，也无多视角线索——相机运动和物体运动的天然耦合（camera-object motion ambiguity）成为更根本的障碍：从二维投影中无法唯一确定相机移动了多少、场景中物体又移动了多少。此前的方法要么依赖专业采集设备（如 Dynamic 3DGS 需要多视角同步捕获），要么局限于特定场景条件（如 D-NeRF 假设相机传送/teleportation，HyperNeRF 假设环境动态/ambient dynamics）。RoDynRF 虽然也是无位姿方法，但受限于前向场景或小视角变化。RoDyGS（2024.12）~\cite{2412_03077}（Lee et al., POSTECH）正是面向这一最弱假设场景的端到端框架：从随意拍摄的单目视频出发，无需任何外部传感器或手动标定，联合优化动态场景表示和相机运动。

\textbf{动态表示模型：运动基分解。}在介绍 RoDyGS 的正则化策略前，值得先理解其动态表示的具体形式。RoDyGS 沿用 DynMF 的运动基分解（motion basis factorization）框架：不直接为每个高斯参数化时间变位置和旋转，而是将运动建模为共享运动基的线性组合。具体而言，运动基函数 $\boldsymbol{\phi}$ 由瓶颈 MLP（bottlenecked MLP）参数化，接收归一化时间步 $t/T \in [0,1]$ 为输入，输出 $B$ 维平移基 $\boldsymbol{b}^\mu(t)$ 和旋转基 $\boldsymbol{b}^q(t)$。每个高斯基元持有可学习的运动系数 $\boldsymbol{m} \in \mathbb{R}^B$，其时间变属性通过组合得到：

\begin{equation}
\boldsymbol{\mu}(t) = \boldsymbol{\mu}_c + \boldsymbol{m} \cdot \boldsymbol{b}^\mu(t), \quad \boldsymbol{q}(t) = \frac{\boldsymbol{q}_c + \boldsymbol{m} \cdot \boldsymbol{b}^q(t)}{\|\boldsymbol{q}_c + \boldsymbol{m} \cdot \boldsymbol{b}^q(t)\|},
\end{equation}

其中 $\boldsymbol{\mu}_c$ 和 $\boldsymbol{q}_c$（四元数表示）分别为规范空间中的位置和旋转。这种分解的核心优势在于引入了低秩偏置（low-rank bias）：瓶颈 MLP 的有限容量天然限制了运动基的复杂度，而共享基使得少量运动模式即可驱动整个场景的动态——这与非刚性结构恢复（Non-Rigid Structure from Motion）领域的经典发现一致：复杂物体的运动可以在低维轨迹空间中有效近似。

\textbf{静态-动态部件分离。}RoDyGS 的第一个核心设计是\textbf{静态-动态部件分离}（Static-Dynamic Part Separation）。在优化开始前，方法首先利用两个预训练模型获取初始化信息：（1）使用 MASt3R（2024）~\cite{mast3r} 从输入视频的帧对中估计初始相机位姿 $\boldsymbol{E}$ 和初始 3D 点云（作为高斯基元的初始化位置）；（2）使用 TAM（Track Anything Model）从视频序列中生成运动掩码（motion mask），标识哪些区域包含动态物体。基于运动掩码，高斯基元被划分为静态集合 $\mathcal{G}_S$ 和动态集合 $\mathcal{G}_D$ 两组。这种分离的优化意义在于：运动相关约束仅施加于 $\mathcal{G}_D$，而 $\mathcal{G}_S$ 与相机参数 $\boldsymbol{E}$ 联合优化——从而将相机运动与动态场景元素解耦，显著降低了相机-物体运动歧义。值得注意的是，这里的数据驱动初始化（MASt3R 替代 SfM）是必要的：传统 SfM（如 COLMAP）假设静态场景，在动态环境中经常失败；MASt3R 基于大规模互联网数据训练，能在大视角变化和动态遮挡下产生鲁棒的位姿估计。此外，分离还带来了渲染效率的提升：运动估计仅限于 $\mathcal{G}_D$，静态区域无需计算变形。

\textbf{几何一致性正则化。}动态物体在随意拍摄的单目视频中不可避免地被部分遮挡，导致某些时间步缺少观测——这破坏了动态高斯的时空一致性。RoDyGS 对动态部分施加两种互补的几何正则化：（1）\textbf{跨时间距离保持}（distance-preserving across time），对动态高斯 $g \in \mathcal{G}_D$ 及其 $K$ 个最近邻 $\mathcal{N}(g)$，约束任意时间步中邻居间的距离保持稳定：

\begin{equation}
\ell_\text{dp}(g, t) = \frac{1}{K |\mathcal{T}'|} \sum_{g' \in \mathcal{N}(g)} \sum_{t' \in \mathcal{T}'} \|d(g, g', t) - d(g, g', t')\|_2^2,
\end{equation}

其中 $d(g_1, g_2, t) = \|\boldsymbol{\mu}_{g_1}(t) - \boldsymbol{\mu}_{g_2}(t)\|_2^2$ 为两个高斯在时间 $t$ 的中心距离，$\mathcal{T}'$ 为随机采样的时间步集合。这一约束保持了动态高斯集合的刚性（rigidity）：即使某些高斯在某些时间步被完全遮挡，其与邻居的距离关系仍然由其他时间步的观测约束维持。这一思路借鉴了 LaplacianFusion 等基于拉普拉斯坐标的网格重建方法，将几何一致性从单点层面提升到邻域关系层面。（2）\textbf{物体表面平滑}（object surface smoothing），约束每个动态高斯位于其 $K$ 个最近邻的平均位置附近：$\ell_\text{surf}(g, t) = \|\boldsymbol{\mu}_g(t) - \frac{1}{K}\sum_{g' \in \mathcal{N}(g)} \boldsymbol{\mu}_{g'}(t)\|_2^2$，防止单个高斯偏离局部表面。两项损失组合为几何正则化 $\mathcal{L}_\text{geo} = \lambda_\text{dp} \mathcal{L}_\text{dp} + \lambda_\text{surf} \mathcal{L}_\text{surf}$，作用于随机采样的动态高斯子集 $\mathcal{G}_D' \subset \mathcal{G}_D$。

\textbf{稳态运动正则化。}RoDyGS 的第三个核心设计是\textbf{稳态运动正则化}（Steady Motion Regularization）。其基本假设是稳态运动（steady motion）——在短时间间隔内运动保持一致、无漂移。对运动基 $\boldsymbol{\phi}$ 的第 $i$ 个元素，施加运动连续性约束：

\begin{equation}
\ell_\text{sm,i}(t) = \|\phi_i(\tfrac{t}{T}) - \phi_i(\tfrac{t-1}{T})\|_2^2,
\end{equation}

直接约束运动基函数在相邻时间步之间的变化量。然而，统一的平滑约束会压制高频运动细节。为避免过度平滑，方法引入\textbf{加权运动正则化}方案，使用累积指数函数为每个运动基元素分配不同权重：

\begin{equation}
w_i = B \times \frac{1 - \exp(-3i/B)}{\sum_{i'=1}^{B}(1 - \exp(-3i'/B))}, \quad i = 1, 2, \ldots, B,
\end{equation}

完整损失为 $\mathcal{L}_\text{sm} = \frac{1}{B(T-1)}\sum_{t=2}^{T}\sum_{i=1}^{B} w_i \cdot \ell_\text{sm,i}(t)$。权重 $w_i$ 随索引 $i$ 单调递增：低频运动基（$i$ 较小，对应运动模式的主要趋势）受到更强的平滑约束，确保运动的整体连贯性；高频运动基（$i$ 较大，对应细节运动）受到较弱的约束，保留运动的精细结构。作者通过可视化对比表明，均匀权重的运动正则化在动态物体边缘产生鬼影伪影（ghosting artifacts）和模糊，而加权方案有效抑制了这些伪影，保持物体边界的清晰度。

\textbf{深度监督与目标函数。}RoDyGS 使用 DepthAnything 提供的单目深度图引入场景几何先验。由于单目深度估计存在尺度歧义（scale ambiguity），方法采用 Pearson 深度损失（而非 $\ell_1/\ell_2$ 距离损失）——Pearson 相关系数量化渲染深度与估计深度之间的线性相关性，对仿射变换不变，因此不受尺度偏移和缩放影响。完整的静态和动态目标函数分别为 $\mathcal{L}_\text{static} = \mathcal{L}_\text{recon} + \mathcal{L}_\text{depth}$ 和 $\mathcal{L}_\text{dynamic} = \mathcal{L}_\text{recon} + \mathcal{L}_\text{depth} + \mathcal{L}_\text{geo} + \mathcal{L}_\text{sm}$，其中 $\mathcal{L}_\text{recon} = \lambda_{\ell_1}\|\hat{\mathbf{I}}_t - \mathbf{I}_t\|_1 + \lambda_\text{SSIM}\frac{1 - \mathrm{SSIM}(\hat{\mathbf{I}}_t, \mathbf{I}_t)}{2}$ 为标准 3DGS 光度损失。

\textbf{Kubric-MRig 基准数据集。}现有动态新视角合成基准存在多个维度的缺失：Sintel 仅有前向视角，Tanks and Temples 无动态运动，D-NeRF 无背景，HyperNeRF 无真值相机位姿，iPhone 无宽基线视角。RoDyGS 的作者提出了 Kubric-MRig 基准数据集，基于 Kubric（Blender 合成场景生成器）构建，首次同时满足四个关键属性：宽基线视角（wide viewpoints）、大幅物体运动（large motion）、真值相机位姿（GT camera）和完整背景。每个场景包含两个 100 帧单目视频：训练视频由围绕物体移动的相机拍摄，测试视频由固定在训练视频初始位置的相机拍摄——这一评估协议能准确衡量动态新视角合成质量，而非简单的视频插值。

\textbf{实验结果。}在 Kubric-MRig 上，RoDyGS（PSNR 19.44, SSIM 0.7169, LPIPS 0.3401）显著优于无位姿基线 RoDynRF（PSNR 18.36, SSIM 0.6121, LPIPS 0.5984）和 SoM（PSNR 16.83, SSIM 0.6283, LPIPS 0.6682），渲染质量接近使用真值相机位姿的 D-NeRF（PSNR 20.65）和 RoDynRF（PSNR 20.80）。消融实验验证了每个组件的贡献：部件分离是最大的性能增益来源（PSNR 从 16.41 提升到 19.05），几何正则化进一步微调（19.08），稳态运动正则化提供最终增益（19.44）。值得注意的是，SoM 等依赖长程点跟踪先验（CoTracker）的方法在大相机运动和复杂动态场景中表现不佳，因为不准确的点轨迹在优化中引入噪声——这进一步验证了 RoDyGS 在框架内直接优化高斯运动（而非依赖外部对应关系先验）的设计选择。

在 iPhone 数据集上，RoDyGS（PSNR 17.38, SSIM 0.4656, LPIPS 0.4361）取得了与 SoM（PSNR 17.55）相当的结果，且无需已知相机位姿（iPhone 标准评估协议假设已知位姿，RoDyGS 不使用此信息）。此外，RoDyGS 在 Tanks and Temples 静态场景数据集上与 CF-3DGS 取得了可比的渲染质量（PSNR 31.51 vs.\ 31.28, SSIM 0.93, LPIPS 0.08），证明其部件分离和正则化策略在静态场景中也不损害性能。

\textbf{局限性。}RoDyGS 的正则化策略编码了弱刚性约束（weak rigidity constraint），假设动态场景元素不发生高度复杂的非刚性变形——因此难以处理人体运动或宽松衣物的形变。作者指出，引入参数化模型（如 SMPL 用于人体动力学）或其他数据驱动先验是将该方法推广到日常视频的未来方向。此外，与所有 4D 重建方法一样，RoDyGS 只能重建输入视频中实际观测到的部分：完全遮挡或长期遮挡的区域无法恢复，作者建议利用生成式先验（如扩散模型）推断未观测区域作为潜在的解决途径。

\paragraph{渐进式动态场景重建}

ProDyG（2025.09）~\cite{2509_17864}（Chen et al., ETH Zürich）是首个支持在线操作的 RGB-only 动态场景重建系统。论文指出，现有动态 SLAM 方法要么仅滤除动态区域、要么依赖 RGB-D 输入，离线方法无法扩展到长视频序列，而基于 Transformer 的前馈方法则缺乏全局一致性和精细外观。ProDyG 同时满足四个实际需求：（1）在线运行，紧密耦合位姿估计与稠密地图重建以保证可扩展性；（2）全局位姿与地图一致性；（3）使用 3DGS 表示以获取精细的外观和几何；（4）灵活支持 RGB 和 RGB-D 输入。

在\textbf{运动无关的在线相机跟踪}方面，ProDyG 以 Splat-SLAM 为后端，维护一个因子图存储相机外参 $\omega_i$、逐关键帧视差估计 $d_i \in \mathbb{R}^{(H \times W) \times 1}$ 和光流 $\tilde{p}_{ij} - p_i \in \mathbb{R}^{(H \times W \times 2) \times 1}$。跟踪通过稠密捆绑调整（Dense Bundle Adjustment, DBA）实现，核心创新在于\textbf{运动掩码策略}：每次 DBA 迭代后，计算残差光流
\begin{equation}
\hat{r}_{ij} = (\tilde{p}_{ij} - p_i) - (p_{ij} - p_i) = \tilde{p}_{ij} - p_{ij},
\end{equation}

即估计光流减去相机引起的光流（$p_{ij}$ 为根据当前位姿和视差计算的投影流），静态区域残差接近零，动态区域残差较大。对每个关键帧 $i$，在所有连接的目标关键帧 $N_i$ 上计算归一化平均幅值：

\begin{equation}
\bar{r}_i(x,y) = \frac{1}{|N_i|}\sum_{j \in N_i} \frac{\|\hat{r}_{ij}(x,y)\|}{\frac{1}{HW}\sum_{x',y'}\|\tilde{p}_{ij}(x',y') - p_i(x',y')\|},
\end{equation}

该归一化将残差除以全局光流幅值后再对连接关键帧取平均，使得阈值选择与光流大小无关。取 top 20\% 阈值生成粗运动掩码 $C_i$，在 DBA 的 Mahalanobis 距离中将动态区域的置信度权重置零（$\Sigma'_{ij} = \text{diag}(w_{ij}\bar{C_i})$），实现运动无关的鲁棒位姿估计。由于 $\bar{r}_i$ 每次 DBA 迭代都更新，粗掩码在迭代过程中逐步精化。

为提升动态-静态分离精度用于后续重建，ProDyG 利用 SAM2 对粗掩码进行\textbf{语义引导的精细化}，分为初始化和增量预测两阶段。初始化阶段，在经中值滤波（$5 \times 5$ 核）后的粗掩码 $C_i$ 的连通域质心处生成提示点（prompting），驱动 SAM2 生成分割掩码。增量预测阶段则将已有的逐物体掩码扩展到新帧，并基于新的粗掩码添加新分割——通过计数提示点候选数量来验证每个分割以过滤误报。在进入增量阶段前，方法还根据最新精细掩码中的动态像素比例自适应调整阈值（不再固定为 20\%），以适应序列特定的运动规模。精细掩码 $\{M_i\}$ 生成后替换 DBA 中的粗掩码（$\Sigma'_{ij} = \text{diag}(w_{ij}(1 - M_i))$），进一步提升后续 DBA 迭代（包括 local BA、global BA 和 loop BA）的跟踪鲁棒性。

在\textbf{非关键帧位姿与深度获取}方面，由于 SLAM 后端仅操作稀疏关键帧以保证效率，而动态场景重建需要时间稠密的监督以捕捉平滑运动，ProDyG 需要为关键帧之间的非关键帧推导位姿和深度。非关键帧位姿通过类似 DROID-SLAM 的方式获取：在相邻关键帧之间构建临时局部图，执行循环更新算子和 DBA 层优化非关键帧位姿。非关键帧深度则通过重投影与单目深度融合获得：先将相邻关键帧的像素反投影到 3D，重投影到非关键帧相机坐标系得到稀疏重投影深度 $D_t^\text{repro}$（过滤掉动态像素和无效视差点），然后通过最小二乘拟合将单目深度估计 $D_t^\text{mono}$ 对齐到重投影深度：

\begin{equation}
\theta_t, \gamma_t = \arg\min_{\theta, \gamma}\sum_{(u,v)}\bigl((\theta D_t^\text{mono}(u,v) + \gamma) - D_t^\text{repro}(u,v)\bigr)^2, \quad D_t^\text{aligned} = \theta_t D_t^\text{mono} + \gamma_t.
\end{equation}

在\textbf{动态场景表示}方面，ProDyG 将场景地图 $\mathcal{G} = \{\mathcal{G}_\text{s}, \mathcal{G}_\text{d}(t)\}$ 划分为静态高斯集合 $\mathcal{G}_\text{s}$ 和时变动态高斯集合 $\mathcal{G}_\text{d}(t)$。静态部分由 Splat-SLAM 代理深度图重建；动态部分采用 Motion Scaffolds（MoSca）表示。MoSca 是一个结构化图 $(\mathcal{V}_\text{d}, \mathcal{E}_\text{d})$，每个节点 $v^{(m)}$ 编码特定区域的运动轨迹，定义为 $v^{(m)} = ([\mathbf{Q}^{(m)}_1, \mathbf{Q}^{(m)}_2, \ldots, \mathbf{Q}^{(m)}_T], r^{(m)})$，其中 $\mathbf{Q}^{(m)}_t \in SE(3)$ 为每时间步的刚体变换，$r^{(m)}$ 为 RBF 影响半径。关键设计在于：动态高斯不定义在世界坐标系中，而是相对于 MoSca 节点定义。当 SLAM 后端产生新的位姿和深度更新时，MoSca 节点在 3D 空间中被重新锚定（re-anchored），动态高斯随之更新——这确保了在光度优化之前就已达到全局动态地图一致性，同时保证静态与动态部分之间的收敛性和一致性。

高斯从参考时间步 $t^\text{ref}_i$（初始化时间步）变换到查询时间步 $t$ 时使用双四元数混合（Dual Quaternion Blending, DQB）：

\begin{equation}
T_i(t) = \mathcal{W}(\mu_i, \mathbf{w}(\mu_i) + \Delta\mathbf{w}_i; t^\text{ref}_i, t) = \text{DQB}\bigl(\{w_m, \Delta\mathbf{Q}^{(m)}\}_{m \in \mathcal{E}_\text{d}(m^*)}\bigr),
\end{equation}

其中 $\mathbf{w}(\cdot)$ 为基于 RBF 的基础蒙皮权重（由 $\{r^{(m)}\}_{m \in \mathcal{E}_\text{d}(m^*)}$ 参数化，邻域 $\mathcal{E}_\text{d}(m^*)$ 包含最近邻 MoSca 节点 $v^{(m^*)}$ 及其所有连接节点），$\Delta\mathbf{w}_i$ 为可学习的逐高斯蒙皮权重修正，$\Delta\mathbf{Q}^{(m)} = \mathbf{Q}^{(m)}_t(\mathbf{Q}^{(m)}_{t^\text{ref}_i})^{-1}$ 为节点 $m$ 的相对变换。

在\textbf{MoSca 的渐进式构造}方面，SLAM 后端完成初始引导后，ProDyG 按以下流程初始化动态重建：（1）使用 CoTracker3 在精细运动掩码 $\{M_i\}$ 内生成稠密长期 2D 像素轨迹及逐时间步可见性标签；（2）将 2D 轨迹在可见时间步通过估计的相机位姿和对齐深度提升到 3D，不可见时间步则在附近观测之间线性插值，然后采样子集作为 MoSca 节点初始位置；（3）对 MoSca 节点进行几何优化，最小化 ARAP 损失、速度一致性损失和加速度一致性损失以推断不可见节点的旋转和位置；（4）在精细运动掩码内像素的反投影 3D 位置初始化动态高斯；（5）对 MoSca 和动态高斯进行联合光度优化，使用 RGB 损失 $\mathcal{L}_\text{rgb}$、深度损失 $\mathcal{L}_\text{depth}$（将渲染深度 $D^\text{r}$ 与对齐单目深度 $D^\text{aligned}$ 对齐）、轨迹损失 $\mathcal{L}_\text{track}$（确保与 2D 轨迹的运动一致性）以及 ARAP/速度/加速度几何正则化损失。此外，ProDyG 引入新颖的\textbf{运动掩码损失}，惩罚动态高斯在静态像素处的渲染不透明度：

\begin{equation}
\mathcal{L}_\text{mask} = \frac{1}{|\mathcal{P}^\text{static}_t|}\sum_{p \in \mathcal{P}^\text{static}_t} O(p),
\end{equation}

其中 $\mathcal{P}^\text{static}_t$ 为时间步 $t$ 的静态像素集（由运动掩码 $M_t$ 取反得到），$O(p)$ 为动态高斯在像素 $p$ 处的累积不透明度。该损失有效防止动态高斯溢出到静态区域，维持静态与动态组件之间的清晰分离——得益于此项约束，ProDyG 重建的动态物体轮廓比 Shape of Motion 和 MoSca 更为准确。深度渲染使用无偏 RaDe-GS 光栅化器，以射线-高斯交点深度替代 z-buffer 深度。

随着新帧到达，ProDyG 通过\textbf{多阶段点追踪策略}实现动态重建的增量扩展。对于每个新帧批次，CoTracker3 在包含新帧与已处理帧 8 帧重叠的时间窗口内运行：第一轮追踪识别在重叠帧中至少可见 4 帧的"近期可见"轨迹并延伸到新帧批次；第二轮通过球搜索（半径 $r_\text{search} = 0.02$ m）识别尚未被 3D 轨迹覆盖的"新可见"像素，专门针对这些区域运行第二次点追踪；第三轮在整个动态区域内补充追踪密度。所有延伸和新增的 2D 轨迹提升到 3D 后，利用已有 MoSca 的 DQB 变换将新轨迹反向变形到过去时间戳，确保时间一致性。仅在新可见像素处初始化新的动态高斯，然后对扩展后的 MoSca 和动态高斯执行与初始化阶段相同的联合几何-外观优化。

\textbf{实验结果。}在相机跟踪方面，ProDyG 在 Bonn RGB-D Dynamic 数据集和 TUM RGB-D 数据集（动态场景序列）上与最先进的动态 SLAM 方法取得可比性能。在 Bonn 数据集上，ProDyG 的 ATE RMSE 为 3.29 cm，显著优于基线 Splat-SLAM 的 10.6 cm，仅略逊于同时期工作 WildGS-SLAM（2.94 cm）。在 TUM 数据集上平均 ATE RMSE 为 1.89 cm，同样大幅优于 Splat-SLAM 的 2.43 cm。值得注意的是，WildGS-SLAM 仅重建静态背景，可激进地抑制动态区域；而 ProDyG 在保持高精度跟踪的同时还显式重建了一致的全局动态模型。

在新视角合成方面，ProDyG 在 iPhone 数据集上使用预计算位姿时（在线重建、离线位姿），PSNR/SSIM/LPIPS 达到 17.87/0.643/0.377，在 PSNR 和 SSIM 上均优于离线方法 Shape of Motion（17.43/0.591/0.303），并与共享相同运动表示的离线方法 MoSca（18.44/0.666/0.311）仅存在小幅差距——这证明了渐进式动态重建的有效性。当同时在线跟踪和重建时（最严格设定），PSNR 仍达 17.65。在纯 RGB 模式下，PSNR 为 15.41，是所有方法中唯一能在该设定下运行的。方法的局限性在于：对于主要在视锥外运动后重新出现的物体，MoSca 缺乏充分的光度约束，优化被正则化项主导，导致不理想的变形；此外，ProDyG 与大多数单目动态视角合成方法一样需要逐场景测试时优化，限制了实时应用场景。

\paragraph{不确定性驱动 4D 重建}

USplat4D（2025.10）~\cite{2510_12768} 指出从单目输入重建动态 3D 场景本质上是欠约束的：现有方法对所有高斯施加均匀的运动监督（深度、光流、光度一致性），忽略了有些高斯被多帧重复观测、约束充分，而有些高斯仅在少数帧中可见、约束薄弱。这种均匀处理导致运动估计在遮挡区域漂移，在新视角渲染时退化。USplat4D 提出一种模型无关的不确定性框架，核心思想是：当观测不完整时，重建应由高置信度的部分引导，并沿时空结构传播到不确定区域。

首先，USplat4D 推导了\textbf{逐高斯时变不确定性}的解析表达式。考虑体渲染中的颜色混合 $C^h = \sum_{i=1}^{N_g} T_i^h \alpha_i c_i$，其中 $v_i^h := T_i^h \alpha_i$ 为混合权重。在 $\mathcal{L}_2$ 损失 $\mathcal{L}_2 = \sum_h \|\hat{C}^h - C^h\|_2^2$ 下，对高斯 $k$ 的颜色参数 $c_k$ 进行最大似然估计（MLE）。令 $b_k^h := \hat{C}^h - \sum_{j \neq k} v_j^h c_j$ 隔离高斯 $k$ 对像素 $h$ 的贡献，通过求解 $\frac{\partial \mathcal{L}}{\partial c_k}\big|_{\text{MLE}} = 0$，得到最优颜色的闭式解：

\begin{equation}
c_k = \frac{\sum_{h} b_k^h \, v_k^h}{\sum_{h} (v_k^h)^2},
\end{equation}

及其方差（即在单位方差高斯噪声模型下，$c_k$ 的估计不确定性）：

\begin{equation}
\sigma_k^2 = \left(\sum_{h \in \mathcal{P}_{k}} (v_k^h)^2\right)^{-1},
\end{equation}

其中 $\mathcal{P}_{k,t} \subseteq \mathcal{P}_t$ 是高斯 $k$ 在帧 $t$ 中贡献的像素子集。该公式直觉上很清晰：一个高斯被越多的像素以越高的混合权重观测到，其不确定性越低。然而，MLE 的局部极小假设可能不成立（像素尚未收敛），因此引入收敛指示函数 $\mathbb{1}_t(h)$：当 $\|I_t^h - C_t^h\|_1 < \epsilon_c$ 时为 1，否则为 0。最终不确定性为 $u_{i,t} = \mathbb{1}_{i,t} \cdot \sigma_{i,t}^2 + (1 - \mathbb{1}_{i,t}) \cdot U$，其中 $U$ 为大常数，未收敛的高斯被赋予高不确定性。该设计联合调整收敛状态：被充分观测支持的高斯获得低 $u_{i,t}$（高可靠性），不可靠的高斯被赋予高值，从而在优化中使可信基元引导模糊基元。

进一步，考虑到单目设置中深度比图像平面坐标更不可靠，均匀的标量不确定性沿光轴方向会过度自信、导致几何畸变。USplat4D 将标量不确定性传播到 3D 空间，构建各向异性不确定性矩阵 $U_{i,t} = R_{wc} \, U_c \, R_{wc}^\top$，其中 $U_c = \mathrm{diag}(r_x u_{i,t},\, r_y u_{i,t},\, r_z u_{i,t})$，$R_{wc}$ 为相机到世界旋转，$r_x, r_y, r_z$ 为轴向缩放因子（深度方向的缩放更大），从而在深度方向赋予更高的不确定性。

基于逐高斯不确定性，USplat4D 构建\textbf{不确定性编码时空图} $G = (V, E)$，将高斯按可靠性排序并划分为关键节点集 $V_k$ 和非关键节点集 $V_n$。

\textbf{关键节点选择}采用两阶段策略平衡空间覆盖与时间稳定性：第一阶段，在每帧中将场景划分为 3D 体素网格，丢弃仅包含高不确定性高斯的体素，在剩余体素中随机选取一个低不确定性高斯——这种逐体素均匀选择确保空间覆盖并减少冗余；第二阶段，计算每个候选高斯的\textit{显著时段}（significance period），即不确定性低于阈值的最长连续帧数，仅保留显著时段至少 5 帧的候选作为关键节点，确保关键节点具有足够的时间支撑来可靠地贡献运动估计。

\textbf{图边构建}对关键节点和非关键节点分别处理。对于关键节点 $i$，在其最可靠帧 $t^* = \arg\min_t u_{i,t}$ 处，采用\textbf{不确定性感知 $k$NN}（UA-$k$NN）在关键节点集中选择邻居：

\begin{equation}
E_i = k\text{NN}_{j \in V_k \setminus \{i\}} \left(\|p_{i,t^*} - p_{j,t^*}\|_{(U_{w,t^*,i} + U_{w,t^*,j})}\right),
\end{equation}

其中使用 Mahalanobis 距离，对高不确定性方向施加更大权重，使边连接在空间上接近且可靠的关键节点。对于非关键节点 $i$，将其分配到整个序列中最近的关键节点 $j^* = \arg\min_{l \in V_k} \sum_{t=0}^{T-1} \|p_{i,t} - p_{l,t}\|_{U_{w,t,i} + U_{w,t,l}}$，并连接 $E_i = \{j^*\}$。直觉上，每个不确定的非关键节点被附着到与之长期保持接近的最可靠关键节点，使其运动可由稳定锚点正则化。

\textbf{不确定性感知优化}对关键节点和非关键节点设计分离的目标函数。

\textbf{关键节点损失}鼓励低不确定性的锚点高斯保持在预优化位置附近，同时施加运动正则化：

\begin{equation}
\mathcal{L}^\text{key} = \sum_{t=0}^{T-1} \sum_{i \in V_k} \|p_{i,t} - p_{i,t}^o\|_{U_{w,t,i}^{-1}} + \mathcal{L}^{\text{motion},\text{key}},
\end{equation}

其中 $U_{w,t,i}^{-1}$ 沿高不确定性方向降权，确保运动修正主要沿可靠轴进行；上标 $o$ 表示预训练高斯 Splatting 模型的参数作为初始化。运动正则化 $\mathcal{L}^{\text{motion},\text{key}}$ 包含五种约束的组合：等距损失 $\mathcal{L}^\text{iso}$（保持高斯对之间的距离不变）、刚性损失 $\mathcal{L}^\text{rigid}$（统一坐标系并约束相对运动）、旋转损失 $\mathcal{L}^\text{rot}$（显式约束四元数表示的相对旋转）、速度损失 $\mathcal{L}^\text{vel}$ 和加速度损失 $\mathcal{L}^\text{acc}$（约束世界空间中的平滑运动）。其中 $\mathcal{L}^\text{iso}$ 和 $\mathcal{L}^\text{rigid}$ 权重设为 1 以确保几何保持，旋转、速度和加速度损失权重设为 0.01 用于刚体方向约束和运动平滑。

\textbf{非关键节点损失}通过双四元数混合（Dual Quaternion Blending, DQB）从邻近关键节点插值运动，并正则化到初始化和插值轨迹：

\begin{equation}
\mathcal{L}^\text{non-key} = \sum_{t=0}^{T-1} \sum_{i \in V_n} \|p_{i,t} - p_{i,t}^o\|_{U_{w,i}^{-1}} + \sum_{t=0}^{T-1} \sum_{i \in V_n} \|p_{i,t} - p_{i,t}^\text{DQB}\|_{U_{w,i}^{-1}} + \mathcal{L}^{\text{motion},\text{non-key}},
\end{equation}

其中 $(p_{i,t}^\text{DQB}, q_{i,t}^\text{DQB}) = \text{DQB}(\{w_{ij} T_{j,t}\}_{j \in \mathcal{N}_i})$，$w_{ij}$ 为归一化边权重，$T_{j,t} \in SE(3)$ 为关键节点 $j$ 的变换，DQB 通过混合关键节点的变换为非关键节点提供平滑的运动传播。最终总损失为 $\mathcal{L}^\text{total} = \mathcal{L}^\text{rgb} + \mathcal{L}^\text{key} + \mathcal{L}^\text{non-key}$。不确定性在框架中服务于三个目的：(1) 重新加权关键节点的偏差；(2) 引导非关键节点的插值；(3) 平衡两者在总损失中的影响。

\textbf{实验结果。}USplat4D 在 DyCheck-iPhone 数据集（7 个真实动态场景）上以 SoM 和 MoSca 为基础模型进行评估。在标准验证视角下，USplat4D+MoSca 在所有指标上一致优于基线 MoSca。更具意义的是在极端新视角（大角度偏移）下的表现：基线方法在极端视角下出现严重的几何坍缩和扭曲，而 USplat4D 通过不确定性引导保持了稳定的几何结构。在 Objaverse 合成数据集（6 个关节化物体，水平视角偏移 $120^\circ$--$180^\circ$）上，USplat4D 在大视角偏移下的优势最为显著，一致超越 SoM 和 MoSca。消融实验表明：(1) 将关键节点选择中的不确定性替换为均匀 2D 采样导致图无法强调充分约束的高斯，锚点不稳定、传播退化；(2) 将 UA-$k$NN 替换为纯距离 $k$NN 削弱图连通性，边忽略节点可靠性而错误连接不稳定基元；(3) 移除训练损失中的不确定性加权（即去掉 $U$）降低 PSNR/SSIM，因为不可靠高斯被与可靠高斯同等力度更新，导致跨帧漂移；(4) 将 3D 体素网格采样替换为空间随机采样导致性能退化，因为非均匀空间分布既不保证覆盖也不避免冗余。此外，该模型天然支持多物体分割：关键图权重矩阵经社区检测（如谱聚类）后呈近似块对角结构，不同实例的运动节点形成独立簇，非关键节点根据非关键图连接被分配到同一运动群组，实现基于运动的动态实例分割与追踪。尤其在\textit{Camel}等序列中，USplat4D 避免了因深度方向不确定性被低估而导致的身体收缩变形。

\subsubsection{深度先验与几何引导}

本路线介绍引入单目深度估计器、特征匹配等几何先验来约束场景3D结构以弥补单目深度歧义的方法。

变形场与正则化策略提供了基本的重建框架，但单目输入的深度歧义使得变形场容易收敛到错误的 3D 结构。本节介绍引入深度先验和几何引导来弥补这一不足的方法。

单目设置中深度信息匮乏是核心瓶颈。一系列方法引入单目深度估计器、特征匹配或其他几何先验来约束场景的 3D 结构。

\paragraph{深度先验引导单目动态重建}

MoDGS（2024.06）~\cite{2406_00434}（被 ICLR 接收为 poster）针对随意拍摄的单目视频（相机静态或缓慢移动，多视角一致性约束极弱）提出利用深度先验的重建管线。此前的动态视角合成（Dynamic View Synthesis, DVS）方法普遍依赖\textit{teleporting camera motion}（相机在不同视角间大幅跳跃），即在不同时刻使用不同相机拍摄以构造伪多视角视频，从而为 3D 结构学习提供多视角一致性约束。这一要求由 DyCheck（2022）~\cite{gao2022monocular} 首次系统指出：D-NeRF、HyperNeRF、Nerfies 等常用单目 DVS 基准数据集虽然名义上是单目设置，但实际拍摄时相机在不同帧之间大幅移动、物体运动却很小，因此本质上仍是借助大幅相机运动获取多视角约束。然而，普通用户用手机随手拍摄的视频往往相机近乎静止，完全不满足这一条件。MoDGS 首次系统地解决了相机缓慢移动甚至完全静止条件下的单目动态重建问题，将 DVS 从实验室设置推向了随意拍摄的真实场景。

该方法的整体架构由三部分组成：在规范空间中定义高斯基元，使用可逆 MLP（invertible MLP）构建变形场 $T_t: \mathbb{R}^3 \to \mathbb{R}^3$（时间 $t$ 归一化到 $[0,1]$），以及通过 Splatting 渲染图像。可逆 MLP 的设计借鉴了 Omnimotion 和 CaDeX，使 $T_t$ 和 $T_t^{-1}$ 都可直接从网络计算，无需额外求解。值得注意的是，规范空间并不显式对应任何特定时刻，而是一个容纳所有高斯基元规范位置的虚拟空间。

MoDGS 的核心贡献在于两个技术组件。

第一，\textbf{3D 感知初始化方案}：原始 3DGS 依赖 SfM 稀疏点云初始化高斯位置，但随意拍摄的单目视频难以通过 SfM 获得可靠的初始点集。朴素地用第一帧的估计深度初始化高斯同样效果不佳，因为不同帧的单目深度图尺度不一致且变形场未初始化。MoDGS 的初始化分三步进行：

（1）\textit{深度尺度统一}——每帧先用 GeoWizard 单目深度估计器估计深度图 $D_t$，但不同帧的估计深度具有不同尺度。MoDGS 先用 RAFT 光流或 SAM2 分割出静态区域，然后在静态区域上将每帧深度重投影到第一帧，通过最小二乘拟合求解每帧的尺度因子，统一所有深度图的尺度。

（2）\textit{变形场初始化}——在均匀采样的关键帧之间（例如每隔 5 帧采样一对），将两帧的深度图和 2D 光流提升为 3D 流：先将深度图反投影为 3D 点，再由 2D 光流建立对应关系从而得到 3D 流 $\mathbf{F}_{t_i \to t_j}^{3D}$。然后训练变形场 $T$ 以最小化
\begin{equation}
\mathcal{L}_\text{init} = \|T_{t_j} \circ T_{t_i}^{-1}(x_{t_i}) - x_{t_j}\|^2,
\end{equation}

即先将 $t_i$ 时刻的 3D 点 $x_{t_i}$ 通过逆变形场映射到规范空间，再正向变形到 $t_j$ 时刻，与 $t_j$ 的目标点对齐。这使变形场在正式优化前即已捕获大致的 3D 运动。实现中变形场初始化训练 20k 步，变形网络学习率从 $10^{-3}$ 指数衰减到 $10^{-4}$。

（3）\textit{高斯初始化}——将所有帧的深度点通过已初始化的变形场反向变换到规范空间（场景归一化到 $[-1,1]^3$），然后以预定义体素大小（$0.004^3$）均匀下采样得到高斯初始位置。这一步将所有时刻的 3D 点汇聚到统一规范空间，为后续联合优化提供密集且尺度一致的初始化。

第二，\textbf{序数深度损失}（Ordinal Depth Loss）：此前方法使用 Pearson 相关损失来处理单目深度的尺度歧义。Pearson 相关损失最大化渲染深度与估计单目深度之间的相关性：

\begin{equation}
\mathcal{L}_\text{Pearson}(\hat{D}_t, D_t) = -\frac{\text{Cov}(\hat{D}_t, D_t)}{\sqrt{\text{Var}(\hat{D}_t) \cdot \text{Var}(D_t)}},
\end{equation}

该损失通过归一化避免了尺度歧义的负面影响，但隐含假设不同帧的估计深度图经归一化后满足线性关系（$D_{t+1} = aD_t + b$）。然而，由于单目深度估计器精度有限，这一线性假设并不成立，即使归一化后不同帧的深度图仍然不一致，导致 Pearson 相关损失产生矛盾的监督信号。

MoDGS 的关键观察是：尽管不同帧的绝对深度值不一致，但同一场景中像素间的\textit{深度顺序}是稳定的——即若某帧中像素 $u_1$ 的深度大于 $u_2$，那么在其他帧中这一相对关系通常也成立。基于此，定义序数指示函数
\begin{equation}
R(D_t(u_1), D_t(u_2)) = \begin{cases} +1, & D_t(u_1) > D_t(u_2) \\ -1, & D_t(u_1) < D_t(u_2) \end{cases},
\end{equation}

然后强制渲染深度图的深度顺序与估计深度图一致：

\begin{equation}
\mathcal{L}_\text{ord} = \tanh\bigl(\eta(\hat{D}_t(u_1) - \hat{D}_t(u_2))\bigr) - R\bigl(D_t(u_1), D_t(u_2)\bigr),
\end{equation}

其中 $\tanh(x) = \frac{e^x - e^{-x}}{e^x + e^{-x}}$ 将深度差映射到 $[-1, 1]$，$\eta$ 为预定义常数（实现中设为 $\eta = 100$）。实现中随机采样 100k 个像素对 $(u_1, u_2)$ 计算该损失，且深度先归一化到 $[0,1]$，仅考虑深度差大于 0.02 的像素对以过滤噪声。总训练损失为 $\mathcal{L} = \lambda_\text{rgb}\mathcal{L}_\text{rgb} + \lambda_\text{ord}\mathcal{L}_\text{ord}$，其中 $\lambda_\text{ord} = 0.1$。

实验在四个数据集上进行系统评估。（1）DyNeRF 数据集包含 6 个场景（桌面操作场景如切牛肉、倒水等），每个场景有 18--20 台同步相机拍摄 10--30 秒视频，仅用 camera0 单目训练，在 camera5 和 camera6 上评估。（2）NVIDIA Dynamic 数据集包含 8 个更丰富的动态场景（跳跃、玩气球、打伞等），每个场景有 12 台同步相机，仅用 camera4 训练，在 camera3 和 camera5 上评估。（3）作者还收集了 6 段在线视频构建 Monocular Casual Video（MCV）数据集（滑冰、狗进食、瑜伽等），用于定性评估野外泛化能力。（4）DAVIS 数据集用于补充评估。需要指出的是，作者没有使用 DyCheck 所采用的 iPhone 数据集，原因是经过验证该数据集的相机位姿精度不足（通过对极几何验证发现：标注位姿计算的对极线偏离实际对应点较远，且尝试用 COLMAP 重跑也无法获得合理结果）。由于单目输入存在尺度歧义，评估时手动标注训练图像与真实新视角图像之间的对应点，通过优化深度尺度因子对齐后再计算 PSNR、SSIM 和 LPIPS。

在定量对比方面，MoDGS 在 DyNeRF 和 NVIDIA Dynamic 数据集的所有场景和所有指标上均超越基线方法，包括 NeRF 类方法（HexPlane、RoDynRF）和 3DGS 类方法（Deformable-GS、SC-GS）。值得强调的是，HexPlane 在纯单目输入下无法产生合理结果，需要额外视角辅助训练，而 MoDGS 完全不依赖额外视角。RoDynRF 虽然也使用了单目深度估计器 GeoWizard 进行深度监督，但其 Pearson 相关深度损失在帧间深度不一致时效果不佳。基线方法失败的主要原因在于：单目相机几乎静止，无法提供足够的多视角一致性来重建高质量 3D 几何；SC-GS 的静态初始化在前景大幅运动时会忽略动态区域（例如人物从画面左侧滑到右侧时，静态初始化会将动态前景视为背景噪声丢弃），而 MoDGS 的 3D 感知初始化从第一步就考虑了动态性。在并发工作方面，Shape-of-Motion 和 Mosca 也探索了类似问题，但它们依赖显式的像素长程对应关系（视频追踪先验），而 MoDGS 仅使用 RAFT 估计的帧间光流作为输入，无需强假设的长期像素对应。DG-Marbles 采用高斯弹珠和分层学习策略，同样面向随意拍摄的单目视频。

消融实验进一步验证了两个核心组件的有效性。在 3D 感知初始化方面，将变形场替换为随机初始化后（仍保留高斯从深度点反向变换到规范空间的初始化），动态前景人物出现明显伪影，验证了 3D 流引导的初始化为后续优化提供了关键的良好起点。在序数深度损失方面，三个设置进行了对比：无深度损失时重建的 3D 几何包含大量噪声和伪影；使用 Pearson 相关损失有所改善但内部区域深度仍然嘈杂；使用序数深度损失在保持边界锐利的同时实现了内部区域的平滑重建。此外，论文还展示了 MoDGS 与简单深度图变形（depth warping）的本质区别：深度图变形仅对单帧深度进行 2D 投影变换，无法聚合不同时刻的信息；而 MoDGS 通过规范空间中的高斯和变形场，能够跨时刻累积 3D 信息，从而重建出比任何单帧深度估计都更完整的场景（尤其在部分可见区域）。

训练效率方面，完整训练过程在 NVIDIA RTX A6000 GPU 上约需 3.5 小时（2 小时初始化 + 1.5 小时联合优化），其中初始化训练 20k 步，联合优化训练 20k 步，使用约 14 GB 显存，渲染速度约 75 FPS。变形网络学习率在联合优化阶段从 $10^{-4}$ 指数衰减到 $10^{-6}$，高斯部分的学习率完全沿用原始 3DGS 的设置。深度渲染使用高斯中心点坐标（而非 $\alpha$-混合深度），这是 Deformable-GS 中引入的做法。

MoDGS 的局限性也值得关注。首先，该方法只能重建可见的 3D 部分，无法想象被遮挡的区域，在渲染包含未见部分的新视角时会产生伪影——引入 3D 相关扩散生成模型（如 LucidDreamer、SyncDreamer）可能是一个有前景的方向。其次，当相机完全静止时，方法高度依赖单目深度估计器的质量，而现有估计器（如 GeoWizard、Depth Anything、ZoeDepth）虽在大规模数据集上训练并在多数场景下预测合理，但仍可能遗漏细节。此外，序数深度损失假设不同帧间的深度顺序一致，尽管这一假设在主流深度估计器中普遍成立，但在极端情况下可能不满足。快速运动场景会导致相机位姿估计不准确，高光反射和低光照条件也是单目设置下尚未解决的难题。训练效率方面，每场景约需数小时的训练时间仍是该方法的瓶颈，如何高效重建动态场是一个值得探索的未来方向。

\paragraph{高斯泼溅与神经表面联合动态重建}

现有动态场景重建方法往往只能在动态新视角合成和 3D 几何重建两个任务之一上表现优异，难以兼顾。DGNS（2024.12）~\cite{2412_03910} 的核心动机来自静态场景建模中的观察：鲁棒的几何引导能够提升而非损害渲染质量。基于此，该方法提出混合表示框架，将变形高斯 Splatting（DGS）模块与动态神经表面（DNS）模块统一优化，两个模块相互增益。DGS 模块基于 D3DGS 框架，使用 8 层 MLP（256 通道）学习变形场，主要负责外观重建。DNS 模块在规范空间（canonical space）中使用符号距离函数（SDF）表示动态表面，规范空间采用哈希网格编码与 4 层 MLP 的混合表示以加速训练；变形场使用双射映射网络（bijective mapping network），将动态观测空间映射到规范空间，并严格保持可逆性以实现帧间的循环一致性约束。训练时两个模块通过三条路径相互促进。

第一，DGS 渲染的 $\alpha$-混合深度图 $d^{\alpha} = \sum_{i}^{N} d_i \alpha_i \prod_{j=1}^{i-1}(1-\alpha_j) / \sum_{i}^{N} \alpha_i \prod_{j=1}^{i-1}(1-\alpha_j)$ 用于约束 DNS 模块的光线采样范围：射线从相机中心 $\vec{\mathbf{o}}$ 沿方向 $\vec{\mathbf{v}}$ 发射，采样点取在 $\vec{\mathbf{o}} + d^{\alpha}\cdot\vec{\mathbf{v}}$ 附近，采样区间为 $[\vec{\mathbf{o}} + (d^{\alpha} - s|d|)\cdot\vec{\mathbf{v}},\; \vec{\mathbf{o}} + (d^{\alpha} + s|d|)\cdot\vec{\mathbf{v}}]$，其中 $d = \mathcal{F}(\mathcal{H}(\vec{\mathbf{o}} + d^{\alpha}\cdot\vec{\mathbf{v}}, t))$ 为该点经变形场映射后在规范空间中的 SDF 值，$s$ 为缩放因子（训练前期设为 3 进行粗搜索，后期设为 1 进行精搜索），区间由 SDF 值自适应缩放，避免在空白区域浪费计算。

第二，针对 $\alpha$-混合深度在物体边缘区域噪声较大的问题，DGNS 提出深度滤波方法：同时计算中值深度 $d^{m}$（累积透射率 $T_k = \prod_{j=1}^{k-1}(1-\alpha_j)$ 降至阈值 $\tau_d=0.6$ 时对应的高斯基元中心深度），仅当 $|d^{\alpha} - d^{m}| < \tau_f$ 时才将滤波深度 $d^f = (d^{\alpha} - d^{m})/2$ 作为可靠的深度监督信号，不满足条件的位置深度置零（即不施加监督），滤波后的深度点用于 SDF 损失 $\mathcal{L}_{\text{sdf}} = \sum_{d^f \in \mathcal{D}} \|\mathcal{F}(\mathcal{H}(\vec{\mathbf{o}} + d^f\cdot\vec{\mathbf{v}}, t))\|_1$。消融实验证实，滤波深度优于直接使用 $\alpha$-混合深度（在 Bird 场景上 Chamfer Distance 从 0.746 降至 0.697）。

第三，DNS 学习到的 SDF 零水平集用于 DGS 模块的表面感知密度控制（surface-aware density control）。对于 DGS 规范空间中的每个高斯基元 $\mathbf{x}_g$，通过变形场确定其在观测空间中的位置 $\mathbf{x}_g + \delta\mathbf{x}_g$，并计算其 SDF 距离 $d_g = \mathcal{F}(\mathcal{H}(\mathbf{x}_g + \delta\mathbf{x}_g, t))$。高斯基元的生长准则定义为 $\epsilon_g = \nabla_{\mathbf{x}_g} + w_g \cdot \phi(d_g)$，其中 $\nabla_{\mathbf{x}_g}$ 为位置的平均梯度，$w_g$ 为几何因素权重，$\phi(x) = \exp(-x^2/2\sigma^2)$ 与 SDF 值成反比——$\phi$ 项使靠近表面的高斯（$d_g$ 小，$\phi$ 大）更容易被分裂或克隆。当 $\epsilon_g$ 超过阈值 $\tau_g$ 时触发高斯生长。剪枝准则则定义为 $\epsilon_p = \sigma_p - w_p(1-\phi(d_g))$，其中 $\sigma_p$ 为 $K$ 次迭代中不透明度之和，$w_p$ 为权重参数——远离表面的高斯（$d_g$ 大，$\phi$ 小，$1-\phi$ 大）使 $\epsilon_p$ 降低，更容易被剪枝。消融实验表明，去除表面感知密度控制后，高斯点更分散地远离表面，导致浮点伪影。

两个模块均引入来自基础模型 Marigold 的单目法线监督。在 DNS 模块中，体渲染法线 $\hat{N}$ 是射线上各采样点法线 $\nabla_{\mathbf{x}_o}\mathcal{F}(\mathcal{H}(\mathbf{x}_o, t))$ 的加权和，通过角度损失和 $L_1$ 损失 $\mathcal{L}_{\text{dn}}^{n} = \frac{1}{N}\sum_{\mathbf{r} \in \mathcal{R}} \|\hat{N}(\mathbf{r}) - \bar{N}(\mathbf{r})\|_1 + \|1 - \hat{N}(\mathbf{r})^\top \bar{N}(\mathbf{r})\|_1$ 与 Marigold 预测的法线图 $\bar{N}$ 对齐。在 DGS 模块中，每个高斯基元的法线通过最小缩放轴方向近似：$\mathbf{n} = \mathbf{R}[k,:] \in \mathbb{R}^3, k = \arg\min([s_1, s_2, s_3])$，其中 $s_1, s_2, s_3$ 为高斯的三个尺度分量，$\mathbf{R}$ 为旋转矩阵；渲染法线 $\hat{N}^g(p) = \sum_{i \in N} \mathbf{n}_i \alpha_i \prod_{j=1}^{i-1}(1-\alpha_j)$ 以同样的角度损失和 $L_1$ 损失正则化。DNS 模块还额外引入 Eikonal 正则化 $\mathcal{L}_{\text{eik}}$ 约束 SDF 梯度的单位长度性质。

训练采用分阶段策略：DGS 模块先独立预热约 10k 次迭代以建立准确的深度预测，此后深度引导和 SDF 监督开始生效；15k 次迭代后，表面感知密度控制开始介入 DGS 的密度管理。两个模块从 15k 迭代开始联合优化，直至 40k 迭代收敛。DGS 模块的图像损失权重 $\lambda_I$ 设为 0.8，法线损失权重 $\lambda_{gn}$ 从 10k 迭代开始生效（值为 0.1）；变形网络在 3k 迭代后激活，密度控制从 500 次迭代开始。最终总损失 $\mathcal{L} = \mathcal{L}_{\text{dg}} + \mathcal{L}_{\text{dn}}$ 联合优化两个模块，其中 DGS 损失 $\mathcal{L}_{\text{dg}} = \lambda_I \mathcal{L}_1 + (1-\lambda_I)\mathcal{L}_{\text{ssim}} + \lambda_{gn}\mathcal{L}_{\text{dg}}^n$，DNS 损失 $\mathcal{L}_{\text{dn}} = \mathcal{L}_1 + \lambda_{\text{sdf}}\mathcal{L}_{\text{sdf}} + \lambda_{nn}\mathcal{L}_{\text{dn}}^n + \lambda_{\text{eik}}\mathcal{L}_{\text{eik}}$。

在 Dg-mesh 数据集上，DGNS 在所有 6 个动态物体上取得了最低的 Chamfer Distance 和 Earth Mover Distance，同时 PSNR 与 D3DGS 持平（如 Duck 场景 37.39 vs 36.76，Horse 场景 35.52 vs 35.66，Bird 场景 30.23 vs 29.38），在 3D 几何重建精度上全面领先 Dg-mesh、D3DGS、NDR 等基线方法。在 D-NeRF 数据集上，DGNS 的渲染质量与 D3DGS 相当，同时提供高质量的 3D 网格重建。该方法的局限性在于 DNS 模块仍需大量计算资源和较长收敛时间，成为整体框架的速度瓶颈；此外，同时训练两个模块增加了显存开销。该论文发表于 ACM Multimedia 2025。

\paragraph{几何一致性稀疏输入动态重建}

GC-4DGS（2025.11）~\cite{2511_23044} 针对稀疏多视角（少至 3 个视角）输入下 4DGS 质量严重退化的问题，基于 4DGS 框架引入两层几何一致性约束。该问题的现实动机来自 AIoT 场景：在智能交通等系统中，成本与隐私约束往往将摄像头限制在少数固定位置，带宽限制和丢包也可能减少可用视角，因此稀疏输入下的动态重建能力至关重要。

第一层为多视角一致重建。该方法利用 MVSFormer（一种具有出色泛化能力的基于学习的 MVS 方法）进行几何初始化与后续结构监督：将所有训练视频按时间切片，以每个视角为参考视角、其余为邻域视角，构建代价体并预测逐像素度量深度流 $D_{\text{mvs}}^{t,n}$。由于可泛化 MVS 在视觉重叠低或噪声显著时估计不准，GC-4DGS 引入动态一致性检查（Dynamic Consistency Checking）策略：对参考视角 $n$ 在时间 $t$ 的像素 $\boldsymbol{p}$，将其投影至邻域视角得到 $\boldsymbol{p}'$，再反投影回参考视角得到 $\boldsymbol{p}''$，分别计算像素空间重投影误差 $\xi_p^{t,n} = \|\boldsymbol{p} - \boldsymbol{p}''\|_2$ 和深度空间重投影误差 $\xi_d^{t,n} = \|D^{t,n}(\boldsymbol{p}) - D^{t,n}(\boldsymbol{p}'')\|_1 / D^{t,n}(\boldsymbol{p})$。为在时空维度动态聚合几何匹配代价，定义动态几何一致性 $c(\boldsymbol{p}) = \frac{1}{N_f}\sum_{t=1}^{N_f}\sum_{n=1}^{N_c}e^{-(\xi_p^{t,n}+\beta\xi_d^{t,n})}$，其中 $\beta$ 为平衡两个重投影度量（实现中设为 $\beta=200$）的参数。对每个像素计算一致性得分后，滤除低于预设阈值（1.8）的不可靠像素，生成几何一致性掩码；再将其与 MVS 模型输出的光度概率图融合，得到最终的动态一致性掩码 $M^{t,n}$。滤波后的多视角度量深度用于 4D 高斯点云初始化（实现中仅使用第一帧深度进行融合）和结构监督损失 $\mathcal{L}_{\text{stru}} = \sum \text{smooth}_{L_1}((D_{\text{ren}}-D_{\text{mvs}}) \odot M)$，其中渲染深度 $D_{\text{ren}}$ 通过可微深度光栅化器按 $\alpha$-混合公式 $\mathcal{D}(t) = \sum_{i=1}^{N}d_i(t)\alpha_i(t)\prod_{j=1}^{i-1}(1-\alpha_j(t))$ 获得，$\text{smooth}_{L_1}$ 为促进零点附近收敛的鲁棒损失函数。

第二层为全局-局部深度正则化，用于补充 MVS 监督不足区域。该方法的关键洞察是：尽管 MDE（实现中使用 DepthAnythingV2）估计的深度缺乏帧间尺度一致性——即使归一化后，相同像素对的深度比在不同帧和视角间仍会变化——但像素间的相对深度排序应当保持一致。据此提出全局排序损失：用 $\mathcal{R}_{\text{ren}}^{t,n} = \text{Sigmoid}(\kappa(D_{\text{ren}}^{t,n}(\boldsymbol{u}) - D_{\text{ren}}^{t,n}(\boldsymbol{u}')))$ 将渲染深度的排序推向 0 或 1（$\kappa$ 为正常数，映射深度差至正/负无穷），并以 MDE 排序指示函数 $\mathcal{R}_{\text{mde}}^{t,n} \in \{0, 1\}$（根据 $D_{\text{mde}}(\boldsymbol{u})$ 与 $D_{\text{mde}}(\boldsymbol{u}')$ 的相对大小判定）为伪标签，定义全局一致排序损失 $\mathcal{L}_{\text{rank}} = \sum_{(\boldsymbol{u},\boldsymbol{u}')} \|\mathcal{R}_{\text{ren}}^{t,n} - \mathcal{R}_{\text{mde}}^{t,n}\|$。实现中每次训练迭代从全视频分辨率随机采样 500k 个像素对计算该损失，使其对 MDE 的尺度不一致先验具有鲁棒性。在局部层面，通过分块归一化保证平滑性：将深度图分割为小块 $\boldsymbol{e}$，对块内每个像素执行局部归一化 $\hat{D}(\boldsymbol{p}) = (D(\boldsymbol{p}) - \text{mean}(D(\boldsymbol{e}))) / (\text{std}(D(\boldsymbol{e})) + \delta)$（$\delta=2\times10^{-4}$ 防止数值不稳定），然后在归一化后的渲染深度与 MDE 深度上施加 $L_2$ 损失 $\mathcal{L}_{\text{patch}} = \sum_{\boldsymbol{e}}\|\hat{D}_{\text{ren}}^{t,n}(\boldsymbol{e}) - \hat{D}_{\text{mde}}^{t,n}(\boldsymbol{e})\|$。块大小根据场景几何复杂度（全局深度范围和局部深度变化）从预定义范围采样（N3DV 数据集为 $[8,16]$，Technicolor 数据集为 $[16,32]$）。

最终总损失 $\mathcal{L}_{\text{total}} = \mathcal{L}_{\text{photo}} + \lambda_1\mathcal{L}_{\text{rank}} + \lambda_2\mathcal{L}_{\text{patch}} + \lambda_3\mathcal{L}_{\text{stru}}$，其中 $\mathcal{L}_{\text{photo}}$ 为 $L_1$ 与 D-SSIM 的组合光度损失，超参数设为 $\lambda_1=0.05$, $\lambda_2=0.02$, $\lambda_3=0.02$。在 N3DV 数据集（3 视角输入）上，GC-4DGS 取得 23.96 dB PSNR，比专为稀疏输入设计的 RF-DeRF 高 2.62 dB，比原始 4DGS 高 1.58 dB，同时保持 115 FPS 实时渲染速度和 1.0h 训练时间（RTX 4090）。消融实验表明三个约束项互补：$\mathcal{L}_{\text{stru}}$ 利用高置信度量深度纠正高斯表示的几何误差，$\mathcal{L}_{\text{rank}}$ 在 MVS 覆盖不足区域通过排序一致性提供全局几何约束（为主导项），$\mathcal{L}_{\text{patch}}$ 保证局部平滑以消除块级伪影。此外，GC-4DGS 可直接部署于资源受限的 IoT 边缘设备（Jetson AGX Orin，1.3 GHz Ampere GPU + 2.2 GHz ARM Cortex-A78AE 处理器），展示了在 AIoT 应用（如数字孪生、远程呈现）中的实际部署潜力。

\paragraph{越野场景空间锚定前馈 4D 重建}

前馈高斯 Splatting（FFGS）在结构化城市场景中已取得成功，但在非结构化越野场景中面临三重叠加的挑战：密集植被产生的高频几何、摇曳树冠和飘散粒子带来的空间弥漫性非刚性动态、以及崎岖地面引起的连续自运动抖动。高频几何放大了抖动敏感度，使细粒度观测错位；弥漫动力学进一步向这些错位信号注入冲突外观，三者互相放大。这些因素在规范高斯空间中引发两类属性冲突：空间偏移冲突来自抖动导致的表面观测散射，时间变化冲突来自弥漫动态导致的同一位置变形状态不一致。现有 FFGS 方法在冲突下呈现两种失败模式：朴素融合（如 NopoSplat）忽略时间相关性，将冲突属性平均为模糊渲染；逐高斯时序建模（如 DGGT）独立评分每个高斯而无空间参照，当查询时间远离所有上下文帧时，置信度均匀衰减导致不透明度坍缩为结构空洞。两者的共同根源在于时间相关性在空间上下文之外独立求解，迫使时间选择性和空间占用陷入固有矛盾。

Ground4D（2026.05）~\cite{2605_04435} 的关键洞察是将时间竞争限制在局部空间邻域内，消除时间选择性与空间占据之间的矛盾。在局部邻域内，时间最主导的高斯同时达到最高时间相关性和保证空间占用，两个目标互相强化而非矛盾。该方法利用冻结的 VGGT 骨干网络处理 $T$ 帧上下文图像，产生三路互补的 token 流：聚合跨视角 token $F_t^{\text{agg}}$、图像级 token $F_t^{\text{img}}$ 和 DINOv2 语义 token $F_t^{\text{dino}}$。骨干网络在微调时保持冻结，仅训练预测头和时序融合模块。相机头从 $F_t^{\text{agg}}$ 解码相机位姿，DPT 深度头预测逐像素深度 $D_t$，通过反投影 $\mathcal{P}_t(u,v) = R_t^\top(K_t^{-1}D_t(u,v)[u,v,1]^\top - T_t)$ 获得 3D 点作为高斯基元中心。高斯头在 $F_t^{\text{img}}$ 和 $F_t^{\text{dino}}$ 上预测 11 通道的密集高斯属性图（颜色、不透明度、尺度和旋转），将像素反投影至共享规范高斯空间 $\mathcal{G}^s = \bigcup_t \mathcal{G}_t^s$。

体素锚定的时序高斯聚合分三步。首先，将规范空间划分为体素边长 $\rho$ 的均匀网格，每个高斯基元中心 $\mu_i$ 通过 round 函数量化为整数坐标 $v_i = \text{round}(\mu_i / \rho) \in \mathbb{Z}^3$，共面高斯归入同一体素，共产生 $M \ll N$ 个非空稀疏体素，将全部基元划分为可变大小组。体素化步骤不可微，仅用于分组，梯度通过聚合阶段的高斯属性回传。然后，对每个体素 $m$ 内的共面高斯执行查询条件的时序注意力：每个基元携带特征 $f_i^{\text{gs}}$（高斯属性先编码为特征嵌入，再与高斯头最后一层的残差特征拼接），经特征编码器 $\phi_f$ 投影至共享隐空间；时间戳 $\tau_i$ 经多频率正弦编码 $\gamma(\cdot)$ 和时间 MLP $\phi_t$ 在多尺度上捕捉时间动态，二者加性融合为 $h_i = \phi_f(f_i^{\text{gs}}) + \phi_t(\gamma(\tau_i))$；查询时间 $\tau^*$ 类似编码为 $h^* = \phi_t(\gamma(\tau^*))$，轻量注意力 MLP 输出标量对数 $a_i = \text{MLP}_{\text{attn}}([h_i; h^*])$。最后，在每个体素内独立进行温度缩放 softmax 归一化 $w_i = \exp(a_i/\beta) / \sum_{j \in \mathcal{V}_m}\exp(a_j/\beta)$，确保每个非空体素无论查询时间距上下文帧多远都能产生有效基元。将 softmax 限制在体素内部赋予两个互补特性：（i）时间选择性——局部最主导的高斯驱动融合输出并抑制冲突观测；（ii）空间占用——由于归一化限于每个体素，每个非空体素保证产生有效基元，确保结构完整性。不同属性采用语义感知融合：位置和颜色作为欧氏量用加权算术均值；不透明度将加权均值与体素内最大值按混合系数 $\lambda$ 混合（避免瞬态遮挡体素的密度坍缩）；尺度在对数空间加权几何均值；旋转在单位四元数流形上加权均值后 $L_2$ 重归一化。

表面法线通过双路径设计提供几何先验。法线头解码 $F_t^{\text{img}}$ 和 $F_t^{\text{dino}}$ 为逐像素单位法线图 $\mathbf{n}_t$，其中间特征注入高斯头直接约束尺度和旋转预测。高斯衍生法线通过 soft-argmin 从三个尺度轴估计：$\tilde{\mathbf{n}} = R_\mathbf{q} \cdot \text{softmax}(-\eta \mathbf{s}) / \|\text{softmax}(-\eta \mathbf{s})\|_2$，最小尺度轴获得最高权重，标识盘状高斯的法线方向。法线监督包含两项余弦距离损失：预测法线损失 $\mathcal{L}_{\text{norm}}^{\text{pred}}$ 直接监督法线头输出，高斯衍生法线损失 $\mathcal{L}_{\text{norm}}^{\text{gs}}$ 将估计法线投影至相机坐标系后与真值对齐，两者耦合的几何感知信号强正则化高斯朝向。

显式动态物体（车辆、行人等）由动态置信头 $\Pi_t^{\text{dyn}}$ 产生逐像素动态分数检测，每个高斯不透明度按 $(1-\Pi_t^{\text{dyn}})$ 加权抑制动态区域，动态高斯由 TAPIP3D 在关键帧间跟踪稀疏 3D 点并线性插值中心至查询时间后与精化高斯合成。弥漫性的高频非刚性变形无法由此显式方式解决，而是由体素锚定时序注意力聚合隐式处理。天空背景由参数化模型 $\mathcal{M}^{\text{sky}}$ 以视角方向为条件回归，最终合成为 $\hat{I}_t = \mathbf{A}_t \odot \hat{I}_t^{3D} + (1-\mathbf{A}_t) \odot \mathcal{M}^{\text{sky}}(\mathcal{I}, E_t, K_t)$，其中 $\mathbf{A}_t$ 为 3DGS 光栅化的累积 $\alpha$。

训练端到端监督，总损失为六项之和 $\mathcal{L} = \mathcal{L}_{\text{rgb}} + \mathcal{L}_{\text{lpips}} + \mathcal{L}_{\text{dyn}} + \mathcal{L}_{\text{sky}} + \mathcal{L}_{\text{norm}}^{\text{pred}} + \mathcal{L}_{\text{norm}}^{\text{gs}}$：$L_1$ 光度损失 $\mathcal{L}_{\text{rgb}}$、LPIPS 感知损失 $\mathcal{L}_{\text{lpips}}$、动态置信 BCE 损失 $\mathcal{L}_{\text{dyn}}$、天空 $\alpha$ 监督损失 $\mathcal{L}_{\text{sky}}$，以及前述两项法线损失。训练时引入随机帧丢弃（概率 $p_{\text{drop}} = 0.7$）——查询时间 $\tau^*$ 的帧以概率 $p_{\text{drop}}$ 从体素化空间中移除，但其渲染输出保留为光度监督目标——迫使网络从时间不连续的上下文推断场景状态，正则化对瞬态变形的过拟合。关键超参数包括体素尺寸 $\rho = 0.002$、时间编码 10 个频率带宽隐维度 64、不透明度混合系数 $\lambda = 0.3$、可学习温度 $\beta$ 初始化为 1.0、softmin 温度 $\eta = 10.0$。

消融实验揭示了关键设计原则。体素锚定聚合的三个组件缺一不可：仅加时序注意力而无局部约束反而降低性能（因为每个高斯独立评分时间相关性，无法在空间区域内强制主导选择，导致空间占用坍缩为结构空洞）；仅体素化而无时序注意力缺乏分辨力，无法解决快速变形下的属性歧义；体素化加时序注意力但无体素内归一化同样无效，因缺少局部竞争的参数化。三者结合将 spatio-temporal 交互从竞争变为协调：时序注意力识别最优时间状态，体素内归一化保持几何完整性。体素尺寸 $\rho$ 需平衡体素内竞争强度与空间连贯性：最优值 $\rho = 0.002$ 将规范高斯从约 600K 压缩至 370K（1.6 倍减少），更小体素（$\rho = 0.001$）使共面高斯落入单一体素内竞争变得平凡，更大体素（$\rho = 0.005$）过度合并空间不同表面导致 2.85 dB PSNR 暴跌。上下文帧数在 10 帧达到峰值后饱和，进一步增加时间窗引入不同变形状态的观测，超出体素内聚合的协调能力。法线双路径设计在光度模糊的越野区域提供不可替代的几何约束，去除法线损失后 PSNR 从 23.89 降至 23.45。该方法在 ORAD-3D 上实现最高 1.48 dB PSNR 增益（23.89 vs 22.41），并零样本泛化至 RELLIS-3D，无需微调即在 SSIM 和 LPIPS 上达到所有方法最优。

\paragraph{ToF 间接深度优化动态辐射场}

连续波飞行时间（C-ToF）相机通过主动照明为单相机动态场景重建提供深度估计。C-ToF 相机以调制频率 $f$ 发射连续正弦光 $\sin(2\pi f t)$，光返回后与参考正弦信号进行相关运算，在四个相位偏移 $\phi \in \{0, \frac{\pi}{2}, \pi, \frac{3\pi}{2}\}$ 下采集原始四分量帧 $\{Q_0, Q_{\pi/2}, Q_\pi, Q_{3\pi/2}\}$，从中恢复深度 $d_{\text{ToF}} = \frac{\mathrm{c}}{4\pi f}\arctan\!\left(\frac{Q_0 - Q_\pi}{Q_{\pi/2} - Q_{3\pi/2}}\right)$。在透射式辐射场模型下，深度并非直接优化量——传感器测量值是各高斯基元信号的加权叠加，而深度仅通过重建这些测量值被间接推断。由于正弦/余弦函数的凸凹交替特性，即使沿射线存在多个密度峰，渲染的传感器信号仍可高度匹配真实测量值，但均值深度（各高斯深度的加权平均）与传感器导出深度之间产生系统性偏差（Jensen 不等式的有限形式不保证等号成立），只有当所有高斯沿射线聚集在同一深度时偏差才消失。这使得基于 NeRF 的 C-ToF 方法虽然因 MLP 的低熵偏差而鲁棒，但优化速度极慢；而更快的 GS 方法因高斯的任意放置自由度，在欠约束设置下容易产生高熵多峰分布而变得脆弱。

Time of the Flight of the Gaussians（2025.05）~\cite{2505_05356} 致力于同时解决速度和稳定性两个问题。速度方面，该方法用 3D 高斯 Splatting 替代 NeRF 的 MLP 场景表示，将优化和渲染速度提升 100 倍。稳定性方面，论文深入分析了 C-ToF 辐射场优化中的经验偏差（empirical biases），揭示了两层核心问题：深度仅通过传感器测量值的重建被间接优化，优化对初始化和超参数高度敏感；而 NeRF 的 MLP 天然偏向低熵（简单）解——沿射线的单一高密度尖峰比多个中等密度值具有更低熵——GS 却无此内在偏差，高斯可以沿射线任意放置来拟合目标。

\textbf{C-ToF GS 图像形成模型。}在标准 GS 渲染中，像素颜色是各高斯颜色的加权叠加。为适配 C-ToF 信号，该方法将场景辐射率替换为红外光传输模型：每个高斯的颜色 $c_k$ 被替换为光源强度 $s$ 与视角相关表面反射率 $r_k$（用 16 个球谐系数参数化）的乘积，再乘以逆平方衰减 $1/d_k^2$（$d_k$ 为高斯中心到相机光心距离），并加入相量调制 $W(d_k) = \exp(\mathrm{i}\psi_k)$，其中相位偏移 $\psi_k = 4\pi d_k f / \mathrm{c}$。由此得到相量形成模型：像素 $\mathbf{x}$ 处的相量为 $\mathbf{p}(\mathbf{x}) = \mathbf{p}^{\text{bg}} T_N + \sum_{k=1}^{N}\frac{2sr_k}{d_k^2}W(d_k)\,o_k\,\mathcal{G}^{\text{2D}}_k(\mathbf{x})\,T_k^2$，其中透射率被平方以表示光在辐射场中往返一次的衰减。进一步，为处理原始四分量帧的异步采集（四帧并非同时拍摄），用调制函数 $\phi(d_k) = \{\sin\psi_k, \cos\psi_k, -\sin\psi_k, -\cos\psi_k\}$ 替换相量，得到原始帧形成模型 $\mathbf{q}(\mathbf{x}) = \mathbf{q}^{\text{bg}} T_N + \sum_{k=1}^{N}\frac{sr_k}{d_k^2}\phi(d_k)\,o_k\,\mathcal{G}^{\text{2D}}_k(\mathbf{x})\,T_k^2$。关键矛盾在于：场景深度 $d(\mathbf{x}) = \sum_{k=1}^{N}d_k\,o_k\,\mathcal{G}^{\text{2D}}_k(\mathbf{x})\,T_k$ 是各高斯深度的加权均值，但即使完美重建原始帧 $\mathbf{q}(\mathbf{x})$，也不保证 $d(\mathbf{x})$ 与传感器导出深度一致。

\textbf{两种启发式策略。}与依赖额外学习型深度先验的方法不同，该方法提出的两种策略完全基于对优化过程的经验分析，不引入外部假设。然而它们并非直觉上的正则化——标准的深度畸变损失 $\text{DD}(\mathbf{x}) = \sum_{k=1}^{N}\sum_{l=1}^{N}\omega_k\omega_l\|d_k - d_l\|^2$ 虽可推动高斯聚向表面，却隐式引入不透明表面假设，与建模复杂光传输的目标相矛盾，且倾向于过度平滑深度。论文的两种策略避开了这一方向。

\textit{启发式 1：占用偏差（occupancy bias）。}在重建参数空间（反射率、位置、不透明度、致密化）中，"物体在哪里"（即占用）是主要方差来源，反射率变化应比占用变化更少发生。据此，该方法将反射率的学习率降低 10 倍，使优化更倾向于通过调整高斯位置和不透明度来拟合数据，而非调整反射率。这一简单调整鼓励高斯从一开始就更贴合真实场景结构，且不强制高斯彼此靠近（与深度畸变损失不同）。

\textit{启发式 2：低反射率偏差（low-reflectivity bias）。}初始反射率值对低反射率区域的优化结果至关重要。若初始化为较高值（如 0.5），结合占用偏差后，逆平方衰减项会促使低反射率区域的高斯被错误放置在远离相机处——在远距离高斯的反射率被纠正之前，空间重叠的近处高斯会干扰纠正过程，导致多峰分布。若初始反射率较低（0.01--0.1），低反射率区域的高斯可被更正确地放置（较低初始化更接近真实场景），而高反射率区域的高斯因逆平方衰减被吸引至更近位置后，会遮挡后方场景从而减少对其他高斯的干扰，获得更多迭代来优化反射率和位置至真实深度。实验证实，结合两种偏差对低反射率和高反射率区域均有效，而仅使用占用偏差会损害低反射率区域。

\textbf{流水线实现。}系统输入为原始四分量 C-ToF 视频序列和可选的颜色序列。动态建模采用变形 MLP（8 层、256 神经元，位置编码频率 $L_x = L_t = 10$），以时间 $t$ 和高斯 3D 位置 $\mathbf{x}_k$ 为输入，输出位置偏移 $\delta\mathbf{x}_k$，将高斯从规范空间隐式变形到各时刻。由于四帧异步采集，关键设计是局部线性假设：仅在整数时间步变形高斯，中间分数时间步 $j$ 通过最近两个整数步的线性插值获得 $\mathbf{x}_j = (i_2 - j)\mathbf{x}_{i_1} + (j - i_1)\mathbf{x}_{i_2}$。为应对快速运动，引入 RAFT 光流弱监督：将 3D 运动偏移投影至 2D 与估计光流计算 $L_2$ 损失。总损失为 $\mathcal{L} = \alpha\mathcal{L}_q + \mathcal{L}_c + \beta\mathcal{L}_f$，其中 $\mathcal{L}_q$ 为原始帧重建损失（$L_2$ + SSIM），$\mathcal{L}_c$ 为可选颜色损失，$\alpha=5$（合成场景为 1），$\beta=0.0008$。此外还有两个实用技巧：（1）随机背景——对低反射率敏感区域，低背景贡献已足够重现 C-ToF 信号而导致高斯贡献过低，因此在每次迭代随机改变背景信号（均匀分布在 $[-1, 1]$）；（2）静态预热——前 2K 次迭代假设场景静止，帮助将大量高斯放置到有效位置。高斯在相机视锥内随机初始化，反射率初始为 0.1，MLP 最后一层权重以标准差 $10^{-5}$ 初始化以确保变形从极小偏移开始稳定优化。

\textbf{实验结果。}在 fToF 合成数据集（七个 Blender 渲染场景，包含轴向/侧向/旋转运动及遮挡）上，该方法在七个场景中的五个取得最优深度 MSE。在真实数据集上（ToRF 数据集：五个室内慢运动场景，含移动相机和 RGB 图像；fToF 数据集：五个静止相机的快运动场景，包含挥动棒球棒、下落枕头等），该方法与 NeRF 基线 fToF 精度相当，在静态区域（地面、推车等）的重建甚至更优。计算时间从 NeRF 方法的 60--72 小时缩短至 40--60 分钟（单 NVIDIA 3090），渲染速度达 100+ Hz，提升约 100 倍。已知局限包括：静态相机下的深度包裹伪影（本质上是病态问题）、强非线性运动（如旋转扇叶）难以处理、薄而快速运动的物体因致密化在规范空间进行而难以维持结构。该方法为 CVPR 2025 口头报告。

\paragraph{细粒度非刚性表面跟踪}

从单目 RGB 视频重建高度可变形表面（如布料、纸张）是 Shape-from-Template（SfT）问题的核心挑战。先前基于网格的物理 SfT 方法存在固有局限：网格分辨率难以预设，有限元方法在不同分辨率下产生不一致的模拟结果阻碍由粗到精策略，逐帧优化易陷入局部极小值，且基于网格的可微渲染器不支持动态细节所需的复杂重网格化。

Thin-Shell-SfT（2025.03）~\cite{2503_19976} 的核心思想是用自适应的连续表面和变形表示替代离散网格，紧密耦合可微物理先验与 3DGS 可微渲染。

\textbf{连续变形模型。}表面表示为连续时空神经场。假设模板以网格 $M$ 形式给出，首先通过共形参数化将其映射到 2D 参数空间 $\Omega$。神经参考场（NRF）$x(\xi;\theta)$ 学习模板表面的连续参数化——在参数空间和顶点空间中随机采样重心坐标构造训练对，用 $L_1$ 几何损失将 MLP 拟合到模板网格参数化。神经变形场（NDF）$u(\xi, t;\phi)$ 预测逐时刻变形偏移，完整的位置映射为 $x(\xi, t) = x(\xi) + u(\xi, t)$，其中 $u(\xi, 1) = 0$（首帧无变形），后续时刻遵循因果递推 $u(\xi, t) = u(\xi, t-1) + \mu \cdot \Delta F(\xi, t)$，标量 $\mu > 0$ 通过使未来状态的变形沿先前变形方向延续来鼓励动量守恒。两个 MLP 均使用 SIREN 正弦激活函数（NRF 频率 $\omega_0 = 5$，NDF 频率 $\omega_0 = 30$），以表示高频褶皱信号并支持物理损失所需的高阶自动微分——两者均为 5 层、256 隐藏单元。NRF 训练约 2 分钟收敛；核心的 NDF 训练在 NVIDIA A100 上需 30 分钟至 1 小时。

\textbf{表面诱导高斯。}高斯基元通过 Poisson 盘采样从模板表面均匀采样的 $N_g \approx 90\text{k}$ 个点初始化，每个点携带参数坐标 $\xi_i$、对应位置 $\bar{x}_i$ 和颜色 $c_i$。旋转矩阵取自初始表面（NRF）的局部协变基向量 $R_i = [\mathbf{a}_1(\xi_i), \mathbf{a}_2(\xi_i), \mathbf{a}_3(\xi_i)]$，其中 $\mathbf{a}_\alpha = \partial x / \partial \xi^\alpha$ 为切向基，$\mathbf{a}_3 = \mathbf{a}_1 \times \mathbf{a}_2$ 为法向。沿法线方向的尺度 $s_3$ 固定为小值 $\epsilon = 10^{-5}$，仅优化切线方向尺度 $s_1, s_2$——利用单目设置下初始表面近似平面且完全可见的先验，避免学习各向异性协方差导致的歧义。关键设计是高斯的共享参数（协方差、不透明度、颜色）仅从第一帧更新，后续帧的反向传播仅更新变形场参数 $\phi$——这避免了变形序列中几何与外观的歧义导致错误的高斯重建。动态跟踪高斯 $\mathcal{G}_t$ 的参数化如下：位置来自变形模型 $x(\xi_i, t; \theta, \phi)$，旋转和尺度来自模板帧的 NRF，不透明度 $o_i$ 和颜色 $c_i$ 跨所有表面状态共享。

\textbf{数据损失。}采用 3DGS 风格的 $L_1$ 光度损失。若分割掩码可用，额外添加剪影损失加速训练。数据损失为：

\begin{equation}
\mathcal{L}_d(s_1, s_2, o, c, \phi) = \mathcal{R}(\mathcal{G}_1, I_1; s_1, s_2, o, c) + \sum_{t=2}^{T} \mathcal{R}(\mathcal{G}_t, I_t; \phi) + \mathcal{R}(\mathcal{G}_t, M_t; \phi)
\end{equation}

其中 $\mathcal{R}(x; \cdot)$ 为高斯光栅化损失。第一帧的 splatting 更新共享高斯属性，其余帧反向传播至 NDF 参数 $\phi$。

\textbf{Kirchhoff-Love 薄壳物理先验。}仅用光度损失优化 NDF 可完美拟合输入图像但因单目歧义导致物理上不合理的表面。物理先验基于 Kirchhoff-Love 薄壳理论，在连续表面上计算非线性膜应变 $\varepsilon = [\varepsilon_{\alpha\beta}]$ 和弯曲应变 $\kappa = [\kappa_{\alpha\beta}]$。具体而言，在每个训练迭代随机采样 $N_p = 100$ 个参数点，在初始状态（NRF）和变形状态（NDF）上执行微分几何运算：首先提取局部协变基 $\{\mathbf{a}_\alpha = \partial x / \partial \xi^\alpha, \mathbf{a}_3 = \mathbf{a}_1 \times \mathbf{a}_2\}$ 及表面度量张量 $a_{\alpha\beta}$、曲率张量 $b_{\alpha\beta}$，然后通过 NDF 的变形梯度计算非线性膜应变 $\varepsilon_{\alpha\beta} = \frac{1}{2}(u_{\alpha|\beta} + u_{\beta|\alpha} + u_{3,\alpha} u_{3,\beta})$ 和弯曲应变 $\kappa_{\alpha\beta}$，其中竖线表示协变导数、逗号表示偏导数，并采用 Einstein 求和约定。采用线性各向同性本构模型，面内刚度 $D = Eh/(1-\nu^2)$、弯曲刚度 $B = Eh^3/12(1-\nu^2)$、弹性张量 $H^{\alpha\beta} = \nu a^{\alpha\beta} a^{\gamma\delta} + \frac{1}{2}(1-\nu)(a^{\alpha\gamma} a^{\beta\delta} + a^{\alpha\delta} a^{\beta\gamma})$。与 NeuralClothSim 的正向模拟不同，逆向设置中外力和边界条件未知，因此将材料参数设为典型值（布料/纸张：$E = 5000\text{Pa}$，$\nu = 0.25$，$h = 1.2\text{mm}$），省略外力势能，图像损失充当外力角色引导变形。物理损失最小化内部超弹性能量：

\begin{equation}
\mathcal{L}_p(\phi) = \frac{1}{2N_p T} \sum_{i=1}^{N_p} \sum_{t=2}^{T} \frac{D \cdot \varepsilon^\alpha H_{\alpha} \varepsilon + B \cdot \kappa^\alpha H \kappa}{\sqrt{a(\xi_i)}}
\end{equation}

其中 $\sqrt{a} = |\mathbf{a}_1 \times \mathbf{a}_2|$。与网格上的离散正则化不同，该连续先验在每次迭代重新采样 $N_p$ 个点，实现对连续表面域的自适应探索——在布料褶皱等精细细节处提供物理监督，同时保持内存高效（固定参数量的全局 MLP，无需高分辨率网格）。

\textbf{全局时空联合优化。}优化采用全局策略——每轮迭代更新所有帧而非随机选取单帧，这对单目设置至关重要：缺乏多视角线索时，随机帧优化会导致变形场 MLP 陷入折叠结构的局部极小值。因果递推结构使梯度从未来状态反向传播至早期帧，实现全局时空优化的同时保持表面动态的因果性。高斯数量固定（不使用自适应密度控制），因为高斯参数仅从第一帧学习，变形细节的展开方式无法预知。总目标 $\mathcal{L} = \lambda_d \mathcal{L}_d + \lambda_p \mathcal{L}_p$，其中 $\lambda_d = 5$，$\lambda_p = 1$。

\textbf{实验结果。}在 Kaireidataset 基准上评估，包含 9 个高难度布料变形 RGB 视频（附参考深度）。与 SfT 方法（Kaireidataset、DDD、Ngo2015、IsMo-GAN、Stotko）和 NRSfM 方法（Neural NRSfM、Diff-NRSfM）比较，使用 Procrustes 对齐后计算 Chamfer 距离。Thin-Shell-SfT 在大多数序列上优于所有现有方法，在所有序列的平均指标上显著领先，尤其能重建先前方法无法处理的细褶皱和折叠。在严重自遮挡的多层折叠序列上，物理先验为遮挡区域提供合理推断，而 Kaireidataset 因内存需求过高而失败。与 K-Planes 和 Deformable Gaussians 等动态视角合成方法比较，后者虽在新视角合成方面有效，但无法从单目静态相机视频中恢复时空一致的表面几何。消融实验量化了各模块贡献（Chamfer 距离）：无物理先验 $34.25 \to$ 无表面诱导高斯 $14.0 \to$ 无法线尺度约束 $3.75 \to$ 完整模型 $3.46$——物理先验和表面诱导高斯初始化是准确跟踪的关键。与先前 SotA Kaireidataset 比较，法线误差从 $L_2 = 0.013$、余弦 $= 0.041$ 降低至 $0.009$ 和 $0.034$。测试时可按任意时空分辨率连续查询 NRF 和 NDF，获取带时间对应关系的网格（通过球轴网格化）。

\subsection{退化处理与多模态融合}
\label{ss:06_3b}

本节介绍两条技术路线：（1）运动模糊与退化处理；（2）事件相机与多模态融合。

\subsubsection{运动模糊与退化处理}

本路线介绍处理实际单目视频中运动模糊和散焦模糊等退化问题的4D重建方法。

深度先验缓解了几何约束不足的问题，但实际单目视频还面临另一个现实挑战：运动模糊和散焦模糊。实际单目视频常因相机抖动和物体快速运动产生运动模糊，或因镜头对焦不准产生散焦模糊。这些退化不仅降低视觉质量，更严重的是会误导变形场的优化方向——模糊图像作为监督信号时，变形网络可能学会拟合模糊模式而非真实几何。

\paragraph{模糊单目视频运动去模糊 4D 重建}

前面讨论的方法都假设输入视频是清晰的。然而在实际拍摄中，相机抖动和物体快速运动会导致运动模糊，使每一帧都成为曝光时间内多个瞬态的叠加。标准的 4DGS 变形场做法是：在清晰视频上，将所有帧的高斯变形回一个共享的\textbf{规范高斯}空间（canonical Gaussians），通过变形矩阵将静态规范高斯变换到各时刻状态，再渲染并与清晰帧计算光度损失——监督信号本身就是干净的，变形场容易学到正确的几何。但当输入帧本身是模糊的，变形网络会拟合模糊模式而非真实运动：如果模型只输出与输入同样模糊的渲染结果，重建损失就已经很低了，优化失去了向"清晰"方向推进的驱动力。

\textbf{核心改动。}Deblur4DGS（2024.12）~\cite{2412_06424} 将上述"直接从模糊帧监督"的范式，改为"先建模曝光时间内的连续动态，再合成模糊帧进行监督"。直觉上，一张模糊照片是曝光时间内多张清晰照片的平均——如果能重建出这些隐含的清晰瞬态，再取平均与模糊输入比较，就能在拟合模糊外观的同时迫使模型学到清晰的底层表示。关键洞察来自 3DGS 的显式运动建模：不需要逐帧估计全部动态高斯，只需估计一组逐高斯的曝光时间参数 $\mathbf{w}_t$，再通过插值获得曝光内的连续状态。这与早期的 NeRF 方法 DyBluRF 形成对比——DyBluRF 需要为曝光内每个时刻建立独立的连续运动模型，优化复杂且难以产生高质量结果；Deblur4DGS 利用 3DGS 的显式运动表示，将连续动态估计简化为曝光时间估计。同期工作 BARD-GS 采用了类似策略，但由于优化约束不足（尤其对大运动场景），效果不及 Deblur4DGS。

\textbf{实现原理——静态区域。}对于第 $t$ 帧模糊图像 $B_t$ 及其动态区域掩码 $M_t$（由 SAM2 提取），将曝光时间均匀划分为 $N$ 个时间戳。静态高斯 $S$ 在所有时间戳共享，只需估计连续相机位姿：优化曝光起止时刻的位姿 $P_{t,1}, P_{t,N}$，通过李代数线性插值获得中间位姿 $P_{t,i} = P_{t,1} \exp(\frac{i-1}{N-1} \log(P_{t,1}^{-1}P_{t,N}))$。为更稳定地优化位姿，Deblur4DGS 部署一个轻量 MLP 作为相机运动预测器，先用静态重建损失 $\mathcal{L}_{\text{rec}}^s$ 预训练该 MLP 和静态高斯 $S$，其中 $\mathcal{L}_{\text{rec}}^s = (1-\lambda)L_1(\hat{B}_t^s, B_t^s) + \lambda\mathcal{L}_{\text{ssim}}(\hat{B}_t^s, B_t^s)$，$\hat{B}_t^s = (1-M_t) \odot \hat{B}_t$ 和 $B_t^s = (1-M_t) \odot B_t$ 分别为合成模糊图像和输入模糊图像的静态区域。对于每个时间戳，用对应的 $P_{t,i}$ 和 $S$ 渲染清晰图像，将 $N$ 幅图像取平均得到合成模糊图像 $\hat{I}_t$，与输入计算重建损失 $\mathcal{L}_{\text{rec}} = (1-\lambda)L_1(\hat{I}_t, B_t) + \lambda_{\text{ssim}}\mathcal{L}_{\text{ssim}}(\hat{I}_t, B_t)$，其中 $\lambda=0.2$，与标准 3DGS 设置一致。

\textbf{实现原理——动态区域的模糊感知可变规范高斯。}动态区域的核心挑战在于规范高斯的选择。前面读者已了解，变形场方法通常在整个视频中维护一个共享规范高斯 $C$，所有时刻的高斯都变形自这个 $C$。但模糊视频常伴随大运动和低帧率，单一 $C$ 难以表达跨度很大的变形。Deblur4DGS 提出\textbf{模糊感知可变规范高斯}（blur-aware variable canonical Gaussians）：将视频均匀划分为 $K$ 个片段，每个片段维护自己的规范高斯 $C_k$，$C_k$ 仅负责邻近时间戳的变形，从而降低大运动建模的难度。选择 $C_k$ 的策略与模糊程度直接关联——按模糊度指标 $b_t = \sum_{(u,v) \in M_t}(\Delta B_t(u,v) - \bar{\Delta B}_t)^2$ 度量每帧的清晰程度，其中 $\Delta B_t$ 为模糊图像的拉普拉斯算子，$\bar{\Delta B}_t$ 为其均值（$b_t$ 越大表示拉普拉斯响应越强、帧越清晰），在每个片段中选择最清晰帧对应的高斯作为 $C_k$。为使片段的起止帧尽量清晰，还在起止帧周围 $H$ 帧内搜索最清晰帧重新定义片段边界。变形方式沿用 Shape-of-Motion 的刚体变换：$x_t = A_t x_c + E_t$，$r_t = A_t r_c$，其中 $A_t, E_t$ 为旋转平移矩阵，$x_c, r_c$ 为规范空间中的位置和旋转。

\textbf{实现原理——曝光内连续动态高斯。}有了整数时间戳上的动态高斯 $D_t$，下一步是在曝光时间内获得连续的 $D_{t,i}$。直接为每个时间戳学习独立的变形参数是不稳定的。Deblur4DGS 的关键简化是：假设曝光时间内物体运动近似均匀，因此 $D_{t,i}$ 可以通过最近整数时间戳的动态高斯插值获得。对于曝光前半段（$i \in [1, N/2]$），在 $D_{t-1}$ 和 $D_t$ 之间插值：$D_{t,i} = \mathbf{w}_{t,i} D_{t-1} + (1-\mathbf{w}_{t,i}) D_t$；对于后半段（$i \in [N/2, N]$），在 $D_t$ 和 $D_{t+1}$ 之间插值。$\mathbf{w}_{t,i}$ 是 $D_{t,i}$ 与 $D_{t,N/2}$ 之间的归一化时间间隔。实际实现中，只需估计曝光起止时刻的权重 $\mathbf{w}_{t,1}$ 和 $\mathbf{w}_{t,N}$，中间时刻通过线性插值获得：$\mathbf{w}_{t,i} = (1-\frac{i-1}{N-1})\mathbf{w}_{t,1} + \frac{i-1}{N-1}\mathbf{w}_{t,N}$。由于假设曝光内运动均匀，$\mathbf{w}_{t,1}$ 和 $\mathbf{w}_{t,N}$ 的绝对值相等且等于半曝光时间 $\mathbf{w}_t$，因此上式可简化为 $\mathbf{w}_{t,i} = (1-\frac{i-1}{N-1})\frac{\mathbf{w}_t}{2} + \frac{i-1}{N-1}(-\frac{\mathbf{w}_t}{2})$。这里 $\mathbf{w}_t$ 是逐高斯的可学习参数（而非单一标量），每个动态高斯独立估计自己的半曝光时间。规范高斯、变形模块和 $\mathbf{w}_t$ 被联合优化。

\textbf{三项正则化防止平凡解。}仅靠重建损失存在多重解：最理想的情况是每个 $D_{t,i}$ 都清晰，但最平凡的解是所有 $D_{t,i}$ 退化为与输入相同的模糊帧——此时平均后仍然匹配模糊输入，重建损失同样很低。Deblur4DGS 通过实验验证了这一机制：将视频帧数减少为单帧以消除帧间一致性后，静态区域也出现模糊，说明帧间一致性本身对静态区域起到了隐式正则化作用。动态区域因帧间一致性弱，需要显式约束。Deblur4DGS 引入三项正则化：（1）\textbf{曝光正则化} $\mathcal{L}_e = \max(0, \epsilon - \mathbf{w}_t)$，约束 $\mathbf{w}_t$ 不趋近于零，否则所有 $D_{t,i}$ 退化为同一状态；（2）\textbf{多帧一致性} $\mathcal{L}_{\text{mfc}} = \frac{1}{N-1}\sum_{i=2}^{N}(\|M_{t,i}(\hat{I}_{t,i-1\to i} - \hat{I}_{t,i})\|_1 + \|M_{t,1}(\hat{I}_{t,i\to 1} - \hat{I}_{t,1})\|_1)$，通过预训练光流网络 PWC-Net 将相邻帧对齐到当前帧、将各帧对齐到首帧后约束内容相似性；（3）\textbf{多分辨率一致性} $\mathcal{L}_{\text{mrc}} = \|(M_{t,i})_\downarrow \odot ((\hat{I}_{t,i})_\downarrow - \text{sg}(\hat{I}_{t,i}^l))\|_1$，利用低分辨率模型（模糊程度更低、更易去除，以降采样视频为监督预训练）的输出监督高分辨率模型，其中 $\text{sg}$ 为停止梯度，$\downarrow$ 为下采样操作，$\odot$ 为逐像素乘法。总正则化损失为 $\mathcal{L}_{\text{reg}} = \lambda_e \mathcal{L}_e + \lambda_{\text{mfc}} \mathcal{L}_{\text{mfc}} + \lambda_{\text{mrc}} \mathcal{L}_{\text{mrc}}$，权重 $\lambda_e=0.1, \lambda_{\text{mfc}}=2, \lambda_{\text{mrc}}=1$。此外，Deblur4DGS 还沿用了 Shape-of-Motion 的额外正则化项 $\mathcal{L}_{oth}$ 以辅助三维运动重建。除新视角合成外，Deblur4DGS 还可通过调整相机位姿和时间戳实现去模糊、帧插值和视频稳定化——具体而言：使用模糊视频的相机位姿渲染得到去模糊结果；使用插值的相机位姿和时间戳渲染得到帧插值结果；使用平滑后的相机位姿渲染得到视频稳定化结果。

\textbf{实验结果。}在合成数据（$288 \times 512$ 分辨率）上，Deblur4DGS 的新视角合成 PSNR 达到 27.66 dB，相比最强基线 BARD-GS（26.91 dB）提升 0.75 dB，SSIM 从 0.923 提升至 0.935，LPIPS 从 0.077 降至 0.060。在高分辨率（$720 \times 1080$）上，PSNR 为 27.16 dB（BARD-GS 为 26.34 dB），同样显著领先。在真实世界数据上，Deblur4DGS 的 CLIPIQA 达到 0.356、MUSIQ 达到 36.756，均优于 BARD-GS（0.288/35.505）和其他方法。在去模糊、帧插值、视频稳定化等下游任务上，Deblur4DGS 与任务专用的监督视频处理模型相比也展现了竞争力的性能，同时保持了实时渲染速度。

\paragraph{动态 GS 去模糊：散焦与运动模糊统一处理}

现有去模糊方法要么只处理运动模糊（通过曝光时间内多帧加权求和），要么只处理散焦模糊（通过分层景深体渲染），尚无方法能同时应对两种退化。Dynamic Gaussian Splatting from Defocused and Motion-blurred Monocular Videos（2025.10）~\cite{2510_10691} 的核心观察是：尽管运动模糊与散焦模糊的物理成因截然不同，两者都可以近似为像素与其邻域像素的加权和。据此，该方法将模糊图像的生成统一建模为
$\hat{B}(x) = \sum_{x_i} I(x_i) \, k_x(x_i)$，其中 $k_x$ 为像素 $x$ 的逐像素模糊核，满足权重归一化约束 $\sum_{x_i} k_x(x_i) = 1$。
关键难点在于为每个像素估计合理的模糊核 $k_x$。该方法设计了模糊预测网络（BP-Net），这是一个四层 CNN，以三类信息为输入：可学习的相机视角嵌入 $e(i)$、由渲染图像/深度图/运动掩码经三层 CNN 编码的场景特征 $f_{\text{scene}}$，以及像素坐标位置编码 $p(x)$。BP-Net 同时输出模糊核 $k_x$ 和模糊强度 $m_x \in [0,1]$，输出模糊像素通过锐利与模糊分量的混合计算：$\hat{B}(x) = (1 - m_x) I(x) + m_x \tilde{B}(x)$，其中 $\tilde{B}(x)$ 为卷积模糊结果。

该混合机制确保模糊主要由网络显式建模，而渲染图像 $I$ 在模糊强度 $m_x$ 较小的区域被直接约束向输入模糊图的清晰部分靠拢，逐步引导 3DGS 表示趋向锐利。
为防止不合理模糊核过度模糊轻微退化区域，该方法引入模糊感知稀疏约束。具体而言，模糊核中心权重 $k_x(c)$ 反映核的集中程度（值越小表示越分散），而设计中心权重目标值 $c_x = \sigma(\text{scale} \cdot (1 - \text{sg}(m_x)))$（scale 设为 5，sg 为停止梯度操作），使得模糊强度 $m_x$ 越大，目标中心权重 $c_x$ 越小，从而自适应约束核的分散程度与实际模糊程度一致。训练损失为 $L_{\text{spa}} = L_1(c_x, k_x(c))$。
在场景表示方面，该方法采用 Shape-of-Motion 初始化动静分离的高斯，动态高斯的运动由 $N_b$ 个 $SE(3)$ 运动基底的加权组合建模。为弥补仅使用规范帧可见跟踪点导致动态区域高斯不完整的问题，引入动态高斯致密化策略：在前 $N_d$ 轮训练后执行一次，从各观测帧的动态区域随机采样像素，利用深度反投影和前景重映射将新高斯补充到规范空间。

此外，该方法通过在训练视角周围生成平行和垂直方向的未见视角（垂直方向偏移距离为 $[0.5, 1]$ 归一化单位），利用训练视角的颜色和深度信息投影获得伪监督，缓解单目视角不足导致的过拟合问题。

\textbf{BARD-GS：显式解耦相机与物体运动模糊。}前面提到

\paragraph{BARD-GS：2025.03 同期工作}

BARD-GS（2025.03）~\cite{2503_15835} 采用了与 Deblur4DGS 类似的"先建模曝光内动态再合成模糊帧"策略，但其独特贡献在于将运动模糊\textbf{显式解耦}为相机运动模糊和物体运动模糊两种独立成分。该方法的核心观察是：静态区域仅受相机抖动影响，而动态区域则同时承受相机和物体运动的双重退化，二者耦合在一起会导致优化方向混乱。BARD-GS 由此设计了两阶段训练流程。第一阶段专注于\textbf{相机运动去模糊}：将曝光时间内的相机轨迹建模为一组虚拟相机位姿序列，从这些虚拟位姿渲染静态区域图像并取平均，重建出与模糊输入对齐的合成模糊图像，通过匹配约束优化静态高斯和虚拟位姿。关键点在于此阶段\textbf{仅优化静态区域}，以干净地分离相机运动的影响。第二阶段处理\textbf{物体运动去模糊}：学习一个时间条件变形场（time-conditioned deformation field）跟踪动态高斯的时空位移，类似于相机运动建模，将曝光时间内动态高斯在不同时间戳的位置渲染出来并取平均，重建物体引起的模糊。最终将完整场景（静态+动态）的合成模糊图与输入对齐，同时实现两类运动模糊的消除。这种显式解耦策略在动态区域获得了显著优于先前耦合式方法的渲染质量。此外，BARD-GS 还收集了一个真实世界动态场景运动模糊数据集，为该方向的评测提供了新的基准。

\textbf{MoBGS：模糊自适应潜在相机估计。}

\paragraph{MoBGS：2025.04 同期工作}

MoBGS（2025.04）~\cite{2504_15122} 从另一个角度切入运动模糊问题：不直接估计虚拟相机位姿序列，而是通过神经 ODE 建模连续的潜在相机轨迹。其核心模块是\textbf{模糊自适应潜在相机估计}（Blur-adaptive Latent Camera Estimation, BLCE），利用一个模糊自适应神经 ODE 求解器从模糊图像中推断曝光时间内的连续相机运动轨迹。与传统离散虚拟位姿方法相比，神经 ODE 提供了更平滑、更灵活的相机运动建模，能更好地适应不同速度和模式的相机抖动。此外，MoBGS 提出\textbf{潜在相机引导曝光估计}（Latent Camera-induced Exposure Estimation, LCEE）方法，确保全局相机运动和局部物体运动的去模糊一致性：利用估计的潜在相机轨迹推导曝光参数，将相机运动建模的精度传递到动态物体的去模糊流程中。MoBGS 在 Stereo Blur 数据集和真实模糊视频上均显著超越先前方法（AAAI 2026 接收）。

\textbf{DBMovi-GS：稀疏控制点的模糊感知动态重建。}

\paragraph{DBMovi-GS：2025.06 同期工作}

DBMovi-GS（2025.06）~\cite{2506_20998} 针对模糊单目视频中的动态新视角合成，提出基于稀疏控制点的高斯运动控制策略。该方法的核心挑战在于：模糊视频中 COLMAP 等经典 SfM 工具因特征匹配困难而无法提供可靠的初始点云，导致高斯初始化质量极差。DBMovi-GS 通过深入分析高斯初始化过程，设计了一种从稀疏点云生成密集 3D 高斯的方法：在稀疏 COLMAP 点云基础上，通过深度约束和几何扩散策略补充高斯密度，确保场景的完整覆盖。在运动建模方面，DBMovi-GS 使用一组稀疏控制点驱动动态高斯的变形，控制点通过可学习的 SE(3) 变换跟踪物体运动，周围高斯通过插值获得各自的运动参数。该方法能够同时处理三种常见模糊类型：相机运动模糊、物体运动模糊和散焦模糊，在动态模糊场景的新视角合成中建立了新的性能基准。

\textbf{UPGS：统一位姿感知框架。}

\paragraph{UPGS：2025.09 同期工作}

UPGS（2025.09）~\cite{2509_00831} 指出现有方法的一个关键瓶颈：它们普遍采用"先估计相机位姿、再优化高斯"的两步流水线，而模糊伪影会严重破坏位姿估计精度，位姿误差又会累积到后续重建中。UPGS 提出将相机位姿作为可学习参数纳入 3DGS 属性的统一优化框架，实现端到端的去模糊重建。具体而言，该方法将相机运动和物体运动统一建模为每个高斯基元的逐基元 SE(3) 仿射变换，构建统一的优化目标。为保持训练稳定性，UPGS 设计了三阶段训练策略：第一阶段固定位姿、仅优化高斯属性；第二阶段固定高斯、仅优化相机位姿；第三阶段联合优化所有可学习参数。这种交替优化策略有效避免了位姿和高斯的耦合退化。在 Stereo Blur 数据集和真实世界序列上，UPGS 在重建质量和位姿估计精度上均显著超越先前方法。

\subsubsection{事件相机与多模态融合}

本路线介绍将事件流与RGB图像融合以辅助动态重建的多模态方法。

运动模糊的根源之一是标准相机的帧率限制。事件相机（event camera，又称动态视觉传感器 DVS）具有微秒级时间分辨率，天然适合捕捉快速运动，但缺少颜色信息。将事件流与 RGB 图像融合为动态重建提供了新的可能性——事件相机提供精确的运动时序，RGB 相机提供丰富的颜色纹理。此外，视频扩散模型作为另一种模态先验，也被引入以提升渲染质量和视角覆盖。\paragraph{事件引导人体-场景联合重建}

现有人体重建方法（如 HUGS）使用分离的高斯集合分别建模人体和场景，需要外部人体掩码作为预处理，且难以应对快速运动产生的模糊。ExFMan 虽然利用事件数据进行动态人体重建，但仅关注人体本身，不具备场景建模能力。Event-guided 3D Gaussian Splatting（2025.09）~\cite{2509_18566} 的设计思路是：与其分离表示，不如用一组统一的 3D 高斯同时编码人体与场景，通过可学习语义属性实现自动分化。
具体而言，每个高斯 $G_o$ 携带一个可学习语义参数 $s_i \in \mathbb{R}$，其初始化基于初始点云可用的语义标签 $s_i = \mathcal{L}_i \in \{0, 1\}$。训练过程中，$s_i$ 经 sigmoid 得到软掩码值 $m_i = \sigma(s_i)$，再以 0.5 为阈值二值化为硬掩码 $\hat{s}_i \in \{0, 1\}$。只有被分类为人体的高斯（$\hat{s}_i = 1$）才进入变形网络进行动画化，场景高斯保持静态。

变形过程沿用 3DGS-Avatar 的范式：先由非刚性变形网络 $\phi_{\text{nr}}$ 在规范空间建模衣物褶皱等细节变化 $\{\mathcal{G}_d\} = \Phi_{\psi_{nr}}(\{\mathcal{G}_c\}, z_p)$，再通过线性混合蒙皮（LBS）变换到观测空间：
$\{\mathcal{G}_{o,i}\} = \hat{s}_i \cdot \Phi_{\psi_r}(\{\mathcal{G}_d\}; \{\mathbf{B}_b\}_{b=1}^B)$，其中 $z_p$ 为编码的姿态向量，$\mathbf{B}_b$ 为骨骼变换，$\Phi_{\psi_r}$ 为蒙皮 MLP 预测位置 $\mathbf{x}_d$ 处的蒙皮权重。语义属性还被集成到高斯致密化过程中，新高斯继承父节点的语义属性，确保分裂和克隆操作不会破坏已建立的人体-场景分化。在颜色建模上，人体和场景高斯分别使用独立的颜色 MLP，人体高斯颜色 $c_h$ 和场景颜色 $c_{\text{bg}} = \mathcal{F}_{\psi_{bc}}(f_{\text{bg}}, \gamma(d))$ 由各自的 MLP 独立预测，其中场景颜色 MLP 将可学习特征向量 $f_{\text{bg}}$ 与球谐方向基 $\gamma(d)$ 联合作为输入，将神经网络的非线性表达能力与球谐函数的方向编码相结合，在噪声或运动模糊等挑战性条件下比纯 SH 表示更具鲁棒性。
为对抗运动模糊，该方法利用事件相机的高时间分辨率引入事件引导损失。事件相机与传统帧式相机不同：每个像素独立异步工作，当像素对数亮度 $L(u_k, t_k) = \log I(u_k, t_k)$ 自上次事件以来的累积变化超过对比度阈值 $C$ 时触发一个事件 $e_k = (u_k, t_k, p_k)$，即 $L(u_k, t_k) - L(u_k, t_{k-1}) = p_k \cdot C$。

据此，该方法计算连续渲染帧 $I_{t_k}$ 与 $I_{t_{k+1}}$ 之间的对数亮度变化作为模拟事件：
$\Delta \mathcal{L} = \log(I_{t_{k+1}}^{2.2} + \epsilon) - \log(I_{t_k}^{2.2} + \epsilon)$，其中图像先从 sRGB 空间转至线性空间（按 IEC 61966-2-1 标准做 gamma 2.2 校正），$\epsilon$ 防止数值不稳定。事件损失定义为归一化 $L_1$ 距离：
$L_{\text{event}} = w_{\text{ev}} \left| \frac{\Delta \mathcal{L}}{\|\Delta \mathcal{L}\|_F} - \frac{E_{\text{gt}}}{\|E_{\text{gt}}\|_F} \right|_1$，其中 Frobenius 范数归一化确保损失对不同光照条件和事件密度具有鲁棒性，$w_{\text{ev}}$ 为权重因子。
该论文还贡献了两个运动模糊基准数据集：ZJU-MoCap-Blur（基于 ZJU-MoCap 的 6 个序列，使用 Super-SloMo 模拟模糊）和 MMHPSD-Blur（基于 MMHPSD 的 6 个序列）。在 ZJU-MoCap-Blur 数据集上，该方法达到 31.91 dB PSNR，相比 HUGS 的 26.70 dB 提升 +19.5\%，SSIM 从 0.9366 提升至 0.9736（+3.95\%），LPIPS 从 0.1075 降至 0.0726（$-32.5\%$）。

在 MMHPSD-Blur 数据集上，相比 HUGS 的 25.23 dB 提升 0.68 dB，SSIM 从 0.8447 提升至 0.9118。值得注意的是，论文的消融实验揭示了一个方法论洞见：直接使用 E2VID 将事件流转换为视频作为额外输入（30.25 dB）不如事件损失监督（31.81 dB）有效，级联 RGB+Event 去模糊方法（31.44 dB）也不如端到端事件损失集成。这表明将事件信息以损失函数形式直接约束渲染过程，比将其作为预处理步骤更能发挥事件相机的时序精度优势。

\paragraph{视频扩散感知 4D 重建}

ViDAR（2025.06）~\cite{2506_18792} 的关键洞察是：正则化方法（如 MoSca）通过几何约束获得紧凑的几何表示，但渲染质量往往偏低；而生成式方法利用扩散模型实现高视觉质量，却难以保持时空一致性。ViDAR 将扩散模型作为外观先验注入重建流程，同时通过扩散感知策略控制其几何不一致性。据作者所知，ViDAR 是首个以几何一致方式将扩散先验融入单目视频重建的方法。
具体流程分为三步。首先，使用 MoSca 作为基线获得初始 4D 表示，将场景分解为静态高斯 $\mathcal{G}_s$ 和动态高斯 $\mathcal{G}_d$（$\mathcal{G} = \mathcal{G}_d \cup \mathcal{G}_s$）。值得注意的是，MoSca 原始的对极误差分析在动/静态区域分类中会将背景误划为动态区域，导致浮点伪影（floater）。ViDAR 在此阶段引入 Track Anything 获取动态掩码 $D_t$ 替代 MoSca 的对极误差分类，更精确地约束动态高斯的位置，从而改善后续伪多视角样本的生成质量。

然后用 DreamBooth 策略对预训练的 Stable Diffusion XL（SDXL）进行场景级个性化微调，将输入视频视为图像集合进行训练，使特定文本标记触发该场景的外观。个性化模型对渲染的退化多视角图像进行 $k$ 步去噪增强：将渲染图 $R_{m,t}$ 编码为潜空间 $x_0 = \mathcal{E}(R_{m,t})$，使用 Discrete Euler 噪声调度器添加 $k$ 步噪声后执行 $k$ 步去噪，得到增强图像 $E_{m,t} = \mathcal{D}(\hat{x}_0)$。ViDAR 使用多步去噪而非一步 SDS，因为多步过程能保留更好的结构保真度。
在视角采样策略上，ViDAR 为每个时间步 $t$ 采样 $M$ 组新视角。具体而言，从已有视角中选择两个视角——一个随机视角和一个最具挑战性的视角（与均值的距离最远）——并以加权线性组合的方式生成新视角。为引入多样性的难度递增，采样过程中逐渐增加最具挑战性视角的混合权重，同时引入递增幅度的噪声。ViDAR 选择一次性生成整个多视角数据集而非在线增强，这样可以复用样本并降低 GPU 内存开销。

然而，扩散增强的图像在细粒度纹理上存在时空闪烁，在某些情况下扩散模型甚至会在训练时未见的新视角中生成错误的粗几何结构。直接将这些图像用于监督会导致模型收敛到平均辐射值（产生模糊渲染）或过拟合到单帧（产生时序不一致的重建）。为此 ViDAR 提出两个核心策略。（1）动态区域聚焦：利用 Track Anything 获取动态掩码 $D_{m,t}$，仅对动态区域施加扩散监督 $E_{m,t}^{\text{dyn}} = E_{m,t} \odot D_{m,t}$，静态区域已通过多帧自然形成多视角约束，无需额外伪监督反而能避免质量下降和浮点伪影的引入。动态区域的损失函数结合 $L_1$、基于预训练 VGG 网络的感知损失和 SSIM 损失：
$L_{\text{dyn}} = \|E_{m,t}^{\text{dyn}} - \hat{I}_{m,t}^{\text{dyn}}\|_1 + \lambda_p \|E_{m,t}^{\text{dyn}} - \hat{I}_{m,t}^{\text{dyn}}\|_{\text{vgg}} + \lambda_s \|E_{m,t}^{\text{dyn}} - \hat{I}_{m,t}^{\text{dyn}}\|_{\text{ssim}}$，其中 $\lambda_p$ 和 $\lambda_s$ 均设为 0.1。感知损失的引入有效抑制了空间错位伪 GT 导致的训练模糊。该损失反向传播更新场景表示 $\mathcal{G}$ 和输入相机位姿 $\mathcal{C}_{\text{inp}}$。（2）采样相机位姿联合优化：采样相机位姿 $\mathcal{C}_{\text{sample}}$ 由输入位姿 $\mathcal{C}_{\text{inp}}$ 插值而来，带有噪声且受扩散图像不一致性影响，需要在线精炼。

与动态区域仅使用掩码区域不同，相机位姿优化使用整幅图像 $E_{m,t}$ 作为监督信号，因为尽管存在细粒度纹理闪烁，静态区域的粗结构提供了比动态区域更一致的空间定位信号。相机位姿损失 $\mathcal{L}_{\text{cam}}$ 采用与 $\mathcal{L}_{\text{dyn}}$ 相同的 $L_1$ + 感知 + SSIM 组合形式，但该损失仅反向传播更新 $\mathcal{C}_{\text{sample}}$，不影响场景表示 $\mathcal{G}$。
在实验评估方面，ViDAR 在 DyCheck（iPhone）数据集上进行测试——这是目前唯一广泛使用的、具有极端视角变化的单目 4D 重建基准。在共可见性掩码区域（co-visibility masked）的半分辨率评估中，ViDAR 达到 19.69 dB PSNR、0.7126 SSIM 和 0.2231 LPIPS，全面超越所有基线方法。

在动态区域专用评估（Dynamic Region Benchmark，ViDAR 首次提出）中，ViDAR 以 16.46 dB PSNR-D 和 0.2793 LPIPS-D 显著领先于基线 MoSca（15.63 dB PSNR-D、0.2904 LPIPS-D），充分证明了扩散先验对运动丰富区域重建的增益。此外，ViDAR 还指出 DyCheck 标准评估中仅约 26\% 的共可见性像素属于动态区域（部分场景如 apple 和 paper-windmill 低至 4\%），因此全场景指标可能掩盖动态区域的实际性能差异。

\paragraph{视频扩散增强城市场景动态建模}

VDEGaussian（2025.08）~\cite{2508_02129} 针对动态城市场景的视频扩散增强框架，其出发点是一个关键观察：相邻输入帧之间的光流越大（通常由近距离或快速运动的物体引起），对应的新视角中间帧误差越高。这种时序不连续性的根源在于动态内容的欠采样，而视频扩散模型天然具备时序一致性先验。然而，直接将扩散生成的插值帧作为伪视角引入 3DGS 重建会面临两个难题：像素级 3D 不一致性，以及缺乏精细位姿控制。
VDEGaussian 采用两阶段架构。第一阶段为场景适配：虽然预训练视频扩散模型在大规模数据上学习了时序一致性先验，但其生成分布与自动驾驶数据存在偏差，且通常只能生成固定数量的帧（如 16 帧），过量帧生成会加剧 3D 不一致性并降低效率。为此，VDEGaussian 基于 DynamiCrafter 构建，在其冻结的空间注意力和时序注意力层上引入可学习 LoRA 适配器，同时联合微调视觉上下文注入模块以改善自动驾驶场景与潜在空间概念的对齐。

适配过程使用从训练视角提取的约 50 个 3 帧 clip，仅需 1000 次迭代微调即可让模型灵活生成目标帧数的插值视频，甚至泛化到训练时未见过的帧数（如 4 帧）。
第二阶段为 4DGS 重建增强。场景表示基于 Periodic Vibration Gaussian（PVG），其中动态高斯通过速度向量 $v$ 和生命期参数 $\beta$ 建模振荡运动：
$\tilde{\mu}(t) = \mu + \frac{l}{2\pi} \cdot \sin\left(\frac{2\pi(t-\tau)}{l}\right) \cdot v$，$\tilde{o}(t) = o \cdot e^{-\frac{1}{2}(t-\tau)^2 \beta^{-2}}$，其中 $\tau$ 为峰值时间，$l$ 为周期长度（场景先验超参数），$\tilde{o}(t)$ 为从峰值衰减的振动不透明度，$\beta$ 控制高斯活跃的时间范围。该表示通过生命期机制实现静态与动态高斯的快速分解，无需标注 3D 边界框。VDEGaussian 的实现基于 PVG 代码库，并集成 iComMa 的位姿反向传播以支持相机位姿在可微高斯光栅化器中的梯度回传。
为从视频扩散模型中有效提取先验，VDEGaussian 引入两个核心策略。（1）联合时间戳优化：由于视频扩散模型无法直接接收相机位姿控制信号，该方法转而主动对齐渲染图像的位姿以匹配生成结果。

插值位姿由相邻帧 $P_t = (R_t, T_t)$ 和 $P_{t+1} = (R_{t+1}, T_{t+1})$ 通过四元数球面线性插值计算。由于相邻帧时间间隔极短，四元数变化缓慢（点积接近 1），球面插值可简化为线性近似 $q_{\text{mid}} \approx (1-t_{\text{mid}})q_t + t_{\text{mid}} q_{t+1}$，既可微又高效。其中引入可学习偏移项 $\delta t$：$t_{\text{mid}} = t + \sigma(\delta t)$，$\sigma$ 为 sigmoid 函数限制偏移范围，$\delta t$ 初始化为零。平移参数同理插值：$T_{\text{mid}} = (1-t_{\text{mid}})T_t + t_{\text{mid}} T_{t+1}$。该偏移通过可微光栅化器反向传播图像损失进行端到端优化，确保渲染帧与生成帧的位姿对齐。消融实验表明，$\sigma(\delta t)$ 在训练中稳定收敛，且对不同初始化值均能收敛到相似结果，验证了方法的鲁棒性；同时发现生成帧的位姿并非严格位于相邻帧中点，暗示视频扩散模型在大规模预训练中习得了固有位姿偏置。（2）不确定性蒸馏：为自适应地从生成图像中提取可靠内容同时保护已良好重建的区域，VDEGaussian 引入可学习不确定性图 $\beta_e$（每帧伪视角一个），其训练目标为
$L_{\text{ca}}(p) = \omega_f\left(\beta_e(p)\left(\hat{I}_{t_{\text{mid}}}(p) - I_{\text{pseudo}}(p)\right)^2 - \lambda_f \beta_e^2(p)\right)$，其中 $\hat{I}_{t_{\text{mid}}}$ 为下采样至生成图像分辨率的渲染视角。

对 $\beta_e$ 求导令其等于零可得最优解 $\beta_e(p) = \frac{(\hat{I}_{t_{\text{mid}}}(p) - I_{\text{pseudo}}(p))^2}{2\lambda_f}$，即不确定性与渲染误差正相关：误差大的区域（如快速运动物体）获得更大权重，误差小的区域（如静态背景或慢速物体）受到更少干扰。可视化验证了该理论——不确定性图成功为快速运动物体分配高权重，为静态背景分配低权重。为缓解过拟合和高频噪声，对不确定性图施加全变差（TV）正则化 $L_{\text{tv}} = \omega_{tv} \text{TV}(\beta_e)$。在实现中，不确定性蒸馏每 4 次迭代执行一次以确保每个输入视角的平衡训练，超参数设为 $\omega_f = 1$、$\omega_{tv} = 0.001$、$\lambda_f = 1$，$\delta t$ 和 $\beta_e$ 使用 Adam 优化器（学习率 0.01）。
消融实验揭示了各组件的贡献：（a）缺少场景适配时，因视频生成模型无法产生正确帧数的视频，PSNR 从 28.31 暴跌至 9.65；仅有适配而无后两个策略时，性能与 PVG 基线持平（28.73 vs 28.31）；（b）联合时间戳优化带来小幅但稳定的提升（28.77 vs 28.73）；（c）不确定性蒸馏贡献最大，单独引入即使 PSNR 提升近 2 dB（30.49 vs 28.77）。

定性分析进一步表明，不确定性蒸馏有效解决了远处运动卡车的严重扭曲和模糊问题。整体框架在 Waymo 的 6 个挑战性动态场景上实现近 2 dB PSNR 提升（30.49 vs PVG 的 28.31），在 NOTR dynamic32 数据集上同样一致优于基线（29.33 vs 28.73）。此外，深度渲染对比表明 VDEGaussian 在新时间戳上展现出更优的几何一致性，有效缓解了基线方法在运动物体周围产生的深度伪影。

\paragraph{扩展受限视角的动态 NeRF}

ExpanDyNeRF（2025.12）~\cite{2512_14406} 针对动态 NeRF 系统在大视角偏差下渲染不稳定的问题，提出利用高斯 Splatting 先验和伪真值生成策略的单目 NeRF 框架。该方法的设计动机源于一个实验发现：在主视角下视觉质量相当的 NeRF 与 3DGS，在侧视角下表现截然不同——3DGS 倾向于将远处天空错误重建为近处结构，遮挡远处山脉和灌木等背景元素；而 NeRF（如 Instant-NGP）能正确保持天空的远距均匀特性，实现更高的保真度和更丰富的场景细节。需要指出，这并非断言 NeRF 普遍优于 3DGS，而是在大视角旋转这一特定设定下 NeRF 更为稳定，因此 ExpanDyNeRF 选择 NeRF 为主干而 3DGS 仅作为先验生成器。遵循 DeRF 的架构，模型采用双分支设计：静态背景分支 $\phi_b$ 使用基于分布的表示编码静态场景组件，改善相机投影与背景几何的对齐，简化低方差静态特征的渲染；动态前景分支 $\phi_f$ 使用三帧滑动窗口捕获时序变化，整合空间坐标、观察方向和时间戳 $t$。

\textbf{伪真值生成策略。}伪真值生成的核心流程如下。对每帧 $I_t$，使用 FreeSplatter 生成前景物体的 3D 高斯先验，以其为中心定义半径为 $R_d$（等于主视角到物体的距离）的穹顶采样空间。主视角位姿 $P_t$ 在穹顶上的位置记为 $(e, \text{azimuth}=0, R_d)$，其中 $e$ 为主视角的仰角。在该空间内系统性地采样新视角：方位角从 $-45^\circ$ 到 $45^\circ$（步长 $5^\circ$），仰角为 $0^\circ$、$15^\circ$、$30^\circ$ 三个层级，所有相机位姿 $P_{\text{nv}}^{(d)}$ 的前向向量始终指向穹顶中心，确保各视角下物体的一致构图。渲染这些视角获得 RGB 图像和形状掩码作为伪真值，编码了前景物体在不同观察角度下的期望密度和颜色分布，提供了从真实单目采集中无法获得的全面监督。为将伪真值从高斯坐标系映射到 NeRF 坐标系，通过刚性对齐矩阵进行变换：设前景 NeRF 坐标系下的主视角位姿为 $P_t^{(n)}$，高斯先验坐标系下的对应位姿为 $P_t^{(d)}$，则变换矩阵为 $T = P_t^{(n)} (P_t^{(d)})^{-1}$，将其应用于每个新视角 $P_{\text{nv}}^{(d)}$ 即可将所有相机位置转移至前景 NeRF 坐标系：$P_{\text{nv}} = T \cdot P_{\text{nv}}^{(d)}$。

\textbf{训练策略与损失函数。}ExpanDyNeRF 优先优化前景表示（因仿射畸变导致近处物体在视角旋转时模糊更严重）。时序连续性约束为 $L_{\text{cont}} = \|\sigma_f(t+1) - \sigma_f(t)\|^2$。主视角重建损失 $L_{\text{rec}} = \sum_{i=1}^N \sum_{r \in R} \|\hat{C}(r) - C(r)\|_2^2$ 优化 $\phi_b$ 和 $\phi_f$，使主视角渲染接近真值。新视角损失分解为颜色损失和密度损失两部分：

\begin{equation*}
L_{\text{nv}} = L_{\text{nv}}^c + L_{\text{nv}}^\sigma = \sum_{r \in R_{\text{nv}}} \left(\|\hat{C}(r_{\text{nv}}) - C(r_{\text{nv}})\|_2^2 + \|\hat{\sigma}(r_{\text{nv}}) - \sigma(r_{\text{nv}})\|_2^2\right)
\end{equation*}

其中 $C(r_{\text{nv}})$ 和 $\sigma(r_{\text{nv}})$ 分别为伪真值的颜色和密度。每次训练迭代中，每帧随机采样两个对称新视角，从相机位姿采样射线集合 $R_{\text{nv}}$，评估前景 NeRF 网络 $\phi_f$ 获得颜色和密度预测。颜色损失权重 $\lambda_{\text{nv}}^c = 1.0$、密度损失权重 $\lambda_{\text{nv}}^\sigma = 0.1$。超分辨率损失通过预训练 VGG-19 的多层特征增强纹理细节，同时包含重建损失和感知损失：

\begin{equation*}
L_{\text{sr}} = \sum_{k=1}^K \|p_k - Q_k\|_1 + \sum_{k=1}^K \sum_l \lambda_l \|F_{\text{vgg}}^l(p_k) - F_{\text{vgg}}^l(Q_k)\|_1
\end{equation*}

其中 $p_k$ 和 $Q_k$ 分别为超分辨率预测和参考图块，$F_{\text{vgg}}^l$ 为 VGG-19 第 $l$ 层特征，$\lambda_l$ 为该层神经元数量的倒数，$\lambda_{\text{sr}} = 0.5$。总损失为 $L = L_{\text{cont}} + L_{\text{rec}} + L_{\text{sr}} + L_{\text{nv}}$，其中 $L_{\text{nv}}$ 延迟若干轮后才引入以防止梯度爆炸。

\textbf{SynDM 数据集。}该方法还贡献了 SynDM（Synthetic Dynamic Multiview）数据集——基于 GTA V 构建的首个具有显式侧视监督的动态场景合成多视角数据集。核心挑战在于 GTA V 仅支持单一活动视口：在 60Hz 刷新率下，每次相机切换至少产生 16.7ms 延迟，随着相机数量增加累积延迟导致运动错位和时间伪影。为此，作者开发了自定义插件，在游戏图形状态半冻结（允许渲染和物理引擎继续运行）期间循环切换相机视角，将每次切换延迟从 16.7ms 降至 0.2ms，实现 22 个同步相机的动态场景采集。数据集包含 9 个场景，涵盖人类、车辆和动物三类，分布于乡村和城市环境。22 台相机的配置为：19 台沿水平方向以 $5^\circ$ 间隔分布于 $-45^\circ$ 至 $45^\circ$ 范围（含中央锚定相机），3 台在垂直方向上分布于 $-45^\circ$、$0^\circ$、$45^\circ$。所有帧以 $1920 \times 1080$ 分辨率渲染，水平视场角 $90^\circ$、垂直视场角 $59^\circ$。现有数据集（DyNeRF、NVIDIA Dynamic Scene 等）均缺乏偏离视角的真值（deviated view ground truth），严重限制了大视角偏差下的定量评估，而 SynDM 唯一同时具备多视角数据、偏离真值、全场景表示和相机运动四项关键特性。

\textbf{实验验证。}ExpanDyNeRF 在三个数据集上进行评估，每个场景使用 2 块 A100 训练 30 万迭代（约 15 小时）。在 SynDM 数据集上（训练 24 帧，在 $-30^\circ$ 至 $+30^\circ$ 间以 $5^\circ$ 间隔均匀采样 12 个新视角评估），ExpanDyNeRF 取得 PSNR 20.86 的最高分，LPIPS 和 FID 分别比第二名低 38\% 和 36\%，在感知对齐和分布一致性方面显著领先。值得注意的是，PSNR 在评估感知质量时存在固有局限：尽管 ExpanDyNeRF 产生更清晰、更细节的重构，像素级误差热图中的局部高密度误差可能导致 PSNR 反而偏低；而基线方法的模糊结果因平滑过渡反而产生更低的 MSE。因此 LPIPS 和 FID 等互补感知指标对全面质量评估至关重要。各基线表现出不同的失败模式：DecNeRF 渲染细节清晰但深度感知差，产生类似纸板的平面外观；D3DGS 在 $30^\circ$ 视角变化后出现前景断裂；RoDynRF 难以保持一致的物体位置，常产生浮游伪影或身体部位错位。

在 DyNeRF 数据集上（从中央相机跨时间戳构建合成单目序列，保留外侧偏离相机 cam01 和 cam10 作为测试），ExpanDyNeRF 在 Coffee 场景取得 FID 132.4、PSNR 30.32、LPIPS 0.189，在 Beef 场景取得 FID 135.8、PSNR 34.92、LPIPS 0.195，全面领先所有基线。在 NVIDIA 数据集上（12 台相机呈紧凑矩阵排列，视角偏差有限），ExpanDyNeRF 在 Skate 场景取得 FID 90.83、PSNR 28.91、LPIPS 0.079，在 Truck 场景取得 FID 69.37、PSNR 30.60、LPIPS 0.034，稳定获得第二名——这反映该数据集主要面向主视角优化，DetRF 因专注于前向相机深度估计而在此占优，而 ExpanDyNeRF 的均衡表现意味着在训练视角之外具有更稳定的泛化行为。

消融实验验证了各损失组件的贡献：移除新视角监督后物体形状模糊不清；仅使用颜色损失 $L_{\text{nv}}^c$ 因缺乏形状约束而产生白色伪影；仅使用密度损失 $L_{\text{nv}}^\sigma$ 虽保持物体形状但颜色褪色或失真；三者协同（$L_{\text{nv}}^c + L_{\text{nv}}^\sigma + L_{\text{sr}}$）方能实现最优感知质量。单独使用超分辨率损失 $L_{\text{sr}}$ 无法产生可接受的渲染，它主要提升感知锐度而非引入新的几何约束，缺乏新视角损失建立的几何一致性基础时倾向于放大现有结构误差。

\textbf{局限性。}ExpanDyNeRF 在处理极端视角（超过 $45^\circ$）和未见背景生成方面仍存在不足，且需要逐场景优化，计算成本高于前馈或纯高斯方法。

\subsection{泛化模型与位姿估计}
\label{ss:06_3c}

本节介绍四条技术路线：（1）前馈与泛化模型；（2）位姿估计与自标定；（3）专用场景与特殊条件；（4）高效表示与运动建模。

\subsubsection{前馈与泛化模型}

本路线介绍通过大规模预训练实现任意输入视频秒级推理的前馈动态重建方法。

正则化、深度先验、运动模糊处理和事件融合分别从不同角度改善了单目重建质量。然而，这些方法仍然需要逐场景优化——每个视频训练数十分钟到数小时。前馈方法通过在大规模多视角/4D 数据上预训练，实现对任意输入视频的秒级甚至毫秒级推理，是走向实际部署的关键方向。

\paragraph{子弹时间前馈多视角重建}

BTimer（2024.12）~\cite{2412_03526} 是首个运动感知前馈动态场景重建模型，其核心创新在于``子弹时间''（bullet-time）公式：给定一组上下文帧及其相机位姿和时间戳，模型在指定的子弹时间戳 $t_b$ 重建完整场景的 3D 高斯 Splatting 表示——即将场景``冻结''在 $t_b$ 时刻的 3D 状态。迭代设置 $t_b$ 为视频中每个时间戳即可高效并行重建完整 4D 视频。关键设计是每个输入帧携带两种时间嵌入：上下文时间嵌入 $\tau_i^{\text{ctx}}$（编码该帧自身的时间戳 $t_i$）和子弹时间嵌入 $\tau_i^{\text{bullet}}$（编码目标重建时间 $t_b$，所有上下文帧共享）。两种时间戳先经标准正弦位置编码（sinusoidal positional encoding），再分别通过两层线性层获得对应特征。两者相加形成最终时间特征 $\tau_i^{\text{time}} = \tau_i^{\text{ctx}} + \tau_i^{\text{bullet}}$，使模型在处理每帧时同时知道``这是什么时候的''和``要重建什么时候的''，从而学会根据目标时间戳聚合上下文信息——动态区域在 $t_b$ 附近的帧被赋予更高权重，远离 $t_b$ 的帧提供静态上下文。

这一公式天然统一了静态和动态重建：当所有时间戳相等时即退化为静态场景。
模型架构基于 24 层 ViT Transformer，首尾各附加 LayerNorm。输入帧被切分为 $8 \times 8$ 图块，经线性投影后与 Plucker 相机位姿嵌入和时间特征相加形成输入 token $\mathbf{f}_{ij} = \mathbf{f}_{ij}^{\text{rgb}} + \mathbf{f}_{ij}^{\text{pose}} + \mathbf{f}_i^{\text{time}}$。所有上下文帧的 token 拼接后通过 Transformer（实现中采用 FlashAttention-3 和 FlexAttention 加速），每个输出 token 经单层线性层解码为 3DGS 参数（RGB 颜色 $c$、缩放 $s$、旋转四元数 $q$、不透明度 $o$、射线距离 $d$，共 12 维），3D 位置通过像素对齐反投影 $\mu = o + d \cdot d_{\text{ray}}$ 获得，其中 $o$ 和 $d_{\text{ray}}$ 来自相机位姿 $P_i$ 的射线原点和方向。渲染使用 gsplat 以支持大规模 3DGS 光栅化。损失函数为 MSE 与 LPIPS 的加权和 $\mathcal{L}_{\text{RGB}} = \mathcal{L}_{\text{MSE}} + 0.5 \cdot \mathcal{L}_{\text{LPIPS}}$。

训练中采用两种互补的监督策略：（1）\textbf{上下文内监督}（in-context supervision）：监督时间戳随机选自上下文帧，鼓励模型精确定位和重建已观测时刻；多视角视频数据中还可纳入额外视角作为监督。（2）\textbf{插值监督}（interpolation supervision）：监督时间戳位于两个相邻上下文帧之间，迫使模型学会插值动态部分同时保持静态区域一致。消融实验表明，缺少插值监督时模型会陷入局部最优——将 3D 高斯放置在靠近上下文视角但被其他视角遮挡的位置。
为处理快速运动场景中的中间时刻重建（$t_b \notin \mathcal{T}$），BTimer 引入 Novel Time Enhancer（NTE）模块。NTE 复制 BTimer 的 ViT 架构，但输入 token 的时间特征仅使用上下文时间嵌入（$\mathbf{f}_i^{\text{time}} = \mathbf{f}_i^{\text{ctx}}$），额外拼接编码目标时间戳和位姿的目标 token，并使用 QK-norm 稳定训练。NTE 直接在图像空间预测目标时刻的 RGB 图像，通过注意力掩码实现 KV-Cache 加速推理，推理开销可忽略不计。

预测的中间帧随后作为上下文帧送入主模型，弥补像素对齐 3DGS 预测在复杂运动下的插值不足。单独使用 NTE 进行新视角合成的质量不如 BTimer，两者级联才能发挥最大效果。
训练采用三阶段课程学习策略，在 32 块 NVIDIA A100 GPU 上进行，三个阶段分别使用余弦退火学习率 $4 \times 10^{-4}$、$2 \times 10^{-4}$、$1 \times 10^{-4}$。（1）\textbf{静态预训练}（90K + 90K + 50K 迭代）：在 Objaverse、RE10K、MVImgNet、DL3DV 等约 39 万样本上从 $128 \times 128$ 低分辨率 4 视角开始，逐步提升至 $512 \times 512$，不使用时间嵌入。（2）\textbf{动态共训练}（10K 迭代）：引入 Kubric、PointOdyssey、DynamicReplica、Spring 等 4D 数据集，开启时间嵌入并与静态数据联合训练。（3）\textbf{长上下文微调}（5K 迭代）：将上下文帧数从 4 增加至 12 以覆盖更长视频。

此外，从 PANDA-70M 数据集中截取约 20 秒短视频，用 Segment Anything Model 掩码动态物体后通过 DROID-SLAM 估计相机位姿，再以重投影误差过滤低质量数据，最终获得超过 4 万段带高质量相机轨迹的互联网视频作为补充训练数据。最终模型在 12 个上下文帧、$256 \times 256$ 分辨率下仅需 150ms 单 GPU 推理（4 视角仅需 20ms），12 视角 $512^2$ 推理约 1.55s，显存占用低于 10GB，输出的 3DGS 以 115 FPS 实时渲染。BTimer 在 DyCheck 和 NVIDIA Dynamic Scene 基准上即使与逐场景优化方法相比也达到领先性能——在 NVIDIA 数据集上 PSNR 比显式 3DGS 方法高出约 5\%（25.82dB vs 24.57dB），且重建速度从数小时级压缩至亚秒级。

\paragraph{可泛化循环单目动态视角合成}

现有可泛化静态 NVS 模型（如 ConvGLR）通过平面扫描体（PSV）编码相机相对位姿，直接预测目标视角图像，无需中间 3D 表示。GRVS（2026.03）~\cite{2603_29734} 将这一范式扩展至动态场景，核心动机在于：单目动态视频中相机运动与场景运动高度纠缠，而 PSV 的同态投影恰好能将相机位姿信息从场景内容中分离出来，使模型集中精力预测动态元素。这一设计选择也是对当前两类主流范式的反思——场景特定优化方法（如 D-3DGS、SC-GS、4DGS、MoSca）虽能良好重建静态区域，但每场景需 20--40 分钟优化且在动态元素上表现不佳；扩散方法（如 GCD 通过直接注入相机参数到 Stable Video Diffusion、Gen3C 通过 ViPE 预测点云再作为扩散模型输入）虽能产生视觉合理的视频，但几何精度不一致，且依赖庞大的预训练模型（Gen3C 达 7B 参数）。GRVS 的核心主张是：大多数目标视角中的视觉内容已在输入视角中可见，生成式需求仅限于小范围遮挡区域，因此一个轻量级的可泛化前馈模型即可胜任，无需调用大规模扩散模型。

为此，GRVS 引入三个关键设计。（1）异步映射：将动态 NVS 建模为输入视频时间步 $t$ 与目标视频时间步 $i$ 之间任意的映射 $i \mapsto t_i$，突破 GCD 等方法 $t_i = i$ 的同步约束，支持子弹时间等自由视角 4D 渲染。映射仅需满足时序约束 $t_i \in \{t_{i-1}-1, t_{i-1}, t_{i-1}+1\}$。子弹时间是异步映射的特殊情形：$t_i$ 对所有 $i$ 保持常数，动态场景在输出视频中呈现冻结状态。（2）循环潜变量渲染：模型在连续预测目标帧时传递隐藏状态 $Z_i$，形成循环回路 $\hat{J}_i, Z_i = M(I_i, P_i, Q_i, Z_{i-1})$，其中前一帧的潜变量通过同态变换 $H_{Q_{i-1} \to Q_i}$ 重投影至当前目标视角，实现渲染结果的高效复用。循环机制在三个互补层面发挥作用：通过持续复用已生成渲染结果提升效率；通过控制输入采样率同时利用全局和局部时间上下文；通过帧间隐藏状态传递确保生成视频的时间一致性。循环架构的灵感来自 RAFT 等迭代光流估计、DROID-SLAM 等递归 SLAM 以及 BasicVSR 等视频去噪/超分方法。

模型外围还引入 Patchify 和 Unpatchify 模块：Patchify 以步长为 $F$ 的卷积核将 PSV 划分为 patch（$F=2$ 时 FLOP 削减约 75\%），Unpatchify 以 $1 \times 1$ 卷积扩展至 $3 \times F \times F$ 通道后还原为全分辨率输出，二者协同使高分辨率输入和训练时的多次循环展开在 GPU 内存上可行。（3）迭代多尺度采样：每个目标帧执行三次迭代，膨胀因子 $d$ 从 $5 \to 3 \to 1$ 递减——大膨胀提供全局上下文以重建静态区域，小膨胀提供局部上下文以捕获动态细节。消融实验表明，无循环模型在 $d$ 减小时全局性能下降、仅动态区域提升；而循环模型在 $d$ 减小时全图和动态区域性能同步提升——第三次迭代时循环模型在全图上比无循环模型高出约 1 dB，证明隐藏状态有效传播了全局时间信息。PSV 配置为 $D=32$ 个深度平面、$V=15$ 个输入视角、通道数 $C=256$、下采样因子 $F=2$，模型以高效 3D U-Net 在低分辨率潜空间完成全局渲染。

训练方面，GRVS 采用标准 VGG 损失在 $320 \times 320$ 像素 patch 上端到端训练，最后 10\% 训练步切换为 L1 损失以避免网格伪影；使用 Adam 优化器训练 100k 步，批量大小为 4，4 块 32GB GPU 上约 4 天完成。训练时各迭代的循环隐藏状态在迭代间 detach 以防止 GPU 内存溢出。对每个目标时间步 $t_i$，独立采样一组输入图像 $I_i$，并以其中间帧与目标帧配对进行监督。

GRVS 在 UCSD 数据集和新提出的 Kubric-4D-dyn 数据集（5000 训练场景，100 测试场景，$512 \times 512$ 分辨率，含动态相机）上均优于 Deformable 3DGS 等场景特定方法和 GCD 等扩散方法。UCSD 实验中，D-3DGS、SC-GS、4DGS 三种场景特定方法因合成序列仅 10 个视角的稀疏性而在静态和动态区域均表现不佳，而 GRVS 无需场景特定优化即与 DynIBaR 性能持平，但推理速度数量级更快（DynIBaR 在 10 秒视频上需 8 块 A100 约 2 天，且依赖深度图、光流、动静态掩码等多种显式先验，而 GRVS 仅需 RGB 图像和相机参数）。Kubric-4D-dyn 实验中，四种 GS 方法（D-3DGS、SC-GS、4DGS、MoSca）在静态区域重建良好但动态区域随目标视角距离增大而显著退化，MoSca 在动态物体上表现更优但全图指标略降；GCD 和 Gen3C 两种扩散方法生成视觉合理的动静态物体但 3D 几何、纹理和相机位置精度不足。

GRVS 在三项指标（PSNR/SSIM/LPIPS）上均超越全部基线，全图和动态区域分别最优。值得注意的是 GRVS 模型仅 40M 参数、3 FPS 推理，而 Gen3C 达 7B 参数、0.3 FPS，凸显了可泛化前馈模型在效率上的巨大优势。Kubric-4D-dyn 数据集的设计也回应了此前 Kubric 变体的局限：DyST 数据集仅 $128 \times 128$ 分辨率且每场景仅单个物体，GCD 数据集（3000 场景，$576 \times 384$）虽有更多物体但仅含静态相机，而 GRVS 引入的动态相机使 4D 重建方法能够利用多视角线索估计场景几何。每个场景包含 40 个静态物体和 5 个具有初始速度的动态物体，摩擦力设为最低以保持全程运动，3 台相机中 1 台为沿虚拟半球运动的动态输入相机、2 台为静态目标相机。

\paragraph{分钟级 4D 高斯泼溅}

现有 4DGS 方法通常依赖耗时的 COLMAP 预处理，且在单目设置下各向异性高斯容易过拟合。Instant4D（2025.10）~\cite{2510_01119} 的设计动机正是解决这两个瓶颈：用深度视觉 SLAM 替代 SfM 实现快速几何初始化，用各向同性 4D 高斯基元替代复杂参数化以提升数值稳定性。该工作已被 NeurIPS 2025 接收。

整个流程分为几何恢复和运动建模两阶段。几何恢复阶段，Instant4D 首先利用 MegaSAM 深度视觉 SLAM 估计相机轨迹 $\{\xi_i \in SE(3)\}$ 和内参 $K$，并通过 Levenberg-Marquardt 算法将单目深度（由 Depth Anything 初始化、Unidepth 提供焦距初值）精炼为时序一致深度图 $\mathcal{D}$；随后将每帧像素反投影为稠密 3D 点云。为分离动静态区域，Instant4D 利用深度 SLAM 管线中间输出的低分辨率运动概率图 $\hat{m} \in \mathbb{R}^{H/8 \times W/8}$，插值至逐像素分辨率后以 Otsu 阈值法生成二值掩码。该策略的一个关键工程细节是：对于时间采样间隔较大的序列，首尾帧的运动估计往往不可靠，因此 Instant4D 在序列两端引入合成伪帧以提升运动一致性。

体素网格剪枝策略将世界空间划分为正则体素网格，自适应边长为 $S_v = \lambda_s \cdot \frac{1}{N}\sum_i \bar{D}_i / f$，其中 $\lambda_s$ 为用户定义的缩放因子。每个被占据体素仅保留内点质心，且颜色、时间戳、时间缩放 $s_t$ 和运动概率 $\hat{m}$ 等属性也取均值聚合；支撑不足的体素作为离群点被丢弃以抑制噪声。在 NVIDIA Dynamic Scene 基准上，该策略将模型内存从 10.7 GB 降至 0.83 GB（削减 92\%），训练时间从 181 秒降至 42 秒（4 倍加速），渲染帧率从 154 FPS 提升至 981 FPS。

运动建模阶段，Instant4D 采用原生 4D 高斯基元，每个高斯由 4D 均值 $\mu = (\mu_x, \mu_y, \mu_z, \mu_t)$、对角缩放 $s = (s_{xyz}, s_t)$、标量不透明度 $\alpha$ 和旋转矩阵 $R$ 参数化。与先前 4DGS（2023.10）~\cite{2310_08528} 依赖两个纠缠四元数和高阶4D球面柱谐函数（4DSH）不同，Instant4D 直接用简单 RGB 值建模外观，而非高阶球谐函数。这一简化设计将每高斯参数量减少超过 60\%，并在实验中有效减轻了单目设置下的过拟合倾向。在渲染时刻 $t$，4D 高斯被条件化为 3D 高斯：

\begin{equation*}
\mu_{xyz|t} = \mu_{1:3} + \Sigma_{1:3,4}\,\Sigma_{4,4}^{-1}(t - \mu_4), \quad
\Sigma_{xyz|t} = \Sigma_{1:3,1:3} - \Sigma_{1:3,4}\,\Sigma_{4,4}^{-1}\,\Sigma_{4,1:3}
\end{equation*}

这一条件化形式将连续运动隐式编码在 4D 协方差中，无需显式存储轨迹。为进一步提升稳定性，Instant4D 采用各向同性设定（$R = I$），协方差仅由两个标量——共享空间缩放 $s_{xyz}$ 和时间缩放 $s_t$——参数化，在充当隐式正则化的同时降低内存开销。运动感知设计的核心在于不透明度的时间衰减：$o_t = o \times \mathcal{N}(t; \mu_4, \Sigma_{4,4})$，其中 $\Sigma_{4,4} = s_t^2$ 是各向同性设定下唯一影响时间属性的项。对静态区域，时间缩放设为视频全长 $s_t = l_{\mathrm{video}}$，使静态高斯在整个时间范围内保持可见；对动态区域，$s_t = 2/\mathrm{fps}$，高斯仅在对应时刻附近活跃。渲染时，距离当前时刻较远的高斯因不透明度低于阈值而被剔除。模型提供两种配置：Lite 版在静态和动态区域均使用体素尺寸 $\lambda_s = \lambda_d = 4$；Full 版对静态区域使用 $\lambda_s = 1$，但跳过动态区域的网格剪枝以保留细节。

在性能方面，Instant4D 在单张 NVIDIA A6000 GPU 上完成完整训练管线：Lite 模型在 NVIDIA 数据集上平均 2 分钟完成重建，最短序列（235 帧 ``paper-windmill''）仅需 96 秒/988 MB 峰值内存，最长序列（379 帧 ``apple''）需 131 秒/1147 MB。在 DyCheck 数据集上平均 7.2 分钟完成，相比当前最优方法提升 29\%。推理速度超过 400 Hz。整体实现 30 倍重建加速和 90\% 内存缩减，体素剪枝阶段每场景不到 5 秒。

\subsubsection{位姿估计与自标定}

本路线介绍联合优化相机位姿和场景表示以摆脱对COLMAP依赖的自标定方法。

大多数动态 3DGS 方法依赖 COLMAP 提供的相机位姿和稀疏点云。然而，对于动态场景，COLMAP 经常因运动物体干扰而失败，因此联合优化相机位姿和场景表示成为一个重要研究方向。

\paragraph{自标定单目相机 4D 重建}

\textbf{改动概要。}本方法修改了动态 3DGS 流水线的位姿获取环节——用"结构点提取 + 相机参数联合优化"替代"依赖 COLMAP 预计算位姿"——下游的规范场-变形场表示框架保持不变。直观地说，COLMAP 就像一位只认得静态世界的测量员：场景中有运动物体时，它会因为分不清"是物体动了还是相机动了"而给出错误的位姿。SC-4DGS 的做法是不再请这位测量员，而是自己从视频中挑出最可靠的静态锚点，根据这些锚点反推相机在哪里。

\textbf{Baseline 的问题与设计动机。}前述变形场方法（如 Deformable 3DGS、4D-GS）都依赖 COLMAP 提供相机外参和稀疏点云作为初始化。这一依赖在动态场景中构成严重瓶颈。COLMAP 的 SIFT 特征提取-匹配-重建流水线本质上为静态场景设计，无法处理随时间变化的 3D 参数。实验表明，在 DAVIS 数据集上，COLMAP 在超过 80\% 的场景中完全失败——即使提供了真实运动掩码帮助排除动态区域，COLMAP 仍然崩溃。进一步分析发现，COLMAP 的 SIFT 特征提取步骤本身是成功的，失败始于穷尽特征匹配（exhaustive feature matching）步骤。现有的自标定方法同样受限：NeRFmm 要求相机前向拍摄且旋转范围限制在 $\pm 20°$；Nope-NeRF 和 CF-3DGS 需要预先提供相机内参，且依赖 MiDaS 单目深度估计作为额外监督，单场景优化耗时约 30 小时；RoDynRF 虽然能处理动态场景，但需同时引入 MiDaS 深度和 RAFT 光流两个外部模型，50--80 帧视频的训练时间超过 28 小时，且在超过 800 帧的长视频上完全失败。SC-4DGS（2024.06）~\cite{2406_01042} 的设计目标是同时克服上述八项限制：不准确的场景特征提取（L1）、训练时间过长（L2）、仅适用于静态场景（L3）、无法处理极端几何条件（L4，即大物体运动+微小相机平移+巨大相机旋转）、需要预设相机内参（L5）、需要外部 3D 先验模型（L6）、需要多模型监督（L7）、长视频失败（L8）。

\textbf{实现原理。}方法分为三步，核心思路是"先找可靠的静态锚点，再用锚点反推相机，最后用相机做 4D 重建"。

\textit{第一步：结构点提取（Structural Point Extraction, SPE）。}此步的目的是从视频中提取可靠的静态 2D-3D 对应关系，作为后续相机标定的锚点。为什么不直接用 CoTracker 追踪随机选取的像素？论文的消融实验（论文 Fig.6）给出了直观证据：在 NeRF-DS 的 basin 场景中，即使最先进的密集追踪模型 CoTracker 也无法准确追踪低纹理区域（如白色墙壁，因局部特征高度相似而漂移至错误位置）和高纹理但特征混淆的区域（如树叶，众多点具有混淆性的相似特征）中的点。相比之下，SPE 通过三层过滤策略确保仅保留最可靠的候选点。

首先，计算 RGB 三通道梯度幅值 $Grad_i^{magni}$ 以定位高频像素（梯度幅值越大，纹理越丰富，追踪越可靠）。然后，通过 Canny 边缘检测 $E_i^{gray}$ 排除低纹理区域像素——尽管梯度幅值的局部极大值能指示高频像素，但低纹理区域中仍存在误检点，Canny 检测器根据灰度强度进一步剔除这些点。两层过滤取交集后再经 Maximum Filter（窗口尺寸 $\mathcal{W}$）在局部邻域内提取极大值点，得到最具代表性的候选像素。最后，用运动掩码 $M_i^{motion}$ 过滤变形物体上的点，因为动态区域的点不适合用于估计相机参数。三层过滤后的候选池记为 $SP_i^{Pool} \in \mathbb{R}^{B \times 2}$（$B$ 因帧而异）。随后用 CoTracker 逐帧追踪候选点的可见性 $Pred^{vis}$ 和位置 $Pred^{pos}$：当某点在后续帧中被遮挡（$Pred_p^{vis} = 0$）或落入运动掩码区域（$M_p^{motion}[Pred_p^{pos}] = 0$）时，标记为缺失；当剩余可用点数首次降至阈值以下时，从上一帧剩余候选点中随机选取补充，分配新的全局索引。如此逐帧迭代直至所有帧的 $SP_i^{2D} \in \mathbb{R}^{\alpha \times 2}$ 和 $SP_i^{Index}$ 全部填满（默认 $\alpha = 100$ 个 2D 结构点），最终构建完整的跨帧 2D-3D 对应关系。SPE 初始化时设 $H = 0$，随着新 3D 结构点的引入逐步递增 $H$，因此总参数量 $7N + 1 + H$ 也是动态增长的。

\textit{第二步：相机参数与 3D 结构点联合优化。}此步用 SPE 提取的 2D-3D 对应关系同时优化相机外参和 3D 结构点坐标。假设整段单目视频共享固定焦距 $f$，每帧定义四元数旋转 $Quat_i \in \mathbb{R}^4$（可转化为旋转矩阵 $R_i \in \mathbb{R}^{3 \times 3}$）、平移 $T_i \in \mathbb{R}^3$，以及世界坐标系中的共享 3D 结构点 $SP^{3D} \in \mathbb{R}^{H \times 3}$。核心约束是将 3D 结构点投影至图像平面后应与 SPE 提取的 2D 点对齐，投影过程为：

\begin{equation*}
SP_i^{Cam} = \text{Homo}(SP^{3D}[SP_i^{Index}]) \cdot W2C_i^T \cdot PP^T, \quad SP_i^{2D_{proj}} = SP_i^{Cam}[:, :2] \,/\, SP_i^{Cam}[:, 3]
\end{equation*}

其中 $\text{Homo}(\cdot)$ 将 3D 点转换为齐次坐标（拼接 $1$），$W2C_i \in \mathbb{R}^{4 \times 4}$ 为世界到相机变换矩阵，$PP \in \mathbb{R}^{4 \times 4}$ 为透视投影矩阵。投影损失 $L_{proj} = \sum_{i=0}^{N-1} \text{MSE}(SP_i^{2D_{proj}} - SP_i^{2D})$ 用均方误差衡量投影偏差。在此基础上，距离误差损失 $L_{de} = \sum_{i=0}^{N-1} \text{MSE}(\text{Dist}(SP_i^{2D_{proj}}) - \text{Dist}(SP_i^{2D}))$ 约束投影点对之间的距离与真实 2D 点对之间的距离一致（保持 3D 结构的相对几何关系），深度正则化损失 $L_{dr} = \sum_{i=0}^{N-1} \text{ReLU}(-SP_i^{Cam}[:, 3])$ 用 ReLU 惩罚相机坐标系中深度为负的 3D 结构点（防止点跑到相机后方）。三者构成总标定损失：

\begin{equation*}
L_{Cali} = L_{proj} + L_{de} + L_{dr}
\end{equation*}

共优化 $7N + 1 + H$ 个参数（$N$ 帧各 4 个四元数 + 3 个平移 = $7N$，共享焦距 1 个，$H$ 个 3D 结构点各 3 坐标）。这一优化无需任何深度先验或光流监督，仅依赖 SPE 提取的对应关系——这也是 SC-4DGS 不受外部模型精度影响（L6、L7）的关键。

\textit{第三步：4D 场景表示优化。}此步将优化后的 3D 结构点 $SP^{3D_O}$ 作为初始点云（而非 COLMAP 的稀疏点云或随机密集点云），沿用 §\ref{sec:06_1} 中介绍的规范场-变形场框架进行动态重建。具体而言，建立规范场 $G_C$ 学习均值位置 $x$、四元数 $r$、缩放 $s$、不透明度 $\sigma$，以及变形场 $G_D$ 学习时变偏移 $\Delta x, \Delta r, \Delta s$：

\begin{equation*}
x_i, r_i, s_i = G_D(P(X_{G_C}), P(T_i)), \quad I_i^{Render} = \mathcal{R}(x + \Delta x_i, r + \Delta r_i, s + \Delta s_i, \sigma)
\end{equation*}

其中 $X_{G_C} = G_C(SP^{3D_O})$ 为规范场从结构点学习到的高斯中心，$P$ 为位置编码（positional encoding），$T_i$ 为帧时间，$\mathcal{R}$ 为 3DGS 的可微高斯泼溅渲染器。渲染损失采用标准的 RGB 损失 $L_{RGB} = (1-\lambda)L_1 + \lambda \mathcal{L}_{D\text{-}SSIM}$（$\lambda = 0.2$）。

\textbf{性能与验证。}SC-4DGS 在三个标准基准上验证了自标定 + 4D 重建的联合效果。在 NeRF-DS 数据集（长视频、小物体运动、大相机运动）上，自标定位姿的精度与 COLMAP 相当（以 COLMAP 位姿为真值计算的 ATE 和 RPE 指标），RoDynRF 则在此数据集上估计失败。优化后的点云空间分布更集中、几何连续性更好——论文从几何角度解释了浮游点的成因：当相机位姿存在误差时，同一 3D 点在多帧中的观测射线无法准确三角化（triangulate），RGB 损失最小化会驱使自适应密度控制（Adaptive Density Control）在每帧中添加额外高斯以补偿位姿偏差，这些高斯远离相机时尤为密集，形成浮游伪影。在 DAVIS 数据集（COLMAP 失败率超 80\% 的极端几何场景）上，SC-4DGS 和 RoDynRF 均能提供合理的位姿估计，两者估计的位姿在大多数场景中相似且均产生高质量渲染，但在出现差异的少数场景中，SC-4DGS 的渲染保真度更高。

与 RoDynRF 相比，SC-4DGS 无需 MiDaS 深度和 RAFT 光流监督，训练时间大幅缩短（RoDynRF 对 50--80 帧视频需 28 小时以上），且不受长视频限制。在 NeRF-DS 的渲染评测中，SC-4DGS 与使用 COLMAP 位姿的 Deformable 3DGS 达到可比的 PSNR/SSIM/LPIPS——例如 plate 场景中 Deformable 3DGS 的 PSNR 仅高 0.005 dB，但 SC-4DGS 能渲染出更精细的光照和阴影细节，而 RoDynRF 的 PSNR 比两者低约 10.00、SSIM 低约 0.20、LPIPS 高约 0.35。

\paragraph{运动感知无位姿动态重建}

直接将 MonST3R 的位姿估计结果喂入 4DGS 往往因位姿精度不足导致渲染质量低，尤其在运动物体占据画面大部分时，RANSAC 将动态区域误判为离群点，位姿估计严重退化。一个直接的强基线组合——先用 MonST3R 估计位姿再送入 SC-GS 重建——在 HyperNeRF 数据集上仅获得 20.4 dB PSNR 和 0.697 MS-SSIM，而使用 COLMAP 位姿的 SC-GS 可达 25.3 dB / 0.841，差距高达 4.9 dB。这一巨大鸿沟表明位姿估计与场景重建在动态场景中存在深度耦合，无法独立处理。4D3R（2025.11）~\cite{2511_05229}（NeurIPS 2025，新加坡国立大学与武汉大学）的核心洞察是：运动掩码和场景重建应当联合优化而非分开处理，使动态-静态分解信息在位姿估计和高斯重建之间双向流动。

方法分三个模块。第一个模块是 4D 感知信息提取：利用 MonST3R 的 ViT 编码器-解码器提取场景坐标图 $X_t \in \mathbb{R}^{H \times W \times 3}$ 和置信度图 $W_t$，同时引入 SEA-RAFT 光流提供跨帧对应关系。通过双重过滤 $S = \{p_i \mid W_t(p_i) > \tau_c \wedge D_t(p_i) < \tau_d\}$ 筛选高置信度点——高置信度分确保几何估计可靠，深度阈值排除无穷远处的不可靠点——再从置信度图中采样 Top-K 点作为 SAM2 的多点提示，利用 SAM2 的语义分割能力生成像素级动态物体分割掩码 $M_t$，其边界精度显著优于纯运动估计器的粗糙输出。

第二个模块是运动感知光束平差（MA-BA）：传统光束平差在动态区域失效，MA-BA 利用掩码 $M_t$ 将对应点匹配限制在静态区域 $P_{static} = \{p_i \in S \mid M_t(p_i) = 0\}$，优化目标为 $E(R_t, T_t) = \sum_{p_i \in P_{static}} \|(\pi(R_t p_i + T_t)) - p'_i\|^2$。RANSAC 阶段仅在静态区域建立 2D-3D 对应，避免动态物体引入错误匹配。随后通过可微密集光束平差（DBA）层进一步精炼位姿和深度，其优化目标在标准 DBA 基础上加入运动掩码加权：

\begin{equation*}
E_{DBA}(C_t', d_t') = \sum_{(i,j) \in \mathcal{E}} (1 - M_t(i)) \|p_{ij}^* - \pi_e(C_{ij}' \circ \pi_e^{-1}(p_i, d_i'))\|_{\Sigma_{ij}}^2
\end{equation*}

其中运动掩码项 $(1 - M_t(i))$ 确保仅静态区域参与优化，$p_{ij}^* = p_{ij} + r_{ij}$ 为光流预测的对应点，$\|\cdot\|_{\Sigma_{ij}}$ 为置信度加权的马氏距离。DBA 层通过稀疏矩阵法方程求解位姿增量 $\Delta\xi_t$ 和深度增量 $\Delta d_t$，实现端到端可微的位姿精化。

第三个模块是运动感知高斯 Splatting（MA-GS）：用 512 个控制点 $P = \{(p_i, \sigma_i)\}$ 参数化动态变形，每个控制点有规范空间坐标 $p_i$ 和 RBF 核半径 $\sigma_i$；变形场 MLP $\Theta$ 以控制点坐标和时间 $t$ 为输入，预测时变 SE(3) 变换 $[R_t^k \mid T_t^k]$（旋转用单位四元数 $r_t^k \in \mathbb{H}$ 表示以保证数值稳定性和连续插值）。高斯点通过线性混合蒙皮（LBS）加权变形，位置和旋转分别更新：

\begin{equation*}
\mu_j' = \begin{cases} \mu_j & M_t(\mu_j) = 0 \\ \sum_k w_{jk}(R_t^k(\mu_j - p_k) + p_k + T_t^k) & \text{otherwise} \end{cases}
\end{equation*}
\begin{equation*}
q_j' = \Big(\sum_{k \in \mathcal{N}_j} w_{jk}\, r_t^k\Big) \otimes q_j
\end{equation*}

其中 $\otimes$ 为四元数乘法，$\mathcal{N}_j$ 为高斯 $j$ 的 $K$ 近邻控制点集合。LBS 权重由归一化指数核计算：$\tilde{w}_{jk} = \exp(-d_{jk}^2 / 2\sigma_k^2)$，$w_{jk} = \tilde{w}_{jk} / \sum_k \tilde{w}_{jk}$，其中 $d_{jk}$ 为高斯中心 $\mu_j$ 到控制点 $p_k$ 的欧氏距离。静态区域的运动掩码约束 $L_{rigid} = \sum_j (1 - M_t(\mu_j))\|\mu_j' - \mu_j\|^2$ 确保静态点不动，ARAP 正则化 $L_{arap}(p_i, t_1, t_2) = \sum w_{ik}\|(p_i^{t_1} - p_k^{t_1}) - R_i(p_i^{t_2} - p_k^{t_2})\|^2$ 维持控制点间的局部刚性。

训练采用两阶段策略：第一阶段仅用 $L_{control} = \sum_{p_i \in P} M_t(p_i) \cdot L_{render}(p_i)$ 优化动态区域的控制点，运动掩码作为二值选择器将计算资源集中在需要变形建模的区域；第二阶段联合优化高斯参数与变形场，总损失为 $L = L_{render} + \lambda_{arap} L_{arap} + \lambda_{rigid} L_{rigid}$，其中 $L_{render}$ 使用 $L_1$ 距离和 DSSIM 度量。第二阶段中变形变换通过梯度分离（detached gradient）路径实施：MA-GS 的变换应用于高斯位置时断开梯度回传，防止变形场 MLP 和高斯参数之间的优化目标竞争，实现更稳定的收敛。此外引入自适应控制点策略：计算每个控制点 $k$ 对重建质量的影响梯度 $g_k = \sum \tilde{w}_j \|\partial L / \partial \mu_j\|^2$，在梯度高的区域增密控制点，自适应地匹配不同运动复杂度的需求。

\textbf{性能与验证。}在 HyperNeRF 数据集上，4D3R 以 25.6 dB PSNR 和 0.844 MS-SSIM 超过使用 COLMAP 位姿的 SC-GS（25.3 dB / 0.841），同时完全不需要预计算位姿；相比另一无位姿方法 RoDynRF（23.8 dB / 0.820），PSNR 提升 1.8 dB。在 DyNeRF 数据集上取得 19.6 dB PSNR 和 0.755 MS-SSIM，同样达到最优。训练效率方面，4D3R 在单块 RTX 3090 上仅需 50 分钟（含位姿估计），而 COLMAP 依赖方法通常需要 4--5 小时，RoDynRF 更需 28 小时以上，加速比达 5 倍。推理速度为 45 FPS，存储仅需 80 MB，远低于 MonST3R+SC-GS 的 153 MB 和 RoDynRF 的 200 MB。消融实验验证了三个核心组件的贡献：移除运动感知掩码后 PSNR 暴跌 5.2 dB（至 20.4 dB），证实动态-静态分解是处理大动态物体场景的基础；移除 SAM2 细化后下降 1.8 dB（至 23.8 dB），验证了 Transformer 运动先验与基础模型分割的协同效应；移除 MA-GS 模块后下降 2.1 dB（至 23.5 dB）且训练时间翻倍，确认紧凑控制点表示在效率和品质上的双重优势。

\paragraph{运动自适应样条插值动态 3DGS}

隐式变形场（MLP）和网格分解模型难以精确捕捉动态场景结构细节，多项式轨迹则受固定阶数限制无法灵活表达复杂运动——三阶多项式建模能力不足，十阶多项式则因高次幂变量引发数值不稳定。SplineGS（2024.12）~\cite{2412_09982}（CVPR 2025）借鉴经典计算机图形学中的样条曲线建模思想，提出用三次埃尔米特样条精确表示每个动态 3D 高斯的连续轨迹，仅需少量控制点即可建模任意复杂的运动路径。场景被分解为静态高斯集合 $\{G_i^{\text{st}}\}$ 和动态高斯集合 $\{G_i^{\text{dy}}\}$，静态区域移除时间编码特征仅保留漫反射和高光特征，动态区域则通过样条建模连续位置轨迹，旋转 $q$ 和尺度 $s$ 作为时间相关可学习参数。

核心是运动自适应样条（MAS）：每个动态 3D 高斯的均值位置 $\mu$ 被替换为一组 $N_c$ 个可学习控制点 $P = \{p_k \in \mathbb{R}^3 \mid k \in [0, N_c-1]\}$，时刻 $t$ 的位置由分段三次函数计算：

\begin{equation*}
\mu(t) = S(t, P) = (2t_r^3 - 3t_r^2 + 1)p_{t_s} + (t_r^3 - 2t_r^2 + t_r)m_{t_s} + (-2t_r^3 + 3t_r^2)p_{t_s+1} + (t_r^3 - t_r^2)m_{t_s+1}
\end{equation*}

其中 $t_r = t_s - \lfloor t_s \rfloor$ 为段内归一化时间，$t_s = t_n(N_c - 1)$，$t_n = t/(N_f-1)$ 为归一化帧时间，$m_k = (p_{k+1} - p_{k-1})/2$ 为控制点处的近似切线。注意 $N_c$ 可以不同于帧数 $N_f$，使样条以远少于逐帧存储的控制点压缩轨迹表示。控制点初始化利用 CoTracker 2D 轨迹和 UniDepth 度量深度先验：先将 2D 轨迹通过反投影函数 $\mathcal{W}_t(\varphi_t^{\text{tr}}) = \hat{R}_t^\top \pi_{\hat{K}}^{-1}(\varphi_t^{\text{tr}}, d_t(\varphi_t^{\text{tr}})) - \hat{R}_t^\top \hat{T}_t$ 提升到 3D 空间获得初始轨迹点，再以最小二乘拟合 $\min_P \sum_t \|\mathcal{W}_t(\varphi_t^{\text{tr}}) - S(t, P)\|_2^2$ 使样条逼近初始 3D 轨迹。

运动自适应控制点剪枝（MACP）确保每个高斯的控制点数量与运动复杂度匹配：每次高斯致密化后，先用最小二乘拟合出减少一个控制点的新样条 $S(t, P')$（$P'$ 包含 $N_c-1$ 个控制点），若投影误差 $E = \frac{1}{N_f}\sum_t \|\pi_{\hat{K}}(\hat{R}_t S(t,P) + \hat{T}_t) - \pi_{\hat{K}}(\hat{R}_t S(t,P') + \hat{T}_t)\|^2 < \epsilon$ 则接受剪枝并更新 $P \leftarrow P'$。简单运动的轨迹自动精简至最少控制点（如人体主体运动），复杂运动（如气球）则保留足够密度。实验显示，使用 MACP 的 SplineGS 在渲染质量上优于固定 $N_c = N_f$（因过拟合）和 $N_c = 4$（建模能力不足）的配置，同时保持更高的变形计算效率。

相机参数估计不依赖 COLMAP，而是用浅层 MLP $F_\theta$ 从位置编码的帧时间预测外参 $[\hat{R}_t \mid \hat{T}_t] = F_\theta(\gamma(t))$，焦距 $\hat{f}$ 作为跨帧共享的可学习参数。通过两类一致性约束联合优化相机参数：光度一致性损失 $L_{pc}$ 利用运动掩码 $M_{t,t_{\text{ref}}} = M_t \cdot M_{t_{\text{ref}}}$ 排除动态区域后，在静态区域对齐参考帧与目标帧颜色；几何一致性损失 $L_{gc}$ 在 3D 空间中约束反投影点对齐。运动掩码由预计算获得，静态区域为 0、动态区域为 1。整个流程采用两阶段训练：暖身阶段（1K 迭代）仅用 $L_{\text{total}}^{\text{warm}} = \lambda_{pc} L_{pc} + \lambda_{gc} L_{gc}$ 优化相机参数；主训练阶段（20K 迭代）联合优化 3D 高斯属性和相机参数，总损失 $L_{\text{total}}^{\text{main}} = \lambda_{rgb} L_{rgb} + \lambda_d L_d + \lambda_M L_M + \lambda_{pc} L_{pc} + \lambda_{\text{d-pc}} L_{\text{d-pc}} + \lambda_{gc} L_{gc}$，其中 $L_{rgb}$ 和 $L_d$ 为渲染图像和深度的 $L_1$ 损失，$L_{\text{d-pc}}$ 用渲染深度替代度量深度先验引导联合优化，$L_M$ 为 binary dice 损失——通过 $\alpha$-混合 渲染动态高斯的运动掩码 $\hat{M}_t$ 并与预计算掩码比较，促使动态-静态高斯更好地分离。MACP 每 100 次迭代执行一次控制点剪枝。

SplineGS 在 NVIDIA 动态场景数据集上以 1.1 dB PSNR 优势超越次优方法，LPIPS 在所有场景上均取得最优，同时实现比 RoDynRF 快 890 倍、比 DynNeRF 快 8000 倍的渲染速度。在 DAVIS 野外视频数据集上，COLMAP 无法为多数场景恢复可靠相机参数，SplineGS 作为无位姿方法在渲染细节上显著优于同为无位姿的 RoDynRF。值得注意的是，样条的连续轨迹建模使得 SplineGS 天然支持未观测中间时刻的新颖时空视角合成（novel view and time synthesis）：在排除训练时段的测试时间索引上，SplineGS 的时序一致性指标 tOF 远优于 D3DGS、4DGS、STGS 等方法，动态物体投影轨迹追踪也与真实运动高度吻合——而对比方法在未观测时刻产生严重退化和不可靠的运动追踪。消融实验验证了各变形模型的对比：MLP 和网格模型计算开销大（每高斯变形延迟高），三阶多项式建模能力受限，Bezier 曲线因递归计算延迟高于 MAS，而 MAS 在渲染质量和效率上同时最优。

\paragraph{自由位姿手术场景动态重建}

手术场景具有组织形变（呼吸、心跳、肠蠕动）和器械交互等独特动态特性，且表面纹理稀少、光照波动使传统 SfM 方法极易失效。此前基于 GS 的动态手术重建方法大多局限于固定相机配置，无法应对相机运动导致的场景扩展。Free-DyGS（2024.09）~\cite{2409_01003} 的设计动机正是解决移动内窥镜下的无位姿动态重建问题。

方法采用逐帧序贯优化策略，分四个阶段运行。第一阶段，场景初始化：引入可泛化高斯参数化模块（GGP），用 UNet 结构的高斯回归器网络（GRN）从首帧 RGBD 数据预测高斯属性。GRN 的关键设计是颜色校正项 $\Delta C_x$ 和深度校正项 $\Delta D_x$，将下采样特征图上像素的属性从原始值修正为 $c = SH(I_x + \Delta C_x)$、$\mu = \pi^{-1}(x, D_x + \Delta D_x)$，弥补下采样导致的信息损失。GRN（论文中也称 Sparse Gaussian Regressor, SGR）在 2000 帧手术场景数据上预训练后冻结，验证集 PSNR 分别达到 34.17 和 32.63，为后续逐帧优化提供高质量的场景初始几何。第二阶段，联合学习：对第 $i$ 帧，基于过去 $2l$ 帧的位姿估计平均速度，在李代数域预测初始位姿 $T_i^0 = T_{i-1}^k \cdot \exp(\sum_{j=1}^l (\hat{\xi}_{i-j}^k - \hat{\xi}_{i-l-j}^k) / l^2)$，然后与变形模型联合优化 $k$ 次迭代（默认 $k=10$）。联合学习的损失函数为 $\mathcal{L}_{JL} = \frac{1}{|M^i|} \sum_{p \in M^i} |I_i(p) - \hat{I}_i(p)| + \lambda_1 \Psi(t_i; \Theta)$，其中 $M^i$ 为器械掩码（排除手术器械的遮挡区域），$\Psi(t_i;\Theta) = \|\frac{1}{N}\sum_{n=1}^N \Phi_n^\mu(t_i)\|^2$ 为变形正则化项，约束所有高斯的平均位移保持较小以防止过拟合。

第三阶段，场景扩展：以不透明度阈值 $\beta = 0.8$ 定义不可见掩码 $M^{ivsb}_i$，标记相机视野中尚未被学习的区域，利用 GRN 在这些区域生成补充高斯 $G_i^{add}$ 并与已有模型合并。第四阶段，回溯学习：变形建模采用部分激活柔性变形模块（PFDM，论文中也称 Temporal Decoupling Deformation Model, TDDM），将变形函数定义为 $B$ 个高斯基函数的加权求和 $\Delta(t; \Theta) = \sum_{j=0}^B \Theta_j \cdot \phi_j(t; \tau_j)$，其中 $\phi_j(t; \tau_j) = w_j \exp(-\frac{1}{2\sigma_j^2}(t - \tau_j)^2)$；基函数的偏置 $\tau_j^{init}$ 在时间跨度上均匀分布，训练时仅激活满足 $|t - \tau_j^{init}| < t_{max} \cdot m / (B-1)$ 的子集，大幅减少梯度计算量。回溯阶段从宽度 $\gamma = 100$ 的历史时间窗口中随机采样历史帧，重新优化变形模块以平衡新旧帧的变形信息。回溯学习的损失函数为 $\mathcal{L}_{RRL} = \frac{1}{|M^{his}|} \sum_{p \in M^{his}} |I_{his}(p) - \hat{I}_{his}(p)| + \lambda_1 \Psi(t_{his}; \Theta)$，确保历史帧的变形精度不随训练推进而退化。对于长视频序列，当累计帧数超过阈值 $\kappa = 100$ 时，系统自动重新初始化一个新的变形模型实例，形成多模型表示策略以应对长时间跨度的复杂动态。训练采用 Adam 优化器，变形参数学习率为 $1.6 \times 10^{-4}$，旋转参数 $1 \times 10^{-3}$，平移参数 $0.05$，每帧联合学习 10 次迭代后回溯学习 40 次迭代。

在 StereoMIS 和 Hamlyn 数据集上的实验验证了该方法在重建质量和训练效率上的优越性：PSNR 分别达到 31.90 和 30.01，相机轨迹绝对轨迹误差（ATE）为 3.15 mm，每帧训练仅需约 0.7 秒，总训练时间分别为 11.9 分钟和 6.6 分钟，渲染速度超过 100 FPS。相比之下，Flex 需要 20 小时训练，RoDyNeRF 需要超过 24 小时，Free-DyGS 在效率上实现了数量级的提升，同时保持竞争力的重建质量。消融实验表明，GRN 初始化、场景扩展、回溯学习及变形正则化项各自对最终性能均有显著贡献。

\subsubsection{专用场景与特殊条件}

本路线介绍面向航拍、人体、越野等特定场景或特殊成像条件的定制化动态重建方法。

除通用动态场景外，一系列工作针对航拍、人体、越野等特定场景或特殊成像条件提出了定制化方案。

\paragraph{航拍单目 4D 场景重建}

AeroDGS（2026.02）~\cite{2602_22376} 针对无人机单目视频的 4D 重建提出物理引导框架。无人机成像面临窄基线、大深度范围和小型快速运动目标等多重挑战，单目深度估计和运动恢复严重病态。该方法的核心洞察是：虽然航拍图像的深度线索天然模糊，但城市场景的几何并非任意——建筑物、道路和动态物体遵循稳定的空间关系，构成低维可行配置流形。

为此，AeroDGS 设计了两阶段管线。第一阶段是单目几何提升（Monocular Geometry Lifting）：利用零样本基础模型获取粗糙先验——UniDepth 估计伪深度图 $D_t$ 和相机内参 $K$，SAM2 与跟踪模型产生跨帧实例对应关系，CoTracker3 的长期背景特征轨迹经局部 bundle adjustment 恢复尺度一致的稀疏关键点和相机位姿 $P_t$。为修正单目深度的尺度漂移，在跟踪点上计算几何深度与预测深度的局部比率场并插值得到精化深度 $\hat{D}_t$，每个像素反投影为 3D 点。对每个 2D 实例分组拟合有向包围盒，通过 PCA 估计 3D 中心 $c_{o,t}$ 和底面尺寸 $(w, l)$，高度 $h$ 则由预训练 MLP 预测（单目条件下无法从深度推断）。3D 空间中沿相机射线将物体定位到物理合理位置并强制轨迹平滑，解决 2D 跟踪的 ID 切换和遮挡问题。

在场景表示方面，AeroDGS 采用静态背景与动态前景统一的高斯基元表示。每个高斯的中心位置 $\mu_i \in \mathbb{R}^3$、协方差 $\Sigma_i$、不透明度 $o_i$ 和外观 $A_i$ 构成其属性。针对航拍影像中因阳光方向变化、表面反射和相机运动引起的大幅外观变化，外观建模为连续场 $A_i = f_\theta(\mu_i, d, t, e_o)$，由三重互补编码联合参数化：空间坐标经场景边界归一化后通过 hash grid 嵌入以捕获局部空间变化，方向依赖由球谐基表达以处理视角相关外观变化，时间信息由正弦位置编码保持时序排序并支持平滑外观演化。这种联合参数化无需为每个高斯存储独立球谐系数，在增强时序连贯性的同时降低密集高斯表示的内存开销。

动态物体的编码在规范物体空间中进行。每个运动实例 $o$ 对应一组高斯子集 $G_o$，定义在以 $(w, l, h)$ 为边界的轴对齐框内。运动建模为 SE(3) 李群上的连续 6-DoF 轨迹 $T_{o,t} = \exp(\xi_o(t))$，其中 $\xi_o(t) \in \mathbb{R}^6$ 为含旋转和平移分量的时变扭转向量，通过样条插值保证时间连续性，并附加可优化残差修正 $\Delta T_{o,t}$ 以吸收小偏差和位姿不确定性，最终变换为 $T'_{o,t} = T_{o,t} \Delta T_{o,t}$。每个高斯中心变换至世界坐标系 $\mu_{i,t} = T'_{o,t} \mu_i$。在时刻 $t$，整个场景表示为静态高斯集与各动态子集的并集 $G(t) = G_{\text{static}} \cup_{o} T'_{o,t} G_o$，通过标准 $\alpha$-合成可微光栅化渲染。

第二阶段是物理引导优化，总损失为 $L = \lambda_{\text{photo}} L_{\text{photo}} + \lambda_{\text{sup}} L_{\text{support}} + \lambda_{\text{upr}} L_{\text{upright}} + \lambda_{\text{traj}} L_{\text{traj}}$，将光度监督与三类物理先验耦合：光度项 $L_{\text{photo}}$ 结合 $L_1$ 范数和 SSIM 损失测量渲染图像与输入帧的颜色一致性，训练初期静态和动态区域等权以稳定优化，待静态背景收敛后动态区域权重递增以精化细粒度外观和运动。(1) 地面支撑约束 $L_{\text{support}} = \mathbb{E}_{o,t}[\rho(c_{o,t} - c_{o,t}^{*,g})]$，测量物体中心沿相机视线到局部地面投影的有符号距离，并将 $c_{o,t}$ 向上偏移半个物体高度以确保车辆底部贴近地面，同时容忍小的重建噪声；(2) 竖直稳定性 $L_{\text{upright}} = \mathbb{E}_{o,t}[1 - |u_{o,t} \cdot v_{o,t}|]$，将物体竖直轴 $u_{o,t}$ 对齐参考方向（刚体对齐地面法线 $n_t$，非刚体对齐重力方向 $g$），将 3-DoF 旋转约束为绕竖直轴旋转以抑制不真实倾斜；(3) 轨迹平滑性 $L_{\text{traj}} = \mathbb{E}_{o,t}[\|a_{o,t}\|^2]$，二阶平滑约束抑制高频抖动并允许物体以自然动量离开场景。训练策略方面，动态掩码权重在 7K 至 15K 迭代间从 1.0 线性递增至 1.3，同期物理引导约束衰减 50\%，总训练 30K 迭代，损失权重设为 $\lambda_{\text{photo}}=1.0$、$\lambda_{\text{sup}}=0.05$、$\lambda_{\text{upr}}=0.1$、$\lambda_{\text{traj}}=0.02$。

该方法还构建了 Aero4D 真实世界无人机数据集，涵盖不同高度（50--100\,m 飞行高度）和运动条件的城市场景，采集分辨率 2K 至 4K、帧率 15\,fps，覆盖面积 2000--18000\,m$^2$，地面采样距离 3--5\,cm/px，每帧标注相机位姿、密集点云以及动态实例的像素级掩码和 6-DoF 轨迹。在合成 UAV3D 和真实 Aero4D 数据集上的实验中，AeroDGS 在动态区域 PSNR 上超越先前方法最高 4\,dB。消融实验表明：用 COLMAP SfM 替换单目几何提升模块会因平面偏差导致动态 PSNR 显著下降；三类物理约束各自独立贡献于重建保真度和时序连贯性；移除动态掩码加权则削弱小型运动物体的重建质量。该方法目前的局限在于：3\,m 运动阈值用于动静分离可能将局部微小运动物体误判为静态，且高空视角下行人像素覆盖极低，当前不对其重建。

\paragraph{无人机人体场景动态渲染}

UAV4D（2025.06）~\cite{2506_05011} 提出面向无人机拍摄场景的动态神经渲染框架。无人机场景的独特挑战在于：硬件限制导致只能使用单目相机，动态行人尺寸小且数量多（可多达 50 人），现有依赖深度估计、光流和点跟踪的先进方法在此条件下失效。与面向手机拍摄或同步多视角设置的主流动态渲染方法不同，无人机俯视视角下场景覆盖面积大、行人像素占比极低（可小至 $10 \times 10$ 像素），导致现有大型动态场景数据集中极少包含此类场景，预训练模型在此泛化性不足。

UAV4D 的管线分为三个阶段。首先是初始化：利用 VGGT 3D 基础模型估计相机位姿 $P_t$（含内参 $K$ 和外参 $R_t, T_t$）及逐像素点图 $X_t$（含对应置信度图 $C_t$）；用 SAM2 以边界框为提示生成精确的行人实例掩码，送入 HMR2.0 预测相机坐标系下的 SMPL 人体网格参数 $(\theta_t^k, \beta_t^k)$。为保证初始质量，引入网格精化机制：将每个人体网格投影至图像空间计算包围盒 $\bar{B}_t^k$，剔除比值 $\bar{B}_t^k / B_t^k$ 超过阈值 $\eta_{\text{box}}$ 的低质量预测；对相邻帧中膨胀掩码存在重叠的缺失帧，利用前帧 $(\theta_{t-1}^k, \beta_{t-1}^k)$ 和后帧 $(\theta_{t+1}^k, \beta_{t+1}^k)$ 的四元数姿态参数和形状向量进行线性插值补全。

其次是联合人体-场景重建。背景几何重建方面，将所有点图 $\{X_t\}_{t=1}^N$ 堆叠为场景 3D 几何，利用置信度图以 $\eta_{\text{conf}}$ 百分位数过滤噪声动态区域和复杂结构，再通过 Poisson 表面重建生成干净的背景网格 $\mathcal{M}_{\text{bg}}$——过滤后的背景网格为平坦的地面表面，避免了噪声人体区域和复杂结构对接触点检测的干扰。核心难题是 3D 基础模型的尺度歧义——背景网格与 SMPL 人体网格处于不一致的尺度空间。由于 SMPL 人体网格定义在度量尺度（metric scale）下，UAV4D 可以相对于人体网格对齐背景网格。具体而言，引入全局尺度参数 $\sigma$ 调整点图、背景网格和相机位姿，从骨骼长度角度优化 $\sigma$：通过 ViTPose 检测 2D 关键点，利用缩放后的点图将其提升为 3D 关节 $\hat{J}_t^k = \sigma X_t(i,j)$，计算损失 $L(\sigma) = \sum_t \sum_k \sum_{(p,c) \in \mathcal{J}_{\text{body}}} \|\hat{d}^k_{(p,c),t} - d^k_{(p,c),t}(\sigma)\|_2$，其中 $\hat{d}$ 来自 SMPL 网格的骨骼长度，$d(\sigma)$ 来自缩放点图，通过 L-BFGS 算法（更新规则 $\sigma_{n+1} = \sigma_n - \alpha \mathcal{H}_n \nabla L(\sigma_n)$，其中 $\mathcal{H}_n$ 为逆 Hessian 近似）优化 $\sigma$。确定最优 $\sigma$ 后，将其应用于背景网格 $\mathcal{M}_{\text{bg}}$ 和相机位姿 $P_t$ 以完成尺度统一。然后利用无人机俯视视角的优势，从人体包围盒底部检测 2D 接触点 $(x_c, y_c) = ((x_{\min}+x_{\max})/2, y_{\max})$，从背景网格渲染深度图 $D_t$，通过反投影公式 $\psi_t = \sigma R^T[K^{-1}D_{x_c,y_c}[x_c,y_c,1]^T] - \sigma R^T T$ 将 2D 接触点提升为世界坐标系下的 3D 地面接触点 $\psi_t$，完成人体在世界坐标系中的定位。

最后是 3D 高斯场景建模：背景用静态高斯 $g_{\text{bg}}$ 表示——高空场景直接用全部多点图初始化，低空场景则利用置信度图过滤并用人体掩码剔除人体区域的点图后初始化；人体用 SMPL 规范网格顶点初始化的高斯 $g_h$ 表示（受 GART 启发），通过 LBS 蒙皮变换驱动变形——$g_h(t) = (R_A(t) \cdot \mu + T_A(t) + \psi_t, R_A(t) \cdot r, s, c, o)$，其中 $R_A(t), T_A(t)$ 来自 LBS 权重与 SMPL 姿态参数的蒙皮变换 $LBS(g_h, \theta_t)$，$\psi_t$ 为 3D 地面接触点偏移。$\theta_t$ 的时序变化驱动关键关节变换更新，通过线性插值获得不同时间步的变形位置 $\mu(t)$ 和旋转 $r(t)$。两者合并为统一图元 $g_{\text{all}} = g_{\text{bg}} + g_h$ 送入光栅化器渲染，总损失为 $L_{\text{total}} = (1-\lambda_{\rho})L_1 + \lambda_{\rho} L_{\text{ssim}} + \lambda_o L_o + L_{\text{smpl}}$，其中 $L_1$ 和 $L_{\text{ssim}}$ 为渲染全图与真值的光度损失，$L_o$ 为不透明度正则化，$L_{\text{smpl}}$ 包含人体高斯和人体区域相关损失项。优化过程中联合更新所有高斯属性 $g_{\text{bg}}, g_h$、所有帧的 SMPL 参数 $\theta_t, \psi_t$ 及对应的蒙皮权重。

实验在三个各具特色的无人机数据集上评估：VisDrone（中国，40--50m 高空，密集行人 30--50 人）、Okutama-Action（日本，10--45m，显著相机变化和多视角）和 Manipal-UAV（印度，10--50m，最高海拔场景）。以每 8 帧为测试集进行新视角合成评估。全图指标上，UAV4D 在三个数据集上均达到最优或接近最优，PSNR 超越先前方法最高 1.5 dB（VisDrone 上 26.03 dB 对比 3DGS 的 24.59 dB），SSIM 和 LPIPS 全面领先。人体区域单独评估同样全面最优。值得注意的是，4DGS 和 Deformable-GS 在 VisDrone 和 Okutama-Action 上训练不稳定——场景中多个小运动物体的存在表明单一变形场无法准确捕获所需变形，导致训练过程高度不稳定；TK-Planes 作为唯一专为无人机动态场景设计的方法，在各动态场景上更为鲁棒但渲染模糊。在 Manipal-UAV 上，Deformable-GS 在密集植被和摇曳树叶区域表现更好，暴露了 UAV4D 的一个局限：静态背景高斯无法表示摇曳树木等弥漫性非刚性变形。

消融实验验证了各组件的必要性（VisDrone 上）：移除人体 SMPL 高斯（仅保留背景，相当于将人体冻结为静态或忽略运动行人）导致 PSNR 下降 1.7 dB 至 24.33 dB；移除尺度优化（固定 $\sigma = 40$）导致人体网格与场景错位，PSNR 暴跌 3.01 dB 至 23.02 dB；移除网格精化则使 PSNR 下降 0.24 dB 至 25.79 dB。尺度优化是贡献最大的组件。此外，人体-背景分解天然支持场景编辑：可以独立移除、平移或修改场景中的人体，为无人机下游应用（如人员检测和行为识别的合成数据生成）提供灵活工具。UAV4D 的局限性包括：(1) 依赖预训练模型（HMR2.0 的网格参数仍存在误差，导致渲染伪影）；(2) 自行车、摇曳树木等非人体动态成分无法由静态背景高斯有效渲染。

\paragraph{越野场景前馈 4D 重建}

Ground4D（2026.05）~\cite{2605_04435} 针对非结构化越野场景提出空间锚定的 4D 前馈重建框架。越野场景呈现三个互相放大特性的挑战：密集植被和松散地形带来的高频几何、摇曳树冠和飘散颗粒形成的空间弥散非刚体动力学、崎岖地面引起的连续自运动抖动。现有前馈高斯 Splatting 方法在此条件下的根本失败在于：规范高斯空间中同位置的高斯之间发生属性冲突——空间偏移冲突来自自运动抖动散射的表面观测，时间变异冲突来自弥散动力学产生的不一致变形状态。朴素融合（如 NoPoS）导致属性模糊，逐高斯时序建模（如 DGGT）则导致空间占用的结构空洞。两者共同的根源在于时间相关性在空间上下文之外独立求解，迫使选择性和空间占用陷入固有矛盾。

Ground4D 的核心洞察是：将时间竞争限制在局部空间邻域内即可消除这一矛盾。在局部邻域内，时间最主导的高斯同时达到最高时间相关性和保证空间占用，两个目标相互强化。具体实现为体素锚定的时序高斯聚合：首先通过可微空间体素化将规范空间 $G^s$ 按体素边长 $\rho$ 划分为空间体素，每个高斯中心 $\mu_i$ 通过 round 函数量化为整数坐标 $v_i = \text{round}(\mu_i / \rho)$。然后在每个体素 $m$ 内，查询条件时序注意力机制对共位置高斯评分——每个基元携带特征 $f_i^{\text{gs}}$（高斯属性编码与高斯头残差特征拼接），经编码器 $\phi_f$ 投影至共享隐空间，时间戳 $\tau_i$ 经多频率正弦编码和时间 MLP 编码后加性融合为隐空间表示 $h_i = \phi_f(f_i^{\text{gs}}) + \phi_t(\gamma(\tau_i))$，查询时间 $\tau^*$ 类似编码为 $h^* = \phi_t(\gamma(\tau^*))$，注意力 MLP 产生标量对数 $a_i = \text{MLP}_{\text{attn}}([h_i; h^*])$。体素内温度缩放 softmax 归一化 $w_i = \text{softmax}(a_i / \beta)_{j \in \mathcal{V}_m}$ 确保每个非空体素都产生有效图元。

属性融合采用语义感知估计器：位置和颜色用加权算术平均；不透明度将加权均值与体素内最大值按混合系数 $\lambda$ 混合以避免瞬态遮挡体素的密度坍缩；尺度在对数空间聚合实现加权几何平均；旋转在单位四球面上加权平均后 $L_2$ 重归一化。显式动态物体（车辆、行人）由动态置信头 $\Pi_t^{\text{dyn}}$ 检测并用 TAPIP3D 插值运动。弥漫性高频非刚性变形则由体素锚定时序注意力隐式处理。天空背景由参数化模型 $\mathcal{M}^{\text{sky}}$ 以视角方向为条件建模，最终合成为 $\hat{I}_t = \mathbf{A}_t \odot \hat{I}_t^{3D} + (1-\mathbf{A}_t) \odot \mathcal{M}^{\text{sky}}(\mathcal{I}, E_t, K_t)$。表面法线通过双路径设计提供几何先验：法线头中间特征注入高斯头约束预测，高斯衍生法线 $\tilde{\mathbf{n}} = R_\mathbf{q} \cdot \text{softmax}(-\eta \mathbf{s}) / \|\text{softmax}(-\eta \mathbf{s})\|_2$ 通过 soft-argmin 从尺度轴估计并由余弦距离损失监督对齐真值。训练端到端监督，总损失为六项之和：$L_1$ 光度损失、LPIPS 感知损失、动态置信 BCE 损失、天空 $\alpha$ 监督损失，以及预测法线和高斯衍生法线两项余弦距离损失。训练时引入随机帧丢弃（概率 $p_{\text{drop}} = 0.7$）迫使网络从时间不连续的上下文推断场景状态。关键超参数包括体素尺寸 $\rho = 0.002$、不透明度混合系数 $\lambda = 0.3$、可学习温度 $\beta$ 初始化为 1.0。

实验结果表明，体素锚定聚合的三个组件（体素化、时序注意力、体素内归一化）缺一不可，三者结合将 spatio-temporal 交互从竞争变为协调。体素尺寸需平衡竞争强度与空间连贯性：最优值 $\rho = 0.002$ 将高斯从约 600K 压缩至 370K（1.6 倍减少），过大体素（$\rho = 0.005$）过度合并空间不同表面导致 2.85 dB PSNR 暴跌。上下文帧数在 10 帧达到峰值后饱和。该方法在 ORAD-3D 上取得最高 1.48 dB PSNR 增益（23.89 dB），并零样本泛化至 RELLIS-3D——无需微调即在 SSIM 和 LPIPS 上达到所有方法最优，说明体素内时序锚定学到了可迁移的场景先验：空间局部化聚合保持的几何连贯性跨越野域泛化。此外，在 nuScenes 城市驾驶基准上微调后达到 27.23 dB PSNR 和 0.814 SSIM，确认该方法可扩展至结构化城市场景。

\paragraph{稀疏观测下的动态重建}

Guess The Unseen（2024.04）~\cite{2404_14410} 提出 3DGS 表示同时重建静态世界和多个动态人体的方法，重点解决野外视频中人体被部分遮挡、裁剪、稀疏观测的极端挑战。整个场景统一用 3DGS 表示——背景 $G^{\text{BG}}$ 由 SfM 稀疏点云初始化（当固定相机输入下 SfM 不适用时，假设背景为以人体平均位置为中心的大型 3D 球体并贴附背景纹理），每个人体 $G^{h_j}$ 在 SMPL 规范空间（A-pose）中用 SMPL 网格顶点 $\mathcal{V}(\beta_j, \theta_c)$ 初始化（颜色灰度 0.5、不透明度 0.9），其中 SMPL 形状参数 $\beta_j$ 由单目 3D 姿态回归器估计并在训练期间固定不变。场景所有高斯可组合为 $\mathcal{G}^{\text{All}} = \{\mathcal{G}^{BG}\} \cup \{\mathcal{G}^{h_j}\}_{j=1}^{N}$，使用与原始 3DGS 完全相同的 $\alpha$-混合渲染公式，无需任何修改即可实现多人加背景的联合渲染——这是相比 NeRF 表示的重要优势，后者组合多人渲染需要复杂的 ST-NeRF 管线。高斯致密化基于累积位置梯度 $\sum_{N_{\text{den}}}\nabla\mu_i$ 超过阈值时执行克隆或分裂操作。通过 LBS 蒙皮变形将规范空间高斯映射到姿态空间：每个高斯中心 $\mu_{i,p}$ 和协方差 $\Sigma_{i,p}$ 的变形公式为 $\mu_{i,p} = \sum_{k=1}^{N_{\text{joint}}} w_k(\mu_{i,c})(R_k \mu_{i,c} + T_k)$，$\Sigma_{i,p} = R_{\text{wei}} \Sigma_{i,c} R_{\text{wei}}^T$，其中 $R_{\text{wei}} = \sum_k w_k R_k$ 为蒙皮旋转权重，$R_k$、$T_k$ 为第 $k$ 个 SMPL 关节的旋转和平移。LBS 权重通过对最近 30 个 SMPL 顶点的逆距离加权（IDW）计算并预存于体素网格中，渲染时通过三线性插值查询而非逐次搜索最近顶点以加速计算。

为从稀疏观测中恢复完整 3D 人体，该方法引入扩散引导重建作为核心创新。直觉上，3D 人体模型的质量可以通过评估新视角渲染图像的真实感来衡量——直接从稀疏遮挡视角优化 3DGS 会产生类似"刺猬"的透明伪影。对每次迭代，随机采样虚拟相机 $v$（以人体为中心的球面上）和姿态 $\theta \in \{\theta_t\}_{t \in [1 \ldots T]} \cup \{\theta_c\}$（从已观测姿态或规范 A-pose 中选取），渲染人体高斯得到 $R_v^h$，计算 SDS 损失作为扩散模型引导的虚拟监督信号。为使生成的外观与特定个体一致而非随机人体，利用 CustomDiffusion 进行文本反演（Textual Inversion）：从输入帧裁剪目标人体图像，反演为文本 token <person-j> 并微调扩散模型的注意力和文本嵌入层，得到人物专属扩散模型 $\Phi_j$。同时引入 OpenPose ControlNet 约束：将 3D SMPL 关节 $J_{\text{smpl}}(\theta, \beta_j)$ 投影到虚拟相机视角并转换为 OpenPose 格式输入 ControlNet，保证扩散生成的人体姿态与高斯模型一致。

为增强人体细部纹理，还随机采样聚焦面部、上半身和下半身的局部相机 $v_j$。噪声时间步从 $\mathcal{U}[0.2, 0.98]$ 采样（前 2000 次迭代），随后最大时间步从 0.98 衰减至 0.3（再经 2000 次迭代），逐步从粗到精地优化。此外加入视角增强语言提示以稳定优化。总训练目标为 $L_{\text{tot}} = \lambda_{\text{recon}} L_{\text{recon}} + \lambda_{\text{sds}} L_{\text{sds}} + \lambda_{\text{hard}} L_{\text{hard}}$，其中重建损失 $L_{\text{recon}} = \lambda_{\text{rgb}}\text{MSE}(R_t, I_t) + \lambda_{\text{ssim}}\text{SSIM}(R_t, I_t) + \lambda_{\text{lpips}}\text{LPIPS}(R_t, I_t)$，SDS 损失对每个人体独立应用 $L_{\text{sds}} = \sum_{j=1}^{N} L_{\text{SDS}}(\mathcal{G}^{h_j}, \Phi_j)$，硬表面正则化 $L_{\text{hard}} = -\log(\exp^{-|\alpha|} + \exp^{|\alpha|}) + \text{const}$（$\alpha$ 为人体高斯的渲染 alpha 图）参考 LOLNeRF 设计以消除透明伪影。训练策略采用三阶段：（1）先用人体掩码后的纯背景图像单独训练背景高斯；（2）训练人体高斯前 1000 次迭代不使用 SDS；（3）背景和人体高斯同时优化并启用 SDS 共计 10k 次迭代。

实验结果验证了该方法在极端稀疏场景下的有效性。在 Panoptic 多人数据集（6 人 ultimatum 序列，仅用单视角输入模拟稀疏观测）上，该方法达到 23.60 dB PSNR，大幅超越 HumanNeRF（19.59 dB）和 Shuai et al.（15.79 dB）。在 Hi4D 近距离双人交互数据集上，pair00-dance 序列 PSNR 23.76 dB vs HumanNeRF 18.79 dB，pair01-hug 序列 PSNR 25.14 dB vs HumanNeRF 21.40 dB——在从未观测到面部的极端情况下，扩散模型仍能"幻觉"出合理的人脸外观。在 ZJU-MoCap 单人充足观测场景下，该方法与 HumanNeRF 质量相当（30.13 vs 30.23 dB PSNR），表明 3DGS + SDS 管线不会损害充分观测场景的重建质量。

消融实验（ZJU-MoCap 377 主体，下身完全遮挡+仅使用 10\% 帧）显示：完整方法 PSNR 26.03 dB，去除 SDS 降至 23.43 dB（下身无法重建且上半身质量也因帧不足而下降），去除文本反演降至 24.51 dB（能重建完整形体但外观与目标个体不一致），证实 SDS 损失与文本反演各自不可替代。渲染效率方面，15K 高斯在 512$\times$512 和 1024$\times$1024 分辨率下分别达到 361 FPS 和 277 FPS，VRAM 仅 1496MiB，速度和内存均显著优于 InstantAvatar（18.03 FPS、4972MiB）。重建后支持任意视角渲染、场景编辑（移除单个人体或替换运动）和人体分离，每个个体的姿态可独立控制，兼容任意运动捕捉数据驱动的动画重定向。

\paragraph{人体动态合成与运动生成}

HumanGenesis（2025.08）~\cite{2508_09858} 提出集成几何和生成建模的模块化框架，通过四个协作智能体解决人体视频生成中的几何不一致、粗糙重建、运动泛化差和场景不协调四个核心问题。该方法的核心洞察是：纯 2D 姿态引导的生成方法缺乏对底层 3D 几何的理解，导致跨帧几何不一致；而纯 3D 重建方法虽然几何先验更强，但难以泛化到新姿态和新场景。HumanGenesis 通过将几何建模与生成建模统一到同一框架中，使两类方法的优势互补。

重建器（Reconstructor）从单目视频分离人体和场景：用 COLMAP 获取 SfM 点云初始化场景高斯 $S_G$，SMPL 规范网格（Da-pose）顶点采样 50K 点初始化人体高斯 $H_G$。变形建模将人体变形分解为非刚性变形和刚性变形两步：首先，特征编码器 $\bm{F} \in \mathbb{R}^{3 \times h \times w \times d}$ 为规范空间中任意点 $\bm{x}_c$ 生成特征嵌入 $f^c \in \mathbb{R}^{64}$，非刚性变形解码器 $\mathcal{D}_{nr}$（3 层 MLP，128 维）预测位置、尺度和旋转偏移量 $(\Delta x, \Delta s, \Delta \bm{q})$，得到变形后的规范高斯——位置 $\bm{x}_d = \bm{x}_c + \Delta x$、尺度 $\bm{s}_d = \bm{s}_c \cdot \exp(\Delta s)$、旋转 $\bm{q}_d = \bm{q}_c \cdot [1, \Delta q_1, \Delta q_2, \Delta q_3]$。随后，刚性变形通过 LBS 蒙皮实现——给定姿态参数 $\bm{\theta}$，每个高斯点 $\bm{x}_d$ 通过 $\bm{x}_p = \sum_{i=1}^{n_k} w^i \bm{B}_i \bm{x}_d$ 刚体变换到姿态空间，其中 $\bm{B}_i$ 为第 $i$ 个关节的骨骼变换矩阵。LBS 权重以 SMPL 权重初始化，并通过 4 层 MLP 解码器 $\mathcal{D}_r$（128 维）预测偏移量以降低计算开销。优化完成后将 LBS 权重与所有 3D 高斯属性一并存储，运行时无需再评估三平面和解码器，整个预处理和训练管线仅需单 GPU 约 20 分钟。重建损失为 $L_{\text{recon}} = L_1 + \lambda_1 L_{\text{mask}} + \lambda_2 L_{\text{SSIM}} + \lambda_3 L_{\text{LPIPS}}$（$\lambda_1=0.5$, $\lambda_2=\lambda_3=0.01$），训练 10k 迭代约 10 分钟。

姿态引导器（Pose Guider）不仅将 SMPL 姿态参数序列编码为时间感知嵌入以驱动人体高斯变形与运动对齐，还直接在像素空间捕捉人体的视觉外观。具体地，给定参考帧和由预训练 SMPL 估计器提取的姿态参数序列，变形后的人体高斯 $H_{\mathcal{G}}'$ 通过 SMPL 模型映射——$\mathrm{SMPL}({H_{\mathcal{G}}'}^{i}) = (\beta_{H_{\mathcal{G}}}, \bm{\theta}_m^i)$，其中 $\beta_{H_{\mathcal{G}}}$ 为体型参数。与 Champ 等方法使用多种 2D SMPL 投影表示不同，Pose Guider 直接捕获像素空间中的人体视觉外观，并通过紧凑的时间感知编码器为下游 Video Harmonizer 提供条件输入。基于 3DGS 的显式表示，HumanGenesis 还支持场景融合（Scene Fusion）：将人体高斯 $H_{\mathcal{G}}$ 叠加到不同场景高斯 $S_{\mathcal{G}}$ 上（配合场景平移 $T_s$）后进行体渲染，即可实现将同一人体嵌入多种背景场景的视频生成。

批判智能体（Critique Agent）实现空间精细重建：首轮从训练视角渲染图像 $\mathcal{I}^{(i)}$，利用多模态 LLM（Qwen-VL）识别模糊、肢体交叉、脚地融合等常见 3DGS 重建缺陷，输出边界框 $b_j^{(i)}$ 和文本反馈 $f_j^{(i)}$——$(b_j^{(i)}, f_j^{(i)}) = \mathrm{MLLM}(\mathcal{I}^{(i)})$。对负面反馈区域启动针对性重训练，迭代多轮直到所有区域达到满意标准或超过最大轮数（实验中第 3 轮后性能趋于稳定）。这一机制模拟专家评估过程，弥补全局图像监督无法关注局部缺陷的问题。在区域感知训练中，被 MLLM 标记为高重要性的局部区域（如面部、关节）被赋予显著放大的损失权重——$L_{\text{final}} = L_{\text{recon}} + \omega \cdot \sum_{(i,j) \in \mathcal{R}_{\text{local}}}[L_{\text{color}}^{(i,j)} + \lambda_1 L_{\text{mask}}^{(i,j)} + \lambda_2 L_{\text{SSIM}}^{(i,j)} + \lambda_3 L_{\text{LPIPS}}^{(i,j)}]$，其中 $\omega \gg 1$，使模型在感知关键区域获得更精细的重建性能。

视频协调器（Video Harmonizer）通过混合渲染管线消除场景不协调：去噪 U-Net 以 Stable Diffusion V2.1 初始化并将 2D 卷积扩展为 3D 以捕获时间动态（在 8 块 A100-80G 上训练，batch size 48），引入时间注意力机制和 3D 卷积块建模跨帧运动。Pose Guider 编码的 3D 渲染条件通过空间特征变换（SFT）注入去噪网络。运动引导扩散采样策略利用光流对齐跨帧特征保证时间一致性。VAE-Decoder 微调引入时间 3D 残差块，混合损失为 $L_{\text{total}} = L_1 + L_{\text{LPIPS}} + \alpha L_{\text{diff}} + \beta L_{\text{GAN}}$（$\alpha=0.5$, $\beta=0.025$），其中帧差损失 $L_{\text{diff}} = \sum \|(\hat{I}_{i+1} - \hat{I}_i) - (I_{i+1}^{gt} - I_i^{gt})\|_1$ 约束连续帧间的动态一致性，并通过可控特征变形（CFW）模块从 VAE 编码器注入信息。为保持预训练 VAE 的完整性，仅更新时间卷积和 CFW 模块的参数，冻结编码器的空间层。推理时采用渐进式分块聚合采样策略，50 步采样，支持任意分辨率输出。

Back-to-4D 迭代反馈循环将几何重建和生成质量相互促进：从初始人体-场景 3DGS 出发，每轮外迭代（共 $E$ 轮）依次执行三步——(1) 基于当前 3DGS 以新相机姿态/轨迹生成合成视角 $\mathcal{V}_t$；(2) 用训练好的 Video Harmonizer 去噪增强得到 $\hat{\mathcal{V}}_t$；(3) 以增强视角作为监督进行 $T=2.5k$ 步内层优化，更新高斯的位置、颜色、不透明度和协方差属性。与 Instruct-NeRF2NeRF 的蒸馏过程不同，该方法确保输出保持人体-场景 3DGS 的一致性和时间连贯性。

在实验评估中，HumanGenesis 在 HumanVid 数据集上与 AnimateAnyone、AniGS、MIMO、Champ 等方法对比，在 PSNR、SSIM 上取得更高值，同时显著降低 LPIPS 和 FID-VID 分数。消融实验验证了各组件的有效性：Critique Agent 的多轮反思在第 3 轮达到最优后趋于稳定；Video Harmonizer 的时间卷积（TC）、时间微调（TFT）和迭代扩散增强重建（IDR）均带来一致的性能提升。训练数据包含 MoVi 数据集（225 个室内视频）和从 HumanVid 策划的 1K 单人视频，覆盖不同年龄和拍摄视角。

\paragraph{周期性环境运动建模}

Modeling Ambient Scene Dynamics（2024.06）~\cite{2406_09395} 针对具有周期性环境运动（如摇曳的树叶、水波纹、蜡烛火焰、风铃摆动、飘动音符）的场景提出单目自由视角合成方法。该方法的核心动机是：自然界中的环境运动具有天然的周期性，可以被周期函数有效建模，从而实现运动外推至未观测时段。论文同时贡献了 Forest 数据集，包含 10 个不同环境下的动态场景视频（10--30 秒，1080p/30fps），涵盖多种植物和自然景观。

方法采用三阶段管线。预处理阶段用 COLMAP 获取相机位姿和稀疏点云初始化高斯。作者发现 3DGS 的重建质量严重依赖初始点云覆盖率——COLMAP 在覆盖不足区域无法提取特征点，导致缺失高斯和错误几何。为此引入深度正则化策略：用 Instant-NGP 拟合场景（仅 10 分钟，10k 迭代），渲染深度图 $D$，然后通过高斯 Splatting 渲染深度 $\hat{D}$ 并施加 $L_2$ 深度正则化损失将高斯拉向正确几何。静态重建阶段（80k 迭代）前 9k 次使用分裂和克隆操作增强高斯密度，之后切换为掩码剪枝策略有效减少高斯数量。

动态重建阶段将每个高斯的规范空间位置 $\mu$ 和旋转 $q$ 扩展为时间相关量 $\mu_t = \mu + \Delta\mu_t$，$q_t = (q + \Delta q_t)/\|q + \Delta q_t\|$。关键创新是不直接预测变形量，而是先建模运动轨迹再从中采样时间步。轨迹表示采用离散余弦变换（DCT）基函数：$v(t) = \sqrt{\frac{2}{K+1}}\sum_{k=1}^{K} \phi_{v,k} \cos(\frac{\pi}{2T}(2t+1)k)$，其中 $\phi_{v,k}$ 为 MLP 预测的第 $k$ 个系数，$v(t)$ 可以是平移 $\Delta\mu_t$ 的任一坐标轴分量。该公式灵感来自 DCT-NeRF，利用体素编码器（基于 HexPlane 架构）将高斯规范位置编码后经 MLP 预测 DCT 系数。实验设 $K = \lceil T/4 \rceil$，将变形参数量压缩至四分之一。DCT 的自然特性使得在有观测监督时拟合运动，在无直接监督时（如高斯被遮挡或超出相机视锥）自然复制邻近运动模式，实现运动外推。DCT 轨迹方法受启发于并行研究（2023）~\cite{katsumata2023efficient}中关于基函数表示变形可提升存储效率的发现。此外，作者发现直接逐高斯优化 DCT 系数（不经过 MLP）会导致运动闪烁，因此 MLP 路径对平滑轨迹预测不可或缺。

场景归一化采用基于相机位置范围而非点云范围的策略，保证前景的高分辨率运动建模同时对远处背景施加运动正则化。作者实验了 Mip-NeRF 360 的 $\infty$-范数场景收缩策略，但发现背景仍因运动预测不准而模糊。基于相机位置范围的归一化将范围外的远距离高斯映射到边界以降低分辨率，既保留前景运动细节，又对距相机较远（不太可能有显著运动）的区域施加正则化。

为解决相邻高斯空间不一致变形导致的不自然运动，引入刚性正则化：每个高斯 $G$ 的邻居应遵循 $G$ 坐标系的刚体变换，且邻居间保持相似的旋转变化以维持时空连贯性。GPU 内存方面，设计两阶段渲染：第一遍不缓存特征获取变形参数并识别参与渲染的高斯；第二遍仅处理这些高斯并分块（每块最多 500k 样本）缓存特征进行梯度累积，解决大范围无界场景下的显存溢出问题。动态重建阶段共 30k 迭代联合优化 MLP 和 3D 高斯属性，在单块 A100-40G 上训练，总训练时间 1.5--8 小时（取决于场景复杂度）。

在实验评估中，与 RoDyNeRF 和 4D-GS 对比。RoDyNeRF 受限于隐式表示和光流约束无法产生清晰图像；4D-GS 能较好捕获静态场景但在远处区域引入明显伪影，且无法准确表示环境运动。用户研究（24 名用户，14 组对比）表明该方法在几何外观质量和运动重建质量上均大幅优于基线。定量评估在每场景保留 3 段各 30 帧的测试集上进行，PSNR/SSIM/LPIPS 全面超越基线。消融实验验证了五个组件的贡献：深度正则化（修复覆盖不足区域的错误几何）、DCT 轨迹（对比直接预测 $\Delta x_t$ 的 4D-GS 方案，运动更平滑）、轨迹 MLP（消除逐高斯独立优化时的运动闪烁）、刚性正则化（保证相邻高斯运动一致）和场景归一化（保证前背景质量均衡）。

Compensating Spatiotemporally Inconsistent Observations（2025.05）~\cite{2505_01235} 从另一角度解决在线动态 3D 高斯泼溅中的时空不一致观测问题。该方法的核心洞察是：真实世界录制中的传感器噪声等误差使观测偏离理想信号，而在线重建因逐帧访问限制会过拟合这些时变误差，导致静态区域出现闪烁伪影。方法将观测建模为理想信号与误差之和 $\tilde{I}_t = I_t + M_t = \hat{I}_t + \hat{M}_t$，引入可学习残差图 $\hat{M}_t^v \in \mathbb{R}^{3 \times H \times W}$ 对每个相机视角建模观测误差，与 3D 高斯联合优化。优化时用 $L_1$ 正则化约束残差图以防止过拟合视角相关颜色。训练流程为：首帧初始化高斯和零值残差图，残差图冻结至高斯开始致密化后再参与优化；后续帧将前帧的高斯（包括新增高斯）传播至当前帧。该方法可即插即用地应用于 Dynamic3DG、3DGStream、HiCoM 等在线重建基线，在 Neural 3D Video 和 MeetRoom 数据集上一致提升 PSNR、SSIM 并降低静态区域掩码总变差（mTV），同时减少高斯数量（初始高斯降至 0.69 倍，新增高斯降至 0.28 倍），缩短训练时间。

\subsubsection{高效表示与运动建模}

本路线介绍SE(3) B样条连续运动表示、高级初始化策略、层次化诱导流引导等高效表示与运动建模方法。

以下介绍 6 个高效表示与运动建模方法，涵盖 SE(3) B-spline 连续运动表示、高级初始化策略、层次化诱导流引导、增强先验重建、局部隐式特征解耦以及各向同性约束等方向。

\paragraph{SE(3) B 样条连续运动表示}

Learning Explicit Continuous Motion Representation（2026.03）~\cite{2603_25058} 提出用 SE(3) B 样条运动基显式建模动态高斯的连续位置与方向变形轨迹。

\textbf{Baseline 与问题。}现有单目动态高斯方法分为两类：一类（如 Deformable-GS（2023.09）~\cite{2309_13101}、SoM（2024.07）~\cite{2407_13764}）用 MLP 学习从规范空间到观测空间的仿射变换，不保证高斯轨迹在时间上连续；另一类（如 SplineGS（2024.12）~\cite{2412_09982}、NodeGS（2025.10）~\cite{2510_02732}）用 Hermite 样条建模连续位置轨迹，但无法保证方向的连续变化——这导致动态高斯的姿态出现非光滑跳变，在复杂运动区域产生明显伪影。HiMoR（2025.04）~\cite{2504_06210} 虽然通过层次化运动树结构将运动分解为粗细不同层级，但每个层级仍然使用仿射变换建模变形，未能保证逐层变形在时间维度上的连续性。

\textbf{核心改动。}作者用数学上保证连续的 SE(3) 累积 B 样条（cumulative B-spline）统一建模位置和方向的连续变形，并设计自适应控制机制动态增减运动基与控制点。SE(3) B 样条运动基由可学习控制点构造。初始化时，遵循 MoSca 的预处理流程获取 3D 轨迹片段的位姿状态 $Q = [R, \mathbf{t}]$（$\mathbf{t}$ 为空间位置，$R$ 为方向），以此作为控制点。相邻控制点之间的相对位姿变换通过对数映射转换到李代数空间：$\boldsymbol{\xi}_i = \log(\Delta Q_i)$，其中 $\boldsymbol{\xi}_i$ 表示相邻轨迹片段之间位姿变换的``速度''。运动基的时间插值定义为：

\begin{equation}
    T(t) = \Big(\prod_{i=0}^{N_c-1} \exp(\beta_i(t)\,\boldsymbol{\xi}_i)\Big)\,T_0,
\end{equation}

其中 $\beta_i(t)$ 为 B 样条基函数，$N_c$ 为控制点数量，$T_0$ 为首帧轨迹片段的位姿状态。该公式将所有变形统一在 SE(3) 流形上，保证位置和方向同时连续变化——直观地说，就像用一条数学上光滑的曲线同时描述物体``走到哪里''和``朝向哪里''，而不是在离散帧之间跳变。

\textbf{场景初始化与动态高斯变形。}为避免静态区域的变形干扰和不必要的计算开销，场景同时构建静态和动态高斯：静态高斯通过跨帧背景深度反投影初始化，与标准 3DGS 参数相同；动态高斯通过跨帧 2D 轨迹深度反投影初始化，额外携带参考时间参数 $t_{\text{ref}} \in \mathbb{R}$（即初始化该高斯的深度图时间戳）。动态高斯跟随其最近 SE(3) B 样条运动基的变形：先用最近邻算法找到高斯在参考时间 $t_{\text{ref}}$ 处最近的运动基 $T(t)^g$，再用 K 近邻找到所有时刻 $T(t)^g$ 周围的 $K$ 个运动基，最后通过 Dual Quaternion Blending（DQB）融合参考空间到观测空间的相对位姿变换 $\Delta Q$，得到高斯级变形 $\Delta Q^g = \text{DQB}(\{(w_i, \Delta Q^i)\}_{i=1}^K)$，其中权重 $w_i$ 遵循 MoSca（2024.05）~\cite{2405_17421} 的加权策略。变形后的位置和方向为 $\mu'_g = \Delta Q^g \mu_g$，$R'_g = \Delta Q^g R_g$。

\textbf{自适应控制机制。}过多控制点既增加计算量又易导致过拟合，且场景不同区域运动复杂度差异大。因此引入两类操作：（1）\emph{剪枝}——每 $N_{\text{prune}}=500$ 次迭代尝试移除一个控制点 $Q$，每次选择使轨迹误差最小化的最优剪枝点 $\hat{Q} = \arg\min_Q \sum_{t=0}^{N_T}\|T(t)^{Q} - T(t)\|_2^2$，若移除后误差 $E = \sum_{t=0}^{N_T}\|T(t)^{\hat{Q}} - T(t)\|_2^2$ 小于阈值 $\tau_{\text{prune}}=5.0$，则确认移除；（2）\emph{加密}——每 $N_{\text{densify}}=500$ 次迭代在复杂运动区域增加运动基密度。具体做法是：先用渲染误差 $m_{\text{error}}^i = \text{True}$（当 $|\hat{I}^i - I^i| > \epsilon_{\text{error}}=0.5$ 时）与动态区域掩码 $m_d^i$ 取交集 $m^i = m_{\text{error}}^i \cap m_d^i$ 定位复杂运动区域，构成掩码序列 $M=[m^0, m^1, \ldots, m^{N_T-1}]$；再将运动基在各时刻的 3D 位置投影到图像平面，若超过 50\% 的投影落入掩码 $M$ 内，则对该运动基进行加密——复制控制点参数并添加随机扰动。此外，\emph{软片段重建策略}动态调整高斯不透明度 $o' = \text{sigmoid}(\text{scale} \cdot (1 - |t_{\text{ref}} - t_{\text{obs}}|)) \cdot o$（scale 设为 5.0），降低参考时刻与观测时刻间隔较大的动态高斯贡献，缓解长间隔变形引起的干扰——即``时间上隔得远的高斯可信度打折''。

\textbf{扩散先验与正则化。}由于单目视频缺乏多视角线索，方法用多视角扩散模型（Zero-123-xl）提供 SDS 损失补充不可见区域的几何先验：

\begin{equation}
    \mathcal{L}_{\text{sds}} = \mathbb{E}_{t,\boldsymbol{\epsilon}}\big[\|\omega(t)(\boldsymbol{\epsilon}_{\phi}(z_t, t, P_t P_s^{-1}, I_s) - \boldsymbol{\epsilon})\|_2^2\big],
\end{equation}

其中 $P_s$、$P_t$ 分别为源视角和目标视角的外参，$I_s$ 为源视角的真实图像。与传统背景翘曲（background warping）只能获取共可见区域不同，SDS 损失能为训练视角中被遮挡的不可见区域提供先验。此外，针对单目视频数据集中相机参数通常不准确的问题，将相机外参设为可学习参数联合优化，并引入相机平滑损失 $\mathcal{L}_{\text{smo}} = \sum_{t=0}^{N_T-1}\|P_t^{-1}P_{t+1}\|_2^2$ 约束相邻帧相机位姿的时间连贯性。运动平滑正则方面，沿用 MoSca（2024.05）~\cite{2405_17421} 的多尺度 ARAP 损失 $\mathcal{L}_{\text{arap}}$ 和跟踪损失 $\mathcal{L}_{\text{track}} = \|p_{t \to t'} - \hat{p}_{t \to t'}\|_2^2$（将高斯 3D 坐标光栅化为 XYZ 图，与基础模型推断的 2D 轨迹监督对齐）。总体损失为 $\mathcal{L} = \lambda_{\text{rec}}\mathcal{L}_{\text{rec}} + \lambda_{\text{geo}}\mathcal{L}_{\text{geo}} + \lambda_{\text{sds}}\mathcal{L}_{\text{sds}} + \lambda_{\text{arap}}\mathcal{L}_{\text{arap}} + \lambda_{\text{track}}\mathcal{L}_{\text{track}} + \lambda_{\text{smo}}\mathcal{L}_{\text{smo}}$，其中 $\lambda_{\text{rec}}=1.0$，$\lambda_{\text{geo}}=0.075$，$\lambda_{\text{sds}}=0.01$，$\lambda_{\text{arap}}=1.0$，$\lambda_{\text{track}}=1.0$，$\lambda_{\text{smo}}=0.01$。

\textbf{对 2D 先验误差的鲁棒性。}与 MoSca、SoM 等方法对 2D 跟踪先验误差高度敏感不同，SE(3) B-spline 的连续性天然具有缓解 2D 跟踪抖动的能力——B 样条的光滑插值将离散控制点之间的噪声平滑化。消融实验证实，即使对 2D 跟踪先验施加 $[-15, 15]$ 像素范围的随机扰动，iPhone 数据集上 mPSNR 仅从 20.17 降至 20.11（下降 0.06 dB），NVIDIA 数据集上 PSNR 从 27.81 降至 27.74（下降 0.07 dB），表现出极强的鲁棒性。

\textbf{运动表示的消融对比。}论文将 SE(3) B-spline 运动基与两种替代表示进行对比：（1）SoM 的 SE(3) 姿态变换（逐帧独立预测仿射变换，无连续性保证）；（2）MoSca（2024.05）~\cite{2405_17421} 的运动脚手架（稀疏图节点 + DQB 插值）。在 iPhone 数据集上，SE(3) B-spline 达到 mPSNR 20.17 / mSSIM 0.729 / mLPIPS 0.274，相比 SoM 姿态变换（18.17 / 0.675 / 0.328）和 MoSca 运动脚手架（19.26 / 0.698 / 0.296）均有显著提升——这表明连续运动表示对复杂动态场景的建模优势不仅来自 SE(3) 群结构本身，更来自 B 样条对时间维度连续性的显式保证。

\textbf{实验结果。}在 iPhone 数据集（5 个真实动态场景，用两台静态相机测试）和 NVIDIA 动态数据集（7 个场景，12 相机阵列）上，方法在所有指标上超越 SOTA。iPhone 数据集上达到 mPSNR 20.17、mSSIM 0.729、mLPIPS 0.274，相比 MoSca（19.33 / 0.718 / 0.274）在 mPSNR 上提升 0.84 dB，相比 SplineGS（15.52 / 0.483 / 0.371）提升 4.65 dB。NVIDIA 数据集上达到 PSNR 27.81、SSIM 0.871、LPIPS 0.049，超越 MoSca（26.76 / 0.854 / 0.070）和 SplineGS（27.12 / 0.872 / 0.052）。模型参数量 134.7M，训练时间 30 分钟（RTX 4090），渲染速度 45.124 FPS。对应点迁移精度（PCK-T）达到 0.833，优于 MarbleGS（0.803）、HiMoR（2025.04）~\cite{2504_06210}（0.809）和 MoSca（0.821）。局限性方面，大幅非刚体变形（如布料剧烈翻折）和快速运动模糊场景仍是挑战。\paragraph{高级初始化策略复杂动态重建}

MOSAIC-GS（2026.01）~\cite{2601_05368}（CVPR 2026，ETH Zürich \& Google）提出全显式且计算高效的动态场景重建框架，核心洞察是：仅凭光度优化推断场景运动既低效又不可靠，高质量初始化对重建保真度和收敛速度至关重要。

\textbf{Baseline 与问题。}先前单目动态高斯方法通常在光度优化阶段从零学习运动，初始点在动态区域覆盖不足时会被大量剪枝，导致训练时间延长和重建质量下降。实验观察表明，光度优化对场景初始化高度敏感：缺乏准确运动估计时，大部分动态区域点在早期即被剪枝，模型需在后续迭代中逐步恢复丢失数据。此外，先前方法多采用逐帧变形模型（如 MoSca、Shape of Motion），存储开销大且难以扩展到长序列。

\textbf{核心改动。}MOSAIC-GS 设计了四阶段预处理管线，在光度优化之前就利用多种几何线索建立可靠的初始运动：（1）用 RAFT 光流结合 Sampson 极线误差 $e_{\text{epi}}(\mathbf{x}) = \frac{(\mathbf{x}'^\top F\mathbf{x})^2}{(F\mathbf{x})_1^2+(F\mathbf{x})_2^2+(F^\top\mathbf{x}')_1^2+(F^\top\mathbf{x}')_2^2}$ 检测动态区域，仅将超过阈值 $\tau_{\text{epi}}$ 的像素标记为动态，有效过滤因光流估计误差或相机参数不准导致的误检；（2）用 SAM2 进行语义分割与跟踪，获得跨帧一致的动态物体实例掩码 $M_t^{(j)}$，具体做法是先在动态区域提取边界框作为 SAM2 提示词，对已检测到的物体掩码做减法后提取新实例，并通过反向传播 pass 将物体掩码扩展到物体尚处于静止状态的更早帧，低置信度掩码（$p_{\text{conf}}(M_t^{(j)}) < \tau_{\text{mask}}$）被过滤；（3）用点跟踪器（BootsTAPIR）采样 $N_p = 10000$ 个查询点生成 2D 轨迹，根据与分割掩码 $M_t^{(j)}$ 的多数重叠将每条轨迹分配到对应物体实例 $j$，利用深度图提升到 3D 场景流，再用 Kabsch+RANSAC 估计每个实例的刚体变换 $\hat{P}_{t+1}^i := R_t^j P_t^i + \mathbf{t}_t^j$ 进行去噪和补全——先前向处理所有帧，再后向处理以恢复早期不可见但对应物体已出现的点，剩余不可见点通过前后最近可见帧线性插值补全；（4）将精化后的场景流直接拟合为时间相关的 Poly-Fourier 曲线，作为动态高斯的初始变形参数。

\textbf{静态--动态场景解耦。}MOSAIC-GS 将场景表示为两个独立的高斯集合 $\mathcal{G} = \mathcal{G}_s \cup \mathcal{G}_d$，静态分量 $\mathcal{G}_s$ 和动态分量 $\mathcal{G}_d$ 各自初始化并遵循不同的致密化策略，最终在渲染时联合光栅化。这种解耦提升了参数效率——静态区域无需存储时间变形参数，动态区域无需覆盖背景——并简化了后续优化。

\textbf{初始化细节。}对于动态区域，每个轨迹查询点 $\mathbf{q}_t^i$ 从对应 RGB 帧提取颜色，尺度通过对数尺度高斯拉普拉斯（Laplacian of Gaussian, LoG）范数结合深度图估计，使高斯大小匹配局部细节密度。对于静态区域，采用在时间间隔采样帧上采点的策略，确保广泛覆盖同时避免平坦区域过采样。

\textbf{Poly-Fourier 轨迹编码。}对每条 3D 轨迹 $\{P_i(t_1),\ldots,P_i(t_T)\}$，求解线性方程组 $Ax = y$ 获得系数，其中设计矩阵行向量 $\Phi(t) = [1, t, t^2, \sin(\omega t), \cos(\omega t)]^\top$，系数 $x = [a_0, a_1, a_2, b_1, c_1, \ldots]^\top$。时间 $t$ 的位置偏移量由 Poly-Fourier 函数给出：

\begin{equation}
    \Delta\mu(t) = \sum_{k=1}^{d_p} a_k t^k + \sum_{k=1}^{d_F}\big(b_k\sin(k\omega t) + c_k\cos(k\omega t)\big),
\end{equation}

观测空间位置为 $\mu(t) = \mu_0 + \Delta\mu(t)$。旋转偏移通过类似 Poly-Fourier 输出加上单位四元数后归一化为有效单位四元数 $\Delta q(t)$，最终旋转 $q(t) = \Delta q(t) \otimes q_0$。该编码将逐帧变形模型压缩为少量系数，显著减少参数量和训练时间——相当于用``多项式+三角函数''的紧凑公式替代逐帧存储每个高斯的位置偏移，同时保持与非刚体变形方法相当的重建质量。与 Gaussian Flow 不同，MOSAIC-GS 不建模时间变化的颜色变形，避免模型通过人工颜色变化补偿运动不准确。

\textbf{光度优化与深度损失。}整体损失为 $\mathcal{L} = (1-\lambda_{\text{ssim}})\mathcal{L}_{\text{L1}} + \lambda_{\text{ssim}}\mathcal{L}_{\text{SSIM}} + \lambda_{\text{depth}}\mathcal{L}_{\text{depth}}$，其中深度损失采用 Pearson 相关损失而非直接 L1 距离，该设计对时序不一致的深度尺度变化保持不变性，仅约束相对几何关系，有效应对单目深度估计中帧间尺度漂移的问题。

\textbf{实验表现。}在 DyCheck 基准上，MOSAIC-GS 在 LPIPS 上达到新 SOTA（0.255 vs MoSca 的 0.264），PSNR 仅次于 MoSca；在 NVIDIA Dynamic Scene 原始版本上同样 LPIPS 最佳、PSNR 第二；在 Gaussian Marbles 修改版本上则以严格共可见性评估协议同时取得 PSNR 和 LPIPS 的 SOTA。值得注意的是，MOSAIC-GS 的训练和渲染速度均显著优于所有对比方法，这得益于场景解耦、紧凑变形表示和准确的运动初始化三者协同减少了优化所需的迭代步数。

\paragraph{层次化诱导流引导动态重建}

HAIF-GS（2025.06）~\cite{2506_09518}（被 NeurIPS 2025 接收）提出基于稀疏运动锚点的统一框架，同时提升动态场景重建的效率与精度。

\textbf{Baseline 与问题。}此前方法（如 SC-GS、SP-GS）引入稀疏控制点通过插值驱动高斯变形，但存在两个瓶颈：（1）变形场使用简单 MLP，表达能力有限，难以建模关节运动或空间非均匀运动；（2）训练仅依赖图像重建损失，缺乏显式运动监督，导致时序不一致和伪影。

\textbf{核心改动。}HAIF-GS 以稀疏 3D 运动锚点作为基本变形单元，通过三个互补模块增强变形质量。运动锚点通过最远点采样（FPS）在规范空间中初始化以确保均匀空间覆盖，每个锚点 $a_i$ 参数化为 3D 位置 $x_i \in \mathbb{R}^3$ 和可学习影响尺度 $\rho_i \in \mathbb{R}^+$。对于每个高斯 $g_j$，通过 KNN 检索其 $K$ 个最近锚点，并用归一化高斯核计算插值权重：

\begin{equation}
    \omega_{ij} = \frac{\exp\left(-\|\mu_j - x_i\|^2 / \rho_i^2\right)}{\sum_{x_k \in \mathcal{A}_j}\exp\left(-\|\mu_j - x_k\|^2 / \rho_k^2\right)},
\end{equation}

其中 $\mu_j$ 为高斯中心，$\mathcal{A}_j$ 为其近邻锚点集合。该权重将用于后续所有锚点到高斯的变形传播。

\emph{锚点滤波器（Anchor Filter）。}为每个锚点 $a_i$ 预测动态置信度 $\sigma_i \in [0,1]$，由轻量 MLP 以锚点空间位置编码和时间编码为输入。训练分两阶段：第一阶段所有锚点以软权重 $\sigma_i$ 参与变形，梯度回传至滤波器；第二阶段对 $\sigma_i$ 硬阈值化，仅保留动态锚点。为规范动态-静态分解，引入两个正则化损失：\emph{稀疏性损失} $\mathcal{L}_{\text{sparsity}} = \mathbb{E}_i[\sigma_i]$ 鼓励用尽可能少的动态锚点解释场景（遵循奥卡姆剃刀原则），\emph{熵损失} $\mathcal{L}_{\text{entropy}} = \mathbb{E}_i[\sigma_i(1-\sigma_i)]$ 惩罚接近 0.5 的不确定分数，促使滤波器做出置信的二元决策。

\emph{诱导流引导变形（IFGD）。}给定锚点位置 $x$ 和时间 $t$，诱导流 MLP 自监督地预测前后场景流 $\Phi_{\text{flow}}(x,t) = (\mathbf{d}^{t-1}, \mathbf{d}^{t+1})$，据此构造三个时移查询 $q_{t-1}=(x+\mathbf{d}^{t-1}, t{-}1)$，$q_t=(x, t)$，$q_{t+1}=(x+\mathbf{d}^{t+1}, t{+}1)$。共享 Deformable MLP 将每个查询映射为特征嵌入 $\hat{f}^{t+i}$，聚合为时间一致特征：

\begin{equation}
    \bar{f}^t = \lambda\,\hat{f}^{t-1} + (1 - 2\lambda)\,\hat{f}^t + \lambda\,\hat{f}^{t+1}, \quad \lambda \in (0, 0.5),
\end{equation}

其中 $\lambda=0.25$。该聚合将场景流预测与场景重建桥接，实现联合优化，无需外部流标签。直观上，它让锚点同时``看到''前一帧、当前帧和后一帧的特征，并以前后帧为参考来校准当前帧的运动预测。聚合特征经多个变换头预测锚点级平移 $\Delta T_i^t \in \mathbb{R}^3$ 和旋转 $\Delta R_i^t \in \mathbb{R}^3$，再通过 KNN 找到每个高斯的 $K$ 个最近锚点，用线性混合蒙皮（LBS）插值得到高斯级变形：

\begin{align}
    \Delta\mu_j^t &= \sum_{i \in \mathcal{A}_j} \omega_{ij}\left(\Delta R_i^t(\mu_j - x_i) + x_i + \Delta T_i^t\right), \\
    \Delta r_j^t &= \left(\sum_{i \in \mathcal{A}_j} \omega_{ij} \cdot \Delta R_i^t\right) \otimes r_j,
\end{align}

其中 $\otimes$ 为四元数乘法。该插值使每个高斯的变形被局部稀疏锚点变换驱动，实现对密集变形场的紧凑建模。

\emph{诱导流时间自监督。}为鼓励场景流的双向一致性，引入循环一致性损失，约束前向流后接后向流回到原始位置：

\begin{equation}
    \mathcal{L}_{\text{cycle}} = \mathbb{E}_i\left[\|\mathbf{d}_i^{t-1} + \text{Flow}(x_i + \mathbf{d}_i^{t-1}, t{-}1)\|_2^2 + \|\mathbf{d}_i^{t+1} + \text{Flow}(x_i + \mathbf{d}_i^{t+1}, t{+}1)\|_2^2\right],
\end{equation}

其中 $\text{Flow}$ 即诱导流 MLP。该自监督信号对学习连贯运动模式至关重要，无需显式运动标注。最终总损失为 $\mathcal{L} = \lambda\mathcal{L}_1 + (1{-}\lambda)\mathcal{L}_{\text{D-SSIM}} + \lambda_1\mathcal{L}_{\text{cycle}} + \lambda_2\mathcal{L}_{\text{entropy}} + \lambda_3\mathcal{L}_{\text{sparsity}}$，其中 $\lambda{=}0.8$、$\lambda_1{=}0.01$、$\lambda_2{=}0.2$、$\lambda_3{=}0.5$。

\emph{层次化锚点加密（HAD）。}计算每个锚点在 $N_t=16$ 个采样时刻的平移方差 $\text{var}(a_i) = \frac{1}{N_t}\sum_{t=1}^{N_t}\|\Delta T_i^t - \overline{\Delta T_i}\|_2^2$，方差超过阈值 $\tau$ 的锚点被标记为需要加密，通过复制并添加小空间偏移生成更细分辨率的子锚点，形成由粗到细的层次结构，捕获细粒度非刚体变形。每个子锚点编码自身的空间位置、时间戳和父节点的平移量，实现跨层运动传播。所有层共享变形 MLP，但使用轻量的层级特定特征提取器以保持效率。对于每个高斯的最终变换，从各层独立检索近邻锚点并在层内做加权插值，然后用层级特定可学习置信度权重融合各层结果——这实现了多尺度运动分解，且只有需要细粒度建模的区域才承担额外锚点层级的计算开销。

\paragraph{增强先验的随意视频动态重建}

Prior-Enhanced Gaussian Splatting（2025.12）~\cite{2512_11356}（SIGGRAPH Asia 2025）发现动态高斯 Splatting 系统的重建瓶颈已从场景表示转移到 2D 先验质量，提出全自动三阶段管线（初始化→提升到 3D→动态场景重建），从基础模型中提取更高保真度的深度、掩码和轨迹先验。该方法不设计新的场景表示，而是构建在 MoSca 的运动支架（Motion Scaffold）表示之上，将增强后的先验直接注入 MoSca 的优化流程。

\textbf{Baseline 与问题。}此前工作（如 MoSca、Shape of Motion）依赖基础模型提供的深度图、动态掩码和 2D 轨迹作为监督信号，并在此基础上设计更丰富的场景表示或优化策略。但在细结构（如肢体）和快速运动区域，这些先验往往质量不足：掩码未能紧密跟随细部、深度图过度平滑丢失细节、轨迹在遮挡后断裂导致重建出现空洞。Prior-Enhanced GS 的核心洞察是：与其改进场景表示，不如先升级喂给系统的先验质量。

\textbf{核心改动。}Prior-Enhanced GS 围绕三项先验增强和两项新增损失展开。

\emph{物体级掩码。}将视频分割片段 $S_t^{(l)}$（由 SAM2 等视频分割模型生成）与极线误差掩码 $E_{t=1}^T$（由光流 RAFT 推导）取交集，保留覆盖运动面积超过 $\tau_{\text{salient}}=0.05$ 的片段，同时剔除外观变化（如移动阴影）导致的误检（自身运动面积占比须超过 $\tau_{\text{appearance}}=0.2$），得到紧密跟随细结构的物体级掩码 $M_t^{(o)}$。该掩码不仅覆盖运动部分，还扩展到整个动态物体，为后续深度精化和轨迹恢复提供精确的空间约束。

\emph{物体深度精化。}标准视频深度（Mega-SAM）在快速运动的细结构上过度平滑，例如冰球棍在深度图中被背景吞没。改进方案为：首先将 Mega-SAM 的默认单目深度替换为 MoGe 预测的仿射不变深度图，然后对每个动态物体掩码裁剪单目深度 $\tilde{D}_t^{(o)} = M_t^{(o)} \odot \tilde{D}_t$，在掩码内做尺度-位移对齐 $\alpha_t^{(o)}\tilde{D}_t^{(o)} + \beta_t^{(o)}$（通过掩码限制的尺度-位移拟合），并将物体深度损失注入 Mega-SAM 的一致性优化阶段：

\begin{equation}
    L_{\text{depth}}^{\text{object}} = \frac{1}{T\,|\Omega|}\sum_{o=1}^{O}\sum_{t=0}^{T}\sum_{p \in \Omega} M_t^{(o)}(p)\,\big|D_t(p) - (\alpha_t^{(o)}\tilde{D}_t(p) + \beta_t^{(o)})\big|,
\end{equation}

其中 $\Omega$ 为图像定义域。该损失仅作用于物体像素，在锐化细结构的同时保持全局尺度和时间一致性——相当于``只在运动物体区域加强深度监督，不动背景''。

\emph{骨架采样与轨迹恢复。}均匀采样轨迹偏向大块刚体区域（如躯干），严重欠采样细高运动区域（如肢体）。改进策略为：从物体掩码中提取 2D 中轴骨架（OpenCV 实现），将骨架膨胀 5 像素，按距掩码边界距离的倒数加权采样轨迹点。在实践中共 19384 条轨迹中，约 3000 条（1/6）来自骨架采样，其余均匀采样，确保细高运动区域被充分覆盖。轨迹传播使用 SpatialTracker 完成，轨迹帧统一缩放至长边 512 像素以保证一致性。对于遮挡后重新出现的轨迹点，若其仍在同物体掩码内则重新标记为可见，恢复的轨迹 $u_t^{(n)}$ 重新进入优化，填补遮挡造成的重建空洞。

\emph{自遮挡过滤。}同一物体的不同部分之间的遮挡（如手臂遮挡躯干）会使轨迹点虽仍在物体掩码内但实际已被遮挡。Prior-Enhanced GS 对每条恢复的轨迹在短窗口 $[t-2, t+2]$ 内重新采样并追踪，若新轨迹 $\{\hat{u}_{t'}\}$ 与原始轨迹 $\{u_{t'}\}$ 的距离超过 $\tau_{\text{self-occ}}=10$ 像素，则判定为自遮挡并保留``不可见''标记——``如果重追的路径跑偏了，说明点被别的表面盖住了，不算恢复''。

此外，引入两个新损失将增强先验传播到高斯云和运动支架：\emph{虚拟视角深度损失} $L_{\text{depth}}^{\text{virtual}} = \|\hat{D}_t^{\text{virtual}} - D_t^{\text{virtual}}\|_1$，通过随机平移训练相机（最大偏移 $0.18\,\text{median}(D_t)$）生成虚拟视角深度来惩罚漂浮物——``从旁边看一眼，如果那里有漂浮的碎片就罚''；\emph{支架投影损失} $L_{\text{track}}^{\text{scaffold}} = \|v_{2D} - u_t^{(n)}\|_2$，将运动支架节点投影回 2D 并约束其与对应轨迹点对齐，防止节点在细结构和快速运动区域漂移脱离物体。两个损失分别解决了``漂浮物''和``支架漂移''两类典型伪影，使重建结果在 DyCheck iPhone、NVIDIA 多视角和 DAVIS 真实视频数据集上均显著优于 MoSca 和 Shape of Motion。

\paragraph{局部隐式特征解耦动态建模}

LocalDyGS（2025.07）~\cite{2507_02363}（ICCV 2025）提出将复杂动态场景分解为基于种子的局部空间阵列，在每个局部空间内解耦静态与动态特征进行运动建模，是首个将动态高斯重建扩展到更大尺度高动态场景的方法。其核心思想借鉴 ScaffoldGS 的锚点-神经高斯架构，将静态场景中的锚点（anchor）推广为动态场景中的种子（seed），每个种子不再输出静态高斯，而是生成随时间变化的时间高斯（Temporal Gaussians）。

\textbf{Baseline 与问题。}先前多视角动态方法可分为三类：变形场方法（如 4DGS）将规范场高斯映射到变形场；轨迹跟踪方法（如 SpaceTimeGS）用多项式或傅里叶级数建模每个高斯的运动轨迹；4D 高斯扩展方法需要大量高斯点拟合，训练慢、存储大。它们共同的问题是对每个高斯显式建模长期运动轨迹，在大尺度复杂运动下容易出现闪烁和模糊。同时，单一时空模型难以同时承载大量静态和动态信息——作者实验验证，不进行特征解耦时动态和静态区域均出现模糊。

\textbf{核心改动。}LocalDyGS 不再跟踪每个高斯的长期运动，而是用种子将 3D 空间分解为独立局部空间，在每个局部空间内生成时间高斯建模运动。其设计理念是允许一个运动物体被多个种子分别建模：当手臂经过某个局部空间时，对应种子激活并生成时间高斯表示手臂运动；手臂离开后，时间高斯自动停用。多帧 SfM 点云融合（在时间域均匀采样 $N$ 帧）初始化种子位置 $\mu$，保证覆盖所有动态物体出现区域。每个种子定义为 $G_{\mathcal{S}}\{\mu, f_s, v\}$，其中 $\mu$ 为种子位置，$f_s$ 为静态特征，$v \in \mathbb{R}^3$ 为局部空间尺度参数，初始化为种子到最近三个邻居的平均距离——种子稀疏的区域局部空间更大。

\emph{特征解耦时空场。}每个局部空间分配独立的静态特征 $f_s \in \mathbb{R}^{64}$（跨时间步共享，捕获静态外观，初始化为零向量），同时所有局部空间共享一个动态残差场 $F_d$（由多分辨率 4D 哈希编码 + 浅层 MLP 构建），为每个种子-时间对 $(\mu, t) \in \mathbb{R}^4$ 提供动态残差特征。具体地，4D 哈希编码在 $L$ 个分辨率级别上分别查询哈希表（表大小 $2^{17}$，设置与 Instant-NGP 一致），得到编码特征 $f_h(\mu, t) = [h_{4d}(\mu,t;1), \ldots, h_{4d}(\mu,t;L)]$，再经浅层 MLP $\phi$ 交叉融合各分辨率特征：$f_d = F_d(\mu, t) = \phi(f_h(\mu, t))$。运动往往具有局部相似性，哈希编码的网格结构确保相邻种子共享相似特征，提供紧致的局部自适应表示。自适应权重场 $F_w$（浅层 MLP）预测平衡权重 $(w_s, w_d) = F_w(\mu, t)$，加权融合得到最终特征：

\begin{equation}
    f_w = w_s\,f_s + w_d\,f_d.
\end{equation}

静态特征承载大部分场景信息（其权重 $w_s$ 通常远大于 $w_d$），动态残差通常趋近零——仅在运动发生时提供时间特定的修正。论文的消融实验可视化显示，仅用静态特征渲染可得到清晰但不包含运动的场景，仅用动态残差渲染则呈现纯运动残差，两者加权融合后准确重建完整动态场景。该解耦设计使动态残差场的负担大幅减轻，避免了单一模型的信息过载。

\emph{时间高斯生成与去激活。}每个种子用特征 $f_w$ 通过独立的浅层 MLP（2 层 ReLU，输出经 Sigmoid 或归一化激活）解码出 $k{=}10$ 个时间高斯，其均值由 $\{\mu_t^i\}_{i=0}^{k-1} = \mu + v \cdot F_\mu(f_w)$ 给出（$v$ 为局部空间尺度参数，$F_\mu$ 输出 $k \times 3$ 向量），不透明度由 $\{\sigma_t^i\} = \text{Sigmoid}(F_o(f_w, \mathbf{d}))$ 解码，其中 $\mathbf{d} = (\mu - \mu_c)/\|\mu - \mu_c\|_2$ 为种子到相机中心的单位方向向量，旋转 $\{q_t^i\}$ 和缩放 $\{s_t^i\}$ 同理由各自 MLP 预测。关键的去激活机制：对于仅在特定时间 $t_a$ 有运动物体的局部空间，在其他时间 $t_b \neq t_a$ 大部分时间高斯的不透明度极低。设置阈值 $\tau_\alpha = 0.01$，将不透明度低于阈值的时间高斯停用，显著减少渲染时的计算量而不影响视觉质量。

\emph{自适应种子生长。}针对 SfM 点云覆盖不足（特别是弱纹理和有限观测区域）的问题，提出基于误差的种子生长策略（Adaptive Seed Growing, ASG）。在训练第 3000 至 15000 迭代之间，每 100 次迭代执行一次种子生长：记录每个时间高斯在所有查询时间上的最大 2D 投影梯度 $\nabla_{\max}^i$ 及其对应的 3D 位置 $\mu_{\max}^i$（即 $\{\nabla_{\max}^i, \mu_{\max}^i\} = \max_{t \in T}\{\nabla_t^i, \mu_t^i\}$），若 $\nabla_{\max}^i > \tau_g = 0.001$，则在 $\mu_{\max}^i$ 处添加新种子，逐步补全覆盖盲区。该梯度驱动策略使新种子精准补充高误差区域，增强模型对初始点云质量的鲁棒性。

\emph{体积正则化。}为鼓励每个时间高斯体积小、仅负责对应局部空间，引入体积正则化损失 $L_v = \sum_{i=1}^{M} \text{Prod}(s_t^i)$，其中 $M$ 为所有活跃时间高斯数量，$\text{Prod}(\cdot)$ 取缩放向量各分量的乘积。总损失为 $L = (1 - \lambda_{\text{SSIM}})L_1 + \lambda_{\text{SSIM}}L_{\text{SSIM}} + \lambda_v L_v$，其中 $\lambda_{\text{SSIM}} = 0.2$，$\lambda_v = 0.001$。

\emph{实验表现。}LocalDyGS 在细尺度数据集（N3DV、MeetRoom）上与 SOTA 方法竞争力相当，同时首次在大尺度 VRU 篮球场数据集（34 相机、$1920 \times 1080$、真实篮球赛场景）上验证了动态高斯方法的可行性，在大运动尺度场景下相比 SpacetimeGS 和 3DGStream 有更清晰的渲染质量和更少的闪烁。训练 30000 迭代，全部实验在单张 RTX 3090 上完成。

\paragraph{各向同性弹珠高斯动态建模}

Dynamic Gaussian Marbles（2024.06）~\cite{2406_18717}（Stearns et al., SIGGRAPH Asia 2024）针对随意拍摄的单目视频，通过简化高斯表示、分治优化策略和几何正则化三方面改进，引导优化走向泛化性更好的运动解。该工作的核心发现是：现有 4D 高斯方法在单目设置下严重失败，根本原因是单目场景极度欠约束——而通过精心设计的优化策略、现成的深度和运动估计以及几何正则化，可以恢复足够的约束条件。

\textbf{Baseline 与问题。}标准 3DGS 使用各向异性高斯（含旋转矩阵 $R \in \text{SO}(3)$ 和 3D 缩放向量 $\mathbf{s} \in \mathbb{R}^3$），在多视角静态场景下表达力强。但在缺乏多视角约束的单目动态场景中，额外的自由度使高斯能够以不泛化的方式拟合训练图像——例如，一个薄片状高斯可以从正面完美拟合像素，但从侧面看则完全失真。作者用一个简单实验直观验证了这一现象：在单张单目 RGBD 图像上分别用各向异性高斯和各向同性弹珠训练 100K 迭代，两者训练视角重建均完美，但各向异性高斯在新视角下产生明显伪影，而弹珠则泛化良好。

\textbf{核心改动一：各向同性``弹珠''。}将每个高斯的旋转固定为单位矩阵 $R = I$，缩放简化为标量 $s \in \mathbb{R}^1$。每个弹珠还附带语义实例标签 $y \in \mathbb{N}$（由 SAM 驱动的 TrackAnything 模型提供时间一致分割）和时间轨迹 $\Delta\mathbf{X} \in \mathbb{R}^{T \times 3}$（记录从初始位置到每帧的平移序列）。这种简化将每个高斯从 9 个几何自由度（3 位置 + 3 旋转 + 3 缩放）减少到 4 个（3 位置 + 1 标量缩放），大幅缩减欠约束优化的解空间——用``小球''代替``任意形状的椭球''，让模型不容易作弊。初始化时，对每帧视频获取单目（或 LiDAR）深度图（使用 DepthAnything），反投影为点云后进行离群点移除和下采样；每个点坐标 $p$ 初始化为一个弹珠，其均值 $\mu = p$、颜色取对应像素色、实例标签取分割预测，缩放和不透明度遵循标准 3DGS 协议初始化，轨迹 $\Delta\mathbf{X} = [\mathbf{0}]$ 即长度为 1 的零平移序列。

\textbf{核心改动二：分治运动估计。}不做全局一次性优化，而是将视频分割为短子序列，每帧独立初始化一组轨迹长度为 1 的弹珠集合 $\mathcal{G}_{ii}$，然后自底向上反复合并相邻子序列。每次合并包含三个阶段：（1）\emph{运动估计}——将弹珠轨迹从 $K$ 帧扩展到 $2K$ 帧，用等速假设初始化新帧平移 $\Delta\mathbf{X} = [\Delta\mathbf{x}_1, \Delta\mathbf{x}_2, \Delta\mathbf{x}_3^{\text{init}}]$，然后渲染到目标帧并反向传播更新轨迹，每帧优化 $\eta$ 步（DyCheck 上 $\eta{=}80$，Nvidia 上 $\eta{=}128$）；（2）\emph{合并}——取相邻子序列弹珠集的并集 $\mathcal{G}_{ij} = \mathcal{G}_{ij}^a \cup \mathcal{G}_{ij}^b$，剪除低不透明度和小缩放的冗余弹珠，并执行随机下采样以维持集合大小恒定；（3）\emph{全局调整}——对合并后的集合联合优化所有弹珠属性（颜色、缩放、不透明度、轨迹平移），每帧优化 $\beta$ 步（DyCheck 上 $\beta{=}32$，Nvidia 上 $\beta{=}48$），通过随机采样帧并渲染来更新。该策略的设计动机是：运动估计阶段受益于逐帧添加的局部性和平滑性（类似 Dynamic 3D Gaussians 的逐帧优化），全局调整阶段则贡献全局一致性（类似 4D Gaussians 的全局优化），两者交替可获得各自优势。分治课程在前景子序列长度达到 8（DyCheck）或 32（Nvidia）时停止，背景子序列最长可达 512 帧。

\textbf{核心改动三：跟踪损失与几何正则。}用 CoTracker 估计 2D 点轨迹 $\mathcal{P}$（时间窗口 $w{=}12$ 帧内的轨迹），对每个轨迹点 $p$ 找到其在图像平面上最近的 $K{=}32$ 个弹珠，惩罚它们与轨迹点之间距离的变化：

\begin{equation}
    L_{\text{track}} = \sum_{p \in \mathcal{P}}\sum_{g \in \mathcal{N}(p_i)} o'_i\;\Big|\,D_i\|\boldsymbol{\mu}'_i - p_i\|_2 - D_j\|\boldsymbol{\mu}'_j - p_j\|_2\Big|,
\end{equation}

其中 $\boldsymbol{\mu}'_i$ 和 $D_i$ 分别为弹珠中心 $\mu_i$ 在图像平面上的投影位置和深度，$o'_i$ 为弹珠对像素 $p_i$ 的不透明度贡献，$\mathcal{N}(p_i)$ 为像素 $p$ 的近邻弹珠集合。该损失约束弹珠在 3D 空间中与被跟踪点保持一致距离——``不管从哪帧看，弹珠到跟踪点的深度距离应该不变''。除跟踪损失外，每次优化迭代还包含多种渲染损失：对渲染图像、视差图和分割图分别计算 L1 损失（与真值图像、初始深度估计和分割预测对比），以及图像的 LPIPs 感知损失。

几何正则包括两个组件：（1）\emph{等距损失}——由局部等距和实例等距两部分组成。局部等距损失约束弹珠近邻在帧间保持刚体运动：当渲染到帧 $j$ 时，选取相邻源时间步 $i \in \{j{-}1, j{+}1\}$，计算
\begin{equation}
    L_{\text{iso-local}} = \sum_{g^a \in \mathcal{G}}\sum_{g^b \in \mathcal{N}(g^a)} \big|\|\mu_i^a - \mu_i^b\| - \|\mu_j^a - \mu_j^b\|\big|,
\end{equation}

实例等距损失则引导同一语义实例内的弹珠做近似等距运动：选取随机帧 $i$，在同类标签弹珠间施加同样的距离不变约束 $L_{\text{iso-instance}} = \sum_{g^a \in \mathcal{G}}\sum_{g^b \in Y(g^a)} \big|\|\mu_i^a - \mu_i^b\| - \|\mu_j^a - \mu_j^b\|\big|$，其中 $Y(g)$ 表示与 $g$ 具有相同语义标签的弹珠集合；（2）\emph{Chamfer 对齐损失}——在合并阶段将两组弹珠按来源帧随机分组（将每组再按原始帧拆分为子集后随机打乱，取前 25\% 和次 25\% 分别组成 $\mathcal{G}^1$ 和 $\mathcal{G}^2$），计算双向 Chamfer 距离：

\begin{equation}
    L_{\text{chamfer}} = \sum_{g^1 \in \mathcal{G}^1}\min_{g^2 \in \mathcal{G}^2}\|\mu^1 - \mu^2\|_2 + \sum_{g^2 \in \mathcal{G}^2}\min_{g^1 \in \mathcal{G}^1}\|\mu^1 - \mu^2\|_2,
\end{equation}

该随机分组引入随机性，避免两组仅因观测到场景不同部分就产生错误的推拉效应，推动两组在 3D 空间中对齐，避免``云雾状''伪影。

\begin{takeaway}
单目动态场景重建是最具挑战性但最实用的设定。核心难题是缺乏多视角约束导致的深度歧义，解决方案包括语义驱动的弹珠对齐、基于单目深度的几何先验注入、以及变形场的运动分解策略。该方向正从纯优化方法走向结合语义和物理先验的结构化建模。
\end{takeaway}

\emph{实验表现。}在 DyCheck iPhone 数据集（有/无相机位姿两个设置）和 Nvidia Dynamic Scenes 数据集上，Gaussian Marbles 在 PSNR/LPIPs 上显著优于 Dynamic 3D Gaussians 和 4D Gaussians 等高斯基线（Nvidia 上平均 PSNR 22.32 / LPIPs 0.129，DyCheck 有位姿 PSNR 16.72 / LPIPs 0.418、无位姿 PSNR 15.79 / LPIPs 0.428），且与 Nerfies、HyperNeRF 等 NeRF 方法性能持平，同时保持高斯表示的快速渲染（推理超过 200Hz on RTX-3090）、跟踪能力和可编辑性优势。在跟踪任务上，Gaussian Marbles 在 DyCheck 上的 PCK-T@5\% 指标甚至略微优于 CoTracker 的原始预测——在深度估计准确的场景中，弹珠能融合 3D 信息改善跟踪。

在 Total-Recon 数据集（室内猫/人体长视频）上，Gaussian Marbles 显著优于通用动态 NeRF 方法 HyperNeRF 和 D2NeRF，但与使用人体/四足动物先验的领域专用方法 Total-Recon 仍有差距。训练时间和显存随视频长度增长：DyCheck 上单张 A5000 训练 5--9 小时（峰值 13--22GB 显存），Nvidia 上单张 V100 训练 3.5--8 小时（峰值 6--10GB 显存）。消融实验表明：各向同性弹珠是最关键的设计选择，去除后性能下降最为显著；跟踪损失和 Chamfer 损失各自贡献了明显的提升；分治策略也优于一次性全局优化。

%% ============================================================

\subsection{近期方法脉络：Generative completion for dynamic Gaussian reconstruction、Foundation-model grounded 4D Gaussian reconstruction}

这一组近期工作延续 Dynamic 3DGS、Dynamic Gaussian Splatting，主要把改动集中在生成路线、增量改进等环节。下面只保留相对既有路线的变化，便于和前文方法直接对照。

\paragraph{World from Motion}
World from Motion（2026.07）~\cite{2607_01202} 可作为“Generative completion for dynamic Gaussian reconstruction”中的新锚点候选。相较于 Dynamic 3DGS，主要改变是：视频模型条件于外观、几何和三维运动渲染，并把生成结果蒸馏回动态 3DGS。高斯基元仍承担可渲染动态场景表示。
其设计动机是：构建 aligned multiview video pairs 和 dynamic 3DGS 数据集，训练模型沿输入和目标相机轨迹补全外观、几何和运动。
具体而言，初始单目重建会产生伪影和缺失区域，方法用 dense pixel-aligned renderings 条件化视频模型生成修正视图，并将新观察区域和运动重新蒸馏回动态 Gaussian 表示。
这一锚点判断来自当前分类中的机制关系：候选锚点：生成路线，且相关性 0.95；后续方法可围绕 将视频模型生成结果蒸馏回单一动态 3D Gaussian 表示。 继续比较和扩展。
适用边界是：依赖训练数据覆盖和视频模型先验；生成内容可能引入幻觉，需要蒸馏阶段约束。

\paragraph{Ground4D}
Ground4D（2026.06）~\cite{2606_28828} 延续 Dynamic Gaussian Splatting 在“Foundation-model grounded 4D Gaussian reconstruction”上的基本设定，主要改变是：3D foundation model 几何初始化结合动态 Gaussian 渲染优化。高斯基元仍承担可渲染动态场景表示。
其设计动机是：弥合动态 Gaussian Splatting 的高质量渲染与 3D foundation model 几何一致性的优势。
具体而言，第一阶段使用 VGGT 等 3D 基础模型以训练自由方式恢复几何和相机，第二阶段将这些几何先验接入动态 Gaussian 表示进行可渲染 4D 优化。
适用边界是：受基础模型几何质量和单目遮挡限制，快速运动和非刚性变形仍可能带来漂移。
\section{运动建模与跟踪}
\label{sec:06_4}

运动建模与跟踪是动态场景理解的核心任务，旨在从视频中提取场景元素的运动规律并持续追踪其状态。该任务面临多重挑战：（1）\textbf{遮挡}（occlusion）——运动物体被其他物体或自身部分遮挡后，需要推断其不可见状态的运动；（2）\textbf{快速运动}（fast motion）——大位移导致帧间对应关系难以建立；（3）\textbf{非刚体变形}（non-rigid deformation）——可变形物体（布料、软体、人体）的不同部分可能经历截然不同的运动模式；（4）\textbf{拓扑变化}（topology change）——物体出现、消失或分裂导致运动场不连续。本节按照方法类型将这些工作分为若干子主题进行讨论。

\subsection{持久跟踪与运动建模}
\label{ss:06_4a}

本节介绍两条技术路线：（1）持久跟踪与动态视图合成；（2）运动场与变形建模。

\subsubsection{持久跟踪与动态视图合成}

本路线介绍利用3DGS表示同时实现高质量动态视图合成和密集稀疏持久跟踪的方法。

动态场景重建和运动跟踪是同一硬币的两面——重建场景的 3D 表示自然产生运动轨迹，而运动轨迹约束又反过来提升重建质量。本节的方法正是利用这一紧密耦合关系，通过 3DGS 表示同时实现高质量的动态视图合成和密集/稀疏的持久跟踪。

\paragraph{持久动态视图合成与跟踪}

\textbf{Baseline 与问题。}先前动态场景重建主要分为两条路径：Eulerian 表示（如 HexPlane、K-Planes、SuDS）在固定 4D 时空网格上记录场景变化，无法建立帧间对应关系；Lagrangian 表示（如 ParticleNeRF）跟踪粒子穿越时空，但点模型缺乏旋转自由度，既无法建模 6-DOF 运动，也无法施加基于方向的物理先验。变形场方法（如 Nerfies）通过规范帧+形变场实现前后向映射，但需要昂贵的求根操作来获取正向对应，且单一规范帧严重限制动态表达能力。两类方法还有一个共同短板：跟踪要么需要预计算光流作为输入（如 OmniMotion 每对帧之间的密集光流，复杂度随帧数平方增长），要么根本不产生对应关系。

\textbf{核心改动。}Dynamic 3D Gaussians（2023.08）~\cite{2308_09713} 的核心思想极为简洁：在标准 3DGS 中，将所有高斯的颜色、不透明度和尺寸在时间轴上完全锁定，只允许每个高斯的位置和旋转自由变化——使高斯成为``随时间做刚体运动的定向粒子''。设计动机有两层：（1）如果每个高斯始终代表物理空间中同一块物质，那么只要通过可微渲染正确重建每一帧图像，密集 6-DOF 跟踪就从``持久动态视图合成''中自然涌现——跟踪不需要作为独立任务来解；（2）仅固定外观属性是不够的，在颜色均匀的大面积区域（如灰色衬衫与背景墙），高斯会自由滑动到相似颜色区域而不受惩罚，因此必须引入物理先验正则化来约束运动合理性。

\emph{动态高斯参数化。}每个动态高斯 $G_i$ 在时刻 $t$ 拥有如下参数：时变中心 $\mu_i^t \in \mathbb{R}^3$、时变旋转四元数 $q_i^t$（构造旋转矩阵 $R_{i,t}$）、固定的 3D 缩放 $s_i \in \mathbb{R}^3$、固定颜色 $c_i \in \mathbb{R}^3$、固定不透明度 $o_i \in \mathbb{R}$，以及固定背景对数几率 $bg_i$。每个高斯共 $7t + 8$ 个参数。场景通常由 200k--300k 个高斯表示，其中仅 30k--100k 个为非静态前景。高斯对空间中任意点 $p$ 的影响定义为：

\begin{equation}
    f_{i,t}(p) = \sigma(o_i) \cdot \exp\!\Big(-\tfrac{1}{2}(p - \mu_{i,t})^\top \Sigma_{i,t}^{-1}(p - \mu_{i,t})\Big),
\end{equation}

其中协方差矩阵 $\Sigma_{i,t} = R_{i,t}\,S_i\,S_i^\top R_{i,t}^\top$ 由固定缩放和时变旋转共同决定。该影响函数具有两个关键特性：局部性强（表示空间中一小块区域），但理论上无限延伸——即使高斯当前在错误的 3D 位置，远处的渲染误差仍能通过梯度将其推向正确位置，这正是``可微渲染实现跟踪''的梯度传播基础。

\emph{在线逐帧优化流程。}整个系统采用在线（online）逐时间步优化策略，无需离线训练任何额外数据。第 0 帧（初始帧）用稀疏深度采样初始化点云，然后像标准 3DGS 一样优化全部参数并执行自适应致密化（10000 次迭代，约 4 分钟）。从第 1 帧起，尺寸、颜色、不透明度、背景对数几率全部固定，致密化关闭，仅优化位置和旋转四元数（每帧 2000 次迭代，约 50 秒）。每帧开始时用前向速度估计初始化运动参数：$\mu_i^{t} \leftarrow \mu_i^{t-1} + (\mu_i^{t-1} - \mu_i^{t-2})$，旋转用归一化四元数做类似前推，并重置 Adam 的一阶/二阶动量参数——这一初始化步骤对结果至关重要，它为优化提供了正确的``起始猜测''。对于前景/背景颜色相似的区域（如灰色衬衫与背景墙），额外渲染前景/背景掩码，与``空场景帧''差分得到的伪真值掩码做 $L^{\text{Bg}}$ 损失，增强前景与背景的对比度。为处理多相机间白平衡、曝光等差异，为每台相机每通道独立优化一个缩放和偏移参数，仅在第 0 帧优化后固定。

\emph{三项物理正则化。}为保证物理合理性，引入三项局部正则化损失。最关键的是\textbf{局部刚性损失}，它要求每个高斯的近邻在其自身坐标系下做刚体跟随：

\begin{equation}
    L_{i,j}^{\text{rigid}} = w_{i,j}\,\big\|(\mu_{j,t-1} - \mu_{i,t-1}) - R_{i,t-1}^\top R_{i,t}(\mu_{j,t} - \mu_{i,t})\big\|_2,
\end{equation}

其中权重 $w_{i,j} = \exp(-\sigma_w \|\mu_{j,0} - \mu_{i,0}\|_2^2)$ 是非归一化各向同性高斯核，$\sigma_w = 2000$（对应标准差约 $\sqrt{2/2000} \approx 2.2$ cm），基于初始帧距离计算后所有时间步固定。$R_{i,t}$ 由四元数 $q_{i,t}$ 构造。直觉理解：如果高斯 $i$ 旋转了，近邻高斯 $j$ 必须像刚体一样围绕 $i$ 的中心做相应的旋转平移；如果 $i$ 平移了，$j$ 也要做一致平移。由于 $i$、$j$ 的角色是对称的，该损失双向施加，隐式迫使近邻具有相同的旋转变化——但显式约束收敛更快，因此引入\textbf{局部旋转相似损失}：

\begin{equation}
    L^{\text{rot}} = \frac{1}{k|S|}\sum_{i \in S}\sum_{j \in \text{knn}_i(k)} w_{i,j}\,\big\|\hat{q}_{j,t}\,\hat{q}_{j,t-1}^{-1} - \hat{q}_{i,t}\,\hat{q}_{i,t-1}^{-1}\big\|_2,
\end{equation}

其中 $\hat{q}$ 为归一化四元数表示，$k{=}20$ 为近邻数。$L^{\text{rigid}}$ 和 $L^{\text{rot}}$ 仅作用于相邻时间步（短期约束），长时间累积可能导致元素漂移，因此补充\textbf{长期等距损失}：

\begin{equation}
    L^{\text{iso}} = \frac{1}{k|S|}\sum_{i \in S}\sum_{j \in \text{knn}_i(k)} w_{i,j}\,\Big|\,\|\mu_{j,0} - \mu_{i,0}\|_2 - \|\mu_{j,t} - \mu_{i,t}\|_2\,\Big|,
\end{equation}

该损失比刚性弱（仅约束距离不变而非位置不变），但跨越任意长时间间隔，防止局部元素逐渐漂离。三项正则化仅作用于前景高斯，背景高斯施加独立的静止约束。消融实验表明：刚性损失是最关键的组件，单独使用即可取得良好结果；旋转损失和等距损失在定量指标上改善较小，但视觉上显著提升重建的物理连贯性。

\emph{跟踪的涌现机制。}训练完成后，3D 对应关系的获取非常自然：对于空间中任意点 $p$，找到对其影响最大（$f_{i,t}(p)$ 最高）的高斯 $i$，将 $p$ 转换到高斯 $i$ 在目标时间步 $t'$ 的局部坐标系下，即得到对应点——这是一个定义良好且可逆的一一映射。对于 2D 像素跟踪，先渲染深度图（用高斯中心深度替换颜色做体渲染），反投影到 3D，确定影响最大的高斯后跟踪到目标帧再投影回去。整个系统无需光流、骨架或任何对应关系作为输入，在 CMU Panoptic Studio 数据集（27 个训练视角、4 个测试视角、150 个时间步）上，单张 RTX 3090 训练约 2 小时，渲染 850 FPS；3D 跟踪中值误差仅 2.21 cm（约一个手腕宽度），2D 跟踪中值误差 1.57 归一化像素，比之前的最佳 2D 跟踪器 PIPs 好一个数量级（10x）。该方法还自然支持多种创意编辑：在单帧修改颜色后自动传播到所有帧、在不同场景间组合/移除动态对象、将虚拟相机``附着''到任意高斯实现第一人称视角跟随。值得注意的是，该方法存在两个设计层面的局限：（1）只能跟踪在初始帧中可见的场景部分——如果某个物体在第 10 帧才进入场景，整个管线将完全无法重建该物体，因为致密化只在第 0 帧执行；（2）需要多相机同步拍摄，无法直接应用于单目视频——这正是后续工作 DynOMo 所要解决的核心问题。此外，为简化优化，实验中关闭了球谐函数（view-dependent color），所有高斯仅使用 RGB 常量颜色。

\paragraph{自适应外观几何复杂度物体跟踪}

GTR（2025.05）~\cite{2505_11905} 从另一个角度利用 3DGS 实现跟踪：不同于 Dynamic 3D Gaussians 对整个动态场景进行全景重建，GTR 聚焦于\textbf{单个未知物体}的 6-DoF 位姿跟踪和高保真 3D 重建。问题设定为：从单目 RGBD 视频中跟踪并重建场景中运动的未知物体，仅需第一帧提供三个查询点、边界框或掩码之一即可启动。设计动机是：现有方法（如 BundleSDF）在面对轴对称物体、复杂几何和多样外观时容易失败，而 3DGS 的可微渲染能力可以在纹理缺乏时提供精确的位姿精调。

核心思想是\textbf{根据物体外观与几何复杂度自适应选择跟踪策略}。整个管线包含三个关键模块：

\emph{关键点检测与跟踪。}对新帧首先检测 SIFT 特征，根据特征密度量化外观复杂度。若特征密度低于阈值（低纹理区域），则使用 LoFTR 进行 detector-free 特征匹配；否则将 SIFT 特征输入 TAP 跟踪器。获得帧间对应关系后，使用 TEASER++ 快速点云配准估计粗略 6-DoF 相对位姿。

\emph{基于可见性与测地距离的关键帧选择。}同时优化所有帧的 3DGS 和位姿代价过高，因此维护一个关键帧池。关键帧选择基于两个条件：旋转测地距离是否超过阈值、关键点跟踪可见率是否低于阈值——任一条件满足即添加为新关键帧。低纹理帧总是被选为关键帧以增加约束。

\emph{自适应位姿精调。}对于几何复杂度低的物体（如轴对称物体），粗略关键点位姿估计不足以支撑高质量跟踪，此时利用 3DGS 的 render-and-compare 方法进一步精调位姿。具体而言，先使用 ICP 进行密集点云配准，然后在几何复杂度低于阈值时启动 3DGS 位姿精调——渲染颜色图和深度图，与观测值比较计算联合损失 $L_{GS} = \lambda_{\text{color}} L_{\text{color}} + \lambda_{\text{depth}} L_{\text{depth}} + \lambda_{\text{reg}} L_{\text{reg}}$，其中颜色损失为 $L_1$ + SSIM 加权组合，深度损失为 $L_1$。关键洞察是：几何复杂度高的物体（如带丰富表面细节的雕塑）仅靠 ICP 即可获得精确位姿，此时使用 3DGS 反而可能引入额外误差，跳过 3DGS 精调还能显著提升运行效率。当检测到回环（利用测地距离和 TAP 可见率判断）时，执行位姿图优化（Pose Graph Optimization）以消除累积漂移。最终使用 TSDF 融合从所有视角恢复完整纹理网格。

GTR 还提供了一个涵盖对称物体、复杂几何和多样材质的 GTR3D 评测基准，包含合成与真实两个数据集，支持完整 360° 视角覆盖的物体重建评估。该方法在单传感器开放世界 3D 重建中展现了 3DGS 在物体级位姿跟踪中的独特价值——自适应地``按需''使用可微渲染，而非像 Dynamic 3D Gaussians 那样将其作为唯一的运动驱动来源。

\paragraph{在线单目重建与跟踪}

DynOMo（2024.09）~\cite{2409_02104} 面向的核心任务是从无位姿单目相机的在线 2D 和 3D 点跟踪（Tracking Any Point, TAP）。其 baseline 是 Dynamic 3D Gaussians (D-3DGS)，但 D-3DGS 要求每帧 27 个多视角图像，直接搬到单目设定下性能急剧下降（D-3DGS-Mono 的 2D 中值平移误差增加 121.9\%，3D 增加 308.8\%）。DynOMo 的核心洞察是：通过精心设计的多模态重建监督和语义引导的正则化，时序对应关系可以从静态 2D 编码器中自然涌现（emergence），无需任何光流或对应关系级别的监督。为此 DynOMo 提出三项关键技术改动：(1) 为每个高斯增加语义实例标签 $g_i$ 和 DINOv2 视觉特征 $f_i \in \mathbb{R}^D$，将重建信号从 RGB 扩展到特征空间、深度和语义的多模态监督；(2) 利用语义标签将前景和背景高斯分离，使相机位姿优化仅依赖静态背景区域中已观测到的像素；(3) 将 D-3DGS 的刚性正则化权重从 3D 空间距离改为特征空间余弦相似度，并用实例 ID 过滤近邻选择。设计动机是：3D 空间中近邻的高斯可能分属不同物体（如手和腿），它们不应被约束为一起运动，而特征空间能更好地区分语义归属。

具体实现上，单目动态高斯定义为 $G_i^{\tau} = (\mu_i^{\tau}, q_i^{\tau}, s_i, c_i^{\tau}, o_i, g_i, f_i^{\tau})$，其中 $g_i$ 不随时间变化，$c_i^{\tau}$ 和 $f_i^{\tau}$ 允许时变但受时序平滑约束。整个在线管线基于 SplaTAM 实现，每帧依次执行：相机位姿优化（仅在静态背景已观测像素上施加重建损失）、新区域高斯增密化（探索场景中新出现的区域）、全部高斯联合优化、高斯和相机的前向传播。总重建损失为多种信号的加权和：

\begin{equation}
\mathcal{L}_{\text{rec}} = \lambda_I \mathcal{L}_I + \lambda_F \mathcal{L}_F + \lambda_D \mathcal{L}_D + \lambda_B \mathcal{L}_B
\end{equation}

其中 $\mathcal{L}_I$、$\mathcal{L}_F$、$\mathcal{L}_D$、$\mathcal{L}_B$ 分别是 RGB 重建（$\ell_1$）、DINOv2 特征重建（$\ell_2$）、深度重建（$\ell_1$）和背景掩码损失。消融实验表明特征重建和深度是两个最关键的监督信号：移除 $\mathcal{L}_F$ 导致平均跟踪精度从 45.4 降至 35.9，移除 $\mathcal{L}_D$ 降至 37.2。3D 正则化包括 D-3DGS 的三项物理约束（局部刚性 $\mathcal{L}^{\text{rigid}}$、旋转相似性 $\mathcal{L}^{\text{rot}}$、长期等距 $\mathcal{L}^{\text{iso}}$）和一项新增的时序平滑 $\mathcal{L}^{\text{sm}}$（约束颜色、特征和背景高斯位置的帧间一致性）。关键改动在于：近邻选择 $\mathcal{N}_i$ 仅在相同实例 ID 的高斯中选取（实例引导近邻），正则化权重 $w_{i,j} = s_{i,j}$ 改用特征余弦相似度而非 3D 欧氏距离。消融表明特征引导权重与刚性约束同等关键，二者协同工作：单目深度初始化的高斯位置噪声较大，特征引导提供了不确定性度量来抑制离群高斯对的影响。前向传播时，位置更新利用语义加权 kNN 常速假设：$\mu_i^{\tau+1} = \mu_i^{\tau} + \sum_{j \in \mathcal{N}_i} \text{softmax}(s_{i,j})(\mu_j^{\tau} - \mu_j^{\tau-1})$，旋转更新采用常旋转假设 $q_i^{\tau+1} = \Delta q_i^{\tau-1 \to \tau} * q_i^{\tau}$。

轨迹估计方面，对于查询像素 $p$，在可见高斯（可见度 $v_i^{\tau} > 0.5$，通过 $\alpha$ 合成累积不透明度计算）中选择投影距离最小的高斯 $G_p$，其 3D 轨迹直接由各时间步的均值 $\mu_p^{\tau}$ 给出，2D 轨迹经相机投影 $W^{\tau}\mu_p^{\tau}$ 获得。这是首个为无位姿单目相机在线点跟踪建立基线的工作。在 Panoptic Sports 数据集上，DynOMo 的 2D 中值平移误差为 6.3 像素（D-3DGS-Mono 为 23.3），3D 为 26.1 cm（D-3DGS-Mono 为 56.0），令人惊讶地在 2D 跟踪上甚至超越了使用 27 个摄像机的 D-3DGS（10.5 像素）。在 TAPVid-DAVIS 上，DynOMo 的 Average Jaccard 为 45.8，与离线方法 OmniMotion（51.7）可比，且无需任何运动监督或预训练；相比在线基线 SplaTAM（AJ 13.0）提升显著。推理速度约 45 秒/帧（RTX 3090），远快于 OmniMotion 的约 18 小时/序列（A100）。

\paragraph{点跟踪增强鲁棒重建}

现有基于可微渲染的动态 3DGS 方法（如 D-3DGS）通过光度损失间接驱动高斯运动。TrackerSplat（2026.04）~\cite{2604_02586}（SIGGRAPH Asia 2025）的核心诊断是：3DGS 的位置梯度计算被限制在有限局部邻域内，当物体经历大帧间运动并移出该区域时，位置梯度无法正确对齐物体运动方向，导致高斯褪色或错误变色（而非仅仅是光度损失不够强）。核心改动是：在训练之前，先用现有点跟踪模型（论文评估了 DOT、CoTracker 和 TAPIR，最终选用 DOT，因为其精度相当但速度显著更高）提取像素轨迹，通过并行加权增量最小二乘（PWI-LS）将 2D 像素轨迹直接转化为 3D 高斯的仿射运动，再经多视角三角化更新高斯的位置、旋转和缩放。设计动机是：显式的像素跟踪提供跨帧对应关系，比可微渲染的光度梯度更可靠地引导大位移运动，使得高斯在梯度优化之前已接近正确位置，训练过程更快更稳定。

具体实现上，对于首帧中的每个 2D 高斯（协方差 $\Sigma_{2D}$、均值 $\mu_{2D}$），定义其运动为仿射变换 $[A|b]$，将第一帧的 2D 高斯映射到后续帧。PWI-LS 通过最小化加权最小二乘求解：

\begin{equation}
[A|b] = V_1^{-1} V_2, \quad V_1 = \sum_{i=1}^{n} w_i P_i, \quad V_2 = \sum_{i=1}^{n} w_i Q_i
\end{equation}

其中 $w_i = \alpha_i T_i$ 是像素 $i$ 对该高斯的不透明度-透明度权重（$\alpha_i$ 为该高斯在像素 $i$ 处的不透明度，$T_i$ 为 $\alpha$-混合 中的累积透明度；边界像素权重低，避免被无关运动误导），$P_i$ 和 $Q_i$ 是像素级别的增量统计量，可在 GPU 上并行累加。为避免数值不稳定，论文丢弃覆盖像素少于 3 个、$V_1$ 行列式低于 $10^{-12}$ 或累积 alpha 低于 $10^{-3}$ 的高斯运动。求解后得到 2D 仿射运动，再通过 SVD 三角化恢复 3D 均值 $\mu_{3D}'$（至少需要 2 个视角，为数值稳定性丢弃少于 3 个视角观察的高斯），通过多视角约束线性系统求解 3D 协方差 $\Sigma_{3D}'$（每个视角的 $\Sigma_{2D}'$ 提供 3 个约束，$\Sigma_{3D}'$ 有 6 个未知参数，至少需要 2 个不同视角），最后特征值分解提取更新后的旋转矩阵和缩放向量。特征值分解有两个需要注意的问题：(1) 负特征值无法开方，因此丢弃含负特征值的 $\Sigma_{3D}'$；(2) 特征值默认按大小排序会破坏旋转和缩放参数在帧间的顺序一致性，导致相邻区域出现看似旋转相似但旋转矩阵差异巨大的不稳定现象，因此论文按首帧的量级顺序重排特征值和对应特征向量。

为进一步抑制离群运动，引入中值滤波和运动传播的正则化：对每个高斯在 $k$ 近邻中取位置、旋转、缩放位移的中值来平滑异常值。未被可靠确定的高斯（低可见度、少于 2 个视角观察、或特征值分解失败）从其近邻传播运动——取位置位移均值、缩放位移均值和旋转位移（Euler 角形式）均值来更新。运动传播前先检测静态高斯：覆盖超过 9 个像素且 90\% 以上像素在至少 2 个视角中位移小于 1 像素的高斯标记为静态，丢弃其计算出的运动并排除在运动传播之外。最后用输入视频做 1000 次迭代的精调训练（标准 3DGS 训练，不做 densification），此时高斯已接近正确位置，训练过程快速且稳定。

整个管线的三阶段均可并行化：点跟踪按视角独立并行，运动补偿和精调按帧独立并行（仅依赖点跟踪结果和首帧初始化）。论文在 8 块 NVIDIA A100 GPU 上实现并行管线，在 Meeting Room、N3DV、Dynamic3DGS、st-nerf 和 RH20T 五个数据集的短片段和长视频实验中，TrackerSplat 在多 GPU 并行设置下保持更高且更稳定的视觉质量，有效处理大帧间位移，显著减少褪色和变色伪影。首帧初始化使用 InstantSplat（2024.03）~\cite{2403_20309}（基于 COLMAP 点云优化 10,000 次迭代），为后续帧提供稳定的起始点。已知局限包括：小/薄物体（如手指）的点跟踪容易遗漏细微运动；低纹理区域的跟踪歧义和 3DGS 表示本身的时空不一致性会导致时序抖动；首帧误差会沿帧传播累积；首帧被遮挡的物体在后续帧出现时无法建立正确轨迹。

\paragraph{3D 高斯自监督跟踪}

现有视频自监督学习（SSL）方法采用时空 patch 预测目标，但该目标可以在不真正理解像素级跨帧对应关系的情况下被优化，因此学到的表示不适合点跟踪。Video-GMAE（2025.12）~\cite{2512_22489} 的核心改动是：训练一个 Masked Autoencoder 风格的编码器-解码器架构，将 $k=16$ 帧视频预测为``随时间移动的高斯集合''。设计动机是：将视频编码为 4D 高斯集合施加了天然的归纳偏置——2D 视频本质上是动态 3D 场景的一致投影，高斯必须跨帧保持身份一致，因此 SSL 任务被迫学习长期对应关系。与经典视频 MAE 方法（如 VideoMAE、MAE-ST）相比，Video-GMAE 的关键区别在于用``对应感知的高斯解码器''替代了朴素的 patch 重建解码器：即使编码器架构相同，高斯解码目标迫使模型学到更适合跟踪的特征表示。

编码器端，Video-GMAE 以预测图像级 3D 高斯的 GMAE 为初始化基础，在此基础上扩展为视频模型——在帧级 $16 \times 16$ patch 化后，添加可学习的时间轴位置嵌入，使 ViT 编码器能处理时序信息。输入采用 95\% 的高掩码率（与 MAE-ST 一致），仅 5\% 的可见 patch 进入编码器，极大降低计算开销。解码器接收编码器隐层和 $k \times n$ 个可学习查询 token，为第一帧预测 $n=256$ 个独立高斯基元 $G_0 = \{g_1, ..., g_n\}$（每个 $g_j \in \mathbb{R}^{14}$ 包含位置 $\mu$、缩放 $s$、旋转四元数 $\phi$、颜色 $r$ 和不透明度 $o$），为后续帧仅预测增量高斯并积分。为训练稳定和显式引入时序对应关系，增量只包含位置 $\Delta\mu_j$ 和颜色 $\Delta r_j$，其他参数在积分时保持不变：

\begin{equation}
G_t = G_{t-1} + \Delta G_{t-1}, \quad \Delta g_j = \{\Delta\mu_j, \Delta r_j\}
\end{equation}

直觉理解：第一帧建立高斯身份（包含完整的几何和外观参数），后续帧只调整位置和颜色，每个高斯的身份在 $k$ 帧内被强制保持，对应关系自然涌现。所有帧的高斯集合通过可微 Gaussian Splatting 渲染，以重建损失端到端训练。

零样本跟踪算法将预测的 4D 高斯轨迹直接转化为 2D 点轨迹。具体流程如下：首先将每个高斯的 3D 均值 $\mu_i^{(t)}$ 递归积分后投影到图像平面 $x_i^{(t)} = \Pi(K[R|t], \mu_i^{(t)})$，计算帧间 2D 位移 $\Delta x_i^{(t+1)} = x_i^{(t+1)} - x_i^{(t)}$。然后将位移编码为伪 RGB 值 $c_i^{(t)} = (\Delta x_{i,x}^{(t)}, \Delta x_{i,y}^{(t)}, 0)$，通过标准 $\alpha$-混合渲染出稠密流场：

\begin{equation}
F^{(t)}(u) = \sum_{i=1}^{n} \alpha_i^{(t)}(u) \cdot \Delta x_i^{(t)}
\end{equation}

其中 $\alpha_i^{(t)}(u)$ 是高斯 $i$ 在像素 $u$ 处的可微 splatting 可见度。对任意查询点 $p^{(0)}$，沿流场 $p^{(t+1)} = p^{(t)} + F^{(t)}(p^{(t)})$ 即可获得零样本轨迹。为处理遮挡，Video-GMAE 引入 anchor set 机制：在 $t=0$ 时刻为每个点选择 top-$k$ 个最相关高斯（$k=8$），定义 anchor mass $\omega^{(t)} = \sum_{i \in \mathcal{S}} \alpha_i^{(t)}(p^{(t)})$ 作为可见度指标。当 $\omega^{(t)} \geq \tau_{\text{vis}}$（阈值 $\tau_{\text{vis}}=0.5$）时，更新位置混合纯流场推进与 anchor 高斯加权位移（混合系数 $\beta=0.3$）；当 anchor mass 低于阈值时，完全依赖 anchor 高斯位移。若更新位置越出图像边界则标记为遮挡。这种设计利用了 Gaussian Splatting 渲染器的软像素分配特性，在无任何标注的情况下实现了遮挡感知的鲁棒跟踪。

监督微调阶段，Video-GMAE 采用交叉注意力读取头架构评估编码器特征对跟踪的迁移能力。具体而言，对编码器特征施加层归一化和可学习时间嵌入后，用 64 维 Fourier 位置查询通过 16 头交叉注意力机制读取隐层特征，经残差 MLP（隐藏维度为 token 特征维度的 4 倍，GeLU 激活）和 sigmoid 输出层预测 2D 跟踪坐标及遮挡概率。在 TAP-Vid 协议下评估（Average Jaccard、$\delta_{\text{avg}}^x$、Occlusion Accuracy 三项指标），支持冻结编码器（仅训练读取头）和端到端微调两种设置。

预训练使用 Kinetics、Kubric 和 Davis 的无标注视频，在 64 块 V100 上以 128 batch size 训练 90 个 epoch（AdamW 优化器，学习率 $10^{-3}$，2000 步预热加余弦衰减，梯度裁剪阈值 2.0）。零样本跟踪在 Kubric 和 TAP-Vid Davis 上与最强自监督方法 GMRW-C 可比，在 TAP-Vid Kinetics 上超过所有自监督方法；定性分析表明 Video-GMAE 产生的轨迹时间一致性更好、遮挡处理更稳健、运动估计更平滑。微调后在 Kinetics 和 Kubric 数据集上分别取得 34.6\% 和 13.1\% 的提升，超过所有监督和自监督基线。当前限制在于仅使用 $n=256$ 个高斯基元建模整个场景，空间分辨率有限，对大幅运动中的细小视觉细节跟踪精度不足。

\subsubsection{运动场与变形建模}

本路线介绍从隐式神经变形场到显式运动轨迹的各种运动场建模方法及其在动态3DGS中的应用。

如何高效且物理合理地建模场景运动场是动态 3DGS 的核心问题。从隐式神经变形场到显式运动轨迹，各种表示在不同场景中展现出各自优势。

\paragraph{显式运动引导高斯变形}

现有基于变形场的动态 3DGS 方法（如 Deformable 3DGS）仅依靠外观重建损失间接监督高斯变形，缺乏显式的运动引导。当物体运动不规则（如突然加速、急停）时，仅凭光度一致性容易陷入局部最优。一个自然的方案是用预训练光流网络提供 2D 运动先验，但光流同时包含相机运动和物体运动的贡献——静态物体在光流中也有非零值——两者混在一起不仅引入噪声，甚至可能误导变形优化。MotionGS（2024.10）~\cite{2410_07707} 的核心改动是引入光流解耦模块，将原始光流 $F_{t \to t+1}$ 分解为相机流 $F_{t \to t+1}^C$（仅由相机自身运动引起）和运动流 $F_{t \to t+1}^M$（仅由物体运动引起），只用运动流显式约束高斯变形。设计动机是：纯物体运动的运动流可以准确定位动态区域，直接约束 3D 高斯变形，避免相机运动造成的错误梯度。消融实验证实了这一点的关键性——直接用原始光流监督高斯变形反而比基线方法更差，正是因为相机运动和物体运动的混合产生了固有的歧义性。

具体实现上，光流解耦分三步完成。第一步，相机流通过深度图和相机位姿解析计算：先用深度图 $D_t$ 将像素 $p_t$ 反投影到 3D 空间 $\mathbf{x}_t = T_t^{-1} K_t^{-1} D_t \tilde{p}_t$，假设该 3D 点不动再投影到下一帧 $p_t^{t+1} = \pi(K_{t+1} T_{t+1} \mathbf{x}_t)$，相机流即为 $F_{t \to t+1}^C = p_t^{t+1} - p_t$。这里深度图 $D_t$ 直接由 3DGS 光栅化器渲染得到，消融实验表明使用离线深度估计器（如 MiDaS）效果较差，原因在于离线估计存在固有的尺度歧义（scale ambiguity），而 3DGS 渲染的深度图天然保持了尺度与几何一致性。第二步，运动流通过差值得到：

\begin{equation}
F_{t \to t+1}^M = F_{t \to t+1} - F_{t \to t+1}^C = p_{t+1} - p_t^{t+1}
\end{equation}

直觉理解：运动流就是假设相机不动、只观察物体运动时看到的光流。一个重要优势是，利用现成分割网络（如 SAM）提取动态物体掩码后，可以直接遮蔽静态区域的运动流，获得更清晰的运动监督信号——如果直接用原始光流则无法享受这一好处，因为静态物体在光流中也有非零值。消融实验也证实了运动掩码的作用：去掉掩码后性能下降，因为背景区域的不准确光流会引入错误的运动引导，导致高斯变形出现偏差。第三步，定义``高斯流'' $F_{t \to t+1}^G$ 来建立高斯变形与运动流之间的对应关系。高斯流将变形网络预测的 3D 高斯从时刻 $t$ 到 $t+1$ 的变形，直接在 CUDA 光栅化器中投影到相机 $C_{t+1}$ 的固定视角下，得到 2D 流场。高斯流的具体计算过程是：先将 3D 点 $\mathbf{x}_t$ 变换到时刻 $t$ 的规范高斯空间 $\hat{\mathbf{x}}_t = \Sigma_{i,t}^{-1/2}(\mathbf{x}_t - \boldsymbol{\mu}_{i,t})$，再将其映射到时刻 $t+1$ 的高斯状态 $\mathbf{x}_{i,t+1} = \Sigma_{i,t+1}^{1/2}\hat{\mathbf{x}}_t + \boldsymbol{\mu}_{i,t+1}$，第 $i$ 个高斯的流贡献为 $F_{i,t\to t+1}^G = \mathbf{x}_{i,t+1} - \mathbf{x}_t$，最终通过 $\alpha$-混合 聚合所有高斯贡献 $F_{t\to t+1}^G = \sum_{i=1}^K w_i F_{i,t\to t+1}^G$。

需要特别注意的是，这里的高斯流与 GaussianFlow（2024）~\cite{gao2024gaussianflow} 不同：GaussianFlow 计算的是相机从 $C_t$ 变换到 $C_{t+1}$ 视角下高斯变形的 2D 投影，对应的是完整光流；而 MotionGS 的高斯流固定在 $C_{t+1}$ 视角下计算，专门用于匹配运动流。运动流与高斯流之间用 $L_1$ 损失约束：$L_{\text{flow}} = \| \text{sg}(F_{t \to t+1}^M) - F_{t \to t+1}^G \|_1$，其中 $\text{sg}(\cdot)$ 为停止梯度（同时停止高斯流计算中时刻 $t$ 所有变量的梯度，提高训练效率）。光流网络方面，MotionGS 选择 GMFlow 作为默认光流估计器，消融实验对比了 FlowFormer 和自监督方法 MDFlow，结果表明光流估计的准确性直接影响最终重建质量——当光流先验本身不准确时，运动约束甚至可能产生负面效果。

此外，MotionGS 提出相机位姿精炼模块。在单目动态场景中，由于运动的复杂性和观测的稀疏性，即使 COLMAP 也难以准确估计相机位姿，而大多数现有动态 3DGS 方法忽略了这个问题。现有面向静态场景的位姿优化方法（如 CF-3DGS）无法直接用于动态场景，因为动态物体的存在使得相邻帧之间缺乏几何一致性。MotionGS 在相邻帧的相对位姿 $T$ 上引入可学习的 SE(3) 残差 $\Delta T$，使梯度可以反向传播到相机位姿。训练时采用交替优化策略（类似 Bundle Adjustment）：先冻结相机位姿、优化 3D 高斯参数；再冻结所有高斯属性、利用光度一致性损失反向传播梯度仅更新相机位姿。旋转增量和平移增量的学习率分别设为 $3 \times 10^{-3}$ 和 $10^{-1}$。交替进行直至收敛，用优化后的相对位姿更新 COLMAP 初始化的全局位姿，实现端到端的位姿精炼。冻结高斯优化相机的策略确保了训练稳定性和鲁棒性。

值得注意的是，该方法仍然依赖 COLMAP 的初始位姿估计作为起点，因此要求场景中存在足够的静态特征（如背景区域）——这在大多数真实世界场景中通常是满足的。此外，精炼后的相机位姿反过来提高了相机流的准确性，从而进一步改善了运动引导的质量，形成正反馈循环。运动约束和位姿精炼均仅在高斯开始变形和移动后才激活，以保证初始化阶段的训练稳定性。整个管线端到端可微，总训练损失为 $L = L_{\text{photo}} + \lambda L_{\text{flow}}$，其中 $L_{\text{photo}}$ 为基线方法的光度重建损失，$\lambda$ 为流损失权重（在 NeRF-DS 上设为 0.5，在 HyperNeRF 上设为 0.1）。消融实验表明，过大的 $\lambda$ 会干扰原有的光度重建优化过程，而过小的 $\lambda$ 则导致运动引导不足。

在 NeRF-DS 数据集上（8 个包含高速运动高光泽表面物体的日常场景），MotionGS 在快速运动和高复杂度场景中显著超越基线方法：例如在 plate 场景中准确渲染了移动盘子的反射和锐利边缘并消除了浮动伪影，在 basin 场景中有效建模了光滑的盆底表面。在 HyperNeRF 数据集上（使用智能手机拍摄的野外场景，包含拓扑变化且相机运动不规则、位姿不准确），MotionGS 同样取得一致的性能提升：在 chicken 场景中捕捉到红色外壳上的细微凹凸，在 broom 场景中准确渲染了快速运动的扫帚细节。整个训练在单张 RTX 3090 上完成（20,000 次迭代），在 NeRF-DS 大多数场景上保持实时渲染能力（FPS $>$ 30）。相较于基线方法，MotionGS 增加了高斯流渲染和相机位姿精炼的计算开销，训练时间和峰值显存均有所上升。此外，值得注意的是这些改进（光流解耦和位姿精炼）对具体的变形网络设计是无关的（agnostic），可以推广应用于其他基于变形的 3DGS 方法。在 DyNeRF 数据集上，MotionGS 遇到了局限性：该数据集的固定稀疏相机视角阻碍了准确的深度渲染，导致后续的相机流计算不准确并产生伪影——这并非光流估计本身的失败，而是由于 DyNeRF 并非连续单目视频而是稀疏视角的动态场景，也不利于规范的 3D 高斯初始化。

\paragraph{速度场渲染动态视角合成}

基于变形的动态 GS 方法（如 4DGS）虽然渲染质量好，但其变形网络预测的高斯轨迹缺乏局部一致性，无法准确对齐真实物体运动。这一问题在复杂运动和显著尺度变化时尤为突出：变形网络往往过拟合到不规则的高斯轨迹，导致次优的视频重建质量。此外，继承自静态 3DGS 的梯度致密化策略依赖光度损失梯度，实验发现该策略在动态区域表现不足，无法有效捕获存在显著运动的场景内容。FlowGaussian-VR（2025.07）~\cite{2507_23704}（ISMAR 2025）针对这两方面问题，引入速度场渲染（VFR）管线和流辅助自适应致密化（FAD）策略。核心改动是为每个高斯分配 2D 速度属性，通过可微光栅化渲染出逐像素速度场，然后用光流作为真值监督。设计动机是：直接渲染速度场比先渲染相邻帧再用光流估计更有效，因为速度场通过 $\alpha$-混合天然聚合了多个高斯的运动贡献，且理论分析和实验均验证了这种基于 Splatting 的方法优于先渲染再估计光流的替代方案。

具体实现上，给定 COLMAP 重建的初始点云 $\mathbf{P}_0$ 和两个相邻时刻 $t$、$t+\delta t$，变形网络将点云变形至两个时刻，分别投影到 2D 图像平面后计算位移。每个高斯在时刻 $t$ 的 2D 速度 $v_i$ 即为其在相邻帧的投影位置差分。速度场渲染类比颜色渲染，用 $\alpha$-混合公式：

\begin{equation}
V(p) = \sum_{i=1}^{N} \alpha_i \, v_i \, \prod_{j=1}^{i-1} (1 - \alpha_j)
\end{equation}

直觉理解：像素速度是其背后所有高斯速度的加权平均，权重由不透明度和遮挡关系决定。该公式的前提假设是像素颜色主要由贡献其颜色的高斯决定，而这些高斯的运动也携带了像素的辐射变化，因此速度场天然地估计了逐像素光流。渲染出的速度场 $V(p)$ 与预训练光流模型（如 RAFT）估计的真实光流对齐。所有损失函数均采用时序滑动窗口机制——在每次训练迭代中采样 $\tau+1$ 个视频帧，生成 $\tau$ 个额外的速度场渲染，确保跨时间戳的速度场预测具有充分监督和一致性。

设计了三项基于速度场的损失：（1）窗口速度误差 $L_{\text{win}} = \|\hat{\mathbf{V}}_i(\tau) - \tilde{\mathbf{V}}_i(\tau)\|_1$ 在滑动窗口内的多个时刻将渲染速度场与真值光流直接对齐，确保预测速度与真实场景运动一致。消融实验表明，增大窗口尺寸（从 2 帧到 8 帧）带来更强的时序一致性和更准确的速度估计，因为相邻迭代的窗口重叠为训练过程提供了稳定性。（2）流变形误差 $L_{\text{warp}} = \|W_{t+\delta t \to t}(\hat{\mathbf{I}}_{t+\delta t}, \hat{\mathbf{V}}_{t+\delta t \to t}) - \tilde{\mathbf{I}}_t\|_1$ 利用渲染速度场将 $t+\delta t$ 时刻的渲染 RGB 图像变形回 $t$ 时刻，与真值图像比较，从另一个途径对齐预测速度场与真实运动。（3）动态渲染损失 $L_{\text{dyn}}$ 用 SAM-v2 分割动态前景掩码 $\mathbf{M}_{\text{dyn}}$，仅在动态区域计算掩码重建误差。损失图可视化表明，$L_{\text{win}}$ 和 $L_{\text{warp}}$ 的损失值主要集中在前景动态像素，而传统的 $L_1$ 光度损失在动态和静态区域分布均匀，这使得流损失天然适合引导动态场景的致密化。

流辅助自适应致密化（FAD）策略解决了速度场渲染仅能优化已有高斯、无法处理动态区域高斯缺失的问题。与传统致密化克隆和分裂已有高斯不同，FAD 直接在动态区域的规范空间中添加新高斯。具体流程：首先利用 $L_{\text{win}}$ 的损失图和梯度图，选取满足 $\mathcal{L}_{\text{win}}(\mathbf{p}) > \epsilon$ 且 $\nabla\mathcal{L}_{\text{win}}(\mathbf{p}) > \epsilon$ 的像素位置；然后用 SAM-v2 排除背景像素，仅保留动态前景候选。接着将候选像素利用渲染深度图和相机内外参反投影到 3D 世界空间，用 FPS 以比例 $r$（Nvidia-long 设 0.01，Neu3D 设 0.001）下采样。新点需映射回规范空间：$\mathcal{G}_c^{\text{cano}} = \mathcal{G}_c - D(\mathbf{P}_0, t)$，其中 $D$ 为变形网络。属性通过 kNN 加权插值从邻近高斯获取，并引入空间半径约束 $r_c$：仅当邻近高斯位于候选高斯的空间范围内时才参与插值，以缓解属性的空间不连续性。

此外，FlowGaussian-VR 还探索了时序速度精炼（TVR）后处理策略，利用扩展卡尔曼滤波（EKF）对训练后的高斯轨迹进行全局时序平滑。具体而言，状态转移函数 $\mathbf{f}(\mathbf{x}_{t-1}) = D(D^{-1}(\mathbf{x}_{t-1}, t-1), t)$ 通过变形网络建模高斯的时序演进，观测函数 $\mathbf{h}(\mathbf{x}) = \text{VRaster}(\mathbf{x})$ 将高斯位置映射为速度场渲染输出。EKF 在预报步和同化步之间迭代：当累积协方差增大时对状态施加更强的修正，从而平滑轨迹并减少锯齿状抖动。由于变形网络不可逆，状态转移的 Jacobian 通过变形网络的 Jacobian 之比近似：$\mathbf{J}_{\mathbf{f}}(\mathbf{x}_{t-1}) \approx \mathbf{J}_D(\mathbf{x}_t) \cdot \mathbf{J}_D(\mathbf{x}_{t-1})^{-1}$。TVR 包含两个关键技术设计：（1）表面高斯过滤——排除投影深度大于渲染深度的被遮挡高斯（这些高斯不贡献光流值，跳过同化步仅执行预报步）；（2）精确光流定位——从第二帧起，光流观测的位置通过累加所有先前帧的流值计算 $z_k = F(p_0 + \sum_{i=0}^{k-1} z_i)$，而非直接使用当前高斯位置，避免因定位偏差导致的灾难性误差累积。实验表明，TVR 在 Balloon1 场景上使 PSNR 略降 0.18 dB、DPSNR 降 0.37 dB（因为 EKF 仅优化中心坐标而无法优化其他高斯属性），但高斯轨迹的物理一致性和可跟踪性显著改善。

在实验评估上，FlowGaussian-VR 基于 4DGS 基线在 Nvidia-long（7 个动态场景，90--210 帧，12 视角）和 Neu3D（6 个场景，300 帧，15--20 视角，分辨率 $1352 \times 1014$）数据集上进行测试。整体 PSNR 相比 4DGS 基线平均提升约 2.5 dB，动态区域 PSNR（DPSNR）分别提升 2.4 dB（Nvidia-long）和 3.14 dB（Neu3D），表明速度场建模对动态内容重建的增益尤为突出。消融实验验证了各组件的贡献：VFR + FAD + $L_{\text{warp}}$ + $L_{\text{win}}$ 的完整管线在 Balloon1 场景上从 20.51 dB 提升至 24.50 dB，最终高斯数量为 141k（基线 214k）。所有训练在单张 RTX 4090 上完成。

\paragraph{高斯泼溅 Lucas-Kanade 速度场分析}

现有动态 3DGS 方法在相机运动不足、场景高度动态的条件下表现不佳，通常依赖数据驱动的语义和几何先验作为正则化项，但这些先验偏向训练数据且忽略更广泛的运动动力学。Gaussian Splatting Lucas-Kanade（2024.07）~\cite{2407_11309}（ICLR 2025）摒弃了这一做法，转而将经典 Lucas-Kanade 光流方法适配至动态高斯 Splatting 框架。该方法的核心直觉是：一个准确的变形场应当产生连续的时空运动——既然所有高斯从同一个规范空间出发，那么速度场可以直接从变形网络的参数中解析地推导出来，而无需引入额外的可学习参数或外部网络。具体而言，前向变形场网络 $\mathcal{W}_{c \to}(\mathcal{G}, t)$ 将规范空间的高斯 $\mathcal{G}$ 变形至时刻 $t$ 的位置。对变形场关于时间求偏导，利用 Lucas-Kanade 式的链式法则分解，将速度表达为变形场对高斯参数的 Jacobian 与变形场对时间偏导的乘积：

\begin{equation}
v(\mathcal{G}; t) = J_{\mathcal{G}}^{+} \frac{\partial \mathcal{W}_{c \to}(\mathcal{G}; t)}{\partial t}
\end{equation}

其中 $J_{\mathcal{G}}^{+}$ 是变形场输出对高斯参数的 Jacobian 矩阵 $\boldsymbol{J}_{\mathcal{G}}$ 的 Moore-Penrose 伪逆（在 $\det(\boldsymbol{J}_{\mathcal{G}}) < \epsilon$ 时用伪逆替代逆矩阵以避免数值不稳定）。公式的含义是：速度等于变形场对时间的偏导数经 Jacobian 逆映射后得到的世界坐标系下的瞬时运动。该设计利用了前向变形场的关键性质——所有高斯从同一个规范空间出发、帧间具有双射对应关系——因此速度场可以直接从网络参数中解析获得，无需额外训练。值得注意的是，该方法同时计算位置速度 $v(\boldsymbol{p}; t) \in \mathbb{R}^3$ 和旋转速度 $v(\boldsymbol{q}; t) \in \mathbb{R}^4$（四元数表示），而对缩放和球谐系数不做速度推导，因为相邻帧间这些属性的变形量极小。在得到速度场之后，采用 Runge-Kutta 数值积分器进行时间积分以计算相邻帧间的场景流：

\begin{equation}
\mathcal{G}(t+\Delta t) = \mathcal{G}(t) + \int_t^{t+\Delta t} v_\theta(\mathcal{G}(t); t)\, dt
\end{equation}

该方法进一步将场景流解耦为光流（面内运动，optical flow）和深度变化（面外运动，depth variation）两个正交分量。从数学上看，一个 3D 点 $\boldsymbol{p}(t)$ 在相机 $i$ 下的投影为 $\boldsymbol{u}_i(t)$，对其运动做分解得到：

\begin{equation}
\frac{d\boldsymbol{p}}{dt} = \frac{\partial \boldsymbol{p}}{\partial \boldsymbol{u}_i}\frac{d\boldsymbol{u}_i}{dt} + \frac{d\boldsymbol{p}}{dt}\bigg|_{\boldsymbol{u}_i}
\end{equation}

其中第一项 $\partial \boldsymbol{p} / \partial \boldsymbol{u}_i \cdot d\boldsymbol{u}_i/dt$ 描述面内运动（由光流捕捉），第二项 $d\boldsymbol{p}/dt|_{\boldsymbol{u}_i}$ 描述沿视线方向的面外运动（由深度变化捕捉）。相应的场景流正则化损失分解为光流损失和深度排序损失的加权和：

\begin{equation}
\mathcal{L}_{scene} = \alpha \, \mathcal{L}_{o_{t \to t+1}} + \beta \, \mathcal{L}_{d_{t \to t+1}}
\end{equation}

其中光流损失 $\mathcal{L}_o$ 对渲染光流图与 RAFT 预测的参考光流取 $L_1$ 范数，深度损失 $\mathcal{L}_d$ 使用局部深度排序损失（将 Marigold 单目深度的排序知识迁移到渲染深度图上）而非直接 $p$-范数约束，以规避单目深度的尺度歧义问题。光流图和深度图通过与颜色图相同的 $\alpha$-混合 体渲染方程同时光栅化输出。实验表明，该方法在四个具有不同相机运动量的动态场景数据集（Dynamic Scenes、Plenoptic Video、HyperNeRF、DAVIS）上均取得领先结果。在相机运动极小、场景高度动态的条件下（如 Dynamic Scenes 数据集中 EMF 最低的场景），该方法在 PSNR 上超过 DeformableGS 约 1.5--4 dB，同时能够精确分离静态与动态区域并产生平滑、可解释的高斯运动轨迹。其局限性在于对 SfM 初始化质量敏感（动态区域点云不足时几何退化）以及伪逆运算的计算复杂度为 $O(N \times d)$（$N$ 为高斯数量，$d$ 为参数维度）。

\paragraph{运动轨迹场建模}

现有动态 3DGS 方法通常使用规范空间（canonical space）加 MLP 变形场的方式建模动态场景，但这种表示在处理长时间跨度的复杂运动时容易产生漂移且优化不稳定。3D Gaussian Representations with Motion Trajectory Field（2025.08）~\cite{2508_07182} 的核心改进是将运动建模从"逐帧变形"转为"轨迹空间分解"。具体做法是：为每个 3D 高斯定义一个全局参考点 $x_i^{\star} \in \mathbb{R}^3$（区别于规范空间参考点——全局参考点是轨迹上的任意位置，可以是起点或中点，而非所有时间步共享的单一规范坐标），用相对位置 $\Delta x_i^t = x_i^t - x_i^{\star}$ 建模运动，并将相对轨迹分解为 $k=40$ 个共享运动轨迹基的线性组合：

\begin{equation}
\Delta x_i^t = \sum_{j=1}^{k} \sigma_j(x_i^{\star}) \, \theta_j^t, \quad x_i^t = x_i^{\star} + \sum_{j=1}^{k} \sigma_j(x_i^{\star}) \, \theta_j^t
\end{equation}

其中 $\theta_j^t \in \mathbb{R}$ 是第 $j$ 个轨迹基在时刻 $t$ 的值（用 DCT 基初始化并在优化中微调，因为固定的 DCT 基往往无法捕捉真实世界中多样的运动模式），$\sigma_j(x_i^{\star}) \in \mathbb{R}^3$ 是由 MLP $\mathcal{F}_{tc}$ 编码的时间不变运动系数（对 $x$、$y$、$z$ 三个维度独立预测），轨迹系数由带位置编码的 MLP 生成：$\{\sigma_j(x_i^{\star})\}_{j=1}^k = \mathcal{F}_{tc}(\mathcal{G}(x_i^{\star}))$，其中 $\mathcal{G}$ 采用 $L=12$ 的正弦-余弦位置编码 $\mathcal{G}(x_i^{\star}) = (\sin(2^k \pi x_i^{\star}), \cos(2^k \pi x_i^{\star}))_{k=0}^{L-1}$，基于场景运动倾向于低频这一先验假设。全局参考点输入 MLP 时梯度被截断（detached），不参与运动场梯度的反向传播。公式的含义是：每个高斯的运动轨迹由一组全局共享的基轨迹加权组合而成，权重由该高斯的空间位置决定。这一设计的动机基于轨迹空间的对偶性——正如形状可以用少量基形状线性表示（2008）~\cite{akhter2008nonrigid}，运动轨迹也可以用少量基轨迹紧凑逼近。全局参考点 $x_i^{\star}$ 的引入有三个好处：(1) 用相对位置表示运动减少全局偏移；(2) 参考点可灵活选择或学习以适应不同运动类型；(3) 有利于自适应密度控制中剪枝透明高斯和加密高频区域。

\textbf{协方差的轨迹分解。}仅用时间变化的高斯中心不足以有效建模动态场景，因为 3D 高斯具有各向异性协方差，不同时刻可能需要不同的旋转来捕捉动态几何和外观变化。因此，缩放和旋转也分别通过独立的轨迹场建模为时间变化量：

\begin{equation}
s_i^t = s_i^{\star} + \sum_{j=1}^{l} \beta_j(x_i^{\star}) \, \lambda_j^t, \quad r_i^t = r_i^{\star} + \sum_{j=1}^{m} \gamma_j(x_i^{\star}) \, \eta_j^t
\end{equation}

其中 $s_i^{\star}$ 和 $r_i^{\star}$ 是全局缩放和旋转参考值，$\beta_j$、$\gamma_j$ 是由各自 MLP 编码的位置依赖系数，$\lambda_j^t$、$\eta_j^t$ 是缩放和旋转的轨迹基。而球谐系数（颜色）和不透明度保持时间不变，从而实现外观与运动的解耦——协方差和中心点已经解释了视角随时间的变化。

\textbf{静态/动态分离。}当使用隐式变形场建模完整动态场景时，静态部分会错误地获得运动轨迹，这不仅损害静态区域的新视角合成质量，还在训练运动场时引入噪声监督。为解决此问题，该方法为每个高斯分配一个连续状态参数 $p_i \in [0,1]$，表示其属于动态前景的概率。理想情况下，动态前景高斯应有 $p_i = 1$，静态背景高斯应有 $p_i = 0$。将渲染方程中的颜色 $c_i$ 替换为 $p_i$，即可通过体渲染输出静态/动态分割图 $\hat{\mathcal{M}}$。伪分割掩码 $\mathcal{M}$ 由 Segment and Track Anything 生成，分割损失为 $\mathcal{L}_m = \|\hat{\mathcal{M}} - \mathcal{M}\|^2$。渲染分割图时，高斯的位置、协方差和不透明度均被截断不参与梯度传播，使该损失仅优化状态概率 $p_i$。选择静态高斯的阈值设为 0.8（较高阈值），因为运动场可以建模静态背景但反向不行。此外，在训练后期引入逐点二值化损失 $\mathcal{L}_{3mr} = \frac{1}{k}\sum_{i=1}^{k} -\bigl(p_i \log p_i + (1-p_i)\log(1-p_i)\bigr)$，强制每个高斯收敛到确定的状态。

\textbf{运动正则化。}轨迹场中每个点仅提供位置依赖的系数，所有点共享同一组全局运动基函数，这仅提供了对平滑轨迹的弱先验。为注入基于物理的先验，方法引入两类正则化。第一，as-rigid-as-possible（ARAP）损失 $\mathcal{L}_{arap}$ 假设每个高斯与其 $k$ 近邻在相邻时刻之间应满足刚体变换，通过随机采样点集 $\mathcal{S}$ 并对每对邻居施加等距约束：

\begin{equation}
\mathcal{L}_{arap} = \frac{1}{k|\mathcal{S}|} \sum_{i \in \mathcal{S}} \sum_{j \in \text{knn}_i} w_{i,j} \| (x_j^{t-1} - x_i^{t-1}) - r_i^{t-1}(r_i^t)^{-1}(x_j^t - x_i^t) \|_2
\end{equation}

其中 $w_{i,j} = \exp(-\rho_w \|x_j(t) - x_i(t)\|_2^2)$ 是各向同性高斯加权因子。第二，空间平滑损失 $\mathcal{L}_{sp}$ 在全局参考点上施加微小扰动 $\epsilon = 0.001 \times \text{scale}$，鼓励邻近位置的轨迹系数一致：

\begin{equation}
\mathcal{L}_{sp} = \|\sigma_j(x_i^{\star}) - \sigma_j(x_i^{\star} + \epsilon)\|_2 + w_\beta \|\beta_j(x_i^{\star}) - \beta_j(x_i^{\star} + \epsilon)\|_2 + w_\gamma \|\gamma_j(x_i^{\star}) - \gamma_j(x_i^{\star} + \epsilon)\|_2
\end{equation}

上述所有运动正则化仅应用于动态高斯点（由前景/背景掩码选择）。由于动态区域光度损失大且同时受运动正则化约束，动态点梯度远大于静态点，若对两类点施加相同密度控制机制会导致动态区域过度密集，因此方法对动态点设置更高的梯度阈值（$8 \times 10^{-4}$，静态点为 $4 \times 10^{-4}$）。整体损失函数为 $\mathcal{L} = \mathcal{L}_{pho} + \mathcal{L}_m + \mathcal{L}_{3mr} + \lambda_a \mathcal{L}_{arap} + \lambda_s \mathcal{L}_{sp}$，其中 $\lambda_a = 0.3$，$\lambda_s = 0.6$。训练共 50k 次迭代：前 5k 次仅训练全局高斯点（不引入轨迹场以保证训练稳定），之后联合优化轨迹运动场直到收敛，密度控制在 25k 次迭代后停止。实验在合成数据集 D-NeRF（8 个场景，$800 \times 800$）和真实数据集 HyperNeRF（下采样至 $540 \times 960$）上进行。在 D-NeRF 上，该方法在大部分场景中取得最优 PSNR、SSIM 和 LPIPS，尤其在 Stand Up（46.28 dB）和 Jumping Jacks（40.01 dB）上大幅领先 D-3DGS。在 HyperNeRF 上，平均 PSNR 达 26.9 dB、SSIM 达 0.871，超越 Gaussian-Flow（26.3 dB / 0.862）和 TiNeuVox（24.3 dB / 0.837）。消融实验表明：(1) 静态/动态分离将 PSNR 从 25.4 提升至 26.9 dB，同时将高斯数量从 740k 降至 480k、渲染速度从 13 FPS 提升至 26 FPS；(2) 运动正则化虽略微降低渲染指标（PSNR 从 27.4 降至 26.9 dB），但产生更平滑、物理合理的运动轨迹。

\paragraph{逐高斯运动推理与迁移}

在线 4D 重建要求从流式多视角输入中低延迟地学习动态场景。现有方法仅用光度损失优化逐高斯运动，导致高斯"追逐像素残差"而非跟随真实 3D 运动。MoRGS（2026.03）~\cite{2603_25042} 的核心改进是引入光流作为显式运动线索来正则化逐高斯运动。具体做法分为三层：首先，在稀疏关键视角上用预训练网络估计相邻帧间的光流 $F_v^{\text{flow}}(x) = F^{\text{flow}}(I_{v,t}, I_{v,t-1})$ 作为 2D 运动先验；然后，将每个高斯的 3D 位移 $\delta_{i,t} = (\mu_{i,t} - \mu_{i,t-1})$ 通过相机投影的一阶线性化 $\pi_v$ 投影到图像平面，用 $\alpha$-混合 渲染出逐像素高斯运动图 $F_v^G(x) = \sum_{i=1}^{N} w_i(x) \, \pi_v(\delta_{i,t})$，与光流对齐训练：

\begin{equation}
L_{\text{flow}} = \sum_{v} \| F_v^{\text{flow}}(x) - F_v^G(x) \|_2
\end{equation}

该公式约束每个高斯在图像平面上的投影运动与观测到的像素级光流方向一致。其次，由于光流仅在稀疏关键视角上计算且存在视角偏差，方法引入可学习的逐高斯运动偏移场 $O_{i,t} \in \mathbb{R}^3$，将最终运动建模为流引导基础运动与学习偏移之和 $\delta_{i,t} = \delta_{i,t}^{\text{flow-guided}} + O_{i,t}^{\text{learnable}}$。偏移场从所有视角聚合梯度，弥补稀疏光流监督的不足并协调多视角一致性。最后，引入逐高斯运动置信度 $m_i \in [0, 1]$（通过 sigmoid 激活），将 2D 运动分割掩码提升到高斯域来区分动态与静态高斯，并以 $m_i$ 对属性残差进行加权 $A_{i,t} = A_{i,t-1} + m_i \odot R_{i,t}$，抑制静态区域的冗余运动更新。这三个组件协同工作，使 MoRGS 在保持可流式性能的同时实现在线方法中最优的重建质量和运动保真度。

\paragraph{抛体运动物理重建}

现有动态重建方法主要局限于短时、小尺度变形，对物理一致性的考虑有限。PMGS（2025.08）~\cite{2508_02660} 专注于通过 3DGS 重建大时空跨度的抛体运动，其核心改进是将牛顿力学先验引入位姿估计。该方法还贡献了一套面向复杂刚体运动建模的数据集，包含 6 个合成模型（Blender 中设置恒定重力场，随机初速度与自转，固定相机 120 FPS 采集 $1024 \times 1024$ 图像）和 3 个真实物体（固定相机，1/2500s 曝光，60 FPS，$4128 \times 2752$ 分辨率），所有运动严格遵循真实物理定律——这是许多现有合成数据集未曾考虑的。工作流程分两阶段：\textbf{目标建模}和\textbf{物理增强运动恢复}。

\textbf{目标建模阶段}旨在将动态场景转化为等价的静态物体中心化场景。首先用预训练的 SAM 分离背景与动态目标，然后通过建立物体中心化归一化坐标（2025.08）~\cite{2508_02660} 解耦复杂抛体运动：物体的自转被转化为伪多视角观测，平移位移则被消除。随后采用 3DGS 显式重建物体外观与几何。由于几何质量直接影响后续运动恢复，方法提出改进的点密度控制策略，以解决高斯表示的两个典型结构缺陷：过大的边界高斯体（导致锯齿边缘和尖刺伪影）和远离目标的离散化高斯（扭曲质心计算并引入系统误差）。具体而言，轴向约束剪裁主轴过大的高斯 $\max(\text{eig}(\Sigma_i)) > \tau_L$（$\tau_L$ 为控制形状各向异性的预设阈值），离群值去除剔除距最近邻过远的高斯 $\min_{g_j \in \mathcal{G}} \|\mu_i - \mu_j\|_2 > \tau_D \cdot D_{\text{avg}}$（其中 $D_{\text{avg}}$ 为所有高斯间的平均距离，$\tau_D$ 为距离滤波因子）。经过此策略处理，获得中心化尺度下的高斯场 $\mathcal{G}_0$。

完成中心化重建后，还需一个\textbf{注册（Registration）}步骤将静态高斯 $\mathcal{G}_0$ 从中心化尺度对齐回原始动态场景，以支撑后续运动恢复。注册通过一组可学习仿射变换 $T_{\text{reg}} = [\mathbf{r}, \mathbf{t}, s]$ 实现，对齐后的高斯场 $\mathcal{G}_1 = T_{\text{reg}} \odot \mathcal{G}_0$。除标准渲染损失 $\mathcal{L}_{\text{GS}}$ 外，注册阶段还引入 VGG 感知损失 $\mathcal{L}_{\text{VGG}}$ 和网格化全变分损失 $\mathcal{L}_{\text{TV}}$ 以提升对齐精度，总损失为 $\mathcal{L}_{\text{Align}} = \lambda_1 \mathcal{L}_{\text{GS}} + \lambda_2 \mathcal{L}_{\text{VGG}} + \lambda_3 \mathcal{L}_{\text{TV}}$，优化目标为最小化 $\mathcal{G}_1$ 渲染图像与原始动态场景首帧之间的差异。

\textbf{运动恢复阶段}通过学习逐帧 SE(3) 变换 $T_n = [\mathbf{r}_n, \mathbf{t}_n]$ 恢复完整运动序列。该阶段包含三个物理增强组件。

第一，\textbf{加速度一致性约束}将物体在重力场中的恒定加速度条件转化为位姿学习的正则项。首先用各高斯位置 $\mu_i = [x_i, y_i, z_i]$ 和半径 $R_i$ 计算质心（由于高斯已经过各向同性处理，半径立方 $R_i^3$ 作为质量权重）：

\begin{equation}
\sigma = \frac{\sum_{i=1}^{N} R_i^3 \cdot \mu_i}{\sum_{i=1}^{N} R_i^3}
\end{equation}

然后通过相邻帧质心的二阶差分计算加速度：

\begin{equation}
\mathbf{a}^t = \frac{\sigma^{t+\Delta t} - 2\sigma^t + \sigma^{t-\Delta t}}{(\Delta t)^2}
\end{equation}

将加速度分解为平行和垂直于重力方向 $\mathbf{g}$ 的分量：$\mathbf{a}_{\parallel}^t = (\mathbf{a}^t \cdot \mathbf{g})\mathbf{g}$，$\mathbf{a}_{\perp}^t = \mathbf{a}^t - (\mathbf{a}^t \cdot \mathbf{g})\mathbf{g}$。加速度一致性损失为：

\begin{equation}
\mathcal{L}_{\text{Acc}} = \| \mathbf{a}_{\parallel}^{(t+\Delta t)} - \mathbf{a}_{\parallel}^t \|_2^2 + \| \mathbf{a}_{\perp}^t \|_2^2
\end{equation}

该约束的物理含义是：沿重力方向的加速度应恒定（抛体运动特征），垂直于重力的分量应为零。这一约束不要求绝对的重力加速度值，即使在尺度模糊的情况下也能有效正则化逐帧平移分量的学习。

第二，\textbf{光流平滑约束}强调长期跟踪中运动变化的连续性。除评估真实光流 $\hat{F}$ 与渲染光流 $F$ 之间的相似度外，还计算光流场水平与垂直分量的梯度，并对相应区域施加不同权重的平滑惩罚，以抑制运动模糊和纹理退化区域的不准确光流估计带来的负面影响：

\begin{equation}
\mathcal{L}_{\text{Smooth}} = \lambda_1 \left[\frac{1}{N}\sum_{i=1}^{N}\left(\Delta F_{i,x} \cdot e^{-\frac{\Delta I_{i,x}}{10}} + \Delta F_{i,y} \cdot e^{-\frac{\Delta I_{i,y}}{10}}\right)\right] + \lambda_2 \|\hat{F} - F\|_1
\end{equation}

其中 $\Delta F$ 和 $\Delta I$ 分别为光流场和渲染图像的梯度，指数项使平滑惩罚在图像梯度大的区域（即边缘处）自动减弱。

第三，\textbf{动态模拟退火（DSA）策略}根据位移-时间的线性关系自适应调度学习率。对于初速度 $v_0$、恒定加速度 $a$ 的运动物体，从时刻 $t$ 到 $t+\Delta t$ 的位移可推导为 $\Delta s_t = (a\Delta t)t + [\frac{1}{2}a(\Delta t)^2 + v_0\Delta t]$。由于 $v_0$、$a$、$\Delta t$ 均为常数，因此 $\Delta s \propto t$，即位移与时间戳成线性关系。这一时空关系要求学习率与位移成比例缩放：大位移时使用较大学习率以快速追踪，小位移时使用较小学习率以精细调整。具体实现中，前一时刻学习到的位移控制下一时刻的初始学习率，随后使用指数衰减逐步精调。此外，对位移较大的时间戳适当增加迭代次数，以更好地寻找最优解。

运动恢复的逐帧跟踪总优化目标为最小化复合损失：

\begin{equation}
T_n = \arg\min_{T_n} \left(\lambda_1 \mathcal{L}_{\text{GS}}(\mathcal{R}(T_n \odot \mathcal{G}_n), \hat{I}_{n+1}) + \lambda_2 \mathcal{L}_{\text{Acc}} + \lambda_3 \mathcal{L}_{\text{Smooth}}(F(T_n \odot \mathcal{G}_n), \hat{F}_{n+1})\right)
\end{equation}

最后，PMGS 引入\textbf{跨模态卡尔曼融合}方案，将基于加速度一致性的位移预测和基于光流反投影的位移观测进行卡尔曼滤波融合，动态平衡不同观测源的权重、抑制多模态误差累积。与现有方法将物理约束仅作为正则化组件不同，PMGS 将物理模型作为系统预测的核心。卡尔曼融合包含四个步骤：\textbf{定义}——系统模型为基于加速度一致性约束的位移预测，两个观测源分别为光流反投影位移 $z_{\text{flow}}^t$（噪声方差 $\sigma_{\text{flow}}^2$）和网络学习位移 $z_{\text{learn}}^t$（噪声方差 $\sigma_{\text{learn}}^2$）；\textbf{预测}——利用状态转移矩阵 $\mathbf{F} = \begin{bmatrix} 1 & \Delta t \\ 0 & 1 \end{bmatrix}$ 和控制输入矩阵 $\mathbf{B} = \begin{bmatrix} \frac{1}{2}(\Delta t)^2 \\ \Delta t \end{bmatrix}$（控制输入 $a_t$ 来源于加速度计算），基于前一时刻的最优估计预测当前状态；\textbf{校正}——用两个观测源通过观测矩阵 $\mathbf{H}$ 校正预测状态；\textbf{更新}——基于卡尔曼增益 $\mathbf{K}_t = \mathbf{P}_{t|t-1}\mathbf{H}^T(\mathbf{H}\mathbf{P}_{t|t-1}\mathbf{H}^T + \mathbf{R})^{-1}$ 自动调整各源权重，输出最优估计。整个融合过程将物理模型、光流观测和网络输出充分整合，卡尔曼增益自适应地最小化多模态观测源的误差累积。

实验在合成（6 个模型）和真实（3 个物体）数据集上进行，从视频重建和运动恢复两个维度评估。视频重建使用 PSNR、SSIM、LPIPS 指标，运动恢复使用 IoU、ATE、RMSE 指标。对比方法包括静态方法（2DGS、InstantSplat）、动态方法（4DGS、DynamicGS、MotionGS）以及深度辅助基线 CFGS。PMGS 在所有跟踪指标上表现优异，特别是在运动恢复上 IoU 接近 1.0、ATE 显著低于竞争方法。4DGS 和 MotionGS 在大尺度运动下出现明显的轨迹断裂和伪影，DynamicGS 在某些实例上泛化能力有限，CFGS 即使在提供预训练 3D 模型的情况下也会在纯光度监督下将空间位置驱动到无穷大。消融实验验证了各模块的有效性：移除点密度控制导致几何退化并引发跟踪误差累积；移除加速度约束和 DSA 策略导致空间位置漂移和旋转误差；移除卡尔曼融合在更具挑战性的场景（如大的时变力干扰）下稳定性明显下降。

\subsection{语义与物体级运动}
\label{ss:06_4b}

本节介绍两条技术路线：（1）语义引导与层次化运动；（2）物体级跟踪与动力学学习。

\subsubsection{语义引导与层次化运动}

本路线介绍语义引导的运动控制与层次化运动表示方法。语义信息为运动建模提供了高层语义约束，层次化表示则为复杂运动提供了结构化的建模框架。

\paragraph{语义引导运动控制}

现有稀疏控制点方法（如 SC-GS）通过将百万级高斯压缩为数千个控制节点来降低计算量，但它们按几何均匀分配控制点，导致一个根本性的资源错配——静态区域浪费大量控制点，而真正需要精细建模的动态区域反而控制点不足。From Tokens to Nodes（2025.10）~\cite{2510_02732} 的核心改进是建立一套运动自适应的控制点分配与轨迹表示框架，使控制密度与运动复杂度精确对齐，从而同时解决"静态冗余"和"动态不足"两个问题。

该方法包含三个紧密衔接的组件。第一，\textbf{补丁-令牌-节点生成}（Patch-to-Node Generation）。选取一组关键帧，将每帧图像划分为固定大小的 patch，用冻结的视觉基础模型（VFM）为每个 patch 提取语义令牌嵌入 $z_{t,p}$ 并估计深度图。每个 patch 中心 $u_{t,p}$ 经反投影得到 3D 坐标 $x_{t,p}$，形成候选节点集 $\{(x_{t,p}, z_{t,p})\}$。关键在于每个节点同时携带几何位置和语义描述符，为后续压缩保留了外观信息。第二，\textbf{运动自适应节点压缩}（Motion-Adaptive Node Compression，MANI）。候选节点数量仍然过大，需要压缩。方法采用迭代体素化加二部图软匹配的策略：从较小的体素尺寸 $v_{\text{init}}$ 开始，每次迭代中将同一体素内的节点进行成对匹配，保留相似度最高的 $r\%$ 对中的代表节点，然后逐步放大体素尺寸直到节点数降至目标阈值。压缩的核心是定义联合相似度：

\begin{equation}
\text{sim}(N_i, N_j) = \cos(z_i, z_j) - \lambda_{\text{fg}}(N_i, N_j)
\end{equation}

其中 $\cos(z_i, z_j)$ 衡量令牌级外观相似度——静态区域的 VFM 令牌跨视角一致，而运动导致令牌变化从而降低相似度，因此余弦相似度天然充当了动静判别线索；$\lambda_{\text{fg}}(N_i, N_j)$ 是冻结分割模型预测的前景先验，防止动态区域被过早合并。更重要的是，方法为每个体素簇计算动态倾向评分 $p_{\text{dyn}}(C)$（综合前景先验均值和簇内成对相似度）：

\begin{equation}
p_{\text{dyn}}(C) = \sigma\!\left(\alpha \cdot \frac{1}{|\mathcal{U}_C|}\sum_{N_k \in \mathcal{U}_C} m(N_k) - \beta \cdot \frac{1}{|\mathcal{M}_C|}\sum_{(N_i, N_j) \in \mathcal{M}_C} \text{sim}(N_i, N_j)\right)
\end{equation}

其中 $\mathcal{U}_C$ 是簇 $C$ 中的节点集合，$\mathcal{M}_C$ 是簇内成对匹配的集合，$\sigma$ 是 sigmoid 函数，$\alpha$ 和 $\beta$ 控制前景先验与相似度两个信号的相对权重。该评分的直觉含义是：前景先验均值越高、簇内节点相似度越低，说明该簇更可能处于动态区域。基于此评分自适应调节压缩率：

\begin{equation}
r\%(C) = r_{\min} + (1 - p_{\text{dyn}}(C)) \cdot (r_{\max} - r_{\min})
\end{equation}

直觉理解：静态簇的 $p_{\text{dyn}}$ 低，被高压缩率大幅缩减；动态簇的 $p_{\text{dyn}}$ 高，被低压缩率保留足够密度的节点。这一设计从根本上解决了控制点密度与运动复杂度的匹配问题。第三，\textbf{样条参数化的节点轨迹}（Spline-Parameterized Node Trajectories）。在时间维度上，每个控制节点 $N_i$ 携带一个 SE(3) 轨迹 $T_i(t)$ 和一个 RBF 核函数半径 $\rho_i$。轨迹不是用 MLP 参数化，而是用三次 Hermite 样条在少量关键帧之间插值：

\begin{equation}
c(t) = h_{00}(s)\,P_k + h_{10}(s)\,(t_{k+1}-t_k)\,m_k + h_{01}(s)\,P_{k+1} + h_{11}(s)\,(t_{k+1}-t_k)\,m_{k+1}
\end{equation}

其中 $s = (t - t_k)/(t_{k+1} - t_k)$ 是归一化时间参数，Hermite 基函数分别为 $h_{00}(s) = 2s^3 - 3s^2 + 1$、$h_{10}(s) = s^3 - 2s^2 + s$、$h_{01}(s) = -2s^3 + 3s^2$、$h_{11}(s) = s^3 - s^2$，$P_k$ 是关键帧 $t_k$ 处的节点位置，$m_k$ 是对应的速度（切线）。样条构造同时保证位置和一阶导数的连续性，提供紧致且可微的轨迹表示。为稳定早期优化，样条轨迹不是从随机参数初始化，而是利用预训练 2D 跟踪模型（TAPIR）提取的 2D tracklet，通过深度反投影得到 3D 轨迹点：$x_t = R_t^\top \pi_K^{-1}(u_t, D_t(u_t)) - R_t^\top T_t$，其中 $\pi_K^{-1}(\cdot)$ 是基于内参矩阵 $K$ 的从图像到相机空间的反投影，$(R_t, T_t)$ 是估计的相机外参。再用最小二乘法拟合 Hermite 样条：

\begin{equation}
\min_{\{P_k\}_{k=1}^K} \sum_{t=0}^{N_f-1} \|x_t - \xi(t)\|_2^2
\end{equation}

其中 $\{P_k\} \subset \mathbb{R}^3$ 是关键帧处的可学习节点位置。这一几何驱动的初始化策略将轨迹锚定在观测到的运动模式上，避免了 MLP 变形场常见的优化不稳定问题。旋转分量在所有时刻初始化为单位矩阵 $R^{\text{node}}(t) = I_3$，留待联合优化阶段精调。

高斯 $G_j$ 的变形由其 $K$ 个最近邻节点通过双四元数混合（DQB）加权插值获得。节点 $N_i$ 到高斯 $G_j$ 的绑定权重为：

\begin{equation}
w_{ij} = \frac{\exp\!\bigl(-\frac{\|x_j - c_i\|^2}{2\rho_i^2}\bigr)}{\sum_{k \in V(G_j)} \exp\!\bigl(-\frac{\|x_j - c_k\|^2}{2\rho_k^2}\bigr)}, \quad \hat{Q}_j(t) = \frac{\sum_{i \in V(G_j)} w_{ij}\, Q_i(t)}{\|\sum_{i \in V(G_j)} w_{ij}\, Q_i(t)\|}
\end{equation}

其中 $Q_i(t)$ 是节点 $N_i$ 的 SE(3) 变换对应的双四元数表示（将旋转四元数 $q_r$ 和平移纯四元数 $q_d$ 编码为 $Q = q_r + \varepsilon q_d$，其中 $\varepsilon^2 = 0$ 是对偶单位），$\hat{Q}_j(t)$ 是归一化后的混合结果，再通过 $DQ2SE3(\cdot)$ 映射回 SE(3) 得到高斯的刚体变形。归一化保证混合结果始终为单位双四元数，确保刚体变换的物理合理性。相比直接对旋转矩阵或欧拉角插值，DQB 能更好地处理大角度旋转的插值。

整体优化使用复合损失 $L_{\text{total}} = \lambda_{\text{rgb}}L_{\text{rgb}} + \lambda_{\text{mask}}L_{\text{mask}} + \lambda_{\text{depth}}L_{\text{depth}} + \lambda_{\text{track}}L_{\text{track}} + \lambda_{\text{arap}}L_{\text{arap}}$，其中光度损失 $L_{\text{rgb}}$ 保证渲染一致性，掩码损失 $L_{\text{mask}}$ 用分割模型监督前景区域，深度损失 $L_{\text{depth}}$ 用单目深度估计（DepthCrafter）约束几何，跟踪损失 $L_{\text{track}}$ 用 2D 跟踪模型（TAPIR）的轨迹约束节点的时序一致性，ARAP 损失 $L_{\text{arap}}$ 惩罚局部邻域的非刚性变形以防止不真实的拉伸。

实验在 Hyper-NeRF（双 Pixel 3 手持拍摄，用其中一相机训练、另一相机评估）和 Neural 3D Video（N3DV，18--20 台同步相机，仅用 cam0 单目训练，在 cam5/cam6 评估）两个真实数据集上进行。在 Hyper-NeRF 上，该方法在所有场景和指标上均超越基线，包括 NeRF 方法（TiNeuVox、Hyper-NeRF、HexPlanes）和 3DGS 方法（D-3DGS、4DGS、ED3DGS、MoDec-GS、Grid4D、SC-GS、MoDGS）。值得注意的是，将 MANI 模块直接嵌入 SC-GS（记为 SC-GS+MANI）即可带来显著提升——例如在 Broom 和 Chicken 场景中重建出更完整的动态区域细节——验证了运动自适应初始化的通用性。在 N3DV 单目设置下，该方法达到 23.31 dB 的平均 PSNR，为当前最优。

消融实验进一步表明：运动自适应节点初始化优于 FPS、体素化、tracklet 聚类等替代策略；样条轨迹参数化优于 MLP 和网格变形场，也优于直接使用 tracklet 聚类的变形（后者虽受益于运动先验但受预测噪声干扰，优化不够稳定）。关键消融数据（Hyper-NeRF 场景，以 PSNR/SSIM/LPIPS 衡量）：基线（类似 SC-GS 架构）为 22.35/0.613/0.335，仅加入 MANI 提升至 23.89/0.635/0.315，仅加入样条轨迹（含 tracklet 初始化）提升至 24.51/0.658/0.278，样条轨迹若不加 tracklet 初始化则降至 24.13/0.639/0.284，完整方法达到 25.78/0.722/0.242。这表明两个核心模块（MANI 和样条轨迹）各自贡献显著，且 tracklet 初始化对样条轨迹的增益约为 0.38 dB。全部实验在单张 V100（32 GB）上完成。

\paragraph{图基运动层次学习}

现有方法在建模层次化运动时通常依赖人工定义或启发式的层次结构，并使用固定的运动基元，限制了跨任务的泛化能力。HEIR（2025.10）~\cite{2510_26786} 的核心改进是将层次推理形式化为一个可微图学习问题，从数据中直接学习结构化的运动父子依赖关系。具体做法是：给定 $N$ 个运动元素在时刻 $t$ 的绝对速度 $\Delta^t \in \mathbb{R}^{N \times d}$，构建一个邻近有向图 $G_0 = (V, E_0)$，用图注意力网络学习边权重 $w_{ij}$（表示顶点 $j$ 作为顶点 $i$ 父节点的概率），通过 Gumbel-Softmax 技巧采样离散的层次矩阵 $H \in \{0,1\}^{N \times N}$。核心分解公式为：

\begin{equation}
\Delta^t = H \, \Delta^t + \delta^t, \quad \text{即} \quad \Delta^t = (I - H)^{-1} \delta^t \approx \sum_{l=0}^{L_{\max}} H^l \, \delta^t
\end{equation}

其中 $\Delta^t$ 是绝对运动，$H\Delta^t$ 是从父节点继承的运动，$\delta^t$ 是局部运动残差。公式的含义是：每个运动元素的全局运动等于其局部运动加上从祖先链继承的运动之和，这正是截断 Neumann 级数的物理意义——$H^l$ 对应层次深度 $l$ 处的祖先关系。为保证层次结构合法，$H$ 须满足三个约束：每行恰好一个非零元素（单父）、对角线为零（非自父）、无环（DAG）。训练目标同时优化重建误差和残差稀疏性 $W^* = \arg\min_W \sum_t \|D(G(\Delta^t; W), H) - \Delta^t\|_1 + \lambda \sum_t \|G(\Delta^t; W)\|_1$，其中残差正则项迫使模型尽量通过父级继承解释运动，避免退化为星形拓扑的平凡解。在 3DGS 场景变形应用中，每个高斯作为图的顶点，学习到的层次矩阵与 ARAP（as-rigid-as-possible）求解器结合，用户只需拖动少量高斯，变形便沿层次结构从父节点传播到子节点，实现可解释的层次化场景编辑。

\paragraph{语义 3D 迁移学习}

传统运动迁移依赖骨架绑定（rigging），无法处理没有预定义关节结构的一般物体；2D 视频域的语义运动迁移（如 MotionClone）虽然能将运动与外观解耦，但缺乏 3D 一致性。Gaussian See, Gaussian Do（2025.11）~\cite{2511_14848} 的核心改进是将 2D 条件反演提升为多视角一致的 3D 运动迁移管线——不要求源和目标共享骨架，也不需要文本描述运动。具体做法分两阶段。第一阶段（结构化多视角运动反演）：从 $N$ 视角的源视频 $V_S$ 中提取运动嵌入。直接为每个视角独立优化嵌入虽然可行，但实验表明为某个特定角度训练的嵌入不一定适合其他角度的生成，且嵌入数量随视角线性增长。为此，作者提出锚定视角感知嵌入机制：只优化 $K$ 个锚定嵌入 $\{m_k^*\}_{k=1}^K$，均匀分布在源视角角度范围上；对任意中间角度 $\phi \in [\phi_i, \phi_j]$，通过球面线性插值获得嵌入：

\begin{equation}
m_\phi = \mathrm{slerp}(m_i^*, m_j^*, \frac{\phi - \phi_i}{\phi_j - \phi_i})
\end{equation}

公式的含义是：角度 $\phi$ 处的运动嵌入由其两侧最近锚点的嵌入平滑混合而成，插值参数控制两个锚点的权重比例。这样每次反向传播同时更新两个锚点，$K$ 个锚点即可覆盖全部视角，相比逐视角优化大幅降低存储和计算开销。第二阶段（视角感知语义迁移与 4D 整合）：渲染目标静态 3DGS 的随机视角图像，用对应的插值运动嵌入条件化冻结的视频扩散模型，生成目标物体在该视角下的 2D 监督视频。然后用多视角监督视频驱动 SC-GS 变形场重建 4D 场景。由于扩散模型生成的监督视频存在空间不一致性和噪声，直接使用 MSE 损失会导致变形场继承这些伪影，因此作者将像素级损失替换为 LPIPS 感知损失以增强鲁棒性；同时引入 ARAP 旋转机制——不使用 MLP 预测的旋转 $R_k^t$，而是用满足刚体假设的最优旋转 $\hat{R}_k^t$ 替代，直接在训练中而非后处理阶段约束旋转合理性：

\begin{equation}
\hat{R}_k^t = \arg\min_{R \in SO(3)} \sum_i w_{ik} \| (p_i^t - p_k^t) - R(p_i - p_k) \|^2
\end{equation}

公式含义：对每个控制点 $p_k$，寻找最优旋转使变形后的邻域点相对位置与原始位置尽可能一致。$w_{ik}$ 是高斯 $i$ 与控制点 $k$ 之间的蒙皮权重，由距离高斯核 $w_{ik} = \exp(-d_{ik}^2 / 2\sigma^2)$ 计算。该方法建立了首个语义 3D 运动迁移基准，在运动保真度和结构一致性上显著优于基线方法。

\paragraph{无位姿图像对动态重建}

现有 3DGS 方法通常假设静态场景和已知位姿。当仅有两幅无位姿图像时，动态场景的几何约束严重不足，传统光束平差（BA）无法处理运动物体。DynSUP（2024.12）~\cite{2412_00851} 的核心改进是将场景分解为分段刚性组件，在无位姿条件下联合恢复相机运动和物体运动。具体做法分两步。第一步（物体级密集光束平差）：用动态物体分割网络和 SAM-2 获取运动分割掩码，将场景像素划分为静态背景区域 $\mathcal{P}^{(0)}$ 和 $n$ 个动态物体区域 $\mathcal{P}^{(1)}, \ldots, \mathcal{P}^{(n)}$。对每个区域估计一个刚性变换 $T^{(i)} \in SE(3)$ 描述其相对运动，通过最小化重投影损失和深度正则化联合优化所有变换和逐像素深度：

\begin{equation}
\mathcal{L}_{\text{ba}} = \sum_{\mathcal{P}^{(i)}} (\lambda_1 \mathcal{L}_{\text{reproj}} + \lambda_2 \mathcal{L}_{d})
\end{equation}

公式含义：对每个刚体区域，重投影损失确保变换后的 3D 点投影到光流追踪的像素位置，深度正则化约束优化后的深度不偏离单目深度估计（通过可学习的缩放 $\alpha$ 和偏移 $\beta$ 参数）。静态区域的运动即相机运动 $T_{\text{cam}}$，动态物体在世界坐标系中的运动通过 $T_{\text{obj}}^{(i)} = T_{\text{cam}} (T^{(i)})^{-1}$ 恢复。第二步（SE(3) 场驱动高斯渲染）：在物体级变换基础上为每个高斯分配独立的 $SE(3)$ 变换 $T_{\text{fine}} \in SE(3)$，构成密集运动场。相比共享物体变换，逐高斯变换提供更灵活的优化空间；相比 MLP 或体素网格隐式建模，显式 $SE(3)$ 场更可解释且支持运动插值。为稳定优化，旋转采用连续 6D 表示（而非不连续的四元数），并通过平移和旋转正则化约束同一物体区域内高斯变换的一致性。该方法仅用两幅图像即可实现动态场景高保真新视角合成，在合成和真实数据集上显著优于针对静态场景或已知位姿设计的方法。

\subsubsection{物体级跟踪与动力学学习}

本路线介绍将物理规律引入物体级运动表示以实现可预测和可控制运动建模的方法。

从场景级运动建模深入到物体级跟踪和动力学学习，这些方法将物理规律引入运动表示，实现可预测和可控制的运动建模。

\paragraph{图神经网络驱动高斯动力学}

人类通过观察物体如何响应外部动作就能直觉地预判其未来运动，这种能力对机器人操控规划至关重要。然而，从真实视频中学习物体动力学面临一个根本性瓶颈：视频是 2D 像素的，而物体运动发生在 3D 空间；传统方法要么依赖物理仿真器（存在 sim-to-real 间隙），要么用 2D 视频预测（无法获得 3D 运动信息）。现有动态 3DGS 方法（如 Dyn3DGS）虽然能拟合给定视频的动态场景，但本质上只是一种事后重建——它们记录了物体``曾经''的运动，却无法预测物体``将要''如何运动。

Dynamic 3D Gaussian Tracking（2024.10）~\cite{2410_18912} 的核心改进是将动态重建与动力学学习统一为一条完整管线：先从多视角机器人交互视频中提取密集 3D 对应关系（回答``物体怎么运动的''），再在其上训练图神经网络（GNN）动力学模型（回答``给定新动作，物体将会怎么运动''），最终实现动作条件下的未来运动预测和操控规划。整个框架不依赖物理仿真器，直接从真实数据中学习动力学，从而天然避免了 sim-to-real 间隙。

具体做法分三步。

\textbf{第一步：密集 3D 高斯跟踪。}作者沿用 Dyn3DGS 的框架但做了三项关键调整：（1）固定每个高斯的颜色、不透明度和缩放不变，仅允许位置和朝向随时间变化。这一约束不仅减少了优化参数量，更重要的是保证了跨帧对应的一致性——尤其是当视频中出现部分遮挡时，颜色一致的椭球体更容易被正确关联；（2）仅保留高不透明度的高斯参与训练，剔除噪声粒子；（3）针对仅有 4 个视角的稀疏观测条件，加强了局部刚性和等距约束的权重。此外，利用 GroundingDINO 和 Segment Anything 获取机器人交互物体的分割掩码，过滤背景干扰，提升优化效率。跟踪输出为逐帧高斯位置序列 $\mu_i^t$（$i=1,\ldots,N$），其中索引 $i$ 跨帧严格一致，形成密集 3D 对应关系。

\textbf{第二步：图神经网络动力学。}从密集高斯中通过最远点采样提取 $n$ 个稀疏控制粒子 $\hat{\mu}_i^t$（$i=1,\ldots,n$）作为图顶点，机器人末端执行器位置 $a_t$ 也作为特殊顶点加入。图的构建规则为：任意两个顶点间 3D 距离低于阈值 $d_e$ 时连接双向边。顶点编码器的输入包含每个粒子前 $k=3$ 步的运动历史（$\mu_i^t - \mu_i^{t-1}$）、到桌面平面的距离，以及一个 one-hot 向量指示该顶点属于物体还是机器人末端执行器——这一区分对建模``动作如何施加到物体上''至关重要。边编码器的输入为两个顶点的 3D 位置差和一个 one-hot 向量指示边的关系类型（物体--物体或物体--机器人）。编码后的图通过 $p=3$ 步消息传递迭代更新顶点潜特征，共享解码器将最终的顶点特征映射为 3D 运动预测。整个前向过程可表示为 $\hat{\mu}_{t+1} = \mu_t + f_\theta(\mu_{t-k:t}, a_t)$。

训练损失为多步递推预测误差：

\begin{equation}
\mathcal{L}_{\text{pred}} = \sum_{i=1}^{L} \|\hat{\mu}_{t+i} - \mu_{t+i}\|^2
\end{equation}

公式含义：从随机采样的起始时间步 $t$ 出发，将模型预测的 $\hat{\mu}_{t+1}$ 递推作为下一步的输入，向前预测 $L=5$ 步，最小化每一步预测位置与真实位置的均方误差。取 $L=5$ 是在预测稳定性和计算成本之间的平衡。此外，作者还添加了两项辅助正则化：（1）边长正则化，惩罚相邻帧间图边长度的均方变化，维持物体局部结构的时序稳定性；（2）刚性正则化，对刚体或近刚体物体，惩罚预测运动与最优刚体变换（由 Umeyama 方法估计）之间的偏差，显式注入刚体先验。

\textbf{第三步：密集运动插值与规划。}稀疏控制粒子的预测运动需要插值到全部高斯，才能渲染完整的预测视频。插值分为两个子步骤。首先，为每个图顶点求解 6-DoF 变换：平移直接取 GNN 预测的运动向量；旋转则通过顶点邻域运动的最优拟合求解：

\begin{equation}
R_i^t = \arg\min_{R \in SO(3)} \sum_{j \in \mathcal{N}(i)} \| R(\hat{\mu}_j^t - \hat{\mu}_i^t) - (\hat{\mu}_j^{t+1} - \hat{\mu}_i^{t+1}) \|^2
\end{equation}

公式含义：对每个控制粒子 $i$，在旋转群 $SO(3)$ 上寻找使邻域粒子相对位置变换前后最一致的旋转矩阵——邻域粒子的运动``投票''决定该顶点的最优旋转。然后，通过线性混合蒙皮（LBS）将稀疏顶点的 6-DoF 变换插值到每个高斯：高斯中心更新为 $\mu_i^{t+1} = \sum_{b=1}^n w_{ib}^t (R_b^t(\mu_i^t - \hat{\mu}_b^t) + \hat{\mu}_b^t + T_b^t)$，旋转更新为四元数加权平均 $q_i^{t+1} = (\sum_{b=1}^n w_{ib}^t r_b^t) q_i^t$。混合权重与 3D 距离成反比：$w_{ib}^t = \|\mu_j^t - \hat{\mu}_b^t\|^{-1} / \sum_{b=1}^n \|\mu_j^t - \hat{\mu}_b^t\|^{-1}$，确保空间距离更近的控制粒子对该高斯有更大影响。从初始状态 $t=0$ 出发递推计算，即可得到全时序的密集高斯运动序列。

实验在绳子、布料和毛绒玩具三种可变形材料（共 8 个物体实例）上进行，使用 4 个相机以 15 FPS 采集多视角 RGBD 数据。在 3D 跟踪任务上，该方法在轨迹中位误差（MTE）和位置精度上均超越所有基线，包括使用深度作为 3D 先验的 SpatialTracker，且在机器人遮挡和物体形变等挑战条件下仍保持精确的密集对应。在动作条件视频预测任务上，与基于物理仿真器的方法（MPM 和 FleX，均通过 CMA-ES 优化仿真参数）相比，该方法在 3D Chamfer 距离、Earth Mover 距离、2D $J\&F$ 分割分数和 LPIPS 感知指标上全面领先——这表明直接从真实数据学习动力学比仿真参数优化更有效，因为前者天然避免了 sim-to-real 间隙。在模型预测控制（MPC）规划任务上，该方法成功应用于绳子拉直和毛绒玩具搬运，而 MPM 和 FleX 因计算效率过低无法在真实场景中运行。

\paragraph{高斯表面跟踪与重建}

现有动态 3DGS 方法要么追踪固定拓扑（外观和几何在新姿态下退化），要么逐帧独立重建（丧失时序一致性和跟踪能力）。GauSTAR（2025.01）~\cite{2501_10283} 的核心改进是将高斯绑定到网格面上构建"高斯表面"（Gaussian Surface），在拓扑一致区域保持跟踪，在拓扑变化区域自适应解绑并重建新表面，兼顾渲染质量、几何精度和时序跟踪。具体做法分四步。第一步（高斯表面表示）：在每个三角面上均匀分布 $N$ 个高斯，中心通过重心坐标 $\mathbf{p} = b_1 \mathbf{v}_1 + b_2 \mathbf{v}_2 + b_3 \mathbf{v}_3$ 由顶点位置计算，z 轴朝向面法线，z 缩放固定为小值 $\epsilon$ 以贴合表面。第二步（场景流变形初始化）：为应对帧间大幅度运动，将 2D 光流通过深度图反投影到 3D 空间获得场景流，并在网格连通邻域上做表面感知平滑：

\begin{equation}
F'(v) = \frac{1}{|\mathcal{N}(v)|} \sum_{u \in \mathcal{N}(v)} w(u,v) F(u)
\end{equation}

公式含义：对每个顶点 $v$，其场景流由网格邻域顶点的场景流按距离加权平均获得，确保变形在网格表面上的连续性。第三步（自适应高斯解绑）：在拓扑不变区域以固定拓扑优化重建；在拓扑变化区域，通过位置梯度和重建误差定义解绑权重 $W(f) = G_{\text{pos}}(f) + \lambda_{\text{rgb}} \mathcal{L}_{\text{rgb}}(f) + \lambda_{\text{depth}} \mathcal{L}_{\text{depth}}(f)$，权重高的面允许其高斯通过额外的旋转 $\Delta R$ 和平移 $\Delta \mathbf{t}$ 偏离网格面。解绑正则化 $\mathcal{L}_{\text{unb}}(g) = (1 - W(f_g))(\|R(g) - I\|_1 + \lambda_t \|\mathbf{t}(g)\|_1)$ 确保仅在拓扑变化处发生解绑。第四步（表面重网格化）：基于解绑高斯的位置从多视角渲染深度图，用 TSDF 融合重建新网格，然后通过边界顶点对应将新网格与原始网格连接。该方法在渲染指标上匹配或超越 SOTA，同时提供高分辨率可跟踪的显式网格表示。

\paragraph{物体中心重建与跟踪}

传统方案将 3D 重建（如 NeRF）和 6-DoF 位姿跟踪分开处理，使用不同的内部表示且通常需要物体的先验知识或预训练模型。Object-centric Reconstruction and Tracking（2024.05）~\cite{2405_20104} 的核心改进是用统一的 3D 高斯表示同时完成增量重建和位姿跟踪，整个框架不使用任何神经网络组件（纯基于可微渲染的梯度优化），无需预训练且不假设物体或运动先验。具体做法：将目标物体表示为一组 3D 高斯椭球，每个高斯由中心 $\boldsymbol{\mu}^i \in \mathbb{R}^3$、协方差 $\Sigma^i$、颜色 $c^i$ 和不透明度 $o^i$ 参数化（省略球谐系数以加速）。核心优化在重建和跟踪之间交替执行。重建步骤固定相机位姿，在包含当前帧和历史关键帧的滑动窗口 $W$ 内优化高斯参数（选取最大视角间距的 $W_{\max}$ 帧）：

\begin{equation}
\mathcal{L}_{\text{recon}} = (1 - \lambda)(\mathcal{L}_1)_{\text{color}} + \lambda (\mathcal{L}_{\text{SSIM}})_{\text{color}} + (\mathcal{L}_1)_{\text{depth}}
\end{equation}

公式含义：重建损失结合 L1 像素误差、结构相似性损失和深度 L1 误差，前者捕捉颜色差异，SSIM 衡量结构一致性，深度项约束几何精度。跟踪步骤固定高斯参数，仅优化当前帧的 SE(3) 位姿 $\mathbf{x}_{CO} = [\mathbf{t}_{CO}, \mathbf{q}_{CO}]^T$，使用 $\mathcal{L}_{\text{track}} = (\mathcal{L}_1)_{\text{color}} + (\mathcal{L}_1)_{\text{depth}}$。每帧的初始位姿通过线性外推 $\mathbf{t}^{t+1} = 2\mathbf{t}^t - \mathbf{t}^{t-1}$，$\mathbf{q}^{t+1} = \mathbf{q}^t (\mathbf{q}^t (\mathbf{q}^{t-1})^{-1})$ 预测。此外，通过比较渲染深度图和观测深度图自动添加新暴露区域的高斯点，并剪枝不透明度低于 0.6 的高斯以平衡表示精度与计算开销。该方法在 10 种未知航天器的合成序列上验证了重建和跟踪的有效性。

\paragraph{拓扑感知动态网格建模与跟踪}

现有 4D 高斯方法在重建动态网格时逐帧独立生成网格，导致帧间拓扑不一致（顶点数量和连接关系不同），使得骨架绑定和关键点跟踪等下游任务无法进行。TagSplat（2025.12）~\cite{2512_01329} 的核心改进是为 3D 高斯引入显式拓扑结构，通过拓扑感知的致密化和剪枝保持流形一致性，再配合时间正则化约束实现拓扑一致的动态网格序列。具体做法分四步。第一步（拓扑感知初始化）：从第一帧高质量多视角重建网格初始化高斯，每个高斯继承对应顶点的位置、颜色和连通关系，将拓扑先验注入高斯表示。第二步（拓扑保持致密化与剪枝）：致密化时，对投影梯度超过阈值的三角面插入新高斯点，属性取三个顶点高斯参数的均值，同时将原三角面细分为三个子三角面 $\{(v_0, v_1, v_{\text{new}}), (v_1, v_2, v_{\text{new}}), (v_2, v_0, v_{\text{new}})\}$，保持流形结构。剪枝时采用受边折叠（edge collapse）启发的策略，对需删除的高斯计算其与 1-环邻域高斯之间的折叠代价 $C = \lambda_g \|\mathbf{p}_i - \mathbf{p}_j\| + \lambda_a \|\mathbf{c}_i - \mathbf{c}_j\|$（结合几何距离和颜色差异），选择代价最小的边进行折叠合并。第三步（时间一致性约束）：为保证帧间拓扑一致，设计三个基于 1-环邻域的正则化项——边长一致性损失 $\mathcal{L}_{\text{len}} = \sum_i \|l_{t,i} - l_{t-1,i}\|$ 保持相邻帧边长稳定；刚性损失 $\mathcal{L}_{\text{rigid}}$ 约束邻域点相对位置在刚体旋转下一致：

\begin{equation}
\mathcal{L}_{i,j} = \|(\boldsymbol{\mu}_{j,t-1} - \boldsymbol{\mu}_{i,t-1}) - R_t(\boldsymbol{\mu}_{j,t} - \boldsymbol{\mu}_{i,t})\|
\end{equation}

公式含义：对每对相邻高斯 $(i,j)$，约束前后帧的相对位置关系仅相差一个旋转 $R_t = R_{i,t-1} R_{i,t}^{-1}$，即邻域应做近似刚体运动。旋转一致性损失 $\mathcal{L}_{\text{rot}}$ 约束相邻高斯的旋转变化一致。第四步（网格导出与跟踪）：训练完成后，高斯位置作为顶点、学到的拓扑定义连接关系，直接导出每帧的拓扑一致网格。由于拓扑帧间不变，任何 3D 点可用重心坐标和面索引参数化，其轨迹可跨帧跟踪。该方法在动态人体序列上重建出顶点数量和连接关系完全一致的网格序列，支持精确的 3D 关键点跟踪。

\paragraph{光线追踪驱动动态镜面场景}

传统方法（如 3DGS）使用球体近似法估计法线，在镜面反射场景中会引入不可忽视的误差——法线偏差会直接传播到反射方向 $\mathbf{d}_{\text{out}} = \mathbf{d}_{\text{in}} - 2(\mathbf{d}_{\text{in}} \cdot \mathbf{n})\mathbf{n}$ 的计算中，导致反射细节严重失真。TraceFlow（2025.12）~\cite{2512_10095} 的核心改进是选择 2D 高斯 Splatting 作为基础表示，因为 2DGS 天然具有无近似的精确法线，但原始 2DGS 只适用于静态场景且缺乏材质建模能力。为此，TraceFlow 设计了残差材质增强 2DGS（Residual Material-Augmented 2DGS）：为每个 2D 高斯引入可学习的镜面色调参数 $\mathbf{s}_{\text{tint}} \in [0,1]$（控制镜面反射的颜色），并设计时间条件残差网络 $\mathcal{F}_{\theta_G}$ 预测残差 $\Delta \mathbf{G}^t = \mathcal{F}_{\theta_G}(\mathbf{p}, \mathbf{t})$，以残差方式更新位置、缩放、旋转、不透明度和镜面色调等全部参数，使其能够捕捉动态变形。在渲染阶段，TraceFlow 引入动态环境高斯 $\mathbf{G}_{\text{env}}$ 替代传统的环境贴图（后者分辨率有限且无法建模近场反射），并通过混合渲染管线将颜色分解为漫反射和镜面分量：

\begin{equation}
\mathbf{C} = (1 - \alpha_{\text{spec}})\mathbf{C}_{\text{diffuse}} + \alpha_{\text{spec}}\mathbf{C}_{\text{specular}}
\end{equation}

其中 $\alpha_{\text{spec}} = \sum_{i} \mathbf{s}_{\text{tint},i}\,\alpha_i \prod_{j<i}(1-\alpha_j)$ 是由镜面色调加权得到的混合权重。漫反射分量 $\mathbf{C}_{\text{diffuse}}$ 通过标准光栅化获得；镜面分量 $\mathbf{C}_{\text{specular}}$ 则通过物理准确的光线追踪，沿反射光线查询动态环境高斯并累积至多 $k$ 个交点的贡献来合成。此外，TraceFlow 采用从粗到细的训练策略：先只训练漫反射分支建立稳定的几何基础，再逐步引入镜面分支，避免漫反射和镜面分量在训练早期相互纠缠导致收敛不稳定。

\begin{takeaway}
运动建模与跟踪将 3DGS 从静态重建扩展到运动场估计和物体跟踪。超点刚体变形提供了结构化的运动分解方案，语义驱动的弹珠对齐通过刚性约束和 Chamfer 对齐实现多视角时序一致，TraceFlow 在动态镜面场景中展示了 2DGS 精确法线的重要性。
\end{takeaway}

%% ============================================================


\section{视频重建与时序一致性}
\label{sec:06_5}

本节将前述动态表示技术推向实际应用——面向自由视角视频和体积视频的高效压缩、流式传输与时序一致性重建。

\subsection{问题背景：面向自由视角视频的4D重建}

前述各节讨论了变形场（\S\ref{sec:06_1}）、单目重建（\S\ref{sec:06_3}）和运动建模（\S\ref{sec:06_4}）等基础表示与重建问题。当这些技术面向实际应用——特别是自由视角视频（Free-Viewpoint Video, FVV）和体积视频（Volumetric Video）——时，涌现出一组新的系统级挑战：如何在保持实时渲染的同时压缩存储、如何在流式传输场景中保证时序一致性、如何处理长序列和大幅度运动、如何面向制作级应用提供足够的外观保真度。

核心矛盾在于：4DGS 的渲染速度极快（通常数百 FPS），但其百万级高斯基元的属性数据量巨大，单帧动辄数百 MB，使得存储和传输成为瓶颈。此外，逐帧优化范式在帧间引入闪烁和几何不一致，破坏观看体验。本节围绕"高效压缩与流式传输"和"时序一致性与高保真重建"两条主线组织，涵盖 47 篇论文。具体分为六条技术路线：（1）超高速渲染与表示效率；（2）压缩与编码；（3）流式传输与在线重建；（4）体积视频与高保真重建；（5）时序一致性与视频增强；（6）特殊应用与生成辅助。同时期还有

\paragraph{VeGaS：2024.11 同期工作}

VeGaS（2024.11）~\cite{2411_11024} 提出视频高斯泼溅方法。

\paragraph{GeoRect4D：2026.04 同期工作}

GeoRect4D（2026.04）~\cite{2604_20784} 从几何兼容角度改进稀疏视角动态重建。

\subsubsection{超高速渲染与表示效率}

本节聚焦如何从表示层面提升 4DGS 的渲染效率与存储效率，涵盖时空冗余剪枝、协方差解耦、时序敏感优化与结构级联变形四项技术。\paragraph{千帧率 4D 高斯渲染}

标准 4DGS 的渲染速度虽快，但在实际部署中面临存储膨胀和计算浪费的双重瓶颈。4DGS-1K（2025.03）~\cite{2503_16422} 首先对 4DGS 的冗余来源做了系统性分析，识别出两类互补的时间冗余，然后针对性地提出全局剪枝与局部滤波两阶段压缩管线，最终在现代 GPU 上实现超过 1000 FPS 的动态场景渲染。

\textbf{冗余分析。}在标准 4DGS 的训练结果中，作者发现了两个关键现象。第一，大量短寿命高斯（Short-Lifespan Gaussians）反复闪烁出现又消失——实验统计表明 70\% 的高斯时间方差 $\Sigma_t < 0.25$，且这些短寿命高斯集中在运动物体的边缘区域。它们消耗了存储空间却对长期渲染质量贡献甚微。第二，渲染时所有高斯都参与 4D$\to$3D 时间分解和光栅化，但每帧实际有贡献的高斯仅占约 15\%（其余为不活跃高斯，Inactive Gaussians）。更关键的是，相邻帧的活跃高斯重叠度（IoU）超过 80\%，说明活跃高斯集合在时间上高度稳定——这意味着可以通过帧间共享掩码来避免对不活跃高斯的冗余计算。

\textbf{第一阶段：基于时空变化评分的全局剪枝。}针对短寿命高斯造成的存储膨胀，4DGS-1K 设计了时空变化评分（Spatial-Temporal Variation Score）作为新型全局剪枝准则。直觉上，一个高斯如果频繁出现又消失（时间不稳定），或者对渲染贡献极低（空间不重要），就应当被剪除。评分由空间和时间两部分组成。空间分数 $\mathcal{S}_i^S$ 度量高斯对所有像素的光线贡献，沿用 $\alpha$ 合成的累积公式：

\begin{equation}
\mathcal{S}_i^S = \sum_{k=1}^{NHW} \alpha_i \prod_{j<i}(1-\alpha_j)
\end{equation}

公式含义：对高斯 $g_i$，在所有像素的光线中沿深度排序累积其 $\alpha$ 贡献，值越大说明该高斯对最终渲染图像的贡献越大。时间分数 $\mathcal{S}_i^T$ 则衡量高斯的寿命长短。4DGS 中每个高斯的时间不透明度服从高斯分布 $p_i(t) \sim \mathcal{N}(t;\mu_t,\Sigma_t)$，作者计算其二阶导数：

\begin{equation}
p_i^{(2)}(t) = \left(\frac{(t-\mu_t)^2}{\Sigma_t^2} - 1\right)\frac{1}{\Sigma_t} p_i(t)
\end{equation}

二阶导数的绝对值大意味着不透明度函数变化剧烈，对应短寿命、不稳定的高斯。通过 $\tanh$ 映射到 $(0,1)$ 区间并在时间轴上积分，得到不透明度变化评分 $S_i^{TV} = \sum_{t=0}^{T} \left(0.5\,\tanh(|p_i^{(2)}(t)|) + 0.5\right)$。此外还需考虑高斯的体积——体积大的高斯即使寿命短也可能覆盖较大区域，因此最终时间分数为 $S_i^T = S_i^{TV} \cdot \text{Norm}(V_i^{4D})$，其中 $V_i^{4D}$ 是高斯的 4D 体积并按 LightGaussian 的方法归一化。综合两部分，时空变化评分定义为：

\begin{equation}
\mathcal{S}_i = \sum_{t=0}^{T} \mathcal{S}_i^T \cdot \mathcal{S}_i^S
\end{equation}

公式含义：对每个 4D 高斯，将其时间维度上的全局重要性（时间分数）与渲染贡献（空间分数）相乘并在所有时间步上求和，得到综合评分。评分低的高斯被全局剪除，剩余高斯经过若干轮微调优化以补偿剪枝带来的轻微质量损失。

\textbf{第二阶段：基于关键帧的时间滤波。}针对不活跃高斯造成的计算浪费，4DGS-1K 引入基于关键帧的时间滤波（Key-frame based Temporal Filtering），在渲染时跳过不活跃高斯。具体做法是：以固定间隔 $\Delta t$ 等间隔选取稀疏关键帧时间戳集合 $\{t_i\}_{i=0}^T$，对每个关键帧 $t_i$，从所有 $N$ 个训练视角渲染图像并计算每个高斯的可见性掩码 $m_{i,j}$（第 $j$ 个视角在第 $t_i$ 时刻的可见性列表），最终取所有视角的并集 $\bigcup_{j=1}^{N} m_{i,j}$ 作为该关键帧的活跃高斯掩码。渲染时，对任意查询时间 $t_{\text{test}}$，找到其两侧最近关键帧 $t_l$ 和 $t_r$，仅对掩码 $\bigcup_{j=1}^{N} m_{l,j} \cup \bigcup_{j=1}^{N} m_{r,j}$ 标记为活跃的高斯执行 4D$\to$3D 分解和光栅化。由于相邻帧活跃高斯重叠度超过 80\%，关键帧掩码可有效覆盖周围多帧的活跃集合，掩码共享机制大幅减少了逐帧重新计算掩码的开销。若间隔过大可能遗漏某些高斯导致质量下降，因此对掩码记录的高斯额外微调以补偿。

\textbf{实验结果。}综合以上两阶段压缩——全局剪枝消除短寿命冗余、时间滤波跳过不活跃高斯——4DGS-1K 在复杂动态场景（如 Technicolor 和 Neural 3D Video 数据集）上实现了 41 倍存储压缩和 9 倍光栅化加速，在现代 GPU 上达到超过 1000 FPS 的渲染速度，同时保持与原始 4DGS 相当的渲染质量。

\paragraph{解耦时空 4D 高斯}

先前的 4DGS 方法（如 RealTime4DGS、Rotor4DGS）将时间和空间维度统一建模为 4D 协方差矩阵 $\Sigma_{4D} = R_{4D}S_{4D}S_{4D}^T R_{4D}^T$，在渲染时需要先对 4D 高斯做条件概率切片得到 3D 高斯，切片后得到的 3D 均值 $\mu_{3D} = (\mu_x,\mu_y,\mu_z)^T + (t-\mu_t)\frac{\mathbf{V}}{\mathbf{W}}$ 依然耦合了时空分量，且 4D 协方差矩阵的投影不是仿射变换，需要计算复杂的雅可比行列式来近似——这两者都带来冗余计算。Disentangled 4DGS（2025.03）~\cite{2503_22159} 的核心改进是将 4D 协方差矩阵分解为三个物理含义清晰的独立组件：3D 协方差矩阵 $\Sigma_{3D} = \mathbf{U} - \frac{\mathbf{V}\mathbf{V}^T}{\mathbf{W}}$、均值速度 $\mathcal{V}_{3D} = \frac{\mathbf{V}}{\mathbf{W}}$（表示 3D 均值随时间的变化率）、时间缩放 $s_t = \sqrt{\mathbf{W}}$。这一解耦表示只需 15 个浮点数（传统方法需 16 个），且时间相关变量（$\mathcal{V}_{3D}$ 和 $s_t$）是向量/标量形式，可直接沿正交于时间轴的方向快速投影。直觉上，这种解耦将"高斯在空间中长什么样"和"高斯随时间如何运动"分为两个独立问题，从而避免了 4D 矩阵运算的冗余。在此基础上，Disentangled4DGS 提出"先投影再切片"（Projection-First Splatting）流程：先在 3D 空间完成协方差投影和均值变换，再通过时间偏移 $dt = t_0 - t$ 计算当前时刻的射线空间位置，避免了传统方法中 4D 矩阵运算的冗余。

此外，引入流梯度引导一致性损失 $L_{\text{fg}} = \lambda_{\text{flow}} \cdot \frac{1}{N}\sum_{x,y} \|\nabla M(x,y)\| \cdot (1 - \|\nabla I(x,y)\|)$，惩罚光流梯度与图像边缘不对齐的区域以改善运动连贯性；并将分裂策略解耦为时间分裂和空间分裂两部分独立执行，以更好地建模突然出现/消失等非线性动态。在 RTX 3090 上以 1352$\times$1014 分辨率达到 343 FPS，存储需求减少至少 4.5\%。

\paragraph{时序敏感剪枝与运动分组}

DeformableGS 使用神经变形场预测每个高斯的运动，精度高但推理开销大——每帧需为所有 $N$ 个高斯逐一调用变形网络。SpeeDe3DGS（2025.06）~\cite{2506_07917} 在此基础上引入三个互补模块。首先是时序敏感剪枝（TSP）：不同于静态 3DGS 中仅考虑空间贡献的剪枝，TSP 定义时间敏感性剪枝分数 $\tilde{U}_{\mathcal{G}_i} \approx \sum_{\phi,t} (\nabla_{g_i} I_{\mathcal{G}_t}(\phi))^2$，即在所有训练视角和时间步上累积高斯投影贡献 $g_i$ 对 $L_2$ 重建损失的二阶敏感度，聚合了时空维度的冗余信息以识别运动依赖的冗余基元。

其次是时序敏感采样（TSS）：TSP 仅在训练时间步上评估梯度，可能遗漏在训练帧表现稳定但在未观测时间戳漂移的"浮点伪影"（floaters）。TSS 在变形函数的时间输入上注入线性退火的高斯扰动 $\mathcal{X}(i) = \mathcal{N}(0,1) \cdot \beta \cdot \Delta t \cdot (1 - i/\tau)$，在训练早期以较大扰动暴露时间不稳定的高斯（使其获得低分并被剪枝），后期退火至零以专注精确重建。

最后是 GroupFlow：将学习到的神经变形场蒸馏为分组共享的 SE(3) 刚体变换（即三维旋转加平移的六自由度变换），通过最远点采样选取 $J$ 个控制点，将轨迹最相似的高斯分配到同一组，每组用 Umeyama 对齐（一种基于奇异值分解的闭式刚体配准算法）估计刚体变换 $[R_j^t|T_j^t]$，将推理开销从 $N$ 降为 $J$。组内高斯的位置更新为 $\mu_i^t = R_j^t(\mu_i^0 - h_j^0) + h_j^0 + T_j^t$，旋转通过四元数乘法同步更新。在 50 个动态场景上，TSP+TSS 实现平均 6.78 倍渲染加速且使用 10 倍更少基元；加入 GroupFlow 后进一步提升至 13.71 倍渲染加速和 2.53 倍训练加速。

\paragraph{结构级联 4D 优化}

先前的在线动态重建方法（如 Dynamic3DGS）为每个高斯独立预测变形，忽略了现实世界中物体的层次结构——同一刚体部件的高斯共享相似运动。SCas4D（2025.10）~\cite{2510_06694} 的核心思想是利用这种结构先验，设计从粗到细的级联变形函数——类似于多分辨率图像金字塔，最粗层捕获整体运动趋势，最细层精修局部细节。具体而言，SCas4D 首先对高斯质心做 KMeans 聚类得到最细层（1280 个簇），再通过层次聚类合并为中间层（320 个簇）和最粗层（64 个簇）。在每一层，同一个簇内的高斯共享参数化的变形函数：

\begin{equation}
\mathbf{x}_d = (R_j(\mathbf{x} - \mathbf{p}_j^c) + \mathbf{t}_j) \cdot (\tanh(\mathbf{c}_j^\top(\mathbf{x} - \mathbf{p}_j^c) + s_j) + 1)
\end{equation}

其中 $R_j$ 为旋转矩阵（以四元数存储）、$\mathbf{t}_j$ 为平移向量、$\tanh$ 项控制缩放因子在 $(0,2)$ 范围内以防止训练中出现 NaN。多层级联组合为 $\mathbf{p}_t = \phi_{K,j_K}(\cdots\phi_{2,j_2}(\phi_{1,j_1}(\mathbf{p}_{t-1}))\cdots)$（实现中使用 $K=3$），最粗层快速逼近整体变形，最细层在此基础上精修局部细节。此外，每个高斯还学习额外的精修参数 $\Delta\mathbf{p}$、$\Delta R$、$\Delta\mathbf{s}$ 以补充簇级变形无法捕捉的个体差异。训练中引入尺度损失 $L^{\text{scale}} = \frac{1}{N}\sum_i \max(0, \text{scale}_{i,t} - \text{max\_scale})$ 防止变形过程中产生过大或过度拉伸的高斯。SCas4D 在每个时间帧仅 100 次迭代内即可收敛，以二十分之一的训练迭代达到与现有方法可比的结果，同时展示了在自监督铰接物体分割、新视角合成和密集点跟踪任务中的有效性。

\subsubsection{压缩与编码}

本路线介绍围绕帧间预测、熵编码和率失真优化的4D高斯压缩与编码方法。

上文通过剪枝和解耦提升了表示效率，但仍需解决 4D 高斯数据本身的存储与传输问题。本节围绕帧间预测、熵编码和率失真优化，介绍六种面向 4D 高斯的压缩方法。

\paragraph{率感知 4D 高斯压缩}

直接逐帧独立存储高斯基元会忽略帧间相似性，造成巨大的时间冗余。更深层的问题在于，现有的动态 3DGS 方法通常将表示构建和压缩编码作为两个独立步骤先后执行——先训练高斯表示，再对训练结果做压缩——这种分离式流程在训练阶段完全不感知最终码率，导致学到的表示天然具有高熵，压缩潜力受限。4DGC（2025.03）~\cite{2503_18421} 面向自由视点视频（Free-Viewpoint Video, FVV）的流式传输场景，核心思路是将表示与压缩端到端联合优化，使模型在训练时就感知码率代价，从而学出低熵、高压缩友好的 4D 高斯表示。

\textbf{运动感知动态表示。} 直觉上，4DGC 的帧间建模与视频编码中的 I帧/P帧 思路一脉相承：以第一帧的完整 3DGS 表示作为关键帧（I 帧）$\mathbf{G}_1$，后续每帧（P 帧）仅存储紧凑的运动网格 $\mathbf{M}_t$ 和稀疏补偿高斯 $\Delta\mathbf{G}_t$。运动估计阶段使用 $L$ 级（默认 $L=3$）多分辨率运动网格 $\mathbf{M}_t = \{\mathbf{M}_t^l\}_{l=1}^L$，以捕获不同速度和方向的多尺度运动。对于前一帧的每个高斯 $\boldsymbol{\mathcal{G}}_{t-1}$，先对其中心位置 $\boldsymbol{\mu}_{t-1}$ 做多频位置编码 $\mathbf{P}_{t-1}^l = \{\sin(2^l\pi\boldsymbol{\mu}_{t-1}), \cos(2^l\pi\boldsymbol{\mu}_{t-1})\}_{l=1}^L$，在各级网格上做三线性插值后拼接多尺度特征，再分别送入两个共享轻量 MLP $\Phi_{\boldsymbol{\mu}}$ 和 $\Phi_{\mathbf{R}}$ 预测平移增量 $\Delta\boldsymbol{\mu}_t$ 和旋转增量 $\Delta\mathbf{R}_t$：

\begin{equation}
  \Delta\boldsymbol{\mu}_t = \Phi_{\boldsymbol{\mu}}\!\Bigl(\bigoplus_{l=1}^L \text{Interp}(\mathbf{P}_{t-1}^l, \mathbf{M}_t^l)\Bigr), \quad \Delta\mathbf{R}_t = \Phi_{\mathbf{R}}\!\Bigl(\bigoplus_{l=1}^L \text{Interp}(\mathbf{P}_{t-1}^l, \mathbf{M}_t^l)\Bigr),
\end{equation}

其中 $\bigoplus$ 表示拼接操作。变换后的当前帧高斯为 $\mathbf{G}_t' = \hat{\mathbf{G}}_{t-1}(\boldsymbol{\mu}+\Delta\boldsymbol{\mu}_t,\; \Delta\mathbf{R}_t \cdot R_{t-1},\; C)$，其中 $C$ 表示固定不变的属性（球谐系数、尺度、不透明度）。这种设计的关键优势在于：运动网格仅需编码帧间的运动参数，而非重新存储每帧的完整高斯属性。

\textbf{高斯补偿。} 刚体变换只能描述已出现物体的平移和旋转，无法处理两类特殊情况：（1）新出现的物体或场景边缘，表现为梯度显著变化；（2）快速变换中的大型高斯基元，刚体近似不足以捕捉其多视角差异。针对第一类，当某位置梯度超过阈值 $\tau_g$ 时，克隆该处高斯基元作为 $\Delta\mathbf{G}_t^g$；针对第二类，当某高斯的平移幅度 $|\Delta\boldsymbol{\mu}_t|$ 或旋转幅度 $|\Delta\mathbf{R}_t|$ 分别超过阈值 $\tau_\mu$、$\tau_R$ 时，从中分裂出两个尺度缩小为原来的 $\frac{1}{100}$ 的新高斯，以更精细地捕捉局部动态。新补偿高斯的中心以 $\boldsymbol{\mu}_t^m \sim \mathcal{N}(\boldsymbol{\mu}, 2\boldsymbol{\Sigma})$ 采样，在第二训练阶段优化。当前帧的完整重建为 $\hat{\mathbf{G}}_t = \mathbf{G}_t' + \Delta\mathbf{G}_t$，并存入参考缓冲区以供下一帧的运动估计使用。

\textbf{端到端率感知压缩。} 4DGC 的压缩方案将表示学习和熵编码联合优化，核心创新在于可微分量化与隐式熵模型的结合。由于标准量化（取整操作）不可微分，4DGC 采用模拟量化策略：在量化值上叠加均匀噪声 $u \sim U(-\frac{1}{2q},\frac{1}{2q})$，其中 $q$ 为量化步长，以近似量化效果并保持梯度回传。码率估计采用一个在线学习的隐式熵模型（区别于图像压缩中基于大规模数据集离线训练的熵模型），通过可微分地估计量化值的概率质量函数来计算码率：$P_{\text{PMF}}(\hat{y}) = P_{\text{CDF}}(\hat{y}+\frac{1}{2}) - P_{\text{CDF}}(\hat{y}-\frac{1}{2})$。训练分两阶段进行：第一阶段（400 次迭代）联合优化运动网格 $\mathbf{M}_t$ 及其熵模型，损失函数为
\begin{align}
  \mathcal{L}_{s1} &= \mathcal{L}_{\text{color}} + \lambda_1 \mathcal{L}_{\text{rate}}^{\text{ME}}, \\
  \mathcal{L}_{\text{color}} &= (1-\lambda_2)\|\mathbf{c}_{\text{gt}} - \hat{\mathbf{c}}\|_1 + \lambda_2 \mathcal{L}_{\text{D-SSIM}}, \quad
  \mathcal{L}_{\text{rate}}^{\text{ME}} = -\frac{1}{N}\sum_{\hat{y} \in \hat{\mathbf{M}}_t} \log_2 P_{\text{PMF}}^1(\hat{y}),
\end{align}

其中 $\lambda_1$ 控制率失真权衡，$\lambda_2$ 平衡 $L_1$ 重建损失与 D-SSIM 感知损失；第二阶段（100 次迭代）固定已训练的运动网格，联合优化补偿高斯的球谐系数（SH 系数占单基元存储的绝大部分）及其熵模型：$\mathcal{L}_{s2} = \mathcal{L}_{\text{color}} + \lambda_1 \mathcal{L}_{\text{rate}}^{\text{MC}}$。通过调节 $\lambda_1 \in \{0.0003, 0.0001, 0.00005, 0.00001\}$ 实现可变比特率，支持从高保真到极低码率的灵活配置。

\textbf{实验表现。} 在 N3DV、MeetRoom 和 Google Immersive 三个真实世界数据集上，4DGC 一致取得最优率失真性能。以 N3DV 数据集为例，在达到与 3DGStream 可比重建质量的前提下，4DGC 仅需 0.5 MB/帧（3DGStream 需 8.1 MB/帧），实现约 16 倍压缩。相对 TeTriRF，4DGC 在三个数据集上的 BDBR（Bjontegaard Delta Bit-Rate，率失真曲线下的码率差异指标）分别降低 68.59\%、40.71\% 和 59.99\%。在计算效率上，4DGC 每帧训练仅需 0.83 分钟（ReRF 需 42.73 分钟、TeTriRF 需 1.04 分钟），渲染仅需 0.006 秒/帧，编码和解码分别为 0.72 秒和 0.09 秒，远快于同类方法，适合 FVV 的实时流式传输场景。消融实验验证了运动估计、高斯补偿和端到端联合优化三个模块各自的必要性：移除任一组件均导致率失真性能下降。

\paragraph{轻量紧凑 4D 高斯生成}

Light4GS（2025.03）~\cite{2503_13948} 针对可变形 4DGS 中高维潜嵌入和大量基元导致的存储膨胀问题，提出基于显著性剪枝与深度上下文模型的轻量压缩框架。

\textbf{Baseline 与改进。} vanilla 4DGS 使用六面体平面（hexplane）——一种将 4D 时空特征分解为六个 2D 平面（xy、xz、yz、xt、yt、zt）的紧凑编码结构——来编码时空变化，但存在两个缺陷：六面体网格通常在 SfM 点云包围盒内初始化，与无界 3D 空间不匹配；标准坐标轴不一定对齐场景主方向，非主成分投影阻碍潜特征学习。Light4GS 引入 PCA 引导的六面体查询：先用非线性空间收缩将高斯中心映射到有界空间，再对空间坐标做 PCA 旋转对齐主方向，然后用变换后坐标进行多尺度六面体双线性插值。

\textbf{时空显著性剪枝（STP）。} vanilla 4DGS 仅通过不透明度剪枝低贡献基元，忽略了基元在不同帧和视角下的全局渲染影响。直觉上，STP 要回答的问题是：如果移除某个高斯，所有帧和所有视角的渲染质量会受到多大影响？影响越小，该高斯越冗余。STP 定义第 $j$ 个可变形基元的全局显著性得分：

\begin{equation}
  \Sigma_j = \sum_{t=1}^{T}\sum_{i=1}^{MHW} \mathbf{1}(G_j(f_h, t),\, r_{i,t})\; T_{j,i,t}\; \alpha_j\; V(\sigma_{j,t}),
\end{equation}

即对所有时间戳和视角，累加该基元与视线相交时的透射率加权不透明度与归一化体积贡献。训练中周期性按 $\Sigma_j$ 排名并按预设比例剪除得分最低的基元，消除超过 64\% 的可变形基元。剩余基元在优化中自动补偿质量损失。

\textbf{多尺度六面体上下文压缩（MHCM）。} 将帧内帧间预测与超先验集成为由粗到细的上下文结构：最低尺度六面体采用棋盘格建模，锚点用超先验编码，非锚点结合超先验与棋盘格上下文；时空平面解码后的时间均值提供跨平面上下文，替代额外超先验存储；高尺度六面体利用低尺度已解码特征的双线性插值作为跨尺度上下文。自适应量化为不同平面分配可学习量化步长 $q_c$，总熵损失为：

\begin{equation}
  R_{\text{Hex}} = -\sum_{l}\sum_{c} \log_2 N\!\left(R_c^l \ast U\!\left(-\tfrac{1}{2},\tfrac{1}{2}\right) q_c\right).
\end{equation}

每个六面体特征被建模为以深度上下文为条件的高斯分布，分布参数由 MLP 预测，最终用算术编码（Arithmetic Coding，一种接近熵极限的无损压缩方法）进行熵编码。对占单基元存储 80\% 以上的球谐系数，采用全因子化熵模型直接压缩。整体实现超过 120 倍压缩，渲染 FPS 提高 20\%。

\paragraph{预测式 4D 高斯压缩}

P-4DGS（2025.10）~\cite{2510_10030} 受视频压缩中帧内/帧间预测范式的启发，将动态 3D 高斯表示类比于视频编码结构——规范空间相当于 I 帧（关键帧），变形场相当于运动矢量：规范空间充当参考帧，变形场类比运动矢量。

\textbf{Baseline 与改进。} 现有 4DGS 框架（如 D3DGS）使用可变形 MLP 学习每个规范高斯的变形，但忽略了高斯基元间的空间相关性与时间冗余。P-4DGS 设计了基于 3D 锚点的时空预测模块：空间预测借鉴 Scaffold-GS 的锚点结构（Anchor-based structure，即以稀疏体素网格中的锚点为根节点，每个锚点关联若干神经高斯），每个锚点 $x_a$ 生成 $k$ 个邻近高斯基元，其位置由可学习偏移与缩放计算：

\begin{equation}
  \{x_0, \ldots, x_{k-1}\} = x_a + \{O_0, \ldots, O_{k-1}\} \odot l_a,
\end{equation}

其余属性（不透明度、颜色、旋转、尺度）由共享 MLP 根据锚点特征 $f_a$、相对距离 $\delta_{vc}$ 和视角方向 $d_{vc}$ 预测。时间预测沿袭规范空间+变形场范式，通过位置编码和变形 MLP 预测位移：

\begin{equation}
  (\Delta x, \Delta s, \Delta r) = \Phi_d\!\left(\text{concat}(E(x), E(t))\right),
\end{equation}

将变形向量加到规范属性上得到时刻 $t$ 的高斯属性。

\textbf{自适应量化与上下文熵编码。} 训练后对锚点属性做标量量化，量化步长 $q$ 根据锚点位置的哈希特征自适应调整：$q = Q_0(1 + \text{sigmoid}(\Phi_q(h)))$，其中 $Q_0$ 为基础步长，$h$ 由包含一个 3D 网格和三个 2D 网格的二值哈希网格查询。量化后的属性假设服从高斯先验，其分布参数 $(\mu, \sigma)$ 由 MLP 基于哈希特征预测，再通过算术编码压缩。总训练损失为率失真联合优化：$L_{\text{total}} = \lambda_{\text{rate}} L_{\text{rate}} + L_{\text{render}}$。在合成和真实数据集上实现约 1MB 平均存储，分别达到 40 倍和 90 倍压缩。

\paragraph{优化最小 4D 高斯}

OMG4（2025.10）~\cite{2510_03857} 基于 Real-Time4DGS 基线（百万级 4D 高斯、存储超 1GB），提出渐进式三阶段优化管道将模型压缩至约 3MB。

\textbf{阶段一：高斯采样。} 先预训练 Real-Time4DGS 模型，然后提出静动态联合评分（SD-Score）量化每个高斯的重要性：

\begin{equation}
  \text{SD}^{(i)} = S_{\text{grad}}^{(i)} \cdot T_{\text{grad}}^{(i)}, \quad S_{\text{grad}}^{(i)} = \sum_{j=1}^{N} \nabla_{u_{i,j}} L_j \big\|_2, \;\; T_{\text{grad}}^{(i)} = \sum_{j=1}^{N} \nabla_{t_i} L_j,
\end{equation}

其中 $S_{\text{grad}}$ 累积视图空间梯度（静态区域高斯稀疏，位置扰动影响大，得分高），$T_{\text{grad}}$ 累积时间梯度（动态区域对时间变化敏感，得分高）。两者互补，仅保留约 20\% 的采样高斯即可保持接近无损的重建质量。

\textbf{阶段二：高斯剪枝。}
采样阶段后仍有冗余高斯。对 $S_{\text{grad}}$ 和 $T_{\text{grad}}$ 双阈值过滤：阈值为各自分布的 $p$-分位数，移除两项得分均低的高斯。两阶段之间插入优化步骤，使剩余高斯重新调整以补偿剪枝损失。

\textbf{阶段三：高斯合并。}
前两阶段逐个评估高斯，未考虑高斯间的相似性。定义相似性得分 $S(G_i, G_j) = -\|x_i - x_j\|_2^2 + \beta\|f_i - f_j\|_2^2$（空间邻近+外观相似），在 4D 时空网格内聚类，将同簇高斯融合为加权代理高斯。合并过程迭代 $M$ 次，逐步增大网格尺寸。

\textbf{属性压缩。} 集成隐式外观压缩（小 MLP 根据空间特征和时间坐标生成颜色与不透明度），并将子向量量化（SVQ）推广到 4D 属性：将旋转四元数和时间轴尺度等属性向量分割为子向量，每个子向量用小型码本量化。采用分阶段 SVQ 方案——先仅对 3D 属性量化并优化，再激活 4D 属性的量化——避免同时量化导致的优化不稳定。最终经 Huffman 编码和 LZMA 压缩，模型大小减少超过 60\%。

\paragraph{时间压缩 3D 高斯}

Temporally Compressed 3D Gaussian Splatting（2024.12）~\cite{2412_05700} 以 Dynamic 3DGS（Luiten et al.）为基线，该基线逐帧独立优化高斯位置和旋转，参数量随帧数线性增长，存储代价高昂。TC3DGS 从三个维度压缩动态表示。

\textbf{时间一致性掩码剪枝。} 核心直觉是：一个真正无用的高斯应该在所有帧都无用，而非仅在某些帧被淘汰、在另一些帧又被保留。静态 3DGS 的掩码剪枝方法在动态场景中失效，因为高斯在运动过程中相对于其他高斯的位置不断变化，其透射率 $T_i$ 也随之改变，导致同一高斯在不同帧的掩码值剧烈波动。TC3DGS 引入时间相关掩码 $m_{n,t}$ 并施加一致性正则：

\begin{equation}
  \mathcal{L}_{\text{mc}} = \sum_{n=1}^{N} |m_{n,t} - \text{sg}(m_{n,t-1})|,
\end{equation}

其中 $\text{sg}(\cdot)$ 为停止梯度操作。这迫使掩码在帧间保持稳定，从而可靠地识别全局不重要的高斯。最终对所有帧的掩码取平均，低于阈值则剪除。

\textbf{梯度感知混合精度量化。} 不同参数对重建质量的敏感度差异巨大。定义每个参数 $\theta$ 的敏感度系数为：

\begin{equation}
  S(\theta) = \frac{1}{\sum_{k=1}^{K} N_k} \sum_{k=1}^{K} |\nabla_\theta Q_k|,
\end{equation}

即所有训练图像像素强度对 $\theta$ 的平均梯度幅度。敏感度高的参数分配更高比特精度，低的则压缩至 4-bit。使用可学习缩放因子的参数量化技术，直接嵌入优化过程而非后处理。

\textbf{关键点轨迹插值。} 受 Ramer-Douglas-Peucker 算法启发，提出自适应关键点选择：仅保留轨迹变化显著的时间点作为关键点，帧间线性插值恢复完整轨迹。支持同时控制误差容限 $\epsilon$ 和最大关键点数 $max\_kp$，对静态背景高斯仅需一个关键点。实现高达 67 倍压缩且视觉质量几乎无损。

\paragraph{时间平滑感知率失真优化}

Temporal Smoothness-Aware Rate-Distortion Optimized 4DGS（2025.07）~\cite{2507_17336} 基于 Ex4DGS 基线，提出首个面向动态 4D 高斯表示的比特级率失真优化压缩框架。

\textbf{Baseline 与动机。} Ex4DGS 将场景分解为静态和动态高斯，动态高斯仅在关键帧存储参数，中间时刻通过插值获得。但关键帧处的位置仍独立存储，未能利用时间冗余。核心观察是：每个动态点沿平滑轨迹运动，不需要记录每一帧的独立坐标。

\textbf{小波轨迹压缩。} 对动态高斯的位置轨迹 $\mathbf{p}_d = [p_1, p_2, \ldots, p_T]^T$（$p_t \in \mathbb{R}^3$），沿时间维度施加单层 Haar 小波变换（一种将信号分解为低频近似和高频细节的线性变换）：

\begin{equation}
  F = W \mathbf{p}_d, \quad F = F_a \sqcup F_d, \quad \tilde{\mathbf{p}}_d = W^T(F_a \sqcup \mathbf{0}),
\end{equation}

其中 $W$ 为 Haar 小波变换矩阵，将轨迹分解为低频近似系数 $F_a$ 和高频细节系数 $F_d$。关键设计是直接丢弃 $F_d$（置零），仅保留 $F_a$ 并通过逆变换重建轨迹。这利用了运动轨迹通常平滑、高频细节对感知质量影响有限的先验，显著减少存储。

\textbf{分层压缩策略。} 静态和动态高斯的公共属性（球谐系数、旋转、尺度）采用熵约束向量量化（ECVQ，一种在量化码率与失真之间寻找最优平衡的编码方法）压缩。学习型掩码剪枝冗余高斯和球谐系数，掩码通过直通估计器（STE）二值化。对动态不透明度的时间参数（中心和方差），精细分析量化敏感度：方差参数对量化伪影特别敏感，因此避免激进量化。总损失为端到端率失真联合优化：

\begin{equation}
  L_{\text{total}} = L_{\text{dist}} + \lambda_R L_{\text{rate}} + \lambda_{\text{reg}} L_{\text{reg}},
\end{equation}

其中 $L_{\text{dist}}$ 结合 L1 损失和 D-SSIM，$L_{\text{rate}}$ 聚合 ECVQ 码率、剪枝稀疏性和小波压缩的存储开销。在 Ex4DGS 基线上实现高达 91 倍压缩，模型大小降至原始的约 1.1\%。

\subsubsection{流式传输与在线重建}

离线压缩解决了存储问题，但面向实时应用还需在线逐帧重建与流式传输的能力。本节介绍十种流式重建方法，核心矛盾是：如何在极低延迟内为每帧生成紧凑的 4D 表示，同时避免逐帧优化带来的误差累积。

\paragraph{量化高效编码流式自由视角视频}

QUEEN（2024.12）~\cite{2412_04469} 以 3DGStream（Sun et al.）为基线，解决在线流式自由视角视频的三个矛盾需求：低延迟编码、低内存带宽、高重建质量。核心改动是：移除所有结构约束，对所有高斯属性自由学习帧间残差，并设计量化-稀疏框架将残差压缩到极低存储，使模型大小随场景动态程度自适应变化。

\textbf{Baseline 与改进。} 3DGStream 使用 Instant-NGP 哈希编码结构缓存位置和旋转残差，但存在两个关键限制。第一，仅更新部分属性（位置和旋转），其余属性（不透明度、球谐系数、尺度等）保持不变，模型表达力受限。第二，施加固定的结构约束——每帧使用相同大小的 Instant-NGP 编码，无法根据场景动态程度调整存储开销。在静态场景中浪费带宽，在剧烈运动场景中又容量不足。QUEEN 的核心洞察是：不施加任何结构约束，对所有高斯属性自由学习帧间残差，同时通过量化-稀疏框架按需压缩残差，使存储开销与实际场景动态成正比。这样，变化小的帧自动获得更小的模型尺寸，变化大的帧获得更多比特。

\textbf{属性残差量化。}
直觉上，如果场景在相邻两帧间变化很小（如人物静止站立），那么绝大多数高斯的属性几乎不变，残差应集中在零附近。量化正好利用这一分布特性：将连续残差映射到少量离散整数值，大幅减少存储。具体而言，对每个时间步 $t$，QUEEN 学习全部属性残差：$\mathcal{A}_t = \mathcal{A}_{t-1} + \mathcal{R}_t$。对旋转、尺度、不透明度、球谐系数等属性残差，维护每个高斯的量化整数潜码 $\mathbf{l}_i \in \mathbb{Z}^L$，通过共享线性解码器 $D$ 恢复残差：

\begin{equation}
  \mathbf{l}_i = \text{STE}(\hat{\mathbf{l}}_i), \quad \mathbf{r}_i = D(\mathbf{l}_i; \mathbf{D}) = \mathbf{D} \cdot \text{float}(\mathbf{l}_i),
\end{equation}

其中 $\hat{\mathbf{l}}_i$ 为连续潜码，STE（Straight-Through Estimator，直通估计器）在前向传递中将连续值取整为整数（$\text{STE}(\cdot) = \text{round}(\cdot)$），在反向传递中直接透传梯度（$\frac{\partial \text{STE}}{\partial \hat{\mathbf{l}}_i} \approx 1$），使不可微的取整操作变得可微。$L$ 为潜码维度（实验中旋转取 $L=4$，尺度取 $L=3$，球谐系数取更高维度），$\mathbf{D} \in \mathbb{R}^{d_{\text{attr}} \times L}$ 为可学习的解码矩阵。四类属性（旋转 $r$、尺度 $s$、不透明度 $o$、球谐系数 $h$）分别维护独立的解码器参数。

训练后对量化整数潜码做熵编码（如算术编码），将离散整数值压缩为分数比特。实验表明，残差分布通常接近单峰拉普拉斯分布或高斯分布，熵编码可高效利用这一统计特性。例如在 N3DV 数据集的 Sear Steak 场景中，所有量化属性平均仅需 0.68 bits/元素，对应约 0.5 MB/帧。量化加熵编码相比原始浮点残差实现约 10 倍压缩。

\textbf{位置残差门控。}
位置属性对量化高度敏感——论文消融实验表明，对位置残差做量化会导致 PSNR 下降超过 0.2 dB，且增加训练轮次也无法完全弥补这一损失。直觉上，位置偏移即使很小也直接改变像素着色，量化舍入误差被放大到图像空间。因此位置残差不适合量化，改用学习型稀疏化策略：大部分高斯在相邻帧间位置不变（静态部分），只有少数高斯需要更新位置（动态部分）。具体做法是 $\Delta \mathbf{p}_i = g_i \cdot \mathbf{l}_{p_i}$，其中 $g_i \in [0, 1]$ 为 Hard Concrete 门控变量（一种可微的稀疏化分布，允许精确的 $L_0$ 正则化），$\mathbf{l}_{p_i}$ 为全精度位置残差。

通过拉伸 Concrete 分布到包含端点 $\{0, 1\}$ 的区间，使门控在训练结束后可以精确取到零值（而非近似零）：

\begin{equation}
  g_i = \min(1, \max(0, \hat{g}_i \cdot (\gamma_1 - \gamma_0) + \gamma_0)), \quad \hat{g}_i = \text{Sigmoid}\!\left(\frac{\log \alpha_i}{\tau}\right),
\end{equation}

其中 $\alpha_i > 0$ 为可学习参数，$\tau$ 为温度超参（控制分布的锐度），$\gamma_0 < 0$ 和 $\gamma_1 > 1$ 为拉伸端点。例如在 N3DV 数据集上设置 $\gamma_0 = -0.1$、$\gamma_1 = 1.1$（对静态场景强调稀疏性），在 Immersive 数据集上设置 $\gamma_0 = -0.2$、$\gamma_1 = 1.1$（允许更多动态高斯）。

配合 $L_0$ 正则化损失 $\lambda_{\text{reg}} \cdot \sum_i \sigma(\log \alpha_i - \tau \log \gamma_0 / (\gamma_1 - \gamma_0))$ 鼓励门控为零。训练后以 COO（坐标格式）稀疏矩阵存储非零位置残差，仅保留门控值为正的高斯索引及其对应的全精度残差向量。由于大多数高斯在静态场景中门控为零，实际存储的位置残差极为紧凑。

\textbf{视图空间梯度差分离动静内容。}
量化-稀疏框架的核心优势是自适应性：帧大小随场景动态程度自动调整。但如何快速识别哪些区域发生了变化？QUEEN 利用相邻帧间视图空间梯度差 $\mathbf{d}_t$ 作为动静分离信号：

\begin{equation}
  \mathbf{d}_t = \frac{1}{V}\sum_{v=0}^{V-1}\left[\frac{\partial L_t^{(v)}}{\partial \mathbf{p}_{t-1}^{\prime(v)}} - \frac{\partial L_{t-1}^{(v)}}{\partial \mathbf{p}_{t-1}^{\prime(v)}}\right],
\end{equation}

其中 $\mathbf{p}_{t-1}^{\prime(v)}$ 为第 $v$ 个视角下高斯投影到屏幕空间的位置，$L_t^{(v)}$ 和 $L_{t-1}^{(v)}$ 分别为当前帧和前一帧的渲染损失。梯度差范数大的区域意味着该高斯在新帧中的误差信号发生了显著变化——即场景内容在该区域发生了运动或光照变化。

这一梯度差信号被用于两个目的。第一，用梯度差范数初始化门控参数 $\alpha_i$（而非随机初始化），使动态高斯的门控偏向开启、静态高斯的门控偏向关闭，显著加速收敛。第二，根据梯度差范数识别动态图像区域，构建掩码进行局部渲染——仅在动态区域渲染和计算损失，跳过静态区域。掩码通过 $48 \times 48$ 核膨胀以覆盖边缘区域的大高斯影响范围。N3DV 数据集上掩码训练占总迭代 30\%，Immersive 数据集上占 65\%（因运动更剧烈）。

\textbf{单目深度初始化。}
COLMAP 生成的初始点云在纹理缺失区域或视角覆盖不足区域往往稀疏或不完整，影响后续残差学习的准确性。QUEEN 利用单目深度估计网络（如 MiDaS）预测稠密深度图，通过最小二乘拟合将其与 COLMAP 稀疏点云对齐（解决尺度-偏移歧义），然后仅在 COLMAP 点云缺失的区域补充新点。实验表明，在视角覆盖不足的区域，这一初始化可将该视角的 PSNR 提升高达 3 dB，且几乎不增加训练时间。

\textbf{实验表现。} 在 N3DV 和 Immersive 两个基准上，QUEEN 在所有指标上超越先前的在线 FVV 方法（3DGStream、TeTriRF 等）。在 N3DV 数据集上，模型平均大小仅 0.6 MB/帧（3DGStream 为 7.8 MB/帧，约 13 倍压缩），训练约 5 秒/帧（3DGStream 约 10 秒/帧），渲染速度 350+ FPS。在 Immersive 数据集的室内外混合场景上，模型大小约 1.49 MB/帧。帧大小随场景动态自适应变化——高动态帧（如人物出入画面）自动分配更大存储，静态帧则极小。用户研究中，76.67\% 的参与者更偏好 QUEEN 而非 3DGStream（N3DV），97.14\% 更偏好 QUEEN（Immersive）。

\paragraph{即时高斯流式传输}

Instant Gaussian Stream（2025.03）~\cite{2503_16979} 瞄准流式重建的核心瓶颈——逐帧优化导致的高延迟和误差累积，提出首个可泛化的流式框架。其核心思想是：与其对每帧都做耗时的逐点优化，不如训练一个泛化网络，单次前向推理即可估计全部高斯的帧间运动。

\textbf{锚点驱动高斯运动网络（AGM-Net）。}
Baseline 做法（如 3DGStream）需要为每帧单独训练一个 InstantNGP 来查询每个高斯的位移，导致每帧训练耗时 10 秒以上。IGS 将问题重新建模为"锚点驱动的运动估计"：先用最远点采样（FPS）从 $N$ 个高斯中均匀抽取 $M=8192$ 个锚点 $\mathcal{C}=\text{FPS}(\{\boldsymbol{\mu}_i\}_{i=1}^N)$，再通过投影感知机制将多视角 2D 光流特征提升到 3D 空间：
$$\mathbf{f}_i=\frac{1}{V}\sum_{j=1}^V \text{Bilinear}\bigl(\pi_j(\mathcal{C}_i),\,F_j\bigr),$$
其中 $\pi_j(\mathcal{C}_i)$ 是锚点在第 $j$ 个视角的投影坐标，$F_j$ 是该视角的 2D 运动特征图。锚点的 3D 运动特征经过 Transformer 自注意力进一步聚合：$\mathbf{z}_i=\text{Transformer}(\{\mathbf{f}_i\}_{i=1}^M)$。之后，每个高斯的运动通过对 $K$ 近邻锚点的特征做距离加权插值得到：
$$\hat{\mathbf{z}}_i=\frac{\sum_{k\in\mathcal{N}(i)}e^{-d_k}\mathbf{z}_k}{\sum_{k\in\mathcal{N}(i)}e^{-d_k}},\quad d_k=\|\boldsymbol{\mu}_i-\mathcal{C}_k\|_2,$$
最后通过线性层解码出位移和旋转增量 $\Delta\boldsymbol{\mu}_i,\,\Delta\text{rot}_i=\text{Linear}(\hat{\mathbf{z}}_i)$，更新高斯属性：$\boldsymbol{\mu}'_i=\boldsymbol{\mu}_i+\Delta\boldsymbol{\mu}_i$，$\text{rot}'_i=\text{norm}(\text{rot}_i)\cdot\text{norm}(\Delta\text{rot}_i)$。这样，8192 个锚点的运动特征即可驱动数十万个高斯的运动，单次推理即完成帧间变换。

\textbf{关键帧引导流式策略。}
仅靠 AGM-Net 的推理只调整了位置和旋转，无法处理非刚性形变和新物体出现，长期使用会导致误差累积。IGS 引入关键帧机制：每隔 $w$ 帧设一个关键帧 $K_{iw}$，在关键帧处对变形后的高斯做有界精化（max-points bounded refinement），精化后高斯数量不会超过上限，避免稀疏视角下的过拟合。非关键帧则直接用 AGM-Net 前向推理变形，无需优化。这一设计使 IGS 在保证时序复杂场景重建质量的同时，将平均每帧重建时间降至约 2.7 秒，且具备跨域泛化能力。

\paragraph{层次一致运动流式重建}

HiCoM（2024.11）~\cite{2411_07541} 发现流式动态场景重建面临两个基础问题：一是动态场景通常由有限数量的相机拍摄，导致初始 3DGS 容易过拟合且高斯数量失控；二是已有方法（如 3DGStream）使用隐式神经网络预测运动场，与 3DGS 显式离散的本质不匹配，学习效率和存储开销都不理想。

\textbf{扰动平滑初始化。}
在标准 3DGS 训练中，稀疏视角往往使高斯过拟合到特定训练视角，且高斯数量可能无限增长，拖慢后续帧的学习速度。HiCoM 借鉴 dropout 思想，在训练时给每个高斯的位置添加高斯噪声扰动：$\hat{\boldsymbol{\mu}}=\boldsymbol{\mu}+\epsilon\cdot\boldsymbol{\mu}_{\text{noise}},\;\boldsymbol{\mu}_{\text{noise}}\sim\mathcal{N}(0,1)$，其中 $\epsilon$ 控制扰动幅度。随机扰动阻止模型死记硬背有限训练视角，收敛后高斯数量显著减少，为后续帧学习提供紧凑稳健的基础。

\textbf{层次一致运动机制。}
3DGStream 使用单个 InstantNGP 对每个高斯独立建模运动，既忽视了高斯间的局部一致性，又引入大量冗余参数。HiCoM 的关键观察是：优化后的 3DGS 呈现非均匀分布——复杂区域密集排布小高斯，简单区域稀疏分布大高斯，许多区域完全为空。基于此，HiCoM 将场景划分为等大小区域，每个区域 $r$ 由公式 $r=\lfloor(\boldsymbol{\mu}+0.5)/e\rfloor\cdot e$ 确定（$e$ 为区域边长），区域内所有高斯共享一组运动参数（平移 $\mathbf{t}\in\mathbb{R}^3$ 和旋转 $\mathbf{q}\in\mathbb{R}^4$）。更细的子区域再分配独立的运动参数，每个高斯的最终运动是多个层次贡献的累加：
$$\mathbf{t}_g=\sum_{l=1}^L\mathbf{t}^l,\quad \mathbf{q}_g=\sum_{l=1}^L\mathbf{q}^l,$$
其中 $L$ 是层次总数，$\mathbf{t}^l$ 和 $\mathbf{q}^l$ 分别是第 $l$ 层的运动贡献。这一从粗到细的设计既保证了局部运动一致性，又大幅减少了需要优化的参数量，使学习效率提高约 20\%、存储减少 85\%，并行训练时每帧墙钟时间降至不到 2 秒。

\textbf{持续精化。}
运动机制主要处理位置和旋转的偏移，难以捕捉所有细节变化。HiCoM 在运动学习期间监控梯度累积，对梯度超过阈值的区域克隆新高斯并进一步优化。新高斯不会丢弃而是合并到初始 3DGS 中，作为下一帧学习的基础；同时选择性移除等量低不透明度高斯以保持紧凑。

\paragraph{动态感知流式 4D 重建}

Dynamics-Aware Gaussian Splatting Streaming（DASS）（2024.11）~\cite{2411_14847} 指出流式重建的一个关键被忽视的事实：场景中绝大多数高斯的帧间偏移极小（$(0, 0.01]$），只有少量动态物体的高斯有显著位移。用同一个变形场统一处理所有高斯，既浪费计算又拖慢收敛。

\textbf{三阶段管道。}
DASS 将每个时间步的 3D 表示分为两组：数量固定的基线高斯 $\mathbf{G}_t^{\text{base}}$（从第 0 帧初始化，贯穿所有时间步）和数量可变的致密化高斯 $\mathbf{G}_t^{\text{var}}$（捕捉每帧独特的新物体）。每个时间步经历三个阶段：

（1）\textbf{选择性继承。}3DGStream 直接丢弃上一帧的致密化高斯 $\mathbf{G}_{t-1}^{\text{var}}$，浪费了已学到的场景先验。DASS 为 $\mathbf{G}_{t-1}^{\text{var}}$ 中每个高斯分配一个可学习参数 $m$，经 sigmoid 和量化后生成二值掩码：$o_r=\sigma(m)\odot o$，$s_r=\sigma(m)\odot s$，其中 $o,s$ 是原始不透明度和缩放。掩码为零的高斯贡献被消除，从而自适应保留有益高斯、丢弃冗余高斯。掩码通过保真度损失 $L_1+\lambda_{\text{D-SSIM}}L_{\text{D-SSIM}}$ 和正则项 $\lambda_{\text{inher}}\|m\|$ 联合优化。

（2）\textbf{动态感知迁移。}利用光流估计和 2D 分割结果，在变形前先估计每个高斯的动态掩码，将高斯分为动态组和静态组。两组分别使用不同容量的哈希编码 MLP：高容量网络 $H_{\text{dyn}}$ 处理动态组的复杂变形（位置偏移 $\Delta\boldsymbol{\mu}\in\mathbb{R}^3$ 和旋转偏移 $\Delta\mathbf{q}\in\mathbb{R}^4$），轻量网络 $H_{\text{st}}$ 处理静态组的微小抖动。这种双网络策略显著加速收敛。

（3）\textbf{误差引导致密化。}传统致密化仅依赖视图空间位置梯度，对新增物体和严重失真区域关注不足。DASS 额外收集 2D 误差图 $D_{\text{err}}^c=\mathbb{1}[e_i^c>\tau_{\text{err}}]$，将高误差像素区域投影回 3D 空间，精准定位需要致密化的高斯子集。这一策略加速了新增物体的优化收敛，实现了最快的即时训练速度和优异的重建质量。

\paragraph{紧凑高斯流式传输}

Compact Gaussian Streaming（2025.05）~\cite{2505_16533} 的出发点是存储效率——已有在线方法（如 QUEEN）每秒仍需传输超过 20 MB 数据，远不能满足实时带宽需求。ComGS 核心洞察有二：一是动态场景中大量区域是静态的，只需建模运动区域的属性残差；二是同一物体的多个高斯运动模式一致，可用共享表示消除冗余。

\textbf{运动敏感关键点选取。}
ComGS 沿用视图空间梯度差思路识别动态高斯。首先分别用上一帧和当前帧的真值图像计算渲染损失对高斯位置的梯度 $\mathbf{G}_{t-1}$ 和 $\mathbf{G}_t$，然后计算动态显著性分数：
$$\mathbf{s}_t=\frac{1}{V}\sum_{v=1}^V|\mathbf{G}_t^{(v)}-\mathbf{G}_{t-1}^{(v)}|,$$
选取分数最高的 $k=200$ 个高斯作为关键点 $K_t$——相比全部高斯（约 200K），仅需极少量关键点即可覆盖运动区域，且复杂运动区域自动分配更多关键点。

\textbf{自适应运动驱动机制。}
不同于 KNN 搜索的固定邻域，ComGS 为每个关键点初始化一个类高斯的空间影响场：参数化一个四元数 $\mathbf{q}_{\text{adap}}^i\in\mathbb{R}^4$ 和缩放向量 $\mathbf{s}_{\text{adap}}^i\in\mathbb{R}^3$，构建 $3\times 3$ 影响协方差。邻域高斯 $G_j$ 到关键点 $K_t^i$ 的距离为 $\mathbf{d}_{ij}=\boldsymbol{\mu}_j-\boldsymbol{\mu}_{K_t^i}$，影响权重为：
$$w_{ij}=\exp\!\bigl(-\tfrac{1}{2}\mathbf{d}_{ij}^T(\boldsymbol{\Sigma}_{\text{adap}}^i)^{-1}\mathbf{d}_{ij}\bigr),$$
若 $w_{ij}>\tau_{\text{adap}}$，则 $G_j$ 被该关键点控制。每个关键点携带可学习的平移偏移 $\Delta\boldsymbol{\mu}_K^i\in\mathbb{R}^3$ 和旋转 $\Delta\mathbf{q}_K^i\in\mathbb{R}^4$，被多个关键点控制的高斯通过加权聚合获得运动：
$$\Delta\boldsymbol{\mu}_t^j=\sum_{i\in\mathcal{I}_t^j}w_{ij}\,\Delta\boldsymbol{\mu}_K^i,\quad \Delta\mathbf{q}_t^j=\sum_{i\in\mathcal{I}_t^j}w_{ij}\,\Delta\mathbf{q}_K^i.$$
每个关键点仅需 14 个参数（3 个平移 $+$ 4 个旋转 $+$ 4 个影响四元数 $+$ 3 个影响缩放），200 个关键点共 2800 个参数即驱动整个场景运动，存储相比 3DGStream 减少超过 159 倍、相比 QUEEN 减少 14 倍。

\textbf{误差感知纠正器。}
关键点驱动的运动只支持刚性变换，对非刚性形变和新物体出现的误差会累积。ComGS 每隔 $s$ 帧设一个关键帧，仅对检测到误差的高斯进行属性残差精化：引入可学习掩码 $m_i$，经 STE 二值化后 $m_i^{\text{hard}}\in\{0,1\}$，残差以 $\boldsymbol{\theta}_i^t=\boldsymbol{\theta}_i^{t-1}+m_i^{\text{hard}}\cdot\Delta\boldsymbol{\theta}_i^t$ 形式叠加。优化后仅存储 $m_i^{\text{hard}}=1$ 的残差和掩码，配合 $\lambda_{\text{error}}=0.001$ 的稀疏正则化，最小化传输开销。

\paragraph{演化式体积视频}

EvolvingGS（2025.03）~\cite{2503_05162} 的核心洞察是：传统 GoP（Group of Pictures）分段处理导致段间闪烁，但连续帧之间共享大量几何和外观信息，大多数高斯只需微调属性值即可复用。基于此，EvolvingGS 提出一种无需 GoP 分段的演化式两阶段框架。

\textbf{变形阶段：稀疏控制点粗对齐。}
给定已精调好的 $t$ 帧高斯模型 $\mathcal{S}_t$，EvolvingGS 目标是将其变形为 $\hat{\mathcal{S}}_{t+1}$ 以粗略匹配 $t+1$ 帧。Baseline 逐点优化变形的方式在大规模高斯场景下计算代价高。EvolvingGS 引入 $M$ 个均匀采样的稀疏控制点，每个控制点 $C_j$ 定义一个局部仿射变换（位置 $\mathbf{c}_j$、四元数 $\mathbf{q}_j$、平移 $\mathbf{t}_j$）。高斯 $G_i$ 的变形通过最近 $K^g$ 个控制点的加权混合实现：
$$\hat{\boldsymbol{\mu}}_i=\sum_{j\in\mathcal{N}(G_i)}w_j(G_i)\bigl[\mathbf{q}_j(\boldsymbol{\mu}_i-\mathbf{c}_j)+\mathbf{c}_j+\mathbf{t}_j\bigr],$$
混合权重 $w_j(G_i)=(1-\|\boldsymbol{\mu}_i-\mathbf{c}_j\|/d_{\max})^2$ 并归一化为和为 1。变形阶段保持高斯数量不变，且所有外观特征（SH 系数、缩放、不透明度）固定不动，只优化位姿相关特征。为加速收敛，EvolvingGS 引入光流引导损失：
$$E_{\text{flow}}=\frac{1}{|\mathbb{1}|}\sum_{i=1}^{N_1}T_{i,q}\|\mathbf{x}_i^{2D}-(\hat{\mathbf{x}}_i^{2D}+F_q(\mathbf{x}_i^{2D}))\|,$$
其中 $F_q$ 是预测的光流图，$\mathbb{1}$ 是过滤静止背景的指示函数（光流值低于 $\tau_f$ 时置零），$T_{i,q}$ 是高斯透射率。总目标 $E=E_{\text{data}}+\lambda_1 E_{\text{flow}}+\lambda_2 E_{\text{arap}}+\lambda_3 E_{\text{rot}}$ 还包含 ARAP 刚性约束和四元数归一化约束，确保变形平滑且物理合理。

\textbf{精化阶段：双流优化与贡献裁剪。}
粗对齐无法处理新物体出现和拓扑变化。精化阶段维护两套高斯：继承自变形阶段的参考点 $G^{\text{ref}}$（外观特征固定，仅微调位姿，保证时序稳定）和新增的扩展点 $G^{\text{ext}}$（无约束优化，用于重建新物体）。由于不透明度被固定，传统基于不透明度的裁剪策略失效。EvolvingGS 改用"贡献度"作为裁剪指标：
$$\text{contrib}_{i,q}=\sum_{\text{pix}}\alpha_{\text{pix},i}^{2D}\prod_{j=1}^{i-1}(1-\alpha_{\text{pix},j}^{2D}),$$
即高斯在某个视角下对所有像素的透射贡献之和。若最大贡献低于阈值 $\tau_c$ 或所有视角贡献总和低于阈值 $\tau_s$，则裁剪该高斯。增量演化表示还天然适配压缩：参考点的位姿残差由变形阶段的有效初始化而数值极小，扩展点与祖先的特征差异也很小，经量化和熵编码后实现超过 50 倍压缩率。

\paragraph{运动感知高斯流式重建}

MGStream（2025.05）~\cite{2505_13839} 指出已有流式方法（如 3DGStream、Dynamic 3DGS）对所有高斯统一施加变形和优化，导致两个问题：一是静态区域的重建被不必要地修改，引发闪烁伪影；二是所有高斯的变形参数都需要存储，效率低下。

\textbf{动静分离：运动掩码与聚类凸包算法。}
MGStream 首先通过光流和时序差分联合生成运动掩码：$\mathcal{M}=(\mathcal{O}(I_t,I_{t-1}),\;\text{DILATE}(\text{ERODE}(\mathbb{1}((I_t-I_{t-1})^2)))$，其中 $\mathcal{O}$ 是光流模型，阈值为 1。然后通过最大 $\alpha$-混合权重的高斯 ID 图（GIM）将 2D 掩码反投影回 3D 高斯：$G_o=\bigcup_{i=1}^V(\text{GIM}_i\cap\mathcal{M}_i)$。但反投影只获取了物体表面的高斯，遗漏了内部的运动高斯。MGStream 因此提出聚类凸包算法：先用 DBSCAN（一种基于密度的聚类算法，无需预设簇数）对 $G_o$ 做位置聚类（阈值设为 2），再对每个聚类用 Delaunay 三角剖分构建凸包，凸包内部的所有高斯 $G_i$ 也被纳入运动相关高斯：$G_m=(G_o,G_i)$。

\textbf{刚性变形与注意力优化。}
对运动相关高斯 $G_m$ 施加刚性变形：$\Delta\mathbf{u}_t,\Delta\mathbf{q}_t=F_d(\mathbf{u}_{G_m})$，$\mathbf{u}_t=\mathbf{u}_{t-1}+\Delta\mathbf{u}_t$，$\mathbf{q}_t=\mathbf{q}_{t-1}\cdot\Delta\mathbf{q}_t$，其中 $F_d$ 是哈希网格实现的映射函数。变形能捕捉运动但无法预测新出现物体的颜色（因为 SH 系数沿用了上一帧的值）。MGStream 通过注意力图定位新物体：$M_a^{i,j}=\mathbb{1}[L_1(\hat{I}_d,I)^{i,j}>\tau]$（$\tau$ 取 $L_1$ 误差的 99 百分位），然后用渲染 GIM 与注意力图求交得到负责新物体的高斯子集 $G_{\text{new}}$，仅对这些高斯优化 SH 系数：$\text{sh}_t=\text{sh}_{t-1}+\Delta\text{sh}_t$，$\Delta\text{sh}_t=F_c(\mathbf{u}_{G_{\text{new}}})$。由于变形和优化仅在运动相关高斯上进行，静态区域的高斯保持不变，MGStream 有效避免了闪烁伪影并显著提高了存储效率。

在此基础上，StreamSplat 进一步发展了流式重建。

\paragraph{StreamSplat：2025.06 同期工作}

StreamSplat（2025.06）~\cite{2506_08862} 从未标定视频流实现在线动态 3D 重建。

\paragraph{StreamSTGS：2025.11 同期工作}

StreamSTGS（2025.11）~\cite{2511_06046} 通过流式时空高斯网格实现实时自由视角视频。\paragraph{细节层次流式视频重建}

StreamLoD-GS（2026.01）~\cite{2601_18475} 面向稀疏视角（少至 4 个相机）下的流式自由视点视频重建，针对标准 3DGS 在稀疏视角下易过拟合、产生冗余高斯并膨胀存储的问题，提出三项相互配合的技术。

第一，在初始化层面，标准 3DGS 使用 COLMAP 点云初始化，在稀疏视角下可靠性骤降。StreamLoD-GS 改用 VGGT 估计初始点云，并将高斯中心组织为 $L$ 层八叉树 LoD 结构，层级由场景尺度 $(d_{\min}, d_{\max})$ 决定：$L = \lfloor\log_2(d_{\max}/d_{\min})\rfloor + 1$。每层 $l$ 的锚点放置在体素大小为 $v_l = v_0 / 2^l$ 的网格中，每个锚点关联特征 $\phi$ 和 $K$ 个带有可学习偏移的神经高斯。为抑制稀疏视角下的过拟合，引入层次化高斯丢弃策略，dropout 率随 LoD 层级递增：$r_m^{(l)} = \alpha^{(l)} \cdot \frac{m}{M}$，其中缩放因子 $\alpha^{(l)} = 0.1 + 0.05l$ 使得更高层级的高斯被更积极地丢弃。同时用共享 MLP $\theta_{\text{mlp}}$ 替代 Scaffold-GS 的多个独立网络预测高斯属性，减少参数加速训练：$G = \theta_{\text{mlp}}(\phi, d_c, v_c)$。

第二，利用帧间时间一致性设计 GMM 运动分区。具体地，计算当前帧与规范帧（首帧）渲染结果的屏幕空间梯度差 $g_{\mu}^{\text{cano}} = \frac{1}{N}\sum_{i=1}^{N}\nabla_{\mu}^{\text{cano}}\|I_{t=c}^{(i)} - \hat{I}_{t=0}^{(i)}\|^2$，以该梯度为特征用高斯混合模型聚类，后验概率超过 $\tau = 80\%$ 的锚点被标记为动态锚点 $a_{\text{dyn}}$，其余为静态锚点 $a_{\text{static}}$ 并冻结更新，从而大幅减少后续帧的训练开销。

第三，量化残差精化框架仅对动态锚点学习量化残差 $\Delta = (\Delta_{\text{anchor}}, \Delta_{\text{location}})$，将锚点特征和位置偏移的潜在维度统一压缩到 12 维，以极低存储开销逼近全精度效果。在 N3DV 和 Meet Room 数据集的 3--6 视角设置下，StreamLoD-GS 在 PSNR、存储效率和渲染速度上均超越 3DGStream、QUEEN 等方法。

\paragraph{片段流高斯重建}

ClipGStream（2026.04）~\cite{2604_13746} 针对"逐帧流式"与"分段联合"两种范式各自的缺陷——前者帧间抖动、后者内存爆炸——提出 Clip-Stream 混合范式，提出 Clip-Stream 混合范式。

Baseline 的核心问题是：逐帧方法（如 3DGStream）通过帧间传递模型实现可扩展性，但单帧优化无法保证长序列时间一致性，导致帧间闪烁；分段方法（如 4DGS）在数百帧窗口内联合优化，时间一致性好但 GPU 内存随序列长度线性增长，且段间拼接出现边界伪影。

ClipGStream 将长视频分为 $N$ 个短片段（Clip），每段含 $M$ 帧多视角图像。第一个片段 $\text{Clip}_0$ 作为基础片段（Base Clip），其余 $\text{Clip}_n$ 为衍生片段（Source Clip）。在片段内层面（Intra-clip），每个锚点同时维护静态特征 $f_s \in \mathbb{R}^{64}$ 和动态特征 $f_d \in \mathbb{R}^{64}$，后者由片段独立的时空场（STF）生成。具体地，$\text{STF}_n$ 实现为 4D 哈希网格 $h_n$ 后接全融合 MLP $\psi_n$：$f_{d,n} = \psi_n(h_n(a_n, t))$。拼接两类特征后输入解码器生成时间高斯：$G_{t,n} = d([f_{s,n}; f_{d,n}])$。这种动静特征解耦设计使得静态信息可以在片段间传递而动态信息独立建模。对于大幅运动场景，基础片段的锚点经空间位移后难以仅靠变形场捕获，因此引入残差锚补偿策略：将基础片段继承的锚点集 $A_0$ 与新检测到的残差锚点 $A_{\text{res}}$ 合并为 $A_n = A_0 \cup A_{\text{res}}$，增强空间覆盖与运动适应性。

在片段间层面（Inter-clip），每个衍生片段从基础片段继承锚点位置、静态特征和解码器，继承组件在优化过程中冻结，仅训练残差锚点和独立的 STF。这种"冻结继承 + 局部精化"策略使得任意长度的动态序列都能保持结构一致性，不会因累积误差产生边界伪影。在 N3DV 和 VRU 数据集上，ClipGStream 在大幅运动和长序列场景下的 PSNR 和时间一致性指标均优于现有方法。

\paragraph{长程无闪烁 4D 运动建模}

MoRel（2025.12）~\cite{2512_09270} 面向长程（数分钟级）4D 运动建模，这是流式方法难以覆盖的场景尺度，解决现有方法的三类问题：全局联合优化（all-at-once）导致 GPU 内存爆炸、分段独立优化（chunk-based）破坏段间时间一致性、滑动窗口方法复杂度高且仅局部修复。

MoRel 采用锚点表示（类似 Scaffold-GS），以稀疏体素网格的锚点定义规范空间，核心设计为锚点中继双向混合（ARBB）机制，分为两个阶段、四个训练步骤。

锚点中继阶段首先训练全局规范锚（GCA）$A^{\text{Global}}$，通过稀疏采样整个视频的点云（而非密集逐帧点云）确保全局一致性，并为每个锚点 $a_k^{\text{Global}}$ 分配基于特征方差的层级 $L_{a,k}^{\text{Global}}$。随后在周期性时间索引 $t_n$ 处优化关键帧锚（KfA）$A_n^{\text{Key}}$，所有 KfA 从层级化 GCA 初始化而非从头训练，保证外观的全局一致。每个 KfA 在时间容忍度 $\epsilon$ 范围内优化，覆盖 $[t_n - \epsilon, t_n + \epsilon]$ 的局部时间邻域。为平衡内存与质量，引入特征方差引导的层次致密化（FHD）方案，根据每个锚点的局部频率特性自适应增删锚点，利用预分配的层级 $L_{a,k}^{\text{Global}}$ 抑制冗余高斯的产生。

双向混合阶段分为两步。渐进窗口变形（PWD）训练让每个关键锚在滑动时间窗口内独立学习前向和后向变形场，控制内存使用并防止段间干扰。中间帧混合（IFB）训练则学习时间融合模型，通过可学习的时间不透明度控制自适应过渡相邻 KfA 的影响力。具体地，对于时刻 $t$ 在两个相邻 KfA（$A_n^{\text{Key}}$ 和 $A_{n+1}^{\text{Key}}$）之间的中间帧，其渲染结果为双向变形输出的加权混合，权重由可学习的时间不透明度参数决定，从而实现平滑无闪烁的运动过渡。

KfA 的周期性放置（类似视频编码的 GOP 结构）提供了自然的时间随机访问点，仅需加载对应的 KfA 及其变形场，无需传输整个模型。MoRel 还构建了包含更大动态运动幅度的数据集 SelfCap$_{\text{LR}}$，实验表明该方法在长程序列上实现了内存可控的无闪烁重建。

\subsubsection{体积视频与高保真重建}

流式方法关注效率，而体积视频面向制作级应用则要求更高的视觉保真度。本节聚焦人体捕获、制片重建和沉浸视频编码，介绍九种追求极致质量的 4D 重建方法。

\paragraph{双重高斯沉浸式人体体积视频}

Robust Dual Gaussian Splatting（2024.09）~\cite{2409_08353} 针对人体体积视频制作中运动与外观高度耦合、导致跟踪不稳且存储冗余的问题，提出双重高斯（DualGS）表示：约 15,000 个运动感知关节高斯捕获全局运动，约 180,000 个外观感知皮肤高斯表达视觉细节。体积视频的现有制作流程依赖显式网格重建与非刚性跟踪，在遮挡区域容易产生孔洞和噪声，纹理质量差且需大量人工清理。隐式方法如 HumanRF 在快速运动下产生模糊，NeuS2 难以恢复高频细节；基于高斯的方法如 Spacetime Gaussian 容易过度平滑丢失衣物褶皱，HiFi4G 依赖网格重建则可能在变形失败时产生严重伪影。DualGS 的目标是无需模板的鲁棒跟踪与高保真渲染，同时实现极高压缩效率以适配 VR 头显等低内存设备。

Baseline 的做法是在标准 3DGS 中，每个高斯同时优化位置（运动）和颜色（外观），导致高斯倾向于改变颜色而非移动到正确位置来拟合光度损失。DualGS 的关键设计借鉴人体骨骼-皮肤的层级结构，将高斯属性分为运动参数（位置 $p$、旋转 $q$）和外观参数（球谐系数 $C$、不透明度 $\alpha$、缩放 $s$），分别由"关节高斯"（数量少，负责运动）和"皮肤高斯"（数量多，负责外观）承担。这种解耦带来两个核心优势：其一，皮肤高斯在局部空间中的相对位置保持稳定——因为它们始终由同一组关节高斯的运动驱动，从而增强时空一致性；其二，减少后的运动参数大幅有利于后续压缩。初始化时先训练关节高斯（30000 次迭代），通过各向同性约束 $E_{\text{iso}} = \frac{1}{N}\sum_{i=1}^{N}(\exp(\max(s_i) - \min(s_i)) - r)$（$r=4$）防止过度扁平的高斯核产生伪影，同时施加尺寸约束 $E_{\text{size}} = \sum_{i=1}^{N}\|s_i - \frac{1}{N}\sum_{i=1}^{N}\text{sg}[s_i]\|$ 惩罚超大的高斯核以避免局部过重建。关节高斯在 15000 次迭代后下采样至固定数量，之后仅优化属性值。再以关节高斯为基础初始化皮肤高斯。每个皮肤高斯 $p_i^s$ 通过 KNN（$k=8$）绑定到最近的 8 个关节高斯，混合权重为 $w(p_i^s, p_k^j) = \exp(-\|p_i^s - p_k^j\|_2^2 / l^2)$，其中 $l$ 为影响半径。这种绑定关系在初始化后固定，类似 SMPL 模型但无需预定义骨架或拓扑——实验表明该方法可处理广泛的动态序列，包括快速乐器演奏等具有挑战性的运动。

优化采用由粗到精策略（粗对齐和精细优化各 10000 次迭代）。粗对齐阶段仅优化关节高斯的运动属性，固定外观属性，施加局部尽可能刚性（ARAP）正则约束 $E_{\text{smooth}}$ 保证运动平滑。ARAP 约束的核心思想是：相邻关节高斯在运动中应保持近似刚性的相对关系，即 $E_{\text{smooth}} = \sum_i\sum_{k \in N(i)} w_{i,k}\|R(q_{i,t}^j \cdot q_{i,t-1}^{j^{-1}})(p_{k,t-1}^j - p_{i,t-1}^j) - (p_{k,t}^j - p_{i,t}^j)\|_2^2$，其中 $N(i)$ 为关节高斯 $i$ 的邻居集合，$R(\cdot)$ 将四元数转换为旋转矩阵。同时引入速度预测模块应对快速运动——为每个高斯维护速度属性，用最近两帧间的位置变化做加权更新，在粗对齐前先估计新帧位置，再用 ARAP 约束限制不合理运动。精细优化阶段在粗对齐基础上，通过预计算的 KNN 混合权重将关节高斯的运动插值到皮肤高斯：$q_{i,t}^s = \sum_{k \in N(p_{i,1}^s)} w(p_{i,1}^s, p_k^j) q_{k,t}^j$，$p_{i,t}^s = \sum_{k \in N(p_{i,1}^s)} w(p_{i,1}^s, p_k^j)(R(q_{k,t}^j)p_{i,1}^s + p_{k,t}^j)$，反向传播时皮肤高斯的梯度沿计算图传递至关节高斯。同时引入时间正则化 $E_{\text{temp}} = \sum_{a \in \{C, \alpha, s\}}\lambda_a\|a_{i,t} - a_{i,t-1}\|_2^2$ 约束外观属性的帧间变化，有效提升视觉质量并为后续压缩提供更高的压缩率。粗-精两阶段实现高斯属性的显式解耦，显著提升跟踪精度。

压缩方面，利用 DualGS 的运动-外观解耦，分别采用三种策略：关节高斯的运动属性用残差向量量化（R-VQ）配合 11-bit 时序量化，保留每段首帧位置 $R_{i,1} = p_{i,1}^j$，后续帧量化残差 $R_{i,t} = Q(p_{i,t}^j - (R_{i,1} + \sum_{k=2}^{t-1}Q^{-1}(R_{i,k})))$（$t > 1$），再用 RANS 编码器做无损压缩，有效防止误差累积；皮肤高斯的不透明度和缩放按行均值排序后嵌入 2D 查找表（高度对应高斯数量、宽度对应帧数），用图像编解码器（WebP/JPEG）以 8-bit 100\% 质量压缩；球谐系数用持久码本（长度 $L=8192$）编码，对 $d$ 阶（$d=0,1,2,3$）SH 系数做 K-Means 聚类得到四个码本，原始 $n \times f \times 48$ 个浮点参数被压缩为 $n \times f \times 4$ 个整数索引。关键是帧间 SH 索引变化极小（平均仅约 1\% 的皮肤高斯索引在相邻帧间改变），因此仅保存首帧索引和变化位置的 $(t, d, i, k)$ 四元组，经长度编码后实现高达 360 倍的 SH 压缩。消融实验表明码本长度超过 8192 后压缩效率提升甚微而存储消耗显著增加。整体实现高达 120 倍压缩比，每帧仅需约 350KB 存储。

数据集方面，DualGS 使用 81 台 Z-CAM 电影级相机以 3840$\times$2160@30fps 拍摄 8 位演员演奏中西方乐器（小提琴、吉他、钢琴、长笛、琵琶、古筝等），每个乐器包含慢/中/快三种节奏的演奏序列，每段序列以标准姿态起始以缓解"从闭合到展开"问题。所有相机配置 640\textmu s 快门速度以消除快速运动中的运动模糊。实验在单个 RTX 3090 上训练，处理时间约 12 分钟/帧；4K 渲染速度达 77 fps（含约 10 ms 的属性解码、内存拷贝和皮肤高斯运动插值），解码过程利用 CPU 资源与 GPU 光栅化并行执行。DualGS 的 Unity 插件基于 OpenXR 实现，支持将 4D 资产导入标准 CG 引擎并通过 $\alpha$-混合正确处理与场景网格的遮挡关系，已在 PICO 4 VR 头显上展示 9 位演奏者共 2700 帧的沉浸式虚拟音乐会场景。此外还开发了基于 Vulkan 的跨平台渲染器（兼容 Windows/Linux/Android，并通过 MoltenVK 扩展至 iOS/Vision Pro），支持在移动端实时播放。

方法的局限性包括：依赖图像分割提取前景，在细长物体（如琴弦、头发）上可能出现分割错误；关节-皮肤 KNN 关系初始化后固定，无法处理拓扑变化；不支持可驱动动画或动作迁移等下游任务；4D 资产缺少几何和法线信息，无法在不同光照条件下重新打光。

\paragraph{打包高斯 UV 图体积视频}

PackUV（2026.02）~\cite{2602_23040} 面向体积视频的存储与分发难题，核心观察是：3DGS 的点云存储格式既不兼容现有视频基础设施，也难以利用成熟的视频编码技术——直觉上，如果能把所有高斯属性"铺平"成一张 2D 图像，就能直接复用数十年积累的视频编解码技术。PackUV 将所有高斯属性直接在 UV 空间中优化并打包为结构化多尺度 UV 图集序列，实现图像原生的体积视频表示。

Baseline 3DGS 存储每个高斯的位置、旋转、缩放、不透明度和颜色为独立的点云数据，既无法直接用视频编码器压缩，也缺乏结构化的空间排列。PackUV 的做法是将场景高斯直接在 UV 空间中优化，用固定空间分辨率和预设层数 $K$ 的 UV 图表示：$U[u_i, v_i, k] = g_i = \{\rho_i, r_i, s_i, o_i, c_i\} \in \mathbb{R}^D$。UV 图的离散网格结构天然强制高斯稀疏性，不同于 UVGS 等方法的"先拟合 3DGS 再投影到 UV"的有损后处理路径。

金字塔 UV 映射是存储效率的关键。按不透明度排序后发现深层（更大 $K$）包含的可见高斯逐渐减少，因此采用渐进降采样策略：第 $k$ 层分辨率为 $(M_k, N_k)$，奇数层仅缩减 $N$ 维度、偶数层仅缩减 $M$ 维度，产生 $M_0 \times N_0, M_0 \times N_0/2, M_0/2 \times N_0/2, \ldots$ 的金字塔。$K$ 层渐进 UV 图通过类四叉树的递归细分打包为单一图集 $A$，打包效率达 88.5\%。图集尺寸为 $W_A = N_0 + \sum_{k=1}^{K-1} N_k$，$H_A = \sum_k M_k$。

为保证长序列时间一致性，PackUV 设计了光流引导的关键帧模块：对每帧计算光流幅度 $M(t)$，选取幅度最大的 $m-1$ 个峰值（最小间距 $\delta$）作为关键帧边界。关键帧之间为过渡帧，从上一帧初始化后仅需少量迭代精化。光流引导的机制自动将大运动、遮挡/去遮挡和外观突变的帧提升为关键帧，从而在不损失质量的前提下处理任意长度序列。

UV 图集的时间一致性和空间局部性使其可直接用标准 2D 视频编解码器存储：使用 FFV1 无损编码，平均存储率低于 10 Mbps，PSNR 仅下降 0.11 dB。PackUV 是首个将所有 3DGS 属性无损转换为标准视频编码可直接处理格式的统一表示。配合 PackUV-2B 数据集（50+ 同步相机，100 序列，20 亿帧），该方法可扩展至 30 分钟的序列。

\paragraph{细节增强大规模体积捕获}

Detail Enhanced Gaussian Splatting（2025.11）~\cite{2511_21697} 面向影视级 4K 体积视频制作，解决大规模捕获与制作分辨率标准之间的质量鸿沟。核心问题是：在 $6\text{m} \times 8\text{m}$ 的大舞台上用 180 台同步 4K 相机捕获多演员表演时，相机距离主体远，重建的面部细节远不足以支撑特写镜头。

该方法的双捕获系统包括场景摄影棚（Scene Studio）和面部摄影棚（Face Studio）。场景摄影棚配备 90 台静态广角相机（14--42mm 镜头）提供全场覆盖，40 台追踪相机（45--175mm 变焦镜头）动态跟踪演员面部，另有 50 台纵向相机跟踪全身运动。动态相机的标定是难点：焦距变化容易与相机运动混淆。该方法先用 Lidar 扫描和中间帧标定静态相机，再固定静态相机位置单独标定追踪相机，通过拟合焦距变化的平滑函数做帧间正则化，防止突然跳变，使初始点云的识别点数增加近 50\%。

重建阶段采用 Poly4DGS，将每个高斯的运动参数化为时间的多项式展开：$\mu_t = \sum_{i=0}^{n_\mu}\mu_i(t - t_0)^i$，$q_t = \sum_{i=0}^{n_q}q_i(t - t_0)^i$，$s_t = \sum_{i=0}^{n_s}s_i(t - t_0)^i$，$o_t = o_0 e^{-\frac{1}{2}\sum_{i=1}^{n_o}\sigma_i(t - t_0)^{2i}}$。实证发现 $n_\mu = 2$, $n_q = 1$, $n_s = 0$, $n_o = 2$ 在不引入可见开销的情况下带来 +0.2 dB PSNR 提升。相比逐帧 3DGS，这种 4D 参数化允许高斯在帧间共享，实现更平滑的插值和更稳定的时间一致性。此外还引入黑电平优化和曝光优化，消除镜头眩光和曝光差异造成的低对比度和偏色问题。

细节增强阶段利用面部摄影棚（75 台 4K 相机，250cm 高、276cm 直径的捕获体积）拍摄每个演员的高保真面部数据。训练数据生成策略是：对每个演员的表演采样 60 个 8 帧子序列，分别用高质量模型（100 万高斯/帧）和低质量模型（5--20 万高斯/帧）重建，渲染出约 3600 对低-高质量配对序列（$4096 \times 4096$ 分辨率）。增强模型基于 Flux 图像扩散模型，引入三项关键修改：（1）扩展 DiT 首层线性层接受 $(N+1) \times L_{\text{ch}}$ 通道的条件输入，而非仅噪声；（2）用光流重投影前帧输出作为时间条件，通过当前帧与前帧低质量渲染之间的光流对前一帧增强结果进行翘曲，保证时间稳定性；（3）修改架构同时预测 RGB 和 Alpha 通道，便于后期合成。低质量模型的参数范围（5--20 万高斯）精确模拟场景摄影棚裁切演员上半身后的高斯数量，确保训练分布与推理时一致，弥合大规模捕获与影视制作分辨率标准之间的差距。

\paragraph{稀疏多视角动态高斯制片}

Splatography（2025.11）~\cite{2511_05152} 面向制片场景的稀疏多视角动态重建问题。在常规 4DGS 中，hex-plane 变形场 $G' = G + \Delta(G|t)$ 对所有高斯统一建模全部动态参数 $(\Delta x, \Delta r, \Delta s, \Delta c, \Delta\alpha)$，但在稀疏视角下会导致背景过度重建、前景欠重建——因为背景占据更大的图像面积，训练时梯度信号被背景主导。Splatography 的核心改动是将初始点云通过 $t{=}0$ 时刻的一组稀疏掩码 $M^* \in \mathbb{R}^{H \times W}$ 分割为前景 $G_f$ 和背景 $G_b$ 两个独立表示，并分别设计变形场：前景模型 $G'_f = \delta_f(G_f|t)$ 学习完整的 $(\Delta x, \Delta r, \Delta c)$，而背景模型 $G'_b = \delta_b(G_b|t)$ 仅学习 $(\Delta x)$。这一非对称设计源于制片实践的先验假设：表演区域（背景）灯光昏暗且人员运动有限，无需学习旋转和颜色变化；而前景包含丰富的纹理、透明和动态效果，需要更强的表达能力。在规范训练阶段，$G_f$ 和 $G_b$ 都锚定在 $t{=}0$，使用专门的掩码损失函数：前景损失 $L_\alpha = \|({M^*}I^* + (1{-}M^*)B) - (\alpha_f I_f + (1{-}\alpha_f)B)\|$ 通过随机背景色 $B$ 消除浮点和边缘伪影，背景损失则通过边缘模糊约束抑制遮挡泄漏。此外，Splatography 修改了时间不透明度函数以改善透明和动态纹理的重建质量，并提出了无需参考的致密化策略减少前景欠重建。在 3D 和 2.5D 娱乐数据集上，该方法 PSNR 提升超过 3 dB 且模型大小减半，无需密集掩码监督即可生成分段动态重建，包括半透明钥匙扣和火焰等 RTD 纹理。

\paragraph{4D 高斯基沉浸视频}

GIFStream（2025.05）~\cite{2505_07539} 解决变形场方法在高度动态场景中表达能力不足的问题。现有基于变形的 4DGS 方法使用位置编码或低维特征平面驱动变形，难以捕获快速运动。GIFStream 基于 Scaffold-GS 的锚点结构，为每个锚点同时存储时间无关特征 $f^* \in \mathbb{R}^C$ 和时间相关特征流 $f_t \in \mathbb{R}^P$（$t$ 为时间戳），在每个时刻通过两个预测头分别解码高斯属性和运动：属性头 $h_{attr}: [\tilde{f}; \hat{f}_t] \to \{\alpha^i_t, s^i_t, r^i_t, c^i_t\}_{i=1}^K$ 生成 $K$ 个高斯的 opacity、缩放、旋转和颜色；运动头则先用 KNN 聚合邻近锚点特征并引入可学习混合因子 $M_{knn}$ 实现粗-细运动组合，$\tilde{f}_t = (1{-}M_{knn})\bar{f}_t + M_{knn} f_t$，最终映射为运动向量 $R_t, T_t$。特征流采用运动自适应机制：通过可学习门控 $M_{de}$ 缩放 $f_t$（即 $\hat{f}_t = M_{de} \cdot f_t$）并在训练中鼓励静态区域的 $M_{de}$ 趋近零，实验表明在复杂场景中约 30\% 的锚点需要保留时间特征，简单场景仅需 0.3\%。压缩方面，GIFStream 将时间无关和时间相关特征分别排列为两段视频序列，支持 HEVC 等视频编码器和基于自回归熵估计的端到端压缩，配合量化感知训练（quantization-aware training）。在 RTX 4090 上以 30 Mbps 码率提供 1080p 高质量沉浸视频，实现实时渲染和快速解码。

\paragraph{颜色自适应体积视频流}

CAGS（2026.05）~\cite{2605_09279}（Color Adaptive Gaussian Streaming）发现现有高斯属性压缩中的一个关键现象：激进的向量量化（VQ）主要破坏颜色而保留场景结构——因此可以用低分辨率参考图像有效修复颜色失真。基于这一观察，CAGS 提出三阶段颜色自适应方案。第一阶段是可扩展向量量化（SVQ），用于建立多级细节层次（LoD）。传统 KMeans VQ 产生扁平聚类结构，无法支持自适应流；层次聚类（AHC）虽能构建质量层级但计算复杂度为 $O(N^2)$，对每帧超过 100k 高斯的场景不可行。SVQ 先对高斯子集运行 KMeans 得到初始聚类，再以量化误差为距离度量 $d(C_1, C_2) = \sum(-\text{mean}(C_1 \cup C_2)|c \in C_1 \cup C_2)$ 迭代合并最相似聚类形成二叉树。索引分配借鉴 Huffman 编码：截断低位索引自然映射到父节点的更粗 LoD，因此码本被组织为存储高位的基础层和存储低位的增强层，解量化仅需拼接位并查表。第二阶段，服务端根据带宽和预测的客户端视口渲染低分辨率参考图像，与选定 LoD 的瓦片一同传输。第三阶段，客户端渲染接收到的瓦片得到颜色失真图像和深度图，通过后渲染透视对齐（PRPA）将参考图像重投影到实际客户端视口，再用轻量级 SRResNet 网络从对齐后的参考图和失真图恢复准确颜色。CAGS 在波动带宽下 PSNR 优于现有自适应流式系统 5$\sim$20 dB，且兼容差分帧格式，可部署于现有体积视频流基础设施。

\paragraph{实时 4D 高斯流式自由视角视频}

AirGS（2025.12）~\cite{2512_20943} 从训练和流式传输两个环节重新设计 4DGS 管线。在训练阶段，现有方法一旦建立规范空间（canonical space）就固定高斯数量，导致出现新物体或大幅运动时无法通过分裂生成新高斯，重建质量急剧下降。AirGS 提出基于质量的帧分组策略：用质量阈值 $\tau$（取 PSNR = 30 dB）监控重建质量，当某帧质量低于阈值时将其标记为新关键帧并启动新的规范空间。组内帧训练利用时间连续性，用 6 个 MLP（分别对应 $\alpha$, $c$, $SH$, $R$, $S$ 和位置）预测帧间运动 $d_{t+1} = [\text{MLP}_1(GS_t), \ldots, \text{MLP}_6(GS_t)]$，通过 $GS_{t+1} = GS_t + d_{t+1}$ 高效生成下一帧，并引入时间损失 $L_{temp} = \sum_{i=1}^{6}\|\text{MLP}_i(GS_t)\|_1$ 防止帧间跳变。对于关键帧的快速重建，由于新内容通常仅占场景约 20\%，AirGS 复用前一帧的 3DGS 作为初始化，通过高斯致密化捕获新区域，同时引入膨胀损失 $L_{inf} = \max(0, N - U)$（$N$ 为当前高斯数，$U$ 为前一个规范空间的高斯数）抑制冗余高斯增长。在流式传输阶段，AirGS 将每帧 3DGS 编码为多通道 2D 图像（每个属性一张图，每个像素对应一个高斯），关键帧传输完整图像，组内帧仅传输差分 $d_t$，利用帧间稀疏性降低带宽。为应对波动带宽，AirGS 设计了基于使用频率的自适应剪枝策略：预先计算不同剪枝级别下的（质量，大小）曲线，在给定带宽约束下通过二分搜索选择最优剪枝级别——当质量下降出现"悬崖"（即质量骤降超过阈值 $\kappa$ 倍）时自动停止加深剪枝。在 HiFi4G 数据集上，AirGS 每帧传输量从 18.5 MB 降至 1.1 MB，训练加速 6 倍，场景变化时 PSNR 质量偏差减少超过 20\%。

\paragraph{解码端高斯泼溅沉浸视频}

Decoder-Side Gaussian Splatting（2026.05）~\cite{2605_17002}（DSGS）面向沉浸视频传输中的像素率瓶颈问题。现有解码端深度估计（DSDE）方案将深度计算转移至客户端，但在复杂几何、跨视图闪烁和非朗伯反射场景下表现不佳；而直接传输 3DGS splat 数据的带宽成本过高，且与标准视频编码器不兼容。DSGS 的核心思路是原生替换 DSDE 的深度估计阶段为前馈 3DGS 推理：服务端仅传输压缩的 2D 纹理图集（atlas）和少量元数据，解码端从这些压缩数据直接预测高斯 splat 参数并优化体积场景，无需传输任何显式 3D 数据或深度图。该方法的一个关键反直觉发现是：有损压缩对纹理的影响类似于隐式低通滤波，反而稳定了前馈 splat 预测——压缩后的比特流在体积缩小约十倍的同时，重建质量超过无损输入。这是因为无损纹理中的高频细节（噪声、精细纹理）会干扰前馈网络的 splat 参数回归，而有损压缩自然滤除了这些干扰。在极端视角稀疏条件下（仅一个包含 4 个输入视角的 2D 图集），DSGS 实现 +5.79 dB BD-PSNR 和 +0.054 BD-SSIM 增益，同时将跨视图 Delta IV-PSNR 的最大值从 17.2 dB 降至 6.4 dB，显著缩小了传输视口与虚拟视口之间的域偏移。

\subsubsection{时序一致性与视频增强}

本路线介绍从后处理与对齐角度增强视频时间一致性和处理非同步采集的方法。

前面各小节通过改进表示本身来提升时序质量，而本节从后处理与对齐角度出发，介绍增强视频时间一致性和处理非同步采集的七种方法。

\paragraph{3DGS 增强视频编辑时序一致性}

Video-3DGS（2024.06）~\cite{2406_02541} 是一个基于 3DGS 的即插即用视频精化器，用于增强任意零样本视频编辑器的时间一致性。零样本视频扩散模型（如 Text2Video-Zero、TokenFlow、RAVE）虽然无需训练即可编辑视频，但由于对视频场景缺乏结构理解，编辑结果的帧间一致性很差。Video-3DGS 的核心思想是利用 3DGS 的多视角统一表示能力建立帧间时间对应关系，分两个阶段实现。

\textbf{第一阶段：MC-COLMAP 视频重建。}原始 COLMAP 在处理视频时面临两大困难：前景运动物体和杂乱背景导致跨帧匹配失败。Video-3DGS 提出空间-时间双重分解策略。空间分解方面，用开放词汇视频分割网络 $S$（如 DEVA）提取前景掩码 $V^f = S(V, class)$，将前景与杂乱背景分离。时间分解方面，将前景掩码序列划分为 $M$ 个短片段——但并非均匀切分，而是采用渐进式策略：以 $k$ 帧初始化第一个片段，若 COLMAP 处理失败则逐帧扩展直到成功，下一个片段从上一片段的最后一帧开始重复此过程。这种"Masked and Clipped COLMAP"（MC-COLMAP）表示为 $\{(p_j, c_j)\}_{j=1}^M = R_{MC}(S(V, class))$，在全部 28 个 DAVIS 视频上均成功重建，而原始 COLMAP 有 8 个视频失败。

\textbf{前景/背景分离 3DGS（Clip-3DGS）。}对每个片段，MC-COLMAP 输出的前景 3D 点云初始化前景高斯 Frg-3DGS；背景则使用球形随机点云（60k 个点，半径为前景点最大欧氏距离的 3 倍）初始化背景高斯 Bkg-3DGS。两组高斯均配备基于 4D multi-resolution hash encoding 的变形网络，根据片段中心时间和归一化时间预测位置、旋转和尺度的偏移量，从而建模片段内的动态运动。每个片段的 Frg-3DGS 和 Bkg-3DGS 合称为 Clip-3DGS。渲染时，对每帧分别渲染前景图 $I^{Frg}$ 和背景图 $I^{Bkg}$，再通过 2D 可学习参数图 $\theta \in \mathbb{R}^{H \times W}$（初始化为 0.5）逐像素合并：$I^{merged} = \theta \cdot I^{Frg} + (1{-}\theta) \cdot I^{Bkg}$，其中学到的 $\theta$ 自动聚焦于前景区域。训练损失包含 $L_1$ 和 $L_{SSIM}$，分别对前景渲染、背景渲染和合并图像与对应真实图像计算。

\textbf{第二阶段：即插即用时间精化器。}重建好的 Clip-3DGS 被转化为时间精化器：固定位置 $(x, r, s)$ 和变形参数，仅微调球谐系数（SH）和 opacity，以编辑后视频帧作为监督信号。仅 1k 次迭代即可将 3DGS 的结构保持能力迁移到编辑域。更关键的是，Video-3DGS 进一步提出递归与集成精化策略（Recursive and Ensembled Refiner, RE）：将扩散模型的总去噪步数 $N_d$ 分为多个阶段（效率考虑取 $N_r=2$），第一阶段先对 $N_d/2$ 步去噪结果做单次精化，产生的中间帧作为下一阶段 DDIM 逆变换的输入，完成剩余 $N_d/2$ 步去噪后再做最终精化。此外，由于扩散模型中引导参数（如 Text2Video-Zero 的图像引导、TokenFlow 的文本引导、RAVE 的 ControlNet 引导）对编辑质量高度敏感，RE 在每个递归阶段对三种不同引导尺度生成的编辑结果做 $SH$ 和 $\alpha$ 的集成更新，降低对单一尺度的敏感性。在 DAVIS 数据集上以 3k 次迭代训练，Video-3DGS 相比 NeRF 和 3DGS 方法分别提升 +3 和 +7 PSNR，训练效率提高 1.9 倍和 4.5 倍。

\paragraph{连续时间 4D 高斯重定时重建}

现有 4DGS 方法将时间不透明度（temporal opacity）作为自由参数在整数时间戳上优化，缺乏正则化——直觉上，这就像只在整数小时设闹钟，但事件的真正发生在小时之间的任意时刻。结果是时间不透明度过拟合到离散帧上、坍缩为子帧级别的窄支撑，在渲染中间帧时产生半透明重叠结构——即时间混叠（temporal aliasing）导致的鬼影伪影。RetimeGS（2026.03）~\cite{2603_13783} 的核心思想是：用双 sigmoid 定义每个高斯基元的有界时间生命周期，并显式建模其空间轨迹，使表示在连续时间上平滑有意义。

\textbf{时间不透明度设计。}每个基元 $p$ 被赋予时间均值 $\mu_p$、左右偏移量 $\delta_l, \delta_r$。初始化时令 $\mu_p = (t_i + t_{i+1})/2$, $\delta_l = \delta_r = \Delta t/2$，即基元仅存活于相邻两帧之间。时间不透明度由两个 sigmoid 函数的乘积定义：
$$\tau_p(t) = \sigma_l\!\bigl(-\beta(t - (\mu_p - \delta_l))\bigr) \cdot \sigma_r\!\bigl(\beta((\mu_p + \delta_r) - t)\bigr),$$
其中 $\sigma(\cdot)$ 为 sigmoid 函数，$\beta$ 是控制过渡平滑度的超参数。在视频边界处，对应的 sigmoid 被替换为常数 $1$ 以防止可见度下降。直觉上，每组基元在两帧之间中心活跃并被两帧同时监督；在输入帧附近，相邻两组基元分别淡入和淡出，实现无缝过渡。对于被检测为静态的基元，$\delta_l$ 和 $\delta_r$ 周期性扩展至 $(\frac{1}{2} + k)\Delta t$，让其覆盖更长时间范围并减少冗余基元。

\textbf{空间轨迹与光流监督。}仅正则化时间不透明度不足以保证准确的插值——稀疏时间输入下，相邻帧间运动物体的内容重叠极少，RGB 监督无法学到可靠对应关系。因此 RetimeGS 用 Catmull-Rom 样条（一种通过四个控制点定义的平滑插值曲线）建模空间均值 $x_p(t)$，每个基元参数化为伪空间均值 $\bar{x}_p$ 和三个速度分量 $(v_1, v_2, v_3)$。四个控制点为：
$$p_0 = p_1 - \Delta t \cdot v_1, \quad p_1 = \bar{x}_p - \tfrac{\Delta t}{2} v_2, \quad p_2 = \bar{x}_p + \tfrac{\Delta t}{2} v_2, \quad p_3 = p_2 + \Delta t \cdot v_3.$$
其中 $v_2$ 是当前帧间隔的线性速度（来自 3D 对应关系），$v_1, v_3$ 是相邻间隔的速度。这些参数由双向光流（WAFT）生成的 2D 光流通过可微光栅化进行监督，而非直接监督 3D 位置。光流监督的学习率在训练后期逐渐衰减至零，交由 RGB 监督做精细调优。

\textbf{三重渲染。}在输入帧 $t_i$ 处，相邻两组基元各自理论上应能独立解释该帧，但实践中分别渲染时覆盖不均。RetimeGS 对每个内部帧渲染三张图：全部基元合渲染、前一组单独渲染、后一组单独渲染（后两者做时间不透明度归一化），全部用 GT 图像监督，从而迫使每组基元都能完整覆盖帧内容，避免中间帧渲染时出现欠重建。

\paragraph{聚类误差校正流式 4D}

4DGS 在动态区域的重建面临两个难题：（1）标准的基于平均视角梯度的高斯致密化策略无法精确定位动态误差区域，因为 splat 外观随时间变化，且误差受时间变换和可见性变化的干扰；（2）单个 splat 只有一种颜色，在相邻帧间存在多个颜色相近的像素候选对应，时间扩展时容易匹配错误导致闪烁。Grouped 4D Gaussian Splatting（2025.11）~\cite{2511_16112} 提出分组时序建模与椭圆误差聚类两个互补策略来解决这些问题。

\textbf{分组时序建模。}传统 4DGS 中每个 splat 独立拥有时间变换参数，难以保证同一刚体上各 splat 的一致运动。该方法将每个 splat 的变换分解为组级共享动态变换和 splat 级相对静态变换。具体地，每个组拥有 $K$ 个关键帧的组级位置 $x_G \in \mathbb{R}^{3 \times K}$ 和旋转 $r_G \in \mathbb{R}^{4 \times K}$，在时间上插值产生连续运动。每个 splat 有相对位置 $\tilde{x}$、相对旋转 $\tilde{r}$ 和全局位移向量 $d$。最终位置和朝向为：
$$x(t) = x_G(t) + R_G(t)\tilde{x} + t \cdot d, \quad R(t) = R_G(t) \cdot \tilde{R},$$
其中 $R_G(t)$ 和 $\tilde{R}$ 分别是组级和 splat 级旋转矩阵。静态 splat 的组变换固定为单位变换，退化为标准 Ex4DGS 的静态公式。动态分组通过图分裂实现：当 splat 的位移 $d$ 超过阈值时，构建以 splat 为节点的无向图，边连接条件为空间重叠且位移方向余弦相似度超过 $\theta_d$：
$$\|x_i(t) - x_j(t)\|_2 < s_i + s_j, \quad \frac{d_i^\top d_j}{\|d_i\|\|d_j\|} > \theta_d.$$
连通分量形成新的动态组，确保运动一致的 splat 共享变换，显著减少时间抖动。

\textbf{椭圆误差聚类与校正。}每隔数百次迭代，方法随机选一个视角对比渲染结果与 GT，通过像素动态性（相邻帧最大 L1 距离超过阈值 $\tau_D$）和 RGB 误差（绝对阈值 $\tau_a$ 加相对阈值 $\tau_r$）筛选误差像素，再用 DBSCAN 空间聚类后递归分割直至每簇可拟合为椭圆。每个椭圆簇被分类为两类错误：（1）\emph{缺失颜色}——该区域需要新的颜色信息，通过从多视角交叉验证后反投影注入新 splat；（2）\emph{遮挡}——前景遮挡导致背景信息丢失，通过前景分裂策略在遮挡边界处生成新的 splat 组。椭圆形状可直接映射为 3D 椭球高斯，保证了注入 splat 与误差区域的几何一致性。在 Neural 3D Video 和 Technicolor 数据集上，该方法将时间一致性显著提升，PSNR 提高 0.39 dB。

\paragraph{跨视频运动对齐}

多视角 4DGS 重建通常假设所有相机严格同步，但实际拍摄中硬件触发抖动、网络延迟和手动操作会导致毫秒到秒级的时间错位，使得多视角融合时出现严重的鬼影和运动模糊。现有方法如 HCP 依赖人体模板（SMPL）作为时空参考，限制了适用场景。SyncTrack4D（2025.12）~\cite{2512_04315} 不依赖任何预定义校准目标，利用密集 4D 像素轨迹作为跨视频同步和 4DGS 重建的统一线索。

\textbf{密集 4D 特征轨迹估计。}对每个单目视频 $v$，方法提取 2D 像素轨迹（SpatialTracker）、光流（RAFT）和 DINOv3 特征图 $F^v$，结合多视角一致深度，通过扩展 MoSca 的运动样条支架 4DGS 构建 4D 密集轨迹集合 $\mathcal{T}_v = \{\tau_i\}_{i=1}^{N_v}$。每条轨迹 $\tau_i$ 关联一个时间不变的语义特征向量 $f_i$（1024 维 DINOv3），这使得跨视频匹配可以在时间无关的空间中进行。

\textbf{跨视频 4D 轨迹匹配。}给定两个视频 $v_a, v_b$ 的密集轨迹集合，方法采用 Fused Gromov-Wasserstein（FGW）最优传输框架建立对应关系。与纯特征匹配不同，FGW 联合优化特征相似度和结构保真度：
$$\pi^* = \arg\min_{\pi \in \Pi(a,b)} \sum_{\substack{i,k \in [N_a] \\ j,l \in [N_b]}} (C^a_{ik} - C^b_{jl})^2\, \pi_{ij}\,\pi_{kl} + 2\alpha \sum_{i,j} M^{ab}_{ij}\,\pi_{ij},$$
其中 $M^{ab}_{ij} = 1 - f_i^\top f_j$ 是余弦距离特征相似度矩阵，$C^v_{ij} = \max_t \|\tau_{i,t} - \tau_{j,t}\|_2$ 是视频内几何距离矩阵（编码轨迹间的刚体关系），$\alpha$ 平衡结构项和特征项。从软传输矩阵 $\pi^*$ 中用匈牙利算法提取离散匹配。

\textbf{帧级同步与亚帧精化。}
跨视频匹配轨迹建立后，方法通过动态时间规整（DTW）执行全局帧级时间对齐——找到使匹配轨迹运动重叠最大化的帧对帧分配。随后，在多视频联合 4DGS 优化中，每条样条支架锚点轨迹被参数化为 SE(3) 序列，通过光度损失和轨迹一致性损失端到端精化时间偏移，实现亚帧级同步精度。在 Panoptic Studio 数据集上达到 26.3 PSNR，时间误差低于 0.26 帧。

\paragraph{非同步视频动态重建}

4DGS 重建的基本假设是所有相机在相同逻辑时间戳严格同步拍摄。当使用消费级相机（智能手机、GoPro）或无线网络控制的分布式相机系统时，缺乏中央时钟信号加上网络延迟和手动操作，几乎不可避免地引入视频流间的时间偏移。这种偏移导致模型将不同时刻的多视角观测错误地融合为同一时间点的场景，快速运动物体的微小时间差就能产生显著位置差异，引发鬼影伪影和运动模糊，并误导 4DGS 优化将时间误差错误归因于空间几何或外观缺陷。

Dynamic Gaussian Scene Reconstruction from Unsynchronized Videos（2025.11）~\cite{2511_11175} 提出了一种\textbf{由粗到细的时间对齐模块}，将每台相机的未知时间偏移 $\Delta t = \Delta t_c + \tau$ 显式纳入到优化目标中（$\Delta t_c$ 为帧级粗偏移，$\tau$ 为亚帧级精化参数）。

\textbf{粗对齐阶段。}
核心洞察是：当不同视角的视频帧确实拍摄于同一时刻时，动态前景物体等效于一个静态刚体从多个方向被观察——此时密集特征匹配器（LoFTR）会产生大量高置信度匹配。方法首先用视频分割模型提取每帧的前景掩码 $M^t$，然后对每对相机 $(i, j)$，在候选偏移范围 $[-k, k]$ 内搜索使前景内匹配数量最大的整数偏移：
$$\Delta t^*_j = \arg\max_{\Delta t_j \in [-k, k]} \sum_{t_i} N_{\text{inlier-fg}}(I_\text{ref}^{t_i},\, I_j^{t_i + \Delta t_j}),$$
其中 $N_{\text{inlier-fg}}$ 是经前景掩码过滤后用 RANSAC 拟合基础矩阵得到的前景内点数量。高内点数意味着两视角中前景物体的姿态在几何上一致，是正确时间对齐的强证据。

\textbf{亚帧精化阶段。}
粗对齐给出整数帧偏移后，方法为每台相机引入可学习参数 $\tau_j$，表示连续的残差时间偏移。查询时间戳变为 $t' = t + \Delta t^*_j + \tau_j$，$\tau_j$ 在训练中与 4DGS 模型联合优化——由于 4DGS 渲染对时间输入可微，光度重建损失的梯度自然反传至 $\tau_j$，端到端发现并校正亚帧级时间偏差。该模块可作为即插即用组件集成到神经 4DGS（如变形 MLP 网络）或直接 4D 表示（如 RT4DGS 的 4D 协方差参数化）中，无需修改基线架构即可显著增强其对非同步数据的鲁棒性。

\paragraph{颜色编码照明高速重建}

高速相机的帧率受限于传感器到内存的数据带宽，且高速模式下总帧数通常很少（如 256 帧），严重限制捕获时长和需要精确触发。Color-Encoded Illumination（2026.04）~\cite{2604_26920} 提出一种范式转换：用低速相机（60 FPS）配合快速颜色编码频闪照明，将高速时间信息压缩编码到单帧的 RGB 通道中，再通过改进的动态 GS 从编码帧中解码出高速 3D 场景。

\textbf{成像模型。}假设场景由均匀反射率的运动物体组成、背景足够暗。场景被一组共置 RGB LED 照明，在相机曝光期间，LED 以 $N$ 种不同颜色快速顺序频闪。每个频闪间隔内物体近似静止，因此相机捕获的单帧像素 $I(\mathbf{x}, c)$（$\mathbf{x}$ 为像素位置，$c \in \{R,G,B\}$）是 $N$ 个高速帧的线性混合：
$$I(\mathbf{x}, c) \approx \sum_{n=1}^{N} c^n_{\text{RGB}}(c) \cdot I^n(\mathbf{x}),$$
其中 $c^n_{\text{RGB}} = (\alpha_n, \beta_n, \gamma_n)$ 是第 $n$ 个频闪的归一化颜色（由 RGB LED 强度比例决定），$I^n(\mathbf{x})$ 是第 $n$ 个高速帧的灰度强度。直觉上，每个颜色通道捕获的是所有高速帧的加权和，权重由频闪颜色决定。实验中颜色序列在 $(\alpha, \beta, \gamma)$ 空间中沿等间隔圆周采样（等效为三路正弦波、各相隔 120 度），类似 Gray 编码——相邻帧颜色相近，即使解码出错也只是时间上相邻的帧被混淆。

\textbf{动态 GS 重建。}方法基于 Gaussian-Flow 框架，将高斯参数建模为时间的函数 $g(t) = g(0) + d(t)$，变形函数 $d(t)$ 为多项式加傅里叶级数。关键修改是训练目标：不再直接将渲染图与输入图比较，而是用 $N$ 个渲染帧 $\hat{R}_n(\theta_m)$ 估计输入帧：
$$\hat{I}(\mathbf{x}, c, \theta_m) = \sum_{n=1}^{N} c^n_{\text{RGB}}(c) \cdot \hat{R}_n(\theta_m),$$
训练损失为 $\{G(0), D(t)\} = \arg\min \sum_{m=1}^{M} \sum_{c \in \{R,G,B\}} |I(c, \theta_m) - \hat{I}(c, \theta_m)|$。多视角几何约束通过可微渲染自然融入，帮助在欠定方程组中更鲁棒地分解混合帧。实际原型使用 8 台 60 FPS 全局快门相机，LED 频闪 $N=10$ 种颜色，将时间分辨率从 60 FPS 提升至 600 FPS，每个频闪间隔约 0.083 ms。这是首个高速体积场景重建工作，展示了对 Nerf 枪飞镖、旋转圆盘等高速运动物体的多视角自由视点渲染。

\subsubsection{特殊应用与生成辅助}

本节关注 4DGS 在特殊传感模态、极端运动和生成式内容创建中的应用，涵盖八种方法。

\paragraph{物理事件增强大时空运动重建}

大时空跨度（large spatiotemporal span）的快速刚体运动重建面临三个难题：（1）现有变形场范式针对小幅度非刚体运动设计，离散采样框架难以表示强非线性运动；（2）RGB 相机的运动模糊使对应关系估计失效，而事件相机虽具微秒级时间分辨率，但现有事件-NeRF 方法仅限于静态场景或轻微变形；（3）纯光度监督无法约束庞大的 3D 解空间，而引入物理仿真模型本质上是在已有近似上叠加新的近似。PEGS（2025.11）~\cite{2511_17116} 提出两阶段框架——目标建模与运动恢复——并用三层监督方案融合物理约束、事件信号和多源观测，另引入运动感知模拟退火（Motion-aware Simulated Annealing, MSA）策略自适应调度训练过程以增强收敛效率和稳定性。

\textbf{目标建模阶段。}该阶段的目标是从模糊 RGB 输入中重建以目标为中心的高斯表示，流程依次包含事件去模糊、中心化、高斯重建和逆注册四个步骤。

\textit{事件去模糊。}标准 3DGS 依赖 COLMAP 进行初始点云和相机位姿估计，但对模糊图像序列该流程经常失败。基于 EDI 模型和 MAENet，方法将每张模糊图像对应的曝光时长划分为 $N$ 个时间戳，获得 $N-1$ 个对应的事件仓 $\{B_i\}_{i=1}^{N-1}$，进而恢复 $N$ 张去模糊帧 $\{I_i\}_{i=1}^{N}$。

\textit{中心化。}基于去模糊帧，首先用 SAM（Segment Anything Model）将动态物体从静态背景中分离。随后建立以物体为中心的归一化坐标系：消除平移分量，将物体的自转等价为伪多视角观测，再用 SfM 初始化 $N$ 个伪相机位姿。这一过程将复杂的抛体运动分解为独立的平移和旋转分量，减轻复杂运动对伪视角估计精度的负面影响，同时为后续物理约束中的质心计算提供便利。

\textit{高斯重建。}参照 3DGS 框架，用一组高斯核显式重建目标物体。去模糊渲染损失为 $\mathcal{L}_\text{blur} = (1-\lambda)\|\hat{I}_\text{blur} - I_\text{blur}\|_1 + \lambda \mathcal{L}_\text{D-SSIM}$，其中 $\hat{I}_\text{blur} = \frac{1}{N}\sum_{i=1}^{N} \hat{I}_i$ 为 $N$ 张渲染去模糊帧的平均。为改善目标高斯的几何精度（影响后续质心定位的准确性），方法引入改进的点密度控制策略：硬修剪锯齿边缘和尖刺伪影，移除远离目标的孤立点。最终在中心化尺度上重建出具有精细化外观和几何表示的高斯集合 $\Theta_0$。

\textit{逆注册。}在恢复运动轨迹之前，需要将中心化坐标系下的 $\Theta_0$ 注册回原始完整动态场景。方法使用一组可学习仿射变换 $T_\text{reg} = [q, t, s]$（四元数旋转 $q \in \mathbb{R}^4$、平移向量 $t \in \mathbb{R}^3$、常数缩放 $s$）实现这一映射，目标高斯 $\Theta_1 = T_\text{reg} \circ \Theta_0$。注册优化目标为最小化 $\Theta_1$ 的渲染图与原始动态场景首帧去模糊图 $I_0$ 之间的差异：
$$T_\text{reg}^* = \arg\min_{T_\text{reg}} \mathcal{L}_\text{reg}\bigl(\text{Render}(\Theta_1),\; I_0\bigr),$$
其中注册损失 $\mathcal{L}_\text{reg} = \lambda_1 \mathcal{L}_\text{GS} + \lambda_2 \mathcal{L}_\text{VGG} + \lambda_3 \mathcal{L}_\text{TV}$，除标准 GS 重建损失 $\mathcal{L}_\text{GS}$ 外，还引入 VGG 感知损失 $\mathcal{L}_\text{VGG}$ 和网格化全变差损失 $\mathcal{L}_\text{TV}$ 以增强对精细结构的感知和空间平滑性约束。至此，建模阶段从模糊输入中重建出清晰的目标高斯 $\Theta_1$，为后续运动恢复奠定基础。

\textbf{运动恢复阶段——三层监督。}
该阶段通过逐帧估计物体的 SE(3) 位姿变换来恢复完整运动序列。如前所述，在视频序列中，时刻 $t$ 的每张模糊输入对应 $N$ 张去模糊帧（等间隔 $\Delta t$），由此产生 $N-1$ 组待学习的位姿变化 $\{P_t = [R_t, T_t]\}$。除基础模糊损失（以完整非中心化模糊输入为真值监督，形式与建模阶段一致）外，方法引入紧密耦合的三层监督方案：

（1）\textbf{加速度一致性损失}（物理层面）：基于牛顿动力学，物体在恒力场中加速度恒定，因此对加速度施加平滑性约束可确保估计位移符合物理一致性。质心近似为 $\bar{c} = \frac{\sum_{i=1}^{M} R_i^3 \mu_i}{\sum_{i=1}^{M} R_i^3}$（$R_i$ 为高斯核半径，$\mu_i$ 为位置），加速度由位移差分计算：$a_t = \frac{\mu_t - \mu_{t-\Delta t} - (T_t \cdot \Delta t)}{(\Delta t)^2}$。损失最小化相邻时间戳的加速度变化：$\mathcal{L}_\text{Acc} = \sum_{t=0}^{T} \|a_t - a_{t+\Delta t}\|_2^2$。

（2）\textbf{事件流损失}（时间分辨率层面）：利用事件相机的微秒级分辨率，对连续运动恢复施加约束。具体而言，将每张模糊帧对应的 $N-1$ 个事件仓按 EDI 模型积分为事件图 $E_i$（作为真值监督信号），同时从 $N$ 张渲染帧之间合成模拟事件序列 $\hat{E}_i$（通过计算相邻渲染图像的强度变化获得），用 $\mathcal{L}_\text{Event} = \sum_{i=1}^{N-1} \|E_i - \hat{E}_i\|_1$ 约束运动细节。

（3）\textbf{卡尔曼融合损失}（多源融合层面）：单独使用事件信号或物理模型各有局限——事件信号稀疏且含不可避免的噪声，物理模型受空气阻力等真实干扰偏离理想假设。为此，方法用卡尔曼滤波（Kalman Filter, KF）融合两者，获得最优估计作为正则化参考。卡尔曼滤波的完整流程包含四个步骤：

\textit{初始化。}定义状态向量为物体的瞬时位移和速度 $s_t = [T_t, v_t]$；系统模型基于加速度约束从上一时刻的最优估计预测当前位移；观测方面，将事件去模糊图输入 RAFT 光流估计器，从光流中映射出观测位移 $\hat{T}_t$。

\textit{预测。}基于系统模型和上一时刻的最优估计 $s_{t-\Delta t|t-\Delta t}$，预测当前状态向量 $s_{t|t-\Delta t}$。状态转移由控制矩阵驱动，控制项由加速度方程计算；过程噪声 $w_t$（包含位移噪声 $\sigma_T^2$ 和速度噪声 $\sigma_v^2$）表征模型不确定性。预测协方差矩阵通过 $\Sigma_{t|t-\Delta t} = A \Sigma_{t-\Delta t|t-\Delta t} A^T + Q$ 递推更新。

\textit{校正。}引入观测值 $z_t = \hat{T}_t + v_t$（$v_t$ 为事件观测噪声）修正预测状态，其中观测矩阵 $H = [I, 0]$ 将状态空间映射到观测空间。

\textit{更新。}基于预测和校正，动态更新卡尔曼增益 $K_t$ 和后验协方差 $\Sigma_t$：
$$K_t = \Sigma_{t|t-\Delta t} H^T (H \Sigma_{t|t-\Delta t} H^T + R)^{-1}, \quad \Sigma_t = (I - K_t H) \Sigma_{t|t-\Delta t}.$$
最终得到当前时刻的最优估计 $s_t^* = s_{t|t-\Delta t} + K_t(z_t - H s_{t|t-\Delta t})$。取最优估计中的平移分量 $\tilde{T}_t$ 为参考信号，监督学习的 $P_t$：$\mathcal{L}_\text{KF} = \sum_{t=0}^{T} \|P_t - \tilde{P}_t\|_2^2$。本质上，卡尔曼滤波对事件观测做贝叶斯估计，用牛顿力学模型作为先验提供自适应正则化，实现多源观测融合。

运动恢复阶段的总损失为 $\mathcal{L}_\text{Total} = \lambda_1 \mathcal{L}_\text{Blur} + \lambda_2 \mathcal{L}_\text{Acc} + \lambda_3 \mathcal{L}_\text{Event} + \lambda_4 \mathcal{L}_\text{KF}$，四项损失分别从光度一致性、物理一致性、时间分辨率和多源融合四个维度协同约束运动估计。

\textbf{运动感知模拟退火（MSA）。}固定学习率难以适应变速运动：高速阶段收敛不足，低速阶段震荡。方法基于卡尔曼滤波的最优位移估计 $\tilde{T}_t$ 自适应调度下一时间戳的学习率：
$$l^{(n)}_{t+\Delta t} = \left(l_\text{min} + \frac{l_\text{max} - l_\text{min}}{|T_\text{max}|} |\tilde{T}_t|\right) \cdot \beta^n, \quad \beta \in (0, 1),$$
大位移用大学习率快速收敛，小位移用小学习率避免震荡，指数衰减项 $\beta^n$ 随训练步数 $n$ 渐进冷却，使高斯逐步收敛到目标空间位置。此外，方法还贡献了首个面向自然快速刚体运动的 RGB-Event 配对数据集，为该领域的后续研究提供了基准评测资源。

\paragraph{交替曝光单目 HDR 4D 重建}

标准 3DGS 使用 LDR 颜色 $c \in [0,1]$ 表示高斯外观，而 HDR 场景要求表示远超 $[0,1]$ 的辐照度范围。在单目交替曝光场景下，问题更加棘手：相机位姿未知，不同帧的曝光时间 $t_s$（短）和 $t_l$（长）交替变化，标准光度重投影误差无法直接用于位姿优化。

Mono4DGS-HDR（2025.10）~\cite{2510_18489} 针对上述困难设计了两阶段优化流程。第一阶段的核心洞察是：如果不需要相机位姿，就能绕过"交替曝光导致光度误差失效"的瓶颈。为此，该方法借鉴 SaV 的思路，在正交相机坐标空间中学习视频 HDR 高斯。在该空间中，3D 点 $p^v = [x^v, y^v, z^v]$ 的 $(x^v, y^v) \in [-1,1]^2$ 即为投影像素的归一化坐标，$z^v$ 为深度；高斯位置用三次 Hermite 样条函数参数化（含 $N_c$ 个控制点），旋转用三次多项式表示四元数 $r(t) = \sum_{j=0}^{3} a_j t^j$。这种正交空间表示使相机运动和物体运动统一建模为高斯的运动，完全消除了对相机参数的依赖。同时，HDR GS 将 LDR 颜色替换为 HDR 辐照度 $e \in [0, +\infty)$，并通过对数域色调映射器将辐照度转为 LDR 颜色：$c = \tau(\log(e \cdot t))$。

第二阶段将视频高斯变换到世界空间。位置变换为 $\mu_t^w = R_t \Pi_K^{-1}(\mu_t^v) + T_t$，旋转变换为 $R_t^w = R_t R_t^v$，其中 $R_t, T_t$ 是由 bundle adjustment 得到的初始外参。对于静态高斯，位置取时间平均 $\bar{\mu}^w = \frac{1}{N_f}\sum_{t=1}^{N_f}\mu_t^w$；缩放则基于 2D 协方差不变性重新拟合：$\min_{S^w}\sum_{t=1}^{N_f}\|\Sigma'^v_t - \Sigma'^w_t\|_2$，确保变换前后投影 2D 高斯的形状一致。由于视频高斯提供了良好初始化，世界空间优化可快速稳定收敛，并联合精化相机位姿。

交替曝光视频还带来一个特殊问题：动态物体在过曝/欠曝帧中缺乏有效监督，导致动态高斯漂浮在物体表面上方，HDR 辐照度随时间剧烈抖动。为此，该方法提出时间亮度正则化（Temporal Luminance Regularization, TLR），用光流引导相邻帧间的 HDR 辐照度对齐：

$$L_{tlr} = \left|V_t^{t-1} \frac{H_{t-1 \to t} - H_t}{H_{t-1 \to t} + H_t}\right|_1$$

其中 $V_t^{t-1}$ 是基于深度排序的遮挡有效掩码，除以 $H_{t-1 \to t} + H_t$ 消除辐照度绝对尺度的影响。TLR 仅在高斯致密化结束后启用（使用渲染光流做 warp），使良好曝光时刻的监督信号传播到曝光不足/过曝时刻，从而保证 HDR 外观的时间一致性。

\paragraph{高动态范围神经场景流场}

HDR 重建的核心难题在于 LDR 图像与 HDR 辐照度之间的非线性映射关系。传统 NSFF 框架假设输入帧亮度一致，直接用光度一致性约束场景流和外观；但在交替曝光输入下，相邻帧颜色差异极大，传统光流估计严重退化，NSFF 的光度损失也无法直接使用。

HDR-NSFF（2026.03）~\cite{2603_08313} 在 NSFF 框架之上引入三个核心模块来解决这些问题。首先，为桥接 LDR 观测与 HDR 辐照度之间的非线性关系，引入可学习色调映射模块 $T$，将渲染 HDR 辐照度 $E$ 映射到 LDR 域：$C = T(E; \theta) = g_\theta(w(E))$，其中 $w$ 是逐通道白平衡校正（参数 $\theta_w = [w_r, w_g, w_b]^T \in \mathbb{R}^3$），$g_\theta$ 是用 256 个均匀采样点参数化的分段线性相机响应函数（CRF）。为了在极端曝光下稳定优化，CRF 采用泄漏阈值化策略：

$$g_{leaky}(x) = \begin{cases} x, & x < 0 \\ g(x), & 0 \leq x \leq 1 \\ \beta x + (1-\beta), & x > 1 \end{cases}$$

其中 $\beta$ 是阈值化系数，在 $[0,1]$ 之外的值仍保留梯度信号而非直接截断，减少饱和区域的信息损失。同时通过二阶导数惩罚 $L_{smooth} = \sum_{i=1}^{N}\sum_{e \in [0,1]} \ddot{g}_i(e)$ 约束 CRF 保持物理合理性。

其次，针对交替曝光导致传统光流失效的问题，HDR-NSFF 的关键洞察是：虽然像素级颜色随曝光剧烈变化，但物体语义特征保持不变。因此，该方法用 DINOv2 特征替代颜色匹配，采用 DINO-Tracker 作为运动估计骨干网络。为避免长序列跟踪误差累积，每个时间步重新初始化跟踪点，估计双向相邻帧光流；并结合 SAM2 运动掩码将跟踪限定在动态区域，过滤背景噪声。这种语义光流在极端曝光变化下仍能提供稳定一致的运动线索。

第三，单目交替曝光视频的观测信息严重不足（每个时刻只有一个视角，且经常因饱和丢失大面积信息），该方法引入生成先验作为正则化器来弥补信息缺口。具体做法是周期性渲染候选新视角，用生成模型 $G$ 以参考训练视角为条件进行增强：$C^{gen} = G(v, C^{ref})$，将增强结果作为伪观测施加 patch 级感知损失 $\mathcal{L}_{gen} = \sum_{p \in \mathcal{P}} \|\phi(I_p) - \phi(C^{gen}_p)\|_1$。为防止生成先验引入幻觉，该损失仅在预热期 $T_{warm}$ 之后以概率 $p_{gen}$ = 0.1 激活（$T_{warm}$ = 200K 步），确保底层几何已初步建立后才施加先验约束。

总体目标函数整合了色调映射后的光度损失、语义光流损失、深度先验损失、CRF 平滑损失和生成先验损失：$\mathcal{L} = \mathcal{L}_{photo} + \mathcal{L}_{warp} + \lambda_\alpha \mathcal{L}_{flow} + \lambda_\beta \mathcal{L}_{depth} + \lambda_\gamma \mathcal{L}_{smooth} + \mathcal{L}_{gen}$。该方法同时发布了首个面向动态 HDR 场景的真实世界 HDR-GoPro 数据集，使用 9 台同步 GoPro 相机按三档曝光分组采集，包含 12 个含复杂运动的室内外场景。

\paragraph{3D 接地视频稳定化}

传统视频稳定方法在 2D 域建模相机运动（仿射变换、单应矩阵或光流），逐帧独立变换后不可避免地裁剪画面、丢失内容，且缺乏显式 3D 建模容易产生扭曲伪影。GaVS（2025.06）~\cite{2506_23957} 的核心思路是：既然源帧（不稳定）和目标帧（稳定）的相机位姿只差一个微小变换，完全没必要全局重建整个动态场景。因此该方法将视频稳定重新表述为时间一致的"局部重建与渲染"问题——用前馈网络为每一帧预测一组逐像素 GS 基元，再在平滑后的目标位姿上渲染。

具体地，给定源帧 $I_k^s$ 和对应局部重建 $\mathcal{G}_k = f(I_k^s)$（$f$ 为增强深度估计器的 GS 预测头），目标帧通过在平滑位姿 $P_k^t$ 上渲染得到：$\hat{I}_k^t = R(\mathcal{G}_k, P_k^t)$。这种逐帧前馈设计自然统一了静态和动态场景区域的处理，无需像全局 4D 优化那样显式分离动静部分。然而，逐帧预测无法保证相邻重建的时间一致性，且预训练模型在开放域视频上存在泛化问题。

为此，GaVS 引入测试时优化方案，从三个维度约束每个不稳定视频的局部重建质量。第一，多视角动态感知光度监督：将 $\mathcal{G}_k$ 渲染到时间邻域帧 $[k-w, k+w]$ 上，对动态物体通过光流将其扭曲到源帧时间戳的"规范"姿态以补偿对极几何违反，光度损失为 $\mathcal{L}_{RGB} = \frac{1}{|P_k|}\sum_{i \in P_k}\|R(\mathcal{G}_k, P_i^s) - I_i^s\|$，融合 L1 与 SSIM 保持锐度。第二，膨胀时间窗口内的跨帧正则化：通过光流建立帧间密集对应，对匹配的 GS 基元施加属性一致性约束，确保相邻局部重建在重叠区域产生一致的几何和外观。第三，场景外推模块：在源帧边界施加掩码，用视频补全方法 ProPainter 沿光流传播跨帧内容填充边界外区域，再对填充帧预测 GS 基元，从而实现全帧稳定化而不裁剪画面。此外，GaVS 还引入尺度正则 $\mathcal{L}_{scale}$ 惩罚过大的 GS 基元（阈值自适应计算为 $\tau = 70 \cdot \bar{d}/w$，其中 $\bar{d}$ 为先验单目深度、$w$ 为图像宽度），以及偏移正则 $\mathcal{L}_{offset}$ 约束渲染深度与先验深度一致（阈值 $0.2 \cdot \bar{d}$），防止基元漂移。

\paragraph{生成式 4D 场景高斯}

现有视频到 4D 的方法大多假设场景中只有一个物体或物体间遮挡简单，难以处理真实多物体视频中严重遮挡和快速运动带来的挑战。GenMOJO（2025.06）~\cite{2506_12716} 的关键洞察是：场景级渲染损失和物体级生成先验分别擅长不同的事情——前者精确建模跨物体遮挡关系，后者能补全未观察到的物体背面几何。因此该方法提出一种"联合 splatting + 独立 SDS"的组合优化策略。

在表示层面，GenMOJO 用 SAM2 等工具将单目视频分解为 $N$ 个前景物体和 1 个背景，每个物体维护一组独立的可变形 3D 高斯 $\mathcal{G}^{obj}_0$（初始帧的规范状态），通过变形网络 $D_\theta(\mathcal{G}^{obj}_0, t)$ 预测每帧的 10 维变形向量（3D 位移 $\Delta x_t$、四元数旋转 $R_t$、3D 缩放 $s_t$）。变形网络采用 Hexplane 骨干加 2 层 MLP 头，Hexplane 分辨率设为 $[64, 64, 64, 0.8T]$（$T$ 为帧数）。背景高斯位置固定但颜色和透明度可训练。

优化框架的巧妙之处在于双坐标系桥接——场景级联合渲染精确捕获跨物体遮挡，物体级独立 SDS 损失补全未观察到的背面几何。每个物体在自身坐标系下初始化和优化，通过一组可微仿射变换 $W_t^{obj}$ 将物体高斯映射到场景坐标系，使所有物体的高斯可以在同一空间联合 splatting：$\mathcal{L}_{render} = \|R(\bigcup_{obj} W_t^{obj}(\mathcal{G}_t^{obj}), P_{cam}) - I_t\|$。这精确捕获了跨物体遮挡——渲染时近处物体自然遮挡远处物体。同时，每个物体独立渲染到未观察视角，用 Zero-1-to-3 扩散模型计算 SDS 损失 $\mathcal{L}_{SDS}^{obj}$ 补全背面几何。总损失整合渲染损失、SDS 损失、光流损失和深度损失。实验中静态优化 1000 步（SDS 批大小 16），动态优化 $35T$ 步（SDS 每帧采样 4 个新视角，批大小 8），Hexplane 学习率 $6.4 \times 10^{-4}$、MLP 头学习率 $6.4 \times 10^{-3}$，32 帧 3 物体场景在 A6000 上约 74 分钟完成。

\paragraph{稀疏控制高斯-网格混合 4D 生成}

现有基于动态 NeRF 或动态高斯的视频到 4D 方法缺少对物体表面的紧约束，优化空间冗余，阻碍变形学习。DreamMesh4D（2024.10）~\cite{2410_06756} 从现代 3D 动画管线获得启发，采用 SuGaR 高斯-网格混合表示：将扁平高斯通过重心坐标绑定到三角网格面上（每面 $\kappa=6$ 个），使 2D 渲染对网格顶点位置和高斯属性都可微。该方法分静态和动态两阶段优化。

静态阶段先用 Zero123 SDS 优化 3000 步生成粗糙高斯场（前 1500 步做致密化和剪枝），然后引入 SuGaR 正则使高斯对齐物体表面，剪除不透明度低于 $\epsilon=0.5$ 的高斯后用 Poisson 重建提取粗糙网格。在网格上绑定新的扁平高斯后，联合优化顶点和表面高斯属性 2000 步，损失权重 $\lambda_{SDS}^s=0.1$、$\lambda_{ref}^s=1000$、$\lambda_{mask}^s=500$。

动态阶段是该方法的精髓。在精化后的网格表面上按测地距离均匀采样 $N_{node}=1024$ 个稀疏控制节点，每个顶点连接 $N_{neighbor}=4$ 个最近邻节点，构建变形图。变形网络为每个控制节点预测局部变换（形变矩阵 $F_p \in \mathbb{R}^{3 \times 3}$ 和平移 $t_p \in \mathbb{R}^3$），然后通过 DreamMesh4D 提出的自适应混合蒙皮算法驱动所有网格顶点和表面高斯变形。该算法结合了 LBS 和 DQS 的优势：LBS 通过线性混合 $v^{lbs} = \sum_{p \in N(v)} w_p(F_p v + t_p)$ 产生自然动画但有"体积塌缩"伪影，DQS 用单位对偶四元数 $\hat{q}_p = DQ(R_p, t_p)$ 消除该伪影但无法处理非刚性变形。自适应混合蒙皮根据局部变形的刚性程度自动选择两者——刚性区域用 DQS 保持体积，非刚性区域用 LBS 保留弹性。动态优化 2000 步（学习率 0.00032），每步随机采样 8 帧、每帧 2 个视角做 SDS 监督，配合 ARAP 正则和法线一致性约束防止几何退化。

\paragraph{扩散增强 4D 高斯时序一致性}

从单目视频生成 4D 内容面临两个核心矛盾：单视角信息不足以推断完整的 3D 结构，而多视角扩散模型生成的图像序列又存在视角间不一致。Splat4D（2025.08）~\cite{2508_07557} 设计了一条三阶段管道来逐步化解这些矛盾。

第一阶段是粗糙 4D 高斯场的快速构建。Splat4D 先用视频扩散模型生成单视角视频，再用 MV-Adapter 从每帧生成正面、左侧、右侧、背面四个视角的多视角图像序列：$\text{MV-Adapter}(I_t) \to \{I_t, I_t^{left}, I_t^{right}, I_t^{back}\}$。由于 MV-Adapter 生成的图像缺乏细节（输入总落在训练分布之外），Splat4D 引入 ESRGAN 图像增强器 $IE$ 对每帧每视角精化纹理和边缘。增强后的多视角序列送入预训练的 LGM 前馈模型：LGM 用非对称 U-Net 将多视角图像编码为多尺度空间特征，解码为 Splatter Image（逐像素编码高斯参数的位置、大小、颜色、不透明度），再通过相机参数反投影到 3D 空间，快速获得粗糙但完整的 4D 高斯场。

第二阶段是时空一致性的识别与修复。直接使用增强器会引入多视角不一致和错位，因为不同视角的增强结果可能产生矛盾纹理。Splat4D 的解决方案是"不确定性掩码 + 视频扩散精化"：通过渲染粗糙高斯场的多视角图像并与增强结果比较，识别出不可靠区域生成不确定性掩码，然后用视频扩散模型对掩码区域做 inpainting，在保持时空一致性的同时补全缺失内容。推理时使用 50 步 Flow Match Euler 离散调度器。

第三阶段是通用化 3D 高斯场预测器训练。在 Objaverse 80K 子集上微调非对称 U-Net（4 块 V100，批大小 4/卡，bfloat16），每批次随机采样 1 个相机视角、用 MV-Adapter 生成 4 个正交视角，U-Net 预测的 3D 高斯场渲染后与 Objaverse 原始渲染图在 $512 \times 512$ 分辨率下用 MSE 损失监督（LPIPS 在 $256 \times 256$ 计算），AdamW 优化器（学习率 $4 \times 10^{-4}$，余弦退火），50\% 概率施加网格扭曲和相机抖动增强泛化性。

\paragraph{跟踪引导 4D 生成}

Animate3D 等方法用多视角视频扩散模型（MV-VDM）生成动画视频再重建 4D 资产，但仅依赖像素/潜空间中的视频扩散损失，缺少特征级的时间感知监督，导致长时间生成中出现外观漂移（物体逐帧退化或改变）。Track4DGen（2025.12）~\cite{2512_06158} 的核心思想是利用基础点跟踪器 CoTracker3 的密集帧间对应关系作为"运动专家"，将其知识注入扩散特征空间，从而在生成阶段就建立强时间一致性。

第一阶段升级 MV-VDM。Track4DGen 在标准 U-Net 去噪架构上，通过系统性的逐块探测发现：U-Net 解码器中第二个时空上采样块的时间运动模块的特征具有最强的长时间对应能力。因此在每个视角的第一帧图像上叠加 $15 \times 15$ 网格、排除物体外区域后，随机选取 8 个查询点/视角，用 CoTracker3 跟踪其跨帧轨迹，然后将跟踪点映射到该特征空间的对应位置并提取描述子。对应跟踪损失强制特征级点对应一致：$\mathcal{L}_{corr} = \frac{1}{|Q|}\sum_{q \in Q} \sum_{t} [1 - \cos(f_t(p_q^t), f_1(p_q^1))]$，其中 $f_t$ 为第 $t$ 帧的扩散特征、$p_q^t$ 为查询点 $q$ 在帧 $t$ 的 2D 位置。多视角扩散训练用 AdamW（学习率 $10^{-6}$）在 16 块 H100 上训练 20 epoch，批大小 1024（$16 \times 4$ 视角 $\times 16$ 帧），推理时 25 步采样。

第二阶段混合 4D-GS 重建。经典 4D-GS 仅用 Hex-plane 特征表示运动场，Track4DGen 将第一阶段训练的扩散特征（已编码跟踪先验）与 Hex-plane 特征在共置点上拼接，形成混合运动编码。Hex-plane 分辨率 $[100, 100, 8]$、特征维度 16，先用光度/重建损失进行 750 次运动重建（批大小 64），再启用 4D-SDS 和 ARAP 正则额外 250 次迭代。此外，Track4DGen 引入 4D 球谐（SH）建模，使高斯的颜色系数随时间变化，从而捕捉动态光照效果（如旋转物体表面的高光流动），进一步提升动态建模保真度。每个物体的 4D 优化在单块 A100 上约 40 分钟。

\paragraph{异步视频扩散传播 3D 几何}

长范围视角外推（即相机大幅超出原始轨迹生成新内容）面临一个根本矛盾：纯相机位姿编码方法缺乏 3D 场景信息导致 OOD 泛化差，而显式 3D 先验（点云/3DGS 重建）因遮挡和深度估计误差导致变形图像存在空洞和扭曲，标准 inpainting 无法同时处理这两种退化。WorldWarp（2025.12）~\cite{2512_19678} 的核心洞察是：在相机条件设定下，未来帧的几何先验可以通过前向变形一次性获得，因此无需因果/自回归约束，可以采用更强大的双向注意力机制。

该方法的关键设计是"时空变化噪声调度"。训练时，从源帧 $x_s$ 和深度图 $D_s$ 反投影为 3D 点云 $P_s$，再用可微点渲染器投影到所有目标视角，得到变形先验序列 $X_{s \to V} = \{x_{s \to t}\}_{t=1}^T$ 和对应有效掩码 $M = \{M_t\}_{t=1}^T$。在潜空间中构建复合序列：有效区域取变形潜变量、空洞区域取真值潜变量，$z_{c,t} = M_{latent,t} \cdot z_{s \to t} + (1 - M_{latent,t}) \cdot z_t$。噪声调度在时间和空间两个维度同时变化：时间上每帧独立采样噪声水平，空间上对"变形区域"和"填充区域"施加不同噪声——$\sigma_t = M_{latent,t} \cdot \sigma_{warped,t} + (1 - M_{latent,t}) \cdot \sigma_{filled,t}$，即空白区域接收全噪声从零开始生成，变形区域接收部分噪声在已有结构上精化。

推理阶段以自回归方式逐块（chunk）生成视频（每块 49 帧，上下文重叠 5 帧）。核心是在线 3D 几何缓存机制：每生成一个新块，就在最近的高质量历史帧上用 3DGS 优化 500 步构建局部几何缓存（学习率 $1.6 \times 10^{-3}$），再从该缓存渲染下一个块所需的变形先验。与使用静态 RGB 点云相比，在线优化的 3DGS 缓存提供更高质量的变形图像（消融实验中 PSNR 从 11.12 提升到 17.13），有效防止了静态 3D 表示中不可逆的误差累积。强度参数 $\beta=0.8$ 控制变形区域的初始噪声水平，使其保留结构完整性。整个框架基于 Wan2.1-T2V-1.3B 微调，在 8 块 H200 上训练 10K 步。

\paragraph{推理时长视频 4D 生成}

预训练视频扩散模型通常只能在训练长度内（如 $F$ 帧）生成高质量片段，直接超出该长度会导致质量急剧下降。FlowLong（2026.05）~\cite{2605_20910} 的核心动机是将长视频生成重新表述为一个逆问题：将多个重叠片段的采样轨迹对齐到一致序列，而不需要重新训练模型。具体而言，将长视频分为 $K$ 个重叠片段，每个片段独立调用同一个预训练模型，相邻片段共享 $O$ 帧重叠区域。设 $M_k, M'_{k+1}$ 为从第 $k$ 和第 $k+1$ 个片段中提取重叠帧的选择矩阵，则理想的重叠约束为 $M_k \hat{\mathbf{x}}_0^{(k)} = M'_{k+1} \hat{\mathbf{x}}_0^{(k+1)}$，其中 $\hat{\mathbf{x}}_0$ 是通过 Tweedie 公式估计的干净样本。FlowLong 将此硬约束松弛为采样引导损失：$\mathcal{L}_k = \frac{1}{2} \| M_k \hat{\mathbf{x}}_0^{(k)} - M'_{k+1} \hat{\mathbf{x}}_{0|t}^{(k+1)} \|^2$，即在每步去噪时通过梯度修正各片段的潜变量，使重叠区域在数据流形上对齐。在重叠帧上采用线性混合调度 $\lambda_j = (j - (F-O))/(O-1)$，左端完全取自片段 $k$、右端完全取自片段 $k+1$，消除边界跳变。为了防止早期轨迹发散，采样前半段注入随机噪声（SDE 阶段）使各片段充分混合，后半段切换为确定性 ODE 采样保持视觉保真度。以 LTX-2 为骨干时实际配置为 $(F, O, S) = (16, 8, 7)$，即每窗口 16 个潜变量帧、8 帧重叠、步长 7。该方法还扩展到文本到 3DGS 场景生成：利用长视频扩散模型生成多视角一致的视频序列，再通过 3DGS 重建出 4D 场景。

\paragraph{4D 高斯光线追踪物理相机效果}

当前视频生成模型无法理解相机效果参数（如鱼眼畸变、滚动快门），生成的视频经常违反物理规律；而传统 3DGS 的光栅化渲染方式也缺乏模拟这些光学现象的能力。4D-GRT（2025.09）~\cite{2509_10759} 的设计动机正是填补这一缺口：用基于物理的光线追踪替代光栅化——直觉上，光栅化像把高斯"拍扁"到屏幕上，而光线追踪则是从眼睛发出真正的光线与高斯求交，因此能自然模拟鱼眼、景深等光学效果，从而在 4D 动态场景上精确模拟真实相机效果。管道分为两阶段。第一阶段为动态场景重建：沿用 4D-GS 的形变场网络，通过六组多分辨率体素平面 $R_l$ 提取时空特征 $\mathbf{f}_h = \bigoplus_l \text{interp}(R_l(i,j))$，经 MLP $\psi$ 融合后由三个独立 MLP $(\phi_x, \phi_r, \phi_s)$ 预测位置、旋转和缩放的残差，得到变形后的高斯 $G_t$。关键改进在于渲染环节：弃用 3DGS 标准的光栅化（$\alpha$ compositing），转而采用 3DGRT 的可微光线追踪框架，直接在 3D 高斯基元上追踪光线，利用 NVIDIA OptiX 硬件加速和 $k$-buffer 命中检测方案实现高效渲染。

训练损失为 $\mathcal{L} = \mathcal{L}_1(C_{v,t}, \hat{C}_{v,t}) + \mathcal{L}_{TV}$，其中 $\mathcal{L}_{TV}$ 为网格化总变差正则化。第二阶段为相机效果渲染：在重建好的 4D 场景上，通过物理相机模型生成三种效果。（1）鱼眼畸变：采用四次多项式径向畸变模型 $\theta = k_0 + k_1 r + k_2 r^2 + k_3 r^3 + k_4 r^4$，其中 $r$ 为像主点到像素的物理距离，$\theta$ 为入射光线与光轴的极角，由系数 $k_i$ 控制畸变程度。（2）景深：给定对焦距离 $f_z$ 和光圈半径 $r_a$，将理想针孔光线与焦平面的交点 $\mathbf{p} = \mathbf{o} + f_z \mathbf{d}$ 作为锚点，然后在圆形光圈上随机抖动光线起点，多次追踪取平均以模拟散焦模糊。（3）滚动快门：将帧拆分为 $N_c$ 个时间块，每块对应不同的读出时刻，各自用对应时刻的变形高斯 $G_t$ 渲染后拼接。在 8 个动态室内场景的基准上，4D-GRT 在所有相机效果类型中均实现最快渲染速度，PSNR 与动态 NeRF 方法（HexPlane、MSTH）持平或更优。

\paragraph{回溯式动态场景新视角合成}

现有动态 3DGS 方法（如 4DGS、Dynamic 3DGS）普遍采用形变场或规范模板来强制帧间时间一致性，但在同步多视角体育场景中，每个时刻的多视角几何已经通过标定和校准提供了强约束，这种复杂的时间耦合是否真的必要？Zhang 等人（2026.05）~\cite{2605_12437} 正是从这一几何视角出发，提出一种无需显式时间形变的暖启动 3DGS 管道。核心思路是"逐帧独立优化、帧间传播初始化"。具体而言，场景在每一时刻 $t$ 表示为固定数量 $K$ 个各向异性 3D 高斯 $G_t = \{g_{t,k}\}_{k=1}^K$，每个高斯参数化为位置 $\boldsymbol{\mu}_{t,k}$、协方差 $\Sigma_{t,k} = R(q_{t,k}) \text{diag}(\mathbf{s}_{t,k}^2) R(q_{t,k})^\top$、不透明度 $\alpha_{t,k}$ 和球谐系数 $b_{t,k}$。关键设计是禁用高斯致密化，使所有时刻保持 $|G_t| = K$，从而保证可预测的内存占用和稳定的渲染吞吐量。初始化仅在起始时刻 $t_0$ 使用 SfM 点云 $P_{t_0} = \{p_k \in \mathbb{R}^3\}_{k=1}^K$，后续时刻 $t$ 的优化直接以 $G_{t-1}$ 的参数作为暖启动初始值。作者通过消融实验验证了暖启动的有效性来自相邻帧的局部时间平滑性（而非目标函数中的隐式时间耦合）：当用远距离帧初始化时，PSNR 从 42.50 dB 骤降至 26.80 dB（下降 15.70 dB），证明暖启动利用的是"相邻帧之间场景变化小"这一先验。此外，作者还提出基于 Blender 的标准化动态多视角数据集框架，自动生成参数化同步相机阵列（半球、球面、环形和体育场布局），统一坐标约定为 OpenCV/COLMAP 格式，消除不同 NeRF/3DGS 代码库之间格式转换的摩擦。在 D-WS、S-PK、S-MP 三个体育场景基准上，该方法在 1920$\times$1080 分辨率达到 55--66 FPS 实时渲染，同时保持高保真回溯 NVS 质量。

\paragraph{混合高斯静态动态分解}

现有动态 3DGS 方法的共同问题是：对所有高斯基元一视同仁地分配时间参数，导致静态区域也携带大量冗余的时变参数，既浪费存储又容易引入时序伪影。HGS（2025.12）~\cite{2512_14352} 的核心动机正是解决这一参数分配不均衡问题，通过静态-动态分解（SDD）策略显式区分两类基元并采用完全不同的参数化方案。首先，利用 Track Anything Model（TAM）对输入视频进行动静分割，生成二值掩码将高斯基元分为静态和动态两组。对于动态基元，采用时间径向基函数（RBF）显式参数化其时间演化：位置通过三阶多项式插值 $\boldsymbol{\mu}_i(t) = \sum_{k=0}^{3} b_{i,k}(t - t_i^*)^k$，旋转通过一阶多项式 $r_i(t) = \sum_{k=0}^{1} c_{i,k}(t - t_i^*)^k$，不透明度采用高斯形时间窗 $\alpha_i(t) = \alpha_i^* \exp(-s_i^* |t - t_i^*|^2)$，其中 $t_i^*$ 为时间中心，$b_{i,k}$ 和 $c_{i,k}$ 为可训练多项式系数。这种显式参数化避免了形变场等隐式运动建模的计算开销。对于静态基元，所有高斯共享一组时间不变参数（位置、旋转、不透明度、颜色），直接消除了静态区域的时间冗余。初始化采用交错帧策略：仅在时间子采样帧上进行 SfM 重建，限制初始高斯数量以平衡效率与保真度。训练采用两阶段策略：第一阶段仅训练静态分支建立稳定背景；第二阶段联合优化动静分支，抑制边界伪影并改善收敛稳定性。最终模型仅需 6.87 MB 存储，相比 STGS 减少 63\%，相比 MixVoxels 减少高达 98\%，在 RTX 3090 上以 1352$\times$1014 分辨率达到超过 300 FPS，在 RTX 3050 上仍维持 160 FPS，已成功集成到 VR 系统中。

\begin{takeaway}
视频重建与时序一致性面向自由视角视频的实际应用需求，核心挑战是在保持实时渲染的同时压缩存储并保证时序一致。HGS 的静态-动态分解策略将存储减少 63\%--98\%，4D-GRT 用物理光线追踪替代光栅化实现鱼眼和滚动快门等相机效果模拟，暖启动管道在体育场景中实现 55--66 FPS 实时渲染。
\end{takeaway}

%% ============================================================


\section{动态物体与场景分解}
\label{sec:06_6}

前面的章节解决了动态场景的表示、压缩和渲染问题，但都隐含假设场景可以被统一建模。现实中，动态物体和静态背景往往需要完全不同的处理策略——这正是本节要解决的核心问题。

\subsection{问题背景：从混乱中提取秩序}

现实场景极少是纯粹静态的：街道上有行人穿行，房间内有人移动物品，旅游照片中充斥着其他游客。标准 3DGS 假设场景在所有输入图像中保持一致，动态物体（也称干扰物，distractors）会违反这一假设，导致浮动伪影或视角相关效果。更深层次的挑战是将场景分解为有意义的语义实体——静态背景、动态前景物体、不同实例——以支持编辑、跟踪和场景理解。本节涵盖 22 篇论文，围绕静态背景重建、动静分离、实例分解和语义 4D 理解展开。

\subsection{静态背景重建与干扰物去除}

本节从静态场景出发，介绍六种在存在动态干扰物的情况下重建干净静态背景的方法。随后本节介绍两条技术路线：（1）动静分离与运动分解；（2）实例分解与语义4D理解。

\paragraph{去耦动静高斯泼溅}

标准 3DGS 将整个场景建模为单一高斯集合，动态物体（行人、手部交互等）违反了多视角一致性假设，产生浮动伪影。现有解决方案存在两类互补的缺陷：一类是基于掩码抑制的方法（如 SpotlessSplats、WildGaussians），利用残差和语义特征识别瞬态区域并从训练中排除，但依赖精心初始化和稳定优化，在 egocentric 等高动态长视频中因持续人体-场景交互、严重运动模糊和快速光照变化而容易过度或不足分割；另一类是基于显式动静分支的 NeRF 方法（如 NeRF-W、NeuralDiff），在先验稀疏的短视频中动态分支欠分割留下浮影，在高动态 egocentric 视频中又过分割导致静态背景欠表示。DeGauss（2025.03）~\cite{2503_13176} 的核心洞察是：动态高斯（如 4DGS 的形变场）擅长建模时变内容但对新视角泛化差，静态高斯擅长跨视角一致但将动态内容表现为视角相关伪影（即浮影）——两者恰好互补。因此 DeGauss 用两组完全独立的高斯分支分别建模前景动态和背景静态，并通过前景高斯自身渲染出的概率掩码协调两者组合，实现自监督的动静分解。

具体而言，前景高斯 $G_f$ 除了标准属性（位置、旋转、缩放、不透明度、颜色）外，额外附加三个可学习掩码属性：$m_f$（前景概率）、$m_b$（背景概率）和 $b$（亮度控制）。前景高斯通过 Hexplane 编码器 $H$ 和多头解码器 $D$ 实现时空变形：$f_d = H(G_f, t)$，然后 $\Delta G_f = D(f_d)$，其中 $D = \{\phi_x, \phi_r, \phi_s, \phi_\alpha, \phi_c, \phi_{m_f}, \phi_{m_b}, \phi_b\}$ 为独立的 MLP 解码器组，各自预测对应属性的残差。变形后的掩码属性通过标准 $\alpha$ compositing 光栅化得到原始前景和背景概率图：$M_f = \sum_{i=1}^N m_{f,i}^\Delta \alpha_i^\Delta P_i(G_{f,i}^\Delta) \prod_{j=1}^{i-1}(1 - \alpha_j^\Delta P_j(G_{f,j}^\Delta))$，$M_b$ 同理。归一化后 $P_f = M_f / (M_f + M_b + \epsilon)$，$P_b = M_b / (M_f + M_b + \epsilon)$。这一概率公式自然地将掩码推向 0 或 1（避免中间值），产生干净的动静分离。值得强调的是，这一自监督收敛机制之所以有效，关键在于 DeGauss 的解耦渲染策略：前景和背景高斯分别独立渲染得到 $C_f$ 和 $C_b$，\emph{渲染之后}再用概率掩码合成为 $\hat{C} = P_f \cdot C_f + P_b \cdot \tilde{B} \cdot C_b$。相比之下，NeRF-based 方法沿射线混合动静密度和辐射度（渲染中组合），容易导致训练中射线的提前终止和局部最优。解耦设计保证两个分支各自享有完整的梯度回传路径，组合掩码在训练中逐步形成而非被早期收敛锁定。

为处理光照变化导致的动静歧义（强光区域容易被前景分支过度捕获），引入亮度控制掩码 $B$ 和分段线性激活函数：$\tilde{B} = B + 0.5$（$0 \le B \le 0.75$），$\tilde{B} = k(B - 0.75) + 1.25$（$0.75 < B \le 1$），其中 $k=35$ 为过亮度建模系数，增强背景分支对非朗伯效应（如镜面反射、镜头眩光）的建模能力。前景分支中还自然涌现出一类"工具高斯"（utility gaussians）——它们贡献概率掩码 $P_f$、$P_b$ 和亮度控制掩码 $\tilde{B}$，但不贡献前景颜色渲染 $C_f$。这一额外自由度有效减少了无约束高斯运动导致的伪影，同时为掩码提供了更灵活的建模空间。训练后的错误建模元素在密度控制过程中被逐步剪枝，最终产生精确且干净的动静分离。

训练采用两阶段策略。粗训练阶段（coarse stage）禁用前景分支的变形模块，仅训练静态属性，短视频/图像集训练 1000 次迭代，长视频训练与序列长度等量的迭代，目的是建立稳定的背景表示。精训练阶段（fine stage）启用变形模块，端到端联合优化两个分支，短视频训练 20K 次迭代，长视频训练 120K 次迭代。损失函数也做了专门设计：主损失 $\mathcal{L}_{\text{main}} = \lambda_1 \mathcal{L}_1 + \mathcal{L}_{\text{SSIM}} + \mathcal{L}_p + \mathcal{L}_{\text{TV}} + \mathcal{L}_{m_f} + \mathcal{L}_{m_b}$ 和辅助损失 $\mathcal{L}_{\text{uti}} = \mathcal{L}_{\text{SSIM}} + \mathcal{L}_\alpha + \mathcal{L}_B + \mathcal{L}_b$ 分离——两者都参与参数更新，但只有 $\mathcal{L}_{\text{main}}$ 的梯度参与高斯致密化，防止前景高斯因 SSIM 梯度过大而过度扩张到静态区域。此外，每隔一定间隔对前景和背景高斯执行 50\% 的部分不透明度重置（partial opacity reset），这在标准 3DGS 的全量重置会导致训练不稳定的方法（如 SpotlessSplats）中不可行，但 DeGauss 的前景-背景公式化提供了天然稳定性保障，使部分重置既能有效控制高斯密度又能跳出局部最优。

初始化方面，背景高斯从 COLMAP（或传感器感知）生成的点云初始化，前景高斯在场景边界内随机生成。DeGauss 从随拍图像集到长动态 egocentric 视频（2800--5000 帧）均可泛化，在 NeRF-on-the-go、ADT、AEA、Hot3D 和 EPIC-Fields 等基准上一致超越现有方法，特别在 LPIPS 指标上显著优于前 SOTA 方法 SpotlessSplats。在 Neu3D 和 HyperNeRF 等动态场景建模基准上，DeGauss 的组合渲染也持续取得更优的 LPIPS 分数，同时相比 4DGS 展现出更好的测试时渲染效率和更精细的重建细节。

\paragraph{动态感知静态场景重建}

从动态视频中重建静态背景面临一个"鸡生蛋"问题：相机位姿估计依赖准确的场景几何，但动态物体的存在又干扰几何重建。现有方法通常依赖极线误差阈值或单帧语义分割，前者需要精确光流且对阈值敏感，后者假设动态物体即语义前景（实际未必成立）。DAS3R（2024.12）~\cite{2412_19584} 的核心思路是用学习-based 方法直接从图像对预测动态掩码，而非依赖手工阈值或单帧语义——直觉上，人类看到两张不同时间的照片就能判断哪些物体在运动，DAS3R 试图让神经网络也具备这种能力。具体而言，以 MASt3R 为骨干网络，输入两帧存在重叠但时间不同的 RGB 图像 $I_n, I_m$，通过 ViT 编码器和交叉注意力机制在潜空间中分解相机运动与物体运动，直接回归密集 3D 点图和置信度图，同时预测动态掩码 $\mathbf{D} \in \mathbb{R}^{W \times H \times 1}$，标记每个像素是否属于动态物体。

模型在 Point Odyssey 合成数据集上用真实动态掩码监督训练 50 个 epoch（每 epoch 采样 20000 对图像），采用 AdamW 优化器、学习率 $5 \times 10^{-5}$，能泛化到未见过的 DAVIS 和 Sintel 数据集。逐对预测的掩码随后聚合为视频级动态掩码：将所有选定图像对的预测结果组合，并通过与 MASt3R 全局对齐管道相同的后处理流程生成密集预测和相机位姿。关键设计是将全局对齐中的静态置信度替换为模型预测的动态掩码，从而提升场景深度重建和位姿估计精度。在 Gaussian Splatting 训练阶段，DAS3R 引入动态感知优化：将预测掩码作为可优化参数，掩码区域内的像素不参与静态高斯的损失计算，确保只有静态场景的光度误差回传到高斯参数。整个管道无需相机位姿输入、无需 SLAM 点云、无需光流阈值调参。在动态物体占据场景显著比例的 DAVIS 和 Sintel 数据集上，DAS3R 相比次优方法 PSNR 提升超过 2 dB，且仅需 4000 次迭代（其他方法至少需要 30000 次）。

\paragraph{干扰物鲁棒新视角合成}

标准 3DGS 的训练损失对所有像素一视同仁，当场景中存在行人、气球等干扰物时，干扰物区域的渲染误差会持续向场景表示注入错误梯度，导致浮动伪影。Robust 3D Gaussian Splatting（2024.08）~\cite{2408_11697} 的核心思路是：在优化过程中自动生成干扰物掩码，将其从损失计算中剔除——直觉上，如果某个区域在训练中始终产生高误差，那它很可能是动态物体而非重建缺陷。

具体而言，该方法首先从每次迭代的渲染残差 $R^{(i)} = |\hat{I}^{(i)} - I^{(i)}|$ 出发，逐通道计算三类掩码特征：原始残差（定位高误差像素）、经 $3 \times 3$ 方框核平滑的残差 $\mu_1$（捕获局部连续性）、以及 $16 \times 16$ 邻域均值 $\mu_2$（捕获干扰物的连通区域），三者叠加得到聚合特征 $\mu_3 = \mu + \mu_1 + \mu_2$。与 RobustNeRF 使用固定阈值不同，Robust 3DGS 在此基础上训练一个逻辑回归分类器 $m(\mathbf{p}) = \sigma(\mathbf{W}^\top \mu_3(\mathbf{p}) + b)$，动态学习干扰物与非干扰物的决策边界。逻辑回归的优势在于它能随训练进程自适应调整阈值——训练早期模型本身误差较大，分类器会将阈值拉高以避免误杀正常区域。

然而，原始掩码仅基于残差，可能将干扰物的一部分误判为静态内容。为解决此问题，Robust 3DGS 引入物体感知模块：利用 SegmentAnything 获取每张图像的物体分割集合 $\mathcal{M}(I)$，将分割结果与原始掩码取交集——只有当某个物体分割区域与掩码的重叠率超过 40\% 时，才将该物体标记为干扰物。最终损失公式为 $L_{\text{mask}} = L_{\text{GS}}(\hat{I} \odot m^c, I \odot m^c) + \lambda \sum_{i,j} m_{i,j}$，其中 $m^c = 1 - m$ 为互补掩码，正则项 $\lambda = 0.1$ 防止分类器将所有像素标记为干扰物这一平凡解。相比标准 3DGS，PSNR 提升 1.86 dB；相比直接将 RobustNeRF 的鲁棒损失移植到 3DGS 的基线，提升 4.3 dB。

\paragraph{生成辅助静态场景重建}

当动态物体长时间大面积遮挡静态背景时，简单地将干扰区域从损失中移除会导致静态背景因缺乏监督而无法重建。GA-GS（2026.04）~\cite{2604_04331} 的思路是：用扩散模型修复被遮挡区域作为伪真值监督，同时引入可靠性感知机制区分真实像素与生成像素的贡献。

该方法首先利用 Flow-SAM（融合 SAM 的分割能力与 RAFT 光流运动线索）生成精确的动态区域掩码 $\mathcal{M} = \{M_i\}_{i=1}^N$，然后用 DiffuEraser 扩散修复模型填充被遮挡区域。修复后的监督图像 $I_i'$ 仅在掩码区域替换为生成内容：$I_i'(\mathbf{y}) = M_i(\mathbf{y}) \cdot \hat{I}_i^{\text{inpaint}}(\mathbf{y}) + (1 - M_i(\mathbf{y})) \cdot I_i(\mathbf{y})$。关键设计在于：每个高斯基元增加一个可学习标量 $\tau \in [0,1]$，称为真实性指示器，编码该高斯来源于真实观测（而非生成内容）的置信度。初始化时，来自原始图像区域的高斯设为 $\tau = 0.9$，来自修复区域的高斯设为 $\tau = 0.1$。渲染时，有效不透明度被调制为 $\alpha' = \tau \cdot \alpha$，使真实区域的高斯主导渲染，修复区域的高斯仅起辅助作用。相应地，损失函数采用加权 L1 损失 $L_{\text{L1}} = \sum_{\mathbf{y}} w_i(\mathbf{y}) |\hat{I}_i(\mathbf{y}) - I_i'(\mathbf{y})|$，其中权重图 $w_i(\mathbf{y}) = 1.0 - (1 - w) \cdot M_i(\mathbf{y})$，真实像素权重为 1.0，生成像素权重为 $w < 1$。为定量评估遮挡区域的重建质量，作者使用固定路径机器人采集成对数据集 Trajectory-Match——同一路径分别录制有/无动态物体的版本，后者作为真值。

\paragraph{2D/3D 高斯解耦瞬态与静态}

标准 3DGS 用统一的 3D 高斯表示整个场景，但瞬态物体（如偶尔出现的行人）不满足多视角一致性假设——它们只在部分图像中出现，硬塞进 3D 表示会导致视角相关的伪影。HybridGS（2024.12）~\cite{2412_03844} 的核心观察是：瞬态物体本质上是单视角的平面现象，因此应当用 2D 高斯而非 3D 高斯来建模。

具体地，HybridGS 将图像分解为静态部分和瞬态部分：$I = M_t \odot I_t + (1 - M_t) \odot I_s$，其中 $M_t \in [0,1]$ 为瞬态掩码。静态部分 $I_s$ 由一组全局共享的 3D 高斯渲染，瞬态部分 $I_t$ 则由每张图像独立的一组 2D 高斯渲染。2D 高斯的掩码值通过其不确定性（uncertainty）自然产生：$M_t(\mathbf{y}_{2d}) = \sum_{i \in \mathcal{N}} \alpha_{2d_i}'$，其中 $\alpha_{2d_i}' = \alpha_{2d_i} \cdot e^{-\frac{1}{2}(\mathbf{y}_{2d} - \mathbf{x}_{2d})^T \Sigma_{2d}^{-1}(\mathbf{y}_{2d} - \mathbf{x}_{2d})}$，即 2D 高斯在像素 $\mathbf{y}_{2d}$ 处的不透明度加权和。这种设计的动机是：在 3D 高斯建模完静态场景后，2D 高斯只需学习残差部分，残差越大的区域自然对应瞬态物体。

为防止 3D 高斯被瞬态物体的 RGB 损失拉偏，HybridGS 提出多视角调节监督：每次迭代随机采样 $K = 4$ 张图像，只在它们的共可见区域（co-visible frustum intersection）反向传播梯度，利用多视角信息帮助区分哪些是静态的（多视角一致）、哪些是瞬态的（仅单视角可见）。训练采用三阶段策略：1k 步预热让 3DGS 建立初始场景表示，然后交替训练 2DGS（10k 步）和 3DGS（1k 步）分别学习瞬态和静态，最后联合训练 30k 步精调。在 NeRF On-the-go 数据集上，平均 PSNR 比前 SOTA 提升 1.10 dB。

\paragraph{无约束照片集合实时新视角合成}

互联网上的旅游照片通常在不同时间、天气和设备条件下拍摄，存在严重的外观变化和瞬态遮挡。标准 3DGS 缺乏外观建模机制，面对这类数据会产生鬼影伪影和几何错误。Wild-GS（2024.06）~\cite{2406_10373} 提出层次化外观分解策略，将每个 3D 高斯在特定视角下的颜色分解为三个组件。

外观分解的核心公式为：$\text{sh}_i = M_\theta^F(\text{Emb}^g \oplus \text{Emb}^l_i \oplus f^{\text{in}}_i)$，其中 $\oplus$ 为拼接操作，$\text{sh}_i$ 为预测的球谐系数。三个组件各有明确语义：（1）全局外观嵌入 $\text{Emb}^g = M_\theta^G(\text{AvgPooling}(F_{I_R}))$，从参考图像经 UNet 编码和全局平均池化提取，编码场景级的照明和色调映射变化；（2）局部外观嵌入 $\text{Emb}^l_i$，通过三平面表示为每个高斯提供位置感知的局部反射率特征；（3）固有特征 $f^{\text{in}}_i$（维度 32），每个高斯自有的材质属性向量，不随视角或外观条件变化。

三平面的构建是关键创新：将参考图像在静态区域的深度和颜色反投影到 3D 空间得到彩色点云 $\{C_{I_R}, P_{I_R}\}$，然后将这些 3D 点投影到三个正交平面生成颜色三平面 $c_{xy}, c_{yz}, c_{zx}$。每个高斯根据自身位置查询三平面即可获得局部外观嵌入，实现显式的外观控制。此外，Wild-GS 利用 2D 可见性图（UNet 预测的瞬态掩码）排除瞬态物体，掩码阈值 $Th$ 将瞬态物体从三平面构建和训练损失中剔除；深度正则化损失 $L_D$ 约束几何以稳定外观建模。相比之前 SOTA 的 CR-NeRF，Wild-GS 在 Phototourism 数据集上 PSNR 提升约 3 dB，训练时间缩短至 1/200，渲染速度快 10000 倍。

\paragraph{野外 3D 高斯泼溅}

WildGaussians（2024.07）~\cite{2407_08447} 同时解决野外数据的两个核心挑战：外观变化（不同时间、天气拍摄）和遮挡（行人和车辆），分别通过外观建模和不确定性建模两个模块实现。

在外观建模方面，3DGS 将颜色以球谐系数形式显式存储在每个高斯上，这比 NeRF 的隐式 MLP 更难处理外观变化——NeRF 可以简单地在 MLP 输入中拼接外观嵌入，但 3DGS 不能在不损失渲染速度的前提下对每个高斯跑 MLP。WildGaussians 的解决方案是：为每张训练图像分配一个可学习嵌入 $\ell_j \in \mathbb{R}^{32}$（编码该图像的全局外观），同时为每个高斯也分配一个嵌入 $\tau_i \in \mathbb{R}^{24}$（编码局部外观），将两者与高斯的基色 $c_i$（0 阶球谐）一起送入一个共享小 MLP $f_\theta$，输出仿射变换参数 $(\gamma, \beta) = f_\theta(\ell_j, \tau_i, c_i)$，其中 $\gamma_k, \beta_k$（$k=1,2,3$ 为颜色通道）对基色进行色调映射：$\hat{c}_i = \gamma \odot c_i(\mathbf{r}) + \beta$。训练完成后，可以将仿射变换"烘焙"回标准球谐系数，恢复原始 3DGS 的渲染速度。

在不确定性建模方面，WildGaussians 用 DINOv2 的 ViT-S/14 特征替代传统的 MSE/SSIM 残差来检测遮挡。具体地，从训练图像和渲染图像分别提取 DINO 特征 $F$ 和 $D$，不确定性由可训练仿射变换预测：$\sigma = g_\phi(\cos(F, D))$，其中 $g_\phi$ 将余弦相似度映射为逐像素不确定性。训练损失为 $L = L_{\text{3DGS}} + \lambda_{\text{uncert}} \cdot L_{\text{uncert}}$，其中 $\lambda_{\text{uncert}} = 0.5$，$\lambda_{\text{dssim}} = 0.2$。不确定性高的像素（即遮挡区域）在训练中被有效忽略。在 Photo Tourism 数据集上，WildGaussians 在所有 NeRF 和 3DGS 基线中取得最优渲染质量，同时保持 3DGS 的实时渲染速度。

\subsubsection{动静分离与运动分解}

本路线介绍将场景显式分解为静态和动态两类并分别建模的方法。

干扰物去除假设动态内容可以完全丢弃，但许多场景需要同时重建静态和动态元素。本节介绍五种将场景显式分解为动静两类并分别建模的方法。

\paragraph{动静感知解耦 4D 场景重建}

标准 4DGS 将所有高斯统一建模为动态的，即 $\mu_t = \mu_0 + \Delta\mu_t$，即使某个高斯对应的是完全静态的建筑墙面。这种一视同仁的处理方式忽略了动静场景成分之间的本质差异——静态结构（如墙壁、地面）具有稳定的几何与外观，而动态物体（如行人、液体）则表现出显著的时间变化。当两类成分共享同一变形网络时，静态区域的低频信号会吸引优化器的注意力，挤压动态区域的建模资源，导致两类区域都无法获得最优重建。SDD-4DGS（2025.03）~\cite{2503_09332} 是首个基于高斯泼溅实现动静解耦4D场景重建的框架，其核心贡献是为每个高斯引入一个概率动态感知系数 $w$，从理论上实现静动解耦——直觉上，每个高斯需要做出一个"选择"：要么参与动态变形，要么保持静态不动。

\textbf{概率动态感知系数的理论推导。}在标准 4DGS 中，像素颜色公式为 $I(u,v,t) = \sum_i p_i(u,v,t) \alpha_i c_i \prod_{j<i}(1-\alpha_j)$，其中投影概率 $p_i(u,v,t)$ 仅依赖于时间 $t$。SDD-4DGS 在概率密度中引入动态感知系数 $f$，将投影概率改写为 $p_i(u,v,t) = p(u,v; \mu_{i,t}^{2d}, \Sigma_{i,t}^{2d} | t, f_i) \cdot p(f_i)$，其中 $f$ 服从伯努利分布 $f \sim \text{Bernoulli}(w)$。由于动态感知系数 $f$ 是高斯椭球体的内在属性、与时间 $t$ 相互独立，联合概率分布满足 $p(t, f) = p(t) \cdot p(f)$，而 $p(t)$ 如标准 4DGS 一样取常数 1，因此 $p(f_i)$ 成为唯一的额外项。对 $f = 1$（动态）和 $f = 0$（静态）两种互斥状态求期望后，投影概率展开为动态特征与静态特征的加权混合：$p_i(u,v,t) = p_i(u,v \mid t, f_i{=}1) \cdot w + p_i(u,v \mid t, f_i{=}0) \cdot (1{-}w)$。在规范空间中，高斯的静态属性由 $x_0 \sim \mathcal{N}(\mu_0, \Sigma_0)$ 表示，动态属性由变形后的 $x_t \sim \mathcal{N}(\mu_t, \Sigma_t)$ 表示，而最终用于渲染的变形高斯通过上述期望推导为 $\mu_t' = (1{-}w)\mu_0 + w\mu_t = \mu_0 + w\Delta\mu_t$，协方差类似地 $\Sigma_t' = \Sigma_0 + w\Delta\Sigma_t$。当 $w \to 0$ 时，高斯退化为静态（仅使用规范空间参数 $\mu_0, \Sigma_0$）；当 $w \to 1$ 时，完全采用动态变形。这一设计的直觉是：静态物体不需要时间维度上的变形，$w$ 自动将它们固定在规范空间；该参数化无需对现有高斯重建管线做架构改动，即可自适应地完成动静分离。

\textbf{二值约束与渐进调度。}为确保 $w$ 收敛到明确的二值状态（每个高斯要么静态要么动态），SDD-4DGS 引入二元熵正则化损失 $L_{\text{bi}} = -[w\log w + (1{-}w)\log(1{-}w)]$，驱动 $w$ 趋向 0 或 1。然而，过早强制二值化会损害重建质量——消融实验表明，在未使用渐进调度时引入 $L_{\text{bi}}$ 导致 PSNR 下降 0.42 dB，原因是早期训练中场景几何和外观特征尚未收敛，过早的二值约束限制了模型的表达能力。因此 SDD-4DGS 采用渐进约束调度 $\epsilon_{\text{bi}}(t) = 1 - e^{-\gamma t}$（$\gamma = 10^{-4}$），让几何和外观先稳定后再逐步施加分类约束。训练时动态/静态阈值均设为 $\tau_d = \tau_s = 0.5$ 以获得平滑优化，使双峰分布自然涌现；推理时收紧为 $\tau_d = 0.85$ 和 $\tau_s = 0.2$ 以确保精确分类。最终通过阈值过滤将场景分解为动态组件 $G_d = \{G \mid w > \tau_d\}$ 和静态组件 $G_s = \{G \mid w < \tau_s\}$。

\textbf{自动监督信号。}仅凭重建损失优化 $w$ 可能陷入局部最优：静止背景被误分类为动态区域，或缓慢运动物体被误识别为静态结构。为解决这一问题，SDD-4DGS 利用一个广泛观察到的现象——"动态区域在重建过程中通常表现出更高的不确定性"——引入基于不确定性估计的运动掩码 $m$，实现端到端的自动监督：$L_{\text{asg}} = (1{-}m)(L_1(\hat{I}_s, I) + L_{\text{ssim}}(\hat{I}_s, I)) + m(L_1(\hat{I}_d, I) + L_{\text{ssim}}(\hat{I}_d, I))$，其中 $\hat{I}_d$ 和 $\hat{I}_s$ 分别由动态组件 $G_d$ 和静态组件 $G_s$ 独立渲染。该损失对静态区域和动态区域分别施加重建约束，消融实验表明引入 $L_{\text{asg}}$ 后 PSNR 进一步提升 0.19 dB。完整的训练目标为 $L = L_{\text{4dgs}} + \epsilon_{\text{bi}}(t) \cdot L_{\text{bi}} + L_{\text{asg}}$，其中 $L_{\text{4dgs}}$ 为标准 4DGS 重建损失，三个损失项分别负责渲染保真度、动静属性二值化和分类质量校正。

\textbf{实验表现。}SDD-4DGS 在五个基准数据集上一致超越 SOTA：在 Neu3D 多视角真实场景数据集上平均 PSNR 达到 32.39 dB，比之前最优的 31.84 dB 提升 0.55 dB，在日光条件复杂的 Coffee Martini 场景中优势尤为显著；在 D-NeRF 单目合成数据集上 PSNR 达到 34.82 dB（之前最优 34.41 dB），渲染速度为 85.3 FPS；在 HyperNeRF 和 Nerfies 单目真实场景数据集上 7 个场景中的 5 个取得最优结果，对动态纹理和高频细节的恢复能力全面领先；在 KITTI 自动驾驶场景中，SDD-4DGS 能够保持动态物体（如快速移动的车辆）的锐利边缘和纹理细节，而 4DGS 和 SpaceTime 等基线方法在运动车辆周围产生明显模糊和伪影。消融实验进一步验证了各模块的贡献：引入动态感知系数 $w$ 带来 +0.20 dB 的提升（34.14 $\to$ 34.34 dB），渐进二值约束在正确使用时贡献额外 +0.29 dB（34.34 $\to$ 34.63 dB），自动监督信号再提升 +0.19 dB 至最终的 34.82 dB。

\paragraph{动态静态高斯层次运动分解}

DynaSplat（2025.06）~\cite{2506_09836} 针对 3DGS 变形框架中静态与动态高斯共用同一变形网络、导致建模资源浪费且运动精度不足的问题，提出三阶段解决方案。

在动静分离阶段，传统做法仅依赖光流判断运动，DynaSplat 改用双信号融合：先计算每个高斯在所有帧中位置偏移量的方差
$\text{Var}(\delta) = \frac{1}{T}\sum_{t=1}^{T}\|\delta_t - \bar{\delta}\|^2$，
其中 $\bar{\delta}$ 为偏移均值；方差超过阈值 $\tau$ 的高斯进入候选动态集 $G_{\text{variance}}$，再与 2D 光流一致性交叉验证——将高斯投影位置与运动光流图 $F^M$ 比对，确认超过 $\rho = 0.5$ 比例帧存在显著运动后才归入最终动态集 $G_{\text{dynamic}}$。这种 3D 几何信号 + 2D 光流的双校验机制比单一光流判断更可靠。

在层次运动建模阶段，对静态高斯仅用三层轻量 MLP 保留早期训练参数，对动态高斯引入自适应运动网络（Adaptive Motion Networks）。核心思路是将运动分解为粗粒度和细粒度两层：先对邻域高斯取平均得到粗估计 $X_g^{\text{coarse}} = \frac{1}{|N_g|}\sum_{i \in N_g} X_i$，再由 MLP 预测残差 $X_g^{\text{fine}} = \phi(X_g^{\text{coarse}}, f_g)$，最终运动参数为 $X_g = X_g^{\text{coarse}} + X_g^{\text{fine}}$。同时多个运动模式通过可学习混合参数 $\beta_m$ 加权组合：$\delta(t), R(t), S(t) = \sum_{m=1}^{M}\beta_m \cdot \phi_m(\gamma(\mu), \gamma(t), S_c, R_c)$，使网络能灵活捕获不同运动模式。

在不透明度建模阶段，原始 3DGS 使用固定不透明度，DynaSplat 引入基于物理的不透明度计算 $o_g = o_0 \cdot \frac{\cos_g}{\|P_g - P_{\text{camera}}\|_2^2}$，其中 $\cos_g = n_g^\top v_g$ 为高斯法线与视线的夹角余弦，$o_0$ 为初始不透明度常数。该公式使不透明度与视角和距离解耦，更好地处理半透明物体和遮挡。此外，用时间重要性过滤 $w_i = \sum_{I,x,t,T}(\alpha_i(x|t)\prod_{j=1}^{i-1}(1-\alpha_j(x|t)))$ 替代 3DGS 的重置不透明度操作，基于所有训练视角的时间域重要性剪枝浮游高斯。

\paragraph{运动感知分区变形 3DGS}

MAPo（2025.08）~\cite{2508_19786} 的出发点是现有变形方法用单一变形网络建模全部时空变化，在快速或复杂运动区域产生模糊的平均化结果，同时静态区域还在反复参与冗余的变形计算。MAPo 的核心思路是按运动强度分区：让高动态高斯获得专属变形网络，让低动态高斯直接冻结为静态。

首先，MAPo 为每个 3DG 设计综合动态分数，同时度量最大位移和位置方差。对于第 $i$ 个高斯，记录训练过程中的 $m$ 个历史位置 $\mu_{ij}$，计算最大位移 $r_i = \|\max_j \mu_{ij} - \min_j \mu_{ij}\|$（历史位置的轴对齐包围盒对角线长度）和位置方差 $v_i = \frac{\sum_{j=1}^{m}\|\mu_{ij} - \bar{\mu}_i\|^2}{m}$。两个指标分别通过百分位归一化映射到 $[0,1]$ 后，用调和平均融合为最终动态分数 $S_i = \frac{2}{\frac{1}{\hat{r}_i + \epsilon} + \frac{1}{\hat{v}_i + \epsilon}}$（$\epsilon = 10^{-6}$）。调和平均的优势在于要求两个输入同时较高时才输出高分值，能更准确地识别真正的高动态区域。

对高动态 3DG（$S_i$ 超过阈值 $\tau_l$），MAPo 沿时间维度递归二分分区：在分区点 $t_{\text{mid}} = (t_{\text{start}} + t_{\text{end}})/2$ 处将原 3DG 复制为两份，各自保留对应子时间段的变形网络 $F_{[t_{\text{start}}, t_{\text{mid}}]}$ 和 $F_{[t_{\text{mid}}, t_{\text{end}}]}$。这种专门化使每个子网络只需拟合短时间段的运动模式，有效缓解时间平均化效应。对低动态 3DG（$S_i$ 低于阈值 $\tau_{\text{static}}$），在随机时刻用变形网络输出初始化属性后即冻结变形计算，仅保留外观参数可优化，显著降低计算成本。

由于时间分区可能在边界处产生视觉跳变，MAPo 引入跨帧一致性损失 $L_{\text{cross}}$，包含两个分量：自一致性项 $L_{\text{current}} = \|I_t(G_t, V) - I_t(G_t', V)\|_1$ 约束相邻分段对同一帧的渲染一致，真值锚定项 $L_{\text{gt}} = \|I_t(G_t', V) - I^{GT}\|_1$ 用真实图像监督相邻分段的 3DG 学习当前帧的锐利细节，防止自一致性优化导致过度平滑。

\paragraph{多次遍历场景分解}

3DGM（2024.05）~\cite{2405_17187} 面向自动驾驶场景中的多次遍历（multitraverse）建图问题：同一区域被车辆反复经过，每次遇到不同的行人和车辆。传统 3DGS 在此场景下的困境是——短暂物体（行人、车辆）破坏多视角一致性，而遮挡、非朗伯表面等噪声进一步干扰重建。3DGM 的核心洞察是：环境结构在遍历间保持一致，短暂物体则频繁变化，因此可用重复遍历的共识信号作为自监督，自动区分环境与物体。

技术流程分三个阶段。第一阶段（初始化）将多次遍历的图像联合输入 SfM 获取初始高斯点和相机位姿，多遍历的丰富视角天然提升三角化质量。第二阶段（EmerSeg 短暂物体分割）通过鲁棒特征蒸馏实现：在每个高斯上附加 $d$ 维 DINOv2 语义特征 $f_i$，联合优化 RGB 和特征渲染损失 $L = \sum_t(L_{\text{rgb}}(I_t(\xi_t; G), I_t) + L_{\text{feat}}(F_t(\xi_t; G), F_t))$。由于环境像素（内点）在数量上远超短暂物体像素（外点），模型主要被内点梯度驱动去学习环境结构，导致短暂物体区域在特征空间产生高残差。接着利用特征残差挖掘：归一化残差图后，低于阈值 $\tau_1 = 0.3$ 的像素被抑制，再经轮廓检测、小轮廓过滤、天空轮廓去除和凸包提取，最终生成短暂物体掩码 $M_t$——整个过程无需额外训练。第三阶段（EnvGS 鲁棒优化）用掩码屏蔽外点：$L = \sum_t L_{\text{rgb}}(M_t \odot I_t(\xi_t; G), M_t \odot I_t)$，仅从环境像素优化高斯属性，输出干净的环境地图。

\paragraph{视觉几何分解重建动态场景}

SplitGaussian（2025.08）~\cite{2508_04224} 针对现有变形方法（如 Deformable 3DGS）对静态和动态区域统一施加变形网络、导致运动伪影泄漏到静态背景的问题——直觉上，给整个场景套上一件"变形外衣"是不合理的，建筑和行人应该穿不同的衣服，提出在 3D 高斯表示内部显式分解视觉几何。

Baseline（Deformable 3DGS）用单一变形网络 $F_\theta$ 对所有高斯的规范位置施加时间偏移。SplitGaussian 将场景在时刻 $t$ 分解为两组高斯基元：$\mathcal{G}(t) := \{g_s(\mu_s, \Sigma_s, w_s(t))\} \cup \{g_d(\mu_d(t), \Sigma_d(t), w_d(t))\}$。静态组件 $g_s$ 保持几何固定（位置 $\mu_s$ 和协方差 $\Sigma_s$ 不随时间变化），仅通过残差外观网络预测时间依赖的外观变化 $w_s(t) = w_s^{(0)} + \Delta w_s(t)$，其中 $\Delta w_s(t) = \text{MLP}_{\text{app}}^{(s)}([\gamma(\mu_s), \gamma(t)])$，$w_s^{(0)}$ 为冻结的初始球谐系数和不透明度。这使得静态背景能适应光照变化而不会产生几何抖动。

动态组件 $g_d$ 通过变形网络建模几何和外观的时间变化：给定规范位置 $\mu_d(0)$ 和时间 $t$，变形网络输出位置、缩放和旋转偏移 $(\delta, \delta_r, \delta_s) = F_\theta([\gamma(\mu_d(0)), \gamma(t)])$，更新后的几何为 $\mu_d(t) = \mu_d(0) + \delta$、$\Sigma_d(t) = A_d(t)\Sigma_d(0)A_d^\top(t)$。两组使用统一的正弦位置编码 $\gamma(p) = (\sin(2^k p), \cos(2^k p))_{k=0}^{L-1}$，其中空间坐标 $L=10$，时间维度合成场景 $L=6$、真实场景 $L=10$。监督时使用二值掩码 $M$ 分别约束：静态损失 $L_{\text{static}} = L_1(I_s \odot M, I_{\text{gt}} \odot M) + \lambda_{\text{ssim}}\text{SSIM}(I_s \odot M, I_{\text{gt}} \odot M)$，动态损失则在 $(1-M)$ 区域计算。这种掩码引导的分离优化确保每个模块仅从对应区域接收梯度，有效解耦动静表示。

\paragraph{观测完整性量化与标签噪声分离}

OCSplats（2025.08）~\cite{2508_01239} 从观测充分性的认知角度重新审视 3DGS 抗噪声重建。传统方法将抗噪声视为简单的度量分类问题：用重建残差或不确定性值量化噪声，设计固定阈值分离噪声标签。然而，不同场景的噪声类型、比例和观测充分程度差异巨大——远离主视角的区域即使无遮挡，重建误差也显著偏高，导致固定阈值方法频繁误分类。

OCSplats 的第一项核心贡献是观测完整性（Observation Completeness, OC）量化指标。基于三角测量原理，第 $i$ 个高斯在第 $n$ 次迭代时的 OC 定义为 $O_i^n = \alpha_{\text{oc}} O_i^{n-1} + (1 - \alpha_{\text{oc}}) s_i^n \cdot v_i^n$，其中 $\alpha_{\text{oc}} = 0.98$ 为动量系数，$v_i^n$ 衡量观测该高斯的相机位置的方差（反映视角多样性），$s_i^n$ 在位置梯度超过 $10^{-7}$ 时为 1 否则为 0（确认有效观测）。相机位置方差通过迭代公式高效计算，无需存储全部历史位置。

第二项贡献是混合噪声评估与 OC 校正。单独使用残差 $R$ 或不确定性 $\sigma$ 都有缺陷（$R$ 难以检测弱噪声如阴影，$\sigma$ 无法区分难重建区域如远处树林），因此混合评估 $A_h = \lambda_2 R + (1 - \lambda_2)\sigma$（$\lambda_2 = 0.5$）取两者交集增强共有噪声区域。然后用 OC 校正：$A_h^{\text{oc}} = A_h \cdot \text{OCR}$，其中校正比 $\text{OCR} = \frac{1}{1 + \sigma_3(\text{th} - O)} \cdot S$，$\sigma_3 = 3.0$、$\text{th} = 0.3$，$S$ 为 Sobel 算子计算的纹理强度。该公式在观测不足的区域自动放大容忍度，避免因观测不充分导致的误分类。

第三项贡献是动态锚点组合阈值。传统固定阈值无法适应噪声比例差异巨大的场景，OCSplats 改用两个统计锚点计算自适应阈值：Otsu 最大类间方差位置和背景类均值位置。阈值完全由每个场景自身的统计特性决定，无需人工调参即可直接应用于不同场景。

\subsubsection{实例分解与语义4D理解}

本路线介绍将实例分割、语义嵌入与动态3DGS重建联合优化以实现精细场景理解的方法。

动静二分只是最粗粒度的分解；更精细的场景理解需要识别每个独立物体实例并赋予语义标签。本节介绍八种将实例分割、语义嵌入与动态 3DGS 重建联合优化的方法。

\paragraph{实例感知高斯动态场景建模与理解}

Director（2026.04）~\cite{2604_01678} 的核心主张是动态重建与场景理解不应遵循"先重建再理解"的分离流水线，而应在实例一致的高斯基元层面联合优化。现有方法要么只关注光度一致性（如 4DGS、GaussianFlow），要么仅在后期附加语义（如 LERF、4DLangSplat），导致高度动态场景中出现身份漂移、前景-背景纠缠和语义关联不稳定。Director 的解决方案是将实例级身份和语言对齐的语义嵌入直接作为每个高斯的可学习属性，与运动和外观联合优化。

在初始化阶段，Director 利用 SAM3 的记忆库机制，以文本提示（如"basketball player""basketball"）引导分割，在单视频内获取时间一致的实例掩码。为建立跨视角身份一致性，将所有视角的首帧拼接为规范视频统一分割——由于相邻摄像机视点间距小、画面重叠度高，规范掩码天然对齐，再通过 IoU 匹配建立跨视角的一一对应关系，传播至全部帧，获得时空一致的实例掩码 $M^{v,t}_d$。同时用多模态大语言模型为每个实例生成双层语义描述：时间不变的全局描述 $C_i^{\text{global}}$（在首帧用 SAM3 掩码渲染红色轮廓高亮实例，提示 Qwen3-VL 描述整体外观和动作特征）和逐帧动态描述 $C_i^t$（将所有视角在第 $t$ 帧的观测聚合为空间一致性视频剪辑，连同全局描述和时间提示输入 MLLM）。两种描述经 LLM 编码后，再通过轻量级自编码器压缩为紧凑的 6 维句子嵌入 $e$，作为语言对齐的语义监督信号。

在场景分解上，Director 将场景显式分为静态背景 $G_{\text{bg}}$ 和动态前景 $G_{\text{fg}}$。背景层通过跨时间和空间维均匀采样 GT 图像全局优化，损失函数为 $E^{\text{bg}}_{\text{color}} = \|M_{\text{bg}} \odot (\hat{C} - C)\|_1$（仅在背景掩码区域计算），配合等渗正则 $E_{\text{iso}}$ 和尺寸正则 $E_{\text{size}}$ 防止高斯过度拉伸（$\lambda_{\text{iso}} = 0.0005$、$\lambda_{\text{size}} = 0.02$），训练 20000 次迭代。背景层不仅提供全局参考，在后续每帧中还持续优化外观属性以适应光照和阴影变化。

前景层中每个高斯附带 8 维可学习语义特征 $f$（均匀分布随机初始化），通过两个 MLP 解码头分别预测实例身份和语言对齐嵌入：分类头将渲染的语义特征图 $F \in \mathbb{R}^{H \times W \times 8}$ 映射到 $(D+1)$ 个通道经 softmax 输出逐像素类概率 $P$，用交叉熵 $E_{\text{id}} = -\sum_{d=0}^{D} y_d \log P_d$ 监督；嵌入头用 L1 损失 $E_{\text{emb}} = \|\hat{e} - e\|_1$ 监督语言对齐嵌入。双解码头设计确保模型同时捕获低层实例身份和高层语义信息。3D 空间一致性通过 KL 散度正则化实现：$E_{\text{3D}} = \frac{1}{|S|}\sum_{i \in S}\frac{1}{k}\sum_{j \in N_k(i)}D_{\text{KL}}(P_i \| P_j)$（采样 6000 个前景高斯，$k=4$ 近邻），约束同一实例的邻近高斯保持一致的语义特征，使遮挡区域和物体内部也能获得连贯的特征表示。首帧初始化阶段联合优化高斯和两个 MLP 共 30000 次迭代（$\lambda_{\text{id}} = 2$、$\lambda_{\text{emb}} = 10$、$\lambda_{\text{3D}} = 2$）。

在动态优化阶段，Director 采用逐帧策略将光流与 4D 高斯运动桥接，分为显式变形和训练两个子阶段。显式变形阶段用 SEA-RAFT 预测帧间光流，将前景高斯 $G_{\text{fg}}^{t-1}$ 的位置 $p_{t-1}$ 投影到各相机平面查询 2D 光流得到变形后像素坐标 $u_v'$，再通过多视角最小二乘三角化恢复 3D 位置 $p_t = \arg\min_{p \in \mathbb{R}^3}\sum_{v=0}^{V-1}\|\pi_v(p) - u_v'\|_2^2$。由于光流噪声和遮挡导致三角化不准确，引入局部近似刚体（ARAP）正则化 $E_{\text{smooth}} = \sum_i\sum_{k \in N(i)}w_{i,k}^{t-1}\|R(q_{i,t} * q_{i,t-1}^{-1})(p_{k,t-1} - p_{i,t-1}) - (p_{k,t} - p_{i,t})\|_2^2$，其中混合权重 $w_{i,k}^t = \exp(-\|p_{i,t} - p_{k,t}\|_2^2 / l^2)$，$l$ 为影响半径。细化阶段冻结外观和语义属性，仅优化位置和旋转以修正物理不合理的变形（$\lambda_{\text{smooth}} = 0.01$，3000 次迭代）。

训练阶段联合优化运动、外观和语义特征。前景高斯动态分裂和剪枝以适应新观测，克隆高斯继承父属性以保证轨迹连续性。SDF 损失 $E_{\text{sdf}} = \sum_i y_d \cdot d_i^2$ 仅对预测标签匹配实例 $d$ 且落在 2D 轮廓外的高斯施加惩罚（$d_i$ 为到掩码边界的欧氏符号距离），将高斯约束在其所属实例的轮廓内。时间正则 $E_{\text{temp}}^{\text{bg}} = \sum_i\|a_{i,t} - a_i\|_2^2$ 约束背景外观不随时间抖动，$E_{\text{temp}}$ 在连续帧间约束所有属性的突变。训练中冻结分类 MLP，仅继续优化语义 MLP，每 300 次迭代执行克隆/剪枝且每帧修改的高斯总数不超过 5\%。整体训练损失为 $E_{\text{train}} = \lambda_{\text{sdf}}E_{\text{sdf}} + \lambda_{\text{temp}}^{\text{bg}}E_{\text{temp}}^{\text{bg}} + \lambda_{\text{temp}}E_{\text{temp}} + \lambda_{\text{smooth}}E_{\text{smooth}} + \lambda_{\text{id}}E_{\text{id}} + \lambda_{\text{emb}}E_{\text{emb}} + \lambda_{\text{3D}}E_{\text{3D}} + E_{\text{color}}$，其中 $\lambda_{\text{sdf}} = 0.01$、$\lambda_{\text{temp}}^{\text{bg}} = 0.001$、$\lambda_{\text{temp}} = 0.01$、$\lambda_{\text{smooth}} = 0.0001$、$\lambda_{\text{id}} = 1$、$\lambda_{\text{emb}} = 10$、$\lambda_{\text{3D}} = 2$，每帧训练 8000 次迭代。

查询阶段采用两阶段策略：先用身份查询（identity query）通过分类 MLP 检索目标实例 ID，仅渲染该实例的 RGB 和特征图，再用段落查询（segment query）通过余弦相似度在渲染特征图的高维空间中定位最相关的帧片段——这一设计有效过滤了其他动态实例的干扰，实现精确的实例级检索与分析。

\paragraph{实例分解 4D 高斯多视角视频}

在多视角动态场景的实例分割中，一个核心难题是跨视角标签不一致：由于视频分割模型（如 SAM3）独立处理每个视角的视频流，不同摄像机产生的实例标签空间完全不对应——视角 $v_1$ 中的"人 1"在视角 $v_2$ 中可能被编号为"人 3"。现有方法要么回避直接监督（如 TRASE 的对比学习），要么只能使用单视角监督（如 SA4D），但对比学习产生的隐式聚类边界无法保证跨时空的身份一致性。Inst4DGS（2026.03）~\cite{2603_18402} 的核心思路是：既然标签不一致是视角间的映射问题，那就直接学习这个映射。具体地，为每个视角 $v$ 引入一个可学习的置换潜变量 $Z_v \in \mathbb{R}^{K \times K}$（$K$ 为最大实例数），通过可微 Sinkhorn 层（一种将任意方阵投影为双随机矩阵的迭代算法）将其约束为近似置换矩阵：
$$S_v^{(0)} = \exp(Z_v), \quad S_v^{(i)} = T_c(T_r(S_v^{(i-1)})),$$
其中 $T_r$、$T_c$ 分别为行归一化和列归一化操作，迭代后 $S_v$ 趋近一个一一对应矩阵。该矩阵将共享的规范标签 $L^{\mathrm{canon}}$ 映射到各视角的本地标签：$L_v^{\mathrm{local}} = S_v \cdot L^{\mathrm{canon}}$。由于整个过程可微，规范身份特征 $f_i$ 和置换潜变量 $Z_v$ 可联合优化。为避免初始阶段规范标签空间未收敛时的训练歧义，Inst4DGS 采用渐进学习策略：初始仅使用单一参考视角更新身份特征，其他视角只优化置换潜变量；当某视角的置换置信度超过阈值 $S_{\mathrm{th}}$ 时才激活该视角对身份特征的梯度。每个视角仅引入 $K \times K$ 个额外参数，开销极小。

在运动建模方面，传统的 Dynamic3DGS 为每个高斯独立优化轨迹，计算量大且缺乏语义结构。Inst4DGS 利用实例分组信息为每个物体实例初始化专属的运动支架（motion scaffold），每个高斯 $G_i$ 仅绑定同实例标签下最近的 $k$ 个运动基元 $B_j$，其轨迹通过双四元数混合（Dual Quaternion Blending，一种在旋转和平移联合空间中进行插值的方法）计算：$T_i = \mathcal{B}(\{B_j, w_{ij}\})$，其中 $\mathcal{B}$ 为 $SE(3)$ 姿态的双四元数混合，$w_{ij}$ 为混合权重。这意味着只需优化每个实例的少量运动基元轨迹，而非所有高斯的独立轨迹。同时配合实例感知刚性损失 $\mathcal{L}_{\mathrm{rigid}} = \mathcal{L}_{\mathrm{coord}} + \mathcal{L}_{\mathrm{len}} + \mathcal{L}_{\mathrm{iso}}$，约束同一实例内邻近高斯的运动一致性。最终在 Panoptic Studio 数据集上，PSNR 从 Dynamic3DGS 的 26.10 提升至 28.36，实例 mIoU 从 SA4D 的 0.6310 提升至 0.9129。

\paragraph{分割任意 4D 高斯}

将 SAM 的分割能力从静态 3D 扩展到动态 4D 场景时，一个看似直觉的方案是：先在某个时间戳对变形后的 3D 高斯做 3D 分割，再通过变形场传播到其他时间戳。但这一方案会遇到"高斯漂移"（Gaussian Drifting）问题——由于非刚体运动，某些高斯在时间 $t_0$ 之前属于物体 A，之后却漂移到了物体 B 的区域，导致分割身份在时间维度上发生混乱。根本原因在于：规范空间中的 3D 高斯及其变形场 $F$ 都不携带物体身份信息，无法区分不同物体的运动轨迹。

Segment Any 4D Gaussians（2024.07）~\cite{2407_04504} 的核心设计是引入时间身份特征场（Temporal Identity Feature Field），让身份信息随时间变化。具体地，给定时间 $t$ 和规范空间中高斯的中心位置 $X$，一个小型 MLP 网络 $\phi_\theta$ 预测每个高斯的时间相关身份编码：
$$e = \phi_\theta(\gamma(X), \gamma(t)),$$
其中 $\gamma(\cdot)$ 为位置编码，输出为 32 维身份特征向量。该特征通过标准高斯泼溅渲染为 2D 身份特征图 $E_p$，再经微型卷积解码器 $\phi_c$ 和 softmax 得到逐像素的实例概率 $f = \mathrm{softmax}(\phi_c(E_p))$。身份编码维度设为 32，卷积分类器为 32 输入通道、256 输出通道，训练 5000 次迭代即可收敛。之所以让身份特征随时间变化而非使用静态编码，正是为了解决高斯漂移：即使某个高斯在不同时间戳所属的物体发生变化，时间相关的特征场也能自适应地调整其身份表示。训练阶段采用视频目标跟踪模型 DEVA 生成的逐帧身份掩码作为监督信号，损失函数为 $\mathcal{L} = \lambda_{2d}\mathcal{L}_{2d} + \lambda_{3d}\mathcal{L}_{3d}$，其中 $\mathcal{L}_{2d}$ 为标准交叉熵身份损失，$\mathcal{L}_{3d}$ 为 3D 正则化 KL 散度损失（强制空间邻近高斯的身份编码一致），$\lambda_{2d}=1$，$\lambda_{3d}=2$。

推理阶段，SA4D 引入高斯身份表（Gaussian Identity Table）将各训练时间戳的分割结果存储为查找表，推理时只需最近时间戳插值：$M_i = \mathrm{nearest}(M(g, t_i))$，避免了逐帧前向传播网络的开销。在此基础上执行两步精化：首先移除离群高斯，然后利用投影损失 $\mathcal{L}_{\mathrm{proj}}$ 消除物体边界处的模糊高斯，最终在 RTX 3090 上 10 秒内完成交互式 4D 分割，推理渲染帧率从无身份表时的 7.12 FPS 提升至 40.36 FPS，支持物体移除、重着色和组合等动态场景编辑操作。

\paragraph{学习 3D 一致运动与语义}

2D 基础模型（SAM2、TAP、Depth Anything）虽在单帧预测上表现优异，但缺乏 3D 推理能力，即便最先进的 SAM2 也难以避免跨帧的空间错位、时序闪烁和边界伪影。现有 4D 语义方法将 2D 先验与动态 3D 表示分开处理——要么独立于 3D 模型学习特征场，要么将语义理解与运动估计解耦——导致外观和运动的有限建模无法产出一致的预测。Motion4D（2025.12）~\cite{2512_03601} 的核心思路是将多种 2D 先验（语义掩码、2D 点轨迹、单目深度）集成到统一的 4D 高斯表示中，通过迭代精化让 2D 先验和 3D 场景表示互相校准。

具体地，每个 3D 高斯除常规属性外增加运动场和语义场。运动场采用全局运动基的加权组合：对每个高斯 $i$ 定义变换 $T_i^{0 \to t} = \sum_{b=0}^{B} w_i^b \Delta_b^{0 \to t}$，其中 $\{\Delta_b^{0 \to t}\}_{b=1}^B$ 为 $B$ 个全局运动基，$w_i^b$ 为各高斯的基系数，位置和旋转更新为 $\mu_i^t = R_i^{0 \to t} \mu_i^0 + t_i^{0 \to t}$，$R_i^t = R_i^{0 \to t} R_i^0$。语义场则为每个高斯附加特征向量 $f_i^{sem}$，通过标准 $\alpha$-混合 渲染语义图：$I^{sem}(p) = \sum_{i \in H(p)} f_i^{sem} \alpha_i \prod_{j=1}^{i-1}(1-\alpha_j)$。

优化框架由两部分交替执行。顺序优化在短时窗口内分两阶段更新运动场和语义场以维持局部一致性：第一阶段为迭代运动精化，引入 3D 置信度图控制 2D 跟踪先验的监督权重——每个高斯附加不确定性标量 $u_i$，渲染为逐像素置信度 $w(p)$，跟踪损失 $\mathcal{L}_{track} = \frac{1}{|I_t|} \sum_{p \in I_t} w(p) \| \hat{U}_t^{t'}(p) - U_t^{t'}(p) \|$ 自动降低不可靠先验的影响；同时，自适应重采样模块根据每像素的 RGB 误差和语义误差识别欠表示区域并插入新高斯。第二阶段为语义精化，利用 SAM2 的可提示特性，将 3D 渲染的语义掩码作为额外提示（边界框和提示点）反馈给 SAM2，迭代更新 2D 语义先验。全局优化则在整个序列上联合优化所有属性，确保长期一致性。在 DyCheck-VOS 基准上，Motion4D 在视频目标分割、点跟踪和新视角合成三项任务上均显著优于 2D 基础模型和 3D 表示方法。

\paragraph{双层次优化语义 4DGS}

将静态场景的语义理解方法（如 LangSplat）直接扩展到动态场景有两种朴素策略：一是先用 4DGS 重建动态场景再冻结几何、单独训练语义特征，但这样会导致特征提取不完整；二是联合优化几何与语义，但高斯倾向于模糊物体纹理以追求语义完整性。两种策略都无法同时实现高保真渲染和准确语义。DHO（2025.03）~\cite{2503_19332} 的核心洞察是：静态背景和动态前景需要不同的优化自由度，分治策略才能兼顾渲染质量与语义精度。

DHO 首先扩展 4DGS 表示，为每个高斯引入 $n$ 维语义特征 $f \in \mathbb{R}^n$（实验中 $n=8$），通过自编码器将 512 维 CLIP 特征压缩为紧凑表示。渲染时语义特征与颜色同理进行 $\alpha$-混合：$F(v) = \sum_i^N f_i \alpha_i \prod_{j=1}^{i-1}(1-\alpha_j)$。语义损失跨三个尺度（子部件、部件、整体）监督：$\mathcal{L}_s = \sum_{l \in \{s,m,l\}} \sum_{t=1}^T \mathcal{L}(L_t^l(p), F_t^l(p))$。

方法由两个核心模块组成。层次高斯流（Hierarchical Gaussian Flow, HGF）对不同世代的高斯施加差异化约束。训练期间每次致密化产生的高斯标记为第 $i$ 代，越早产生的高斯受到越强的锚点约束：$\mathcal{L}_{anchor}^s = \sum_{i=0}^n \lambda_i[(S_{i,c} * \beta_c + S_{i,f} * \beta_f) - \theta_i]^2$，其中 $S_{i,c}$ 为粗阶段的几何属性（$\beta_c=1, \beta_f=0$），$S_{i,f}$ 为精阶段的语义特征（$\beta_c=0, \beta_f=1$），$\theta_i$ 为锚点状态。粗阶段（前 5k 迭代）锁定静态背景几何，精阶段（后 15k 迭代）释放动态前景自由度但约束其语义特征，实现静态与动态的有效分离。

层次高斯引导（Hierarchical Gaussian Guidance, HGG）利用语义先验引导渲染关注动态前景。对噪声背景中的场景，通过 CLIP 获取前景掩码 $M$，将图像分解为掩码内外两部分并施加自适应权重：$\mathcal{L}_{mask}^v = \beta_m \mathcal{L}(I_m - \hat{I}_m) + (1-\beta_m)\mathcal{L}(I_{\bar{m}} - \hat{I}_{\bar{m}})$，其中 $\beta_m = \sigma(\alpha(\bar{c} - \tau))$ 基于场景平均纹理复杂度 $\bar{c}$ 自适应调节（$\alpha=10, \tau=0.01$），纹理越复杂的场景对前景赋予越大权重，将高斯资源从背景噪声中解放出来。在 HyperNeRF 和 D-NeRF 数据集上，DHO 在动态前景区域渲染质量和语义分割精度上均显著优于顺序训练、联合训练和并行工作 DGD，且仅需单张 RTX 3090（8GB 显存）即可训练。

\paragraph{语义感知 4DGS 刚体约束}

物理驱动 4DGS 方法（如 TRACE）为每个高斯独立预测平移-旋转动力学参数，再通过二阶 Runge-Kutta 积分外推未来帧。这种逐高斯方案存在两个问题：推理时需对所有 $N$ 个高斯逐一调用 MLP 和积分器，计算开销大（$O(N)$）；且独立演化的高斯缺乏语义一致性，无法被聚合为有意义的物体实例以支持下游感知任务。R5DGS（2026.05）~\cite{2605_25909} 的设计出发点正是解决这两个问题——用语义分组实现物体级理解，用刚体约束降低推理成本。

R5DGS 首先为每个高斯附加 16 维可学习身份向量 $e_i \in \mathbb{R}^{16}$，通过 $\alpha$-混合 渲染为身份图 $E_{id} = \sum_i^N e_i \alpha'_i \prod_{j=1}^{i-1}(1-\alpha'_j)$。身份图由分类器 $f_{cls}$ 映射到 $K$ 个物体类别的交叉熵损失监督，同时施加 3D 空间正则化（$k=5$ 近邻的 KL 散度惩罚）强制空间邻近高斯的身份编码一致。总训练目标为：$\mathcal{L}_{5DGS} = \mathcal{L}_{render} + \lambda_{obj}\mathcal{L}_{obj} + \lambda_{3d}\mathcal{L}_{3d}$，其中 $\mathcal{L}_{render} = (1-\lambda_{dssim})\mathcal{L}_1 + \lambda_{dssim}(1-\mathrm{SSIM})$。

核心创新是刚体推理约束。直觉上，刚体上所有点的运动可以由一个代表点的运动完全确定——因此只需对每个物体实例的代表高斯做推理，其余高斯通过刚体变换继承结果。设物体 $k$ 包含高斯集合 $G_k$，选取距几何质心最近的代表高斯 $r_k$，预计算规范空间偏移 $o_i = x_i - x_{r_k}$。外推时仅对 $K$ 个代表高斯运行 TRD 动力学预测，获得未来位置 $x_{r_k}^{vel}$ 和四元数 $q_{r_k}^{vel}$，然后刚体传播到组内所有高斯：
$$\Delta x_i = (x_{r_k}^{vel} - x_{r_k}^{def}) + (R(q_{r_k}^{vel})o_i - o_i), \quad q_i^{out} = q_{r_k}^{vel} \otimes q_i^{def},$$
其中 $R(\cdot)$ 将四元数转为旋转矩阵，$\otimes$ 为 Hamilton 乘积。该公式在构造上保持了组内高斯间的相对距离和朝向，同时将 MLP 和积分器调用从 $O(N)$（$N \approx 40000$）降至 $O(K)$（$K \approx 9$--$12$），外推速度提升 11 FPS。

此外，R5DGS 构建离线 CLIP 物体查找表：对每个物体组提取代表性掩码视角并编码为文本对齐嵌入 $\tau_g$，推理时通过余弦相似度 $g^* = \arg\max_g \cos(\phi_{text}(q_{text}), \tau_g)$ 实现开放词汇文本提示检索。在 Dynamic Indoor Scene 数据集上，R5DGS 在保持物理合理运动轨迹的同时，分割精度几乎不受刚体约束影响。

\paragraph{部件感知运动引导动态场景重建}

大多数动态 3DGS 方法将高斯视为独立基元进行逐点变形，忽略了物理运动的基本规律——真实场景由部件（part）构成，部件内部保持刚性一致，部件之间才发生相对运动。同时，2D 光流作为天然的逐像素运动线索，能提供比光度损失更直接的运动先验，但现有方法要么完全忽略光流，要么仅做简单的像素级监督，在光流估计不准确时反而干扰优化方向。PaMoSplat（2026.05）~\cite{2605_10307} 基于这两个观察，将高斯部件作为变形的基本单元，用多视角光流引导部件运动。

在初始时间戳 $t=0$，PaMoSplat 并行执行两条管线。第一条遵循标准 3DGS 生成初始高斯场 $S_0$（10000 次迭代加致密化）。第二条用 SAM 生成的多视角 2D 分割掩码 $\{m_v^i\}$，通过深度引导的跨维度对应提升到 3D：对每个高斯 $g^j$ 和视角 $v$，若渲染深度与高斯深度的差小于阈值 $\epsilon_{depth}$（实际场景设为 0.1m），则标记该高斯在此视角可见。然后构建全连接加权图 $G=(V, \omega)$，边权定义为两个高斯在同一掩码内共现次数与共可见视角数之比：
$$\omega(g^j, g^k) = \frac{\sum_v \mathbb{1}_{g^j \in V_v[m_v^i]} \cdot \mathbb{1}_{g^k \in V_v[m_v^i]}}{\sum_v \mathbb{1}_{g^j \in V_v} \cdot \mathbb{1}_{g^k \in V_v}},$$
再用 Louvain 算法对图做聚类，每个簇即一个高斯部件 $P^\pi$，赋予部件 ID $\pi$。这种基于图聚类的方法对分割不精确具有鲁棒性。

对后续时间戳 $t>0$，PaMoSplat 利用 RAFT 计算多视角光流，通过差分进化算法估计每个部件的刚体运动（6DoF），为高斯场提供鲁棒的暖启动。差分进化在连续空间中搜索最优刚体变换，将多视角光流作为适应度函数的监督信号，避免了简单像素级光流损失在估计不准时干扰优化。在暖启动基础上，PaMoSplat 引入可学习的部件内刚性损失确保部件内部运动一致性，以及光流监督的渲染损失增强动态建模质量。此外，根据暖启动质量自适应调整每个时间戳的优化迭代次数，兼顾效率与精度。在 PanopticSports 数据集（含大幅不规则运动）上，PaMoSplat 在渲染和密集跟踪精度上均优于所有基线，同时其部件感知能力自然支持 4D 场景编辑等下游应用。

\paragraph{多人多物体 3D 高斯渲染}

多人多物体（Multi-Human Multi-Object, MHMO）渲染面临两个耦合难题。一是稀疏视角下每个实例的表示不完整：标准 NeRF/3DGS 管线逐视角独立处理信息，缺乏让信息充足视角指导遮挡视角重建的显式机制，导致同一实例从一个视角看完整、另一视角却模糊或碎片化。二是实例间依赖关系的建模：MHMO 场景远非独立实体的简单叠加，接触区域的几何歧义、近距离的颜色渗透和阴影等细微效果都需要实例间的协同推理，而现有方法只能建模单个人-物对的关系。MM-GS（2026.04）~\cite{2604_02996} 提出从粗到精的层次框架，分别解决这两个问题。

框架首先建立鲁棒的初始表示：对每个行人用 SMPL 模型变形到目标姿态，通过调制线性混合蒙皮（Modulated LBS）获得更精细的初始几何。具体地，姿态化中心为 $\mu_i^0 = \sum_{k=1}^K w_k(R_k \mu_i^c + t_k) + b$，其中蒙皮权重 $w_k = \mathrm{softmax}(w_k^{SMPL} + m_k)$ 通过小型 MLP 预测的调制向量 $m_k$ 修正标准 SMPL 权重。对每个物体则直接用刚体变换将模板网格顶点变换到场景中。所有网格顶点初始化为 3D 高斯，形成初始状态 $G_i^0$。

逐实例多视角融合模块（Per-Instance Multi-View Fusion）解决表示一致性问题。对每个实例 $i$，先用 2D CNN 提取各视角图像的特征图 $F_j$，将高斯中心投影到图像平面采样视觉特征 $f_{i,j}^{vis}$，与初始属性拼接构成视角依赖特征 $f_{i,j}^{vd}$。核心融合操作让不同视角的特征互相通信：$f_i^{fused} = \frac{1}{N_{ctx}} \sum_{j=1}^{N_{ctx}} \frac{1}{Z_{i,j}}(f_{i,j}^{vd} + \gamma \sum_{p \in V_{i,p \neq j}} f_{i,p}^{vd})$，其中 $\gamma=0.1$ 为融合因子。这使得信息充足的视角能有效补充被遮挡视角的缺失信息，建立每个实例的完整一致表示。

场景级实例交互模块（Scene-Level Instance Interaction）在全局场景图上建模所有实例间关系。用 2 层图注意力网络（GAT，一种在图结构数据上应用注意力机制的神经网络，4 头注意力，隐层维度 64，dropout 0.1）连接所有交互实例，聚合的 64 维 GAT 特征映射为 7 维调制向量更新高斯的颜色和不透明度。该模块能精化接触面的高斯属性以获得清晰边界，并渲染近距离实例间的细微视觉效果。在 HOI-M$^3$ 数据集上，MM-GS 在 PSNR、SSIM 和 LPIPS 上均显著超过 NeRF 基线 NeuralHOIFVV-MM 和 3DGS 基线 GTU-MM。

\paragraph{资产驱动多人多物体交互语义重建}

从单目视频重建多人多物体交互场景是一个极具挑战的问题。DreamScene4D 等纯 GS 优化方法依赖 2D 扩散模型生成多视角图像，但这些图像往往不一致，导致重建质量低下；在严重遮挡下逐帧 GS 优化更无法维持物体结构的完整性。Asset-Driven Semantic Reconstruction（2025.12）~\cite{2512_00547} 的核心思路是"生成-变形-优化"三阶段混合策略：先用 3D 生成模型为每个场景元素创建高保真网格，再用语义感知的变形将其映射到每一帧，最后用 GS 渲染损失精化对齐。

第一阶段为每个物体选择遮挡最少的标准帧（canonical frame），用预训练 3D 生成模型 TRELLIS 生成多视角一致的高保真网格 $\{O_c\}_{i=1}^N$。TRELLIS 结合多视角 2D 外观特征与 3D 稀疏网格，克服了纯 2D 扩散模型的多视角不一致问题。但该模型只能处理单张图片输入，无法保证视频序列中不同帧的时序一致性，因此仅用于生成标准帧的网格。

第二阶段将场景元素分为非刚性（人）和刚性（物体）两类分别处理变形。对人类，利用 SMPL 模型的线性混合蒙皮实现时序一致的变形：$x_v^d = w \hat{B}_v^d \cdot (w \hat{B}_v^c)^{-1} x_v^c$，其中 $w$ 为从 SMPL 映射到生成网格顶点的蒙皮权重，$\hat{B}_v^c$ 和 $\hat{B}_v^d$ 分别为标准帧和目标帧的骨骼变换矩阵。这一公式将 LBS 变形从 SMPL 的 T-pose 桥接到自定义的标准帧，使生成模型预测的高质量网格能正确变形到每一帧。对刚性物体，先用单目深度估计模型 VGGT 生成逐帧深度图并反投影为点云，通过 ICP 配准将标准帧网格与场景对齐，再计算相邻帧点云间的刚体变换 $PC_t^c \to PC_{t \pm 1}$，逐帧传播到整个序列。

第三阶段用每帧的网格顶点初始化高斯，通过渲染损失和深度损失联合优化场景元素的 6DoF 位姿，精化个体元素在场景中的对齐。这种混合方法的关键优势在于：即使某个物体在某些帧中被严重遮挡，其结构完整性仍由标准帧的 3D 生成网格保证，不会因逐帧优化中的遮挡信息缺失而崩溃。在 HOI-M$^3$（目前唯一包含多人多物体交互的数据集）上，该方法在距离精度、完整度、F-score 和 Chamfer 距离四项指标上均超越 SOTA。

\begin{takeaway}
动态物体与场景分解将场景拆分为独立运动的实例，核心挑战是在严重遮挡下维持物体结构完整性。MM-GS 用图注意力网络建模实例间交互，Asset-Driven 方法则用 3D 生成模型创建高质量网格再变形到每帧，后者即使物体被严重遮挡也能保证结构完整性。从逐帧优化走向"生成-变形-优化"的混合策略是当前的关键趋势。
\end{takeaway}

%% ============================================================
\section{其他动态场景方法}
\label{sec:06_7}

前五节覆盖了动态场景重建的主流方向。本节将视野扩展到特殊传感模态、非典型动态过程（如植物生长）以及 SLAM、超分辨率等辅助任务，展示 4DGS 的广泛适用性。

\subsection{问题背景：动态4D的多元化拓展}

前述各节聚焦于动态场景的核心表示、重建和应用问题。此外，3DGS/4DGS 的动态能力正在向多个方向拓展：4D 语义理解、特殊传感模态（雷达、热红外、事件相机）、植物生长等非典型动态、SLAM 系统、超分辨率和特征匹配等。本节涵盖 20 篇论文，展示 4DGS 的广泛适用性。具体分为五条技术路线：（1）4D语义与语言场；（2）HDR与光照动态；（3）特殊模态与场景；（4）SLAM与不确定性建模；（5）超分辨率与匹配。

在此基础上，UniMGS 进一步统一了网格与高斯表示。

\paragraph{UniMGS：2026.01 同期工作}

UniMGS（2026.01）~\cite{2601_19233} 统一网格和 3D 高斯泼溅为单次光栅化管线。

\paragraph{DeMapGS：2025.12 同期工作}

DeMapGS（2025.12）~\cite{2512_10572} 同时实现网格变形和表面属性映射。

\subsubsection{4D语义与语言场}

本路线介绍将语言场扩展到4D动态场景以支持时间敏感的开放词汇查询的方法。

动态场景的理解不仅需要识别物体，还需要用自然语言描述它们的运动状态。本节介绍两种将语言场扩展到 4D 动态场景的方法。\paragraph{多模态大语言模型 4D 语言高斯}

4D LangSplat（2025.03）~\cite{2503_10437} 致力于解决动态场景中 4D 语言场的构建问题，支持时间无关（如"狗"）和时间敏感（如"正在奔跑的狗"）两类开放词汇查询。现有方法 LangSplat 将 CLIP 特征与 SAM 分割掩码结合提升到 3D 高斯空间，但直接将其扩展到 4D 场景存在根本缺陷：CLIP 设计用于静态图文匹配，无法捕获视频中物体的状态变化、动作和时序语义。一种直接思路是用视觉模型裁剪感兴趣物体并提取 patch 特征，但裁剪会引入背景噪声，而用精确掩码裁剪又使视觉模型无法区分物体运动与相机运动。因此，4D LangSplat 绕过视觉特征，转而从 MLLM 生成的物体级视频描述文本中直接学习语言特征。

具体而言，4D LangSplat 为每个高斯点同时构建两种语义场。时间无关语义场沿用 LangSplat 的 CLIP 特征方案，在 SAM 定义的三个语义粒度层级上学习静态语义嵌入，编码"人""杯子""狗"等不随时间改变的属性。值得注意的是，尽管每个高斯点的位置和形状在 4D-GS 变形场中随时间动态变化，其时间无关语义特征始终保持静态——位置在动，但"是什么"不变。时间变化语义场则是本文核心创新，采用多模态物体级视频提示方法：对于视频 $V = (I_1, I_2, \ldots, I_T)$，先用 SAM 加 DEVA 跟踪获取 $n$ 个物体的时间一致掩码 $(M_1, M_2, \ldots, M_n)$；然后为每个物体 $i$ 在每帧 $t$ 构建视觉提示 $P_{i,t} = \mathcal{R}(M_{i,t}) \circ \mathcal{G}(\bar{M}_{i,t}) \circ \mathcal{B}(\bar{M}_{i,t})$，分别对目标物体加红色轮廓、背景灰度化、背景高斯模糊，在聚焦目标的同时保留场景上下文——这一设计确保 MLLM 能够关注目标物体的动作和状态，同时通过保留背景参照来维持空间关系和环境上下文的理解；接着先让 MLLM 生成视频级运动描述 $D_i$，再以 $D_i$ 为上下文结合帧级视觉提示生成逐帧物体描述 $C_{i,t}$，实现从粗到细的时序语义标注。最终用微调 LLM（2023）~\cite{wang2023improving} 提取句子嵌入 $e_{i,t}$，对于掩码 $M_{i,t}$ 内的每个像素 $(x, y)$，赋予 $f_{x,y,t} = e_{i,t}$，形成像素对齐的 2D 监督信号。选择 LLM 提取嵌入而非 CLIP 的关键优势在于：LLM 在共享语言空间中处理自由文本的能力更强，使嵌入能够自然地对齐用户查询的自然语言表达，从而高效支持开放词汇查询。

在学习时间变化 4D 场时，若直接用类似 4D-GS 的变形场让每个高斯点的语义特征自由变化，会导致任意语义状态跳跃，增加学习难度且破坏时间一致性。现实中物体通常在有限状态集合间平滑过渡（如人从"站立"到"行走"再到"奔跑"）。为此，4D LangSplat 提出状态变形网络——直觉上，物体的语义状态（如"站立""行走""奔跑"）是有限的，语义变化应在这些状态之间平滑插值而非任意跳变。具体地，将高斯点 $i$ 在时刻 $t$ 的语义特征表示为 $K$ 个状态原型特征的线性组合 $f_{i,t} = \sum_{k=1}^{K} w_{i,t,k} \cdot S_{i,k}$，其中权重满足 $\sum_{k=1}^{K} w_{i,t,k} = 1$。MLP 解码器以 HexPlane 提取的时空特征为输入，预测每个时刻的权重系数 $w_{i,t,k}$，MLP 解码器与逐高斯的状态原型 $\{S_{i,1}, S_{i,2}, \ldots, S_{i,K}\}$ 联合训练。这一约束确保语义特征只能在预定义状态间平滑插值，而非任意跳转。

在推理阶段，4D LangSplat 通过两种语义场的协作实现开放词汇 4D 查询。对于时间无关查询（如"找到狗"），仅在时间无关语义场上渲染特征图像，计算其与查询文本的相关性分数并获得分割掩码。对于时间敏感查询（如"找到正在奔跑的狗"），则联合使用两种语义场：先用时间无关语义场定位目标物体的空间位置，获得初始掩码；再在掩码区域内计算时间变化语义场渲染特征与查询文本的余弦相似度，以全视频平均相似度为阈值筛选出目标动作活跃的时间段，最终保留对应时间段的掩码作为查询结果。这种两阶段协作机制使系统能同时回答"哪里有这个物体"和"何时在做这件事"两个层次的问题。

\paragraph{运动-语言高斯时序场景理解}

4D Synchronized Fields（2026.03）~\cite{2603_14301} 的出发点是一个关键洞察：现有 4D 语言方法（如 4D LangSplat）在运动优化完成后才蒸馏语言特征，语义场因此不包含物体如何运动的结构化知识，只记录了"那里有什么"而非"那里在做什么"。要真正理解动态场景，运动、语义和重建必须同步涌现。该方法分五个阶段依次完成：可变形 4D 高斯重建、物体分配、环内运动分解、反退化正则化和运动条件化语言场，其中运动结构必须先于语言学习（阶段顺序不可交换）。

在运动分解阶段，每个高斯 $i$ 的轨迹被分解为共享物体运动加隐式残差。具体而言，先用外部分割器（SAM 3）通过多视角投票为每个高斯分配物体标签 $o(i)$，然后学习一个按物体 $k$ 和时间 $t$ 参数化的 SE(3) 变换：$\hat{\mu}_i(t) = R_{o(i)}(t) \mu_i^0 + t_{o(i)}(t)$，即物体的刚体运动由所有属于该物体的高斯共享。隐式残差定义为 MLP 预测位移减去共享变换位移：$\delta_i(t) = \Delta\mu_i(t) - \hat{\mu}_i(t)$。由于 MLP 的表达能力严格大于 SE(3) 模型，不加约束时所有运动会坍缩到残差中。为此，方法引入五种正则化，核心是残差自适应调制：通过 EMA（$\alpha\!=\!0.95$）跟踪每个高斯的残差幅度，对残差大的边界或关节处高斯降低惩罚权重 $m_i = \mathrm{sigm}\!\left(\frac{\bar{r}_i - \mu_r}{\sigma_r + \epsilon}\right)$，允许非刚体区域保留形变自由度。总损失 $\mathcal{L} = \mathcal{L}_{\text{rgb}} + \lambda_{\text{res}} \mathcal{L}_{\text{res}}^{\text{mod}} + \lambda_{\text{share}} \mathcal{L}_{\text{share}} + \lambda_{\text{vel}} \mathcal{L}_{\text{vel}} + \lambda_{\text{smooth}} \mathcal{L}_{\text{smooth}}$ 在运动冻结后开始训练语言场。

语言场的核心创新在于运动条件化脊图。对于每个物体 $k$，构造 28 维运动学特征向量 $\phi_{k,f}$（编码核心运动、时间导数、变换诊断和关系上下文），然后拟合从运动学到语义残差的脊回归 $\hat{\psi}_{k,f} = \bar{\psi}_k^{\text{static}} + \phi_{k,f} W_k$，其中 $\bar{\psi}_k^{\text{static}}$ 为该物体外观锚点（所有观测帧 SigLIP 嵌入的均值）。开放词汇查询通过混合静态和时序相似度评分：$S(k; q) = w_s \langle \bar{\psi}_k^{\text{static}}, e_q \rangle + w_t \langle \bar{\psi}_{k,f}^{\text{sync}}, e_q \rangle$，权重 $w_s\!=\!0.35$，$w_t\!=\!0.65$。在 HyperNeRF 上实现 28.52 dB 平均 PSNR（语言接地和运动感知基线中最高），时间状态检索准确率达 0.884，大幅超越 4D LangSplat 和 LangSplat。

\subsubsection{HDR与光照动态}

动态场景中的光照变化——强光、弱光和交替曝光——给重建带来额外挑战。本节介绍一种 HDR 动态新视角合成方法。

\paragraph{高动态范围动态新视角合成}

HDR-4DGS（2025.09）~\cite{2509_21853} 首次提出 HDR 动态新视角合成（HDR DNVS）任务：从稀疏、时变的 LDR 输入重建时间相干的 HDR 辐射场和动态几何。现有 DNVS 方法局限于 LDR 图像，在强光或弱光场景中丢失大量信息；而现有 HDR NVS 方法仅适用于静态场景，无法处理运动物体和时变光照。核心挑战在于如何在动态场景中维持 HDR-LDR 转换的时间辐射一致性。

方法基于 4DGS 框架，将颜色表示空间从 LDR 扩展到 HDR。HDR-4DGS 的关键创新是动态色调映射器（DTM），灵感来源于人眼视网膜感光细胞对环境亮度的动态适应机制——直觉上，DTM 就像眼睛的瞳孔，随环境亮度自动调节，确保 HDR 内容在不同曝光条件下都能正确显示。DTM 首先构建辐射库：存储每帧的均值 HDR 颜色统计 $r_t^h = \frac{1}{N}\sum_{i=1}^{N} c_{i,t}^h \in \mathbb{R}^3$，然后用滑动窗口收集 $k$ 帧历史辐射序列 $\{r_{t-k:t}^h\}$，送入动态辐射上下文学习器（DRCL，可用 LSTM/GRU/Transformer 实现）生成辐射上下文嵌入 $f_t \in \mathbb{R}^d$。最终色调映射函数为 $c_t^l = g_\theta([\log\|c_t^h + \log e_t, f_t])$，即将 HDR 颜色取对数后与曝光时间 $e_t$ 和辐射上下文 $f_t$ 拼接，通过逐通道函数 $g_\theta$ 转换为 LDR 颜色。HDR 颜色乘以曝光时间的操作符合相机响应函数原理（场景辐射与观测强度的映射是曝光依赖的），而辐射上下文 $f_t$ 则确保色调映射随场景光照变化而动态调整。

训练目标为 $\mathcal{L}_{\text{total}} = \mathcal{L}_{\text{ldr}} + \lambda \mathcal{L}_{\text{hdr}}$，其中 LDR 损失同时监督直接渲染的 LDR 图像 $I_{t,3D}^l$ 和经 DTM 色调映射的 LDR 图像 $I_{t,2D}^l$，双路监督防止色调映射器过拟合。HDR 损失通过 $\mu$-law 压缩对齐 HDR 和 LDR 域进行对比：$\hat{I}^h = \frac{\log(1 + \mu \cdot \text{norm}(I^h))}{\log(1 + \mu)}$。方法还发布了 HDR-4D-Syn（8 个合成场景）和 HDR-4D-Real（4 个真实场景）两个基准数据集。

\subsubsection{特殊模态与场景}

本路线介绍面向热红外、雷达、植物生长等非典型场景的动态高斯方法。

3DGS/4DGS 的动态能力不仅适用于可见光场景，还可扩展到热红外、雷达、植物生长等特殊模态和应用。本节介绍七种面向非典型场景的动态高斯方法。

\paragraph{夜间热红外动态重建}

热红外成像具有全天候工作能力，可在无光、雾霾等恶劣条件下非接触测量物体表面温度，广泛用于安防巡检、建筑监测和能源管理。然而，现有的 3D 热场景重建方法——无论是基于 NeRF 的 ThermoNeRF，还是基于 3DGS 的 Thermal Gaussian Splatting 和 Thermal3DGS——都只针对单一时刻做静态重建，无法建模或预测场景温度随时间的变化。NTR-Gaussian（2025.03）~\cite{2503_03115} 首次将动态 4D 高斯表示引入热红外领域，目标是实现夜间户外场景的多视角、多时刻温度重建与新视角温度图合成。

该方法的核心挑战在于：夜间没有太阳辐射输入，场景温度完全由散热过程决定，且温度随时间持续变化。在标准 4DGS 中，位置和协方差随时间变形，而颜色 $c$ 和不透明度 $\alpha$ 保持常数。NTR-Gaussian 将颜色属性替换为红外辐射温度 $T$ 后，温度也必须成为时变量——否则无法表达夜间持续降温的过程——这正是 NTR-Gaussian 需要建模的关键量。方法的输入为带时间戳的多视角热红外（TIR）图像及对应的合成可见光图像。

方法分两阶段训练。第一阶段为粗粒度温度学习：Temp\_Net（MLP）建立离散时间到高斯温度的映射，输入为时间编码（对应热监督图像的采集时刻）和高斯位置编码（来自倾斜摄影模型的点云），输出每个高斯的辐射温度 $\text{Temp}(t)$。预测温度与高斯参数（位置、旋转、缩放、不透明度）联合渲染温度图，与真实温度图计算损失 $\mathcal{L} = \lambda_1 \mathcal{L}_{\text{Temp}} + \lambda_2 \mathcal{L}_{\text{SSIM}}$（$\lambda_1\!=\!0.8$，$\lambda_2\!=\!0.2$），同时优化 Temp\_Net 和高斯参数。

第二阶段引入热力学先验，是整个方法的核心创新。夜间物体的温度变化遵循物理散热规律，可以用热力学微分方程约束高斯的温度预测。夜间场景的热量释放由两部分组成：对流换热 $Q_{\text{con}} = C \cdot (T(t) - T_m)$ 和辐射散热 $Q_{\text{rad}} = E \cdot \sigma \cdot (T(t))^4$，其中 $C$ 为对流传热系数（取决于物体材质和辐射温度）、$E$ 为发射率（取决于物体材质）、$\sigma$ 为 Stefan-Boltzmann 常数、$T_m$ 为环境温度。温度变化率满足微分方程 $T'(t) = (Q_{\text{rad}} + Q_{\text{con}}) / H$，其中 $H$ 为热容。Thermal\_Net（另一 MLP）预测每个高斯的三个热力学参数 $E(t)$、$H(t)$、$C(t)$，其输入包含四路信息：Temp\_Net 在第一阶段输出的高斯温度、Feature\_GS 提取的高斯语义特征（使物理量的空间分布合理而非散乱）、从 $t_0$ 到目标时刻等分的 $N$ 段时间编码、以及倾斜摄影模型提供的位置编码。关键设计在于，第一阶段训练好的 Temp\_Net 被冻结，其输出的 $T(t)$ 直接代入微分方程作为被积函数；Thermal\_Net 通过数值积分 $T_{\text{integral}}(t_N\!=\!k) = T(t_0) + \sum_{n=1}^{k} \Delta T(t_N)$ 得到任意时刻的温度预测，并与 Temp\_Net 的输出 $T(t_N\!=\!k)$ 联合渲染温度图，通过一致性约束确保积分结果与直接预测结果对齐。训练目标在第一阶段损失基础上增加平滑损失 $\mathcal{L} = \lambda_1 \mathcal{L}_{\text{Temp}} + \lambda_2 \mathcal{L}_{\text{SSIM}} + \lambda_3 \mathcal{L}_{\text{Smooth}}$（$\lambda_3\!=\!10$），确保温度沿时间轴无突变。消融实验表明：两阶段训练显著优于单阶段（后者网络和参数同时训练，PSNR/SSIM 均下降且温度 MAE 增大）；语义特征对物理量预测的空间合理性至关重要（缺少时物理量值散乱不真实）；平滑损失有效提升未知时刻的温度预测精度。

方法还发布了首个面向夜间热成像的动态数据集 NTR，包含 4 个户外场景（2 个城市建筑场景、2 个郊区农田/池塘场景），由大疆 M300 RTK 无人机搭载 H20T 热红外相机在夜间 4 个时间段采集，共 3052 张 TIR 图像，首尾时段用于训练、中间时段用于测试。实验与 4DGS（动态对比）、Thermal Gaussian Splatting 和 3DGS（静态对比）进行基准测试，NTR-Gaussian 在 PSNR、SSIM 和温度 MAE 上均取得最优，预测温度误差控制在 1 摄氏度以内。值得注意的是，静态方法在某些特定时刻可能表现不错（因为学习了该时段的平均温度），但在其他时刻表现较差，而 NTR-Gaussian 基于物理先验的积分预测在所有时刻都保持稳定精度。

方法的局限在于：(1) 依赖倾斜摄影模型提供高斯初始位置编码，需要额外的三维重建预处理；(2) 热力学方程的简化假设（仅考虑对流和辐射，忽略风冷、蒸发等复杂因素）可能在极端天气条件下引入误差；(3) 训练需要覆盖足够时间跨度的热红外图像，数据采集成本较高。

\paragraph{热红外观角依赖变形高斯}

Veta-GS（2025.05）~\cite{2505_19138} 针对热红外新视角合成中两个特殊挑战：热图像随观角变化（传输效应、周围环境反射的热辐射）以及热红外传感器固有的低分辨率和噪声。直接将 Deformable 3DGS 应用于 TIR 图像无法捕获观角相关的热变化。

方法的核心改动是观角依赖变形场。标准 Deformable MLP 仅输入高斯位置和时间：$(\Delta\mu, \Delta r, \Delta s) = F_\Delta(\mu(\text{sg}(\mu)), \gamma(t))$，Veta-GS 在此基础上增加相机位置 $x$ 和观察方向 $v$：$(\Delta\mu, \Delta r, \Delta s) = F_\Delta(\mu(\text{sg}(\mu)), \gamma(t), x, v)$，使变形场能够建模因视角不同而产生的热辐射变化。同时引入视锥掩码策略：仅对落在当前相机视锥内的高斯施加变形，而非对所有高斯变形。这一设计基于 TIR 场景中通常不存在大尺度刚体运动，而只需处理微妙热变化的观察结果，既加速训练又提升质量。

为应对 TIR 图像的噪声和伪影，方法提出热特征提取器（TFE）和 MonoSSIM 损失。TFE 从三个角度用独立 CNN 提取特征：外观特征 $f_{\text{App}}$（原始 TIR 图像的纹理细节）、边缘特征 $f_{\text{Edg}}$（经 Canny 检测后的物体轮廓）、频率特征 $f_{\text{Frq}}$（经傅里叶变换后的高频纹理）。MonoSSIM 损失 $\mathcal{L}_{\text{Mono}} = (1 - \mathcal{L}_{\text{App}})(1 - \mathcal{L}_{\text{Edg}})(1 - \mathcal{L}_{\text{Frq}})$ 用乘积形式（而非原始 SSIM 的 $\mathcal{L}_L \mathcal{L}_C \mathcal{L}_S$）防止某一维度接近零时其他维度被低估，确保外观、边缘和频率三方面均衡贡献。

\paragraph{时空相干高斯雷达序列预测}

STC-GS（2025.02）~\cite{2502_14895} 将 3D 高斯表示引入气象雷达序列预测，这是一个与日常场景重建截然不同的问题：雷达观测的云团持续运动变化，不遵循刚体或形态约束。传统 4DGS 假设场景中大部分内容静态、仅少数主体动态，而雷达数据中所有区域都是高动态的。STC-GS 不采用 4D 高斯，而是在首帧用一组 3D 高斯重建后作为锚点，通过跟踪这些锚点在后续帧中的变化来捕获运动趋势。

STC-GS 的关键是双向重建方案，保证每个高斯在时间轴上的一致跟踪。直觉上，这类似于先用"未来"的信息校正"过去"的重建，再用校正后的"过去"正向传播。后向重建从时刻 $T$ 到 $0$ 逆序进行，仅优化高斯位置，将未来帧的信息回传到初始帧；前向重建从时刻 $0$ 到 $T$ 顺序进行，从后向重建结果和初始帧数据均匀采样初始化新高斯集合，依次优化。两个尺度约束确保时空相干：局部细节约束利用 RAFT 估计的 2D 光流融合为伪 3D 光流 $u_{\text{flow}}$，约束高斯位移与之对齐 $L_{\text{local}} = \|\Delta\mu - u_{\text{flow}}\|^2$；全局趋势约束用高斯核平滑原始雷达数据并阈值抑制估计最优位置分布 $P^*(x)$，通过能量函数 $L_{\text{global}} = \|R_{\text{reused}}(G) - P^*\|^2$ 约束高斯群的宏观分布。

重建完成后，雷达预测问题转化为高斯参数沿时间轴的外推。GauMamba 将 GRU 记忆机制集成到 Mamba 框架中：第 $l$ 层 MambaGRU 块在时刻 $t$ 接收上一层输出 $C_t^{l-1}$ 和上一时刻记忆 $H_{t-1}^l$，前者用 Mamba 的线性复杂度高效处理大量高斯 token，后者用 GRU 学习相邻帧高斯群之间的时空关系。训练目标预测高斯差分 $\Delta G_t = G_t - G_0$，对于预测阶段（$t > T_{\text{in}} + 1$）使用前一时刻的预测结果 $\hat{G}_{t-1}$ 加初始状态 $G_0$ 估计输入，实现自回归预测。空间分辨率比现有 3D 表示方法高 16 倍以上。

\paragraph{植物生长 4D 重建}

传统 4D 高斯方法在处理动态场景时，要么从规范空间学习离散变形（如 4D-GS（2024.04）~\cite{2404_00510}），要么在不同时间步引入独立高斯（如

\paragraph{4DGS：2024.03 同期工作}

4DGS（2024.03）~\cite{2403_15350}），两者都不保证跨时间步的几何对应关系。对于植物生长这类需要追踪每片叶子和枝条发育轨迹的任务，这种时间不一致性是致命缺陷。

\paragraph{GrowFlow：2026.02 同期工作}

GrowFlow（2026.02）~\cite{2602_08958} 的核心思路是将植物生长建模为一个连续动力学系统：每个高斯基元沿一条由神经 ODE 定义的向量场平滑演化，而非在离散时间步之间"跳跃"。

具体而言，每个高斯的几何参数——中心 $\boldsymbol{\mu}_i(t)$、旋转 $\boldsymbol{q}_i(t)$ 和缩放 $\boldsymbol{s}_i(t)$——被建模为时间的连续函数，通过积分一个学习到的流场来演化：
$$\boldsymbol{\theta}_i(t) = \boldsymbol{\theta}_i(t_0) + \int_{t_0}^{t} F_\theta(\boldsymbol{\theta}(\tau), \tau)\, d\tau$$
其中 $F_\theta$ 是由 HexPlane 架构加 MLP 解码器参数化的神经 ODE。该积分过程天然地约束了高斯轨迹必须沿着一致的向量场运动，提供了非约束变形场所缺乏的归纳偏置。外观参数（颜色和透明度）在时间上保持不变，因为受控捕获条件下光照恒定。

一个关键难题是如何随时间引入新几何——3DGS 的致密化是离散、不可微的操作。GrowFlow 巧妙地将问题反转：先在成熟态 $t=T$ 重建完整植物（此时结构最完整），然后从 $T$ 向 $t=0$ 学习"逆生长"过程，让多余的高斯通过缩放到零或被遮挡来"消失"。这一过程是完全可微的。训练分三个阶段：静态重建阶段用 $30{,}000$ 次迭代优化成熟植物的高斯；边界重建阶段从 $T$ 逐步反向积分 ODE，每个时间步 $300$ 次迭代并缓存高斯状态；全局优化阶段重新初始化 ODE 网络，随机采样时间步训练 $30{,}000$ 次迭代，使 ODE 能在未见时间步间插值。该方法在合成和真实植物数据集上均显著优于 Dynamic 3DGS 和 4D-GS，几何精度（Chamfer Distance）和图像质量（PSNR/LPIPS）均达到最优。\paragraph{植物时间数字孪生}

GrowSplat（2025.05）~\cite{2505_10923} 面向的场景与 GrowFlow 不同：它处理的是工业级高通量表型平台（荷兰 NPEC 温室系统，15 台固定相机三层排列）采集的真实植物多时序数据，目标是构建跨天级别的时间数字孪生。核心挑战在于植物生长同时包含刚体变换（传感器移动引起的平移旋转）和非刚体变形（叶片展开、茎秆弯曲），需要对齐模型能同时处理两者。

方法分为三步。首先，使用 Splatfacto-MCMC 训练策略（将 3DGS 视为马尔可夫链蒙特卡洛采样）在每个时间步独立重建 3D 高斯场景，仅 15 张输入图像的场景通过分割遮罩和初始点云辅助，训练速度约 3 秒/100 次迭代。然后对高斯点云执行三重过滤：丢弃过大或过小尺度的异常高斯（通过检查对数尺度范围）、移除过度拉长的高斯（通过尺度比检测）、验证旋转四元数的模以消除数值误差。过滤后估计表面法线并计算快速点特征直方图（FPFH，一种编码局部几何形状的旋转不变特征描述子）描述子。

配准阶段采用两阶段策略。粗对齐使用 FPFH 特征匹配加 RANSAC 鲁棒外点剔除，将参考点云 $P_{\mathrm{ref}}$ 与时序点云 $P_{t_k}$ 之间的刚体变换初始化。精对齐使用迭代最近点（ICP）优化变换，同时建模非刚体变形场 $\boldsymbol{d}_{t_k}$，变换模型为：
$$T_{t_k}(P_{t_k}) = \mathbf{T}_{t_k} P_{t_k} + \boldsymbol{d}_{t_k}$$
其中 $\mathbf{T}_{t_k} \in SE(3)$ 是刚体对齐，$\boldsymbol{d}_{t_k} \in \mathbb{R}^{n_{t_k} \times 3}$ 是生长引起的非刚体变形。该框架在藜麦（55 个时间点/76 天）和 Sequoia（40 个时间点）等长周期数据上验证了 4D 植物模型的有效性。

\paragraph{龙卷风 3D 重建}

该工作（2025.06）~\cite{2506_18677} 的贡献主要在数据采集和实验验证方面，方法上使用标准 3DGS 流水线。团队搭建了一个小型实验室龙卷风模拟装置：透明亚克力箱体四角设有倾斜缝隙使空气逆时针流入，顶部风扇持续抽气产生垂直拉伸涡旋，加热水盘上的干冰升华冷却水蒸气形成可见凝结雾滴，使涡旋结构可视化。8 台 Raspberry Pi 相机从不同视角同时采集影像，使用 Colmap 进行 SfM 初始化（特征提取、匹配、相机三角化），然后按原始 3DGS 的标准流程训练。由于高斯渲染固有的"漂浮物"问题——不属于物理空间的高斯伪影——团队使用 SuperSplat 工具进行手动后处理来清除这些伪影。这项工作的意义在于首次验证了 3DGS 对透明、高速旋转流体结构的重建可行性，为后续将风暴追踪者拍摄的真实龙卷风多视角视频应用于 3D 重建奠定了基础。

\paragraph{4D 天气合成实时重建渲染}

现有方法要么无法同时重建场景和模拟天气（如扩散模型仅改变外观但不产生真实的天气元素），要么无法实时渲染（如基于 NeRF 的 ClimateNeRF）。Weather-Magician（2025.05）~\cite{2505_19919} 的策略是"不在天气场景上重建，而在晴天场景上叠加天气"——先用增强版 3DGS 重建清晰场景，再通过高斯操作和后处理合成天气效果。

几何增强方面，该方法在初始化点云上手动添加上半球点覆盖以辅助天空的 RGB 和深度渲染，并引入深度+法线监督优化几何。深度监督使用 Depth Anything Model 生成伪真值 $D_{\mathrm{DAM}}$，通过硬/软深度正则化分别优化高斯位置和透明度：
$$L_{\mathrm{hard/soft}} = \|D_{\mathrm{ref}}^{\mathrm{GN}} - D_{\mathrm{DAM}}^{\mathrm{GN}}\|_2 + \|D_{\mathrm{ref}}^{\mathrm{LN}} - D_{\mathrm{DAM}}^{\mathrm{LN}}\|_2$$
法线一致性损失延迟到 $6000$ 次迭代后才引入，避免过早影响透明度优化。

天气合成分为三类。静态天气（雾/霾/烟尘）利用深度图的指数衰减关系生成模糊遮罩，公式为 $\beta_d = \min(1, 1 - e^{-I_d \cdot D_d})$，再用加权混合 $C_{\mathrm{fog}} = C_{\mathrm{fog}} \beta_d + C_{\mathrm{render}}(1-\beta_d)$ 实现距离相关的模糊效果。动态天气（降雨/降雪）将雨滴建模为拉长高斯、雪花建模为三个 $60^\circ$ 交叉的拉长高斯；关键设计是将原始场景 $G$ 和噪声粒子场景 $G_n$ 分开渲染再加权混合 $C_{\mathrm{fall}} = \sum_i f_l^i (C_{\mathrm{noise}}^i \cdot m^i) + C_{\mathrm{render}}$，并引入亮度自适应透明度因子和深度遮罩处理遮挡。累积天气（积雪）利用高斯最短轴方向计算法线，选择朝上的表面，通过局部平面估计和插值生成积雪层。整个框架在单张消费级 GPU（如 RTX 3060）上实现实时渲染，同时保持视角一致性和时间一致性。

\subsubsection{SLAM与不确定性建模}

本节将 4DGS 扩展到同时定位与建图（SLAM）和不确定性建模场景，介绍两种方法。

\paragraph{动态感知 4DGS SLAM}

现有 3DGS SLAM 方法（如 SplaTAM、MonoGS）在静态环境中表现优异，但遇到动态物体时会出现相机位姿估计漂移和地图伪影。根本原因在于它们采用静态 3DGS 地图表示，只能"过滤掉"动态物体而无法重建完整场景。D4DGS-SLAM（2025.04）~\cite{2504_04844} 转而"拥抱"动态性——用 4D 高斯作为地图表示，将时间维度纳入场景建模，使其能同时重建静态和动态元素。直觉上，如果说静态 3DGS SLAM 是用一张"快照"描述整个场景，那么 D4DGS-SLAM 则是用一段"时空录像"描述场景如何随时间演变，动态物体不再是需要剔除的噪声，而是地图的有机组成部分。

\textbf{4D 高斯地图表示。}每个 4D 高斯的概率密度定义为：
$$p(\mathbf{x}|\boldsymbol{\mu}, \boldsymbol{\Sigma}) = e^{-\frac{1}{2}(\mathbf{x}-\boldsymbol{\mu})^T \boldsymbol{\Sigma}^{-1}(\mathbf{x}-\boldsymbol{\mu})}$$
其中 $\mathbf{x} = (x,y,z,t)$ 是 4D 时空点，均值 $\boldsymbol{\mu} = (\mu_x, \mu_y, \mu_z, \mu_t)$ 表示高斯在时空中的中心位置，协方差矩阵 $\boldsymbol{\Sigma} = \mathbf{RSR}^T \in \mathbb{R}^{4\times4}$ 编码高斯的空间范围和时间持续长度。缩放矩阵 $S = \mathrm{diag}(s_x, s_y, s_z, s_t)$ 中前三维控制空间形状，第四维 $s_t$ 控制时间持续范围；旋转矩阵由两个四元数 $q_l, q_r$ 构造（4D 旋转需要两个四元数，类比 3D 旋转需要单个四元数）。

4D 高斯的核心优势在于可分解为条件 3D 高斯 $p(x,y,z|t)$（给定时刻的空间分布）和边缘 1D 高斯 $p(t)$（时间分布）：

\begin{align}
\boldsymbol{\mu}_{xyz|t} &= \boldsymbol{\mu}_{1:3} + \boldsymbol{\Sigma}_{1:3,4}\boldsymbol{\Sigma}_{4,4}^{-1}(t-\mu_t) \\
\boldsymbol{\Sigma}_{xyz|t} &= \boldsymbol{\Sigma}_{1:3,1:3} - \boldsymbol{\Sigma}_{1:3,4}\boldsymbol{\Sigma}_{4,4}^{-1}\boldsymbol{\Sigma}_{4,1:3}
\end{align}

条件均值 $\boldsymbol{\mu}_{xyz|t}$ 实现时空解耦——给定查询时间 $t$，高斯的空间位置根据时空协方差自适应偏移，自然地建模了物体的运动轨迹。在渲染时，像素 $(u,v)$ 在时刻 $t$ 的颜色通过 $\alpha$-混合 混合可见 4D 高斯的贡献，每个高斯的颜色由 4D 球谐函数（4D Spherindrical Harmonics, 4DSH）$c_i(d,t)$ 表示，将视角相关和时间相关的颜色变化统一编码在一个函数中。

\textbf{基于 LEAP 的长期有效点跟踪。}跟踪阶段的输入是 RGB-D 图像序列。系统首先通过 Sobel 算子计算 RGB 图像梯度，经平均池化平滑后，将梯度图划分为多个子区域，在每个子区域中选择梯度幅值最大的点作为锚点。这种基于梯度的分布式采样策略确保锚点在图像中均匀分布且位于特征丰富的区域。

随后将锚点和 RGB 图像序列送入预训练 LEAP（Long-term Effective Any Point Tracking）模块，LEAP 对每个查询点预测三类信息：可见性 $v_s \in \{0,1\}$（该点在当前帧是否可见或被遮挡）、动态标签 $m_d$（该点属于静态还是动态物体，由浅层 MLP 对轨迹特征和锚点特征联合预测）、可靠性 $\rho$（可见性预测的置信度，由多元 Cauchy 分布的尺度矩阵参数推导，$\rho(x_s) = \sigma_a[s,s] + \sigma_b[s,s]$）。与其他成对跟踪器不同，LEAP 能在完整图像序列上进行长期点跟踪，这对于 SLAM 系统维持跨帧一致的跟踪至关重要。

系统通过多重阈值过滤获得稳定的静态跟踪点集：可见性 $v_s \geq \tau_v=0.9$（排除被遮挡点）、动态性 $m_d < \tau_d=0.9$（排除动态物体上的点）、可靠性 $\rho \geq Q(\tau_u=0.8)$（排除预测不确定的点），以及至少具有 $\tau_{\mathrm{track}}=3$ 次有效观测。当跟踪点数量低于预设阈值时，系统从当前帧重新提取锚点补充。

相机位姿通过最小化加权光度-几何残差联合优化：
$$E_{\mathrm{total}} = \lambda_{\mathrm{pho}} E_{\mathrm{pho}} + (1-\lambda_{\mathrm{pho}}) E_{\mathrm{geo}}$$
其中光度残差 $E_{\mathrm{pho}} = \|I_f(\boldsymbol{\theta}, \boldsymbol{\Gamma}) - \bar{I}_f\|_1$ 衡量渲染图像与观测图像的差异，几何残差 $E_{\mathrm{geo}} = \|D_f(\boldsymbol{\theta}, \boldsymbol{\Gamma}) - \bar{D}_f\|_1$ 衡量渲染深度与观测深度的一致性，残差仅在过滤后的稳定静态跟踪点上计算，从根本上避免了动态物体对位姿估计的干扰。

\textbf{动态感知 4DGS 映射。}映射阶段的目标是维护一个连贯的时空结构并优化新插入的高斯。每次映射迭代中，系统使用当前关键帧窗口 $\mathcal{K}$ 重建可见区域，同时随机选择两个历史关键帧 $\mathcal{K}_r$ 以防止全局地图遗忘。

关键创新是动态感知的各向同性正则化。映射阶段利用 LEAP 前端提供的动态标签将高斯分为强动态和弱动态两组。对于强动态高斯，正则化约束时空四个维度的缩放一致性：
$$E_{\mathrm{iso}}^{\mathrm{dynamic}} = \sum_{i=1}^{|\mathcal{G}^{\mathrm{dynamic}}|} \|\boldsymbol{s}_i^{4D} - \overline{S_i^{4D}} \mathbf{1}\|_1$$
对于弱动态高斯，仅约束空间三维的缩放一致性：
$$E_{\mathrm{iso}}^{\mathrm{static}} = \sum_{i=1}^{|\mathcal{G}^{\mathrm{static}}|} \|\boldsymbol{s}_i^{3D} - \overline{S_i^{3D}} \mathbf{1}\|_1$$
这种差异化策略的效果是：静态区域的高斯在时间轴上自然展宽（$s_t$ 不受约束，用更少高斯覆盖整个时间范围），而动态区域的高斯精确匹配物体的时空出现模式（$s_t$ 被约束为紧凑的时间脉冲）。如果对所有高斯施加统一的 4D 各向同性正则化，静态区域会冗余地产生多个高斯；而使用动态感知正则化后，静态和动态区域的高斯在时间轴上呈现出合理的分布差异。

映射的完整优化目标为：
$$\boldsymbol{\theta}^k, \boldsymbol{\mu}^{3D}, \boldsymbol{\Gamma}^k = \arg\min_{\boldsymbol{\theta}^k, \boldsymbol{\mu}^{3D}, \boldsymbol{\Gamma}^k} E_{\mathrm{total}}^k + \lambda_{\mathrm{iso}}^{\mathrm{dynamic}} E_{\mathrm{iso}}^{\mathrm{dynamic}} + \lambda_{\mathrm{iso}}^{\mathrm{static}} E_{\mathrm{iso}}^{\mathrm{static}}$$

关键帧策略沿袭 MonoGS 的共视性选择和帧管理机制，但做了 4D 适配：计算共视性时将时间参数设为当前帧的时间戳，从而获得该时刻的条件 3D 高斯进行空间共视性判断。

\textbf{实验验证。}系统在两类动态场景数据集上评估：室内受控环境（BONN 数据集，含人在内的刚性/非刚性运动）和室外复杂动态环境（TartanAir-Shibuya 数据集，大规模城市街景中的密集行人流）。实验平台为 Intel i7-13700HX CPU + 单张 RTX 4080 GPU，运行速度 1.5 FPS（BONN）/ 0.9 FPS（TartanAir-Shibuya）。

跟踪精度方面，在 BONN 数据集上相比次优方法 DG-SLAM，ATE RMSE 从 5.5 cm 降低至 5.1 cm（降幅 7.3\%）。传统几何方法（如 ORB-SLAM3）在高动态序列中严重失败（ps\_tk2 场景 ATE RMSE 达 77.9 cm），神经隐式方法（NICE-SLAM、Vox-Fusion）因静态场景假设导致灾难性跟踪失败。在 TartanAir-Shibuya 上优势更加显著：相比 SplaTAM，平均 ATE 降低 95.1\%（3.6 cm vs.\ 73.2 cm），特别是在快速相机运动场景 rc06 中保持 2.1 cm 的亚厘米级跟踪精度，而基线 3DGS SLAM 系统的 ATE 超过 80 cm。

映射质量方面，在 BONN 数据集上平均 PSNR 达 28.71、SSIM 0.854、LPIPS 0.259，相比 SplaTAM 和 MonoGS 显著提升（PSNR +10.5、LPIPS $-$0.15），且所有含运动物体（如 ball）的序列上一致领先，表明 4DGS 能无重影地建模时空变化。在 TartanAir-Shibuya 上，PSNR 22.75 相比 SplaTAM 的 12.43 提升 83.1\%，SSIM 从 0.384 提升至 0.705，跟踪与映射质量同步提升的根本原因在于方法显式建模物体运动而非将其视为异常值。

\paragraph{不确定性下的 4DGS 图模型}

现有 4DGS 方法普遍假设数据质量充足——足够多的相机覆盖所有视角、良好的标定和同步。但在实际场景（如交通路口监控）中，这些假设常常不成立。GraphiXS（2026.01）~\cite{2601_19843} 的核心贡献是将 4DGS 重新表述为一个概率图模型，把相机位姿 $\mathcal{C}$、帧时间 $\mathcal{T}$、视频帧 $\mathcal{I}$ 都视为随机变量，而非常量——直觉上，这意味着承认"我不确定相机在哪里、帧是何时拍的"，并将这种不确定性显式纳入优化过程。这使得框架能显式建模视角稀疏、帧缺失、相机异步等多种数据不确定性。

GraphiXS 将 4DGS 的渲染流程重新表述为生成过程：从混合模型 $F = \sum_{i=1}^{M}\psi_i$（$M \to \infty$）出发，给定相机 $\mathcal{C}$ 和时间 $\mathcal{T}$，先采样投影到图像平面的组件子集 $\alpha_i$，再为每个像素的光线采样命中子集 $\beta_j$，最终生成颜色。原始 3DGS 可视为该框架的特例（无 $\mathcal{C}, \mathcal{T}$ 随机性，组件为高斯分布）。关键创新是所有可学习参数 $\Theta$ 被视为潜在变量，通过最大后验（MAP）推断：
$$p(\Theta | \mathcal{C}, \mathcal{T}, \mathcal{I}) \propto p(\mathcal{I}|\Theta, \mathcal{C}, \mathcal{T}) \cdot p(\Theta) \cdot p(\mathcal{C}) \cdot p(\mathcal{T})$$
其中 $p(\mathcal{C})$ 和 $p(\mathcal{T})$ 的引入使得相机位姿和时间戳的不确定性被纳入优化。

该框架可实例化为不同变体：GraphiGS 使用高斯基元 $\phi_{\mathrm{GS}}(\mathbf{x}; \boldsymbol{\mu}, \boldsymbol{\Sigma}) = \exp(-\frac{1}{2}(\mathbf{x}-\boldsymbol{\mu})^T\boldsymbol{\Sigma}^{-1}(\mathbf{x}-\boldsymbol{\mu}))$；GraphiTS 使用 Student-t 分布，额外引入可学习的自由度参数 $\nu \in [1, +\infty)$ 控制尾部厚度 $\phi_{\mathrm{TS}}(\mathbf{x}; \boldsymbol{\mu}, \boldsymbol{\Sigma}, \nu) \propto (1+\frac{1}{\nu})^{-\frac{\nu+3}{2}}$，在 Cauchy 分布（重尾、高不确定性容忍）和高斯分布之间平滑插值。优化采用随机梯度哈密顿蒙特卡洛（SGHMC，一种结合随机梯度下降和哈密顿动力学的采样优化方法），4D 组件的添加/回收通过重要性得分 $s_i^t = \lambda_{sg} g_{sg,i}^t + \lambda_o o_i^t$ 调控。在 N3DV 和 Google Immersive 数据集上的五种不确定性设置下，GraphiXS 在 198 项指标中的 172 项取得最优或次优成绩。

\subsubsection{超分辨率与匹配}

本节介绍如何突破 3DGS 的分辨率上限，以及利用 3DGS 生成的数据改进图像匹配。

\paragraph{自适应学习高斯超分辨率}

标准 3DGS 通过自适应致密化（split/clone）逐步增加高斯数量，但致密化仅在当前训练分辨率下工作，无法跨分辨率外推——要获得更高分辨率的重建只能从头重新训练，计算量随分辨率二次增长。GaussianVAE（2025.06）~\cite{2506_07897} 的核心思路是将高斯在致密化过程中的参数演化视为一条时序轨迹（lineage），用 VAE（变分自编码器，一种学习数据分布并支持采样的生成模型）学习其规律——直觉上，就像语言模型学会了"词的演变规律"后可以预测新词，VAE 学会了"高斯的演化规律"后可以预测更高分辨率下的高斯参数，从而在无需重新训练的情况下外推到更高分辨率。

具体而言，预处理阶段追踪每个高斯在连续致密化步骤中的参数变化，将其参数演化路径组织为时序序列。这一过程类似于元学习——模型不是直接观测某一分辨率的固定快照，而是学习预测高斯在分辨率变化下\emph{如何演化}，从而推断出单次观测中无法显式获取的结构和外观连续性。模型架构为 Hessian 感知 VAE：Transformer 编码器通过自注意力捕捉高斯谱系中的长程依赖（某个高斯在某一分辨率下的重要性可能仅在其跨尺度行为中才显现），解码器将几何属性（位置 $\boldsymbol{\mu}$、缩放 $s$、旋转 $q$）与外观属性（颜色 $c$、不透明度 $o$）分开重建，避免两类属性优化互相干扰。

训练策略采用粗到细（coarse-to-fine）策略——模型先学习低分辨率高斯分布，再逐步精修，改善收敛性并稳定梯度流。同时借鉴语言模型的掩码预测——遮蔽谱系末尾节点，让模型预测缺失状态，迫使其学习高斯动力学的鲁棒表示而非记忆固定分辨率快照。训练目标结合三部分损失：KL 散度损失使用归一化流先验 $p_\theta(z) = \mathcal{N}(0, I)|\det T_{\theta}^{-1}(z)|$（其中 $z = T_\theta(\epsilon)$，$\epsilon \sim \mathcal{N}(0, I)$）和循环退火调度 $\beta(t) = \min(t/R,\, 1)$ 防止后验坍缩，权重设为 $10^{-6}$；MSE 损失（权重 1.0）监督外观属性 $L_{\text{MSE}} = \|c - \hat{c}\|_2^2 + \|o - \hat{o}\|_2^2$；倒角距离（权重 0.01）衡量几何保真度 $L_{\text{chamfer}} = \sum_{p \in P}\min_{q \in Q}\|p-q\|_2^2 + \sum_{q \in Q}\min_{p \in P}\|q-p\|_2^2$。Hessian 加权优化通过分析模型在哪里分配更多高斯来隐式学习感知或结构重要区域，有效复现了 3DGS 自适应致密化行为但无需训练信号；结合动态剪枝降低内存开销而不损失质量。

值得强调的是，GaussianVAE 定位为 3DGS 的后处理插件——不修改核心 splatting 管线，而是在其输出之上预测未训练分辨率下的高斯属性，类似于对缺失信息的"补全"。推理速度 0.015 秒/步，在 600p 训练后外推至更高分辨率时，Ficus 场景几何误差降低 15\%，Mic 场景（镜面物体）降低 14\%，Drums 场景感知质量（CenseoQoE 指标）提升 11\%。外推范围可达 $+1$ 至 $+4$ 步（每步为模型的一次顺序推理），在超出训练分辨率后仍保持一致的性能，而传统 3DGS 上采样则出现明显退化。在均质纹理场景（如 Lego、Chair）上改进有限（约 1\%），原因在于谱系变化不足导致时序注意力机制缺乏足够的演化信号。

\paragraph{选择性超分辨率多视角一致 3DGS}

在低分辨率（LR）图像上训练的 3DGS 模型缺少高频信号，直接将超分辨率（SR）图像均匀应用于所有像素是最直觉的补救方案，但单图 SR 在不同视角间引入不一致的虚构纹理，导致 3D 优化接收矛盾梯度。SplatSuRe（2025.12）~\cite{2512_02172} 的核心观察是：场景中的 3D 区域并非均匀需要 SR——直觉上，就像一幅画中并非每个区域都需要同样精细的笔触——近距 LR 视角可能已包含远距视角所需的高频信息，许多区域实际上已获得充分的高频监督。

为此，SplatSuRe 先用 LR 图像训练一个 3DGS 模型获得稳定几何，然后计算每个高斯 $G^i$ 在所有训练视角中的屏幕空间半径比 $\eta^i = r^i_{\max} / r^i_{\min}$，以此近似该高斯的跨视图采样频率。高比值意味着该高斯在某些视角中被高保真采样，已有充足的高频监督；比值接近 1 则说明该区域均匀低频采样，需要 SR 补充细节。该比值经阈值 $\tau$ 偏移和 sigmoid 映射后得到高斯保真度分数 $s_G^i = \sigma(\eta^i - \tau)$（$k=0.05$ 控制过渡平滑度，阈值 $\tau=1.1$）。训练目标将 LR 真值监督与选择性加权的 SR 监督结合：$L = (1-\beta)L_{\text{LR}} + \beta L_{\text{SR}}$，其中 $\alpha=0.2$、$\beta=0.4$ 控制 SR 损失项的相对权重。高保真度区域（权重低）依靠 LR 监督保持一致性，低保真度区域（权重高）由 SR 提供额外细节。在 Tanks and Temples 数据集上，选择性 SR 策略在几乎所有指标上取得最佳表现。

\paragraph{零样本图像匹配}

图像匹配的训练数据传统上依赖深度传感器或 SfM 重建获取几何监督，但这些数据源在场景和视角多样性上有限。MatchGS（2025.11）~\cite{2511_21265} 观察到 3DGS 天然支持自由视点采样，可作为数据生成管道，但标准 3DGS 的 $\alpha$-混合 深度图因不透明度加权而引入系统性几何偏差，导致对应标签中存在大量错误。

MatchGS 的数据生成管道从两个层面修正这一问题。在深度渲染层面，标准做法是 $\alpha$-混合 加权平均深度 $D = \sum_i d_i \alpha_i \prod_{j=1}^{i-1}(1-\alpha_j)$，MatchGS 改为将每个高斯椭球沿最短轴压平为平面，渲染法线图 $N(p)$ 和距离图 $D(p)$，再通过射线-平面求交获得精确深度 $d(p) = \frac{\hat{D}(p)}{N(p)^{\top}\hat{p}}$，并辅以 $L_1$ 深度正则化抑制浮点伪影。在视角采样层面，采用基于扰动的视角生成器对训练相机外参施加可控抖动，并通过预渲染检查（渲染质量统计指标偏离均值超过 $2\sigma$ 即拒绝）确保新视角质量。这一管道在 245 个场景上生成约 168K 帧，极线误差相比传统方法降低 10--40 倍。

在表示对齐层面，MatchGS 利用高斯基元的显式 3D 属性增强 2D 匹配器。每个高斯的显式特征 $f_i^{\text{gs}} \in \mathbb{R}^{59}$ 由中心位置、不透明度、归一化缩放因子、四元数旋转和球谐系数拼接而成。粗粒度层面，将高斯聚类的 3D 特征与 2D patch 特征对齐；细粒度层面，在像素级匹配上施加 3D 几何和外观一致性约束。以 ELoFTR 为基础匹配器，在 ScanNet 上零样本 AUC 提升 13.9\%，LoFTR 变体提升 17.7\%，甚至超过在域内数据上训练的基线。

\paragraph{刚体感知单目 4DGS}

现有单目 4DGS 方法要么逐帧优化（时间一致性差），要么采用刚体假设（无法建模快速非刚体运动）。RiGS（2026.05）~\cite{2605_23672} 分析了基于刚体假设的高斯表示中动态高斯的时间持续时长分布，发现其呈现显著双峰模式：长时峰对应稳定低频运动，短时峰对应快速高频运动。这启发了将动态内容进一步分解为刚体和瞬态两种高斯基元。

RiGS 引入三种高斯基元。静态高斯 $g^s = (\boldsymbol{\mu}^s, s^s, q^s, o^s, c^s)$ 表示不变背景。刚体高斯 $g^r = (\boldsymbol{\mu}^r, s^r, q^r, o^r, c^r, w, \delta^r, \tau^r)$ 通过 $K$ 个运动基 $\{T_{j,t} \in SE(3)\}$ 建模长期低频运动，查询时间 $t$ 的变换为 $\boldsymbol{\mu}^r(t) = (\sum_{j=1}^{K}w_j T_{j,t})\boldsymbol{\mu}^r$，其中 $\|w\|_2^2=1$ 确保归一化权重。瞬态高斯 $g^t = (\boldsymbol{\mu}^t, s^t, q^t, o^t, c^t, v, \delta^t, \tau^t)$ 通过可学习速度 $v$ 捕获短期高频运动：$\boldsymbol{\mu}^t(t) = \boldsymbol{\mu}^t + v(t - \tau^t)$。两者共享软门控时间衰减函数 $o(t) = o \cdot \sigma(\delta - |t - \tau^r|)$（$\delta$ 控制边界锐度，设为 3.0），使高斯仅在有限时间窗口内活跃。在动态掩码方面，RiGS 提出物体级掩码而非逐像素掩码：对 SAM2 分割的每个物体，计算其帧级运动分数 $s_{i,t}$（结合光流不确定性和 Sampson 误差），经时间滤波 $\tau^{\text{temp}} = 10^{-4}$ 后聚合为整体运动分数 $s_i$，超过动态阈值 $\tau^{\text{dyn}}$ 的整个物体被标记为动态。这种设计避免了同一物体被静态和动态高斯不一致地建模。训练中，刚体高斯可根据 $\delta^r$ 缩短自然退化为瞬态高斯，运动基数量 $K=10$。在 Nvidia Dynamic Scenes 数据集上 PSNR 超越此前 SOTA，渲染速度约 160 FPS。

\paragraph{时间归档相机虚拟化}

4DGS 等方法在动态场景渲染速度上表现出色，但其高斯基元面向渲染质量而非几何精度优化，且时间维度通过变形场或规范空间建模，不支持精确的时间索引回溯。Time-Archival（2026.02）~\cite{2602_15181} 针对体育赛事和视觉表演场景的多相机同步采集需求，主张回归神经隐式体渲染以实现精确的时间归档功能。

该方法的场景表示采用多分辨率哈希网格编码加轻量 MLP，为每个时间步 $t$ 维护独立参数集 $\Theta_t$，场景建模为连续体积函数 $F_t: (\mathbf{x}, \mathbf{d}) \mapsto (\mathbf{c}, \sigma)$。这种独立时间步设计使得用户可精确定位并重访任意历史时刻——这恰是 3DGS 变形场方法难以保证的。体渲染沿相机射线数值积分：$C(\mathbf{r}, t) = \sum_{j=1}^{M} T_j^t (1 - e^{-\sigma_j \delta_j}) \mathbf{c}_j$，其中 $T_j^t = e^{-\sum_{k=1}^{j-1} \sigma_k \delta_k}$ 为透射率。训练损失结合光度重建和时间正则化：$L(\Theta_t) = \sum_{i=1}^{N} \sum_{\mathbf{r} \in \mathcal{R}_i} \|\hat{C}(\mathbf{r}; \Theta_t) - C(\mathbf{r})\|_2^2 + \lambda \|\Theta_{t+1} - \Theta_t\|_2^2$。多相机同步配置保证了充分的空间约束（$N$ 台相机通过 Fibonacci 球面采样均匀分布在半球上），而独立时间建模允许完全并行训练，避免了时间耦合在剧烈运动、多人物交互场景中的漂移问题。在 CMU Panoptic Studio 数据集和新合成的体育/舞蹈数据集上均验证了该方法的时间归档能力。

\paragraph{空地图像特征匹配}

空地图像融合三维建模的核心瓶颈在于两者视角差异巨大，直接特征匹配几乎不可行。现有方法要么依赖粗糙 POS 数据进行点云配准（精度有限），要么通过建筑立面矫正缩小视角差（适用范围窄）。该工作（2025.09）~\cite{2509_19898} 提出 3DGS 中间视图渲染作为桥梁：先用空拍图像通过增量 SfM（COLMAP 默认配置）重建稀疏模型并训练 3DGS 场景，再从空地视角之间的虚拟位置渲染中间图像，使匹配问题转化为两次较小视角差的匹配。

中间视点位置的确定是关键。给定空拍视角位姿 $(R_a, C_a)$，通过射线与场景求交计算地面目标坐标 $(X_t, Y_t, Z_t)$，然后将空拍相机沿垂直方向移动至 $Z_r = Z_a/2$，同时绕 $X$ 轴旋转角度 $\theta$（使主光线重新对准地面目标），得到渲染视点 $(R_r, C_r)$。渲染图像与空拍、地面图像的视角差均大幅缩小，特征匹配成功率显著提高。具体流程为：先对渲染图-空拍图和渲染图-地面图分别进行特征提取与匹配（因视角差较小，SIFT 等传统方法即可胜任），再通过渲染图像传递获得空地之间的间接对应关系。在三个标准空地数据集（CENTER、ZECHE 和 Campus）上，该方法的正确匹配数（NCM）和匹配对数（NMP）均显著优于传统图像矫正方法。

\paragraph{快速体积捕获重建管道}

Fast Volumetric Capture（2025.12）~\cite{2512_15719} 面向的核心问题是：如何用少量廉价相机在不受控环境中快速完成人体体积捕获。传统多视角系统依赖密集相机阵列和受控照明，部署成本高且灵活性差。该工作的出发点是将点云和高斯 splat 两条重建路径统一在同一预处理流程之下，从而在 6--12 台 Orbbec Femto Bolt RGB-D 相机组成的半圆弧形阵列下实现实时预览。

在点云分支中，baseline 做法是将深度图直接反投影到世界坐标系。该系统沿用这一思路，但针对野外噪声增加了半径邻域有效性检验：对于每个世界空间点 $p$，在半径 $r = 20\,\text{cm}$ 内搜索邻居，仅当邻居数 $|N(p)| \geq N_\theta = 30$ 时保留该点，孤立离群点被有效滤除。

高斯 splat 分支是该系统的核心创新。传统 GPS-Gaussian 等前馈回归器将高斯旋转预测在相机坐标系下，却将其当作世界坐标系旋转直接使用。这一错误在圆形、等高布局的相机阵列中不易暴露（因为高度各向异性的"针状"splat 恰好与训练几何对齐），但在任意稀疏布局下会产生严重的视角依赖拉伸伪影。该工作对此进行了关键的坐标系修正：设网络预测的相机坐标系下高斯旋转为 ${_g^c R}$，相机到世界的标定旋转为 ${_c^w R}$，则正确的世界坐标系旋转应为
$$_g^w R = {}_c^w R \cdot {}_g^c R,$$
即相机到世界旋转必须左乘预测的相机坐标系旋转。相应地，世界坐标系下的协方差矩阵由 ${^w \Sigma} = {}_c^w R \cdot {}^c \Sigma \cdot ({}_c^w R)^\top$ 给出，其中 ${^c \Sigma} = {}_g^c R \cdot \text{diag}(s^2) \cdot ({}_g^c R)^\top$ 是相机坐标系下的协方差。这一修正消除了视角依赖的拉伸伪影。

在尺度参数化方面，该系统也做了改进。GPS-Gaussian 原本使用 softplus 激活，该系统改为先对预激活值施加双曲正切 $\tanh$ 有界化，再进行逐轴实例归一化并施加可学习的仿射参数：
$$s_{a,i} = s_\text{max} \cdot \text{softplus}(\gamma_a \hat{z}_{a,i} + \beta_a),$$
其中 $\hat{z}_{a,i} = (z_{a,i} - \mu_a)/\sigma_a$ 是归一化后的预激活值，$\gamma_a, \beta_a$ 为可学习参数。这一改动使尺度参数能够跨不同相机距离和内参迁移。

微调阶段采用了 SSIM--Huber 重建损失与尺度熵正则化的组合：重建损失 $L_\text{rec} = \lambda_{L1} \cdot L_\text{Huber} + \lambda_\text{SSIM} \cdot (1 - \text{SSIM})$，其中 $\lambda_{L1} = 0.8$、$\lambda_\text{SSIM} = 0.2$，SSIM 在 $11 \times 11$ 像素窗口上计算；熵正则化 $L_\text{ent}$ 通过软直方图鼓励尺度分布保持高熵，防止退化的极端各向异性。总损失为 $L = L_\text{rec} + L_\text{ent}$，其中熵正则化权重 $\lambda_\text{ent} = 10^{-2}$。整套系统支持 RGB-D 和纯 RGB 输入（纯 RGB 模式通过标定双目恢复深度 $Z = fB/d$），在线预览达到 5--10\,FPS，完整框架开源。

\paragraph{高保真点云渲染}

PFGS（2024.07）~\cite{2407_03857} 要解决的问题是：给定一个仅有坐标和颜色的稀疏点云 $P = \{(p_k, c_k)\}$，如何从任意视角渲染出真实感图像。传统图形学方法直接将点投影到 2D 平面进行 z-buffer 光栅化，但点的稀疏分布导致出血表面和孔洞伪影；而基于 NeRF 的方法虽然能提供 3D 一致性，但需大量采样点和网络查询，难以实时。PFGS 的核心思路是利用 3DGS 的 splatting 机制将点的神经特征光栅化到 2D 平面，再通过多尺度渐进解码器将其解码为最终图像，从而在保持 3D 一致性的同时实现高效渲染。

具体而言，管道包含四个模块。第一步，多尺度特征提取器对输入点云进行多分辨率编码。该提取器采用多输入单输出的 UNet 架构，核心组件为 Point-Voxel (PV) 卷积模块。原始点云以递减采样率 $r_1 > r_2$ 进行均匀下采样，不同分辨率的点云被拼接为提取器的额外输入：
$$f_{1:K} = E_\theta\big((p_{1:K}, c_{1:K}),\; (p_{1:K}, c_{1:K})_{\downarrow r_1},\; (p_{1:K}, c_{1:K})_{\downarrow r_2}\big),$$
其中 $E_\theta$ 为提取网络，输出 $f_{1:K}$ 为逐点神经描述符。多分辨率输入使网络能够捕获不同感受野范围内的空间特征。

第二步，高斯回归器以逐点方式预测每个点的 3DGS 属性。该回归器将神经描述符 $f_k$、坐标 $p_k$ 和颜色 $c_k$ 拼接后映射为统一特征向量 $F_k = R_\theta(f_k, p_k, c_k)$，再通过三个独立预测头分别输出旋转四元数 $R$、尺度因子 $S$ 和不透明度 $\alpha$：
$$R = \text{norm}(H_r(F_k)),\quad S = \exp(H_s(F_k)),\quad \alpha = \sigma(H_o(F_k)),$$
其中每个预测头由两层核大小为 1 的一维卷积组成，$H_r$ 使用归一化激活，$H_s$ 使用指数激活确保尺度为正，$H_o$ 使用 sigmoid 将不透明度限制在 $[0,1]$。颜色属性 $c$ 直接保留原始点云颜色，不做回归预测。

第三步是特征光栅化。与标准 3DGS 仅光栅化颜色不同，PFGS 同时渲染两张图：一张颜色 splatting 图 $I_\text{gs}$（用点颜色作为 splat 属性）和一张特征 splatting 图 $F_\text{gs}$（用神经描述符 $f_k$ 替换颜色作为 splat 属性）。两者使用完全相同的高斯参数进行光栅化，因此在像素级精确对应。$F_\text{gs}$ 可视为 $I_\text{gs}$ 的补充信息载体。

第四步，多尺度渐进解码器将拼接的 $[I_\text{gs}; F_\text{gs}]$（共 12 维：3 维颜色 + 9 维特征）逐步解码为最终图像。解码器采用 MIMO-UNet 架构，包含三个编码块和三个解码块，以速率 $r_a, r_b$ 对输入进行两级下采样并在对应层级融合多尺度特征。关键是解码器循环运行两次：第一次输出 12 维特征作为第二次的输入，第二次才输出最终的 3 维 RGB 图像。这种循环精炼机制有效恢复了因点云稀疏而丢失的细节。

训练采用两阶段策略。第一阶段仅训练特征提取器和高斯回归器（10K 次迭代），使用 $L_1$ 渲染损失 $L_\text{gs} = \frac{1}{|\Omega|}\sum_u \|I_\text{gs}(u) - I_\text{gt}(u)\|_1$，确保回归器能产生合理的高斯参数。第二阶段联合训练全部模块（60K 次迭代），采用渐进多尺度图像损失 $L_\text{mim}$ 和频域重建损失 $L_\text{mfr}$：
$$L_\text{mim} = \sum_{i=1}^{2} \sum_{s=1}^{3} \frac{w^{(l)}}{|\Omega|}\|I_s^i(u) - I_\text{gt}^{\downarrow r_s}(u)\|_1,\quad L_\text{mfr} = \sum_{i=1}^{2} \sum_{s=1}^{3} \frac{w^{(l)}}{|\Omega|}\|\text{FFT}(I_s^i(u)) - \text{FFT}(I_\text{gt}^{\downarrow r_s}(u))\|_1,$$
其中 $l$ 为循环次数，$w^{(l=1)} = 0.75$、$w^{(l=2)} = 1$ 为权重，FFT 将信号变换到频域以强化高频细节（如边缘和纹理）的恢复。总损失为 $L_\text{total} = \lambda_1 L_\text{gs} + \lambda_2 L_\text{mim} + \lambda_3 L_\text{mfr}$，其中 $\lambda_1 = 0.75$、$\lambda_2 = 1$、$\lambda_3 = 0.25$。

\begin{takeaway}
其他动态场景方法涵盖了自动驾驶、端到端驾驶、动物重建等特定领域的动态 3DGS 应用。这些方法的共性是将动态 3DGS 的通用技术与领域特定的先验（车辆运动学约束、动物骨架模板等）相结合，核心趋势是从通用动态重建走向领域特化的解决方案。
\end{takeaway}

%=====================================================================

%=====================================================================
\section{物理仿真：从粒子动力学到多物理场}
\label{sec:06_8}

把一只高斯泼溅重建的杯子从桌上推下去——它会碎成几片？杯中液体如何飞溅？碎片的运动轨迹是否符合牛顿力学？3D Gaussian Splatting 将场景编码为数十万个离散椭球体，每个椭球只有位置、颜色、不透明度和协方差，既没有质量、也没有速度、更没有内部应力。要让这些"没有灵魂的彩色气泡"遵循物理定律运动，研究者必须在纯视觉表示和物理引擎之间架起桥梁。这一根本矛盾催生了 3DGS 物理仿真领域：如何赋予高斯基元以物理属性，使其在可微分渲染管线中产生物理可信的动态行为？

本节按材料类型和技术路线组织，涵盖物理仿真集成、流体与粒子动力学、可变形物体与布料仿真、声音与多物理场，以及其他物理方法共五个子方向。

\subsection{物理仿真与流体动力学}
\label{ss:06_5a}

本节介绍两条技术路线：（1）物理仿真集成；（2）流体与粒子动力学。

\subsubsection{物理仿真集成}

本路线介绍将物理仿真引入3DGS场景的三种耦合路线：紧耦合、松耦合和感知驱动。

将物理仿真引入 3DGS 场景，首要问题是选择何种耦合方式——高斯基元如何与物理引擎对接，以及物理属性从何而来。这一领域的技术演进经历了三个阶段：从直接将高斯基元作为 MPM 粒子参与仿真的"紧耦合"路线，到通过代理结构桥接高斯与黑盒物理引擎的"松耦合"路线，再到利用大语言模型（LLM）和视觉语言模型（VLM）自动推断物理属性的"感知驱动"路线。

\begin{table}[htbp]
\centering
\caption{3DGS 物理仿真方法对比}
\label{tab:06_physics_methods}
\resizebox{\textwidth}{!}{
\begin{tabular}{lccccc}
\hline
\textbf{方法} & \textbf{耦合策略} & \textbf{物理引擎} & \textbf{材料类型} & \textbf{属性获取} & \textbf{速度} \\
\hline
PhysGaussian & 紧耦合（MPM） & 显式 MPM & 弹性/塑性 & 手动 & -- \\
i-PhysGaussian & 紧耦合（隐式 MPM） & 隐式 Newton-GMRES & 弹性/塑性 & 手动 & 20$\times$ 更大 $\Delta t$ \\
GaussianFluent & 紧耦合（CD-MPM） & CD-MPM & 断裂/流体 & 手动+生成 & 实时 \\
GASP & 松耦合 & 任意黑盒 & 多材料 & 黑盒引擎 & -- \\
Sim Anything & 感知驱动 & MPM & 多材料 & MLLM 零样本 & 2 分钟 \\
PhysFlow & 感知驱动 & 可微分 MPM & 7 种材料 & GPT-4 + 光流 & -- \\
NeuMA & 端到端 & Particle-GS & 弹性 & 物理先验+修正 & -- \\
\hline
\end{tabular}
}
\end{table}

\paragraph{MPM 紧耦合路线}

物质点法（Material Point Method, MPM）是连接 3DGS 与物理仿真最自然的选择——高斯基元本身就是三维空间中的点，可以直接充当 MPM 粒子。MPM 的核心思想是在拉格朗日粒子和欧拉网格之间交替传递信息：粒子携带质量、速度和变形梯度等物理量，在每一步中将信息转移到背景网格上求解动量方程，再将更新后的速度和位置传回粒子。

PhysGaussian（2023.11）~\cite{2311_12198}修改了高斯基元与物理仿真的对接方式——用\textbf{统一的粒子表示}替代传统的\textbf{仿真-渲染分离管线}。在标准 3DGS 中，高斯基元仅是渲染基元，不携带任何物理属性；若要为场景生成物理动态，传统管线需要先将高斯转换为四面体网格或笼状代理（cage mesh），再在仿真器中驱动网格运动，最后将结果映射回渲染表示——这一多阶段流程引入了分辨率不匹配和仿真-渲染间的离散化误差。PhysGaussian 的核心洞察是：高斯基元本身就是三维空间中的点，可以直接充当 MPM 粒子，从而实现``所见即所仿''（what you see is what you simulate, WS$^2$）。

具体而言，PhysGaussian 为每个高斯基元赋予质量 $m_p$、初始体积 $V_p^0$、速度 $v_p$ 和\textbf{变形梯度} $F_p$（$F_p$ 描述该粒子从初始材料状态到当前世界状态的局部仿射变换，包含拉伸、旋转和剪切信息；可进一步分解为弹性部分和塑性部分 $F = F^E F^P$ 以支持永久形变）。MPM 仿真循环中，粒子先将质量和动量通过 B-spline 核函数 $w_{ip}$ 传递到背景欧拉网格，在网格上求解动量方程得到更新速度 $\mathbf{v}_i^{n+1}$，再传回粒子：

\begin{equation}
\frac{m_i}{\Delta t}(\mathbf{v}_i^{n+1} - \mathbf{v}_i^{n}) = -\sum_p V_p^0 \frac{\partial \Psi}{\partial \mathbf{F}}(\mathbf{F}_p^{E,n}) {\mathbf{F}_p^{E,n}}^\top \nabla w_{ip}^n + \mathbf{f}_i^{\text{ext}},
\end{equation}

其中 $\Psi$ 为超弹性势能密度，等号右侧求和项将各粒子的内力（应力）汇聚到网格节点，$\mathbf{f}_i^{\text{ext}}$ 为外力。网格速度回传后，粒子的位置和变形梯度同步更新：

\begin{equation}
\mathbf{x}_p^{n+1} = \mathbf{x}_p^n + \Delta t \, \mathbf{v}_p^{n+1}, \quad \mathbf{F}_p^{E,n+1} = (\mathbf{I} + \Delta t \, \nabla \mathbf{v}_p) \, \mathbf{F}_p^{E,n}.
\end{equation}

这意味着每一步仿真不仅移动了高斯的位置，还更新了其变形状态。

关键创新在于如何将变形梯度映射为高斯的协方差矩阵。直观地说，如果一个椭球在材料空间中是 $A_p$，经过局部仿射变换 $F_p$ 后，它在世界空间中的形状就变成 $F_p A_p F_p^\top$——就像捏住一块橡皮泥的两端拉伸，整个形状随之变形。PhysGaussian 利用这一关系直接推导出仿真后高斯的协方差：

\begin{equation}
\mathbf{a}_p(t) = \mathbf{F}_p(t) \, \mathbf{A}_p \, \mathbf{F}_p(t)^\top,
\end{equation}

其中 $\mathbf{A}_p$ 为初始材料空间中的协方差（即标准 3DGS 重建后的原始椭球形状），$\mathbf{a}_p(t)$ 为仿真 $t$ 时刻的世界空间协方差。该映射的关键优势在于：仿真的输出可以直接用于渲染，无需任何格式转换——MPM 粒子就是高斯基元，变形梯度就是协方差的更新规则。不透明度 $\alpha_p$ 和球谐系数 $C_p$ 在仿真过程中保持不变，视为材料的固有视觉属性。该方法支持弹性体、塑性金属、非牛顿流体和颗粒材料等多种本构模型，并通过回映映射（return mapping）$\mathcal{Z}$ 将超出弹性区域的变形梯度投影回屈服面内以处理塑性流动。

然而，标准 3DGS 重建只在高斯表面放置粒子，内部是空心的——一个空心杯子在物理仿真中会塌陷。因此 MPM 耦合方法都需要\textbf{内部粒子填充}（interior filling）：通过体素化分析判定内部区域，在新填充的体素中初始化高斯粒子，并从最近的表面高斯继承外观参数。

\paragraph{i-PhysGaussian：2026.02 同期工作}

i-PhysGaussian（2026.02）~\cite{2602_17117}从 PhysGaussian 出发，针对其显式时间积分的稳定性瓶颈提出根本改进。显式 MPM 使用步初量直接预测步末状态（前向 Euler），时间步 $\Delta t$ 受到 CFL 条件严格限制，超出阈值即产生数值发散。i-PhysGaussian 将时间积分替换为隐式 Newton-GMRES 求解器：在每一步中，不直接接受显式更新结果，而是求解一个动量残差方程
\begin{equation}
R_I(\mathbf{u}) = f_I^{\text{ext}}(\mathbf{u}) + f_I^{\text{int}}(\mathbf{u}) - m_I \, a_I^{n+1}(\mathbf{u}_I),
\end{equation}

其中 $\mathbf{u}$ 为网格位移增量，$f_I^{\text{int}}$ 和 $f_I^{\text{ext}}$ 分别为内力和外力，$a_I^{n+1}$ 为 Newmark 关系下的步末加速度。该残差通过 Newton 外循环（线性化残差）和 GMRES 内循环（矩阵无关的 Krylov 子空间求解）联合优化至收敛，使步末状态严格满足动量平衡。实际效果是：在高达 20 倍更大的时间步下仍保持稳定，且对材料刚度不敏感——这意味着同样的设置可以处理从软橡皮泥到硬金属等不同刚度的材料，无需针对每种材料反复调试时间步。

\paragraph{GaussianFluent：2026.01 同期工作}

GaussianFluent（2026.01）~\cite{2601_09265}在 MPM 紧耦合路线的基础上，首次在 3DGS 框架中实现了\textbf{脆性断裂}和\textbf{混合材料}仿真。断裂仿真的核心挑战是标准 3DGS 缺乏内部结构——切开一个西瓜需要看到红色的瓤和黑色的籽。GaussianFluent 提出基于生成模型的无训练内部纹理合成管线：先将物体沿 X 轴均匀切片，对每片使用 MVInpainter 和 SD-XL 进行纹理修复，再沿三个正交轴迭代细化以确保三维一致性。在物理引擎方面，GaussianFluent 引入连续损伤 MPM（CD-MPM）替代标准 MPM：为每个粒子跟踪损伤变量 $\kappa$，当变形梯度的体积拉伸 $p$ 和剪切畸变 $q$ 超出屈服面 $y(p, q; p_0, \kappa, M) = 0$ 时，粒子的应力生成能力随 $\kappa$ 递增而衰减。关键创新是\textbf{连续回映映射}（continuous return mapping）：原始 CD-MPM 的屈服椭圆投影在 $p = p_0$ 边界处存在不连续跳跃，一个机器精度波动可能导致完全不同的投影结果；GaussianFluent 引入动态投影中心 $p_c' = p_c + \delta_k(p^{\text{tr}})(p^{\text{tr}} - p_c)$ 替代固定中心，当 $k=2$ 时实现光滑连续的映射。通过为物体的不同部位（如西瓜的籽、瓤、皮）赋予不同 $\kappa$ 值，GaussianFluent 实现了混合材料的差异化断裂行为。

沿同一紧耦合路线，PhysMorph-GS 被提出。

\paragraph{PhysMorph-GS：2025.11 同期工作}

PhysMorph-GS（2025.11）~\cite{2511_16988}修改了 PhysGaussian 的开环仿真范式——用目标驱动的可微分优化闭环替代手动设定物理参数后直接前向仿真。PhysGaussian 中，用户设定材料参数后仿真按自身规律运行，结果如何全凭物理决定；PhysMorph-GS 的目标则是指定一个目标形状（如从球体变形为牛），让物理仿真自动产生走向该目标的变形轨迹。其核心洞察是：渲染梯度应当通过变形梯度 $F$（而非粒子位置 $x$）注入仿真循环。直觉上，MPM 的驱动力来自应力，而应力由 $F$ 决定——直接在 $F$ 空间施加梯度与物理求解器天然对齐，无需跨空间（位置$\to$应力）的任意缩放；相比之下，$L/x$ 仅作用于表面附近粒子，且与物理步的应力更新缺乏量纲对应，容易引发数值不稳定。此外，PhysMorph-GS 引入 Chamfer 引导塑性（Chamfer-guided plasticity）：当弹性回复力威胁抵消已获得的变形时，通过乘法分解 $F = F_e F_p$ 逐步迁移材料静止状态，将物理从"阻力"转化为"盟友"。

\paragraph{Resonance4D：2026.04 同期工作}

Resonance4D（2026.04）~\cite{2604_01994}与 PhysMorph-GS 的目标不同——后者是给定目标形状做主动变形，前者则是从观测视频中反推物理参数并生成动态响应。与 PhysGaussian 的"给定参数$\to$仿真$\to$渲染"单向管线相比，Resonance4D 实现了反向路径：从视觉观测反演物理参数。该方法的关键瓶颈在于运动监督信号的质量。此前方法依赖在线视频扩散模型提供帧级运动引导，峰值 GPU 内存高达 35\,GB。Resonance4D 提出双域运动监督（Dual-domain Motion Supervision, DMS）替代这一重计算开销：空间域使用 SSIM 约束帧级结构一致性；频谱域对帧差序列执行分块 3D FFT，匹配局部频谱的幅度（$\lambda_{\text{mag}}$）和相位（$\lambda_{\text{phase}}$）。帧差设计抑制静态外观、突出运动模式，分块策略避免全局 FFT 将运动信号与背景混合。这一轻量级监督将峰值 GPU 内存从 35\,GB 降至约 20\,GB，首次实现全参数（杨氏模量、泊松比、屈服应力、塑性粘度、内摩擦角、密度共 6 项）的联合优化。Resonance4D 还引入基于拉丁超立方采样（LHS）的仿真驱动初始化策略：在参数空间的对数域采样多组候选，各运行短仿真后按 MS-SSIM 选择最优起点，显著提升多参数联合优化的收敛稳定性。

\paragraph{PIDG：2025.11 同期工作}

PIDG（2025.11）~\cite{2511_06299}走出了一条与 PhysGaussian 完全不同的技术路线：不使用 MPM 等显式物理引擎，而是从连续介质力学出发，将每个高斯粒子视为具有时变本构参数的拉格朗日物质点。PhysGaussian 通过 MPM 的粒子-网格传输求解运动方程，PIDG 则将物理约束直接编码为损失函数：Cauchy 动量残差 $r(x,t) = \rho(\frac{\partial v}{\partial t} + (v \cdot \nabla)v) - \nabla \cdot \sigma$ 衡量每个粒子的动量平衡偏离，其 $L_2$ 范数作为物理正则化项 $L_{\text{CMR}}$ 加入优化。该方法的核心设计是时变材料场（time-evolving material field）：一个多感知机联合预测每个粒子的速度 $v$ 和本构应力 $\sigma$ 的六个独立分量，无需指定具体的本构模型（弹性/流体/刚体），网络自动学习适合当前场景的应力-应变关系。在场景表示层面，PIDG 采用静态-动态解耦的 4D 分解哈希编码（四个 3D 哈希网格分别覆盖 xyz、xyt、yzt、xzt），先联合优化静态几何与动态变形，再冻结静态部分仅精调动态区域。在监督层面，PIDG 提出拉格朗日粒子流匹配：将高斯基元按速度场推进后投影到图像平面，生成"速度流"并与从光流分解得到的运动流对齐，作为数据拟合项加速物理场收敛。整体上，PIDG 实现了物理一致的单目动态重建——不依赖多视角输入，仅凭单目视频即可恢复具有物理意义的粒子运动。

\paragraph{PG-3DGS：2026.05 同期工作}

PG-3DGS（2026.05）~\cite{2605_11266}代表了一条更激进的路线：不再将物理仿真视为 3DGS 的后处理步骤，而是将可微分物理模拟直接嵌入 3DGS 的优化过程中，使生成的场景不仅视觉准确，还满足物理功能目标（如茶壶能倒水、机翼能产生升力）。与 PhysGaussian 的"先重建再仿真"范式不同，PG-3DGS 将物理仿真纳入重建目标函数本身：优化目标 $J(\theta; \mathcal{I}, \mathcal{P}) = L_{\text{vis}}(\theta; \mathcal{I}) + L_{\text{phys}}(\theta; \mathcal{P}) + L_{\text{reg}}(\theta)$ 在视觉损失之外增加物理损失 $L_{\text{phys}} = \mathcal{P}(\Phi(\theta))$，其中 $\Phi(\theta)$ 是从高斯几何到物理仿真网格的转换，$\mathcal{P}$ 是任务相关的物理功能评价（如倒出的液体体积、产生的升力）。直觉上，视觉损失只能约束外观——两个外观相同的茶壶，一个壶嘴连通内部可以倒水，另一个壶嘴只是外部凸起无法出水，仅靠图像无法区分。PG-3DGS 首次在 3DGS 优化中引入这类物理约束，通过可微分流体仿真将物理梯度反传至高斯参数，使优化过程自动调整几何以满足功能需求，生成的场景可直接用于物理交互而无需人工后处理。\paragraph{松耦合与代理结构路线}

紧耦合路线将高斯基元直接作为物理粒子，优势是无需格式转换，但限制为 MPM 一种物理引擎。许多实际场景需要更灵活的物理仿真器——例如刚体碰撞、有限元分析或商业游戏引擎。松耦合路线通过引入代理结构（proxy structure）桥接高斯基元与任意物理引擎。

与紧耦合路线将高斯基元直接作为 MPM 粒子不同，GASP（2024.09）~\cite{2409_05819}修改了高斯基元与物理引擎的对接方式——用代理点结构替代高斯基元直接参与仿真，从而将物理仿真从 MPM 单一引擎解放为任意黑盒物理引擎。紧耦合路线的局限在于：PhysGaussian 等方法需要修改 MPM 内核以适应高斯参数（如用一阶 Taylor 近似传播变形梯度），因此无法接入 Blender、Genesis 等通用仿真软件。GASP 的核心思路是：既然物理引擎只关心点的运动，那就只把点的位置交给物理引擎，再把运动结果映射回高斯参数——引擎是黑盒，GASP 不需要知道引擎内部如何计算。

\paragraph*{代理结构：扁平高斯→三角面片→物理顶点}

GASP 采用 GaMeS 参数化（一种扁平高斯表示，即令高斯在某个轴上缩放为 0），将每个 3D 高斯表示为一个三角面片。直觉上，一个扁平高斯就像一张"扑克牌"，它由中心点 $\mu$ 和两个沿主轴的边长 $s_2, s_3$ 确定，可以自然地转化为三角形的三个顶点：

\begin{equation}
V = [v_1, v_2, v_3], \quad v_1 = \mu, \quad v_2 = \mu + s_2 \, r_2, \quad v_3 = \mu + s_3 \, r_3,
\end{equation}

其中 $r_2, r_3$ 是旋转矩阵 $R$ 的两个列向量。这三个顶点就是 GASP 交给物理引擎的代理点。反之，物理引擎输出变形后的顶点 $\hat{V}$ 后，GASP 通过逆映射 $\phi^{-1}$ 从 $\hat{V}$ 恢复高斯参数——中心取 $v_1$，方向向量由 Gram-Schmidt 正交化计算，缩放取边长。

\paragraph*{黑盒集成与高斯控制}

由于物理引擎只操作离散点，GASP 能与 Genesis、Blender、Taichi Elements 等任意黑盒引擎无缝集成：引擎输出每步的顶点位置变化 $\hat{V}_t = \{(V_i, t)\}$，GASP 将其映射回高斯参数 $\{(\phi^{-1}((V_i, t)), o_i, c_i)\}_{i=1}^P$。然而，当物理引擎独立处理每个顶点时，某些点可能出现速度波动，导致三角面片急剧膨胀，进而产生巨大的高斯基元伪影。GASP 提出高斯控制规则（Gaussian Control Rules）：当变形后顶点间距超过原始间距的 $\lambda$ 倍时（$\lambda$ 为超参数），自动将缩放值钳制为原始值乘以 $\lambda$，即
\begin{equation}
\| \hat{v}_1(t) - \hat{v}_i(t) \| > \lambda \, \| v_1 - v_i \|, \quad i = 2, 3 \quad \Rightarrow \quad s_i \leftarrow s_i \cdot \lambda.
\end{equation}

此外，GASP 引入 Core-Gaussians 层次化策略：通过 BIRCH 聚类将高斯分为 Core-Gaussians（每簇的代表）和 Sub-Gaussians（簇内成员），物理仿真仅在 Core-Gaussians 上执行，然后将运动模式传递给 Sub-Gaussians，显著减少参与仿真的点数。

与 GASP 将高斯固定映射到三角面片顶点的做法不同，MaGS 提出了更灵活的方案。

\paragraph{MaGS：2024.06 同期工作}

MaGS（2024.06）~\cite{2406_01593}放松了高斯与网格之间的刚性绑定，提出网格吸附高斯表示（Mesh-adsorbed Gaussian Splatting），允许 3D 高斯在网格附近游走而非钉死在顶点上，形成互相吸附的网格-高斯混合表示。这一设计的动机在于：GASP 的固定映射在大幅变形时可能产生高斯扭曲或撕裂，而游走式吸附让高斯保持渲染质量的同时跟随网格运动。网格提供物理仿真的结构化拓扑，高斯提供高质量渲染的灵活性。配合 RMD-Net（残差运动解码器）、RGD-Net（残差几何解码器）和 MPE-Net（运动先验编码器）三个专用网络，MaGS 能从视频中学习运动先验并泛化到新的变形。

从 MaGS 的网格吸附策略继续往前，PhysTalk 进一步简化了代理表示。

\paragraph{PhysTalk：2025.12 同期工作}

PhysTalk（2025.12）~\cite{2512_24986}进一步简化了代理表示：GASP 使用三角面片作为代理、MaGS 使用吸附网格作为代理，而 PhysTalk 用更轻量的凸包替代两者，并将控制入口从物理引擎参数提升到自然语言。具体而言，PhysTalk 利用 LLM 将自然语言翻译为物理引擎的可执行代码，通过轻量级凸包代理将物理模拟结果传递给 3DGS。整个过程无需训练：LLM 输出代码驱动物理引擎，凸包代理映射物理状态到高斯位置，渲染管线输出最终画面。这是首个将 3DGS 直接与物理模拟器耦合（无需网格提取）的语言驱动框架。\paragraph{物理属性感知路线}

上述方法假设物理属性（质量、弹性、摩擦系数等）已知或手动设定。然而，从视觉观察中自动推断物理属性是实现"全自动物理仿真"的关键一步。物理属性感知路线利用多模态大语言模型（MLLM，包括视觉语言模型 VLM）的先验知识或可微分物理引擎的梯度反传来自动估计物理参数。

Sim Anything（2024.11）~\cite{2411_12789}修改了3DGS物理仿真中的属性设定环节——用MLLM零样本视觉推理替代手动指定或视频扩散模型迭代优化物理参数，首次实现从静态场景到可交互仿真场景的全自动转换。

前述方法面临一个共同瓶颈：PhysGaussian等紧耦合方法需要手动设定材料参数，PhysDreamer等视频扩散方法虽能自动估计但计算昂贵、可控性差，且仅适用于非刚性小物体。人类仅凭视觉就能推断物体的物理属性（金属重而硬、橡胶轻而弹），Sim Anything 正是利用多模态大语言模型（MLLM）模拟这一推理过程。

其流程分为三个阶段。第一阶段，MLLM-P3（MLLM-based Physical Property Perception）对场景中的物体进行物理属性的均值预测：先通过 BLIP 生成物体描述文本，再由 GPT-4V 从图像和文本中推理出 $K$ 个候选材料，最后通过 CLIP 相似度匹配选择最可信材料，输出均值物理属性向量 $\mathbf{M} = (\rho, E, \nu)$（密度、杨氏模量、泊松比）。

第二阶段，MPDP（Material Property Distribution Prediction）将属性预测从逐点回归重构为概率分布估计。直觉上，即便同一材料的不同区域也存在物理属性的固有差异，而 MLLM 只能给出全局均值。MPDP 训练一个网络 $D_\theta$，以物体的3D点云位置和 MLLM 预测的均值为输入，预测每个高斯粒子的物理属性分布 $P = D_\theta(\mathbf{X})$，再将分布 $P$ 与全局均值逐元素相乘，得到每个点的最终属性值。这种"均值缩放分布"的设计使 MPDP 的计算开销仅为 Physics3D 的约 $2\%$。

第三阶段，PGAS（Physical-Geometric Adaptive Sampling）解决仿真效率问题。3DGS场景通常包含百万级高斯，直接全部仿真计算量巨大。PGAS 的直觉是：软物体和形状复杂的区域需要更多驱动粒子才能精确模拟形变，而硬物体和平坦区域则不需要。PGAS 自适应调整采样半径 $r$：

\begin{equation}
r_{\text{adapt}} = \text{clip}\!\left(r_0 \cdot k \cdot \frac{E}{\kappa},\; V_{\min},\; V_{\max}\right),
\end{equation}

其中 $E$ 为杨氏模量，$\kappa$ 为物体表面曲率（由高斯位置协方差矩阵的特征值计算），$k=10$，$V_{\min}=1$，$V_{\max}=10$。软材料（$E$ 小）和高曲率区域（$\kappa$ 大）获得更小的采样半径，即更密集的驱动粒子。采样后的粒子通过 MLS-MPM 驱动物理仿真，最终整个流程在单 GPU 上约 2 分钟内完成——将静态 3DGS 场景转化为可交互的物理仿真场景。

沿同一方向，PUGS 聚焦于属性维度扩展。

\paragraph{PUGS：2025.02 同期工作}

PUGS（2025.02）~\cite{2502_12231}聚焦于 Sim Anything 未充分探索的属性维度扩展——Sim Anything 仅预测密度、杨氏模量和泊松比三个力学属性，而 PUGS 利用 3DGS 重建物体并结合 VLM 零样本预测多种物理属性（质量、摩擦系数、硬度），引入几何感知正则化损失和区域感知特征对比损失提升形状质量与局部亲和性，在 ABO-500 质量预测基准上达到新 SOTA。

\paragraph{PhysGS：2025.11 同期工作}

PhysGS（2025.11）~\cite{2511_18570}则针对 Sim Anything 的 MLLM 点估计缺乏不确定性量化这一局限，将物理属性估计建模为 3DGS 上的贝叶斯推理问题：使用 Dirichlet-Categorical 模型融合多视角材料标签，Normal-Inverse-Gamma 先验估计连续物理属性并建模认知不确定性和偶然不确定性，质量估计精度提升 22.8\%，硬度误差降低 61.2\%。

与 Sim Anything 依赖 MLLM 零样本推理的范式不同，PIXIE 采用了监督学习路线。

\paragraph{PIXIE：2025.08 同期工作}

PIXIE（2025.08）~\cite{2508_17437}更进一步，提出首个通过监督学习训练通用神经网络预测 3D 场景物理属性的前馈框架——将属性预测从"每次推理"提升为"一次训练、处处适用"。其核心洞察是利用 CLIP 等预训练视觉特征将外观与物理特性关联——金属看起来光滑且反光，橡胶看起来粗糙且哑光。构建了 PIXIEVERSE 数据集（最大的 3D 资产与物理标注配对数据集），比测试时优化方法快数个数量级且效果提升 1.46--4.39 倍。

在 PIXIE 将属性预测从优化范式转向前馈范式之后，后续工作沿加速和输入简化两个方向展开。

\paragraph{FastPhysGS：2026.02 同期工作}

FastPhysGS（2026.02）~\cite{2602_01723}针对 Sim Anything 约 2 分钟的端到端流程，聚焦于加速物理仿真流程：提出实例感知粒子填充（IPF）使用蒙特卡洛重要性采样高效填充内部粒子，双向图解耦优化（BGDO）通过自适应策略快速优化 VLM 预测的材料参数，在 1 分钟内完成高保真物理仿真。

\paragraph{ReconPhys：2026.04 同期工作}

ReconPhys（2026.04）~\cite{2604_07882}与 Sim Anything 需要预建 3DGS 场景再推断属性的两步流程不同，提出首个前馈框架从单目视频联合推断物理属性和 3DGS，采用弹簧-质点系统绑定 3D 高斯，推理速度 $<$1 秒，未来预测 PSNR 达 21.64（vs.\ 优化基线的 13.27）。

上述方法均假设场景中只有单一材料，PIG 将问题推进到多材料物理交互。

\paragraph{PIG：2025.06 同期工作}

PIG（2025.06）~\cite{2506_07657}将问题推进到 Sim Anything 未触及的多材料物理交互：通过 2D 像素到 3D 高斯的快速精确映射实现物体级分割，为不同物体分配独特物理属性进行多材料耦合交互，并提出基于约束尺度的变形梯度方法消除渲染伪影。

\paragraph{Mirage2Matter：2026.02 同期工作}

Mirage2Matter（2026.02）~\cite{2602_00096}则将视野从单物体仿真扩展到完整世界模型——与 Sim Anything 停留在静态场景的物理化不同，Mirage2Matter 将 3DGS 重建与生成模型结合构建物理真实的仿真世界模型，利用标定靶实现重建场景与真实世界的精确尺度对齐，VLA 模型在该仿真数据上训练后实现零样本性能。

\subsubsection{流体与粒子动力学}

本路线介绍将3DGS高斯基元用作拉格朗日粒子进行流体模拟以及从视频观测恢复流体动力学的方法。

流体是物理仿真中最具挑战性的材料类型之一——不可压缩性、涡旋结构和自由表面等特征对空间表示提出特殊要求。3DGS 的离散高斯基元天然适合作为拉格朗日粒子跟踪流体运动，但其各向异性椭球形状与流体的连续速度场之间存在表示鸿沟。本节的方法沿两条路线展开：一条从流体模拟器出发，将高斯基元用作流体粒子；另一条从重建出发，从视频观测中恢复流体动力学。

\paragraph{基于高斯空间表示的流体模拟}

Gaussian Fluids（2024.05）~\cite{2405_18133}首次将 3DGS 的高斯空间表示引入流体模拟领域——用连续高斯核加权和替代传统的欧拉网格或拉格朗日粒子离散化，构建无网格（grid-free）流体求解器。

\textbf{Baseline 问题。}传统流体求解器依赖两种空间表示：欧拉方法在固定网格上存储速度值，受限于网格分辨率且容易引入数值黏性（插值导致非物理耗散）；拉格朗日方法（如 SPH）用粒子跟踪流体运动，但精度有限，难以捕捉精细涡旋结构。混合方法虽试图结合两者优势，但在两种离散化之间传递数据时引入数值误差，且需维护双重表示的额外开销。隐式神经表示（INR）虽提供连续表示，但优化收敛慢、难以施加硬物理约束。

\textbf{核心表示：高斯空间表示（GSR）。}Gaussian Fluids 的关键洞察是：将流速场视为一堆"柔软的高斯斑点"的叠加——每个高斯核覆盖一块局部区域，携带一个速度向量，所有核的贡献加起来构成完整的速度场。直觉上，这就像用一组可变形、可移动的"速度画笔"来描绘流场，而非用固定网格点。具体地，连续流速场建模为多个高斯函数的加权和：

\begin{equation}
\tilde{\mathbf{v}}(\mathbf{x}) = \sum_{i=1}^{N} \mathbf{v}_i \, \hat{G}_i(\mathbf{x}),
\end{equation}

其中 $\hat{G}_i$ 为第 $i$ 个截断高斯核（在远离中心处截断为零以加速计算），$\mathbf{v}_i \in \mathbb{R}^m$ 为该核携带的权重向量。每个高斯核自身有参数 $\Theta = \{\mathbf{v}_i, \boldsymbol{\mu}_i, \theta_i/\mathbf{r}_i, \mathbf{s}_i^{-1}\}$（权重、位置、旋转、逆尺度），通过优化来拟合和演化流场。也就是说，流场在每个空间点 $\mathbf{x}$ 处的速度，等于所有"覆盖"该点的高斯核的贡献之和。

\textbf{求解流程。}Gaussian Fluids 遵循经典的算子分裂框架，每个时间步分三阶段：(1) \emph{重播种}（Reseeding）——将过度拉伸的高斯粒子分裂（最大尺度超过最小尺度的 $r_\mathrm{aniso}=1.5$ 倍时分裂），补充自适应分辨率；(2) \emph{平流初始猜测}——将高斯粒子视为拉格朗日粒子，沿速度场用四阶 Runge-Kutta 平流至新位置，作为优化的初始猜测；(3) \emph{物理引导优化}——以涡度损失 $\mathcal{L}_\mathrm{vor}$（匹配平流后的涡旋场）和散度损失 $\mathcal{L}_\mathrm{div}$（施加不可压缩约束 $\nabla \cdot \mathbf{v} = 0$）为核心，辅以各向异性正则 $\mathcal{L}_\mathrm{aniso}$（限制高斯长宽比 $\leq 1.5$）、体积正则 $\mathcal{L}_\mathrm{vol}$（保持粒子体积均匀）和位置惩罚 $\mathcal{L}_\mathrm{pos}$（防止粒子聚集），联合优化所有参数。一个关键技术细节是梯度投影：当涡度损失和散度损失的梯度方向矛盾时，将两者投影到彼此正交的方向上，避免梯度冲突导致的振荡伪影。

\textbf{GSR 的三重优势。}(1) 内存高效——不需要密集网格，高斯自适应分布到需要精细分辨率的区域；(2) 连续可微——速度场在任意位置都有定义且可直接解析求导，无需自动微分，无网格插值误差；(3) 空间自适应——高斯可以稀疏分布在低涡旋区域、密集分布在需要精细解析的高涡旋区域。相比隐式神经表示，GSR 在稳定性、精度和跨场景泛化性上均更优。该方法在 Taylor 涡旋、Kármán 涡街、3D 涡环碰撞等多种 2D/3D 流体场景中验证了精细涡旋结构的保持能力，且无需逐场景调参。

与 Gaussian Fluids 从\emph{模拟器设计}出发的路线不同，LagrangianSplats 从\emph{流体重建}角度切入。

\paragraph{LagrangianSplats：2026.05 同期工作}

LagrangianSplats（2026.05）~\cite{2605_09299}从\emph{流体重建}的角度提出结构化约束框架：用无散度核（Divergence-Free Kernel, DFK）参数化速度场，\emph{在构造上}保证速度场严格无散度（不可压缩），而非通过惩罚项软约束。直觉上，DFK 将不可压缩性"嵌入"表示本身——无论参数如何变化，速度场始终满足 $\nabla \cdot \mathbf{v} = 0$。高斯基元作为被动示踪物被速度场驱动平流，位置由拉格朗日平流算子确定 $\mathbf{x}_i^t = \mathcal{A}_{0 \to t}(\mathbf{x}_i^0, \mathbf{v})$。为了处理长序列优化中的梯度消失问题，LagrangianSplats 引入滑动窗口策略：将长视频分割为重叠的时间窗口，在每个窗口内联合优化速度场和高斯属性，窗口之间通过连续的平流轨迹保证物理一致性。

与 Gaussian Fluids 和 LagrangianSplats 均以\emph{仿真为核心目标}不同，FluidGaussian 转向\emph{重建}视角。

\paragraph{FluidGaussian：2026.03 同期工作}

FluidGaussian（2026.03）~\cite{2603_21356}转向\emph{重建}视角，提出了一个全新的思路：将流体仿真的\emph{不确定性}引入 3DGS 重建。传统重建方法追求视觉质量，但重建结果在物理上可能不可行（如流体内部存在非零散度）。FluidGaussian 定义基于流体仿真的不确定性度量，将其与主动学习结合，优先选择同时提升视觉和物理保真度的视角。在 NeRF Synthetic、Mip-NeRF 360 和 DrivAerNet++ 上实现 PSNR 提升 8.6\%、流体仿真速度散度降低 62.3\%，开创了"功能智能 3D 重建"方向。

与前述方法均围绕\emph{流体本身}展开不同，FLUIDSPLAT 将思路迁移至广义物理场重建。

\paragraph{FLUIDSPLAT：2026.05 同期工作}

FLUIDSPLAT（2026.05）~\cite{2605_18866}将 3DGS 的空间表示思想迁移至广义物理场重建（非流体仿真）：从稀疏表面传感器重建连续物理流场（速度场、压力场）。提出用 $K$ 个各向异性高斯基元构成 partition-of-unity 支架作为空间显式中间表示，配合状态条件残差解码器提供支架无法捕获的细节。理论上证明了高斯基元数量的最优缩放律 $K^* \sim (N/\sigma^2)^{d/(2s+d)}$，在四个 2D/3D 基准上实现 11--28\% 误差降低。

上述方法从模拟器或重建的角度"正向"处理流体问题——设计求解器或从视频恢复速度场。另一个方向则更为实用：从单目视频直接提取可重仿真的 3D 流体资产。

\paragraph{从视频恢复流体动力学}

3D Dynamic Fluid Assets（2025.03）~\cite{2503_00868}修改了流体重建的表示与优化环节——用\emph{可微分物理仿真}替代纯视觉重建，使输出不仅是视觉资产，而是"仿真就绪"的物理流体资产。传统单目 3D 重建方法（如 NeRF 或标准 3DGS）仅恢复外观和几何，无法提取流体的速度场、密度、黏度等物理属性，因此重建结果不可重仿真，也无法编辑交互。该方法正是为了解决这一问题：从单目视频中提取几何、速度场和仿真参数，输出可直接导入物理引擎的流体资产。

整个流程分为三步，核心思路是\textbf{先用视觉方法获取几何和表面速度，再用物理方法优化内部体积速度和仿真参数}：

\textbf{第一步：生成式 3DGS 获取与预处理流体几何。}利用图像条件的 3DGS 生成模型（如 TriplaneGaussian）从单帧图像生成初始 3DGS，再通过单视角优化提升分辨率。优化损失函数包含五项：图像重建损失 $L_\text{img}$、各向异性损失 $L_\text{aniso}$（抑制过度拉伸）、体积损失 $L_\text{vol}$、尺度损失 $L_\text{scl}$（鼓励更细粒度的高斯基元）和亮度损失 $L_\text{lumi}$（惩罚过饱和）。优化后进行去噪和体素化填充：先剪枝低透明度或过度拉伸的高斯基元，再以网格采样填充流体内部空腔（通过射线-壳体交点奇偶性判断内部点），确保流体体积完整。直觉上，这一步将生成模型产生的稀疏"壳"变为密实的"体"。

\textbf{第二步：从光流到体积速度场。}利用光流估计 2D 屏幕速度 $V_\text{2D}$，结合 3DGS 的深度信息反投影到 3D 表面速度 $V_\text{surf}$。然而，光流在流体局部纹理均匀的区域（如平静水面）往往检测失败，输出接近零速度。为此，该方法提出\emph{主流校正}（Mainstream Correction）策略：(1) 利用几何约束（如河岸走向）确定主流方向，从邻近有效光流区域插值补偿缺失区域；(2) 基于不可压缩约束推导 2D 散度方程 $\nabla_\text{2D} \cdot V_\text{2D} = -(\frac{1}{z}\frac{\partial v_z}{\partial z} + \frac{2v_z}{z}) - v_z \frac{\partial}{\partial u} - v_z \frac{\partial}{\partial v}$，通过投影法保证校正后速度满足物理约束。该方程仅依赖深度方向速度分量 $v_z$（由深度信息确定），在局部平行于相机的水面区域 2D 散度趋于零，物理直觉是：平坦水面的 2D 投影天然接近不可压缩。

获得表面速度后，还需计算整个体积的内部速度场 $V_\text{vol}$。流体底部（如河床）速度未知，该方法利用 CFD 中的壁面函数估算边界层阻尼，采用层流外流的三阶多项式近似：$\frac{|v_\text{base}|}{V_\infty} = \frac{3}{2}\left(\frac{y}{\delta}\right) - \frac{1}{2}\left(\frac{y}{\delta}\right)^3$，其中 $V_\infty$ 为对应表面速度幅值，$\delta$ 为边界层厚度（由流体类型动态决定：河流取较大厚度，烟雾取亚网格尺度），$y$ 为到壁面的距离。该公式直觉含义是：越靠近固体壁面，流体速度越慢，呈现平滑的阻尼过渡。最后，以封闭表面速度为边界条件，通过约束投影法求解无散度 3D 速度场。

\textbf{第三步：可微分仿真参数优化。}将粒子速度场转换到欧拉网格上，利用 Affine Particle-in-Cell（APIC，一种混合拉格朗日-欧拉数值方法，通过额外计算局部仿射速度场提升精度和稳定性）建立可微分仿真管线。两个关键修改使仿真可微分：(1) 用 PeRCNN（物理编码循环卷积网络）替代传统迭代压力求解器（如 Jacobi 迭代），学习卷积核近似空间微分算子（梯度、散度、拉普拉斯），外层循环结构模拟迭代过程，在保持精度的同时实现端到端可微分；(2) 对流项用解析梯度替代不可微的回溯网格索引操作。优化目标为 11 个物理参数：入口/出口体积速度 $v_\text{in}$/$v_\text{out}$、入口速度波动 $\sigma_\text{in}$、密度 $\rho$、黏度 $\mu$、边界弹跳 $b$ 与阻尼 $d$、重力 $g$ 和时间步长 $dt$。损失函数结合速度 $L_2$ 距离和方向一致性：$L = \mathbf{1} \cdot \left[(1 - \frac{v_\text{sim}}{|v_\text{sim}|} \cdot \frac{v_\text{gt}}{|v_\text{gt}|}) + L_2(v_\text{sim}, v_\text{gt})\right]$，按网格类型（表面/占据/空）加掩码。

输出为与源视频物理动态匹配的"仿真就绪"流体资产：3DGS 几何 $+$ 优化后的物理参数。该资产支持几何编辑（修改河道形状后重新仿真）、运动时长扩展（仿真可运行比视频更长时间）、以及与外部场景的物理交互。APIC 仿真中，仿射矩阵 $C$ 近似速度梯度，通过变形梯度 $F_p^{n+1} = (I + dt \cdot C_p^{n+1}) F_p^n$ 更新高斯协方差 $A_p = F_p A_p^0 F_p^\top$，使高斯基元的形状随流体变形自然演化。

\subsection{可变形物体与多物理场}
\label{ss:06_5b}

本节介绍三条技术路线：（1）可变形物体与布料仿真；（2）声音与多物理场；（3）其他物理方法。

\subsubsection{可变形物体与布料仿真}

本路线介绍为3DGS高斯基元注入结构化物理连接关系以实现可变形物体动力学建模的方法。

可变形物体（绳子、布料、毛绒玩具等）的动力学建模是物理仿真的重要子领域。与刚体和流体不同，可变形物体既有整体运动（如布料飘动），又有局部形变（如布料褶皱），且不同部位可能呈现截然不同的材料行为（弹性、黏弹性、塑性等）。3DGS 的挑战在于：高斯基元缺乏拓扑连接关系，无法直接表达可变形物体的内部结构。

\paragraph{可变形物体动力学建模}

PhysTwin（2025.03）~\cite{2503_17973}修改了可变形物体的动力学表示环节——用弹簧-质点图结构替代标准 3DGS 的无拓扑高斯基元集合，为离散高斯点注入结构化的物理连接关系。标准 3DGS 的动态变体（如 Deformable 3DGS、Dynamic 3DGS）通过学习形变场或逐帧优化高斯位置来拟合运动，但缺少物理约束使其无法外推到未见交互场景（如改变抓取方式或施加新的外力）。PhysTwin 的设计动机正是针对这一问题：以弹簧-质点模型作为物理骨架，每个质点通过弹簧与相邻质点连接，弹簧的弹力提供恢复力、阻尼力耗散能量，从而在拓扑层面建立可变形物体的内部结构。

PhysTwin 的核心表示是一个弹簧-质点图 $G=(V,E)$，其中 $V$ 为质点集、$E$ 为弹簧集。质点 $i$ 受到相邻弹簧力和阻尼力的合力驱动：

\begin{equation}
\mathbf{F}_i = \sum_{(i,j)\in E} \left[ k_{ij}(\|\mathbf{x}_j - \mathbf{x}_i\| - l_{ij})\frac{\mathbf{x}_j - \mathbf{x}_i}{\|\mathbf{x}_j - \mathbf{x}_i\|} - \gamma(\mathbf{v}_i - \mathbf{v}_j) \right] + \mathbf{F}_i^{\text{ext}}
\end{equation}

其中 $k_{ij}$ 为弹簧刚度，$l_{ij}$ 为弹簧自然长度，$\gamma$ 为阻尼系数，$\mathbf{F}_i^{\text{ext}}$ 包含重力、碰撞等外力。系统通过显式欧拉积分更新状态 $\mathbf{X}_{t+1} = f_{\boldsymbol{\theta}, G_0}(\mathbf{X}_t, \mathbf{a}_t)$，其中 $\boldsymbol{\theta}$ 编码所有物理参数（刚度、阻尼、碰撞参数），$G_0$ 为初始几何与拓扑，$\mathbf{a}_t$ 为控制点动作。直观而言，弹簧力将质点拉回平衡位置（类比橡皮筋），阻尼力防止永久振荡（类比减震器），二者共同约束物体的弹性变形行为。

PhysTwin 的核心流程分两个阶段。第一阶段聚焦物理与几何优化：先用生成式形状先验（TRELLIS）从单张 RGB 图像恢复完整几何——因为稀疏视角下深度观测仅覆盖部分表面，直接从点云恢复完整拓扑不可行——获得初始网格后，通过配准模块（基于 SuperGlue 的 2D 匹配估计旋转、PnP 精修位姿、as-rigid-as-possible 非刚性形变）对齐到真实观测。随后以稀疏到密集的分层优化策略联合优化拓扑和物理参数：先以零阶采样优化（CMA-ES）估计不可微的拓扑参数（连接半径、最大邻居数）和均匀刚度，再用一阶梯度下降通过自建可微弹簧-质点仿真器精修稠密刚度（单个弹簧独立刚度，参数空间达数万维）。

第二阶段进行外观优化：以线性混合蒙皮（LBS，即根据邻近质点的运动加权插值高斯基元的位置和形状）将高斯绑定到质点网格，使高斯随质点运动自动变形。外观通过 $L_1$ 损失加 D-SSIM 项在首帧优化，高斯被约束为各向同性以避免形变时的尖刺伪影。

PhysTwin 的关键设计是引入控制点——通过弹簧将控制点连接到物体点（参数为连接半径和最大邻居数），使系统支持交互式仿真和基于模型的机器人运动规划。在 22 个真实可变形物体场景（绳索、毛绒玩具、布料、快递盒）上的实验表明，PhysTwin 在重建重仿真、未来预测和未见交互泛化三个任务上均优于 Spring-Gaus 和 GS-Dynamics 基线，且支持实时交互仿真（基于 Warp 加速）。

与 PhysTwin 以弹簧-质点图作为显式物理骨架不同，NeuMA 转而从视觉观察中学习物体动力学。

\paragraph{NeuMA：2024.10 同期工作}

NeuMA（2024.10）~\cite{2410_08257}转而从视觉观察中学习物体内在动力学，提出了一种优雅的分解策略：将材质模型分解为物理先验和可学习修正项 $\sigma = \sigma_{\text{prior}} + \sigma_{\text{correct}}$，其中 $\sigma_{\text{prior}}$ 来自经典本构模型（如 Neo-Hookean 弹性），$\sigma_{\text{correct}}$ 由神经网络预测。PhysTwin 的弹簧参数通过零阶采样和一阶梯度联合优化确定，而 NeuMA 将物理知识嵌入为结构化先验（而非可优化参数），使学习过程聚焦于修正项。这一分解兼顾了物理先验的泛化性（对已知材料行为建模准确）和学习修正的灵活性（捕获偏离理想模型的真实行为）。配合 Particle-GS 可微分渲染器实现端到端优化。

与 NeuMA 将物理先验嵌入为结构化约束的范式不同，3DGSim 提出完全端到端可微的框架。

\paragraph{3DGSim：2025.03 同期工作}

3DGSim（2025.03）~\cite{2503_24009}进一步放弃了显式物理模型的依赖，提出首个完全端到端可微的 3D 物理模拟框架，从多视角 RGB 视频联合训练逆渲染和动力学预测，无需深度信息或粒子追踪等特权信息。PhysTwin 和 NeuMA 均保留了显式物理结构（弹簧-质点图或本构模型），而 3DGSim 用纯神经力场替代物理方程：学习全局变换力场 $F^{\text{global}}$（捕获物体间接触、滑动和重力）和局部应力场 $F^{\text{local}}$（捕获软体变形），通过二阶积分器恢复物理轨迹。该框架能统一建模从刚性到弹性、布料等多种物理行为。

与 3DGSim 用纯神经力场建模动力学不同，Particle-Grid Neural Dynamics 引入混合粒子-网格表示。

\paragraph{Particle-Grid Neural Dynamics：2025.06 同期工作}

Particle-Grid Neural Dynamics（2025.06）~\cite{2506_15680}回归混合表示的思想，提出拉格朗日-欧拉混合表示学习可变形物体动力学：粒子表示携带物体状态（拉格朗日视角，跟踪个体运动），均匀网格注入空间连续性先验（欧拉视角，提供全局空间信息）。其关键组件包括全局点云编码器（PointNet）捕获整体形状信息以克服 GNN 的连通性限制，网格速度编辑（GVE）机制建模物体-地面碰撞和机器人夹持，以及 B 样条插值实现网格到粒子的速度传递。结合 3DGS 的 LBS 绑定实现可变形物体的数字孪生和动作条件视频生成。

上述方法均依赖多视角观测。

\paragraph{MonoPhysics：2026.05 同期工作}

MonoPhysics（2026.05）~\cite{2605_30320}将可变形物体动力学建模推进到更具挑战性的单目设置，实现了首个单目逆物理估计框架，通过三个创新的视觉-物理桥接机制将可微分 MPM 模拟与 3DGS 耦合：(1) 尺度重参数化——在相机空间维护粒子并通过单一可学习标量 $s$ 控制全局尺度，使图像损失通过可微分仿真聚合一致的尺度信号；(2) 物理感知几何精修——使每个粒子的物理体积 $V_i$ 成为可优化参数并通过分区-of-unity 约束自适应更新，配合物理重要性驱动的粒子管理；(3) 可微分位置图——通过重参数化 $\hat{x}_i = x_i^{2D} + (\Sigma_i^{2D})^{1/2} z$ 使像素位置成为高斯参数的可微函数，通过 Sinkhorn 散度的最优传输提供全局方向梯度。在单目设置下达到多目方法的性能水平。

除上述代表性方法外，可变形物体动力学建模还在运动表示设计、物理基元构造和布料动力学三个方向上持续拓展。

在\emph{运动表示}方向，核心挑战是如何在保持计算效率的同时精确捕捉复杂物体的全局和局部运动。

\paragraph{Dynamic 3D Gaussian Tracking：2024.10 同期工作}

Dynamic 3D Gaussian Tracking（2024.10）~\cite{2410_18912}结合动态 3DGS 重建与图神经网络（GNN）从真实多视角视频中学习可变形物体动力学，其关键设计是将密集 3D 高斯重建结果下采样为稀疏控制粒子（control particles），在控制粒子上运行 GNN 动力学模型预测运动，再通过插值将运动传播回密集高斯——这种"稀疏驱动、稠密插值"的策略使动力学模型既能学习到全局的物体间交互模式，又保持了可部署于机器人操控规划的实时推理能力。

\paragraph{HEIR：2025.10 同期工作}

HEIR（2025.10）~\cite{2510_26786}则从运动分解的视角出发，提出一种通用的层次运动建模方法：通过可微图学习自动发现运动的层次结构，将每个运动单元视为图的顶点，有向边编码学习到的父子依赖关系，从而将全局绝对运动显式分解为父级继承模式（parent-inherited pattern，即子节点从父节点继承的运动分量）和局部运动残差（local motion residual，即子节点独有的精细运动）。这种数据驱动的层次分解无需手工定义运动基元，可在不同任务间泛化，且分解出的层次结构本身具有可解释性。

在\emph{物理基元}方向，PhysConvex 跳出高斯/粒子的传统表示范式。

\paragraph{PhysConvex：2026.02 同期工作}

PhysConvex（2026.02）~\cite{2602_18886}跳出高斯/粒子的传统表示范式，提出用物理驱动的 3D 动态凸面基元（dynamic convex primitives）统一视觉渲染和物理模拟。与标准 3DGS 用无物理属性的椭球体表示场景不同，PhysConvex 的凸面基元由连续力学（continuum mechanics）支配，通过边界驱动动态凸面表示（boundary-driven dynamic convex representation）建模变形：顶点和面的动力学演化捕捉空间自适应、非均匀的变形行为以及动态变化的边界。为进一步提升大规模复杂几何和异质材料的仿真效率，PhysConvex 开发了降阶凸面模拟（reduced-order convex simulation）策略——利用神经蒙皮特征模态（neural skinning eigenmodes）作为形状和材料感知的变形基底，在时间变化的参数空间中驱动凸面场的对流（advection），在保持物理精度的同时显著降低计算维度。

在\emph{布料动力学}方向，CloDS 聚焦布料这一特定可变形物体类别。

\paragraph{CloDS：2026.02 同期工作}

CloDS（2026.02）~\cite{2602_01844}聚焦布料这一具有大非线性变形和严重自遮挡的特定可变形物体类别，提出在未知物理条件下从纯视觉观测无监督学习布料动力学。现有方法依赖已知物理属性作为监督或输入，而 CloDS 采用三阶段管线：首先通过视频到几何的接地（video-to-geometry grounding）从多视角视频恢复布料时变几何，然后在地化的网格上训练动力学模型。为应对接地过程中布料的大幅非线性变形和自遮挡问题，CloDS 引入双位置不透明度调制（dual-position opacity modulation）机制——同时考虑绝对位置和相对位置的网格高斯泼溅，实现 2D 观测与 3D 几何之间的双向映射，从而在无需任何物理标注的条件下学习到物理一致的布料运动模式。

\paragraph{布料仿真专用方法}

布料是一类特殊的可变形物体，具有薄壳结构和大变形特征，其状态估计面临无穷维状态空间和严重自遮挡的双重挑战。

\textbf{Baseline 做法与问题。}现有布料状态估计主要依赖两类方法：(1) 基于深度传感器的 3D 表示（如点云、图结构），虽能捕捉几何但缺乏纹理等视觉线索，且仿真到真实的迁移需要耗时的在线微调；(2) 基于 3DGS 的动态跟踪方法（如 MD-Splatting、DeformGS）需要至少 50 个相机的密集观测和逐场景优化，计算开销极大，难以满足实际部署需求。

Cloth-Splatting（2025.01）~\cite{2501_01715}修改了布料状态估计的观测模型环节——用\emph{mesh 约束的 Gaussian Splatting}替代传统的深度传感器或密集多目 3DGS，仅需 3--4 个 RGB 相机即可从图像空间反向传播梯度来修正布料 3D 状态。

\textbf{实现原理。}该方法将布料状态估计建模为贝叶斯滤波问题——类似卡尔曼滤波的``预测-更新''两步循环，其中``预测''由学习到的动力学模型提供先验，``更新''由可微渲染器利用观测修正先验。

预测步骤使用基于图神经网络（GNN）的动作条件动力学模型 $f_\theta$：给定过去 $m$ 帧的布料网格状态 $M_{t-m:t}$ 和机器人动作 $a_t$，预测下一帧状态 $\hat{M}_{t+1} = f_\theta(M_{t-m:t}, a_t)$。网络输出每个顶点的加速度，经前向欧拉积分得到下一帧的速度和位置。

更新步骤的关键创新在于\emph{mesh 约束的 Gaussian Splatting}（$h_{\text{GS}}$）：将每个 3D 高斯绑定到布料 mesh 的一个三角面上，其位置用重心坐标参数化：

\begin{equation}
  \mu = b_1 v^1 + b_2 v^2 + b_3 v^3, \quad b_1 + b_2 + b_3 = 1
\end{equation}

其中 $v^1, v^2, v^3$ 为所属面的三个顶点。直觉上，高斯``粘贴''在 mesh 表面，随布料变形而自动移动——mesh 变形时高斯的重心坐标和外观参数保持不变，仅顶点位置变化驱动高斯重新定位。这使得布料状态空间到图像空间形成完全可微的映射。

为避免直接优化 GNN 导致梯度消失，更新后的状态定义为预测状态加一个学习到的残差偏移：

\begin{equation}
  \tilde{M}_{t+1} = \hat{M}_{t+1} + \delta_{t+1}, \quad \delta_{t+1} = u_\phi(t+1)
\end{equation}

其中残差网络 $u_\phi$ 是一个以归一化时间为输入的 3 层 MLP（宽度 256），输出初始值接近零。通过最小化观测损失 $L_{\text{obs}} = \|Y_{t+1} - h_{\text{GS}}(\tilde{M}_{t+1}, P)\|_2^2$ 优化残差，即让渲染的布料图像与真实 RGB 观测尽可能一致。正则化损失 $L_{\text{reg}} = L_{\text{SSIM}} + L_{\text{iso}} + L_{\text{magn}}$ 分别约束结构相似性、邻域顶点等距（保持布料不可拉伸性）和最小运动（防止数值不稳定）。

\textbf{效果。}相比纯物理基线（仅用 GNN 预测），GS 更新将跟踪误差降低 57\%；相比次优基线 MD-Splatting，收敛速度提升 85\%，且仅需约 4k 个高斯即可重建布料外观。该方法进一步支持闭环折叠操控：通过模型预测控制（MPC）在每次重规划时用 GS 修正布料状态估计，折叠精度接近有真值状态的 Oracle 方法。

SketchRodGS（2026.01）~\cite{2601_02072}从用户交互的角度解决 3DGS 场景中细长几何的提取与动画化问题。其核心思路十分直觉：用户在屏幕上画一条线（草图），系统自动找到对应的 3D 细长结构并将其转化为可仿真的折线网格（polyline mesh）。由于 3DGS 缺乏拓扑连接信息且高斯基元的配置包含大量噪声，直接从点云提取细长结构极为困难。SketchRodGS 的关键创新在于屏幕空间最短路径分析：将用户草图投影到多视角的深度图上，在屏幕空间中通过动态规划高效求解最短路径，从而在不同视角间鲁棒地重建出绳索、布料边缘等细长部件的折线表示。提取的折线可直接接入弹性杆（elastic rod）物理仿真器，使 3DGS 场景中的细长物体能够进行物理驱动的动画——绳索可以摆动、打结，布料边缘可以随风吹动。该方法在多个自然场景（in-the-wild）示例上验证了有效性。

前述方法主要处理单时间步或短序列的物理行为。将物理仿真扩展到完整的时空 4D 场景，则需要解决材料参数的自动推断和跨帧运动一致性两个更根本的问题。

\paragraph{物理驱动的 4D 动态仿真}

将物理仿真扩展到 4D（3D + 时间）需要解决参数推断和运动监督两个核心问题。

PhysFlow（2024.11）~\cite{2411_14423}修改了材料参数优化的监督信号环节——用光流损失（optical flow loss）替代传统的渲染损失（render loss）和分数蒸馏采样损失（SDS loss），实现从任意输入（单图/多视角/视频）自动估计材料属性并进行 4D 物理仿真。

\textbf{Baseline 做法与问题。}在 PhysFlow 之前，基于视频扩散的材料参数优化主要有两条路线。PhysDreamer 使用渲染损失 $L_{\text{render}} = \sum_t \lambda L_1(\hat{I}_t, I_t) + (1-\lambda)L_{\text{D-SSIM}}(\hat{I}_t, I_t)$，直接比较仿真渲染帧与扩散生成帧的像素差异——但由于视频扩散输出的不稳定性与高噪声，该损失仅能处理微小运动，稍大的运动即导致优化发散。另一条路线（如 Physics3D）采用 SDS 损失 $L_{\text{sds}}$，从预训练扩散模型中蒸馏梯度，但 SDS 需要对整个扩散模型做反向传播，显存开销极大，且在结合复杂文本提示的视频扩散模型时更为严重。此外，上述方法均局限于弹性材料，无法处理流体、沙子等复杂材料。

\textbf{光流引导的参数优化。}PhysFlow 的核心洞察是：优化材料参数并不需要逐像素对比图像，而只需对比\emph{运动场}——光流恰好编码了像素级位移信息，是材料参数对运动的"直接传感器"。具体流程为：首先利用 GPT-4 从输入图像推断材料类型 $M$、密度 $\rho$ 和初始参数 $M_p^0$；然后根据输入类型（多视角/视频/单图）分别选择 PGSR、Grid4D 或 Splatt3R 构建 3D 高斯；接着通过视频扩散模型（根据输入类型可选 SVD、CogVideoX、MotionCtrl 等）生成参考视频，并用 RAFT 计算相邻帧光流 $U(I_t, I_{t+1})$；最终在可微分 MPM（Material Point Method，一种混合拉格朗日-欧拉的粒子仿真方法，将材料离散为粒子追踪变形，在背景网格上求解动量守恒）框架下，最小化光流损失：

\begin{equation}
    L_{\text{flow}} = \sum_t \|U(I_t, I_{t+1}) - \hat{U}(\hat{I}_t, \hat{I}_{t+1})\|^2
\end{equation}

其中 $\hat{U}$ 是仿真渲染帧之间的光流。直观地说，该损失迫使仿真运动与参考视频中的运动对齐，而非逐像素对比外观——这避免了图像级噪声的干扰，同时使梯度信号更直接地反映材料参数对运动的影响。

\textbf{材料类型覆盖。}得益于 MPM 的通用性，PhysFlow 支持 7 种材料类型（弹性、橡皮泥、金属、泡沫、沙子、牛顿流体、非牛顿流体），每种材料由不同参数组合定义：弹性体由杨氏模量 $E$ 和泊松比 $\nu$ 控制；橡皮泥和金属额外引入屈服应力 $\sigma_Y$ 以建模塑性变形；沙子由摩擦角 $\phi_{\text{fric}}$ 表征；流体则由粘度 $\eta$ 和体积模量 $\kappa$ 控制。优化后的参数集成回可微分 MPM 框架，驱动 3D 高斯形变并从新视角渲染，完成 4D 重建。在合成数据集上，PhysFlow 在多数材料类型的相对误差（RE）和端点误差（EPE）上优于 PAC-NeRF，与 GIC 表现相当；在真实数据集上，32 人的人类评估中物理真实感和视觉真实感评分均超过 3 分（5 分制），显著高于所有基线方法。

与 PhysFlow 侧重于从视觉观测优化材料参数不同，NGFF 从力场预测的角度重新审视问题。

\paragraph{NGFF：2026.02 同期工作}

NGFF（2026.02）~\cite{2602_00148}（Neural Gaussian Force Field）从力场预测的角度重新审视 3DGS 物理仿真——不依赖 MPM 等传统数值求解器推断材料属性，而是直接学习一个端到端的神经力场。其力场预测器学习全局变换力场和局部应力场，通过神经算子在关系图上建模物体间物理接触，速度比先前高斯仿真器快两个数量级。这一设计的动机在于：PhysFlow 的可微分 MPM 框架虽然物理可解释性强，但计算开销受限于粒子数量与网格分辨率，且材料类型需要预先指定；NGFF 用数据驱动的神经算子替代数值仿真，在牺牲部分物理解释性的前提下换取了数量级的加速和更强的泛化能力。同时提出 GSCollision 数据集，包含超过 64 万渲染物理视频（约 4TB），涵盖多种材料和复杂交互场景。

在 PhysFlow 和 NGFF 分别从材料参数优化和力场预测推进 4D 物理仿真的基础上，后续工作进一步拓展了训练数据来源和感知模态。与 PhysFlow 依赖合成参考视频和光流监督不同，该方法直接从真实世界视频学习。

\paragraph{Learning from Real-world Videos：2026.05 同期工作}

Learning from Real-world Videos（2026.05）~\cite{2605_23845}首次实现从真实世界无标注视频通过渲染监督学习多物体碰撞动力学，直接在密集高斯基元（无启发式子采样）上操作，并贡献约 500 视频的真实数据集——这解决了 PhysFlow 框架中视频扩散模型生成的参考运动可能偏离真实物理分布的问题。

\paragraph{PEGS：2025.11 同期工作}

PEGS（2025.11）~\cite{2511_17116}则从感知模态角度出发：PhysFlow 的光流监督仅利用标准 RGB 帧，无法捕获高速运动中的瞬间状态；PEGS 融合物理约束、事件相机和 3DGS，实现大时空尺度刚体运动的高质量去模糊重建，贡献首个针对自然快速刚体运动的 RGB-Event 配对数据集。

\subsubsection{声音与多物理场}

本路线介绍将3DGS从纯视觉表示扩展到声场、热场等多模态物理场的方法。

物理仿真不仅限于视觉——声音、热量、电磁等多物理场（multi-physics）现象同样可以用高斯基元表示和模拟。本节的方法将 3DGS 从纯视觉表示扩展到多模态物理场。

\paragraph{声音场重建与合成}

AudioGS（2026.04）~\cite{2604_08967}将 3DGS 从视觉域迁移到声学域——用显式"音频高斯"替代隐式神经声场表示，实现不依赖视觉先验的声场重建与双耳合成。在此之前，新视角声学合成（NVAS）的主流做法依赖视觉引导：ViGAS 用跨模态 Transformer 从视频帧提取空间线索，AV-NeRF 用 NeRF 联合渲染 RGB/深度图来指导声学合成，AV-Cloud 则用 SfM 稀疏点云作为声学锚点。然而，这些方法存在两个问题：（1）依赖视觉先验，在纹理缺失、相机标定不准或视觉数据稀疏时性能显著退化；（2）仅重用源视角相位而忽略相位建模，导致耳间时间差（ITD）线索不准，空间感失真。因此 AudioGS 选择完全抛弃视觉依赖，直接从稀疏声学观测显式重建 3D 声场。

直觉上，AudioGS 将声场编码为一组"音频高斯"（Audio Gaussians），每个音频高斯对应 STFT 频谱图上的一个时频 bin（time-frequency bin），配备可学习的 3D 位置、双球谐系数和距离衰减系数。这一设计利用了语音信号在时频域的近似 W-disjoint 正交性——每个时频 bin 的能量通常由单一声源主导——从而可以将每个 bin 独立映射为一个空间基元。双 SH 参数化将方向性能量与空间双耳线索解耦：单耳 SH 系数 $c^{\text{mono}}_{f,t}$ 产生非负的共享能量模式，差分 SH 系数 $c^{\text{diff}}_{f,t}$ 编码左右耳差分线索。渲染时，对目标监听位置 $p$ 和朝向 $R_{\text{cam}}$，AudioGS 依次执行三步操作：（1）方向性调制——通过球谐基函数 $Y(d_{f,t})$ 将 SH 系数投影到从音频高斯指向监听者的方向 $d_{f,t}$ 上，得到单耳掩码 $M(f,t) = 2\,\text{sigmoid}(c^{\text{mono}}_{f,t} \cdot Y(d_{f,t}))$ 和差分掩码 $D(f,t)$；（2）距离衰减——建模频谱幅度随距离的物理衰减：

\begin{equation}
G(f,t) = \left( \frac{\|p_{\text{ref}} - x_{f,t}\|_2 + \epsilon}{\|p - x_{f,t}\|_2 + \epsilon} \right)^{\!\alpha_{f,t}}
\end{equation}

其中 $\alpha_{f,t} > 0$ 是可学习的衰减指数，$\epsilon$ 为数值稳定项，该公式允许模型自适应地从混响场（$\alpha \approx 0$）过渡到直达声衰减（$\alpha \gtrsim 1$）；（3）相位校正——基于刚球模型计算物理 ITD 幅值，并通过有界残差因子进行可学习校正，弥补传统方法直接复用源相位的缺陷。最终将幅度与校正相位组合，经逆 STFT 输出双耳波形。

实验在 Replay-NVAS 数据集上进行，音频重采样至 16\,kHz，STFT 使用 512 点 FFT、400 窗长、160 帧移和 Hamming 窗。训练 60 个 epoch，所有 SH 系数初始化为 0，衰减系数初始化为 1。与视觉引导方法相比，AudioGS 在不使用任何视觉数据的前提下，MAG 误差降低 14\%+，ENV 降低约 6\%，感知质量指标 DPAM 降低约 25\%，耳间能量比误差（LRE）也达到竞争性水平。消融实验证实三个组件缺一不可：移除距离衰减导致空间精度退化最严重，移除 SH 导致重建质量显著下降，移除相位校正则使 LRE 指标明显劣化。MUSHRA 听觉测试（12 名受试者）也表明 AudioGS 的合成音频最接近真实录音。

与 AudioGS 聚焦于声场的空间重建与双耳合成不同，SonicGauss（2025.07）~\cite{2507_19835}关注的是碰撞声合成——不是重建已有声音的空间分布，而是从 3DGS 几何表示主动生成物理碰撞声。AudioGS 的音频高斯编码的是环境中已存在的声场信息，而 SonicGauss 需要从高斯椭球体的几何特征推断材质特性（金属、木质、玻璃等）和空间声学相关性，再通过扩散声音合成模型生成对应的声音。这一差异的设计动机在于：AudioGS 解决的是"听者在不同位置听到什么"的空间声学重建问题，而实际物理仿真场景中更基础的需求是"物体碰撞时发出什么声音"——只有先生成碰撞声，才能进一步讨论其空间传播。SonicGauss 整合 PointTransformer 特征提取器与扩散声音合成模型，支持空间变化的声音响应——同一物体在不同位置被敲击产生不同的回声效果。

声音之外，3DGS 的显式点云表示和可微分渲染管线同样适用于火焰、降雨、风等其他物理现象的建模与仿真。

\paragraph{火焰、雨等多物理场}

FlameGS（2024.12）~\cite{2412_19841}修改了火焰化学发光断层成像（CTC）的三维表示环节——用 3D 高斯基元替代传统体素网格，实现火焰光场的高效重建。传统方法采用 ART（代数重建技术）：将火焰空间离散为体素网格（如 $120^3 = 1{,}728{,}000$ 个体素），通过 2D 投影方程逐体素迭代更新发光强度。然而体素表示存在两个瓶颈：网格分辨率固定，重建区域受预定义边界严格限制；体素数量庞大导致计算和内存消耗极高。FlameGS 的核心思路是：火焰本身是连续的发光体，可以自然地用一组可微的 3D 高斯函数来近似——每个高斯由位置 $\boldsymbol{\mu}$、各向异性协方差 $\boldsymbol{\Sigma} = \mathbf{R}\mathbf{S}\mathbf{S}^\top\mathbf{R}^\top$（$\mathbf{S}$ 为缩放向量，$\mathbf{R}$ 为旋转四元数）和发光强度 $\sigma$ 三个参数描述，整体仅需约 13{,}000 个可训练参数。直觉上，每个高斯对应火焰中一个"发光团"，其位置、形状和亮度共同刻画火焰的三维光场分布。训练时利用 10 台同步相机（$800 \times 1024$ 分辨率）捕获的 2D 投影图像作为监督，最小化 L1 损失与 D-SSIM 损失的加权和来联合优化高斯参数与相机位姿；初始化阶段则通过光线追踪从多视角反向投影，只在真实火焰区域生成稀疏初始点云，有效缓解稀疏视角下的过拟合。相比传统 ART 算法，FlameGS 计算时间减少约 34 倍，内存减少约 10 倍，同时达到 0.96 的 SSIM 和 39.05 dB 的 PSNR。

与 FlameGS 聚焦于\emph{火焰光场的离线重建}不同，FieryGS 从\emph{在线燃烧模拟}角度切入。

\paragraph{FieryGS：2026.05 同期工作}

FieryGS（2026.05）~\cite{2605_00177}从\emph{在线燃烧模拟}角度出发，修改了 3DGS 场景编辑中的物理属性推理与特效生成方式——用 MLLM 驱动的零样本材料推理替代手工标注的可燃性参数，用 Taichi 体积燃烧模拟引擎替代预设视觉特效资产，再通过统一体积渲染器将模拟结果与 3DGS 场景融合。具体而言，框架先将重建的 3DGS 场景分割为语义一致区域，每个区域渲染为 2D 图像后送入 MLLM 推理材料类型、可燃性、热扩散系数和烟雾颜色等燃烧属性；随后 Taichi 引擎基于这些属性执行物理燃烧模拟，处理火焰沿可燃表面传播、烟雾体积扩散和表面碳化等复杂现象；最终统一体积渲染器将燃烧体素与 3DGS 高斯基元在同一管线中渲染，确保火光照明与场景几何的一致性。直觉上，可以理解为三个模块各司其职：MLLM 充当"材料专家"，自动判断场景中每个物体能否燃烧、如何燃烧；物理引擎充当"燃烧模拟器"，负责火焰传播和烟雾扩散的物理计算；渲染器充当"画面合成器"，负责将模拟结果无缝融合到 3DGS 场景中。这是将物理准确燃烧模拟集成到 3DGS 的首个工作，支持火焰强度、气流和点火位置等参数的用户交互控制。

与 FlameGS 和 FieryGS 均围绕\emph{火焰}展开不同，RainyGS 转向\emph{降雨}场景。

\paragraph{RainyGS：2025.03 同期工作}

RainyGS（2025.03）~\cite{2503_21442}转向\emph{降雨}这一多相流动场景，修改了传统雨效果合成中的场景交互建模方式——用基于物理的降雨粒子系统与浅水模拟替代简单的 2D 雨滴叠加或视频风格迁移，将其集成到 3DGS 的实时渲染框架中。降雨仿真的物理原理分为两层：空中部分将雨滴建模为球状 3D 高斯基元，沿重力方向排列形成雨丝，通过粒子系统控制雨滴的生成、运动和消亡；地面部分则在二维交错 MAC 网格上执行浅水模拟，处理雨滴落地后的飞溅、积水和涟漪传播，并通过屏幕空间反射（SSR）技术实时计算湿润表面的镜面反射效果。该方法以超过 30fps 的速度生成真实感强、物理准确的动态雨效果，支持从细雨到暴雨的强度灵活控制，在 Mip-NeRF 360 和 Tanks and Temples 等真实场景及大规模驾驶场景上均表现有效。

与前述方法均针对\emph{单一物理现象}不同，DiffWind 聚焦于\emph{风-物体交互}这一跨模态耦合问题。

\paragraph{DiffWind：2026.03 同期工作}

DiffWind（2026.03）~\cite{2603_09668}聚焦于\emph{风-物体交互}这一跨模态耦合问题，修改了动态场景重建中的运动建模方式——用可微分物理模拟驱动的风场-物体联合优化替代纯数据驱动的变形场学习，提出首个物理信息可微分风-物体交互建模框架。其核心创新在于为风和物体分别选择最适合的表示并进行耦合交互：风被建模为规则的网格物理场，每个网格节点存储密度、速度和力等物理量，用 Lattice Boltzmann Method（LBM）模拟其演化以符合流体力学基本定律；物体则继承 3DGS 的粒子表示，每个高斯基元携带外观、材质和运动属性；两者之间的交互通过 Material Point Method（MPM）实现——MPM 将粒子属性映射到背景网格上计算受力，再将更新后的物理量传回粒子，自然桥接了连续场（风）与离散粒子（物体）两种表示。在重建阶段，通过可微分模拟与可微分渲染的端到端反向传播，联合优化时空风场和物体运动参数；在 LBM 物理约束的保障下，优化得到的风场不仅匹配观测运动，还满足流体守恒律。该框架支持前向模拟（在新风条件下预测物体运动）和风力重定向（将重建的风场效果迁移到新场景），并构建了 WD-Objects 风驱动场景数据集用于系统评测。

\paragraph{物理场可视化}

GS-Surrogate（2026.04）~\cite{2604_06358}修改了集合仿真（ensemble simulation）参数空间探索的空间表示与变化建模方式——用\textbf{可变形显式高斯场}替代传统的\textbf{端到端 2D 图像代理}或\textbf{单统一隐式 NeRF 场}，首次实现同时调整仿真参数和可视化设置的实时交互探索。

\textbf{Baseline 做法与问题。}集合仿真是科学计算中的核心需求：研究者对同一物理系统在多种参数设置下运行仿真（如宇宙学中调整物质密度 $\Omega_m$、重子密度 $\Omega_b$ 和哈勃常数 $h$），以理解参数对结果的影响。传统做法有两类：(1) 保存高分辨率仿真输出需要海量存储和 I/O，in-situ 渲染虽减少存储但固定渲染配置，牺牲后分析灵活性；(2) 可视化代理模型（如 InSituNet、ViSNeRF）直接从参数合成 2D 图像或用 NeRF 建模场景，但前者缺乏几何一致性（无法在不同视角保持结构一致），后者因密集射线采样导致渲染速度慢，且将所有参数驱动的变化编码在单一隐式场中——仿真参数变化与可视化设置变化相互纠缠，学习传递函数的变化可能干扰仿真参数的几何差异。

\textbf{核心思路：规范场 + 两阶段变形。}GS-Surrogate 的核心洞察是：不同仿真成员共享大量几何结构，差异可由高斯基元的位置偏移和属性变化捕捉。因此，先用一组共享的\textbf{规范高斯场}（canonical Gaussian field）$\mathcal{G}$ 捕获所有成员的基础几何，再通过条件变形网络将其适配到特定参数组合。关键是将变形分解为两个\textbf{顺序阶段}，将仿真参数驱动的变化与可视化参数驱动的变化显式分离：

\begin{equation}
\mathcal{G}' = F_{\text{sim}}(\mathcal{G}, P_{\text{sim}}), \quad \mathcal{G}'' = F_{\text{vis}}(\mathcal{G}', P_{\text{vis}}),
\end{equation}

其中 $F_{\text{sim}}$ 以每个规范高斯的位置 $\mu_i$ 和仿真参数向量 $P_{\text{sim}}$ 为输入，预测五类属性偏移量：位置偏移 $\Delta\mu$ 捕捉结构位置变化（如风应力增大使洋流结构扩张），缩放偏移 $\Delta s$ 和旋转偏移 $\Delta q$ 捕捉局部几何变化，不透明度偏移 $\Delta o$ 建模拓扑变化（如宇宙学仿真中某些纤维结构在特定参数下消失），颜色偏移 $\Delta c$ 捕捉标量场值变化导致的外观差异。$F_{\text{vis}}$ 在仿真变形结果上进一步调整以适配可视化任务（如传递函数编辑或等值面提取）。这种模块化设计使同一套规范高斯可通过不同的可视化变形复用于多种分析任务，无需重新训练。

\textbf{设计动机。}这一分解源于两个考量：(1) 规范场与可视化参数解耦，使得新增可视化设置时无需从头重训；(2) 显式高斯基元可由可微分光栅化实时渲染，避免了 NeRF 的逐射线查询开销。在 Nyx 宇宙学（暗物质密度场，$512^3$ 体素）、MPAS-Ocean 海洋环流、XCompact3D 湍流和 CloverLeaf3D 冲击波四个数据集上，GS-Surrogate 在体渲染和等值面渲染两种任务中均实现了实时参数空间探索，渲染质量优于 ViSNeRF（PSNR 提升 1--3 dB），推理速度达到交互级别。

\subsubsection{其他物理方法}

本路线介绍从哈密顿力学、能量守恒等不同角度推进3DGS物理仿真边界的方法。

除了上述按材料类型划分的方法，还有一些工作从不同角度推进 3DGS 物理仿真的边界。

\paragraph{哈密顿力学与能量守恒}

NeHaD（2025.12）~\cite{2512_10424}（Neural Hamiltonian Deformation）修改了动态高斯变形场的解码环节——用哈密顿神经网络（Hamiltonian Neural Network, HNN）替代标准 MLP 解码器，使变形场内嵌物理约束。

\textbf{动机。}标准 4DGS 的做法是：hex-plane 编码器提取时空特征后，由纯数据驱动的 MLP 解码器直接预测高斯属性偏移量（位置、缩放、旋转）。这种 MLP 缺乏物理归纳偏置，常产生物理不合理的运动——高斯突变式出现/消失、轨迹不连续、违反能量守恒。NeHaD 的核心思路是：高斯基元天然存在于相空间（位置-动量空间）中，其演化应遵循哈密顿力学，因此用 HNN 替代 MLP 以注入能量守恒的结构性约束。

\textbf{哈密顿变形场。}直觉上，哈密顿力学要求运动的速度场可分解为两种成分：保守场（能量守恒，沿梯度方向推动运动）和螺线场（体积守恒，产生旋转运动）。NeHaD 不直接学习向量场（高维易模式坍缩且难保证能量约束），而是学习两个标量势函数 $F_1$ 和 $F_2$，通过自动微分生成具有物理保证的向量场：

\begin{equation}
\mathbf{v}_c = \nabla_{h_i} F_1(h_i), \quad \mathbf{v}_s = \nabla_{h_i} F_2(h_i) \, M^\top,
\end{equation}

即保守场取标量势的梯度，螺线场取标量势的梯度再乘以辛置换张量 $M$（辛结构是哈密顿力学的核心几何性质，保证相空间体积不变性）。$\nabla \times \mathbf{v}_c = 0$ 保证保守性，$\nabla \cdot \mathbf{v}_s = 0$ 保证螺线（散度自由）性。变形位移由轻量属性特定适配器（线性层）处理：$\Delta\mu_i = A_\mu \mathbf{v}$, $\Delta s_i = A_s \mathbf{v}$, $\Delta r_i = A_r \mathbf{v}$。

\textbf{时间积分。}为保证时间演化的数值稳定性，NeHaD 采用二阶辛积分器（Position Verlet）更新高斯位置：

\begin{equation}
\mu_i = \mu_i + \Delta t \, v_i + \frac{(\Delta t)^2}{2} F_i,
\end{equation}

相比欧拉法，辛积分器在长时间演化中保持能量守恒，避免数值漂移。

\textbf{玻尔兹曼平衡分解。}此外，NeHaD 引入玻尔兹曼平衡分解（Boltzmann Equilibrium Decomposition, BED）自适应分离静态和动态高斯：仅远离平衡态的高斯参与变形，基于 Boltzmann 分布构建软掩码
\begin{equation}
M_{\text{pos}}^{(i)} = (1-\epsilon) \exp(-\beta E_{\text{st}}^{(i)}) + \epsilon,
\end{equation}

即能量偏差 $E_{\text{st}}^{(i)}$ 越大（越偏离平衡态）的高斯越可能被激活，$\epsilon$ 保证最小激活率防止梯度消失，使计算资源集中在真正运动的高斯上。

前述方法聚焦于动力学建模——让高斯"动起来"。另一个方向则关注结构推断——从视觉观测中"看透"物体内部。

\paragraph{拓扑推断与结构发现}

TopoGaussian（2025.03）~\cite{2503_12343}修改了不透明物体内部拓扑推断的表示与仿真管线——用\textbf{全粒子化可微仿真器}替代传统的\textbf{网格依赖管线}，首次从易获取的照片和视频推断不透明物体的内部拓扑结构，无需 CT/MRI 等昂贵设备。

\textbf{Baseline 做法与问题。}推断不透明物体的内部结构（如平衡鹰玩具的内部空腔、古代欹器的非对称内壁）是工程和制造领域的经典难题。传统方法需要 CT 扫描或 MRI 成像，设备昂贵且操作复杂。近年来的 3D 重建方法（NeRF、3DGS）仅恢复外表面几何，不提供内部拓扑。若使用网格方法（如 PGSR、Gaussian Surfels）从点云提取表面网格再进行拓扑优化，则面临网格质量依赖——低质量网格需要繁琐的修补和填充，且网格提取本身就是计算瓶颈。

\textbf{核心思路：全粒子管线。}TopoGaussian 的核心洞察是：利用 3DGS 的点云输出，将建模、仿真和优化三个阶段统一在同一粒子表示上，完全跳过网格提取步骤。具体流程分为四步：(1) \emph{体积形状生成}——用 3DGS 从多视角图像重建外表面点云，再通过体素化填充内部粒子（遵循 PhysGaussian 的内部填充策略），得到密集体积点云 $\{(\mathbf{x}_i, V_i)\}_{i=1}^N$；(2) \emph{拓扑表示}——在每个粒子上附着拓扑参数，提供三种灵活的表示选项：(a) 粒子表示——为每个粒子分配一个决策变量，通过 sigmoid 映射为二元指示函数（$r(\mathbf{x})=1$ 表示实心，$r(\mathbf{x})=0$ 表示空心），最为灵活；(b) 神经隐式表面——用 DeepSDF 网络编码内部边界的符号距离函数 $s(\mathbf{x})$，将 SDF 通过 sigmoid 映射为拓扑指示，能编码几何先验；(c) 二次曲面——用仅 10 个参数的二次函数 $s(\mathbf{x}) = \mathbf{x}^\top A \mathbf{x} + \mathbf{b}^\top \mathbf{x} + c$ 参数化拓扑边界，适用于刚体场景（已证明对刚体拓扑优化充分）；(3) \emph{可微分仿真}——基于粒子系统构建仿真器，同时支持刚体（Newton-Euler 方程，RK2 积分）和软体（Saint Venant--Kirchhoff 超弹性模型，Leapfrog/隐式 Euler 积分），通过核函数估计变形梯度 $F_i$，还集成了气动执行器模型（气压力 $E_a = \sum_i p |\det F_i| V_i$）和罚函数碰撞模型；(4) \emph{梯度优化}——以仿真运动与参考视频之间的 $L_2$ 损失为目标，通过可微分仿真器反传梯度，用 L-BFGS-B 或 Adam 优化拓扑参数。

\textbf{设计动机。}全粒子管线避免了网格提取这一瓶颈步骤。实验表明，相比网格方法 PGSR，TopoGaussian 在体积形状生成阶段平均快 5.26 倍，且重建质量（表面法向量拉普拉斯光滑度）平均优 2.33 倍。在刚体平衡（悬挂鹰、不倒翁）、软体变形（气动软手）和真实物体（平衡鹰、欹器）等任务上验证了有效性，其中一个优化结果已成功 3D 打印并验证物理行为。

DynamicTree（2025.10）~\cite{2510_22213}是首个为 3DGS 重建的真实树木生成交互式 3D 运动动画的框架。其核心表示创新是引入稀疏体素谱（sparse voxel spectrum）来紧凑编码树的运动——直觉上，就像将树的振动分解为频率模态，类似音乐中的泛音：低频模态对应树干的整体摆动，高频模态对应细枝的颤动。通过模态分析（modal analysis），树的复杂空间运动被投影到一组正交频率基上，每个基函数描述一种典型的振动模式。动画生成时，无需逐步运行物理仿真，而是以前馈方式直接合成各模态的叠加结果；当施加外力（如风力）时，稀疏体素谱还可作为快速模态分析的基底，实现实时交互式响应——树枝随风摆动的方向和幅度由外力频率与各模态固有频率的共振关系决定。该管线先将 3DGS 重建的树转化为网格，利用稀疏体素谱生成网格运动，再将高斯基元绑定到变形网格上完成渲染。为支撑模型训练，DynamicTree 构建了首个大规模 4D 树数据集 4DTree（8786 棵树，每棵 100 帧序列），相比逐步物理仿真加速 100 倍以上。

\paragraph{VR/交互与系统集成}

LIVE-GS（2024.12）~\cite{2412_09176}修改了 VR 物理资产创作中的参数配置环节——用\textbf{LLM 自动化参数预测}替代\textbf{手动逐参数调优}，首次实现从静态 3DGS 资产到可交互 VR 物理场景的全自动转换。

\textbf{Baseline 做法与问题。}VR-GS 等现有 3DGS 物理仿真系统需要用户手动设定每个物体的物理参数（质量、刚度、阻尼等）。作者访谈 7 位物理仿真专家后发现：手动调参既依赖材料科学专业知识又极其耗时，且随场景中交互物体数量增加呈非线性增长——物体间的耦合效应迫使创作者反复修改已完成的参数配置。基于视频扩散模型的自动估计方法（如 PhysDreamer、DreamPhysics）虽可避免手动调参，但计算代价高昂、收敛慢，不适合 VR 资产的即时创作需求。

\textbf{核心思路：LLM 分析 + 实时仿真。}LIVE-GS 的核心洞察是：VR 交互不需要物理精度，只需``看起来对''——而 LLM（GPT-4o）具备足够的视觉-物理常识来推断``看起来对''的参数。系统分两阶段工作：(1) \emph{LLM 驱动的物理分析}——对输入图像，GPT-4o 通过精心设计的提示链（prompt engineering）和函数工具集成（function tool），依次推理物体类别、材料属性和质量，输出结构化 JSON 文件直接对接物理引擎，整个分析过程约 10 秒；(2) \emph{实时物理仿真}——选择 XPBD（Extended Position-Based Dynamics，一种基于位置约束的实时物理求解器，相比 MPM 计算量小两个数量级）作为仿真引擎。对于弹性体，通过 alpha shape 算法从高斯基元重建中间网格表示，在其上生成不可见物理粒子（体素采样密度 $d=15$，内部每个体素中心放一个粒子防塌缩，边界在网格顶点放置粒子保几何），然后为每个高斯基元找到最近的 $m=4$ 个物理粒子，通过距离反比加权的球面线性插值（SLERP）将物理粒子的运动传递到高斯：

\begin{equation}
\mathbf{p}^k = \mathbf{p}_0^k + \underbrace{\sum_{i=1}^{m} w_i \Delta \mathbf{p}_i}_{\text{平移}} + \underbrace{\sum_{i=1}^{m} w_i (q_i \, \mathbf{r}_i \, q_i^{-1} - \mathbf{r}_i)}_{\text{旋转引起的平移}}, \quad q^k = \text{SLERP}_{w_i}(q_i) \cdot q_0^k,
\end{equation}

其中 $w_i$ 为初始化时计算的逆距离权重，$\Delta \mathbf{p}_i$ 和 $q_i$ 为第 $i$ 个物理粒子相对初始状态的平移偏差和旋转偏差，$\mathbf{r}_i$ 为高斯基元到物理粒子的初始距离向量。对于刚体，仿真直接在网格上操作，每个高斯基元绑定到最近网格的局部坐标系；当外力超过阈值时通过 Voronoi 分割实现断裂效果。

\textbf{设计动机与验证。}这一设计平衡了交互质量与预测效率：LLM 提供``近似但可信''的参数（在 ABO 数据集 11 个物体上质量预测的 MAPE 为 15.0\%），XPBD 提供实时仿真反馈（渲染帧率 30 fps）。16 人用户研究（VBench 评估）表明，LLM 自动配置的弹性参数在视觉质量、运动平滑度、美学质量和视频-文本一致性四个维度上均达到满意水平，且创作时间从手动调参的数十分钟缩短至约 10 秒。

在综述方面，"Grounding Creativity in Physics"（2025.02）~\cite{2502_07007}系统综述了物理先验在 AI 生成内容中的应用，按表示类型（视觉、NeRF、3DGS）分类方法，揭示了生成模型与物理真实感之间的差距。

\begin{takeaway}
物理仿真将高斯基元从纯视觉表示推进到具有物理属性的动态实体。核心技术路线是在高斯基元与物理粒子/网格之间建立绑定关系，通过 XPBD、MPM 等物理引擎驱动运动，再映射回高斯参数。LIVE-GS 用 LLM 自动预测物理参数将创作时间从数十分钟缩短至 10 秒，DynamicTree 的稀疏体素谱实现树运动的实时交互式动画。
\end{takeaway}

\section{近期研究（2026 年 8--11 月）}
\paragraph{Struct-GStream：用结构化高斯跟随基本运动，再补新出现区域}
2026.08~\cite{2608_01053}。Struct-GStream 将自由视点视频拆成可随时间更新的结构化高斯和用于补缺的自由高斯。前者沿近似局部刚体更新以复用已有结构，后者补足初始表示不足和之后出现的区域。原文比较跨帧整体补片与每帧独立生成，以区分新物体重建和时间连贯性。

论文在 N3DV、Meet Room 和 Google Immersive 给出对照，并报告所测百万像素场景平均每帧约 3.0 至 5.4 MB、在线训练每帧少于 10 秒、渲染超过 100 FPS。后两项分别是场景更新与显示，不是同一个端到端流延迟。作者指出首帧在稀疏角落和不可见室外区域重建不足会引起抖动，仍未达到实用流系统的全部要求。

\paragraph{D²-4DGS：把两类深度分别用于稀疏动态高斯的监督与补充}
2026.08~\cite{2608_01588}。少量同步相机只靠 RGB 拟合时，动态原语容易被放到只能解释训练视图的位置。D²-4DGS 对齐多视图与单目深度，先用验证过的深度做清理和约束，再将可靠多视图深度及部分单目估计用于提出缺失区域的增密候选。原有时序参数化保留，改动放在几何证据与 RGB-D 联合优化。

原文在 N3DV 六个、Technicolor 五个和 ENeRF-Outdoor 三个场景，分别让 2、3 或 4 路标定相机参与训练，其余有效相机留出。几何证据的益处在两路视图时较明显。结果针对同步标定相机和验证通过的深度，并不能把对齐的单目提议等同于实测深度。

\paragraph{FAST-GS：用傅里叶时间基表达动态高斯运动}
2026.08~\cite{2608_01958}。FAST-GS 以傅里叶分解替换时空高斯中的多项式运动参数，通过各轴的正弦和余弦系数表达时间变化，并加上运动正则以约束振荡和噪声。这是对动态高斯中心运动的表示调整，仍从多视图视频优化用于渲染的场景。

论文在 N3V 留出一台相机测视图泛化，在 Google Immersive 则用前 80\% 帧训练、后 20\% 帧测试。空间留出与时间留出的结果对应不同难度，需按协议分别解读。有限的频率基与有限观测时间也不能保证未见长时间行为的重建。

\paragraph{ASTRA：用二维轨迹帮助估计多相机时间偏移}
2026.08~\cite{2608_02006}。不同相机的帧率或采集起点不一致时，仅用颜色误差可把时间错位吸收到场景形状和运动中。ASTRA 在动态高斯上联合优化每台相机的时间偏移，用三维运动投影与观测轨迹对齐，并用动态和置信掩码减少不可靠匹配的影响。

作者将框架接入 4DGS 与 EDGS，在 Plenoptic 和 Blender 等数据评估 5 至 25 帧的偏移，并另用低纹理数据及其他实录验证。原文指出简单可靠性掩码在严重遮挡下可能失效，真实硬件的滚动快门、时钟漂移及曝光变化也超出当前常量偏移设置。

\paragraph{ACA-GS：让动态锚点按局部需求分配原语和特征}
2026.08~\cite{2608_04581}。锚点网络若为每个锚点生成相同数量的高斯，并分配相同特征容量，静态简单区域与变形复杂区域会消耗类似预算。ACA-GS 在 GIFStream 基础上允许锚点输出数量变化，把冗余高斯转给需要细节的区域，并通过学习到的通道掩码减少无用的锚点特征。

原文在 MPEG、Panoptic Sports 与 N3DV 比较渲染和表示大小，报告所测场景约 3.5 至 7.3 MB 的表示。三套协议的 LPIPS 骨干并不相同，不能跨表直接混算。容量选择和掩码强度也会改变压缩与画质的权衡。

\paragraph{RORA：把视频中的物体转成带关节的可模拟资产}
2026.08~\cite{2608_04842}。RORA 先从单段视频建立高斯和网格混合表示，再做几何分解。使用者参与把分解结果归组为可动部件、确认候选关节，高斯负责外观，网格提供可绑定部件及模拟几何。系统最后导出可交互资产，因此关节恢复是半自动流程。

作者在 PartNet-Mobility 对照 Articulate-Anything 和 ScrewSplat，并展示真实拍摄物体的模拟操作。对比中更新了 Articulate-Anything 的视觉语言模型，这一点影响复现实验定义。模拟演示也不能替代对不可见关节、柔性形变或网格误差下物理精度的验证。

\paragraph{ERF-GS：从不同相机位置的事件与 RGB 恢复快速运动}
2026.08~\cite{2608_08531}。多相机的事件和 RGB 输入不必从同一视点获得。ERF-GS 在动态 RGB 高斯骨干上引入两条事件线索，分别约束渲染及原语增密，使模糊或低帧率 RGB 缺少的运动信息可由另一位置的事件输入补充。这里的事件损失不要求同视点的清晰 RGB 真值。

论文在 Neu3D-dv 和 Nvidia-dv 比较，除整图指标还用 SAM2 标出运动区域计算动态 PSNR。作者说明主要多视图事件和模糊数据来自模拟，源 RGB 帧率限制合成事件的时间分辨率，真实高速运动和多传感器采集的差异仍需验证。

\paragraph{用黎曼流匹配生成和细化动态高斯参数}
2026.08~\cite{2608_18388}。题为 Depth Anything V4 的论文将单目视频特征送入黎曼流匹配模型，在 4D 高斯参数空间以常微分方程生成从粗到细的表示，再用渲染损失作测试时优化。该流程包含预训练先验与测试适配，部署成本不能等同于只作一次深度前馈。

原文用 Dynamic-Sim 和 KITTI Dynamic 分析动态重建，另以 ETH3D、ScanNet 做新视图测量。论文控制实验在相同预训练和测试细化条件下，相对确定性 MLP 报告 F 分数提高 0.044。预训练能否摊销取决于实际部署规模和测试优化开销，不能将实验设定直接外推。

\paragraph{按时间变化匹配四维高斯的空间滤波尺度}
2026.08~\cite{2608_21828}。动态高斯的运动会改变其跨相机的投影大小，使用固定抗混叠尺度可能在放大时显露空洞，在缩小时出现锯齿。该方法为原语估计时间与焦距深度比的联合分布，在给定时间据条件分布选取采样尺度，再对投影高斯施加平滑。

论文以 SARO-GS 为主要底座，在 Plenoptic Video 上测试低分辨率训练后的一到四倍输出，另用 D-NeRF 做缩小实验。原文指出为每个原语保存时间与焦距深度比的离散密度矩阵，内存随原语数和两个维度的分箱增长，大场景需要另计这项成本。

\paragraph{NemoSplat：在水下无标定视频中联合预测几何、瞬态与介质}
2026.08~\cite{2608_22888}。散射和游动物体会同时污染水下相机跟踪与静态地图。NemoSplat 从无标定视频经前馈网络联合估计相机、动态分离、可渲染高斯和介质参数，并以衰减和后向散射解释部分颜色变化，避免直接把这些变化写进静态几何。

作者用六条合成序列评价跟踪和渲染、另用 14 条真实序列检验新视图。合成平均 ATE 为 1.88 米，优于列出的两种前馈几何对照 VGGT 和 StreamVGGT；原文也指出极低可见度的 Deepsea 上光流 SLAM 仍能更稳定跟踪。真实渲染测试不能充当真实位姿误差，语义与介质模块另增加训练显存。

\paragraph{EvoGS：把逐帧形变改成带观测校正的时间递推}
2026.09~\cite{2609_00994}。EvoGS 为每个高斯保存持续形变状态，由历史校正值外推出当前运动，再用轻量网络从当前条件生成补充观测。校正权重取决于历史残差和运动变化，原语克隆与分裂也沿校正后的形变进行，以减少不确定原语继续增密。

论文将模块装到多种形变骨干，并在 NeRF-DS、HyperNeRF 的 vrig 子集和 Neu3D 比较。Neu3D 表中网格骨干 Grid4D 的平均 PSNR 为 30.75 dB，EvoGS 对应配置为 32.44 dB。短时外推假设采样时间间隔近似均匀，并要保留额外状态；这些条件限制了对非均匀长时序的直接推断。

\paragraph{UniFusion：把各相机各时刻的单目深度对齐后再训练动态高斯}
2026.09~\cite{2609_05888}。多个相机或时间的单目预测具有不同尺度与偏移，单帧深度不能直接当成统一的几何监督。UniFusion 用同步相机的稀疏 SfM 参照联合拟合时空深度场，通过低秩时间权重和沿射线预测的深度偏移维持跨帧共享，并将对齐结果用于动态高斯的初始化和优化监督。

论文在 Ego-Exo4D 按时间留出帧，在 EgoHuman 按相机留出视图，两个数据均使用已知相机位姿。深度对齐也有独立于高斯后端的重投影一致性实验。原文展示快速复杂运动仍可受变形后端限制，改善深度先验不等于取消运动模型的表达瓶颈。

\paragraph{FastFlowGS：用多来源对应与时序先验初始化每帧高斯}
2026.09~\cite{2609_16310}。流式动态重建若直接从上一帧复制位置，快速运动会给当前帧优化留下较大误差。FastFlowGS 从不同粒度的匹配、轨迹和光流提取带置信度的观测，将它们通过相机几何结合到三维位置，并以 Kalman 式时间先验给当前帧提供更好的初值，再继续图像优化。

原文在 CMU-Panoptic 六段数据分别用原帧率和五倍抽样测试，并在 Monaco4D 检验大位移室外场景。位置预测不更新旋转和尺度，所以快速旋转或形变仍要在本帧重新拟合；低纹理及沿深度方向运动也会削弱对应。不同指标和序列的改善不应汇成一个通用加速比。

\paragraph{以风驱振子检验动态高斯的物理外推}
2026.09~\cite{2609_17810}。作者将树的高斯按空间位置分为刚性部件，为每个部件学习受实测风速驱动的阻尼振子，再通过时间积分得到绕枢轴的转动。实验分别改变相机、时间与风力，区分同窗拟合、新视角表现和物理外推能力。

在三棵合成树上，学习形变的基线在未见相机上优于刚性振子；中等及密集树上，当时间或风力超出训练分布时，排序反过来。结果支持物理先验可帮助外推，同时作者未可靠恢复阻尼，模型没有分支耦合和部件内弯曲，并要求测试时同步提供风速。

\paragraph{SplashSplat：在同步多视角下分开重建飞溅液体与背景}
2026.09~\cite{2609_20818}。SplashSplat 用带标定和液体掩码的同步图像分别拟合动态液体高斯与静态背景，物理与几何结构只在有观测支撑的位置提供约束。最终渲染将两部分组合得到视频。

评估用同一七相机装置，五台训练、两台保留，测试视点位于训练视点之间，且按各相机固定裁剪比较。结果说明该协议下真实飞溅的新视角效果，不能直接外推自由相机与未观测物理轨迹。

\paragraph{4DGS-Fixer：用密集点初始化和视频先验细化四视图动态重建}
2026.09~\cite{2609_21176}。4DGS-Fixer 用多视图深度估计得到初始的稠密点云，再让视频修复先验迭代细化欠观测区域。视频结果为后续动态高斯优化提供额外约束，真实相机观测仍是检验生成结果的依据。

Neural3DV 测试沿用 4C4D 的协议，选取四个空间分散的训练相机，其余作测试。作者指出生成质量可能影响细化，并且迭代仍然耗费算力；结果不能解释为实时四图推断。

\paragraph{Dynamic Thermal Gaussians：用共享几何组织动态 RGB 和热红外}
2026.09~\cite{2609_24531}。该方法将 RGB 与热红外置于同一规范高斯几何，分别用模态属性、时变形变和结构感知路由渲染。模态专有的辐射变化可以由专门原语和属性解释，同时两类观测共享空间与相机参考。

论文在 DynamicRGBT-Scenes 的十二个真实序列中报告热红外 PSNR 平均 32.24 dB。输入依赖 RGB 视频获得的 COLMAP 位姿及与热帧的人工同步，评测的是配对动态新视角，不是以未标定的纯热红外完成重建。

\paragraph{ChronoFuseGS：融合不同时相并记录每个高斯的持续性}
2026.09~\cite{2609_31339}。ChronoFuseGS 将各时相分别训练的高斯场先配准到同一坐标系，再给每个高斯配置按时相变化的持续性属性。联合精调从各时相共享观测，并约束持续性数值的相邻时相变化，以免提高渲染分数的同时抹掉变化线索。渲染可以用颜色提示某个局部结构在哪些时相出现。

论文比较了独立训练和合并精调结果，整体图像指标的增幅较小，部分局部结构的改善更明显。参与者研究在二维预渲染图中检查变化标示。大范围变化可能使 SfM 对齐失败，单次采集中发生的结构或照明变化也不在当前建模范围内。

\paragraph{OneFixer：用同一次自回归生成轨迹训练街景视频修复}
2026.09~\cite{2609_32175}。OneFixer 将部署时的自回归历史纳入训练，在同一条生成轨迹上同时做流匹配与逐帧感知监督。它用单步生成修复输入的街景高斯渲染，训练目标兼顾帧质量和连续外观，而输出属于视频修复结果。

论文分别用 Waymo 和首尔驾驶数据训练，报告视频和逐帧指标，并说明细结构仍可能出现问题。不同数据域的模型训练分开进行，单步时延和图像一致性也不能直接解释为三维几何重建精度。

\paragraph{Endo-TSR：为内镜高斯的颜色和位移分别补充时间频谱}
2026.09~\cite{2609_32399}。Endo-TSR 在基础形变模型上为颜色和位移分别学习受限的傅里叶残差。颜色支路处理时间光度变化，位移支路由频谱先验约束，减少外观误差被挤进轨迹。多尺度拉普拉斯监督只作用在有效组织区域，以恢复细小纹理。

六段 EndoNeRF 和 StereoMIS 序列沿用既有七比一训练测试划分。定量结果评价渲染质量，消融中颜色建模带来的 PSNR 收益最大。这并不证明内部组织运动已准确恢复，且方法输入依赖组织掩码、深度先验和相机参数。

\paragraph{ProDyGS：以单目深度生成静止相机没有采到的代理视角}
2026.09~\cite{2609_32711}。静止或近乎静止的相机没有跨帧视差，难以用不同帧三角化动态表面。ProDyGS 先用单图模型预测每帧深度，把颜色随深度前向投影到相邻虚拟视角，再与真实相机观测一起约束规范高斯和时间形变网络。

在 DyNeRF，系统用单一相机序列训练，并在未见的真实相机视角评价。预测深度与 COLMAP 坐标的尺度平移对齐是该实验的评价设置。代理图像重投影的是同一帧已有观测，不能提供真实遮挡面证据。

\paragraph{DispFlow-GS：去掉相机运动后再监督形变位移}
2026.09~\cite{2609_36940}。DispFlow-GS 指出把渲染高斯流与二维光流直接拟合会混淆观测域。它用中间帧渲染分解相机与场景运动，在固定视点下用位移幅度约束形变。这样可以把图像质量指标与运动恢复指标分开检验。

论文在 NeRF-DS 观察到所用运动指标改善，但 PSNR 基本不变。HyperNeRF 存在因相机位姿误差而失败的场景。位移量监督也不等于知道逐像素对应，因此复杂运动中的精确轨迹仍难恢复。

\paragraph{EndoPrior-GS：用组织纹理引导稠密高斯生长与时间一致性}
2026.09~\cite{2609_37874}。EndoPrior-GS 以无需另训网络的联合纹理线索初始化内镜高斯，限制失真区域的冗余增密，并按邻近原语的局部纹理关系调整时序约束。在缺少可靠纹理的位置，方法无法获得同等强度的组织结构证据。

作者报告 EndoNeRF 和 SCARED 上的 Flow Error 下降。该指标来自 RAFT 的图像平面光流，并辅以渲染深度诊断，因为数据没有稠密三维组织运动真值。模糊、过曝、烟雾和失准掩码都可能削弱纹理先验。

\paragraph{用固定欧拉速度场和周期高斯输运重建水景}
2026.09~\cite{2609_38622}。该方法从具有相机视差的单目水景视频学习随位置变化、随时间固定的欧拉速度场。动态水体高斯沿固定速度场输运以产生可循环运动，余留的非周期外观变化另由时间残差描述。由掩码识别的静态高斯不随水体运动。

论文使用七段自采视频，将输入轨迹中的连续片段留作定量测试，并安排主观外观与运动比较。缺少同步多视角真值，因此留出片段的成绩不等于对独立角度的三维流速验证。

\paragraph{按表示机制理解动态重建的四维综述}
2026.09~\cite{2609_39960}。这篇综述以表示机制而非应用场景组织路线，分别梳理四维 NeRF 的隐式场、特征体和时间先验，以及高斯的显式时空原语、规范形变与逐帧训练。这个框架适合区分运动是直接写入基元，还是经共享规范坐标恢复。

性能章节汇总了动态场景数据、图像质量、几何和效率指标以及论文中的已报告结果。它能说明不同路线的取舍，但这些表格并非统一硬件、输入相机和训练协议下的重新实验。

\paragraph{以空间、时间和生成先验补全六相机动态场景}
2026.10~\cite{2610_04606}。该工作用前景掩码与背景深度约束不同区域，借助时间插值与光流约束投影运动，再在宽视角间隙内渲染虚拟视图。生成模型参考相邻已采图像修复虚拟渲染，结果作为伪观测循环回到四维高斯训练。

论文报告 SIGGRAPH Asia 2026 稀疏视图竞赛结果，并在验证集分别列出前景和整幅图像的消融分数。附加恢复可以弥补六相机覆盖的缺口，但虚拟视图不是实测，作者还将重建耗时与模型大小留作改进。

\paragraph{OuroWorld：把生成运动蒸馏为严格周期的四维高斯动画}
2026.10~\cite{2610_12461}。OuroWorld 从给定静态高斯场生成循环参考视频，再把跨视角生成的运动信号蒸馏到以傅里叶级数表示的周期形变。单独的漂移项吸收生成视图中不一致的部分，最终按周期形变驱动可渲染的三维基元。

作者在固定相机和绕行轨迹上分析运动活跃度、视觉质量、自然度和接缝，并辅以用户比较。这里没有真实动态真值。生成先验的错误可能进入最终场景，单参考视频在跨度大的新视角也更容易留下不一致。

\paragraph{Make It Flow：从静态高斯重建生成可循环流体动画}
2026.10~\cite{2610_13860}。Make It Flow 从输入参考图或静态高斯渲染图获得二维流动线索，再借高斯深度将不同相机的候选运动提升并融合成固定三维欧拉场。生成的局部形变作用到流体高斯上，通过对称时间处理衔接循环，不需要先采集动态多视图。

作者用七个随手采集的流动场景检查动画观感和成本，缺乏真实流场真值。反光水面或重建深度不一致会使提升的运动出错，静态模型里没表示的流体细节也不会因此成为实测几何。

\paragraph{Demiurge：自回归重渲染时持续合并四维点云记忆}
2026.10~\cite{2610_14685}。Demiurge 把已有单目视频和相机轨迹作为输入，从记忆中投影过去的几何及外观，为新的相机轨迹生成视频，再用补全深度把新观测写回四维点云。高斯重建可以使用这些生成多视图作为下游输入，核心记忆本身是点云。

论文分开评价生成视频的感知质量、跨轨迹同步和几何可重建性。依次执行外观生成、深度恢复与记忆更新，作者报告每条目标轨迹约需五分钟，当前不能视为交互帧率重渲染。

\paragraph{PUGS：将水下高斯位移建模为随场景变化的分布}
2026.10~\cite{2610_17936}。水下动态重建容易将折射及暂态视觉噪声误当作物体形变。PUGS 为每个高斯预测平均位移及方差，训练一条随机探索的渲染支路和一条确定性支路，并以自一致约束协调二者。学习的方差还用于优先补充较难可靠拟合的区域。

作者评价 UVEB、DRUVA 与 PhysGaia 共 22 个场景。强浑浊时图像误差和感知指标可能发生取舍，且不确定性退火仍需手工安排。这里的方差是重建中的可学习分量，不是经外部观测校准的运动误差概率。

\paragraph{SwitchSplat：将物体开关状态从连续时间中分离}
2026.10~\cite{2610_24812}。对可以打开、关闭或切换的物体，时间只是观测顺序，不能单独作为场景状态的索引。SwitchSplat 用共享静态底层储存不变部分，并为每个可交互物体维护可独立切换的二态高斯槽；可微槽表把各帧观测连到相应状态。这使新视图可以组合不同物体的状态，而不必重放某个采集时刻。

实验的常规十个场景与额外两组详细测试区分留出相机和新状态组合。方法当前假定状态为二元，采集约束在单人操作，槽位数量和大致交互窗口还由视觉语言模型给出初值。论文验证的是状态可控渲染，并未建立物理交互模型。

\paragraph{GTSM：在高斯空间支撑上附加连续热状态}
2026.10~\cite{2610_26918}。GTSM 将高斯支撑与其温度演变分开：三维或四维高斯后端描述原语位置和投影，附着在原语上的连续热状态由场景级潜在驱动和原语局部动态共同决定。对标定红外观测进行可微渲染，可以在已观测时间段查询新视角，也能让学习到的状态继续滚动。

实验分别使用真实 RHD 红外序列以及几何、温度共同变化的 D-NeRF 衍生数据，并区分段内重建与时外预测。几何运动会改变原语的空间覆盖，但几何变量不直接进入热演化方程，预测结果不等同于求解任意场景的热传导。

\paragraph{AEON：将视频生成潜空间接到显式四维高斯解码}
2026.10~\cite{2610_28631}。AEON 用一个四维重建模型从单目视频解码深度、高斯外观、动态和相机轨迹。为了让预训练视频扩散模型直接产生可查询的四维表示，另训练域适配器，把生成潜变量转换为重建潜变量。得到的显式高斯可沿指定视角和时间渲染。

论文分别检查动态重建、新轨迹渲染和生成视频，不同协议的重建与感知指标应分别解读。扩散模型序列长度限制当前生成时长，大幅转动相机可能暴露原视频从未生成的区域；跨窗口持续融合属于作者提出的扩展方向。

\paragraph{DUET4D：用当前高斯约束生成，再筛选伪观测回填}
2026.10~\cite{2610_35002}。DUET4D 先从当前四维高斯抽取跨时空对应，让冻结的单视图扩散模型恢复伪观测时沿这些锚点保持一致。生成结果再反向辅助建图，但只有置信度足够的差异才用于增密，与真实观测持续冲突的高斯则会被裁掉。

作者在三个稀疏视角动态数据上比较渲染效果。初始几何严重失真、运动太快或遮挡时间太长时，对应锚点的支持会下降。系统此时退回较独立的单视图恢复，可以限制错误传播，但无法保证重建真实不可见内容。

\paragraph{GaussianSpline：让每个动态高斯学习自己的时间分段}
2026.10~\cite{2610_35224}。GaussianSpline 为每个动态高斯分配分段三次贝塞尔轨迹，联合优化控制点与非均匀时间区间，因此固定段数可以向快运动处集中。接合点使用基于切向反射的几何连续性约束，在保持方向连续的同时允许相邻段速度变化。

论文在公开 N3DV、一个 Google Immersive 场景及两个合成序列评价。短时碰撞可能超出固定小段数的表达能力，而且各高斯轨迹独立学习，没有显式的跨原语刚体或关节约束。拓扑变化靠透明度和增密间接表示。

\paragraph{Topo-4DGS：在拓扑事件前后改变运动分支和表面支撑}
2026.11~\cite{2611_16012}。离散的拆开、合并、打开或闭合会改变物体部件运动及哪些表面可见。Topo-4DGS 按空间位置和事件进度对规范高斯分配分段运动，并在事件涉及的部位激活或停用相应高斯。新暴露的表面可用新原语补足，并绑定到所属部件的运动。

论文在 24 段第一人称视频上使用评审过的事件标注评价，不等同于全自动发现拓扑事件。当前适用对象是事件较局部且可近似分段刚性的部件。强形变、长时间手部遮挡和复合缓慢变化仍是难点。

\paragraph{ReCan-GS：用形变敏感度分离动态场景中的短时干扰}
2026.11~\cite{2611_16259}。ReCan-GS 同时维护可随时间形变的核心高斯和逐视图的干扰高斯。局部形变场的雅可比敏感度指出哪些运动难以由共享形变稳定解释，然后让这部分观测更倾向于进入干扰层，或者用于后处理裁除。

论文测试四个动态场景和四个合成烟雾场景。该分数并非经过校准的概率或语义标签，稳定出现的干扰可能保留，真实但复杂或观测偏少的运动也可能被削弱。作者还指出静态基准效果有限且缺乏专门烟雾移除基线。

\paragraph{OCLDA-GS：让流式高斯锚点沿物体边界传播运动}
2026.11~\cite{2611_23416}。OCLDA-GS 用固定锚点预算从已精调关键帧向后续视频帧传播运动。光流较强的物体组分得更多锚点；锚点插值限于组内，以减少物体边界穿插。每组刚性运动还作为局部非刚性偏移的条件，不要求每帧再拟合全部高斯。

论文在 N3DV、Google Immersive 等测试，域内对比仍有在线基线在 PSNR 上更强。组划分错误和缺少清楚主刚性变换的关节物体会影响传播，当前学习校正只更新位置，不纠正尺度、方向等其他属性。


\section*{本章小结}
\addcontentsline{toc}{section}{本章小结}

\begin{futurebox}{本章总结与展望}{ch06_summary}
本章系统梳理了 3DGS 在动态场景和 4D 建模领域的 250 篇论文，是该方向论文最密集的章节之一。

\textbf{变形场与 4D 高斯表示}（第~\ref{sec:06_1}~节）是动态 3DGS 的技术核心。从规范空间+变形场的基础解耦架构，到原生 4D 高斯（4DGS）直接在时空域建模，再到高阶编码和物理约束变形，时序建模的方法论从简单的线性插值演进到复杂的动力学系统。前馈 4D 重建（DynamicVGGT、DrivingScene）和大规模 4D 表示正在推动动态 3DGS 走向实时应用。

\textbf{单目动态场景重建}（第~\ref{sec:06_3}~节）面对最具挑战性的输入条件——单一 RGB 相机、无深度信息、无相机位姿。变形场正则化、深度先验引导和运动模糊处理是缓解严重不适定性的三条互补路线。事件相机和多模态融合为高速运动场景提供了独特的时间分辨率优势。

\textbf{运动建模与跟踪}（第~\ref{sec:06_4}~节）从像素级运动场演进到物体级运动轨迹，语义引导和层次化运动表示使动态场景的分析从底层光流提升到高层运动理解。

\textbf{视频重建与时序一致性}（第~\ref{sec:06_5}~节）面向自由视角视频的实际应用需求。超高速渲染、4D 压缩编码和流式传输构成了从实验室到产品化的技术栈。时序一致性和体积视频是当前的前沿方向。

\textbf{动态物体与场景分解}（第~\ref{sec:06_6}~节）解决了"从混乱中提取秩序"的根本问题——将混合的动态场景分解为静态背景和独立运动的实例。自监督和无标注的分解方法（$\text{S}^3$Gaussian、DeSiRe-GS）正在取代依赖昂贵 3D 标注的方案。

\textbf{4D 语义、HDR 动态和特殊模态}（第~\ref{sec:06_7}~节）展示了动态 3DGS 的多元化拓展。4D 语义场将语义理解从静态扩展到动态，HDR 动态建模处理了时变光照，SLAM 集成将动态重建推向在线应用。

当前该领域最突出的开放问题包括：
\begin{itemize}
\item \textbf{长时序 4D 一致性}：当前方法在短序列（数秒到数十秒）上效果良好，但长时间跨度的累积漂移、外观变化和运动建模误差使重建质量显著退化。如何实现数分钟甚至数小时的持续 4D 重建仍是核心挑战。
\item \textbf{复杂运动的统一建模}：刚体运动、铰接运动、弹性变形、流体流动等不同类型的运动目前需要专门的表示。统一的运动建模框架能否覆盖所有运动类型，同时保持计算效率？
\item \textbf{4D 数据的存储与传输}：4D 高斯的数据量随时间线性增长，如何在不牺牲质量的前提下实现高效的 4D 压缩和流式传输，是从实验室走向实际部署的关键瓶颈（参见第~\ref{ch:11}~章压缩方法的讨论）。
\item \textbf{动态场景的语义理解}：将第~\ref{ch:10}~章的场景理解能力扩展到 4D 场景——在线、实时、开放词汇的动态语义理解——是一个极具价值但尚处于早期阶段的方向。
\item \textbf{物理仿真的集成}：将第~\ref{sec:06_8}~节的物理仿真能力与动态重建方法更紧密地集成，实现从观察到物理仿真就绪资产的全自动管线，仍是长期目标。
\end{itemize}
\end{futurebox}
